% Glossary entries (used in text with e.g. \acrfull{fpsLabel} or \acrshort{fpsLabel})
%\newacronym[longplural={Frames per Second}]{fpsLabel}{FPS}{Frame per Second}
%\newacronym[longplural={Tables of Contents}]{tocLabel}{TOC}{Table of Contents}
\newglossaryentry{abstraccion}{
	name={Abstracci{\'o}n},
	text={abstracci{\'o}n},
	sort={abstraccion},
	plural={abstracciones},
	first={abstracci{\'o}n},
	firstplural={abstracciones},
	description={Proceso cognitivo y metodol{\'o}gico mediante el cual se identifican y modelan las propiedades esenciales de un sistema, omitiendo detalles irrelevantes para un determinado nivel de an{\'a}lisis. En ingenier{\'i}a de software y sistemas IoT, la abstracci{\'o}n permite separar niveles conceptuales, facilitar el dise{\~n}o modular y reducir la complejidad durante la especificaci{\'o}n y construcci{\'o}n de soluciones.},
	type=main,
	symbol={🔎},
	user1={Reducci{\'o}n de complejidad},
	user2={Separaci{\'o}n por niveles},
	user3={Modelo conceptual},
	see={abstraccionsuperior,modelo,especificaciontecnica}
}

\newglossaryentry{abstraccionsuperior}{
	name={Abstracci{\'o}n superior},
	text={abstracci{\'o}n superior},
	sort={abstraccion superior},
	plural={abstracciones superiores},
	first={abstracci{\'o}n superior},
	firstplural={abstracciones superiores},
	description={Nivel conceptual elevado que permite definir decisiones arquitect{\'o}nicas y propiedades globales antes de comprometerse con detalles tecnol{\'o}gicos espec{\'i}ficos.},
	type=main,
	symbol={🔝},
	user1={Nivel conceptual alto},
	user2={Independencia tecnol{\'o}gica},
	user3={Visi{\'o}n arquitect{\'o}nica},
	see={abstraccion,decisionarquitectonica}
}

\newglossaryentry{accesibilidad}{
	name={Accesibilidad},
	text={accesibilidad},
	sort={accesibilidad},
	plural={accesibilidades},
	first={accesibilidad},
	firstplural={accesibilidades},
	description={La \emph{accesibilidad} se refiere al grado en que un sistema, servicio, dispositivo o entorno digital puede ser utilizado de manera eficaz por cualquier persona, incluidas aquellas con discapacidades temporales, permanentes o situacionales. En el contexto de \gls{iot} y \glspl{sibiot}, la accesibilidad implica dise{\~n}ar interfaces, aplicaciones y \glspl{dispositivointeligente} que permitan una interacci{\'o}n inclusiva mediante adaptaciones como control por voz, retroalimentaci{\'o}n h{\'a}ptica, se{\~n}alizaci{\'o}n visual, configuraciones personalizables y cumplimiento de est{\'a}ndares como WCAG. Garantiza equidad de uso, mejora la \gls{ux} y ampl{\'i}a el alcance de las soluciones tecnol{\'o}gicas.},
	type=main,
	symbol={♿},
	user1={Dise{\~n}o universal e inclusivo},
	user2={Adaptaci{\'o}n multicanal en SIbIoT},
	user3={Cumplimiento de est{\'a}ndares WCAG},
	see={interfaz,interaccion,ux,usuario}
}

\newglossaryentry{acid}{
	name={ACID},
	text={ACID},
	sort={acid},
	plural={propiedades ACID},
	first={atomicidad, consistencia, aislamiento y durabilidad (\textit{Atomicity, Consistency, Isolation, Durability}, ACID)},
	firstplural={propiedades ACID},
	description={El t{\'e}rmino \emph{ACID} describe un conjunto de propiedades fundamentales que garantizan la fiabilidad en la ejecuci{\'o}n de \glspl{transaccion} dentro de una \gls{basedatosrelacional}. Estas propiedades aseguran que las operaciones se completen de manera confiable incluso frente a fallos o concurrencia: (1) \textbf{Atomicidad}: todas las operaciones de una transacci{\'o}n se ejecutan en su totalidad o no se realizan; (2) \textbf{Consistencia}: la base de datos pasa de un estado v{\'a}lido a otro tras la transacci{\'o}n; (3) \textbf{Aislamiento}: las transacciones concurrentes no interfieren entre s{\'i}; y (4) \textbf{Durabilidad}: los cambios persisten incluso ante fallos del sistema. En ecosistemas \gls{iot} y \glspl{sibiot}, las propiedades ACID son esenciales para la integridad de datos cr{\'i}ticos en aplicaciones financieras, industriales o de salud, donde la precisi{\'o}n y confiabilidad son imprescindibles},
	type=main,
	symbol={⚖️},
	user1={Atomicidad: todo o nada},
	user2={Consistencia: validez de datos},
	user3={Durabilidad y confiabilidad de transacciones},
	see={transaccion,basedatosrelacional,oltp,amazonaurora,sibiot}
}

\newglossaryentry{acoplar}
{
	name={Acoplar},
	text={acoplar},
	sort={acoplar},
	plural={acoplar},
	first={acoplar},
	firstplural={acoplar},
	description={\emph{Acoplar} significa unir o encajar entre s{\'i} dos piezas o cuerpos de manera que ajusten perfectamente. En ingenier{\'i}a de software, tambi{\'e}n se emplea para describir la conexi{\'o}n entre componentes de un \gls{sistemainformatico}.},
	type=\glsdefaulttype,
	symbol={🔗},
	user1={Enlazar},
	user2={Unir},
	user3={Conectar},
	see={acoplado}
}

\newglossaryentry{acoplada}
{
	name={Acoplada},
	text={acoplada},
	sort={acoplada},
	plural={acopladas},
	first={acoplada},
	firstplural={acopladas},
	description={\emph{Acoplada} es la forma femenina de acoplado. Se usa como adjetivo para describir una entidad unida o conectada a otra. Ver \gls{acoplar} y \gls{acoplado}.},
	type=\glsdefaulttype,
	symbol={🔗},
	user1={Vinculada},
	user2={Unida},
	user3={Conectada},
	see={acoplar}
}

\newglossaryentry{acoplado}
{
	name={Acoplado},
	text={acoplado},
	sort={acoplado},
	plural={acoplados},
	first={acoplado},
	firstplural={acoplados},
	description={\emph{Acoplado} es el participio del verbo acoplar. En inform{\'a}tica se utiliza para indicar que dos elementos de software trabajan de manera m{\'a}s o menos independiente seg{\'u}n el grado de acoplamiento. Cuanto menor es el acoplamiento, mayor es la modularidad y la flexibilidad del \gls{sistema}.},
	type=\glsdefaulttype,
	symbol={⚙️},
	user1={Conectado},
	user2={Vinculado},
	user3={Relacionado},
	see={acoplar}
}

\newglossaryentry{acoplamiento}{
	name={Acoplamiento},
	text={acoplamiento},
	sort={acoplamiento},
	plural={acoplamientos},
	first={acoplamiento},
	firstplural={acoplamientos},
	description={Grado de dependencia entre los componentes de un sistema, que refleja cu{\'a}n estrechamente est{\'a}n vinculadas sus responsabilidades, interfaces, estructuras de datos o flujos de control. Un alto nivel de acoplamiento implica que los cambios en un componente tienden a propagarse a otros, dificultando la evoluci{\'o}n, el mantenimiento y la escalabilidad del sistema. En el contexto de los sistemas de informaci{\'o}n basados en IoT y los sistemas distribuidos, la reducci{\'o}n del acoplamiento es un principio de dise{\~n}o fundamental, logr{\'a}ndose mediante arquitecturas estratificadas, comunicaci{\'o}n orientada a eventos, uso de interfaces estandarizadas y mecanismos de mensajer{\'i}a as{\'i}ncrona, lo que favorece la interoperabilidad, la tolerancia a fallos y la evoluci{\'o}n independiente de los componentes},
	type=main,
	symbol={🔗},
	see={desacoplamiento,arquitecturaestratificada,orientadoeventos,procesamientodistribuido,interoperabilidad,sibiot}
}

\newglossaryentry{actor}{
	name={Actor},
	text={actor},
	sort={actor},
	plural={actores},
	first={actor},
	firstplural={actores},
	description={Entidad (persona, rol, organizaci{\'o}n o sistema externo) que interact{\'u}a con un \gls{sistema} para alcanzar objetivos espec{\'i}ficos. En el modelado de requisitos, un \emph{actor} representa una fuente de acciones o solicitudes que influyen en el comportamiento del sistema.},
	type=main,
	symbol={👤},
	user1={Rol o entidad externa},
	user2={Interacci{\'o}n con el sistema},
	user3={Modelado de requisitos},
	see={analisisrequisitos,sistemainformacion}
}

\newglossaryentry{actuador}
{
	name={Actuador},
	text={actuador},
	sort={actuador},
	plural={actuadores},
	first={actuador},
	firstplural={actuadores},
	description={Un \emph{actuador} es un dispositivo electr{\'o}nico conectado a un sistema controlador que permite que el sistema interact{\'u}e con el mundo real, ejecutando acciones de acuerdo con la \gls{logicaprogramacion} previamente cargada. En \glspl{sibiot}, los actuadores reciben instrucciones derivadas de datos de \glspl{sensor} y realizan tareas f{\'i}sicas como mover, encender, abrir o ajustar mecanismos.},
	type=\glsdefaulttype,
	symbol={🤖},
	user1={Dispositivo de acci{\'o}n},
	user2={Elemento de salida},
	user3={Controlador f{\'i}sico},
	see={sensor,iot,sibiot}
}

\newglossaryentry{actualizacion}{
	name={Actualizaci{\'o}n},
	text={actualizaci{\'o}n},
	sort={actualizacion},
	plural={actualizaciones},
	first={actualizaci{\'o}n},
	firstplural={actualizaciones},
	description={Proceso mediante el cual se instalan nuevas versiones de software o firmware en \glspl{dispositivo} y \glspl{servicio} para corregir vulnerabilidades, mejorar rendimiento, incorporar funcionalidades y garantizar compatibilidad. En \gls{iot}, las actualizaciones suelen realizarse de forma remota (\emph{OTA}) y requieren mecanismos seguros de distribuci{\'o}n, verificaci{\'o}n de integridad y plan de reversi{\'o}n para preservar la continuidad operativa.},
	type=main,
	symbol={🛠️},
	see={provisionamiento,seguridad,observabilidad,gateway},
	user1={Gest{\'o}n de versiones},
	user2={Parcheo y mejora continua},
	user3={Distribuci{\'o}n segura (OTA)}
}

\newglossaryentry{adaptabilidad}
{
	name={Adaptabilidad},
	text={adaptabilidad},
	sort={adaptabilidad},
	plural={capacidades de adaptabilidad},
	first={adaptabilidad},
	firstplural={capacidades de adaptabilidad},
	description={La \emph{adaptabilidad} es la capacidad de un \gls{sistema}, componente o tecnolog{\'i}a para ajustarse eficazmente a cambios en su \gls{entorno}, \glspl{requisito}, carga de trabajo o condiciones de operaci{\'o}n sin comprometer su funcionalidad ni su \gls{rendimiento}. En ingenier{\'i}a de software y \glspl{sibiot}, la adaptabilidad es clave para garantizar \gls{flexibilidad}, \gls{escalabilidad} y resiliencia frente a la evoluci{\'o}n del contexto o de los usuarios.},
	type=\glsdefaulttype,
	symbol={🔁},
	user1={Resiliencia},
	user2={Flexibilidad operativa},
	user3={Capacidad de ajuste},
	see={flexibilidad}
}

\newglossaryentry{adaptable}{
	name={Adaptable},
	text={adaptable},
	sort={adaptable},
	plural={adaptables},
	first={adaptable},
	firstplural={adaptables},
	description={El t{\'e}rmino \emph{adaptable} califica a un \gls{sistema}, componente o tecnolog{\'i}a que puede ajustarse de manera eficaz a cambios en su entorno, requisitos, carga de trabajo o condiciones de operaci{\'o}n, manteniendo un rendimiento aceptable. A diferencia de la \gls{adaptabilidad}, que es la propiedad general, lo adaptable describe el atributo concreto de un elemento espec{\'i}fico. En \glspl{sibiot}, un dispositivo o servicio es adaptable cuando puede integrarse con diferentes \glspl{protocolocomunicacion}, soportar actualizaciones din{\'a}micas o modificar su configuraci{\'o}n en funci{\'o}n de las necesidades del entorno.},
	type=main,
	symbol={🔧},
	user1={Calificativo de sistemas o componentes flexibles},
	user2={Ajuste a condiciones cambiantes en IoT},
	user3={Vinculado con la propiedad de adaptabilidad},
	see={adaptabilidad,flexibilidad,escalabilidad}
}

\newglossaryentry{adaptacion}{
	name={Adaptaci{\'o}n},
	text={adaptaci{\'o}n},
	sort={adaptacion},
	type={main},
	plural={adaptaciones},
	description={Proceso mediante el cual un sistema, componente o tecnolog{\'i}a ajusta su configuraci{\'o}n, par{\'a}metros o comportamiento para responder adecuadamente a cambios en el entorno, condiciones operativas o requisitos funcionales. En los sistemas basados en IoT, la adaptaci{\'o}n permite mantener la conectividad, el rendimiento y la resiliencia frente a variaciones din{\'a}micas del entorno f{\'i}sico o l{\'o}gico},
	first={adaptaci{\'o}n},
	firstplural={adaptaciones},
	symbol={🔧},
	user1={Proceso t{\'e}cnico de ajuste din{\'a}mico},
	user2={Relacionado con autooptimizaci{\'o}n en sistemas IoT},
	user3={Clave para resiliencia y continuidad operativa},
	user4={Puede implicar cambios autom{\'a}ticos o manuales}
}

\newglossaryentry{adaptaciontecnologica}{
	name={Adaptaci{\'o}n tecnol{\'o}gica},
	text={adaptaci{\'o}n tecnol{\'o}gica},
	sort={adaptacion tecnologica},
	type={main},
	plural={adaptaciones tecnol{\'o}gicas},
	description={Proceso mediante el cual una organizaci{\'o}n, sistema o conjunto de dispositivos adopta, integra o ajusta tecnolog{\'i}as para mejorar su funcionamiento, actualizar sus capacidades o responder a nuevas demandas del entorno. En los sistemas de informaci{\'o}n basados en IoT, la adaptaci{\'o}n tecnol{\'o}gica implica la selecci{\'o}n, ajuste y despliegue de tecnolog{\'i}as de conectividad, hardware, software o protocolos que permitan mantener la eficiencia, la interoperabilidad y la resiliencia ante cambios din{\'a}micos o condiciones adversas},
	first={adaptaci{\'o}n tecnol{\'o}gica},
	firstplural={adaptaciones tecnol{\'o}gicas},
	symbol={⚙️},
	user1={Proceso de integraci{\'o}n y ajuste de nuevas tecnolog{\'i}as},
	user2={Clave para la modernizaci{\'o}n continua en IoT},
	user3={Relacionada con la interoperabilidad y actualizaci{\'o}n del sistema},
	user4={Puede incluir cambios en hardware, software o protocolos}
}

\newglossaryentry{adc}{
	name={Convertidor anal{\'o}gico--digital},
	text={convertidor anal{\'o}gico--digital},
	sort={convertidor analogico digital},
	first={convertidor anal{\'o}gico--digital (\textit{Analog-Digital Converter}, ADC)},
	description={Subsistema de hardware encargado de transformar se{\~n}ales anal{\'o}gicas continuas en valores digitales discretos que pueden ser procesados por el microcontrolador. En el ATmega328P, el \emph{convertidor anal{\'o}gico--digital} es de 10 bits, lo que permite representar valores entre 0 y 1023.},
	type=main,
	symbol={📊},
	user1={Digitalizaci{\'o}n},
	user2={Se{\~n}ales anal{\'o}gicas},
	user3={Entrada de sensores},
	see={sensor,prototipado}
}

\newglossaryentry{adquisicion}{
	name={Adquisici{\'o}n},
	text={adquisici{\'o}n},
	short={adquisici{\'o}n},
	long={Proceso sistem{\'a}tico de captaci{\'o}n y registro de datos},
	sort={adquisicion},
	plural={adquisiciones},
	first={adquisici{\'o}n},
	firstplural={adquisiciones},
	description={Proceso mediante el cual se capturan, registran y digitalizan datos provenientes de fuentes internas o externas, tales como sensores, dispositivos, sistemas de informaci{\'o}n o interacciones humanas. En el contexto de \gls{iot} y \glspl{sibiot}, la adquisici{\'o}n constituye la etapa inicial del \gls{pipeline} de \gls{procesamientodatos}, donde las se{\~n}ales f{\'i}sicas son convertidas en datos digitales mediante \glspl{sensor}, interfaces de \gls{io} o servicios \gls{api}. La calidad de la adquisici{\'o}n influye directamente en la \gls{exactitud}, \gls{integridad} y confiabilidad del sistema.},
	type=\glsdefaulttype,
	symbol={📥},
	user1={Captura de datos},
	user2={Etapa inicial del pipeline},
	user3={Conversi{\'o}n se{\~n}al-dato},
	see={sensor,procesamientodatos,pipeline,exactitud,integridad}
}

\newglossaryentry{agile}{
	name={Agile},
	text={agile},
	sort={agile},
	first={agile},
	description={Enfoque de desarrollo de \gls{software} que prioriza iteraciones cortas, retroalimentaci{\'o}n continua y adaptaci{\'o}n al cambio, mediante entrega incremental y colaboraci{\'o}n frecuente entre los involucrados.},
	type=main,
	symbol={🏃},
	user1={Iterativo e incremental},
	user2={Adaptaci{\'o}n al cambio},
	user3={Entrega continua},
	see={iteracion,refactorizacion,testabilidad}
}

% ===================== Agricultura Inteligente =====================
\newglossaryentry{agriculturainteligente}{
	name={Agricultura inteligente},
	text={agricultura inteligente},
	sort={agriculturainteligente},
	plural={agriculturas inteligentes},
	first={agricultura inteligente},
	firstplural={agriculturas inteligentes},
	description={La \emph{agricultura inteligente} se refiere a la aplicaci{\'o}n de tecnolog{\'i}as digitales y sistemas \gls{iot} en el sector agr{\'i}cola, con el objetivo de optimizar la producci{\'o}n, reducir costos y mejorar la sostenibilidad. Utiliza \glspl{sensor} para medir variables ambientales y del suelo, \glspl{actuador} para controlar riego y fertilizaci{\'o}n, plataformas de \gls{procesamientodatos} para generar an{\'a}lisis predictivos y \gls{tomadecisiones} basadas en datos. Incluye pr{\'a}cticas como agricultura de precisi{\'o}n, monitorizaci{\'o}n remota, mantenimiento predictivo y uso eficiente de recursos h{\'i}dricos y energ{\'e}ticos.},
	type=main,
	symbol={🌱},
	user1={Aplicaci{\'o}n de IoT en el sector agr{\'i}cola},
	user2={Optimiza recursos h{\'i}dricos, energ{\'e}ticos y fertilizantes},
	user3={Base de la agricultura de precisi{\'o}n},
	see={iot,sensor,actuador,procesamientodatos,tomadecisiones}
}

\newglossaryentry{alarma}
{
	name={Alarma},
	text={alarma},
	sort={alarma},
	plural={alarmas},
	first={Alarma},
	description={Mecanismo o se{\~n}al de advertencia que indica la presencia de una condici{\'o}n cr{\'i}tica o de emergencia que requiere intervenci{\'o}n inmediata. 
		En los \glspl{sibiot}, las alarmas se activan autom{\'a}ticamente cuando los valores capturados por \glspl{sensor} superan umbrales predefinidos o cuando se detecta una falla estructural o de seguridad. 
		Su implementaci{\'o}n puede incluir avisos sonoros, luminosos o digitales, y est{\'a} orientada a la mitigaci{\'o}n de riesgos y la protecci{\'o}n de activos, personas y procesos. 
		Las alarmas forman parte del sistema de respuesta autom{\'a}tica y del subsistema de seguridad operacional en arquitecturas inteligentes.},
	symbol={🚨},
	user1={Aviso ante una condici{\'o}n cr{\'i}tica o de emergencia},
	user2={Activaci{\'o}n autom{\'a}tica mediante umbrales de sensores},
	user3={Elemento de seguridad y respuesta inmediata en \glspl{sibiot}},
	see={[Relacionado con:]{alerta,notificacion,seguridad,sensor}}
}

\newglossaryentry{alcance}{
	name={Alcance},
	text={alcance},
	sort={alcance},
	plural={alcances},
	first={alcance},
	firstplural={alcances},
	description={El \emph{alcance} se refiere a la distancia m{\'a}xima o la cobertura efectiva en la que un \gls{dispositivo}, sistema o \gls{red} puede transmitir y recibir informaci{\'o}n de forma fiable. Puede tener distintos contextos: (1) en telecomunicaciones, indica el l{\'i}mite f{\'i}sico de la \gls{transmision} de una se{\~n}al; (2) en gesti{\'o}n de proyectos, se refiere al conjunto de objetivos, funcionalidades y requisitos que deben cumplirse; (3) en seguridad, al grado de protecci{\'o}n o cobertura frente a riesgos. En el {\'a}mbito de \glspl{sibiot}, el alcance de las tecnolog{\'i}as de \gls{comunicacion} (como Wi-Fi, LoRa, Zigbee o LTE-M) es un factor determinante para garantizar \gls{conectividad}, \gls{escalabilidad} y continuidad de servicio.},
	type=main,
	symbol={📡},
	user1={Distancia m{\'a}xima de transmisi{\'o}n o cobertura},
	user2={Definici{\'o}n de objetivos o requisitos en proyectos},
	user3={Grado de protecci{\'o}n o aplicaci{\'o}n en seguridad},
	see={comunicacion,transmision,red,conectividad,sibiot}
}

\newglossaryentry{alerta}
{
	name={Alerta},
	text={alerta},
	sort={alerta},
	plural={alertas},
	first={alerta},
	description={Se{\~n}al o mensaje generado por un \gls{sistema} para advertir sobre una condici{\'o}n an{\'o}mala o un evento que requiere atenci{\'o}n inmediata. 
		En los \glspl{sibiot}, las alertas son emitidas por \glspl{sensor} o m{\'o}dulos de an{\'a}lisis cuando se detectan desviaciones de los par{\'a}metros normales de operaci{\'o}n. 
		Pueden manifestarse de forma visual, sonora o digital, y su objetivo es prevenir fallos, incidentes o riesgos en procesos automatizados. 
		Las alertas se integran en sistemas de monitorizaci{\'o}n inteligente, donde su gesti{\'o}n eficiente permite la respuesta oportuna y la continuidad operativa.},
	symbol={⚠️},
	user1={Advertencia generada ante un evento an{\'o}malo},
	user2={Prevenci{\'o}n de fallos o incidentes en sistemas inteligentes},
	user3={Emisi{\'o}n por sensores o algoritmos en \glspl{sibiot}},
	see={[Relacionado con:]{alarma,notificacion,evento,sensor}}
}

\newglossaryentry{algoritmo}{
	name={Algoritmo},
	text={algoritmo},
	sort={algoritmo},
	plural={algoritmos},
	first={algoritmo},
	firstplural={algoritmos},
	description={Conjunto finito y ordenado de pasos o instrucciones bien definidas que permiten resolver un problema, realizar un c{\'a}lculo o ejecutar una tarea. Un \emph{algoritmo} puede representarse en lenguaje natural, pseudoc{\'o}digo, diagramas de flujo o implementarse en un lenguaje de programaci{\'o}n. En la ingenier{\'i}a de software, los algoritmos constituyen la base de los programas y sistemas, garantizando que las operaciones sean reproducibles y eficientes. En el contexto de los \glspl{sibiot} y de los \glspl{sistemainteligente}, los algoritmos procesan datos provenientes de \glspl{sensor}, aplican t{\'e}cnicas de \gls{aprendizajeautomatico}, \gls{ia} o \gls{modelopredictivo}, y generan acciones o decisiones autom{\'a}ticas en tiempo real},
	type=main,
	see={ia,aprendizajeautomatico,modelopredictivo,sistemainteligente},
	user1={Procedimiento definido paso a paso},
	user2={Base de la programaci{\'o}n y la l{\'o}gica computacional},
	user3={Motor de procesamiento en SIbIoT},
	symbol={🔢}
}

\newglossaryentry{algoritmoeficiente}{
	name={Algoritmo eficiente},
	text={algoritmo eficiente},
	sort={algoritmo eficiente},
	type={main},
	plural={algoritmos eficientes},
	description={Procedimiento computacional que resuelve un problema optimizando el uso de recursos, tales como tiempo de ejecuci{\'o}n, memoria o consumo energ{\'e}tico. En el contexto de los sistemas basados en IoT, un algoritmo eficiente permite ejecutar tareas en dispositivos con recursos limitados, mantener la conectividad, procesar datos en tiempo real y reducir el impacto energ{\'e}tico sin comprometer la precisi{\'o}n ni la funcionalidad},
	first={algoritmo eficiente},
	firstplural={algoritmos eficientes},
	symbol={⚡},
	user1={Optimiza tiempo y recursos computacionales},
	user2={Fundamental para dispositivos con restricciones de hardware},
	user3={Clave en IoT para reducir consumo energ{\'e}tico},
	user4={Puede involucrar t{\'e}cnicas de optimizaci{\'o}n, heur{\'i}sticas o estructuras de datos eficientes}
}

\newglossaryentry{algoritmogestionanchobanda}{
	name={Algoritmo de gesti{\'o}n de ancho de banda},
	text={algoritmo de gesti{\'o}n de ancho de banda},
	sort={algoritmo de gestion de ancho de banda},
	type={main},
	plural={algoritmos de gesti{\'o}n de ancho de banda},
	description={Conjunto de reglas, m{\'e}todos o procedimientos destinados a asignar, regular y optimizar el uso del ancho de banda disponible en una red. En los sistemas basados en IoT, estos algoritmos permiten priorizar tr{\'a}fico cr{\'i}tico, evitar congesti{\'o}n, equilibrar la carga entre dispositivos y garantizar niveles adecuados de calidad del servicio (QoS), especialmente en entornos con recursos de comunicaci{\'o}n limitados},
	first={algoritmo de gesti{\'o}n de ancho de banda},
	firstplural={algoritmos de gesti{\'o}n de ancho de banda},
	symbol={📶},
	user1={Regula y optimiza la utilizaci{\'o}n del ancho de banda},
	user2={Previene congesti{\'o}n y p{\'e}rdida de paquetes},
	user3={Fundamental para garantizar QoS en redes IoT},
	user4={Puede incluir priorizaci{\'o}n, asignaci{\'o}n din{\'a}mica y control adaptativo}
}

\newglossaryentry{algoritmogestionenergia}{
	name={Algoritmo de gesti{\'o}n de energ{\'i}a},
	text={algoritmo de gesti{\'o}n de energ{\'i}a},
	sort={algoritmo de gestion de energia},
	plural={algoritmos de gesti{\'o}n de energ{\'i}a},
	first={algoritmo de gesti{\'o}n de energ{\'i}a},
	firstplural={algoritmos de gesti{\'o}n de energ{\'i}a},
	description={Procedimiento computacional dise{\~n}ado para optimizar el consumo, distribuci{\'o}n y aprovechamiento de energ{\'i}a en un \gls{sistema}, priorizando la \gls{eficienciaenergetica} y la \gls{sostenibilidad}. Un \emph{algoritmo de gesti{\'o}n de energ{\'i}a} puede incluir estrategias de ahorro, balanceo de cargas, uso de fuentes renovables, almacenamiento en bater{\'i}as y predicci{\'o}n de demanda. En el contexto de los \glspl{sibiot}, estos algoritmos permiten prolongar la autonom{\'i}a de \glspl{sensor} y \glspl{actuador}, optimizar el funcionamiento de \glspl{sistemarespaldoenergia}, reducir la \gls{sobrecarga} en \glspl{nodo} y favorecer la integraci{\'o}n con redes el{\'e}ctricas inteligentes, garantizando \gls{eficienciaoperativa} y continuidad de servicio en aplicaciones cr{\'i}ticas},
	type=main,
	symbol={🔋},
	user1={Optimizar consumo y distribuci{\'o}n de energ{\'i}a},
	user2={Prolongar autonom{\'i}a y reducir sobrecarga},
	user3={Aplicaci{\'o}n en IoT y redes inteligentes},
	see={eficienciaenergetica,sistemarespaldoenergia,sostenibilidad,eficienciaoperativa,sibiot}
}

\newglossaryentry{algoritmogestionrecursoseficientes}{
	name={Algoritmo de gesti{\'o}n de recursos eficientes},
	text={algoritmo de gesti{\'o}n de recursos eficientes},
	sort={algoritmo de gestion de recursos eficientes},
	plural={algoritmos de gesti{\'o}n de recursos eficientes},
	first={algoritmo de gesti{\'o}n de recursos eficiente},
	firstplural={algoritmos de gesti{\'o}n de recursos eficientes},
	description={Conjunto de procedimientos computacionales dise{\~n}ados para optimizar el uso de recursos limitados, como energ{\'i}a, memoria, \gls{anchobanda} o capacidad de procesamiento, garantizando el mejor desempe{\~n}o posible en un \gls{sistema}. Los \emph{algoritmos de gesti{\'o}n de recursos eficientes} buscan equilibrar la carga de trabajo, reducir la \gls{latencia}, minimizar la \gls{sobrecarga} y prolongar la vida {\'u}til de los dispositivos. En el contexto de los \glspl{sibiot}, permiten asignar tareas entre \glspl{nodo}, optimizar la comunicaci{\'o}n entre \glspl{sensor} y \glspl{actuador}, y gestionar la energ{\'i}a en sistemas de bajo consumo, contribuyendo a la \gls{eficienciaoperativa}, la \gls{sostenibilidad} y la resiliencia.},
	type=main,
	see={algoritmo,eficienciaoperativa,sostenibilidad,resiliencia,sibiot},
	user1={Optimizar el uso de recursos limitados},
	user2={Minimizar sobrecarga y latencia},
	user3={Soporte a eficiencia en IoT y sistemas distribuidos},
	symbol={⚖️}
}

\newglossaryentry{algoritmointeligente}{
	name={Algoritmo inteligente},
	text={algoritmo inteligente},
	sort={algoritmo inteligente},
	plural={algoritmos inteligentes},
	first={algoritmo inteligente},
	firstplural={algoritmos inteligentes},
	description={Procedimiento computacional que incorpora t{\'e}cnicas de \gls{ia}, \gls{aprendizajeautomatico}, heur{\'i}sticas o modelos adaptativos para resolver problemas complejos de manera aut{\'o}noma, aprendiendo de los datos y mejorando su desempe{\~n}o con el tiempo. Un \emph{algoritmo inteligente} es capaz de razonar, predecir, optimizar y tomar decisiones en entornos inciertos o cambiantes. En el contexto de los \glspl{sibiot}, estos algoritmos se aplican en la \gls{monitorizacion} avanzada de \glspl{variablefisica}, el \gls{mantenimientopredictivo}, la gesti{\'o}n de energ{\'i}a y la personalizaci{\'o}n de servicios, aportando \gls{adaptabilidad}, \gls{eficiencia} y sostenibilidad},
	type=main,
	see={algoritmo,ia,aprendizajeautomatico,mantenimientopredictivo,sibiot},
	symbol={🧠},
	user1={Uso de IA y aprendizaje para resolver problemas},
	user2={Capacidad de adaptaci{\'o}n y mejora continua},
	user3={Aplicaci{\'o}n en toma de decisiones en IoT}
}

\newglossaryentry{almacenamiento}
{
	name={Almacenamiento},
	text={almacenamiento},
	sort={almacenamiento},
	plural={almacenamientos},
	first={almacenamiento},
	firstplural={almacenamientos},
	description={El \emph{almacenamiento} es el proceso y la tecnolog{\'i}a destinada a guardar, conservar y gestionar \glspl{dato} en soportes f{\'i}sicos o virtuales. En el contexto de \gls{iot}, el almacenamiento puede darse en dispositivos locales, \gls{edgecomputing}, \gls{fogcomputing} o en la \gls{nube}, garantizando la disponibilidad, integridad y escalabilidad de la informaci{\'o}n. Es un componente esencial de los \glspl{sibiot}, pues permite la trazabilidad hist{\'o}rica, el an{\'a}lisis posterior y la toma de decisiones basada en evidencia.},
	type=\glsdefaulttype,
	symbol={💾},
	user1={Storage},
	user2={Memoria de datos},
	user3={Repositorio digital},
	see={procesamientodatos}
}

\newglossaryentry{almacenamientoenergetico}{
	name={Almacenamiento energ{\'e}tico},
	text={almacenamiento energ{\'e}tico},
	sort={almacenamiento energetico},
	plural={sistemas de almacenamiento energ{\'e}tico},
	first={almacenamiento energ{\'e}tico},
	firstplural={sistemas de almacenamiento energ{\'e}tico},
	description={Tecnolog{\'i}as que conservan energ{\'i}a el{\'e}ctrica y la liberan cuando se requiere, equilibrando generaci{\'o}n y consumo en el hogar.},
	type=main,
	symbol={🔋},
	see={pv,hems,gestionenergetica},
	user1={Bater{\'i}as residenciales},
	user2={Desplazamiento de carga},
	user3={Gestor de picos}
}

\newglossaryentry{amazonaurora}{
	name={Amazon Aurora},
	text={Amazon Aurora},
	sort={amazon aurora},
	plural={Amazon Aurora},
	first={Amazon Aurora},
	firstplural={Amazon Aurora},
	description={\emph{Amazon Aurora} es un servicio de \gls{basedatosrelacional} administrado en \gls{aws}, compatible con MySQL y PostgreSQL, que ofrece alto rendimiento y disponibilidad. Aurora est{\'a} dise{\~n}ada para combinar la simplicidad de las bases de datos comerciales con la \gls{escalabilidad} y la econom{\'i}a de los servicios en la \gls{nube}. Su arquitectura desacopla el almacenamiento distribuido de la capa de procesamiento, soportando replicaci{\'o}n autom{\'a}tica, backups continuos y recuperaci{\'o}n ante desastres. En ecosistemas \gls{iot}, Aurora se emplea para gestionar datos cr{\'i}ticos de aplicaciones sensibles al tiempo, integrando an{\'a}lisis en tiempo real, seguridad avanzada y escalamiento bajo demanda},
	type=main,
	symbol={☁️},
	user1={Servicio de base de datos relacional en AWS},
	user2={Compatibilidad con MySQL y PostgreSQL},
	user3={Escalabilidad autom{\'a}tica en la nube},
	see={aws,basedatosrelacional,nube,sibiot}
}

\newglossaryentry{amazonglacier}{
	name={Amazon Glacier},
	text={Amazon Glacier},
	sort={amazon glacier},
	first={Amazon Glacier},
	description={Servicio de almacenamiento de bajo costo ofrecido por Amazon Web Services (AWS), optimizado para la conservaci{\'o}n a largo plazo de datos a los que se accede con poca frecuencia. Amazon Glacier est{\'a} orientado al archivado, respaldo y retenci{\'o}n de informaci{\'o}n durante largos periodos, ofreciendo distintos tiempos de recuperaci{\'o}n seg{\'u}n la clase de almacenamiento. En sistemas SIbIoT, Glacier se utiliza para preservar datos hist{\'o}ricos, registros de auditor{\'i}a y evidencias operativas requeridas por normativas o an{\'a}lisis retrospectivos},
	type=main,
	symbol={🧊},
	see={almacenamiento,archivado,seguridad,analiticahistorica,sibiot}
}

\newglossaryentry{amazonredshift}{
	name={Amazon Redshift},
	text={Amazon Redshift},
	sort={amazon redshift},
	first={Amazon Redshift},
	description={Servicio gestionado de almacenamiento y an{\'a}lisis de datos en columnas (\textit{data warehouse}) ofrecido por \gls{aws}, dise{\~n}ado para ejecutar consultas anal{\'i}ticas complejas sobre grandes vol{\'u}menes de datos. \emph{Amazon Redshift} est{\'a} optimizado para cargas de trabajo \gls{olap} y soporta procesamiento masivo en paralelo. En el contexto de \gls{sibiot}, Redshift se emplea para el an{\'a}lisis hist{\'o}rico, la generaci{\'o}n de indicadores, la explotaci{\'o}n de datos agregados y el soporte a la \gls{tomadecisiones} estrat{\'e}gicas},
	type=main,
	symbol={🏢},
	see={analiticadatos,procesamientoparalelo,procesamientobatch,nube,sibiot}
}

\newglossaryentry{amazons3}{
	name={Amazon S3},
	text={Amazon S3},
	sort={amazons3},
	first={Amazon S3},
	description={Servicio de almacenamiento de objetos ofrecido por Amazon Web Services (AWS), dise{\~n}ado para proporcionar alta durabilidad, disponibilidad y escalabilidad pr{\'a}cticamente ilimitada. Amazon S3 permite almacenar y recuperar datos no estructurados, tales como archivos, registros, respaldos y conjuntos de datos hist{\'o}ricos. En arquitecturas SIbIoT, Amazon S3 se emplea como repositorio de datos brutos, almacenamiento hist{\'o}rico y base para lagos de datos (\textit{data lakes}), soportando an{\'a}lisis posteriores y procesos batch},
	type=main,
	symbol={🪣},
	see={nube,almacenamiento,procesamientobatch,analiticadatos,sibiot}
}

% Acr{\'o}nimo completo: Amazon EMR
\newglossaryentry{amazonemr}{
	name={Amazon EMR},
	text={Amazon EMR},
	first={servicio el{\'a}stico de procesamiento distribuido de Amazon (\textit{Amazon Elastic MapReduce}, Amazon EMR)},
	firstplural={Amazon EMRs},
	sort={amazon emr},
	description={\emph{Amazon EMR} es un servicio administrado de \gls{aws} para procesar grandes vol{\'u}menes de \glspl{dato} con cl{\'u}steres escalables que ejecutan marcos como Apache Hadoop, Apache Spark, Apache Hive y Presto, automatizando el aprovisionamiento, la configuraci{\'o}n y el escalado de recursos en la \gls{nube}.
	}
}

\newglossaryentry{amazonrds}{
	name={Amazon RDS},
	text={Amazon RDS},
	sort={amazonrds},
	first={Amazon RDS},
	description={Servicio gestionado de bases de datos relacionales ofrecido por Amazon Web Services (AWS), dise{\~n}ado para simplificar la implementaci{\'o}n, operaci{\'o}n y escalado de motores de bases de datos relacionales en la nube. Amazon RDS soporta sistemas gestores como PostgreSQL, MySQL, MariaDB, Oracle y SQL Server, automatizando tareas administrativas tales como aprovisionamiento, copias de seguridad, replicaci{\'o}n, actualizaciones y gesti{\'o}n de alta disponibilidad. En sistemas de informaci{\'o}n basados en IoT y SIbIoT, Amazon RDS se emplea como capa de persistencia operacional para datos transaccionales y de control},
	type=main,
	symbol={☁️},
	see={basedatos,postgresql,oltp,nube,sibiot}
}


\newglossaryentry{amenaza}{
	name={Amenaza},
	text={amenaza},
	sort={amenaza},
	plural={amenazas},
	first={amenaza},
	firstplural={amenazas},
	description={Evento potencial, circunstancia o agente que puede explotar una vulnerabilidad y afectar confidencialidad, integridad o disponibilidad de activos de informaci{\'o}n. Puede ser de origen humano, t{\'e}cnico, natural u organizacional},
	type=main,
	symbol={⚠️},
	user1={Gesti{\'o}n de riesgos},
	user2={Seguridad de la informaci{\'o}n},
	user3={Ciberseguridad},
	see={riesgo,vulnerabilidad,ataque}
}

\newglossaryentry{amqp}{
	name={AMQP},
	text={AMQP},
	sort={amqp},
	plural={AMQPs},
	first={protocolo avanzado de encolado de mensajes (\textit{Advanced Message Queuing Protocol}, AMQP)},
	description={Protocolo de mensajer{\'i}a estandarizado para \glspl{sistemadistribuido} que define c{\'o}mo enviar/recibir mensajes de forma fiable y eficiente. Suele emplear colas, rutas e intercambios para desacoplar productores y consumidores},
	type=main,
	symbol={},
	user1={Mensajer{\'i}a},
	user2={Sistemas distribuidos},
	user3={Integraci{\'o}n de servicios},
	see={mqtt,protocolocomunicacion,arquitecturadistribuida}
}

\newglossaryentry{analisis}
{
	name={An{\'a}lisis},
	text={an{\'a}lisis},
	sort={analisis},
	plural={an{\'a}lisis},
	first={An{\'a}lisis},
	description={Proceso sistem{\'a}tico de descomposici{\'o}n, interpretaci{\'o}n y evaluaci{\'o}n de datos, componentes o sistemas con el fin de comprender su estructura, funcionamiento o comportamiento. En el contexto de los \glspl{sibiot}, el an{\'a}lisis implica la recopilaci{\'o}n y el procesamiento de informaci{\'o}n proveniente de \glspl{sensor}, permitiendo identificar patrones, detectar anomal{\'i}as y extraer conocimiento valioso para la toma de decisiones. Este proceso puede incluir diversas ramas como el \gls{analisisdescriptivo}, el \gls{analisispredictivo} y el \gls{analisisprescriptivo}, las cuales aplican t{\'e}cnicas estad{\'i}sticas, \glspl{algoritmo} de \gls{aprendizajeautomatico} y modelos de \gls{ia} para convertir los datos en informaci{\'o}n accionable.},
	symbol={🔍},
	user1={Examen sistem{\'a}tico de datos o sistemas},
	user2={Base para la toma de decisiones en entornos inteligentes},
	user3={Incluye an{\'a}lisis descriptivo, predictivo y prescriptivo},
	see={dato,aprendizajeautomatico,ia,analisispredictivo,analisisprescriptivo,sibiot}
}

\newglossaryentry{analisisbigdata}{
	name={an{\'a}lisis Big Data},
	text={an{\'a}lisis Big Data},
	sort={analisisbigdata},
	plural={an{\'a}lisis Big Data},
	first={an{\'a}lisis Big Data},
	firstplural={an{\'a}lisis Big Data},
	description={El an{\'a}lisis Big Data comprende el uso de t{\'e}cnicas avanzadas de \gls{mineriadatos}, \gls{aprendizajeautomatico} y \gls{analisispredictivo} para identificar patrones, correlaciones y tendencias en conjuntos de datos de gran tama{\~n}o y complejidad. En los \glspl{sibiot}, este tipo de an{\'a}lisis permite transformar los datos generados por m{\'u}ltiples \glspl{sensor} en informaci{\'o}n estrat{\'e}gica para la \gls{tomadecisionesinformadas}. El proceso suele apoyarse en infraestructuras distribuidas como Hadoop o Spark, que optimizan el rendimiento del \gls{procesamiento} paralelo y la escalabilidad del sistema.},
	type=\glsdefaulttype,
	symbol={📊},
	user1={Big Data analysis},
	user2={Explotaci{\'o}n masiva de datos para obtener conocimiento estrat{\'e}gico},
	user3={An{\'a}lisis avanzado de datos a gran escala mediante t{\'e}cnicas estad{\'i}sticas y de aprendizaje autom{\'a}tico},
	see={mineriadatos,aprendizajeautomatico,procesamiento,sibiot,tomadecisionesinformadas}
}

\newglossaryentry{analisisdatos}{
	name={An{\'a}lisis de datos},
	text={an{\'a}lisis de datos},
	sort={analisis de datos},
	plural={an{\'a}lisis de datos},
	first={an{\'a}lisis de datos},
	firstplural={an{\'a}lisis de datos},
	description={Proceso sistem{\'a}tico de inspeccionar, limpiar, transformar y modelar datos para descubrir informaci{\'o}n {\'u}til, formular conclusiones y respaldar la \gls{tomadecisiones}. En \gls{sibiot} permite extraer conocimiento a partir de datos de \glspl{sensor} para \gls{monitorizacion}, predicci{\'o}n y optimizaci{\'o}n},
	type=main,
	symbol={},
	user1={Ciencia de datos},
	user2={Apoyo a decisiones},
	user3={SIbIoT y anal{\'i}tica},
	see={mineriadatos,aprendizajeautomatico,sibiot}
}

\newglossaryentry{analisisdescriptivo}
{
	name={An{\'a}lisis descriptivo},
	text={an{\'a}lisis descriptivo},
	sort={analisis descriptivo},
	plural={an{\'a}lisis descriptivos},
	first={an{\'a}lisis descriptivo},
	description={T{\'e}cnica de \gls{analisis} de datos orientada a resumir, caracterizar y representar la informaci{\'o}n recogida de diferentes fuentes para comprender su estado actual. 
		Este tipo de an{\'a}lisis utiliza medidas estad{\'i}sticas, visualizaciones y res{\'u}menes de datos para identificar comportamientos, tendencias y distribuciones sin realizar predicciones. 
		En los \glspl{sibiot}, el an{\'a}lisis descriptivo permite observar el comportamiento de variables provenientes de \glspl{sensor}, facilitando la detecci{\'o}n de patrones normales o an{\'o}malos en sistemas de monitorizaci{\'o}n ambiental, salud conectada, agricultura inteligente e industria 4.0. 
		Constituye la base sobre la cual se desarrollan los modelos de \gls{analisispredictivo} y \gls{analisisprescriptivo}.},
	symbol={📊},
	user1={Resumen e interpretaci{\'o}n de datos existentes},
	user2={Identificaci{\'o}n de patrones y tendencias sin predicci{\'o}n},
	user3={Base para el an{\'a}lisis predictivo y prescriptivo en SIbIoT},
	see={[Relacionado con:]{analisis,analisispredictivo,analisisprescriptivo,dato,sibiot}}
}

\newglossaryentry{analisispredictivo}
{
	name={An{\'a}lisis predictivo},
	text={an{\'a}lisis predictivo},
	sort={analisis predictivo},
	plural={an{\'a}lisis predictivos},
	first={an{\'a}lisis predictivo},
	description={T{\'e}cnica avanzada de \gls{analisis} de datos que emplea m{\'e}todos estad{\'i}sticos, modelos matem{\'a}ticos y \glspl{algoritmo} de \gls{aprendizajeautomatico} para identificar patrones y tendencias en grandes vol{\'u}menes de informaci{\'o}n. 
		Su objetivo es estimar o anticipar la ocurrencia de eventos futuros con base en datos hist{\'o}ricos y variables contextuales. 
		En los \glspl{sibiot}, el an{\'a}lisis predictivo permite anticipar fallos en \glspl{dispositivo}, optimizar el mantenimiento preventivo, gestionar recursos de forma eficiente y mejorar la toma de decisiones mediante modelos de \gls{ia}. 
		Se aplica en sectores como la agricultura inteligente, la salud conectada, la gesti{\'o}n ambiental, la industria 4.0 y los sistemas de energ{\'i}a inteligente.},
	symbol={📈},
	user1={Predicci{\'o}n de eventos futuros mediante datos hist{\'o}ricos},
	user2={Uso de modelos estad{\'i}sticos y \glspl{algoritmo} de aprendizaje},
	user3={Aplicaciones en mantenimiento predictivo e inteligencia operativa},
	see={[Relacionado con:]{analisis,aprendizajeautomatico,ia,dato,sibiot}}
}

\newglossaryentry{analisisprescriptivo}
{
	name={An{\'a}lisis prescriptivo},
	text={an{\'a}lisis prescriptivo},
	sort={analisis prescriptivo},
	plural={an{\'a}lisis prescriptivos},
	first={an{\'a}lisis prescriptivo},
	description={Etapa avanzada del \gls{analisis} de datos que no solo busca comprender o predecir lo que puede ocurrir, sino tambi{\'e}n recomendar acciones espec{\'i}ficas para alcanzar un resultado deseado. Se basa en el uso combinado de \glspl{algoritmo} de \gls{aprendizajeautomatico}, optimizaci{\'o}n matem{\'a}tica, \gls{ia} y simulaci{\'o}n para determinar la mejor decisi{\'o}n posible dentro de un conjunto de alternativas. En los \glspl{sibiot}, el an{\'a}lisis prescriptivo permite ajustar en tiempo real los par{\'a}metros de operaci{\'o}n de \glspl{dispositivo}, coordinar \glspl{actuador} y optimizar procesos industriales, agr{\'i}colas o urbanos de manera aut{\'o}noma. Su aplicaci{\'o}n mejora la eficiencia, reduce costos y potencia la capacidad adaptativa de los sistemas inteligentes.},
	symbol={🧠},
	user1={Recomendaci{\'o}n de acciones basadas en resultados esperados},
	user2={Uso de optimizaci{\'o}n, simulaci{\'o}n e inteligencia artificial},
	user3={Aplicaci{\'o}n en toma de decisiones autom{\'a}tica en SIbIoT},
	see={[Relacionado con:]{analisis,analisispredictivo,analisisdescriptivo,ia,sibiot}}
}

\newglossaryentry{analisisrequisitos}{
	name={An{\'a}lisis de requisitos},
	text={an{\'a}lisis de requisitos},
	sort={analisisrequisitos},
	plural={an{\'a}lisis de requisitos},
	first={an{\'a}lisis de requisitos},
	firstplural={an{\'a}lisis de requisitos},
	description={El an{\'a}lisis de requisitos constituye una de las fases fundamentales de la \gls{ingenieriasoftware}. Su prop{\'o}sito es comprender qu{\'e} debe hacer el sistema y bajo qu{\'e} condiciones debe operar. Incluye actividades como la elicitaci{\'o}n, modelado, especificaci{\'o}n, validaci{\'o}n y gesti{\'o}n de requisitos. En los \glspl{sibiot}, este proceso adquiere especial relevancia, pues permite definir los comportamientos esperados de los \glspl{dispositivoconectado}, la \gls{interoperabilidad} entre componentes y los criterios de \gls{rendimiento}, \gls{seguridad} y \gls{escalabilidad}. Una correcta ejecuci{\'o}n del an{\'a}lisis de requisitos garantiza la alineaci{\'o}n entre las necesidades del usuario y las funcionalidades implementadas.},
	type=\glsdefaulttype,
	symbol={📋},
	user1={Requirements analysis},
	user2={Proceso de identificaci{\'o}n y documentaci{\'o}n de necesidades del sistema},
	user3={Fase de la ingenier{\'i}a de software orientada a definir qu{\'e} debe hacer el sistema},
	see={ingenieriasoftware,usuario,requisito,dispositivoconectado,interoperabilidad}
}

\newglossaryentry{analisisstreaming}{
	name={An{\'a}lisis en streaming},
	text={an{\'a}lisis en streaming},
	sort={analisis en streaming},
	first={an{\'a}lisis en streaming},
	description={Tratamiento anal{\'i}tico de datos conforme llegan como flujo continuo, con latencia reducida y procesamiento incremental. En \glspl{sibiot}, habilita detecci{\'o}n temprana de eventos y respuesta casi en \gls{tiemporeal}.},
	type=main,
	symbol={🌊},
	user1={Flujos continuos},
	user2={Baja latencia},
	user3={Procesamiento incremental},
	see={evento,monitorizacion,websocket}
}

\newglossaryentry{analisistrafico}{
	name={An{\'a}lisis de tr{\'a}fico},
	text={an{\'a}lisis de tr{\'a}fico},
	sort={analisis de trafico},
	type={main},
	plural={an{\'a}lisis de tr{\'a}fico},
	description={Proceso de observaci{\'o}n, medici{\'o}n e interpretaci{\'o}n del flujo de datos que circula a trav{\'e}s de una red con el fin de identificar patrones, evaluar rendimiento, detectar anomal{\'i}as o diagnosticar problemas de conectividad. En sistemas basados en IoT, el an{\'a}lisis de tr{\'a}fico permite comprender el comportamiento de los dispositivos, optimizar el uso del ancho de banda, mejorar la seguridad y garantizar la eficiencia operativa de la red},
	first={an{\'a}lisis de tr{\'a}fico},
	firstplural={an{\'a}lisis de tr{\'a}fico},
	symbol={📡},
	user1={Herramienta clave para evaluar rendimiento y detectar fallos},
	user2={Permite identificar congesti{\'o}n, p{\'e}rdidas o comportamientos an{\'o}malos},
	user3={Fundamental en redes IoT por la alta variabilidad del tr{\'a}fico},
	user4={Apoya la optimizaci{\'o}n de protocolos y la gesti{\'o}n eficiente de recursos}
}

\newglossaryentry{analisistreaming}{
	name={an{\'a}lisis de streaming},
	text={an{\'a}lisis de streaming},
	sort={analisistreaming},
	plural={an{\'a}lisis de streaming},
	first={an{\'a}lisis de streaming},
	firstplural={an{\'a}lisis de streaming},
	description={El an{\'a}lisis de streaming consiste en aplicar m{\'e}todos anal{\'i}ticos sobre datos transmitidos de forma continua desde \glspl{sensor}, \glspl{dispositivo} o fuentes distribuidas. A diferencia del procesamiento por lotes, este enfoque proporciona resultados inmediatos, esenciales para la \gls{tomadecisionesinformadas} en entornos din{\'a}micos. En los \glspl{sibiot}, el an{\'a}lisis de streaming permite detectar patrones, anomal{\'i}as y eventos relevantes, utilizando marcos tecnol{\'o}gicos como Apache Kafka, Spark Streaming o Flink.},
	type=\glsdefaulttype,
	symbol={📡},
	user1={Stream analytics},
	user2={Procesamiento en tiempo real de flujos continuos de datos},
	user3={M{\'e}todo para obtener conocimiento instant{\'a}neo a partir de informaci{\'o}n entrante},
	see={tiemporeal,procesamiento,analisisbigdata,capturadatos}
}

\newglossaryentry{analitica}{
	name={Anal{\'i}tica},
	text={anal{\'i}tica},
	sort={analitica},
	plural={anal{\'i}ticas},
	first={anal{\'i}tica},
	firstplural={anal{\'i}ticas},
	description={Disciplina orientada a la recopilaci{\'o}n, procesamiento, interpretaci{\'o}n y modelado de datos con el fin de obtener informaci{\'o}n significativa que apoye la comprensi{\'o}n de fen{\'o}menos, la generaci{\'o}n de conocimiento y la toma de decisiones basadas en evidencia. La \emph{anal{\'i}tica} integra m{\'e}todos estad{\'i}sticos, algoritmos de \gls{aprendizajeautomatico}, t{\'e}cnicas de visualizaci{\'o}n y modelos matem{\'a}ticos para revelar patrones, tendencias y relaciones. En los \glspl{sibiot}, la anal{\'i}tica constituye el n{\'u}cleo cognitivo que transforma datos provenientes de \glspl{sensor}, dispositivos inteligentes y plataformas distribuidas en valor operativo, favoreciendo la \gls{eficienciaoperativa}, la \gls{tomadecisiones} informada y la \gls{sostenibilidad}.},
	type=main,
	symbol={📈},
	user1={Procesamiento avanzado de datos},
	user2={Modelado estad{\'i}stico y algor{\'i}tmico},
	user3={Soporte estrat{\'e}gico a decisiones en SIbIoT},
	see={eficienciaanalitica,aprendizajeautomatico,eficienciaoperativa,modelopredictivo}
}

\newglossaryentry{analiticadatos}{
	name={Anal{\'\i}tica de datos},
	text={anal{\'\i}tica de datos},
	sort={analiticadatos},
	plural={anal{\'\i}ticas de datos},
	first={anal{\'\i}tica de datos},
	firstplural={anal{\'\i}ticas de datos},
	description={Conjunto de t{\'e}cnicas, m{\'e}todos y herramientas orientadas a examinar, transformar y modelar datos con el fin de extraer informaci{\'o}n significativa, identificar patrones y apoyar la toma de decisiones. En los \glspl{sibiot}, la anal{\'\i}tica de datos permite convertir flujos de datos provenientes de dispositivos distribuidos en conocimiento accionable, integrando enfoques descriptivos, diagn{\'o}sticos, predictivos y prescriptivos sobre datos hist{\'o}ricos y en tiempo real.},
	type=main,
	symbol={📊},
	user1={An{\'a}lisis de informaci{\'o}n},
	user2={Extracci{\'o}n de conocimiento},
	user3={Soporte a decisiones},
	see={analiticahistorica,analisisprescriptivo,analisisstreaming,gestiondatos}
}

\newglossaryentry{analiticahistorica}{
	name={Anal{\'\i}tica hist{\'o}rica},
	text={anal{\'\i}tica hist{\'o}rica},
	sort={analiticahistorica},
	plural={anal{\'\i}ticas hist{\'o}ricas},
	first={anal{\'\i}tica hist{\'o}rica},
	firstplural={anal{\'\i}ticas hist{\'o}ricas},
	description={Enfoque de an{\'a}lisis de datos orientado al estudio de informaci{\'o}n almacenada a lo largo del tiempo, con el objetivo de identificar tendencias, patrones recurrentes y comportamientos pasados. En los \glspl{sibiot}, la anal{\'\i}tica hist{\'o}rica permite evaluar el desempe{\~n}o de sistemas, analizar la evoluci{\'o}n de variables del entorno y servir como base para modelos predictivos y procesos de mejora continua.},
	type=main,
	symbol={🕒},
	user1={An{\'a}lisis temporal},
	user2={Datos hist{\'o}ricos},
	user3={Identificaci{\'o}n de tendencias},
	see={analiticadatos,analisisprescriptivo,gestiondatos}
}

\newglossaryentry{anchobanda}
{
	name={Ancho de banda},
	text={ancho de banda},
	sort={ancho de banda},
	plural={anchos de banda},
	first={ancho de banda},
	firstplural={anchos de banda},
	description={El \emph{ancho de banda} es la cantidad m{\'a}xima de datos que puede transmitirse a trav{\'e}s de un canal de \gls{conectividad} en un tiempo determinado, expresada normalmente en bits por segundo (bps). Es un indicador clave del \gls{rendimiento} de las \glspl{red} y de los \glspl{sistemacomunicacion}.},
	type=\glsdefaulttype,
	symbol={📶},
	user1={Capacidad de transmisi{\'o}n},
	user2={Bandwidth},
	user3={Tasa de datos},
	see={red}
}

\newglossaryentry{analogica}{
	name={Anal{\'o}gica},
	text={anal{\'o}gica},
	sort={analogica},
	first={anal{\'o}gica},
	description={Propiedad de una se{\~n}al o magnitud que puede variar de forma continua dentro de un rango determinado. Una se{\~n}al \emph{anal{\'o}gica} no est{\'a} restringida a valores discretos, sino que puede adoptar infinitos valores intermedios, como ocurre con la temperatura, la presi{\'o}n o el voltaje. En sistemas embebidos e IoT, las se{\~n}ales anal{\'o}gicas suelen ser generadas por sensores y posteriormente convertidas a formato digital para su procesamiento.},
	type=main,
	symbol={〰️},
	user1={Se{\~n}al continua},
	user2={Valores infinitos},
	user3={Magnitudes f{\'i}sicas},
	see={analogico,digital,adc}
}

\newglossaryentry{analogico}{
	name={Anal{\'o}gico},
	text={anal{\'o}gico},
	sort={analogico},
	first={anal{\'o}gico},
	description={T{\'e}rmino que describe dispositivos, se{\~n}ales o sistemas que operan sobre magnitudes continuas. Un elemento \emph{anal{\'o}gico} puede representar directamente valores del mundo f{\'i}sico sin discretizaci{\'o}n previa, como ocurre en sensores de voltaje, corriente o se{\~n}ales biol{\'o}gicas. En microcontroladores, los valores anal{\'o}gicos suelen ser tratados mediante convertidores anal{\'o}gico--digital.},
	type=main,
	symbol={📈},
	user1={Magnitudes continuas},
	user2={Representaci{\'o}n f{\'i}sica},
	user3={Sensores},
	see={analogica,digital,adc}
}

\newglossaryentry{anemometro}{
	name={Anem{\'o}metro},
	text={anem{\'o}metro},
	short={anem{\'o}metro},
	long={Sensor de velocidad del viento basado en rotaci{\'o}n y conteo de pulsos},
	sort={anemometro},
	plural={anem{\'o}metros},
	first={anem{\'o}metro},
	firstplural={anem{\'o}metros},
	description={Un anem{\'o}metro de copa puede codificar la velocidad del viento generando cierres de contacto (o pulsos) a una frecuencia proporcional a la rotaci{\'o}n; el fabricante publica un factor de conversi{\'o}n entre pulsos por segundo y velocidad del viento.},
	see={sensor}
}

\newglossaryentry{anomalia}
{
	name={Anomal{\'i}a},
	text={anomal{\'i}a},
	sort={anomal{\'i}a},
	plural={anomal{\'i}as},
	first={anomal{\'i}a},
	firstplural={anomal{\'i}as},
	description={Una \emph{anomal{\'i}a} es un evento, condici{\'o}n o patr{\'o}n que se desv{\'i}a significativamente del comportamiento normal o esperado en un conjunto de \glspl{dato}, procesos o sistemas. En \glspl{sibiot}, la detecci{\'o}n de anomal{\'i}as es esencial para identificar fallos en \glspl{sensor}, intrusiones de \gls{seguridad}, comportamientos at{\'i}picos de \glspl{dispositivoiot} o condiciones ambientales inusuales. Se utilizan t{\'e}cnicas de \gls{aprendizajeautomatico}, estad{\'i}stica avanzada y an{\'a}lisis de \glspl{patron} para automatizar su detecci{\'o}n.},
	type=\glsdefaulttype,
	symbol={🚨},
	user1={Desviaci{\'o}n},
	user2={Evento at{\'i}pico},
	user3={Outlier},
	see={patron}
}

% Acr{\'o}nimo completo: AODV
\newglossaryentry{aodv}{
	name={AODV},
	text={AODV},
	first={Protocolo de enrutamiento ad hoc bajo demanda basado en vector de distancias, (\textit{Ad hoc On-Demand Distance Vector}, AODV)},
	plural={AODVs},
	firstplural={AODVs},
	sort={aodv},
	description={\emph{AODV} es un protocolo de enrutamiento reactivo para \glspl{red} ad hoc m{\'o}viles: establece rutas bajo demanda cuando se requiere comunicaci{\'o}n, disminuyendo la sobrecarga de la \gls{red}. Emplea mensajes de descubrimiento y mantenimiento de rutas para adaptarse din{\'a}micamente a cambios topol{\'o}gicos, favoreciendo la \gls{movilidad} y reduciendo la \gls{latencia} en escenarios \gls{iot}.}
}
	
	\newglossaryentry{api}{
		name={API},
		text={API},
		sort={api},
		plural={APIs},
		first={interfaz de programaci{\'o}n de aplicaciones, (\textit{Application Programming Interface}, API)},
		firstplural={APIs},
		description={Una \emph{API} es un conjunto de definiciones, protocolos y herramientas utilizadas para desarrollar e integrar componentes de \gls{software} o \glspl{aplicacion}. En los \glspl{sibiot}, las API facilitan la \gls{interoperabilidad} entre \glspl{dispositivo}, \glspl{servicio} y \glspl{plataforma}, permitiendo la extensi{\'o}n, automatizaci{\'o}n y comunicaci{\'o}n eficiente de funcionalidades en entornos distribuidos.},
		type=main,
		user1={Interfaz de programaci{\'o}n de aplicaciones},
		user2={Application Programming Interface},
		user3={Mecanismo de integraci{\'o}n entre sistemas de software},
		user4={Punto de conexi{\'o}n entre servicios o componentes de un sistema},
		see={interfaz,servicioweb,integracionsoftware},
		symbol={🔌}
	}
	
	\newglossaryentry{aplicacion}{
		name={Aplicaci{\'o}n},
		text={aplicaci{\'o}n},
		sort={aplicacion},
		plural={aplicaciones},
		first={aplicaci{\'o}n},
		firstplural={aplicaciones},
		description={Programa inform{\'a}tico dise{\~n}ado para ejecutar tareas o proporcionar servicios espec{\'i}ficos a un usuario o a otro sistema. Incluye \emph{aplicaciones web} y \emph{aplicaciones m{\'o}viles}. En \gls{sibiot} act{\'u}a como interfaz y capa de gesti{\'o}n para visualizaci{\'o}n, control y an{\'a}lisis en \gls{tiemporeal}, as{\'i} como integraci{\'o}n con servicios en la nube},
		type=main,
		symbol={},
		user1={Interacci{\'o}n usuario--sistema},
		user2={Capa de acceso y control},
		user3={Cliente seg{\'u}n plataforma},
		see={aplicacionweb,aplicacionmovil,software,sibiot}
	}
	
	\newglossaryentry{aplicacionempresarial}
	{
		name={Aplicaci{\'o}n empresarial},
		text={aplicaci{\'o}n empresarial},
		sort={aplicacion empresarial},
		plural={aplicaciones empresariales},
		first={aplicaci{\'o}n empresarial},
		firstplural={aplicaciones empresariales},
		description={Una \emph{aplicaci{\'o}n empresarial} es un sistema de software orientado a la gesti{\'o}n, automatizaci{\'o}n o soporte de procesos de negocio de una organizaci{\'o}n. Facilita la integraci{\'o}n de datos, la \gls{interoperabilidad} y la eficiencia operativa en dominios como log{\'i}stica, finanzas, recursos humanos o manufactura.},
		type=\glsdefaulttype,
		symbol={🏢},
		user1={Enterprise application},
		user2={Software corporativo},
		user3={Aplicaci{\'o}n de negocio},
		see={aplicacionweb}
	}
	
	\newglossaryentry{aplicacionmovil}{
		name={Aplicaci{\'o}n m{\'o}vil},
		text={aplicaci{\'o}n m{\'o}vil},
		sort={aplicacion movil},
		plural={aplicaciones m{\'o}viles},
		first={aplicaci{\'o}n m{\'o}vil},
		firstplural={aplicaciones m{\'o}viles},
		description={Programa para dispositivos port{\'a}tiles (tel{\'e}fonos inteligentes, tabletas) optimizado para pantallas t{\'a}ctiles y capaz de aprovechar recursos y sensores del dispositivo (GPS, c{\'a}mara, aceler{\'o}metro, conectividad inal{\'a}mbrica). En \gls{sibiot} habilita monitorizaci{\'o}n, control y gesti{\'o}n remota en tiempo real},
		type=main,
		see={aplicacion,aplicacionweb,sibiot},
		symbol={📱},
		user1={Software m{\'o}vil},
		user2={Interfaz en movilidad},
		user3={Cliente dependiente de plataforma}
	}
	
	\newglossaryentry{aplicacionweb}
	{
		name={Aplicaci{\'o}n web},
		text={aplicaci{\'o}n web},
		sort={aplicacion web},
		plural={aplicaciones web},
		first={aplicaci{\'o}n web},
		firstplural={aplicaciones web},
		description={Una \emph{aplicaci{\'o}n web} es una herramienta de software a la que se accede mediante un \gls{servidorweb} y un \gls{browser}, a trav{\'e}s de Internet o una intranet. Se codifica en un lenguaje interpretable por los navegadores web, a los que se conf{\'i}a su ejecuci{\'o}n.},
		type=\glsdefaulttype,
		symbol={🌐},
		user1={Web app},
		user2={Aplicaci{\'o}n en l{\'i}nea},
		user3={Programa web},
		see={servidorweb}
	}
	
	\newglossaryentry{aprendizajeautomatico}
	{
		name={Aprendizaje autom{\'a}tico},
		text={ML},
		sort={aprendizaje automatico},
		plural={aprendizaje autom{\'a}tico},
		first={aprendizaje autom{\'a}tico (\textit{Machine Learning}, ML)},
		firstplural={ML},
		description={El \emph{aprendizaje autom{\'a}tico} es una rama de la inteligencia artificial centrada en el desarrollo de algoritmos y modelos capaces de identificar patrones en \glspl{dato} y mejorar su rendimiento en tareas espec{\'i}ficas a partir de la experiencia, sin ser programados expl{\'i}citamente para cada caso. Se aplica en clasificaci{\'o}n, regresi{\'o}n, agrupamiento y detecci{\'o}n de anomal{\'i}as, entre otros.},
		type=\glsdefaulttype,
		symbol={🧠},
		user1={Machine Learning},
		user2={ML},
		user3={Aprendizaje de m{\'a}quina},
		see={bigdata}
	}
	
	\newglossaryentry{archivado}{
		name={Archivado},
		text={archivado},
		sort={archivado},
		plural={archivados},
		first={archivado},
		firstplural={archivados},
		description={Proceso de almacenamiento sistem{\'a}tico y organizado de datos, documentos o registros con el fin de preservarlos a largo plazo, facilitar su recuperaci{\'o}n y cumplir requisitos operativos, legales o hist{\'o}ricos. En los \glspl{sibiot}, el archivado resulta esencial para conservar datos generados por sensores y sistemas distribuidos, permitir an{\'a}lisis retrospectivos, respaldar auditor{\'\i}as y servir como base para la anal{\'\i}tica hist{\'o}rica y la mejora continua de los sistemas.},
		type=main,
		symbol={🗂️},
		user1={Almacenamiento hist{\'o}rico},
		user2={Preservaci{\'o}n de datos},
		user3={Recuperaci{\'o}n de informaci{\'o}n},
		see={analiticahistorica,gestiondatos,basedatos,documentacion}
	}
	
	\newglossaryentry{arduino}
	{
		name={Arduino},
		text={Arduino},
		sort={arduino},
		plural={Arduinos},
		first={Arduino},
		firstplural={Arduinos},
		description={\emph{Arduino} es una \gls{placadesarrollo} de hardware libre basada en un \gls{microcontrolador} y acompa{\~n}ada de un entorno de desarrollo integrado (IDE). Facilita la creaci{\'o}n de prototipos electr{\'o}nicos para proyectos de \gls{iot}, automatizaci{\'o}n y educaci{\'o}n tecnol{\'o}gica.},
		type=\glsdefaulttype,
		symbol={🔌},
		user1={Plataforma Arduino},
		user2={Arduino IDE},
		user3={Prototipado electr{\'o}nico},
		see={placadesarrollo,arduinopromini,arduinoduemilanove,atmega328p,arduinouno,arduinonano}
	}
	
	\newglossaryentry{arduinoduemilanove}{
		name={Arduino Duemilanove},
		text={Arduino Duemilanove},
		sort={arduino duemilanove},
		first={Arduino Duemilanove},
		description={Placa de desarrollo cl{\'a}sica de Arduino basada en los microcontroladores ATmega168 o ATmega328P. El \emph{Arduino Duemilanove} representa una generaci{\'o}n temprana del ecosistema Arduino y es frecuentemente utilizado con fines educativos para comprender los fundamentos de la programaci{\'o}n embebida.},
		type=main,
		symbol={📘},
		user1={Placa cl{\'a}sica},
		user2={Enfoque educativo},
		user3={Fundamentos embebidos},
		see={arduino,arduinonano,arduinopromini,atmega328p,arduinouno}
	}
	
	\newglossaryentry{arduinonano}{
		name={Arduino Nano},
		text={Arduino Nano},
		sort={arduino nano},
		first={Arduino Nano},
		description={Placa de desarrollo basada en el microcontrolador ATmega328P, caracterizada por su tama{\~n}o compacto y conectividad USB integrada. El \emph{Arduino Nano} es ampliamente utilizado en prototipos y sistemas embebidos donde el espacio f{\'i}sico es limitado, manteniendo compatibilidad con el ecosistema Arduino.},
		type=main,
		symbol={🔧},
		user1={Microcontrolador ATmega328P},
		user2={Placa compacta},
		user3={Prototipado},
		see={arduino,arduinopromini,arduinoduemilanove,atmega328p,arduinouno}
	}
	
	\newglossaryentry{arduinopromini}{
		name={Arduino Pro Mini},
		text={Arduino Pro Mini},
		sort={arduino pro mini},
		first={Arduino Pro Mini},
		description={Placa de desarrollo Arduino basada en el ATmega328P, dise{\~n}ada para aplicaciones embebidas definitivas. El \emph{Arduino Pro Mini} carece de interfaz USB integrada, lo que reduce su consumo y tama{\~n}o, siendo ideal para integraci{\'o}n directa en sistemas SIbIoT.},
		type=main,
		symbol={🧩},
		user1={Bajo consumo},
		user2={Integraci{\'o}n embebida},
		user3={Dise{\~n}o compacto},
		see={arduino,arduinonano,arduinoduemilanove,atmega328p,arduinouno}
	}
	
	\newglossaryentry{arduinouno}{
		name={Arduino Uno},
		text={Arduino Uno},
		sort={arduino uno},
		first={Arduino Uno},
		description={Placa de desarrollo de prop{\'o}sito general del ecosistema Arduino, basada en el microcontrolador ATmega328P. El \emph{Arduino Uno} integra una interfaz USB para programaci{\'o}n y comunicaci{\'o}n serial, regulador de voltaje y un conjunto de pines de entrada/salida digitales y anal{\'o}gicos, lo que lo convierte en una plataforma ampliamente utilizada para la ense{\~n}anza de sistemas embebidos, el prototipado r{\'a}pido y el desarrollo de aplicaciones IoT y SIbIoT de baja complejidad.},
		type=main,
		symbol={🔌},
		user1={Placa de desarrollo},
		user2={ATmega328P},
		user3={Uso educativo y prototipado},
		see={arduinonano,arduinopromini,arduinoduemilanove,atmega328p}
	}
	
	\newglossaryentry{arm}{
		name={ARM},
		text={ARM},
		plural={ARM},
		first={m{\'a}quina RISC avanzada (\textit{Advanced RISC Machine}, ARM)},
		plural={m{\'a}quinas RISC avanzadas (\textit{Advanced RISC Machines}, ARMs)},
		sort={arm},
		see={risc,microcontrolador,sistemaembebido},
		description={Arquitectura de procesadores basada en el paradigma RISC (Reduced Instruction Set Computer), desarrollada originalmente como \emph{Advanced RISC Machines}. ARM se caracteriza por su alta eficiencia energ{\'e}tica, simplicidad de dise{\~n}o y amplio ecosistema, siendo ampliamente utilizada en microcontroladores, sistemas embebidos, dispositivos IoT y sistemas ciberf{\'i}sicos.}
	}
	
	\newglossaryentry{armcortexm}
	{
		name={Arquitectura ARM Cortex-M},
		text={arquitectura ARM Cortex-M},
		first={arquitectura ARM Cortex-M},
		sort={arm cortex m},
		plural={arquitecturas ARM Cortex-M},
		first={arquitectura ARM Cortex-M},
		firstplural={arquitecturas ARM Cortex-M},
		description={La \emph{arquitectura ARM Cortex-M} es una familia de procesadores basados en ARMv7-M y ARMv8-M, dise{\~n}ados para sistemas embebidos que requieren eficiencia energ{\'e}tica, alto rendimiento y baja \gls{latencia}. Son ampliamente usados en \glspl{microcontrolador} como los STM32, ofreciendo soporte para interrupciones anidadas, operaciones en \gls{tiemporeal}, depuraci{\'o}n avanzada y gesti{\'o}n eficiente de perif{\'e}ricos. Son fundamentales en plataformas \gls{iot}, \gls{automatizacion} industrial y dispositivos port{\'a}tiles.},
		type=\glsdefaulttype,
		symbol={⚡},
		user1={Cortex-M},
		user2={Procesadores ARM embebidos},
		user3={Arquitectura eficiente},
		see={microcontrolador}
	}
	
	\newglossaryentry{arquitectura}
	{
		name={Arquitectura},
		text={arquitectura},
		sort={arquitectura},
		plural={arquitecturas},
		first={arquitectura},
		firstplural={arquitecturas},
		description={La \emph{arquitectura} es la estructura organizativa de un sistema que define sus componentes, relaciones, principios de dise{\~n}o y modos de interacci{\'o}n. En los \glspl{sibiot}, describe la disposici{\'o}n jer{\'a}rquica de \glspl{dispositivo}, servicios, redes y plataformas que permiten la adquisici{\'o}n, \gls{procesamientodatos}, almacenamiento y explotaci{\'o}n de \glspl{dato}, garantizando \gls{escalabilidad}, \gls{interoperabilidad} y eficiencia.},
		type=\glsdefaulttype,
		symbol={🏛️},
		user1={Dise{\~n}o de sistemas},
		user2={Estructura de software},
		user3={Organizaci{\'o}n del sistema},
		see={arquitecturacapas,arquitecturadistribuida,arquitecturareferencia}
	}
	
	\newglossaryentry{arquitecturaabierta}{
		name={Arquitectura abierta},
		text={arquitectura abierta},
		sort={arquitectura abierta},
		plural={arquitecturas abiertas},
		first={arquitectura abierta},
		firstplural={arquitecturas abiertas},
		description={Enfoque de dise{\~n}o de sistemas basado en el uso de est{\'a}ndares abiertos y mecanismos de interoperabilidad que facilitan la integraci{\'o}n de componentes heterog{\'e}neos. En los \glspl{sibiot}, la \emph{arquitectura abierta} es fundamental para garantizar escalabilidad, extensibilidad y acceso universal a la informaci{\'o}n.},
		type=main,
		symbol={🔓},
		user1={Est{\'a}ndares abiertos},
		user2={Interoperabilidad},
		user3={Escalabilidad},
		see={infraestructuracomunicacion,evolucion}
	}
	
	\newglossaryentry{arquitecturacapas}{
		name={Arquitectura en capas},
		text={arquitectura en capas},
		sort={arquitectura en capas},
		plural={arquitecturas en capas},
		first={arquitectura en capas},
		firstplural={arquitecturas en capas},
		description={Estilo que organiza el sistema en niveles jer{\'a}rquicos (capas) con responsabilidades separadas. Cada capa ofrece servicios a la superior y usa los de la inferior, favoreciendo modularidad, reutilizaci{\'o}n y mantenibilidad},
		type=main,
		symbol={📚},
		user1={Separaci{\'o}n de responsabilidades},
		user2={Modularidad},
		user3={Mantenibilidad},
		see={arquitectura,arquitecturaclienteservidor}
	}
	
	\newglossaryentry{arquitecturaclienteservidor}
	{
		name={Arquitectura cliente-servidor},
		text={arquitectura cliente-servidor},
		sort={arquitectura cliente servidor},
		plural={arquitecturas cliente-servidor},
		first={arquitectura cliente-servidor},
		firstplural={arquitecturas cliente-servidor},
		description={La \emph{arquitectura cliente-servidor} es un modelo de sistemas distribuidos en el que las funciones se dividen entre \glspl{servidor} y \glspl{cliente}. Los clientes env{\'i}an solicitudes y los servidores procesan y responden con datos o resultados. En \glspl{sibiot}, esta arquitectura permite estructurar la comunicaci{\'o}n entre \glspl{dispositivoiot} y plataformas de \gls{procesamientodatos}, facilitando la \gls{escalabilidad}, el control remoto y la integraci{\'o}n con la \gls{nube}.},
		type=\glsdefaulttype,
		symbol={🖧},
		user1={Cliente-servidor},
		user2={Modelo CS},
		user3={Arquitectura distribuida},
		see={cliente,servidor,arquitectura}
	}
	
	\newglossaryentry{arquitecturaeventos}{
		name={Arquitectura basada en eventos},
		text={arquitectura basada en eventos},
		sort={arquitectura basada en eventos},
		plural={arquitecturas basadas en eventos},
		first={arquitectura basada en eventos},
		firstplural={arquitecturas basadas en eventos},
		description={Estilo donde los componentes se comunican generando y reaccionando a \emph{eventos} (cambios significativos de estado). En \gls{iot} favorece alto desacoplamiento y \gls{escalabilidad} al reaccionar a lecturas de \glspl{sensor} o comandos},
		type=main,
		symbol={},
		user1={Reactividad},
		user2={Desacoplamiento},
		user3={Escalado por eventos},
		see={arquitecturadistribuida,mqtt,amqp}
	}
	
	\newglossaryentry{arquitecturadistribuida}
	{
		name={Arquitectura distribuida},
		text={arquitectura distribuida},
		sort={arquitectura distribuida},
		plural={arquitecturas distribuidas},
		first={arquitectura distribuida},
		firstplural={arquitecturas distribuidas},
		description={La \emph{arquitectura distribuida} es un modelo en el que los componentes de \gls{hardware} y software se ubican en distintos nodos interconectados por una \gls{red}, colaborando para ejecutar tareas comunes. Favorece la \gls{escalabilidad}, la \gls{toleranciafallos} y la eficiencia en el uso de recursos en \glspl{sibiot}.},
		type=\glsdefaulttype,
		symbol={🌍},
		user1={Distributed architecture},
		user2={Sistemas distribuidos},
		user3={Dise{\~n}o distribuido},
		see={arquitectura,arquitecturaeventos,msoa}
	}
	
	\newglossaryentry{arquitecturaestratificada}{
		name={Arquitectura estratificada},
		text={arquitectura estratificada},
		sort={arquitectura estratificada},
		description={%
			Enfoque arquitect{\'o}nico que organiza un sistema complejo en un conjunto de capas o estratos funcionales claramente definidos, donde cada capa encapsula responsabilidades espec{\'i}ficas y se comunica con las capas adyacentes mediante interfaces bien establecidas. 
			Este tipo de arquitectura promueve la separaci{\'o}n de preocupaciones, la modularidad, la escalabilidad y la mantenibilidad del sistema, al permitir que los cambios en una capa tengan un impacto m{\'i}nimo sobre las dem{\'a}s.
			En el contexto de los sistemas de informaci{\'o}n basados en IoT, una arquitectura estratificada facilita la integraci{\'o}n coherente del mundo f{\'i}sico, el procesamiento en el borde, la mensajer{\'i}a, la gesti{\'o}n de datos, la inteligencia operativa y las aplicaciones de usuario, soportando flujos de datos orientados a eventos y toma de decisiones distribuida.
		},
		first={arquitectura estratificada},
		plural={arquitecturas estratificadas},
		firstplural={arquitecturas estratificadas},
		see={arquitectura,capa,sistemadistribuido,iot}
	}
	
	\newglossaryentry{arquitecturajerarquica}{
		name={Arquitectura jer{\'a}rquica},
		text={arquitectura jer{\'a}rquica},
		sort={arquitectura jerarquica},
		plural={arquitecturas jer{\'a}rquicas},
		first={arquitectura jer{\'a}rquica},
		firstplural={arquitecturas jer{\'a}rquicas},
		description={Modelo de organizaci{\'o}n estructural en el que los componentes de un sistema se disponen en niveles con responsabilidades diferenciadas. En sistemas IoT y \glspl{sibiot}, este enfoque facilita la separaci{\'o}n funcional entre adquisici{\'o}n de datos, comunicaci{\'o}n y procesamiento.},
		type=main,
		symbol={🏗️},
		user1={Organizaci{\'o}n por niveles},
		user2={Separaci{\'o}n funcional},
		user3={Estructura del sistema},
		see={arquitecturajerarquicainformacion,infraestructuracomunicacion}
	}
	
	\newglossaryentry{arquitecturajerarquicainformacion}{
		name={Arquitectura jer{\'a}rquica de informaci{\'o}n},
		text={arquitectura jer{\'a}rquica de informaci{\'o}n},
		sort={arquitectura jerarquica informacion},
		plural={arquitecturas jer{\'a}rquicas de informaci{\'o}n},
		first={arquitectura jer{\'a}rquica de informaci{\'o}n},
		firstplural={arquitecturas jer{\'a}rquicas de informaci{\'o}n},
		description={Estructura de organizaci{\'o}n de la informaci{\'o}n en niveles jer{\'a}rquicos, donde los datos fluyen desde fuentes operativas hacia capas de procesamiento, an{\'a}lisis y decisi{\'o}n. Este enfoque es caracter{\'i}stico de los \glspl{sibiot}, al permitir una gesti{\'o}n escalable y sistem{\'a}tica de datos provenientes del entorno f{\'i}sico.},
		type=main,
		symbol={🗂️},
		user1={Flujo de informaci{\'o}n},
		user2={Capas informacionales},
		user3={Procesamiento jer{\'a}rquico},
		see={integraciondatos,scada,sistemainformacioncorporativo}
	}
		
	\newglossaryentry{arquitecturapoliclinica}{
		name={Arquitectura policl{\'i}nica de persistencia},
		text={arquitectura policl{\'i}nica de persistencia},
		sort={arquitecturapoliclinica},
		plural={arquitecturas policl{\'i}nicas de persistencia},
		first={arquitectura policl{\'i}nica de persistencia},
		firstplural={arquitecturas policl{\'i}nicas de persistencia},
		description={Enfoque arquitect{\'o}nico que combina de manera coordinada m{\'u}ltiples tecnolog{\'i}as de almacenamiento de datos, cada una optimizada para un tipo espec{\'i}fico de carga de trabajo, tales como procesamiento transaccional, anal{\'i}tico, almacenamiento de eventos o archivado. En una arquitectura policl{\'i}nica de persistencia, bases de datos relacionales, NoSQL, data lakes y data warehouses coexisten para maximizar el rendimiento, la escalabilidad y la flexibilidad del sistema. Este enfoque es especialmente relevante en sistemas SIbIoT, donde coexisten datos heterog{\'e}neos, flujos en tiempo real y necesidades anal{\'i}ticas avanzadas},
		type=main,
		symbol={🧬},
		see={basedatos,nosql,datalake,datawarehouse,procesamientodistribuido,sibiot}
	}
	
	\newglossaryentry{arquitecturaprocesador}{
		name={Arquitectura de procesador},
		text={arquitectura de procesador},
		sort={arquitectura de procesador},
		first={arquitectura de procesador},
		plural={arquitecturas de procesadores},
		see={cpu,risc,cisc,arm},
		description={Modelo conceptual y estructural que define la organizaci{\'o}n interna, el conjunto de instrucciones, los modos de direccionamiento, los registros y los mecanismos de ejecuci{\'o}n de un procesador. La arquitectura de procesador establece c{\'o}mo el hardware ejecuta las instrucciones del software y determina aspectos clave como el rendimiento, la eficiencia energ{\'e}tica, la escalabilidad y la compatibilidad, siendo un elemento fundamental en el dise{\~n}o de microprocesadores, microcontroladores, sistemas embebidos y plataformas de computaci{\'o}n modernas.}
	}
	
	\newglossaryentry{arquitecturared}{
		name={arquitectura de red},
		text={arquitectura de red},
		sort={arquitectura de red},
		plural={arquitecturas de red},
		first={arquitectura de red},
		firstplural={arquitecturas de red},
		description={Estructura organizativa y funcional que define la disposici{\'o}n de los elementos de una \gls{red}, as{\'i} como los principios, modelos y mecanismos utilizados para gestionar el flujo de informaci{\'o}n entre \glspl{dispositivo}. La arquitectura de red especifica la topolog{\'i}a, los protocolos, los niveles de servicio, la gesti{\'o}n de recursos, la seguridad y la forma en que los \glspl{sibiot} se comunican y coordinan de manera eficiente y confiable},
		symbol={🌐},
		see={topologia,protocolocomunicacion,enrutamiento},
		type=\glsdefaulttype
	}
	
	\newglossaryentry{arquitecturareferencia}
	{
		name={Arquitectura de referencia},
		text={arquitectura de referencia},
		sort={arquitectura de referencia},
		plural={arquitecturas de referencia},
		first={arquitectura de referencia},
		firstplural={arquitecturas de referencia},
		description={La \emph{arquitectura de referencia} es un modelo abstracto que define la estructura y los componentes fundamentales de un sistema dentro de un dominio espec{\'i}fico, as{\'i} como sus relaciones. Sirve como gu{\'i}a para dise{\~n}o, desarrollo y evaluaci{\'o}n de sistemas reales, promoviendo la reutilizaci{\'o}n, la \gls{interoperabilidad} y la alineaci{\'o}n con \glspl{buenaspracticas}. No prescribe una implementaci{\'o}n espec{\'i}fica, sino que ofrece una base conceptual com{\'u}n.},
		type=\glsdefaulttype,
		symbol={🏗️},
		user1={Reference architecture},
		user2={Modelo de referencia},
		user3={Esquema arquitect{\'o}nico},
		see={arquitectura,buenaspracticas}
	}
	
	\newglossaryentry{artefacto}{
		name={Artefacto},
		text={artefacto},
		sort={artefacto},
		plural={artefactos},
		first={artefacto},
		firstplural={artefactos},
		description={Elemento tangible o intangible producido durante un proceso de desarrollo, que puede incluir modelos, documentos, c{\'o}digo ejecutable, configuraciones o pruebas. Los artefactos constituyen evidencia estructurada del avance y formalizaci{\'o}n de un sistema.},
		type=main,
		symbol={📦},
		user1={Producto intermedio},
		user2={Resultado de desarrollo},
		user3={Elemento verificable},
		see={artefactoauxiliar,ejecutable}
	}
	
	\newglossaryentry{artefactoauxiliar}{
		name={Artefacto auxiliar},
		text={artefacto auxiliar},
		sort={artefacto auxiliar},
		plural={artefactos auxiliares},
		first={artefacto auxiliar},
		firstplural={artefactos auxiliares},
		description={Elemento complementario dentro de un proceso de desarrollo que apoya la especificaci{\'o}n o construcci{\'o}n del sistema sin constituir el producto ejecutable principal. Puede incluir diagramas, documentaci{\'o}n t{\'e}cnica, prototipos parciales o representaciones intermedias.},
		type=main,
		symbol={📎},
		user1={Soporte metodol{\'o}gico},
		user2={Elemento intermedio},
		user3={Apoyo a especificaci{\'o}n},
		see={modelo,artefacto,especificaciontecnica}
	}
	
	\newglossaryentry{aseguramientocalidad}{
		name={Aseguramiento de la calidad},
		text={aseguramiento de la calidad},
		sort={aseguramiento de la calidad},
		plural={aseguramientos de la calidad},
		first={aseguramiento de la calidad},
		firstplural={aseguramientos de la calidad},
		description={Conjunto de actividades planificadas y sistem{\'a}ticas implementadas en un proceso de desarrollo, producci{\'o}n o servicio, con el fin de proporcionar la confianza de que dicho proceso cumplir{\'a} con los requisitos de \gls{calidad} establecidos. El \emph{aseguramiento de la calidad} se centra en la prevenci{\'o}n de defectos, el dise{\~n}o de procedimientos, la definici{\'o}n de est{\'a}ndares y la mejora continua, diferenci{\'a}ndose del \gls{controlcalidad}, que se orienta a la detecci{\'o}n de fallos. En el contexto de los \glspl{sibiot}, el aseguramiento de la calidad abarca desde la selecci{\'o}n de \glspl{sensor} y \glspl{actuador}, hasta la implementaci{\'o}n de \glspl{aplicacion} y servicios en la nube, garantizando \gls{fiabilidad}, \gls{seguridad}, \gls{interoperabilidad} y sostenibilidad},
		type=main,
		see={calidad,controlcalidad,fiabilidad,seguridad,sostenibilidad,sibiot},
		symbol={🛡️},
		user1={Prevenci{\'o}n de defectos y no conformidades},
		user2={Planificaci{\'o}n y aplicaci{\'o}n de est{\'a}ndares},
		user3={Confianza en procesos y productos IoT}
	}
	
	\newglossaryentry{asincrono}
	{
		name={As{\'i}ncrono},
		text={as{\'i}ncrono},
		sort={asincrono},
		plural={as{\'i}ncronos},
		first={as{\'i}ncrono},
		firstplural={as{\'i}ncronos},
		description={Dos se{\~n}ales son \emph{as{\'i}ncronas} cuando sus instantes significativos no coinciden. En sistemas digitales, el modo as{\'i}ncrono permite que procesos o \glspl{dispositivo} operen de manera independiente sin compartir un mismo reloj, siendo com{\'u}n en protocolos de comunicaci{\'o}n, arquitecturas distribuidas y ejecuci{\'o}n de tareas en paralelo.},
		type=\glsdefaulttype,
		symbol={⏳},
		user1={No simult{\'a}neo},
		user2={Independiente},
		user3={Modo as{\'i}ncrono},
		see={asincrona}
	}
	
	\newglossaryentry{asincrona}
	{
		name={As{\'i}ncrona},
		text={as{\'i}ncrona},
		sort={asincrona},
		plural={as{\'i}ncronas},
		first={as{\'i}ncrona},
		firstplural={as{\'i}ncronas},
		description={Forma femenina de \emph{as{\'i}ncrono}. Se refiere igualmente a operaciones, se{\~n}ales o tareas que no est{\'a}n sincronizadas en el tiempo.},
		type=\glsdefaulttype,
		symbol={⌛},
		user1={No simult{\'a}nea},
		user2={Independiente},
		user3={Operaci{\'o}n as{\'i}ncrona},
		see={asincrono}
	}
	
	\newglossaryentry{assets}{
		name={Assets},
		text={assets},
		sort={assets},
		first={assets},
		description={Recursos de una aplicaci{\'o}n (im{\'a}genes, fuentes, archivos de configuraci{\'o}n, sonidos u otros) empaquetados para su uso en tiempo de ejecuci{\'o}n. En desarrollo m{\'o}vil, se gestionan mediante convenciones de estructura y herramientas del entorno.},
		type=main,
		symbol={📦},
		user1={Recursos de aplicaci{\'o}n},
		user2={Empaquetado},
		user3={Ejecuci{\'o}n},
		see={drawable,layout,values}
	}
	
	\newglossaryentry{ataque}
	{
		name={Ataque},
		text={ataque},
		sort={ataque},
		plural={ataques},
		first={ataque},
		firstplural={ataques},
		description={Un \emph{ataque} es una acci{\'o}n deliberada destinada a explotar vulnerabilidades de un \gls{sistemainformatico} o de una red con el prop{\'o}sito de causar da{\~n}o, alterar su funcionamiento, robar informaci{\'o}n o interrumpir servicios. En el contexto de los \glspl{sibiot}, los ataques pueden afectar \glspl{dispositivoiot}, comprometer la \gls{seguridad} de los datos o interrumpir la \gls{conectividad} de la infraestructura.},
		type=\glsdefaulttype,
		symbol={⚔️},
		user1={Amenaza},
		user2={Intrusi{\'o}n},
		user3={Explotaci{\'o}n de vulnerabilidades},
		see={amenaza,vulnerabilidad,riesgo}
	}
	
	\newglossaryentry{atmega328p}{
		name={ATmega328P},
		text={ATmega328P},
		sort={atmega328p},
		first={ATmega328P},
		description={Microcontrolador de 8 bits desarrollado por Atmel (actualmente Microchip Technology), basado en la arquitectura AVR RISC. El \emph{ATmega328P} integra memoria Flash, SRAM y EEPROM, temporizadores, convertidor anal{\'o}gico--digital de 10 bits, interfaces de comunicaci{\'o}n UART, SPI e I\textsuperscript{2}C, y un conjunto de pines de entrada/salida configurables. Es ampliamente utilizado en placas como Arduino Uno, Arduino Nano, Arduino Pro Mini y Arduino Duemilanove, constituyendo una plataforma de referencia en sistemas embebidos educativos, prototipos y \glspl{sibiot} de baja complejidad.},
		type=main,
		symbol={🧠},
		user1={Microcontrolador AVR},
		user2={Arquitectura RISC de 8 bits},
		user3={Base del ecosistema Arduino cl{\'a}sico},
		see={arduino,arduinonano,arduinopromini,arduinoduemilanove}
	}
	
	\newglossaryentry{auditoria}{
		name={Auditor{\'i}a},
		text={auditor{\'i}a},
		sort={auditoria},
		plural={auditor{\'i}as},
		first={auditor{\'i}a},
		firstplural={auditor{\'i}as},
		description={Proceso sistem{\'a}tico, independiente y documentado de evaluaci{\'o}n que tiene como finalidad verificar el cumplimiento de normas, pol{\'i}ticas, requisitos t{\'e}cnicos o criterios previamente establecidos. En el contexto de sistemas inform{\'a}ticos y \glspl{sibiot}, la auditor{\'i}a permite examinar la integridad, seguridad, trazabilidad y desempe{\~n}o de los procesos, datos e infraestructuras, identificando desviaciones, riesgos o incumplimientos. Puede ser interna o externa, y constituye un mecanismo clave para garantizar gobernanza, transparencia, responsabilidad y mejora continua en entornos tecnol{\'o}gicos complejos.},
		type=main,
		symbol={📝},
		user1={Verificaci{\'o}n independiente},
		user2={Control y cumplimiento normativo},
		user3={Trazabilidad y mejora continua},
		see={verificacion,control,seguridad,protecciondatospersonales,trazabilidad}
	}
	
	\newglossaryentry{autenticacion}{
		name={Autenticaci{\'o}n},
		text={autenticaci{\'o}n},
		sort={autenticacion},
		plural={autenticaciones},
		first={autenticaci{\'o}n},
		firstplural={autenticaciones},
		description={Proceso mediante el cual un \gls{sistema} verifica la identidad declarada de un usuario, dispositivo o componente, asegurando que es quien dice ser. La \emph{autenticaci{\'o}n} se realiza mediante credenciales como contrase{\~n}as, certificados digitales, tokens, biometr{\'i}a (ej.\, \gls{reconocimientofacial} o huellas dactilares) o combinaciones de factores (MFA). En el contexto de los \glspl{sibiot}, garantiza que solo entidades leg{\'i}timas puedan interactuar con \glspl{sensor}, \glspl{actuador} o plataformas, contribuyendo a la \gls{seguridad} y a la protecci{\'o}n de la privacidad},
		type=main,
		see={controlacceso,identificacion,autorizacion,seguridad,privacidad},
		symbol={🔑},
		user1={Verificaci{\'o}n de identidad},
		user2={Uso de credenciales o biometr{\'i}a},
		user3={Pilar del control de acceso}
	}
	
	\newglossaryentry{autenticacionmultifactor}{
		name={Autenticaci{\'o}n multifactor},
		text={autenticaci{\'o}n multifactor},
		sort={autenticacion multifactor},
		plural={autenticaciones multifactor},
		first={autenticaci{\'o}n multifactor},
		firstplural={autenticaciones multifactor},
		description={Mecanismo de seguridad que requiere la verificaci{\'o}n de dos o m{\'a}s factores de autenticaci{\'o}n independientes para validar la identidad de un usuario o entidad. Estos factores se clasifican t{\'i}picamente en: algo que el usuario sabe (por ejemplo, contrase{\~n}a o PIN), algo que el usuario posee (como un token f{\'i}sico, dispositivo m{\'o}vil o certificado digital) y algo que el usuario es (biometr{\'i}a como huella dactilar, reconocimiento facial o iris). La autenticaci{\'o}n multifactor incrementa significativamente la seguridad frente a accesos no autorizados, reduciendo el impacto de credenciales comprometidas. En entornos \gls{iot} y sistemas distribuidos, es fundamental para proteger dispositivos, servicios en la nube y plataformas cr{\'i}ticas, garantizando \gls{confidencialidad}, \gls{integridad} y \gls{disponibilidad}.},
		type=main,
		symbol={🔐},
		user1={Mecanismo avanzado de seguridad},
		user2={Combinaci{\'o}n de m{\'u}ltiples factores},
		user3={Protecci{\'o}n contra accesos no autorizados},
		see={seguridad,controlacceso,identificacion,autorizacion}
	}
	
	\newglossaryentry{autointeligente}{
		name={Auto (coche) inteligente},
		text={auto (coche) inteligente},
		sort={auto inteligente},
		plural={autos (coches) inteligentes},
		first={auto (coche) inteligente},
		firstplural={autos (coches) inteligentes},
		description={Veh{\'i}culo con sensores, actuadores, software y conectividad que habilitan navegaci{\'o}n asistida, funciones ADAS o conducci{\'o}n aut{\'o}noma, optimizaci{\'o}n energ{\'e}tica y comunicaci{\'o}n con otros veh{\'i}culos e infraestructuras},
		type=main,
		symbol={🚗},
		user1={Movilidad inteligente},
		user2={Sistemas ciberf{\'i}sicos},
		user3={V2X/IoT vehicular},
		see={transporteinteligente,objetointeligente,iot}
	}
	
	\newglossaryentry{automatizacion}{
		name={Automatizaci{\'o}n},
		text={automatizaci{\'o}n},
		sort={automatizacion},
		plural={automatizaciones},
		first={automatizaci{\'o}n},
		firstplural={automatizaciones},
		description={Dise{\~n}o e implementaci{\'o}n de sistemas capaces de ejecutar tareas o procesos con intervenci{\'o}n humana m{\'i}nima. Emplea sensores, actuadores, controladores y algoritmos para aumentar eficiencia, reducir errores y mejorar seguridad y productividad},
		type=main,
		symbol={🏭},
		user1={Ingenier{\'i}a de control},
		user2={Operaciones},
		user3={Optimaci{\'o}n de procesos},
		see={automatizacionprocesos,sibiot}
	}
	
	\newglossaryentry{automatizacionanalitica}
	{
		name={Automatizaci{\'o}n anal{\'i}tica},
		text={automatizaci{\'o}n anal{\'i}tica},
		sort={automatizacion analitica},
		plural={automatizaciones anal{\'i}ticas},
		first={automatizaci{\'o}n anal{\'i}tica},
		firstplural={automatizaciones anal{\'i}ticas},
		description={La \emph{automatizaci{\'o}n anal{\'i}tica} es la ejecuci{\'o}n sistem{\'a}tica y programada de procesos de an{\'a}lisis sin intervenci{\'o}n manual constante, utilizando flujos de datos y modelos previamente definidos. En un \gls{sibiot}, esta propiedad permite que el sistema reaccione ante eventos detectados y reduzca latencia decisional en procesos industriales, energ{\'e}ticos o log{\'i}sticos.},
		type=\glsdefaulttype,
		symbol={⚙️},
		user1={Orquestaci{\'o}n de procesos},
		user2={An{\'a}lisis continuo},
		user3={Flujo de datos},
		see={procesamientoanalitico,modelopredictivo,retroalimentacionoperacional}
	}
	
	\newglossaryentry{automatizacionindustrial}{
		name={Automatizaci{\'o}n industrial},
		text={automatizaci{\'o}n industrial},
		sort={automatizacion industrial},
		plural={automatizaciones industriales},
		first={automatizaci{\'o}n industrial},
		firstplural={automatizaciones industriales},
		description={Uso de sistemas de control, computadoras y \glspl{dispositivo} programables para operar y supervisar procesos industriales con intervenci{\'o}n humana m{\'i}nima. En el contexto de los \glspl{sibiot}, la \emph{automatizaci{\'o}n industrial} constituye un antecedente clave al integrar sensores, actuadores y sistemas de informaci{\'o}n para la gesti{\'o}n eficiente de operaciones.},
		type=main,
		symbol={🏭},
		user1={Control autom{\'a}tico},
		user2={Procesos industriales},
		user3={Supervisi{\'o}n y control},
		see={scada,redsensor,integraciondatosoperativos}
	}

	\newglossaryentry{automatizacionprocesos}{
		name={Automatizaci{\'o}n de procesos},
		text={automatizaci{\'o}n de procesos},
		sort={automatizacion de procesos},
		plural={automatizaciones de procesos},
		first={automatizaci{\'o}n de procesos},
		firstplural={automatizaciones de procesos},
		description={Uso de tecnolog{\'i}as digitales y sistemas inteligentes para ejecutar actividades o flujos de trabajo sin intervenci{\'o}n directa, con fines de eficiencia, calidad y trazabilidad},
		type=main,
		symbol={},
		user1={BPM/RPA},
		user2={Optimaci{\'o}n operativa},
		user3={Transformaci{\'o}n digital},
		see={automatizacion,workflow}
	}
	
	\newglossaryentry{autonomia}{
		name={autonom{\'i}a},
		text={autonom{\'i}a},
		sort={autonomia},
		plural={autonom{\'i}as},
		first={autonom{\'i}a},
		firstplural={autonom{\'i}as},
		description={Capacidad de un \gls{dispositivoiot} para operar de manera continua durante un periodo determinado sin necesidad de intervenci{\'o}n externa, usualmente condicionada por el consumo energ{\'e}tico, la eficiencia de los componentes y la duraci{\'o}n de la bater{\'i}a. En los \glspl{sibiot}, la autonom{\'i}a constituye un factor cr{\'i}tico que influye directamente en la confiabilidad, disponibilidad y desempe{\~n}o del sistema, especialmente en escenarios remotos o de dif{\'i}cil acceso},
		symbol={🔋},
		see={bateria,consumoenergetico,gestionenergetica,optimizacion},
		type=main
	}
	
	\newglossaryentry{autorizacion}{
		name={Autorizaci{\'o}n},
		text={autorizaci{\'o}n},
		sort={autorizacion},
		plural={autorizaciones},
		first={autorizaci{\'o}n},
		firstplural={autorizaciones},
		description={Proceso que define y regula los privilegios y permisos que un usuario, dispositivo o componente autenticado tiene dentro de un \gls{sistema}. La \emph{autorizaci{\'o}n} determina qu{\'e} recursos pueden utilizarse, qu{\'e} operaciones pueden ejecutarse y en qu{\'e} condiciones. En el contexto de los \glspl{sibiot}, la autorizaci{\'o}n es esencial para proteger datos sensibles, evitar accesos indebidos a \glspl{actuador} o \glspl{sensor}, y garantizar que las interacciones cumplen con pol{\'i}ticas de \gls{seguridad} y normativas vigentes. Se implementa com{\'u}nmente mediante roles, listas de control de acceso (ACL) o pol{\'i}ticas basadas en atributos (ABAC)},
		type=main,
		symbol={🛡️},
		user1={Definici{\'o}n de permisos y privilegios},
		user2={Protecci{\'o}n de recursos y operaciones},
		user3={Complemento de la autenticaci{\'o}n en control de acceso},
		see={controlacceso,autenticacion,seguridad,privacidad}
	}
	
	\newglossaryentry{aws}
	{
		name={AWS},
		text={AWS},
		sort={aws},
		plural={AWSs},
		first={servicio web de Amazon (\textit{Amazon Web Service}, AWS)},
		firstplural={AWSs},
		description={\emph{AWS} es una colecci{\'o}n de servicios de computaci{\'o}n en la \gls{nube} p{\'u}blica que forman una plataforma integral, ofrecida a trav{\'e}s de Internet por Amazon.com. Incluye servicios de \gls{almacenamiento}, \gls{procesamientodatos}, \gls{seguridad}, inteligencia artificial y despliegue de aplicaciones escalables.},
		type=\glsdefaulttype,
		symbol={☁️},
		user1={Amazon Web Services},
		user2={Plataforma cloud de Amazon},
		user3={Servicios en la nube},
		see={nube,arquitecturadistribuida}
	}
	
	% ===================== Microsoft Azure =====================
	\newglossaryentry{azure}{
		name={Azure},
		text={Azure},
		sort={azure},
		plural={Azure},
		first={Microsoft Azure},
		firstplural={Microsoft Azure},
		description={\emph{Microsoft Azure} es una plataforma de \gls{cloudcomputing} desarrollada por Microsoft que ofrece m{\'a}s de 200 servicios para infraestructura, desarrollo de aplicaciones, inteligencia artificial, redes y \gls{seguridad}. En entornos \gls{iot}, Azure proporciona herramientas como \textit{Azure IoT Hub} y \textit{Azure Digital Twins}, que permiten gestionar millones de \glspl{dispositivoiot}, integrar \glspl{protocolocomunicacion}, realizar an{\'a}lisis en \gls{tiemporeal} y garantizar la escalabilidad y resiliencia de soluciones empresariales.},
		type=main,
		symbol={🔷},
		user1={Plataforma de nube de Microsoft},
		user2={Amplio portafolio de servicios para IoT},
		user3={Integraci{\'o}n con ecosistemas empresariales},
		see={aws,gcp,servicionube,iot}
	}
	
	\newglossaryentry{back}
	{
		name={Back},
		text={back},
		sort={back},
		plural={back},
		first={back},
		firstplural={back},
		description={\emph{Back} es un t{\'e}rmino abreviado que remite al concepto de \gls{backend}.},
		type=\glsdefaulttype,
		symbol={↩️},
		user1={Abreviatura de back-end},
		user2={Parte interna},
		user3={Servidor},
		see={backend}
	}
	\newglossaryentry{backend}
	{
		name={Back-end},
		text={back-end},
		sort={backend},
		plural={back-ends},
		first={back-end},
		firstplural={back-ends},
		description={El \emph{back-end} es la parte de un sistema de software que gestiona la \gls{logicanegocios}, el acceso a \glspl{basedatos}, la \gls{seguridad} y la integraci{\'o}n con otros servicios. Es responsable del \gls{procesamientodatos} y del funcionamiento interno de una aplicaci{\'o}n. En arquitecturas \glspl{sibiot}, el back-end recibe informaci{\'o}n desde los \glspl{dispositivoiot}, la almacena, la analiza y provee servicios a \glspl{interfaz} de usuario o sistemas externos mediante APIs.},
		type=\glsdefaulttype,
		symbol={🖥️},
		user1={Servidor de aplicaci{\'o}n},
		user2={Parte interna},
		user3={Lado del servidor},
		see={frontend,api,arquitecturacapas}
	}
	
	\newglossaryentry{bajapotencia}
	{
		name={Baja potencia},
		text={baja potencia},
		sort={baja potencia},
		plural={dispositivos de baja potencia},
		first={baja potencia},
		firstplural={dispositivos de baja potencia},
		description={La \emph{baja potencia} es una caracter{\'i}stica de dise{\~n}o en sistemas electr{\'o}nicos o \glspl{dispositivoiot} que operan con un consumo energ{\'e}tico m{\'i}nimo, optimizando el uso de bater{\'i}as o energ{\'i}a. Es esencial en \glspl{sensor} inal{\'a}mbricos, dispositivos port{\'a}tiles, redes \gls{lpwan} o sistemas embebidos, ya que prolonga la vida {\'u}til, reduce costos operativos y minimiza el impacto ambiental.},
		type=\glsdefaulttype,
		symbol={🔋},
		user1={Low power},
		user2={Consumo eficiente},
		user3={Dise{\~n}o energ{\'e}ticamente optimizado},
		see={eficienciaenergetica,ble,lpwan,iot}
	}
	
	\newglossaryentry{balanceocarga}{
		name={Balanceo de carga},
		text={balanceo de carga},
		sort={balanceo de carga},
		plural={mecanismos de balanceo de carga},
		first={balanceo de carga},
		firstplural={mecanismos de balanceo de carga},
		description={El \emph{balanceo de carga} es la t{\'e}cnica mediante la cual se distribuyen de forma equitativa las tareas, peticiones o flujos de datos entre m{\'u}ltiples recursos disponibles, como servidores, \glspl{nodo} o instancias de servicio. Su objetivo es optimizar el rendimiento, maximizar la utilizaci{\'o}n de recursos y evitar la sobrecarga de un componente espec{\'i}fico. En \glspl{arquitecturadistribuida} y sistemas \gls{iot}, el balanceo de carga resulta fundamental para garantizar \gls{escalabilidad}, \gls{disponibilidad} y \gls{fiabilidad}, apoyando el procesamiento en paralelo y la redundancia. Puede implementarse mediante hardware dedicado, software especializado o algoritmos como round-robin, least-connections o balanceo basado en latencia.},
		type=main,
		symbol={⚖️},
		user1={Distribuci{\'o}n de peticiones entre recursos},
		user2={Optimiza rendimiento y disponibilidad},
		user3={Soporte clave en sistemas IoT distribuidos},
		see={escalabilidad,arquitecturadistribuida,fiabilidad,disponibilidad,sibiot}
	}
	
	\newglossaryentry{barrera}{
		name={Barrera},
		text={barrera},
		sort={barrera},
		plural={barreras},
		first={barrera},
		firstplural={barreras},
		description={Mecanismo de protecci{\'o}n dise{\~n}ado para reducir la probabilidad o el impacto de un riesgo. Una \emph{barrera} puede ser f{\'i}sica (como muros, cerraduras o jaulas de Faraday), l{\'o}gica (como \glspl{cortafuego}, cifrado, autenticaci{\'o}n) u organizativa (como pol{\'i}ticas de seguridad o capacitaci{\'o}n). En sistemas \gls{iot} y \glspl{sibiot}, las barreras son esenciales para garantizar la seguridad, la fiabilidad y la resiliencia de la infraestructura tecnol{\'o}gica},
		type=main,
		symbol={🛡️},
		user1={Medida de protecci{\'o}n},
		user2={Defensa preventiva},
		user3={Mecanismo de seguridad},
		see={riesgo,seguridad,cortafuego}
	}
	
	\newglossaryentry{barrerafisica}{
		name={Barrera f{\'i}sica},
		text={barrera f{\'i}sica},
		sort={barrera fisica},
		plural={barreras f{\'i}sicas},
		first={barrera f{\'i}sica},
		firstplural={barreras f{\'i}sicas},
		description={Dispositivos o estructuras tangibles utilizadas para impedir, retrasar o limitar accesos no autorizados y reducir riesgos. Ejemplos de \emph{barreras f{\'i}sicas} incluyen cerraduras, muros, rejas, puertas blindadas, jaulas de Faraday, sistemas biom{\'e}tricos y controles de entrada. En el contexto de sistemas \gls{iot} y \glspl{sibiot}, complementan las medidas l{\'o}gicas de seguridad al proteger el hardware frente a manipulaciones, sabotajes o intrusiones externas},
		type=main,
		symbol={🏰},
		user1={Defensa material},
		user2={Protecci{\'o}n estructural},
		user3={Control de acceso f{\'i}sico},
		see={barrera,seguridad,controlacceso}
	}
	
	\newglossaryentry{basedatos}{
		name={Base de datos},
		text={base de datos},
		sort={base de datos},
		plural={bases de datos},
		first={base de datos},
		firstplural={bases de datos},
		description={Conjunto organizado de datos almacenados de manera estructurada o no estructurada, administrados mediante un sistema de gesti{\'o}n (DBMS) que permite su definici{\'o}n, manipulaci{\'o}n, consulta y mantenimiento. Una \emph{base de datos} puede adoptar diferentes modelos (relacional, documental, grafos, clave-valor, etc.) seg{\'u}n el tipo de informaci{\'o}n y los requisitos del sistema. En el contexto de los \glspl{sibiot}, las bases de datos almacenan y procesan grandes vol{\'u}menes de datos generados por \glspl{sensor}, \glspl{actuador} y \glspl{dispositivointeligente}, garantizando \gls{fiabilidad}, \gls{disponibilidad}, \gls{seguridad} y escalabilidad en arquitecturas distribuidas y servicios en la nube},
		type=main,
		symbol={💾},
		user1={Almacenamiento estructurado o no estructurado},
		user2={Administraci{\'o}n mediante DBMS},
		user3={Soporte para an{\'a}lisis y toma de decisiones en SIbIoT},
		see={basedatosrelacional,basedatosnosql,basedatosdistribuida,sibiot}
	}
	
	\newglossaryentry{basedatosdistribuida}{
		name={Base de datos distribuida},
		text={base de datos distribuida},
		sort={base de datos distribuida},
		plural={bases de datos distribuidas},
		first={base de datos distribuida},
		firstplural={bases de datos distribuidas},
		description={Conjunto de datos gestionados por un sistema de administraci{\'o}n de bases de datos (DBMS) que se encuentra distribuido f{\'i}sicamente en varios nodos interconectados por una red, pero que se presenta a los usuarios como una {\'u}nica base l{\'o}gica. Una \emph{base de datos distribuida} busca garantizar \gls{transparencia} en el acceso, consistencia de la informaci{\'o}n, \gls{fiabilidad} y \gls{disponibilidad} incluso ante fallos parciales. En el contexto de los \glspl{sibiot}, estas bases permiten almacenar y procesar grandes vol{\'u}menes de datos generados por \glspl{sensor} y \glspl{dispositivointeligente} en distintos entornos, favoreciendo la escalabilidad, la redundancia y la integraci{\'o}n con servicios en la nube y arquitecturas de arquitectura distribuida},
		type=main,
		see={arquitecturadistribuida,fiabilidad,disponibilidad,sibiot},
		symbol={🗄️},
		user1={Almacenamiento distribuido de datos},
		user2={Transparencia y consistencia},
		user3={Soporte a SIbIoT y arquitecturas escalables},
	}
	
	\newglossaryentry{basedatoslocal}{
		name={base de datos local},
		text={base de datos local},
		plural={bases de datos locales},
		description={Sistema de gesti{\'o}n y persistencia de informaci{\'o}n que se ejecuta de forma interna en un dispositivo o entorno cerrado, sin dependencia de servicios remotos, infraestructura en la nube o conexiones externas. Este tipo de almacenamiento ofrece baja latencia, autonom{\'i}a operativa, funcionamiento sin conectividad, aislamiento de datos y control total de la configuraci{\'o}n. Se emplea habitualmente en dispositivos embebidos, sistemas \gls{iot}, aplicaciones m{\'o}viles y soluciones que requieren respuesta inmediata, alta disponibilidad o restricciones de privacidad},
		user1={almacenamiento y persistencia de datos},
		user2={SQLite, LevelDB, RocksDB},
		symbol={},
		sort={basedatoslocal},
		see={basedatosremota,basedatosdistribuida}
	}
	
	\newglossaryentry{basedatosrelacional}{
		name={Base de datos relacional},
		text={base de datos relacional},
		sort={base de datos relacional},
		plural={bases de datos relacionales},
		first={base de datos relacional},
		firstplural={bases de datos relacionales},
		description={Modelo de gesti{\'o}n de datos basado en tablas (relaciones) formadas por filas (tuplas) y columnas (atributos), donde las asociaciones entre entidades se expresan mediante claves primarias y for{\'a}neas. Una \emph{base de datos relacional} utiliza SQL como lenguaje est{\'a}ndar para definir, manipular y consultar datos, asegurando propiedades ACID (atomicidad, consistencia, aislamiento y durabilidad). En el contexto de los \glspl{sibiot}, se emplea para gestionar datos estructurados y cr{\'i}ticos que requieren integridad referencial, como inventarios, configuraciones o historiales de transacciones},
		type=main,
		symbol={🗄️},
		user1={Modelo tabular de datos},
		user2={Soporte de SQL y propiedades ACID},
		user3={Uso en sistemas con integridad referencial},
		see={basedatosdistribuida,basedatosnosql,sibiot}
	}
	
	\newglossaryentry{basedatosremota}{
		name={Base de datos remota},
		text={base de datos remota},
		sort={basedatosremota},
		plural={bases de datos remotas},
		first={base de datos remota},
		firstplural={bases de datos remotas},
		description={Sistema de almacenamiento y gesti{\'o}n de datos que se encuentra ubicado en un servidor externo al dispositivo o aplicaci{\'o}n cliente, accesible a trav{\'e}s de una red de comunicaci{\'o}n. En los \glspl{sibiot}, las bases de datos remotas permiten centralizar informaci{\'o}n proveniente de nodos distribuidos, facilitar el acceso concurrente, soportar an{\'a}lisis avanzados y garantizar persistencia, disponibilidad y escalabilidad de los datos recolectados.},
		type=main,
		symbol={🗄️},
		user1={Almacenamiento externo},
		user2={Acceso distribuido},
		user3={Persistencia de datos},
		see={basedatos,sgbd,redcomunicacion,cloudcomputing}
	}
	
	\newglossaryentry{basedatosnosql}{
		name={Base de datos NoSQL},
		text={base de datos NoSQL},
		sort={base de datos nosql},
		plural={bases de datos NoSQL},
		first={base de datos NoSQL},
		firstplural={bases de datos NoSQL},
		description={Modelo de gesti{\'o}n de datos que prescinde del esquema tabular de las bases de datos relacionales, empleando estructuras flexibles como documentos, grafos, pares clave-valor o familias de columnas. Una \emph{base de datos NoSQL} est{\'a} orientada a la escalabilidad horizontal, al manejo de grandes vol{\'u}menes de datos distribuidos y a la tolerancia a fallos, aunque puede relajar las propiedades ACID en favor de consistencia eventual (teorema CAP). En el contexto de los \glspl{sibiot}, se utilizan para almacenar datos masivos y heterog{\'e}neos generados por \glspl{sensor}, registros en \gls{tiemporeal} o flujos de eventos en arquitecturas distribuidas},
		type=main,
		symbol={🗃️},
		user1={Modelo flexible y no tabular},
		user2={Consistencia eventual y escalabilidad},
		user3={Uso en sistemas distribuidos e IoT},
		see={basedatosrelacional,basedatosdistribuida,sibiot}
	}
	
	\newglossaryentry{basedatosnube}{
		name={Base de datos en la nube},
		text={base de datos en la nube},
		sort={base de datos en la nube},
		plural={bases de datos en la nube},
		first={base de datos en la nube},
		firstplural={bases de datos en la nube},
		description={Servicio de \gls{basedatos} ofrecido y gestionado en plataformas de \gls{nube}, que permite almacenar, procesar y acceder a la informaci{\'o}n de forma escalable, flexible y con disponibilidad global. Una \emph{base de datos en la nube} elimina la necesidad de administraci{\'o}n local de infraestructura, ofreciendo caracter{\'i}sticas como replicaci{\'o}n autom{\'a}tica, \gls{toleranciafallos}, \gls{escalabilidad} bajo demanda y seguridad integrada. En el contexto de los \glspl{sibiot}, estas bases permiten gestionar grandes vol{\'u}menes de datos generados por \glspl{sensor} y \glspl{actuador}, integrarse con \glspl{serviciodistribuido}, soportar an{\'a}lisis en \gls{tiemporeal} y habilitar aplicaciones cr{\'i}ticas en salud, industria, log{\'i}stica y ciudades inteligentes},
		type=main,
		symbol={☁️💾},
		user1={Almacenamiento y gesti{\'o}n de datos en la nube},
		user2={Escalabilidad y disponibilidad global},
		user3={Soporte a sistemas distribuidos e IoT},
		see={basedatos,basedatosdistribuida,nube,serviciodistribuido,sibiot}
	}
	
	\newglossaryentry{basedatostiemporeal}{
		name={Base de datos en tiempo real},
		text={base de datos en tiempo real},
		sort={base de datos en tiempo real},
		plural={bases de datos en tiempo real},
		first={base de datos en tiempo real},
		firstplural={bases de datos en tiempo real},
		description={Tipo de \gls{basedatos} dise{\~n}ada para gestionar, procesar y responder a operaciones de lectura y escritura con \glspl{latencia} extremadamente bajas, garantizando resultados inmediatos o dentro de plazos cr{\'i}ticos. Una \emph{base de datos en tiempo real} soporta aplicaciones que requieren alta velocidad, coherencia y disponibilidad continua, como la monitorizaci{\'o}n de procesos industriales, sistemas financieros o aplicaciones \gls{iot}. En el contexto de los \glspl{sibiot}, estas \glspl{basedatos} permiten almacenar y analizar datos generados por \glspl{sensor} y \glspl{actuador} en \gls{tiemporeal}, facilitando la \gls{tomadecisiones} autom{\'a}tica, la \gls{eficienciaoperativa} y la \gls{fiabilidad} de los sistemas inteligentes},
		type=main,
		symbol={⏱️💾},
		user1={Procesamiento inmediato de datos},
		user2={Alta disponibilidad y baja latencia},
		user3={Soporte a aplicaciones cr{\'i}ticas en IoT},
		see={basedatos,basedatosdistribuida,tiemporeal,eficienciaoperativa,fiabilidad,sibiot}
	}
	
	\newglossaryentry{bateria}{
		name={Bater{\'i}a},
		text={bater{\'i}a},
		sort={bateria},
		plural={bater{\'i}as},
		first={bater{\'i}a},
		firstplural={bater{\'i}as},
		description={Fuente de energ{\'i}a electroqu{\'i}mica para alimentar dispositivos, especialmente nodos IoT de bajo consumo. Su capacidad y perfil de descarga condicionan autonom{\'i}a, dise{\~n}o de muestreo y estrategias de ahorro energ{\'e}tico.},
		type=main,
		symbol={🔋},
		user1={Autonom{\'i}a},
		user2={Energ{\'i}a portable},
		user3={Bajo consumo},
		see={energia,consumoenergetico}
	}
	
	\newglossaryentry{beaglebone}
	{
		name={BeagleBone},
		text={BeagleBone},
		sort={beaglebone},
		plural={BeagleBones},
		first={BeagleBone},
		firstplural={BeagleBones},
		description={\emph{BeagleBone} es una \gls{placadesarrollo} de c{\'o}digo abierto basada en procesadores ARM, utilizada en aplicaciones embebidas y proyectos de \gls{iot}. Ofrece \glspl{puerto} \gls{gpio}, soporte para sistemas operativos Linux y diversas opciones de \gls{conectividad}.},
		type=\glsdefaulttype,
		symbol={🦴},
		user1={Placa BeagleBone},
		user2={BeagleBone Black},
		user3={BeagleBone AI},
		see={arduino,microcontrolador,placadesarrollo}
	}
	
	\newglossaryentry{bi}{
		name={Business Intelligence},
		text={BI},
		sort={business intelligence},
		plural={BIs},
		first={inteligencia de negocios (\textit{Business Intelligence}, BI)},
		firstplural={BIs},
		description={Conjunto de estrategias, metodolog{\'i}as, procesos, aplicaciones y tecnolog{\'i}as que permiten recopilar, depurar, transformar y analizar datos con el objetivo de apoyar la toma de decisiones en una organizaci{\'o}n. El \emph{BI} integra herramientas de \gls{etl}, \gls{olap}, visualizaci{\'o}n de datos y \gls{modelopredictivo}, facilitando la generaci{\'o}n de informes, cuadros de mando y an{\'a}lisis interactivos. En el contexto de los \glspl{sibiot}, el \gls{bi} ayuda a transformar datos masivos en informaci{\'o}n procesable que mejora la eficiencia, la productividad y la innovaci{\'o}n en sectores como salud, agricultura, industria y ciudades inteligentes.},
		type=main,
		symbol={📊},
		user1={An{\'a}lisis de datos para la toma de decisiones},
		user2={Integraci{\'o}n de ETL, OLAP y visualizaci{\'o}n},
		user3={Aplicaci{\'o}n en entornos corporativos e IoT},
		see={etl,olap,modelopredictivo,bigdata,sibiot}
	}
	
	\newglossaryentry{bigdata}
	{
		name={Big Data},
		text={big data},
		sort={big data},
		plural={Big Data},
		first={Big Data},
		firstplural={Big Data},
		description={\emph{Big Data} es el conjunto de t{\'e}cnicas, herramientas y tecnolog{\'i}as destinadas al procesamiento, almacenamiento y an{\'a}lisis de grandes vol{\'u}menes de datos que no pueden ser gestionados eficazmente mediante sistemas tradicionales. En el contexto de los \glspl{sibiot}, permite extraer informaci{\'o}n relevante a partir de flujos masivos de datos generados por \glspl{sensor} y \glspl{dispositivoconectado}, facilitando la \gls{tomadecisiones}, la predicci{\'o}n de eventos y la automatizaci{\'o}n inteligente de procesos.},
		type=\glsdefaulttype,
		symbol={📊},
		user1={Macrodatos},
		user2={Datos masivos},
		user3={An{\'a}lisis de datos a gran escala},
		see={analisisdatos,hadoop,spark,procesamientodatos}
	}
	
	\newglossaryentry{bigquery}{
		name={BigQuery},
		text={BigQuery},
		sort={bigquery},
		first={BigQuery},
		description={Servicio de an{\'a}lisis de datos a gran escala basado en un motor de \textit{data warehouse} distribuido y orientado a consultas anal{\'i}ticas masivas. \emph{BigQuery} permite ejecutar consultas SQL altamente paralelas sobre grandes vol{\'u}menes de datos estructurados y semiestructurados, sin necesidad de gestionar infraestructura subyacente. En arquitecturas \glspl{sibiot}, BigQuery se utiliza como repositorio anal{\'i}tico para el almacenamiento hist{\'o}rico de eventos IoT, transacciones, telemetr{\'i}a y datos agregados, soportando an{\'a}lisis descriptivo, diagn{\'o}stico y predictivo, as{\'i} como procesos de inteligencia operativa y toma de decisiones basada en datos},
		type=main,
		symbol={📊},
		see={datawarehouse,analiticadatos,procesamientodistribuido,olap,sibiot,cloudcomputing}
	}
	
	% ===================== Bigtable =====================
	\newglossaryentry{bigtable}{
		name={Bigtable},
		text={Bigtable},
		sort={bigtable},
		plural={Bigtable},
		first={Bigtable},
		firstplural={Bigtable},
		description={\emph{Bigtable} es una base de datos NoSQL distribuida de Google, orientada a columnas, capaz de gestionar petabytes de informaci{\'o}n en miles de m{\'a}quinas. Sirve de base para servicios como \gls{hbase} y Google Search. En entornos \gls{iot}, \emph{Bigtable} es ideal para almacenar series temporales, datos de sensores y flujos masivos de \gls{bigdata}, ofreciendo escalabilidad horizontal, baja \gls{latencia} y alta \gls{fiabilidad}.},
		type=main,
		symbol={📊},
		user1={Base de datos NoSQL orientada a columnas},
		user2={Escalabilidad horizontal para petabytes de datos},
		user3={Soporte para series temporales e IoT},
		see={hbase,bigdata,basedatosnosql,sibiot}
	}
	
	\newglossaryentry{bios}
	{
		name={BIOS},
		text={BIOS},
		sort={bios},
		plural={BIOSs},
		first={sistema b{\'a}sico de entrada/salida (\textit{Basic Input/Output System}, BIOS)},
		firstplural={BIOSs},
		description={\emph{BIOS} es un firmware almacenado en una memoria no vol{\'a}til que se encarga de iniciar el \gls{hardware} de un sistema al encenderlo y de proporcionar una interfaz b{\'a}sica entre el sistema operativo y los \glspl{dispositivo} f{\'i}sicos. Aunque ha sido reemplazado progresivamente por UEFI en computadoras modernas, sigue presente en sistemas embebidos y \glspl{dispositivoiot}, donde su simplicidad y fiabilidad resultan ventajosas.},
		type=\glsdefaulttype,
		symbol={⚙️},
		user1={Firmware de arranque},
		user2={Sistema b{\'a}sico de entrada/salida},
		user3={Bootloader},
		see={uefi,hardware}
	}
	
	\newglossaryentry{bluetooth}
	{
		name={Bluetooth},
		text={Bluetooth},
		sort={bluetooth},
		plural={Bluetooth},
		first={Bluetooth},
		firstplural={Bluetooth},
		description={\emph{Bluetooth} es una tecnolog{\'i}a de comunicaci{\'o}n inal{\'a}mbrica de corto alcance utilizada para la transmisi{\'o}n de datos entre \glspl{dispositivo}. Opera en la banda de 2.4 GHz y es com{\'u}n en tel{\'e}fonos m{\'o}viles, perif{\'e}ricos de computadoras, auriculares y sistemas \gls{iot}. Fue desarrollada por el grupo \textit{Bluetooth Special Interest Group (SIG)} y es ampliamente usada en aplicaciones de bajo consumo y conectividad continua.},
		type=\glsdefaulttype,
		symbol={🔵},
		user1={Tecnolog{\'i}a inal{\'a}mbrica},
		user2={BT},
		user3={Conectividad de corto alcance},
		see={ble,protocolocomunicacion,conectividad}
	}
	
	\newglossaryentry{ble}
	{
		name={BLE},
		text={BLE},
		sort={ble},
		plural={BLEs},
		first={Bluetooth de bajo consumo energ{\'e}tico (\textit{Bluetooth Low Energy}, BLE)},
		firstplural={BLEs},
		description={\emph{BLE} es una tecnolog{\'i}a de comunicaci{\'o}n inal{\'a}mbrica dise{\~n}ada para aplicaciones de bajo consumo energ{\'e}tico. Utiliza la misma banda de frecuencia que \gls{bluetooth} cl{\'a}sico (2.4 GHz), pero est{\'a} optimizada para mantener \gls{conectividad} intermitente, transmitir peque{\~n}os paquetes de \glspl{dato} y extender la duraci{\'o}n de bater{\'i}a en \glspl{dispositivoiot}, \glspl{sensor} y dispositivos port{\'a}tiles.},
		type=\glsdefaulttype,
		symbol={🔵},
		user1={Bluetooth Low Energy},
		user2={Bluetooth LE},
		user3={Bluetooth de baja energ{\'i}a},
		see={bluetooth}
	}
	
	\newglossaryentry{borde}{
		name={Borde},
		text={borde},
		sort={borde},
		plural={computaciones de borde},
		first={compoutaci{\'o}n de borde (\textit{Edge Computing})},
		firstplural={computaciones de borde},
		description={El borde se refiere al entorno computacional ubicado pr{\'o}ximo a la fuente de generaci{\'o}n de datos, como \glspl{sensor} o \glspl{dispositivo} de campo. En el contexto de los \glspl{sibiot}, el procesamiento en el borde permite ejecutar tareas de an{\'a}lisis, filtrado y respuesta local antes de transmitir los datos hacia la \gls{nube}, reduciendo la latencia, mejorando la seguridad y optimizando el uso del ancho de banda.},
		type=main,
		symbol={📡},
		user1={Edge computing},
		user2={Computaci{\'o}n en el borde},
		user3={Procesamiento local cercano a la fuente de datos},
		user4={Arquitectura distribuida entre el borde y la nube},
		see={nube,edgecomputing,cloudcomputing,sibiot}
	}
	
	\newglossaryentry{browser}
	{
		name={Browser},
		text={browser},
		sort={browser},
		plural={browsers},
		first={navegador o \emph{browser}},
		firstplural={browsers},
		description={Un \emph{browser} es el t{\'e}rmino en ingl{\'e}s que identifica a un navegador web. Es un software, programa o aplicaci{\'o}n que ofrece al usuario acceso a la \gls{www} y permite interactuar con p{\'a}ginas y servicios en l{\'i}nea.},
		type=\glsdefaulttype,
		symbol={🌐},
		user1={Navegador web},
		user2={Explorador de Internet},
		user3={Cliente HTTP},
		see={aplicacionweb,http,www}
	}
	
	\newglossaryentry{bson}
	{
		name={BSON},
		text={BSON},
		sort={bson},
		plural={BSONs},
		first={BSON},
		firstplural={BSONs},
		description={\emph{BSON} es un formato binario de serializaci{\'o}n de documentos basado en JSON, dise{\~n}ado para el almacenamiento y transmisi{\'o}n eficiente en bases de datos orientadas a documentos, como \gls{mongodb}. Combina la flexibilidad de JSON con la eficiencia del c{\'o}digo binario para representar estructuras complejas de datos.},
		type=\glsdefaulttype,
		symbol={📄},
		user1={Binary JSON},
		user2={Formato binario de documentos},
		user3={Serializaci{\'o}n eficiente},
		see={json,mongodb}
	}
	
	\newglossaryentry{buenaspracticas}
	{
		name={Buenas pr{\'a}cticas},
		text={buenas pr{\'a}cticas},
		sort={buenas practicas},
		plural={buenas pr{\'a}cticas},
		first={buenas pr{\'a}cticas},
		firstplural={buenas pr{\'a}cticas},
		description={Las \emph{buenas pr{\'a}cticas} son un conjunto de metodolog{\'i}as, normas, recomendaciones y procedimientos validados que se aplican de manera sistem{\'a}tica para lograr resultados eficaces, seguros y sostenibles en un dominio espec{\'i}fico. En el desarrollo de \glspl{sibiot}, comprenden aspectos como la \gls{seguridad}, la \gls{interoperabilidad}, el consumo eficiente de recursos, la trazabilidad de los datos y el mantenimiento del \gls{codigo} y de los \glspl{dispositivo}.},
		type=\glsdefaulttype,
		symbol={✅},
		user1={Buenas pr{\'a}cticas de desarrollo},
		user2={Recomendaciones t{\'e}cnicas},
		user3={Metodolog{\'i}as validadas},
		see={arquitecturareferencia,interoperabilidad,seguridad}
	}
	
	\newglossaryentry{burndownchart}{
		name={Burndown Chart},
		text={\textit{Burndown Chart}},
		sort={burndown chart},
		plural={\textit{Burndown Charts}},
		first={\textit{Burndown Chart}},
		firstplural={\textit{Burndown Charts}},
		description={Gr{\'a}fico que muestra la cantidad de trabajo pendiente a lo largo del tiempo dentro de un Sprint o una versi{\'o}n del producto. Permite visualizar el ritmo de avance y detectar desviaciones respecto al plan previsto.},
		type=main,
		symbol={📉},
		user1={Seguimiento del trabajo restante},
		user2={Visualizaci{\'o}n del progreso},
		user3={Control de ritmo del Sprint},
		see={sprint,sprintbacklog,velocitychart}
	}
	
	\newglossaryentry{burnupchart}{
		name={Burnup Chart},
		text={\textit{Burnup Chart}},
		sort={burnup chart},
		plural={\textit{Burnup Charts}},
		first={\textit{Burnup Chart}},
		firstplural={\textit{Burnup Charts}},
		description={Gr{\'a}fico que representa el trabajo completado frente al alcance total definido para el producto o Sprint. Permite observar cambios en el alcance y evoluci{\'o}n del progreso acumulado.},
		type=main,
		symbol={📈},
		user1={Trabajo completado acumulado},
		user2={Visualizaci{\'o}n del alcance},
		user3={Seguimiento del progreso total},
		see={burndownchart,productbacklog,sprint}
	}
	
	\newglossaryentry{cadenalogistica}{
		name={Cadena log{\'i}stica},
		text={cadena log{\'i}stica},
		sort={cadena logistica},
		plural={cadenas log{\'i}sticas},
		first={cadena log{\'i}stica},
		firstplural={cadenas log{\'i}sticas},
		description={Conjunto interconectado de procesos, actores, recursos y tecnolog{\'i}as que participan en el movimiento de bienes y servicios desde el proveedor de materias primas hasta el consumidor final. La \emph{cadena log{\'i}stica} incluye actividades de producci{\'o}n, almacenamiento, transporte, distribuci{\'o}n y devoluciones, integrando tanto flujos f{\'i}sicos como informativos. En el contexto de los \glspl{sibiot}, la cadena log{\'i}stica se optimiza mediante el uso de \glspl{sensor} para trazabilidad, \gls{gnss} para geolocalizaci{\'o}n, \glspl{basedatosdistribuida} para gesti{\'o}n de inventarios y \gls{modelopredictivo} para anticipar la demanda, garantizando \gls{fiabilidad}, \gls{eficienciaoperativa} y sostenibilidad},
		type=main,
		symbol={🔗🚚},
		user1={Conjunto de procesos desde proveedor a cliente},
		user2={Integraci{\'o}n de flujos f{\'i}sicos e informativos},
		user3={Soporte mediante IoT y tecnolog{\'i}as inteligentes},
		see={logistica,gnss,basedatosdistribuida,modelopredictivo,eficienciaoperativa,sostenibilidad,sibiot}
	}
	
	\newglossaryentry{calidad}{
		name={Calidad},
		text={calidad},
		sort={calidad},
		plural={calidades},
		first={calidad},
		firstplural={calidades},
		description={Grado en que un producto, proceso o servicio cumple con requisitos, especificaciones o expectativas establecidas. La \emph{calidad} implica atributos como \gls{fiabilidad}, \gls{seguridad}, \gls{eficiencia} y facilidad de uso. En ingenier{\'i}a de software y en los \glspl{sibiot}, la calidad se mide mediante m{\'e}tricas, pruebas y auditor{\'i}as que eval{\'u}an tanto los \glspl{sensor} y \glspl{actuador} como las \glspl{aplicacion}, garantizando un desempe{\~n}o consistente, una buena \gls{ux} y el cumplimiento de normativas t{\'e}cnicas y legales. La gesti{\'o}n de la calidad integra procesos de \gls{controlcalidad}, \gls{aseguramientocalidad} y mejora continua},
		type=main,
		symbol={✨},
		user1={Cumplimiento de requisitos y expectativas},
		user2={Medida de fiabilidad, seguridad y eficiencia},
		user3={Eje de mejora continua en SIbIoT},
		see={controlcalidad,aseguramientocalidad,fiabilidad,seguridad,eficiencia,sibiot}
	}
	
	\newglossaryentry{calidadaire}{
		name={Calidad del aire},
		text={calidad del aire},
		sort={calidad del aire},
		plural={calidades del aire},
		first={calidad del aire},
		firstplural={calidades del aire},
		description={Medida de la composici{\'o}n y condiciones del aire en un entorno determinado, evaluada en funci{\'o}n de la concentraci{\'o}n de contaminantes, gases y part{\'i}culas, as{\'i} como de par{\'a}metros ambientales como temperatura y humedad. La \emph{calidad del aire} es una \gls{variablefisica} crucial para la salud humana, la seguridad laboral y la sostenibilidad ambiental. En el contexto de los \glspl{sibiot}, se monitoriza mediante \glspl{sensor} especializados que detectan compuestos org{\'a}nicos vol{\'a}tiles (COV), di{\'o}xido de carbono (CO\textsubscript{2}), mon{\'o}xido de carbono (CO), part{\'i}culas finas (PM2.5, PM10) y otros contaminantes. Estos datos permiten activar sistemas de ventilaci{\'o}n o purificaci{\'o}n autom{\'a}ticos, generar alertas en \gls{tiemporeal}, cumplir normativas ambientales y apoyar la \gls{eficienciaoperativa} en entornos industriales, dom{\'o}ticos y urbanos},
		type=main,
		symbol={🌫️},
		user1={Medici{\'o}n de contaminantes y condiciones ambientales},
		user2={Variable clave en salud y sostenibilidad},
		user3={Dato monitorizado por SIbIoT},
		see={variablefisica,sensor,monitorizacion,sostenibilidad,sibiot}
	}
	
	\newglossaryentry{calidaddatos}{
		name={calidad de datos},
		text={calidad de datos},
		sort={calidaddatos},
		plural={calidades de datos},
		first={calidad de datos},
		firstplural={calidades de datos},
		description={La calidad de datos representa un aspecto esencial en la \gls{gestioninformacion} y en los \glspl{sibiot}, pues determina la confiabilidad de los resultados anal{\'i}ticos y de la \gls{tomadecisionesinformadas}. Mantener una alta calidad requiere la aplicaci{\'o}n de procesos de \gls{validacion}, \gls{limpieza} y \gls{enriquecimiento}, adem{\'a}s de reglas de negocio que aseguren la \gls{consistencia} y la \gls{integridad}. Los sistemas con buena calidad de datos permiten mayor eficiencia operativa, precisi{\'o}n en los modelos anal{\'i}ticos y confianza en los informes generados.},
		type=\glsdefaulttype,
		symbol={✅},
		user1={Data quality},
		user2={Conjunto de propiedades que garantizan la fiabilidad de los datos},
		user3={Medida de la precisi{\'o}n, consistencia y utilidad de la informaci{\'o}n almacenada},
		see={validacion,limpieza,enriquecimiento,integridad,consistencia,gestioninformacion}
	}
	
	\newglossaryentry{calidadexperiencia}{
		name={Calidad de experiencia},
		text={QoE},
		sort={calidad experiencia},
		plural={m{\'e}tricas de QoE},
		first={calidad de experiencia (\textit{Quality of Expirence}, QoE)},
		firstplural={m{\'e}tricas de calidad de experiencia},
		description={Percepci{\'o}n subjetiva del usuario sobre el desempe{\~n}o general de un servicio o sistema. En \gls{iot}, la \emph{calidad de experiencia} combina par{\'a}metros t{\'e}cnicos (\gls{calidadservicio}) y factores humanos (\gls{usabilidad}), determinando el grado de satisfacci{\'o}n del usuario final.},
		type=main,
		symbol={⭐},
		user1={Satisfacci{\'o}n del usuario},
		user2={Percepci{\'o}n global de calidad},
		user3={Combinaci{\'o}n de QoS y experiencia humana},
		see={usabilidad,calidadservicio,interfaz}
	}
	
	\newglossaryentry{calidadservicio}{
		name={Calidad de servicio},
		text={QoS},
		sort={calidad servicio},
		plural={par{\'a}metros de calidad de servicio},
		first={calidad de servicio (\textit{Quality of Service}, QoS)},
		firstplural={par{\'a}metros de calidad de servicio},
		description={Conjunto de mecanismos y par{\'a}metros que determinan el nivel de desempe{\~n}o de una red o servicio, incluyendo latencia, disponibilidad, p{\'e}rdida de paquetes y ancho de banda. En sistemas \gls{iot}, la \emph{calidad de servicio} garantiza la entrega oportuna y fiable de datos cr{\'i}ticos para el correcto funcionamiento de las aplicaciones.},
		type=main,
		symbol={📶},
		user1={M{\'e}tricas de rendimiento de red},
		user2={Garant{\'i}a de entrega de datos},
		user3={Optimiza experiencia y fiabilidad},
		see={redinteligente,seguridad,fiabilidad,calidadexperiencia}
	}
	
	\newglossaryentry{calidadsoftware}{
		name={Calidad de software},
		text={calidad de software},
		sort={calidad de software},
		first={calidad de software},
		description={Grado en que un producto software satisface requisitos y atributos de calidad (p.\,ej., mantenibilidad, fiabilidad, seguridad, eficiencia y usabilidad) en un contexto de uso.},
		type=main,
		symbol={✅},
		user1={Atributos de calidad},
		user2={Satisfacci{\'o}n de requisitos},
		user3={Evaluaci{\'o}n},
		see={testabilidad,refactorizacion,documentacion}
	}
	
	\newglossaryentry{capa}{
		name={Capa},
		text={capa},
		sort={capa},
		plural={capas},
		first={capa},
		firstplural={capas},
		description={Unidad o nivel funcional dentro de una arquitectura de sistema o red que agrupa componentes, servicios o procesos con responsabilidades afines. Las capas permiten dividir la complejidad de un sistema en secciones l{\'o}gicas, facilitando la modularidad, la escalabilidad y el mantenimiento. En los \glspl{sibiot}, las arquitecturas suelen organizarse en capas como percepci{\'o}n, red y aplicaci{\'o}n, cada una de las cuales cumple un papel espec{\'i}fico: la adquisici{\'o}n de datos, la comunicaci{\'o}n entre dispositivos y la provisi{\'o}n de servicios inteligentes. Este enfoque por capas tambi{\'e}n se aplica a modelos de referencia como el \gls{osi} o el \gls{iotstack}, que estandarizan la interconexi{\'o}n y el flujo de informaci{\'o}n en entornos distribuidos.},
		type=main,
		user1={Nivel funcional dentro de una arquitectura},
		user2={Estructura modular que organiza componentes o servicios},
		user3={Base para la escalabilidad y la interoperabilidad del sistema},
		user4={Ejemplos: percepci{\'o}n, red, aplicaci{\'o}n en SIbIoT},
		see={arquitectura,osi,iotstack,sibiot,protocolo},
		symbol={\ensuremath{\mathcal{L}}}
	}
	
	\newglossaryentry{capaaplicacion}
	{
		name={Capa de aplicaci{\'o}n},
		text={capa de aplicaci{\'o}n},
		sort={capa de aplicacion},
		plural={capas de aplicaci{\'o}n},
		first={capa de aplicaci{\'o}n},
		firstplural={capas de aplicaci{\'o}n},
		description={La \emph{capa de aplicaci{\'o}n} es el nivel superior en arquitecturas orientadas a servicios, redes o \glspl{sibiot}, responsable de proporcionar funcionalidades espec{\'i}ficas al usuario final e integrar la \gls{logicanegocios} con los datos procesados. Interpreta la informaci{\'o}n proveniente de capas inferiores para su visualizaci{\'o}n, an{\'a}lisis o ejecuci{\'o}n de acciones. Incluye \glspl{interfaz}, \glspl{dashboard}, servicios web, sistemas de notificaci{\'o}n y aplicaciones m{\'o}viles o empresariales.},
		type=\glsdefaulttype,
		symbol={🧩},
		user1={Nivel de aplicaci{\'o}n},
		user2={Capa de servicios},
		user3={Application layer},
		see={capatecnologica}
	}
	
	\newglossaryentry{capacidad}{
		name={Capacidad},
		text={capacidad},
		sort={capacidad},
		plural={capacidades},
		first={capacidad},
		firstplural={capacidades},
		description={Magnitud que representa el potencial, volumen o nivel m{\'a}ximo de recursos que un sistema, dispositivo o infraestructura puede procesar, almacenar, transmitir o gestionar de manera eficaz. En el contexto de los \glspl{sibiot}, la \emph{capacidad} abarca dimensiones como el ancho de banda disponible, el rendimiento de c{\'o}mputo, la potencia de procesamiento en nodos de \gls{edgecomputing}, el almacenamiento distribuido y la posibilidad de atender flujos simult{\'a}neos de datos provenientes de \glspl{sensor}. Una adecuada gesti{\'o}n de la capacidad mejora la \gls{eficienciaoperativa}, garantiza la escalabilidad y permite que los modelos de \gls{analitica} y \glspl{modelopredictivo} funcionen de forma continua incluso en escenarios de alta demanda o variabilidad ambiental.},
		type=main,
		symbol={⚙️},
		user1={Recursos m{\'a}ximos de operaci{\'o}n de un sistema},
		user2={Procesamiento y almacenamiento en entornos IoT},
		user3={Base para la escalabilidad y la eficiencia operativa},
		see={rendimiento,eficienciaoperativa,analitica,edgecomputing}
	}
	
	\newglossaryentry{capacidadpredictiva}
	{
		name={Capacidad predictiva},
		text={capacidad predictiva},
		sort={capacidad predictiva},
		plural={capacidades predictivas},
		first={capacidad predictiva},
		firstplural={capacidades predictivas},
		description={La \emph{capacidad predictiva} es la aptitud del sistema para estimar eventos futuros o comportamientos probables mediante modelos estad{\'i}sticos o algoritmos de aprendizaje basados en datos hist{\'o}ricos. En entornos \gls{sibiot}, esta capacidad permite anticipar fallas, variaciones de demanda o fluctuaciones energ{\'e}ticas, fortaleciendo la planificaci{\'o}n operativa y estrat{\'e}gica.},
		type=\glsdefaulttype,
		symbol={🔮},
		user1={Modelado estad{\'i}stico},
		user2={Predicci{\'o}n},
		user3={Aprendizaje supervisado},
		see={modelopredictivo,procesamientoanalitico,analisis}
	}
	
	\newglossaryentry{capacidadprocesamiento}{
		name={Capacidad de procesamiento},
		text={capacidad de procesamiento},
		sort={capacidad de procesamiento},
		type={main},
		plural={capacidades de procesamiento},
		description={Cantidad de operaciones, c{\'a}lculos o tareas que un sistema, dispositivo o microcontrolador puede ejecutar en un intervalo determinado. En los sistemas de informaci{\'o}n basados en IoT, la capacidad de procesamiento determina la habilidad de un dispositivo para analizar datos, ejecutar algoritmos, gestionar comunicaciones y tomar decisiones locales, especialmente en contextos con recursos limitados},
		first={capacidad de procesamiento},
		firstplural={capacidades de procesamiento},
		symbol={🧠},
		user1={Mide el rendimiento computacional de un dispositivo},
		user2={Factor cr{\'i}tico en IoT por las limitaciones de hardware},
		user3={Influye en la ejecuci{\'o}n de algoritmos locales y tareas en tiempo real},
		user4={Puede escalar mediante optimizaci{\'o}n, hardware avanzado o delegaci{\'o}n a la nube},
		see={procesamiento,capacidad,optimizacion,rendimiento}
	}
	
	\newglossaryentry{capaenlace}{
		name={Capa de enlace de datos},
		text={capa de enlace de datos},
		sort={capa enlace de datos},
		plural={capas de enlace de datos},
		first={capa de enlace de datos},
		firstplural={capas de enlace de datos},
		description={Subcapa que define acceso al medio, direccionamiento y control de errores entre nodos adyacentes (p.~ej., IEEE~802.11).},
		type=main,
		symbol={🧶},
		see={wifi,halow,redinalambrica},
		user1={MAC/LLC},
		user2={Control de errores},
		user3={Acceso al medio}
	}
	
	\newglossaryentry{capanegocios}
	{
		name={Capa de negocios},
		text={capa de negocios},
		sort={capa de negocios},
		plural={capas de negocios},
		first={capa de negocios},
		firstplural={capas de negocios},
		description={La \emph{capa de negocios} es el nivel l{\'o}gico de una arquitectura de sistemas encargado de implementar las reglas de negocio, procesos organizacionales y \gls{logicanegocios} central del sistema. Orquesta las operaciones entre la \gls{interfaz} de usuario y los servicios de datos, definiendo el comportamiento del sistema seg{\'u}n objetivos y pol{\'i}ticas de la organizaci{\'o}n. Es esencial para mantener coherencia funcional, reutilizaci{\'o}n de componentes y separaci{\'o}n de responsabilidades.},
		type=\glsdefaulttype,
		symbol={🏢},
		user1={Capa de gesti{\'o}n},
		user2={Business layer},
		user3={Nivel funcional},
		see={logicanegocios}
	}
	
	\newglossaryentry{capapercepcion}
	{
		name={Capa de percepci{\'o}n},
		text={capa de percepci{\'o}n},
		sort={capa de percepcion},
		plural={capas de percepci{\'o}n},
		first={capa de percepci{\'o}n},
		firstplural={capas de percepci{\'o}n},
		description={La \emph{capa de percepci{\'o}n} es el nivel inferior de la arquitectura t{\'i}pica de los \glspl{sibiot}, encargada de la adquisici{\'o}n de datos del entorno f{\'i}sico mediante \glspl{sensor}, dispositivos embebidos o tecnolog{\'i}as de identificaci{\'o}n como \gls{rfid}. Esta capa detecta eventos, condiciones o estados y los transmite a capas superiores para su procesamiento, almacenamiento o an{\'a}lisis.},
		type=\glsdefaulttype,
		symbol={🎛️},
		user1={Perception layer},
		user2={Capa sensora},
		user3={Nivel f{\'i}sico de IoT},
		see={sensor}
	}
	
	\newglossaryentry{capared}
	{
		name={Capa de red},
		text={capa de red},
		sort={capa de red},
		plural={capas de red},
		first={capa de red},
		firstplural={capas de red},
		description={La \emph{capa de red} corresponde al nivel del modelo de referencia OSI responsable del direccionamiento, encaminamiento y entrega de paquetes de datos entre \glspl{dispositivo} en distintas redes. En \glspl{sibiot}, esta capa permite la interconexi{\'o}n eficiente de nodos y el soporte a protocolos como IP, 6LoWPAN o RPL, facilitando la \gls{escalabilidad}, \gls{movilidad} y robustez en la transmisi{\'o}n de datos.},
		type=\glsdefaulttype,
		symbol={🌐},
		user1={Network layer},
		user2={Capa de comunicaci{\'o}n},
		user3={Nivel de interconexi{\'o}n},
		see={sistemacomunicacion}
	}
	
	\newglossaryentry{capasensores}{
		name={Capa de sensores},
		text={capa de sensores},
		sort={capa de sensores},
		plural={capas de sensores},
		first={capa de sensores},
		firstplural={capas de sensores},
		description={Nivel fundamental en la arquitectura de un \gls{sibiot}, encargado de la percepci{\'o}n del entorno mediante la \gls{captura} de \glspl{dato} f{\'i}sicos o l{\'o}gicos. La \emph{capa de sensores} est{\'a} compuesta por \glspl{sensor}, dispositivos de adquisici{\'o}n y elementos de interconexi{\'o}n que permiten transformar \glspl{fenomenofisico} en \glspl{variablefisica} digitales procesables. Su funci{\'o}n principal es proporcionar informaci{\'o}n fiable y en \gls{tiemporeal} para alimentar las capas superiores de procesamiento, an{\'a}lisis y \gls{control}. En el contexto de los sistemas distribuidos, esta capa tambi{\'e}n debe garantizar \gls{precision}, \gls{eficienciaenergetica} y \gls{interoperabilidad} para integrarse en ecosistemas complejos y heterog{\'e}neos},
		type=main,
		symbol={📡},
		user1={Nivel de percepci{\'o}n en SIbIoT},
		user2={Captura de datos mediante sensores},
		user3={Base de la arquitectura IoT},
		see={sensor,variablefisica,fenomenofisico,captura,sibiot}
	}
	
	\newglossaryentry{capatecnologica}
	{
		name={Capa tecnol{\'o}gica},
		text={capa tecnol{\'o}gica},
		sort={capa tecnologica},
		plural={capas tecnol{\'o}gicas},
		first={capa tecnol{\'o}gica},
		firstplural={capas tecnol{\'o}gicas},
		description={La \emph{capa tecnol{\'o}gica} es una representaci{\'o}n abstracta de la arquitectura global de un \gls{sibiot} que integra los distintos componentes tecnol{\'o}gicos involucrados. Incluye el dise{\~n}o de \gls{hardware} (como \glspl{sensor}, \glspl{actuador} y \glspl{placadesarrollo}), as{\'i} como el software que los controla y comunica. En la metodolog{\'i}a TDDM4IoTS, esta capa se define tempranamente para guiar el desarrollo, asegurar la compatibilidad entre tecnolog{\'i}as heterog{\'e}neas y servir como documento de referencia com{\'u}n para todos los equipos de trabajo.},
		type=\glsdefaulttype,
		symbol={🖧},
		user1={Nivel tecnol{\'o}gico},
		user2={Capa de implementaci{\'o}n},
		user3={Arquitectura tecnol{\'o}gica},
		see={sibiot}
	}
	
	
	\newglossaryentry{captura}{
		name={Captura},
		text={captura},
		sort={captura},
		plural={capturas},
		first={captura},
		firstplural={capturas},
		description={Acci{\'o}n o proceso de obtener, registrar y almacenar informaci{\'o}n proveniente de una fuente determinada. En el {\'a}mbito de los \glspl{sibiot}, la \emph{captura} se refiere a la adquisici{\'o}n de datos de \glspl{fenomenofisico} o \glspl{variablefisica} mediante \glspl{sensor}, para su posterior procesamiento, an{\'a}lisis o transmisi{\'o}n. La calidad de la captura est{\'a} determinada por factores como precisi{\'o}n, frecuencia de muestreo, resoluci{\'o}n y fiabilidad del dispositivo. Una captura adecuada de datos es fundamental para garantizar la validez de la \gls{monitorizacion}, el \gls{control} y la \gls{tomadecisiones} en tiempo real dentro de ecosistemas IoT},
		type=main,
		symbol={📥},
		user1={Adquisici{\'o}n de informaci{\'o}n},
		user2={Registro de datos mediante sensores},
		user3={Base de procesos de an{\'a}lisis y control},
		see={sensor,monitorizacion,fenomenofisico,variablefisica,sibiot}
	}
	
	\newglossaryentry{capturadatos}{
		name={captura de datos},
		text={captura de datos},
		sort={capturadatos},
		plural={capturas de datos},
		first={captura de datos},
		firstplural={capturas de datos},
		description={La captura de datos es el punto de inicio en el ciclo de gesti{\'o}n de informaci{\'o}n. Consiste en la obtenci{\'o}n y digitalizaci{\'o}n de informaci{\'o}n procedente del entorno f{\'i}sico o digital mediante \glspl{sensor}, formularios o flujos autom{\'a}ticos. En los \glspl{sibiot}, la captura de datos se realiza frecuentemente en \gls{tiemporeal}, garantizando la disponibilidad inmediata para su procesamiento, an{\'a}lisis o almacenamiento. La precisi{\'o}n y consistencia en esta fase determinan la \gls{calidaddatos} y el valor anal{\'i}tico del sistema.},
		type=\glsdefaulttype,
		symbol={📥},
		user1={Data capture},
		user2={Adquisici{\'o}n o recolecci{\'o}n de informaci{\'o}n desde fuentes diversas},
		user3={Obtenci{\'o}n automatizada de datos en sistemas conectados},
		see={tiemporeal,procesamiento,calidaddatos,analisistreaming}
	}
	
	\newglossaryentry{carga}
	{
		name={Carga},
		text={carga},
		sort={carga},
		plural={cargas},
		first={Carga},
		description={Magnitud o cantidad que representa la acci{\'o}n de aplicar una fuerza, peso o energ{\'i}a sobre un sistema f{\'i}sico, el{\'e}ctrico o inform{\'a}tico. 
			En f{\'i}sica y mec{\'a}nica, la carga se asocia con la fuerza o el peso aplicado sobre una superficie o estructura, originando \gls{deformacion} o \gls{tension}. 
			En el contexto el{\'e}ctrico, se refiere a la propiedad de la materia que origina interacciones electromagn{\'e}ticas, medida en culombios (C), o tambi{\'e}n al consumo de energ{\'i}a de un circuito. 
			En ingenier{\'i}a de materiales e industrial, la carga define el esfuerzo que soporta un componente, siendo esencial en la calibraci{\'o}n de \glspl{celdacarga}. 
			Finalmente, en el campo inform{\'a}tico y de los \glspl{sibiot}, el t{\'e}rmino carga se aplica al nivel de procesamiento o demanda de recursos que un \gls{sistema} o \gls{dispositivo} debe gestionar, tanto en t{\'e}rminos de datos como de tareas concurrentes.},
		symbol={Q},
		see=[Relacionado con:]{deformacion,esfuerzo,masa,energia,procesamiento}
	}
	
	\newglossaryentry{casouso}{
		name={Caso de uso},
		text={caso de uso},
		sort={caso de uso},
		plural={casos de uso},
		first={caso de uso},
		firstplural={casos de uso},
		description={T{\'e}cnica utilizada en el an{\'a}lisis y dise{\~n}o de sistemas para describir las interacciones entre un usuario (o \emph{actor}) y el sistema con el fin de lograr un objetivo espec{\'i}fico. Cada \emph{caso de uso} representa una funci{\'o}n o servicio que el sistema proporciona, definiendo su flujo principal y alternativo de eventos. En ingenier{\'i}a de software y en el desarrollo de \glspl{sibiot}, los casos de uso facilitan la comprensi{\'o}n de requisitos funcionales, la validaci{\'o}n de comportamientos esperados y la comunicaci{\'o}n entre usuarios, analistas y desarrolladores.},
		type=main,
		symbol={📘},
		user1={Describe interacciones sistema-usuario},
		user2={Identifica funciones del sistema},
		user3={Base para requisitos funcionales},
		see={requisito,actor,diagrama,analisis}
	}
	
	\newglossaryentry{cassandra}{
		name={Apache Cassandra},
		text={Apache Cassandra},
		sort={cassandra},
		plural={Apache Cassandra},
		first={Apache Cassandra},
		firstplural={Apache Cassandra},
		description={Sistema de gesti{\'o}n de \glspl{basedatosnosql} distribuido y altamente escalable, dise{\~n}ado para manejar grandes vol{\'u}menes de datos en m{\'u}ltiples \glspl{nodo} sin un {\'u}nico punto de fallo. \emph{Cassandra}, desarrollado inicialmente por Facebook y mantenido por la Apache Software Foundation, utiliza un modelo de almacenamiento basado en familias de columnas y ofrece consistencia configurable seg{\'u}n el teorema CAP. En el contexto de los \glspl{sibiot}, Cassandra es id{\'o}nea para aplicaciones que requieren \gls{disponibilidad}, \gls{toleranciafallos} y rendimiento en tiempo real, como plataformas de telecomunicaciones, an{\'a}lisis de datos de sensores y sistemas de monitorizaci{\'o}n distribuidos},
		type=main,
		symbol={🪶},
		user1={Base de datos NoSQL distribuida},
		user2={Modelo basado en familias de columnas},
		user3={Alta disponibilidad y tolerancia a fallos},
		see={basedatosnosql,basedatosdistribuida,consistencia,disponibilidad,sibiot}
	}
	
	\newglossaryentry{cascada}
	{
		name={Cascada},
		text={cascada},
		plural={modelos en cascada},
		first={modelo en cascada},
		firstplural={modelos en cascada},
		description={El \emph{modelo en cascada} es un enfoque secuencial de desarrollo de software en el que el proceso avanza de manera lineal a trav{\'e}s de fases bien definidas: an{\'a}lisis de \glspl{requisito}, dise{\~n}o, implementaci{\'o}n, pruebas, despliegue y mantenimiento. Cada fase debe completarse antes de pasar a la siguiente, lo que implica baja \gls{flexibilidad} ante cambios durante el ciclo de vida del proyecto.},
		type=\glsdefaulttype,
		symbol={💧},
		user1={Modelo lineal},
		user2={Waterfall},
		user3={Modelo secuencial},
		see={requisito}
	}
	
	\newglossaryentry{catalogoactivos}{
		name={Cat{\'a}logo de activos},
		text={cat{\'a}logo de activos},
		sort={catalogo activos},
		plural={cat{\'a}logos de activos},
		first={cat{\'a}logo de activos},
		firstplural={cat{\'a}logos de activos},
		description={Repositorio estructurado que registra, clasifica y documenta los activos f{\'i}sicos, digitales o l{\'o}gicos de una organizaci{\'o}n o sistema. Incluye informaci{\'o}n como identificadores {\'u}nicos, ubicaci{\'o}n, estado, responsable, ciclo de vida, dependencias y criticidad. En entornos \gls{iot} y \glspl{sibiot}, el cat{\'a}logo de activos permite gestionar dispositivos, sensores, nodos de red, servicios y componentes de software, facilitando la trazabilidad, el mantenimiento, la auditor{\'i}a y la \gls{tomadecisiones} basada en inventarios confiables.},
		type=main,
		symbol={📚},
		user1={Inventario estructurado de recursos},
		user2={Base para trazabilidad y gesti{\'o}n},
		user3={Control de ciclo de vida},
		see={activo,trazabilidad,gestionempresarial,sibiot}
	}
	
	\newglossaryentry{cbor}{
		name={CBOR},
		text={CBOR},
		plural={CBOR},
		sort={cbor},
		first={Concise Binary Object Representation (CBOR)},
		firstplural={Concise Binary Object Representation (CBOR)},
		description={Formato de serializaci{\'o}n binaria definido por el IETF, dise{\~n}ado para representar datos estructurados de forma compacta y eficiente. CBOR es especialmente adecuado para entornos con recursos limitados, como sistemas \gls{iot}, ya que reduce el consumo de ancho de banda y memoria manteniendo compatibilidad conceptual con estructuras de datos similares a JSON. CBOR se utiliza ampliamente en protocolos de seguridad y comunicaci{\'o}n eficiente, incluyendo \gls{cose}, \gls{oscore} y \gls{edhoc}.},
		type=main
	}
	
	\newglossaryentry{celda}
	{
		name={Celda},
		text={celda},
		sort={celda},
		plural={celdas},
		first={Celda},
		description={Unidad f{\'i}sica o l{\'o}gica elemental que forma parte de una estructura mayor. 
			En ingenier{\'i}a electr{\'o}nica puede referirse a un m{\'o}dulo funcional (p.\,ej., una celda en un arreglo de \textit{sensores} o en un circuito integrado); 
			en metrolog{\'i}a aplicada a \glspl{sibiot}, alude a componentes discretos de medici{\'o}n o acondicionamiento de \glspl{senal}. 
			Su significado concreto depende del dominio (electr{\'o}nico, computacional o metrol{\'o}gico) y del rol que desempe{\~n}a dentro del \gls{sistema}.},
		see=[Relacionado con:]{modulo,dispositivo,sistema}
	}
	
	\newglossaryentry{celdacarga}
	{
		name={Celda de carga},
		text={celda de carga},
		sort={celda de carga},
		plural={celdas de carga},
		first={celda de carga},
		firstplural={celdas de carga},
		description={Una \emph{celda de carga} es un tipo especializado de \gls{sensor} utilizado para medir fuerza o peso mediante la conversi{\'o}n de una carga mec{\'a}nica en una se{\~n}al el{\'e}ctrica. Generalmente emplea galgas extensiom{\'e}tricas que detectan la deformaci{\'o}n del material ante la aplicaci{\'o}n de una fuerza. En \glspl{sibiot}, las \glspl{celdacarga} se integran para la monitorizaci{\'o}n remota de peso en procesos industriales, log{\'i}stica, agricultura de \gls{precision} y control de inventarios.},
		type=\glsdefaulttype,
		symbol={⚖️},
		user1={Load cell},
		user2={Sensor de peso},
		user3={Sensor de fuerza},
		see={sensor}
	}
	
	\newglossaryentry{certificadodigital}{
		name={Certificado digital},
		text={certificado digital},
		sort={certificado digital},
		plural={certificados digitales},
		first={certificado digital},
		firstplural={certificados digitales},
		description={Documento electr{\'o}nico emitido por una autoridad de certificaci{\'o}n confiable que vincula una identidad digital con una clave criptogr{\'a}fica p{\'u}blica. El \emph{certificado digital} permite verificar la autenticidad de entidades, establecer canales cifrados y garantizar la integridad de las comunicaciones mediante infraestructuras de clave p{\'u}blica (PKI). En sistemas IoT y \glspl{sibiot}, los certificados digitales se emplean para autenticar dispositivos, servidores y servicios en protocolos seguros como TLS, HTTPS y WSS, proporcionando un alto nivel de confianza, protecci{\'o}n frente a ataques de suplantaci{\'o}n y soporte a pol{\'i}ticas de seguridad de extremo a extremo},
		type=main,
		symbol={📜},
		see={seguridad,tls,https,wss,autenticacion,integridad,pki,sibiot}
	}
	
	\newglossaryentry{cibernetico}
	{
		name={Cibern{\'e}tico},
		text={cibern{\'e}tico},
		sort={cibernetico},
		plural={cibern{\'e}ticos},
		first={cibern{\'e}tico},
		firstplural={cibern{\'e}ticos},
		description={El t{\'e}rmino \emph{cibern{\'e}tico} se refiere al uso de tecnolog{\'i}as digitales y \glspl{sistemacomunicacion}, especialmente en contextos de \gls{seguridad}, ataques o \glspl{sistemainformatico}.},
		type=\glsdefaulttype,
		symbol={🌐},
		user1={Digital},
		user2={Tecnol{\'o}gico},
		user3={Ciber},
		see={ciberseguridad}
	}
	
	\newglossaryentry{ciberseguridad}
	{
		name={Ciberseguridad},
		text={ciberseguridad},
		sort={ciberseguridad},
		plural={estrategias de ciberseguridad},
		first={ciberseguridad},
		firstplural={estrategias de ciberseguridad},
		description={La \emph{ciberseguridad} es el conjunto de pr{\'a}cticas, tecnolog{\'i}as, procesos y pol{\'i}ticas dise{\~n}adas para proteger \glspl{sistemainformatico}, \glspl{red}, \glspl{dispositivo} y \glspl{dato} frente a accesos no autorizados, ciberataques, da{\~n}os o robos. Abarca aspectos t{\'e}cnicos, organizacionales y legales para garantizar la \gls{confidencialidad}, la \gls{integridad} y la \gls{disponibilidad} de la informaci{\'o}n.},
		type=\glsdefaulttype,
		symbol={🔐},
		user1={Seguridad inform{\'a}tica},
		user2={Seguridad digital},
		user3={Protecci{\'o}n de datos},
		see={amenaza}
	}
	
	\newglossaryentry{cicd}{
		name={CI/CD},
		text={CI/CD},
		sort={cicd},
		type={main},
		plural={CI/CD},
		first={integraci{\'o}n continua y entrega/despliegue continuo (CI/CD)},
		firstplural={integraci{\'o}n continua y entrega/despliegue continuo (CI/CD)},
		description={Conjunto de pr{\'a}cticas y automatizaciones para integrar cambios de software de manera frecuente, ejecutar verificaciones (p.\,ej., compilaci{\'o}n, pruebas y an{\'a}lisis) y entregar o desplegar versiones de forma controlada en entornos de ejecuci{\'o}n, reduciendo riesgos y tiempos de liberaci{\'o}n.},
		see={[see also]{devops,infraestructuracodigo}}
	}
	
	\newglossaryentry{ciclo}{
		name={Ciclo},
		text={ciclo},
		plural={ciclos},
		first={ciclo},
		firstplural={ciclos},
		see={},
		sort={ciclo},
		description={Conjunto de etapas que se repiten de forma iterativa para sostener un proceso de mejora continua. En SIbIoT, el ciclo t{\'i}pico integra captura de datos, an{\'a}lisis, decisi{\'o}n, ejecuci{\'o}n, medici{\'o}n de resultados y retroalimentaci{\'o}n para el ajuste de reglas y modelos}
	}
	
	\newglossaryentry{ciclodesarrollo}{
		name={Ciclo de desarrollo},
		text={ciclo de desarrollo},
		sort={ciclodesarrollo},
		plural={ciclos de desarrollo},
		first={ciclo de desarrollo},
		firstplural={ciclos de desarrollo},
		description={Secuencia organizada y repetible de actividades mediante la cual se concibe, construye, valida y evoluciona una unidad funcional o un sistema. Un ciclo de desarrollo suele abarcar, como m{\'i}nimo, el levantamiento y refinamiento de requisitos, el dise{\~n}o (arquitect{\'o}nico y detallado), la implementaci{\'o}n, la prueba, la integraci{\'o}n, el despliegue y la retroalimentaci{\'o}n para mejora continua. En \gls{iot} y \glspl{sibiot}, el ciclo de desarrollo se extiende a la integraci{\'o}n de hardware, firmware, redes y servicios distribuidos, incorporando validaciones de latencia, consumo energ{\'e}tico, fiabilidad y seguridad.},
		type=main,
		symbol={🔁},
		user1={Proceso iterativo completo},
		user2={Fases de construcci{\'o}n y validaci{\'o}n},
		user3={Base para mejora continua},
		see={iterativa,incremental,pruebaautomatizada,verificacion,despliegue}
	}
	
	\newglossaryentry{ciclovalor}
	{
		name={Ciclo de valor},
		text={ciclo de valor},
		sort={ciclo de valor},
		plural={ciclos de valor},
		first={Ciclo de valor},
		description={Secuencia de actividades, procesos y transformaciones que generan valor a lo largo del desarrollo, implementaci{\'o}n y uso de un \gls{sistema} o producto tecnol{\'o}gico. 
			En el contexto de los \glspl{sibiot}, el ciclo de valor abarca desde la generaci{\'o}n y adquisici{\'o}n de datos mediante \glspl{sensor}, su procesamiento e interpretaci{\'o}n, hasta la toma de decisiones automatizada y la retroalimentaci{\'o}n del entorno. 
			Este ciclo permite medir la eficiencia de los flujos de informaci{\'o}n, optimizar recursos, y maximizar el impacto del conocimiento generado en cada etapa del sistema. 
			Su an{\'a}lisis facilita la identificaci{\'o}n de puntos de mejora, la reducci{\'o}n de redundancias y la integraci{\'o}n de valor en procesos industriales, sociales o ambientales basados en tecnolog{\'i}as IoT.},
		symbol={♻️},
		user1={Secuencia de actividades que generan y amplifican valor},
		user2={Incluye captura, procesamiento y aprovechamiento de datos IoT},
		user3={Permite optimizar eficiencia y sostenibilidad en los SIbIoT},
		see={[Relacionado con:]{valor,sibiot,procesamiento,dato,retroalimentacion}}
	}
	
	\newglossaryentry{ciclovida}{
		name={Ciclo de vida},
		text={ciclo de vida},
		sort={ciclo de vida},
		plural={ciclos de vida},
		first={ciclo de vida},
		firstplural={ciclos de vida},
		description={Conjunto de fases o etapas por las que transcurre un producto, sistema o \gls{software} desde su concepci{\'o}n hasta su retiro o sustituci{\'o}n. El \emph{ciclo de vida} abarca actividades como el an{\'a}lisis de requisitos, el dise{\~n}o, la implementaci{\'o}n, las pruebas, el despliegue, el uso, el \gls{mantenimiento} y la disposici{\'o}n final. En el contexto de los \glspl{sibiot}, el ciclo de vida no solo se aplica al desarrollo de \glspl{aplicacion}, sino tambi{\'e}n a \glspl{sensor}, \glspl{actuador} y dispositivos inteligentes, considerando su durabilidad, actualizaciones de firmware, integraci{\'o}n con servicios en la \gls{nube} y gesti{\'o}n sostenible de residuos electr{\'o}nicos. La gesti{\'o}n adecuada del ciclo de vida garantiza \gls{fiabilidad}, \gls{seguridad}, \gls{eficienciaoperativa} y sostenibilidad},
		type=main,
		symbol={🔄},
		user1={Conjunto de fases desde el inicio hasta el retiro},
		user2={Incluye desarrollo, uso y mantenimiento},
		user3={Aplicable a software y hardware IoT},
		see={sistema,software,sibiot,eficienciaoperativa,sostenibilidad}
	}
	
	\newglossaryentry{cifrado}
	{
		name={Cifrado},
		text={cifrado},
		sort={cifrado},
		plural={cifrados},
		first={cifrado},
		firstplural={cifrados},
		description={\emph{\gls{cifrado}} es el proceso de transformar informaci{\'o}n legible en un formato codificado, accesible solo para quienes posean la clave adecuada. En los \glspl{sibiot}, el cifrado garantiza la confidencialidad e integridad de los datos que viajan a trav{\'e}s de \glspl{sistemacomunicacion} y redes, protegiendo la informaci{\'o}n frente a accesos no autorizados.},
		type=\glsdefaulttype,
		symbol={🔒},
		user1={Encriptaci{\'o}n},
		user2={Protecci{\'o}n criptogr{\'a}fica},
		user3={Codificaci{\'o}n de datos},
		see={cifradocontenido}
	}
	
	\newglossaryentry{cifradocontenido}
	{
		name={Cifrado de contenido},
		text={cifrado de contenido},
		sort={cifrado de contenido},
		plural={cifrados de contenido},
		first={cifrado de contenido},
		firstplural={cifrados de contenido},
		description={\emph{\gls{cifradocontenido}} es una aplicaci{\'o}n espec{\'i}fica del \gls{cifrado} que asegura que la informaci{\'o}n almacenada o transmitida solo pueda ser accedida por usuarios autorizados. En \glspl{sibiot}, el cifrado de contenido protege datos sensibles provenientes de \glspl{sensor}, \glspl{dispositivoiot} y \glspl{dashboard}, evitando filtraciones y fortaleciendo la \gls{seguridad} en entornos distribuidos.},
		type=\glsdefaulttype,
		symbol={🛡️},
		user1={Content encryption},
		user2={Protecci{\'o}n de datos},
		user3={Seguridad de informaci{\'o}n},
		see={cifrado}
	}
	
	\newglossaryentry{cisc}{
		name={CISC},
		text={CISC},
		plural={CISCs},
		first={computador con conjunto de instrucciones complejo (\textit{Complex Instruction Set Computer}, CISC)},
		sort={cisc},
		see={risc,arquitecturaprocesador,cpu},
		description={Paradigma de arquitectura de procesadores caracterizado por un conjunto amplio y complejo de instrucciones, dise{\~n}adas para realizar operaciones de alto nivel en una sola instrucci{\'o}n. La arquitectura CISC busca reducir el n{\'u}mero de instrucciones por programa mediante operaciones m{\'a}s expresivas, a costa de una mayor complejidad en el hardware y, en general, un mayor consumo energ{\'e}tico. Hist{\'o}ricamente, este enfoque ha sido utilizado en arquitecturas ampliamente difundidas como x86 y ha tenido un papel fundamental en la evoluci{\'o}n de los microprocesadores de prop{\'o}sito general.}
	}
	
	% ===================== Ciudad Conectada =====================
	\newglossaryentry{ciudadconectada}{
		name={Ciudad conectada},
		text={ciudad conectada},
		sort={ciudadconectada},
		plural={ciudades conectadas},
		first={ciudad conectada},
		firstplural={ciudades conectadas},
		description={Una \emph{ciudad conectada} es un entorno urbano en el que se despliegan tecnolog{\'i}as de \gls{iot}, redes de \gls{comunicacion} y plataformas digitales para interconectar infraestructuras, servicios y ciudadanos. Se diferencia de la \gls{ciudadinteligente} en que enfatiza la infraestructura de \gls{conectividad} como base para habilitar servicios urbanos digitales: transporte p{\'u}blico monitorizado en tiempo real, alumbrado inteligente, gesti{\'o}n energ{\'e}tica, seguridad ciudadana y servicios p{\'u}blicos digitalizados. El objetivo principal es crear ecosistemas urbanos m{\'a}s sostenibles, eficientes y participativos mediante la integraci{\'o}n de \glspl{dispositivoconectado} y plataformas de gesti{\'o}n centralizada.},
		type=main,
		symbol={🏙️},
		user1={Infraestructura urbana digital basada en IoT},
		user2={Servicios urbanos monitorizados en tiempo real},
		user3={Base para la evoluci{\'o}n hacia ciudades inteligentes},
		see={iot,ciudadinteligente,conectividad,comunicacion,dispositivoconectado}
	}
	
	\newglossaryentry{ciudadinteligente}{
		name={Ciudad inteligente},
		text={ciudad inteligente},
		sort={ciudad inteligente},
		plural={ciudades inteligentes},
		first={ciudad inteligente (\textit{Smart City})},
		firstplural={ciudades inteligentes (\textit{Smart Cities})},
		description={Modelo urbano que integra tecnolog{\'i}as digitales, \glspl{sibiot} y soluciones basadas en datos para mejorar la calidad de vida de los ciudadanos, optimizar la gesti{\'o}n de recursos y promover la \gls{sostenibilidad}. Una \emph{ciudad inteligente} emplea \glspl{sensor}, \glspl{actuador}, \glspl{serviciodistribuido} y plataformas en la \gls{nube} para monitorear y gestionar infraestructuras como transporte, energ{\'i}a, agua, residuos, seguridad y servicios p{\'u}blicos. Adem{\'a}s, favorece la \gls{eficienciaoperativa}, la participaci{\'o}n ciudadana y la resiliencia urbana, integrando tecnolog{\'i}as emergentes como \gls{ia}, \gls{visionartificial}, veh{\'i}culos aut{\'o}nomos y redes 5G},
		type=main,
		symbol={🏙️},
		user1={Modelo urbano basado en tecnolog{\'i}a IoT},
		user2={Optimiza recursos y servicios p{\'u}blicos},
		user3={Promueve sostenibilidad y resiliencia urbana},
		see={sibiot,sostenibilidad,eficienciaoperativa,ia,visionartificial,vehiculoautonomo}
	}
	
	\newglossaryentry{correlacion}
	{
		name={Correlaci{\'o}n},
		text={correlaci{\'o}n},
		sort={correlacion},
		plural={correlaciones},
		first={correlaci{\'o}n},
		firstplural={correlaciones},
		description={Una \emph{correlaci{\'o}n} es una relaci{\'o}n estad{\'i}stica que mide el grado de dependencia o asociaci{\'o}n entre dos o m{\'a}s \glspl{dato} o variables. En \glspl{sibiot}, el an{\'a}lisis de correlaciones permite descubrir patrones de interacci{\'o}n entre \glspl{sensor}, detectar tendencias en series temporales y predecir comportamientos futuros. Su uso es com{\'u}n en \gls{analisisdatos}, \gls{bigdata} y modelos de \gls{aprendizajeautomatico} para inferir relaciones causales o de dependencia.},
		type=\glsdefaulttype,
		symbol={🔗},
		user1={Asociaci{\'o}n},
		user2={Dependencia},
		user3={Correlation},
		see={analisisdatos}
	}
	
	\newglossaryentry{cliente}
	{
		name={Cliente},
		text={cliente},
		sort={cliente},
		plural={clientes},
		first={cliente},
		firstplural={clientes},
		description={Un \emph{cliente} es un componente de un \gls{sistemainformatico} o una aplicaci{\'o}n que realiza solicitudes de servicios o recursos a un \gls{servidor}. Puede referirse tanto a software (como un \gls{browser} o aplicaciones m{\'o}viles) como a hardware (como \glspl{dispositivoiot} conectados a una \gls{red}).},
		type=\glsdefaulttype,
		symbol={🖥️},
		user1={Software cliente},
		user2={Terminal},
		user3={Nodo solicitante},
		see={servidor}
	}
	
	\newglossaryentry{clienteservidor}{
		name={Cliente--servidor},
		text={cliente--servidor},
		first={cliente--servidor},
		sort={cliente servidor},
		description={Modelo de \gls{arquitectura} de \glspl{sistemadistribuido} en el cual las responsabilidades se separan entre entidades que solicitan \glspl{servicio} (\emph{clientes}) y entidades que los proveen (\emph{\glspl{servidor}}). En este \gls{enfoque}, los \glspl{cliente} inician peticiones de recursos, datos o funcionalidades, mientras que los \glspl{servidor} procesan dichas solicitudes y retornan las respuestas correspondientes. El \gls{modelo} cliente--servidor favorece la centralizaci{\'o}n de la l{\'o}gica de negocio, la \gls{escalabilidad} del \gls{sistema}, la \gls{interoperabilidad} entre \glspl{componente} \glspl{heterogeneo} y la administraci{\'o}n controlada de recursos compartidos. Es ampliamente utilizado en sistemas de informaci{\'o}n empresariales, aplicaciones web, servicios en la nube y arquitecturas IoT, donde los dispositivos o aplicaciones cliente interact{\'u}an con servicios remotos a trav{\'e}s de protocolos estandarizados como HTTP, REST o WebSockets},
		plural={cliente--servidor},
		first={cliente--servidor},
		firstplural={cliente--servidor},
		see={arquitecturaestratificada,procesamientodistribuido,servicioweb,rest,websocket}
	}
	
	\newglossaryentry{cloudcomputing}
	{
		name={Cloud Computing},
		text={cloud computing},
		sort={cloud computing},
		plural={servicios de cloud computing},
		first={computaci{\'o}n en la nube o Cloud Computing},
		firstplural={servicios de cloud computing},
		description={El \emph{cloud computing} es un modelo de prestaci{\'o}n de servicios inform{\'a}ticos a trav{\'e}s de Internet, que proporciona acceso bajo demanda a recursos como \gls{almacenamiento}, \gls{procesamientodatos} y aplicaciones, sin necesidad de infraestructura local. Es la base de proveedores como \gls{aws}, Azure u Oracle Cloud, y es clave en arquitecturas modernas de \glspl{sibiot}.},
		type=\glsdefaulttype,
		symbol={☁️},
		user1={Computaci{\'o}n en la nube},
		user2={Servicios en la nube},
		user3={Informatica en la nube},
		see={aws}
	}
	
	\newglossaryentry{cloudfirestore}{
		name={Cloud Firestore},
		text={Cloud Firestore},
		sort={cloudfirestore},
		first={Cloud Firestore},
		description={Servicio de base de datos NoSQL orientado a documentos, completamente gestionado y ofrecido por Google Cloud Platform, dise{\~n}ado para aplicaciones distribuidas que requieren sincronizaci{\'o}n en tiempo real, alta disponibilidad y escalabilidad autom{\'a}tica. Cloud Firestore organiza la informaci{\'o}n en colecciones y documentos jer{\'a}rquicos, permitiendo consultas indexadas, actualizaciones reactivas y replicaci{\'o}n geogr{\'a}fica transparente. En sistemas de informaci{\'o}n basados en IoT y SIbIoT, Cloud Firestore se emplea como repositorio operacional o de coordinaci{\'o}n para estados de dispositivos, eventos recientes y datos de aplicaciones m{\'o}viles, complementando arquitecturas policl{\'i}nicas de persistencia junto a bases de datos relacionales y anal{\'i}ticas},
		type=main,
		symbol={🔥},
		see={nosql,procesamientocontinuo,arquitecturapoliclinica,cloudcomputing,sibiot}
	}
	
	\newglossaryentry{cloudspanner}{
		name={Cloud Spanner},
		text={Cloud Spanner},
		sort={cloud spanner},
		first={Cloud Spanner},
		description={Denominaci{\'o}n abreviada y ampliamente aceptada de \gls{googlecloudspanner}. Se utiliza frecuentemente en documentaci{\'o}n t{\'e}cnica y literatura acad{\'e}mica para referirse al sistema de base de datos relacional distribuida de Google con consistencia fuerte, escalabilidad horizontal y soporte de transacciones globales},
		type=main,
		symbol={🗄️},
		see={googlecloudspanner,spanner}
	}
	
	% ===================== Cloud SQL =====================
	\newglossaryentry{cloudsql}{
		name={Cloud SQL},
		text={\textit{Cloud SQL}},
		sort={cloudsql},
		plural={\textit{Cloud SQL}},
		first={\textit{Cloud SQL}},
		firstplural={\textit{Cloud SQL}},
		description={\emph{Cloud SQL} es un servicio de \gls{gcp} que ofrece \glspl{basedatosrelacional} administradas en la nube, con soporte para motores como \gls{mysql}, \gls{postgresql} y \gls{sqlserver}. Este servicio automatiza tareas como copias de seguridad, replicaci{\'o}n, actualizaciones y escalamiento. En proyectos \gls{iot}, \emph{Cloud SQL} permite gestionar datos estructurados generados por \glspl{dispositivoiot}, asegurando \gls{fiabilidad}, \gls{seguridad} y compatibilidad con ecosistemas empresariales.},
		type=main,
		symbol={🗄️},
		user1={Servicio administrado de SQL en GCP},
		user2={Compatibilidad con MySQL, PostgreSQL y SQL Server},
		user3={Soporte a datos estructurados en IoT},
		see={gcp,basedatosrelacional,aws,azure}
	}
	
	\newglossaryentry{cluster}
	{
		name={Cl{\'u}ster},
		text={cl{\'u}ster},
		sort={cluster},
		plural={cl{\'u}steres},
		first={cl{\'u}ster},
		firstplural={cl{\'u}steres},
		description={Un \emph{cl{\'u}ster} es un conjunto de computadoras o nodos interconectados que trabajan de manera coordinada como si fueran un solo sistema, para aumentar la capacidad de \gls{procesamientodatos}, la disponibilidad o el \gls{almacenamiento}. En contextos de \gls{bigdata} y \glspl{sibiot}, los cl{\'u}steres permiten el manejo de grandes vol{\'u}menes de informaci{\'o}n y la ejecuci{\'o}n de algoritmos distribuidos.},
		type=\glsdefaulttype,
		symbol={🖧},
		user1={Grupo de nodos},
		user2={Cluster inform{\'a}tico},
		user3={Sistema agrupado},
		see={bigdata}
	}
	
	\newglossaryentry{coap}
	{
		name={CoAP},
		text={CoAP},
		first={protocolo de aplicaci{\'o}n para entornos restringidos (\textit{Constrained Application Protocol}, CoAP)},
		sort={coap},
		plural={CoAPs},
		firstplural={CoAPs},
		description={\emph{CoAP} es un protocolo de aplicaci{\'o}n dise{\~n}ado espec{\'i}ficamente para entornos de \gls{iot} y \glspl{dispositivoiot} con recursos limitados. Proporciona un mecanismo eficiente para el intercambio de \glspl{dato} en la web, optimizando el uso de ancho de banda y energ{\'i}a en \glspl{red} de baja potencia.},
		type=\glsdefaulttype,
		symbol={📡},
		user1={Constrained Application Protocol},
		user2={Protocolo IoT ligero},
		user3={Protocolo de aplicaci{\'o}n restringido},
		see={lpwan}
	}
	
	
	\newglossaryentry{codigo}{
		name={C{\'o}digo},
		text={c{\'o}digo},
		sort={codigo},
		plural={c{\'o}digos},
		first={c{\'o}digo},
		firstplural={c{\'o}digos},
		description={El \emph{c{\'o}digo} es un concepto amplio que puede adoptar distintos significados seg{\'u}n el contexto: 
			\begin{itemize}
				\item En programaci{\'o}n, se refiere al \gls{codigofuente} escrito en un \gls{lenguajeprogramacion}, que posteriormente se compila o interpreta para ser ejecutado por un \gls{microprocesador} o un \gls{microcontrolador}.
				\item En identificaci{\'o}n, corresponde a secuencias num{\'e}ricas o alfanum{\'e}ricas que garantizan la unicidad de entidades dentro de un sistema, como c{\'o}digos de barras, c{\'o}digos QR o identificadores RFID, fundamentales en trazabilidad, \gls{seguridad} y gesti{\'o}n de activos en entornos \gls{iot}.
				\item En criptograf{\'i}a, designa sistemas o algoritmos utilizados para transformar mensajes en formatos ininteligibles, protegiendo la \gls{confidencialidad} y la \gls{integridad} de los datos. Ejemplos incluyen \gls{cifrado}, \gls{firmadigital} y protocolos seguros como \gls{https}.
			\end{itemize}
			De esta manera, el t{\'e}rmino abarca tanto el \gls{codigofuente} que da vida a las aplicaciones, como los mecanismos de identificaci{\'o}n y seguridad que soportan la infraestructura de los \glspl{sibiot}.},
		type=\glsdefaulttype,
		symbol={💻},
		user1={Programaci{\'o}n (c{\'o}digo fuente)},
		user2={Identificaci{\'o}n (barras, QR, RFID)},
		user3={Seguridad (c{\'o}digos criptogr{\'a}ficos)},
		see={codigofuente,lenguajeprogramacion,criptografia,confidencialidad,identificacion}
	}
	
	% ===================== C{\'o}digo fuente =====================
	\newglossaryentry{codigofuente}{
		name={C{\'o}digo fuente},
		text={c{\'o}digo fuente},
		sort={codigo fuente},
		plural={c{\'o}digos fuente},
		first={c{\'o}digo fuente},
		firstplural={c{\'o}digos fuente},
		description={El \emph{c{\'o}digo fuente} es el conjunto de instrucciones y declaraciones escritas en un \gls{lenguajeprogramacion} que definen el comportamiento de un \gls{sistema} o aplicaci{\'o}n. Este texto, legible por humanos, debe ser compilado o interpretado para ejecutarse en un \gls{microprocesador} o \gls{microcontrolador}. En el contexto de \glspl{sibiot}, el c{\'o}digo fuente integra la l{\'o}gica de negocio, el manejo de \glspl{sensor} y \glspl{actuador}, y la comunicaci{\'o}n con servicios distribuidos.},
		type=\glsdefaulttype,
		symbol={👨‍💻},
		user1={Texto escrito por desarrolladores},
		user2={Necesita compilaci{\'o}n o interpretaci{\'o}n},
		user3={Base del software IoT},
		see={lenguajeprogramacion,codigo,microprocesador,microcontrolador,sibiot}
	}
	
	% ===================== C{\'o}digo de identificaci{\'o}n =====================
	\newglossaryentry{codigoidentificacion}{
		name={C{\'o}digo de identificaci{\'o}n},
		text={c{\'o}digo de identificaci{\'o}n},
		sort={codigo identificacion},
		plural={c{\'o}digos de identificaci{\'o}n},
		first={c{\'o}digo de identificaci{\'o}n},
		firstplural={c{\'o}digos de identificaci{\'o}n},
		description={Un \emph{c{\'o}digo de identificaci{\'o}n} es una secuencia num{\'e}rica, alfab{\'e}tica o alfanum{\'e}rica asignada a un objeto, persona o entidad para garantizar su unicidad dentro de un sistema. Ejemplos incluyen el c{\'o}digo de barras, el c{\'o}digo QR, el IMEI de dispositivos m{\'o}viles o los identificadores RFID, fundamentales en trazabilidad, \gls{seguridad} y gesti{\'o}n de activos en ecosistemas \gls{iot}.},
		type=\glsdefaulttype,
		symbol={🔑},
		user1={Identificaci{\'o}n {\'u}nica de objetos o personas},
		user2={Ejemplos: QR, c{\'o}digos de barras, IMEI},
		user3={Uso en trazabilidad y seguridad IoT},
		see={identificacion,rfid,sistemacomunicacion}
	}
	
	% ===================== C{\'o}digo criptogr{\'a}fico =====================
	\newglossaryentry{codigocriptografico}{
		name={C{\'o}digo criptogr{\'a}fico},
		text={c{\'o}digo criptogr{\'a}fico},
		sort={codigo criptografico},
		plural={c{\'o}digos criptogr{\'a}ficos},
		first={c{\'o}digo criptogr{\'a}fico},
		firstplural={c{\'o}digos criptogr{\'a}ficos},
		description={Un \emph{c{\'o}digo criptogr{\'a}fico} es un conjunto de reglas o algoritmos destinados a transformar mensajes en formatos ininteligibles, con el fin de proteger la \gls{confidencialidad} y la \gls{integridad} de la informaci{\'o}n. Se aplica en \gls{cifrado}, \gls{firmadigital}, protocolos seguros como \gls{https} y sistemas de autenticaci{\'o}n en redes \gls{iot}.},
		type=\glsdefaulttype,
		symbol={🔐},
		user1={Protecci{\'o}n de datos},
		user2={Uso en criptograf{\'i}a},
		user3={Base de seguridad IoT},
		see={cifrado,firmadigital,https,seguridad}
	}
	
	\newglossaryentry{compatibilidad}{
		name={Compatibilidad},
		text={compatibilidad},
		sort={compatibilidad},
		plural={compatibilidades},
		first={compatibilidad},
		firstplural={compatibilidades},
		description={Propiedad de un \gls{sistema}, \gls{software}, \gls{hardware} o \gls{plataforma} que le permite funcionar correctamente en conjunto con otros elementos, sin generar conflictos ni requerir modificaciones sustanciales. La \emph{compatibilidad} puede ser hacia atr{\'a}s (funcionar con versiones anteriores), hacia adelante (funcionar con versiones futuras) o cruzada (funcionar con sistemas distintos). En el contexto de los \glspl{sibiot}, la compatibilidad asegura que \glspl{sensor}, \glspl{actuador}, protocolos de \gls{comunicacioninalambrica} y servicios en la nube trabajen de forma integrada, favoreciendo la \gls{interoperabilidad}, la \gls{escalabilidad} y la extensi{\'o}n del \gls{ecosistemaiot}.},
		type=main,
		symbol={🔗},
		user1={Funcionamiento conjunto entre sistemas o versiones},
		user2={Soporte a integraci{\'o}n sin conflictos},
		user3={Base para interoperabilidad en SIbIoT},
		see={interoperabilidad,escalabilidad,software,hardware,sibiot}
	}
	
	\newglossaryentry{componente}{
		name={Componente},
		text={componente},
		sort={componente},
		plural={componentes},
		first={componente},
		firstplural={componentes},
		description={Unidad b{\'a}sica, f{\'i}sica o l{\'o}gica, que forma parte de un \gls{sistema} y cumple una funci{\'o}n espec{\'i}fica dentro de {\'e}l. Un \emph{componente} puede ser hardware (como \glspl{sensor}, \glspl{actuador}, microcontroladores o dispositivos de red) o software (m{\'o}dulos, bibliotecas, servicios o aplicaciones), y su interacci{\'o}n coordinada garantiza el funcionamiento global. En el contexto de los \glspl{sibiot}, los componentes se integran para capturar, procesar y analizar datos en \gls{tiemporeal}, favoreciendo la \gls{interoperabilidad}, la \gls{modularidad} y la \gls{escalabilidad} de los sistemas.},
		type=main,
		symbol={⚙️},
		user1={Unidad b{\'a}sica de un sistema},
		user2={Puede ser hardware o software},
		user3={Elemento modular en SIbIoT},
		see={sistema,sensor,actuador,software,sibiot}
	}
	
	\newglossaryentry{componenteelectronico}{
		name={Componente electr{\'o}nico},
		text={componente electr{\'o}nico},
		sort={componente electronico},
		plural={componentes electr{\'o}nicos},
		first={componente electr{\'o}nico},
		firstplural={componentes electr{\'o}nicos},
		description={Dispositivo f{\'i}sico que forma parte de un circuito electr{\'o}nico y que posee una funci{\'o}n espec{\'i}fica en el procesamiento, control o transmisi{\'o}n de se{\~n}ales el{\'e}ctricas. Un \emph{componente electr{\'o}nico} puede ser pasivo (resistencias, condensadores, inductores) o activo (diodos, transistores, circuitos integrados, microcontroladores). En el contexto de los \glspl{sibiot}, los componentes electr{\'o}nicos son la base del dise{\~n}o de \glspl{sensor}, \glspl{actuador} y nodos, permitiendo la captura de \glspl{variablefisica}, el procesamiento local de datos y la interconexi{\'o}n con plataformas de \gls{monitorizacion} y \gls{control}. Su calidad, fiabilidad y eficiencia influyen directamente en la \gls{disponibilidad}, \gls{fiabilidad} y \gls{sostenibilidad} de los \glspl{sistemainteligente}.},
		type=main,
		symbol={🔌},
		user1={Elemento f{\'i}sico de un circuito},
		user2={Clasificaci{\'o}n en pasivos y activos},
		user3={Base de sensores, actuadores y nodos en IoT},
		see={componente,sensor,actuador,variablefisica,sibiot}
	}
	
	\newglossaryentry{comportamiento}{
		name={Comportamiento},
		text={comportamiento},
		sort={comportamiento},
		plural={comportamientos},
		first={comportamiento},
		firstplural={comportamientos},
		description={El comportamiento representa la respuesta din{\'a}mica de un sistema ante variaciones internas o externas. En los \glspl{sibiot}, el comportamiento puede referirse tanto a las acciones automatizadas de los \glspl{dispositivo} inteligentes como a la interpretaci{\'o}n de patrones de conducta humana obtenidos mediante \glspl{sensor} y algoritmos de \gls{aprendizajeautomatico}. Su an{\'a}lisis permite la adaptaci{\'o}n del sistema, la detecci{\'o}n de anomal{\'i}as y la mejora continua de los procesos basados en datos.},
		type=\glsdefaulttype,
		symbol={🔄},
		user1={Behavior},
		user2={Respuesta observable ante est{\'i}mulos},
		user3={An{\'a}lisis de patrones en sistemas y usuarios},
		see={patron,aprendizajeautomatico,analisis,usuario,sibiot}
	}
	
	\newglossaryentry{composicionservicios}
	{
		name={Composici{\'o}n de servicios},
		text={composici{\'o}n de servicios},
		sort={composicion de servicios},
		plural={composiciones de servicios},
		first={composici{\'o}n de servicios},
		firstplural={composiciones de servicios},
		description={La \emph{composici{\'o}n de servicios} es la construcci{\'o}n de un proceso organizacional combinando varios \glspl{servicio} mediante reglas de secuencia, condiciones y manejo de fallos, para ejecutar flujos de negocio sin integraciones punto a punto.},
		type=\glsdefaulttype,
		symbol={🧩},
		user1={Orquestaci{\'o}n},
		user2={Flujo de proceso},
		user3={Reutilizaci{\'o}n},
		see={soa,api}
	}
	
	\newglossaryentry{computacionverde}{
		name={Computaci{\'o}n verde},
		text={computaci{\'o}n verde},
		sort={computacion verde},
		first={computaci{\'o}n verde (\emph{Green Computing})},
		description={Enfoque de dise{\~n}o, desarrollo, uso y disposici{\'o}n de sistemas computacionales orientado a minimizar el impacto ambiental a lo largo de su ciclo de vida. La computaci{\'o}n verde promueve la eficiencia energ{\'e}tica, la reducci{\'o}n de emisiones, el uso responsable de recursos y la adopci{\'o}n de tecnolog{\'i}as sostenibles. En el contexto de los \glspl{sibiot}, este enfoque se materializa mediante arquitecturas de bajo consumo, optimizaci{\'o}n de comunicaciones, gesti{\'o}n energ{\'e}tica inteligente y aprovechamiento de energ{\'i}as renovables.},
		type=main,
		symbol={🌱},
		user1={Sostenibilidad},
		user2={Eficiencia energ{\'e}tica},
		user3={Impacto ambiental reducido},
		see={consumoenergetico,energia,arquitecturadistribuida}
	}
	
	
	\newglossaryentry{computador}{
		name={Computador},
		text={computador},
		sort={computador},
		plural={computadores},
		first={computador},
		firstplural={computadores},
		description={Un \emph{computador} es un dispositivo electr{\'o}nico programable capaz de recibir datos, procesarlos mediante operaciones l{\'o}gicas y aritm{\'e}ticas, almacenarlos y generar resultados como salida. Se compone de elementos de \gls{hardware} como la \gls{cpu}, \gls{memoria} principal, \glspl{dispositivo} de entrada/salida y sistemas de almacenamiento, coordinados por \gls{software} que permite ejecutar aplicaciones y gestionar recursos. En el contexto de los \glspl{sibiot}, los computadores funcionan como nodos de procesamiento local, servidores de borde, plataformas de an{\'a}lisis, gateways o centros de control, facilitando el \gls{procesamientodatos}, la visualizaci{\'o}n, la administraci{\'o}n de dispositivos y la integraci{\'o}n con servicios en la \gls{nube}.},
		type=main,
		symbol={💻},
		user1={Dispositivo electr{\'o}nico programable},
		user2={Unidad de procesamiento y almacenamiento},
		user3={Nodo de gesti{\'o}n y control en SIbIoT},
		see={hardware,software,cpu,microprocesador,procesamientodatos,sibiot}
	}
	
	
	\newglossaryentry{comunicacion}{
		name={Comunicaci{\'o}n},
		text={comunicaci{\'o}n},
		sort={comunicacion},
		plural={comunicaciones},
		first={comunicaci{\'o}n},
		firstplural={comunicaciones},
		description={La \emph{comunicaci{\'o}n} es el proceso mediante el cual dos o m{\'a}s entidades (personas, dispositivos o sistemas) intercambian informaci{\'o}n a trav{\'e}s de un medio com{\'u}n. En el contexto tecnol{\'o}gico, puede realizarse de forma anal{\'o}gica o digital, al{\'a}mbrica o inal{\'a}mbrica, y puede involucrar distintos \glspl{protocolocomunicacion} y \glspl{sistemacomunicacion}. En sistemas \gls{iot}, la comunicaci{\'o}n conecta \glspl{sensor}, \glspl{actuador}, \glspl{dispositivoiot} y plataformas en la \gls{nube}, garantizando que los datos capturados sean transmitidos y procesados con \gls{fiabilidad}, baja \gls{latencia} y adecuada \gls{seguridad}.},
		type=main,
		symbol={🔊},
		user1={Intercambio de informaci{\'o}n entre entidades},
		user2={Puede ser anal{\'o}gica, digital, al{\'a}mbrica o inal{\'a}mbrica},
		user3={Base de IoT mediante sensores, actuadores y protocolos},
		see={transmision,protocolocomunicacion,sistemacomunicacion,iot,fiabilidad,seguridad}
	}
	
	\newglossaryentry{comunicacioninalambrica}{
		name={Comunicaci{\'o}n inal{\'a}mbrica},
		text={comunicaci{\'o}n inal{\'a}mbrica},
		sort={comunicacion inalambrica},
		plural={comunicaciones inal{\'a}mbricas},
		first={comunicaci{\'o}n inal{\'a}mbrica},
		firstplural={comunicaciones inal{\'a}mbricas},
		description={Transmisi{\'o}n y recepci{\'o}n de informaci{\'o}n entre dispositivos sin necesidad de cables f{\'i}sicos, utilizando ondas electromagn{\'e}ticas como radiofrecuencia, infrarrojo o microondas. La \emph{comunicaci{\'o}n inal{\'a}mbrica} es esencial en redes modernas por su flexibilidad, movilidad y facilidad de despliegue. En el contexto de los \glspl{sibiot}, permite la interconexi{\'o}n de \glspl{sensor}, \glspl{actuador}, \glspl{dispositivointeligente} y plataformas de \gls{control}, empleando protocolos como \gls{wifi}, \gls{bluetooth}, \gls{zigbee}, \gls{lorawan} o NB-IoT. Su desempe{\~n}o depende de par{\'a}metros como \gls{anchobanda}, \gls{rendimiento}, \gls{latencia} y \gls{fiabilidad}, siendo fundamental para garantizar servicios en \gls{tiemporeal} en aplicaciones cr{\'i}ticas.},
		type=main,
		symbol={📡},
		user1={Transmisi{\'o}n de datos sin cables},
		user2={Base de conectividad en IoT},
		user3={Flexibilidad y movilidad en sistemas inteligentes},
		see={red,anchobanda,latencia,rendimiento,sibiot}
	}
	
	\newglossaryentry{concurrente}{
		name={Concurrente},
		text={concurrente},
		sort={concurrente},
		plural={concurrentes},
		first={concurrente},
		firstplural={concurrentes},
		description={Propiedad de un sistema o proceso que permite la ejecuci{\'o}n de varias tareas de manera solapada en el tiempo, ya sea mediante paralelismo real o alternancia controlada por el sistema. En microcontroladores, la concurrencia suele lograrse mediante temporizaci{\'o}n y control del flujo de ejecuci{\'o}n.},
		type=main,
		symbol={🔀},
		user1={Ejecuci{\'o}n solapada},
		user2={Control de tiempo},
		user3={Gest{i}on de procesos},
		see={tareaconcurrente,sincronizada}
	}
	
	\newglossaryentry{condicion}
	{
		name={Condici{\'o}n},
		text={condici{\'o}n},
		sort={condicion},
		plural={condiciones},
		first={condici{\'o}n},
		firstplural={condiciones},
		description={Una condici{\'o}n es una circunstancia, restricci{\'o}n o conjunto de factores que influyen en el desempe{\~n}o y operaci{\'o}n de un \gls{sistema}. En los \glspl{sibiot}, las condiciones incluyen aspectos ambientales, t{\'e}cnicos y operativos que determinan la \gls{efectividad} y \gls{confiabilidad} del \gls{sistemainformatico}.},
		type=\glsdefaulttype,
		symbol={},
		user1={Circunstancia},
		user2={Restricci{\'o}n},
		user3={Factor},
		see={filtro,filtradodatos,filtrodatos,sibiot}
	}
	
	\newglossaryentry{conectividad}
	{
		name={Conectividad},
		text={conectividad},
		sort={conectividad},
		plural={conectividades},
		first={conectividad},
		firstplural={conectividades},
		description={La \emph{conectividad} es la capacidad que puede tener un dispositivo, desde ordenadores hasta objetos de uso cotidiano, de conectarse y comunicarse con otro con el fin de intercambiar informaci{\'o}n o establecer una conexi{\'o}n digital. En los \glspl{sibiot}, la conectividad asegura la interacci{\'o}n fluida entre \glspl{dispositivo}, \glspl{sensor}, \glspl{actuador} y \glspl{plataforma} de an{\'a}lisis en la \gls{nube}.},
		type=\glsdefaulttype,
		symbol={},
		user1={Interconexi{\'o}n},
		user2={Enlace},
		user3={Comunicaci{\'o}n digital},
		see={conexion}
	}
	
	\newglossaryentry{conexion}
	{
		name={Conexi{\'o}n},
		text={conexi{\'o}n},
		sort={conexion},
		plural={conexiones},
		first={conexi{\'o}n},
		firstplural={conexiones},
		description={Seg{\'u}n la Real Academia Espa{\~n}ola, conexi{\'o}n es la uni{\'o}n que se establece entre dos o m{\'a}s cosas, como aparatos, sistemas, lugares o personas, para que entre ellas haya una relaci{\'o}n o comunicaci{\'o}n. En el contexto de los \glspl{sibiot}, la conexi{\'o}n constituye la base f{\'i}sica o l{\'o}gica que permite la interoperabilidad de los \glspl{dispositivo} y \glspl{servicio}.},
		type=\glsdefaulttype,
		symbol={},
		user1={Enlace},
		user2={V{\'i}nculo},
		user3={Uni{\'o}n},
		see={conectividad}
	}
	
	\newglossaryentry{confiabilidad}
	{
		name={Confiabilidad},
		text={confiabilidad},
		sort={confiabilidad},
		plural={confiabilidades},
		first={confiabilidad},
		firstplural={confiabilidades},
		description={La confiabilidad es la probabilidad de que un \gls{sistema} funcione sin fallos durante un periodo de tiempo determinado y bajo condiciones establecidas. En el contexto de los \glspl{sibiot}, la confiabilidad se refiere a la capacidad de los dispositivos y redes para operar de manera continua, entregar datos precisos y ejecutar tareas cr{\'i}ticas sin interrupciones, lo cual resulta fundamental en aplicaciones sensibles como salud, transporte o automatizaci{\'o}n industrial.},
		type=\glsdefaulttype,
		symbol={},
		user1={Fiabilidad},
		user2={Disponibilidad operativa},
		user3={Resiliencia t{\'e}cnica},
		see={efectividad}
	}
	
	\newglossaryentry{confiable}
	{
		name={Confiable},
		text={confiable},
		sort={confiable},
		plural={confiables},
		first={confiable},
		firstplural={confiables},
		description={Un elemento es confiable cuando mantiene un desempe{\~n}o consistente y estable bajo condiciones definidas, generando confianza en sus resultados. En los \glspl{sibiot}, la confiabilidad implica que los \glspl{dispositivo} y servicios operen sin fallos, mantengan la seguridad de los datos y garanticen la continuidad de las funciones cr{\'i}ticas.},
		type=\glsdefaulttype,
		symbol={},
		user1={Seguro},
		user2={Estable},
		user3={Fidedigno},
		see={confiabilidad}
	}
	
	\newglossaryentry{confidencialidad}{
		name={Confidencialidad},
		text={confidencialidad},
		sort={confidencialidad},
		plural={confidencialidades},
		first={confidencialidad},
		firstplural={confidencialidades},
		description={Propiedad de la \gls{seguridad} de la informaci{\'o}n que garantiza que los datos solo sean accesibles por personas, sistemas o entidades autorizadas, evitando su divulgaci{\'o}n no autorizada. La \emph{confidencialidad} se asegura mediante mecanismos como \glspl{metodocifrado}, autenticaci{\'o}n y pol{\'i}ticas de \gls{controlacceso}. En el contexto de los \glspl{sibiot}, la confidencialidad protege datos sensibles generados por \glspl{sensor} o intercambiados entre \glspl{nodo}, resguardando la \gls{privacidad} de los usuarios y la integridad de los procesos cr{\'i}ticos en aplicaciones industriales, m{\'e}dicas o de \gls{seguridad}.},
		type=main,
		symbol={🔒},
		user1={Protecci{\'o}n de datos frente a accesos no autorizados},
		user2={Pilar de la seguridad de la informaci{\'o}n},
		user3={Esencial en IoT y sistemas distribuidos},
		see={seguridad,privacidad,controlacceso,metodocifrado,sibiot}
	}
	
	\newglossaryentry{configuracion}{
		name={Configuraci{\'o}n},
		text={configuraci{\'o}n},
		sort={configuracion},
		first={configuraci{\'o}n},
		description={Conjunto de par{\'a}metros y opciones que determinan el comportamiento de un sistema, incluyendo ajustes de red, seguridad, almacenamiento y pol{\'i}ticas operativas.},
		type=main,
		symbol={⚙️},
		user1={Par{\'a}metros},
		user2={Ajustes del sistema},
		user3={Operaci{\'o}n},
		see={gestiondatos,controlversiones}
	}
	
	\newglossaryentry{conmutacion}{
		name={conmutaci{\'o}n},
		text={conmutaci{\'o}n},
		sort={conmutacion},
		plural={conmutaciones},
		first={conmutaci{\'o}n},
		firstplural={conmutaciones},
		description={Proceso mediante el cual un sistema de \gls{red} selecciona, cambia o redirige rutas de \gls{comunicacion} para garantizar el flujo eficiente y continuo de datos. En el contexto de los \glspl{sibiot}, la conmutaci{\'o}n permite responder a variaciones de carga, fallos en enlaces o cambios en la calidad de la se{\~n}al, mejorando la disponibilidad del sistema},
		symbol={🔀},
		see={enrutamiento,fallo,conmutacionautomatica},
		type=\glsdefaulttype
	}
	
	\newglossaryentry{conmutacionautomatica}{
		name={conmutaci{\'o}n autom{\'a}tica},
		text={conmutaci{\'o}n autom{\'a}tica},
		sort={conmutacion automatica},
		plural={conmutaciones autom{\'a}ticas},
		first={conmutaci{\'o}n autom{\'a}tica},
		firstplural={conmutaciones autom{\'a}ticas},
		description={Mecanismo mediante el cual un sistema detecta un \gls{fallo} o degradaci{\'o}n en un enlace o componente y cambia autom{\'a}ticamente a una ruta o recurso alternativo sin interrumpir la \gls{conectividad}. Es esencial en \glspl{sibiot} para asegurar la continuidad operativa mediante redundancia, ruteo m{\'u}ltiple y balanceo inteligente de carga},
		symbol={🔁},
		see={fallo,enrutamiento},
		type=\glsdefaulttype
	}
	
	\newglossaryentry{conocimiento}{
		name={Conocimiento},
		text={conocimiento},
		sort={conocimiento},
		plural={conocimientos},
		first={conocimiento},
		firstplural={conocimientos},
		description={El conocimiento representa la fase m{\'a}s avanzada del tratamiento de la informaci{\'o}n, donde los datos y la informaci{\'o}n se integran para generar comprensi{\'o}n, predicci{\'o}n o toma de decisiones. En los \glspl{sibiot}, el conocimiento surge del an{\'a}lisis inteligente de los datos capturados por \glspl{sensor} y procesados mediante t{\'e}cnicas de \gls{analisisdescriptivo}, \gls{analisispredictivo} o \gls{analisisprescriptivo}, aportando valor a la gesti{\'o}n, la automatizaci{\'o}n y la mejora de procesos.},
		type=\glsdefaulttype,
		symbol={📘},
		user1={Knowledge},
		user2={Integraci{\'o}n de informaci{\'o}n para la toma de decisiones},
		user3={Nivel superior del procesamiento de datos},
		see={informacion,dato,analisisdescriptivo,analisispredictivo,analisisprescriptivo,sibiot}
	}
	
	\newglossaryentry{consistencia}{
		name={Consistencia},
		text={consistencia},
		sort={consistencia},
		plural={consistencias},
		first={consistencia},
		firstplural={consistencias},
		description={Propiedad de un \gls{sistema} o de una \gls{basedatos} que asegura que los datos mantienen validez, coherencia y cumplimiento de reglas definidas en todo momento, incluso tras operaciones concurrentes, fallos o procesos de \gls{replicacion}. La \emph{consistencia} implica que todos los usuarios o \glspl{nodo} perciben el mismo estado de la informaci{\'o}n seg{\'u}n las garant{\'i}as del modelo de datos empleado. En el contexto de las \glspl{basedatosdistribuida} y de los \glspl{sibiot}, la consistencia est{\'a} relacionada con el teorema CAP y puede clasificarse como fuerte, d{\'e}bil o eventual, dependiendo de si los cambios se reflejan de inmediato, con tolerancia a retardo, o de forma progresiva. Garantizar consistencia es fundamental para la \gls{fiabilidad}, la \gls{disponibilidad} y la confianza en la \gls{tomadecisiones} en \gls{tiemporeal}.},
		type=main,
		symbol={🧩},
		user1={Coherencia y validez de datos},
		user2={Garant{\'i}a en sistemas distribuidos},
		user3={Clasificaci{\'o}n fuerte, d{\'e}bil o eventual},
		see={replicacion,sincronizacion,fiabilidad,disponibilidad,basedatosdistribuida,sibiot}
	}
	
		\newglossaryentry{consistenciafuerte}{
		name={Consistencia fuerte},
		text={consistencia fuerte},
		sort={consistencia fuerte},
		first={consistencia fuerte},
		description={Garant{\'i}a de que, tras una escritura confirmada, toda lectura posterior devuelve el {\'u}ltimo valor escrito, evitando observaciones divergentes entre nodos. Es deseable en operaciones cr{\'i}ticas, aunque puede implicar mayor latencia o menor disponibilidad.},
		type=main,
		symbol={🔒},
		user1={Lecturas actualizadas},
		user2={Coherencia global},
		user3={Trade-offs},
		see={basedatosdistribuida,sistemadistribuido}
	}
	
	\newglossaryentry{construccion}{
		name={Construcci{\'o}n},
		text={construcci{\'o}n},
		sort={construccion},
		plural={construcciones},
		first={construcci{\'o}n},
		firstplural={construcciones},
		description={Proceso sistem{\'a}tico mediante el cual se materializa una soluci{\'o}n tecnol{\'o}gica a partir de modelos, especificaciones y decisiones arquitect{\'o}nicas previamente definidas. En ingenier{\'i}a de software y sistemas IoT, la construcci{\'o}n implica implementar componentes preservando coherencia estructural, trazabilidad y cumplimiento de restricciones.},
		type=main,
		symbol={🏗️},
		user1={Materializaci{\'o}n de modelos},
		user2={Implementaci{\'o}n controlada},
		user3={Ejecuci{\'o}n conforme a especificaci{\'o}n},
		see={modelo,especificaciontecnica,decisionarquitectonica}
	}
	
	\newglossaryentry{construccioncolaborativa}{
		name={Construcci{\'o}n colaborativa},
		text={construcci{\'o}n colaborativa},
		sort={construccioncolaborativa},
		plural={construcciones colaborativas},
		first={construcci{\'o}n colaborativa},
		firstplural={construcciones colaborativas},
		description={Modalidad de construcci{\'o}n de soluciones en la que participan activamente m{\'u}ltiples actores (equipo t{\'e}cnico, usuarios, responsables de operaci{\'o}n, calidad y seguridad) en la definici{\'o}n, refinamiento, verificaci{\'o}n y evoluci{\'o}n de los artefactos del sistema. Se apoya en mecanismos de coordinaci{\'o}n (revisi{\'o}n por pares, control de versiones, integraci{\'o}n continua, acuerdos de interfaces y criterios de aceptaci{\'o}n), con el fin de reducir ambig{\"u}edades, mejorar la calidad y evitar divergencias entre dise{\~n}o e implementaci{\'o}n. En \glspl{sibiot}, integra adem{\'a}s a perfiles de hardware, redes y datos para asegurar consistencia extremo a extremo.},
		type=main,
		symbol={🤝},
		user1={Co-creaci{\'o}n y revisi{\'o}n},
		user2={Alineaci{\'o}n entre actores},
		user3={Reducci{\'o}n de divergencias},
		see={equipoagil,iterativa,integracion,consistencia,verificacion}
	}
	
	\newglossaryentry{consumoenergetico}{
		name={Consumo energ{\'e}tico},
		text={consumo energ{\'e}tico},
		sort={consumo energetico},
		first={consumo energ{\'e}tico},
		description={Cantidad de energ{\'i}a utilizada por un dispositivo o sistema en un intervalo de tiempo. En IoT, se optimiza mediante modos de bajo consumo, muestreo eficiente y comunicaciones oportunistas.},
		type=main,
		symbol={⚡},
		user1={Energ{\'i}a usada},
		user2={Eficiencia},
		user3={Autonom{\'i}a},
		see={bateria,energia}
	}
	
	\newglossaryentry{contenedor}
	{
		name={Contenedor},
		text={contenedor},
		sort={contenedor},
		plural={contenedores},
		first={contenedor},
		firstplural={contenedores},
		description={Unidad de despliegue que empaqueta una aplicaci{\'o}n y sus dependencias con aislamiento a nivel de sistema operativo, usada para ejecutar servicios de forma uniforme en nube y borde.},
		type=\glsdefaulttype,
		symbol={📦},
		user1={Empaquetado},
		user2={Aislamiento},
		user3={Portabilidad},
		see={kubernetes}
	}
	
	\newglossaryentry{continuidad}
	{
		name={Continuidad},
		text={continuidad},
		sort={continuidad},
		plural={continuidades},
		first={continuidad},
		firstplural={continuidades},
		description={La \emph{continuidad} es la capacidad de un \gls{sistema} para mantener su operaci{\'o}n prevista sin interrupciones, degradaciones severas ni p{\'e}rdida de \glspl{dato}, incluso frente a fallos parciales, mantenimiento programado o variaciones de carga. En \gls{iot}, la continuidad exige estrategias de dise{\~n}o y explotaci{\'o}n como redundancia de \glspl{red} y \glspl{servidor}, replicaci{\'o}n de \gls{almacenamiento}, colas de mensajer{\'i}a y \gls{monitorizacion} proactiva, con el fin de sostener niveles de servicio y tiempos de respuesta acordes al \gls{tiemporeal}. Se relaciona estrechamente con la \gls{estabilidad}, la \gls{confiabilidad} y la disponibilidad del servicio, y se apoya en pr{\'a}cticas como \gls{escalabilidad} el{\'a}stica, conmutaci{\'o}n por error y recuperaci{\'o}n ante incidentes.},
		type=\glsdefaulttype,
		symbol={♾️},
		user1={Operaci{\'o}n ininterrumpida},
		user2={Continuidad de servicio},
		user3={Sostenibilidad operativa},
		see={estabilidad,confiabilidad,escalabilidad,infraestructura}
	}
	
	\newglossaryentry{continuonubeperiferia}
	{
		name={Continuo nube--periferia},
		text={continuo nube--periferia},
		sort={continuo nube periferia},
		plural={continuos nube--periferia},
		first={continuo nube--periferia},
		firstplural={continuos nube--periferia},
		description={Organizaci{\'o}n de funciones de c{\'o}mputo y almacenamiento desde dispositivos y pasarelas hasta servicios de nube, con decisiones de ubicaci{\'o}n guiadas por \gls{latencia}, costo y requisitos de operaci{\'o}n.},
		type=\glsdefaulttype,
		symbol={↔️},
		user1={Distribuci{\'o}n de funciones},
		user2={Ubicaci{\'o}n por requisitos},
		user3={Integraci{\'o}n extremo a extremo},
		see={edgecomputing,cloudcomputing,fogcomputing}
	}
	
	\newglossaryentry{contratoservicio}
	{
		name={Contrato de servicio},
		text={contrato de servicio},
		sort={contrato de servicio},
		plural={contratos de servicio},
		first={contrato de servicio},
		firstplural={contratos de servicio},
		description={Un \emph{contrato de servicio} especifica la interfaz, los mensajes, las reglas de uso, los errores y la versi{\'o}n de un \gls{servicio}, de manera que consumidores y proveedores puedan integrarse sin depender de detalles internos de implementaci{\'o}n.},
		type=\glsdefaulttype,
		symbol={📜},
		user1={Interfaz publicada},
		user2={Versionado},
		user3={Pol{\'\i}ticas},
		see={soa,wsdl,api}
	}
	
	\newglossaryentry{control}{
		name={Control},
		text={control},
		sort={control},
		plural={controles},
		first={control},
		firstplural={controles},
		description={Proceso mediante el cual un \gls{sistema} regula su comportamiento, ajusta sus operaciones o impone restricciones con el fin de alcanzar objetivos definidos o mantener condiciones deseadas. El \emph{control} puede ser autom{\'a}tico, a trav{\'e}s de algoritmos, \glspl{sensor} y \glspl{actuador} en bucles de retroalimentaci{\'o}n, o manual, a trav{\'e}s de la intervenci{\'o}n humana. En ingenier{\'i}a de software, el control tambi{\'e}n implica la gesti{\'o}n de acceso, autorizaci{\'o}n y trazabilidad. En entornos \gls{iot} y \glspl{sibiot}, el control asegura que los dispositivos inteligentes act{\'u}en en funci{\'o}n de datos en \gls{tiemporeal}, garantizando eficiencia, \gls{seguridad} y cumplimiento de pol{\'i}ticas organizacionales.},
		type=main,
		symbol={},
		user1={Regulaci{\'o}n de procesos},
		user2={Mecanismo de gesti{\'o}n de acceso y seguridad},
		user3={Retroalimentaci{\'o}n en SIbIoT},
		see={automatizacion,seguridad,monitorizacion,sibiot}
	}
	
	\newglossaryentry{controlabilidad}{
		name={Controlabilidad},
		text={controlabilidad},
		sort={controlabilidad},
		plural={controlabilidades},
		first={controlabilidad},
		firstplural={controlabilidades},
		description={Propiedad de un proceso f{\'\i}sico o t{\'e}cnico que indica si puede ser conducido hacia estados deseados mediante acciones de control aplicadas por actuadores. En \glspl{cps}, la controlabilidad condiciona qu{\'e} objetivos pueden imponerse (p.\,ej., confort, calidad, seguridad) y con qu{\'e} restricciones de operaci{\'o}n, ya que el software solo puede influir el sistema cuando existen variables manipulables y un margen de acci{\'o}n compatible con la din{\'a}mica del proceso.},
		type=main,
		symbol={🎛},
		user1={Acci{\'o}n de control},
		user2={Variables manipulables},
		user3={Objetivos del proceso},
		see={control,actuador,seguridad}
	}
	
	\newglossaryentry{controlacceso}{
		name={Control de acceso},
		text={control de acceso},
		sort={control de acceso},
		plural={controles de acceso},
		first={control de acceso},
		firstplural={controles de acceso},
		description={Conjunto de mecanismos, pol{\'i}ticas y tecnolog{\'i}as destinados a regular qui{\'e}n o qu{\'e} puede acceder a recursos, instalaciones, sistemas o datos, y bajo qu{\'e} condiciones. El \emph{control de acceso} garantiza la protecci{\'o}n de la confidencialidad, integridad y disponibilidad de los activos, mediante t{\'e}cnicas de \gls{autenticacion}, \gls{autorizacion} y \gls{identificacion}. En el contexto de los \glspl{sibiot}, el control de acceso se implementa tanto en el plano f{\'i}sico (por ejemplo, mediante \gls{reconocimientofacial}, huellas dactilares, tarjetas RFID o c{\'o}digos de acceso) como en el plano l{\'o}gico (control de permisos en plataformas, servicios distribuidos y bases de datos), contribuyendo a la \gls{seguridad} y a la confianza del sistema.},
		type=main,
		symbol={🔐},
		user1={Restricci{\'o}n y regulaci{\'o}n de accesos},
		user2={Mecanismos f{\'i}sicos y l{\'o}gicos de seguridad},
		user3={Aplicaci{\'o}n en entornos IoT y organizacionales},
		see={seguridad,identificacion,autenticacion,autorizacion,reconocimientofacial,sibiot}
	}
	\newglossaryentry{controlcalidad}{
		name={Control de calidad},
		text={control de calidad},
		sort={control de calidad},
		plural={controles de calidad},
		first={control de calidad},
		firstplural={controles de calidad},
		description={Conjunto de procedimientos, actividades y t{\'e}cnicas orientadas a verificar que un producto, servicio o proceso cumple con los requisitos, especificaciones y est{\'a}ndares previamente definidos. El \emph{control de calidad} incluye inspecci{\'o}n, pruebas, auditor{\'i}as y m{\'e}tricas que permiten detectar defectos, prevenir errores y garantizar la satisfacci{\'o}n del usuario final. En ingenier{\'i}a de software y en los \glspl{sibiot}, implica la validaci{\'o}n de \glspl{sensor}, la integridad de datos, la revisi{\'o}n de c{\'o}digo, pruebas de \gls{fiabilidad} y \gls{rendimiento}, as{\'i} como el cumplimiento de \glspl{buenaspracticas}, con el fin de asegurar soluciones seguras, sostenibles y eficientes.},
		type=main,
		symbol={✅},
		user1={Verificaci{\'o}n de conformidad},
		user2={Prevenci{\'o}n y detecci{\'o}n de defectos},
		user3={Garant{\'i}a de calidad en SIbIoT},
		see={fiabilidad,rendimiento,buenaspracticas,sostenibilidad}
	}
	
	\newglossaryentry{controlinteligente}
	{
		name={Control inteligente},
		text={control inteligente},
		sort={control inteligente},
		plural={controles inteligentes},
		first={control inteligente},
		firstplural={controles inteligentes},
		description={El \emph{control inteligente} es un enfoque de control que emplea t{\'e}cnicas de \gls{aprendizajeautomatico}, l{\'o}gica difusa, redes neuronales u otros m{\'e}todos adaptativos para gestionar sistemas complejos, no lineales o inciertos. Su objetivo es mejorar la \gls{tomadecisiones} en \gls{tiemporeal}, optimizar el \gls{rendimiento} del \gls{sistema} y adaptarse a cambios din{\'a}micos en el \gls{entorno} o en el propio proceso.},
		type=\glsdefaulttype,
		symbol={🤖},
		user1={Control basado en IA},
		user2={Control adaptativo},
		user3={Control inteligente artificial},
		see={automatizacion}
	}
	
	\newglossaryentry{controlador}
	{
		name={Controlador},
		text={controlador},
		sort={controlador},
		plural={controladores},
		first={controlador},
		firstplural={controladores},
		description={Un \emph{controller} es un componente que gestiona las solicitudes entrantes en un \gls{sistemainformatico}, coordinando la \gls{logicanegocios} y la respuesta generada. En arquitecturas como \gls{mvc}, el controller sirve de intermediario entre la vista y el modelo, garantizando la separaci{\'o}n de responsabilidades y la \gls{modularidad} del \gls{software}.},
		type=\glsdefaulttype,
		symbol={🎮},
		user1={Controlador},
		user2={Manejador de solicitudes},
		user3={Componente coordinador},
		see={logicanegocios}
	}
	
	\newglossaryentry{controlgranular}{
		name={Control granular},
		text={control granular},
		sort={control granular},
		plural={controles granulares},
		first={control granular},
		firstplural={controles granulares},
		description={Mecanismo de gesti{\'o}n y seguridad que permite definir permisos, pol{\'i}ticas o reglas con un alto nivel de detalle sobre recursos, datos, operaciones o usuarios dentro de un sistema. El control granular posibilita restringir o habilitar acciones espec{\'i}ficas (lectura, escritura, modificaci{\'o}n, ejecuci{\'o}n) sobre elementos concretos, en lugar de aplicar autorizaciones globales. En sistemas \gls{iot} y plataformas empresariales, el control granular es esencial para garantizar \gls{seguridad}, \gls{integridad}, trazabilidad y cumplimiento normativo, especialmente cuando m{\'u}ltiples actores interact{\'u}an con \glspl{recurso} distribuidos.},
		type=main,
		symbol={🔐},
		user1={Gesti{\'o}n detallada de permisos},
		user2={Pol{\'i}ticas de acceso espec{\'i}ficas},
		user3={Seguridad basada en roles y atributos},
		see={seguridad,controlacceso,autenticacionmultifactor,integridad}
	}
	
	\newglossaryentry{controller}{
		name={Controller},
		text={\textit{controller}},
		sort={controller},
		plural={controllers},
		first={\textit{controller}},
		firstplural={\textit{controllers}},
		description={Componente de una arquitectura de software que gestiona solicitudes, orquesta la l{\'o}gica de flujo y coordina la interacci{\'o}n entre vista y modelo o entre clientes y servicios (p.\,ej., en MVC o APIs REST).},
		type=main,
		symbol={🧭},
		user1={Orquestaci{\'o}n},
		user2={Interfaz de entrada},
		user3={MVC/REST},
		see={plataformaweb,servicioweb}
	}
	
	\newglossaryentry{controlversiones}{
		name={Control de versiones},
		text={control de versiones},
		sort={control de versiones},
		first={control de versiones},
		description={Conjunto de pr{\'a}cticas y herramientas para registrar y gestionar cambios en artefactos (c{\'o}digo, configuraciones y documentos), habilitando historial, colaboraci{\'o}n, fusiones y trazabilidad.},
		type=main,
		symbol={🧬},
		user1={Historial de cambios},
		user2={Colaboraci{\'o}n},
		user3={Trazabilidad},
		see={documentacion,refactorizacion}
	}
	
	\newglossaryentry{convencionmarcastiempo}{
		name={Convenci{\'o}n de marcas de tiempo},
		text={convenci{\'o}n de marcas de tiempo},
		sort={convencion marcas tiempo},
		plural={convenciones de marcas de tiempo},
		first={convenci{\'o}n de marcas de tiempo},
		firstplural={convenciones de marcas de tiempo},
		description={Conjunto de reglas y formatos normalizados utilizados para representar fechas y horas en registros de datos, garantizando coherencia temporal en sistemas distribuidos. Incluye est{\'a}ndares como ISO 8601, uso de zona horaria expl{\'i}cita, tiempo universal coordinado (UTC) y precisi{\'o}n en milisegundos o nanosegundos. En \glspl{sibiot} y arquitecturas distribuidas, la convenci{\'o}n de marcas de tiempo es cr{\'i}tica para correlacionar eventos, realizar \gls{procesamientoanalitico}, sincronizar nodos mediante protocolos como \gls{ntp} o \gls{ptp}, y asegurar integridad temporal en auditor{\'i}as y trazabilidad.},
		type=main,
		symbol={⏱️},
		user1={Estandarizaci{\'o}n temporal},
		user2={Sincronizaci{\'o}n distribuida},
		user3={Base para correlaci{\'o}n de eventos},
		see={ntp,ptp,evento,trazabilidad}
	}
	
	\newglossaryentry{cortafuego}{
		name={Cortafuego},
		text={cortafuego},
		sort={cortafuego},
		plural={cortafuegos},
		first={cortafuego},
		firstplural={cortafuegos},
		description={Mecanismo de \gls{seguridad} que controla el tr{\'a}fico de red entrante y saliente en un \gls{sistema} o infraestructura, aplicando un conjunto de reglas definidas para permitir, restringir o bloquear comunicaciones seg{\'u}n criterios establecidos. Un \emph{cortafuego} puede implementarse en hardware, software o de forma h{\'i}brida, y protege frente a accesos no autorizados, ataques externos y fugas de informaci{\'o}n. En el contexto de los \glspl{sibiot}, los cortafuegos son esenciales para resguardar \glspl{sensor}, \glspl{actuador} y \glspl{nodo} conectados a redes inal{\'a}mbricas o distribuidas, garantizando la \gls{confidencialidad}, \gls{integridad} y \gls{disponibilidad} de los datos intercambiados.},
		type=main,
		symbol={🛡️🔥},
		user1={Filtro de tr{\'a}fico en redes},
		user2={Defensa contra accesos no autorizados},
		user3={Componente clave de seguridad en IoT},
		see={seguridad,confidencialidad,integridad,disponibilidad,red,sibiot}
	}
	
	\newglossaryentry{cose}{
		name={COSE},
		text={COSE},
		plural={COSE},
		sort={cose},
		type=main,
		first={Firma y cifrado de objetos CBOR (\textit{CBOR Object Signing and Encryption}, COSE)},
		firstplural={CBOR Object Signing and Encryption (COSE)},
		plural={COSE},
		description={Est{\'a}ndar definido por el IETF para proporcionar mecanismos de firma digital, autenticaci{\'o}n y cifrado de datos estructurados en formato \gls{cbor}. COSE est{\'a} dise{\~n}ado para entornos con recursos limitados y se utiliza ampliamente en sistemas \gls{iot}, seguridad de mensajes, intercambio de claves y protocolos criptogr{\'a}ficos eficientes.},
		see={seguridad,confidencialidad,integridad,cbor}
	}
	
	\newglossaryentry{cps}{
		name={Sistema Ciberf{\'i}sico},
		text={CPS},
		sort={sistema ciberfisico},
		plural={CPSs},
		first={sistema ciberf{\'i}sico (Cyber-Physical System, CPS)},
		firstplural={sistemas ciberf{\'i}sicos (CPS)},
		description={\Gls{sistemaintegrado} que combina componentes f{\'i}sicos, procesos computacionales y redes de comunicaci{\'o}n para interactuar de manera coordinada con el entorno f{\'i}sico en tiempo real. Un sistema ciberf{\'i}sico integra sensores, actuadores, algoritmos de control, infraestructura de red y plataformas de procesamiento, permitiendo la supervisi{\'o}n, an{\'a}lisis y toma de decisiones autom{\'a}tica sobre procesos f{\'i}sicos. Los CPS constituyen la base tecnol{\'o}gica de dominios como la manufactura inteligente, veh{\'i}culos aut{\'o}nomos, redes el{\'e}ctricas inteligentes, salud digital y soluciones avanzadas de \gls{iot}.},
		type=main,
		symbol={⚙️},
		user1={Integraci{\'o}n f{\'i}sico-digital},
		user2={Control en tiempo real},
		user3={Base de Industria 4.0 e IoT},
		see={iot,gemelodigital,retroalimentacion,dispositivoembebido}
	}
	
	\newglossaryentry{cpu}
	{
		name={CPU},
		text={CPU},
		sort={cpu},
		plural={CPUs},
		first={unidad central de procesamiento (\textit{Central Processing Unit}, CPU)},
		firstplural={CPUs},
		description={La \emph{CPU} es el componente principal de un \gls{sistemainformatico}, responsable de ejecutar instrucciones y procesar \glspl{dato}. Coordina y controla las operaciones b{\'a}sicas del \gls{hardware}, siendo esencial para el funcionamiento de cualquier \gls{dispositivo} computacional.},
		type=\glsdefaulttype,
		symbol={🖥️},
		user1={Procesador},
		user2={Unidad de procesamiento},
		user3={Central Processing Unit},
		see={hardware}
	}
	
	\newglossaryentry{criptografia}{
		name={Criptograf{\'i}a},
		text={criptograf{\'i}a},
		sort={criptografia},
		plural={criptograf{\'i}as},
		first={criptograf{\'i}a},
		firstplural={t{\'e}cnicas de criptograf{\'i}a},
		description={Disciplina de la seguridad inform{\'a}tica que estudia y desarrolla m{\'e}todos matem{\'a}ticos para garantizar la confidencialidad, integridad, autenticidad y no repudio de la informaci{\'o}n. La \emph{criptograf{\'i}a} incluye algoritmos de cifrado sim{\'e}trico y asim{\'e}trico, firmas digitales, funciones hash y protocolos de intercambio de claves. En sistemas \gls{iot} y \glspl{sibiot}, la criptograf{\'i}a es esencial para proteger datos sensibles, comunicaciones de dispositivos, control de accesos y cumplimiento normativo frente a ciberamenazas.},
		type=main,
		symbol={🔐},
		user1={Cifrado de datos},
		user2={Seguridad criptogr{\'a}fica},
		user3={Protecci{\'o}n de informaci{\'o}n},
		see={seguridad,privacidad,tls,ssl}
	}
	
	\newglossaryentry{criptografiaasimetrica}
	{
		name={Criptograf{\'i}a asim{\'e}trica},
		text={criptograf{\'i}a asim{\'e}trica},
		sort={criptografia asimetrica},
		plural={criptograf{\'i}as asim{\'e}tricas},
		first={criptograf{\'i}a asim{\'e}trica},
		firstplural={criptograf{\'i}as asim{\'e}tricas},
		description={La \emph{criptograf{\'i}a asim{\'e}trica} es una t{\'e}cnica criptogr{\'a}fica que emplea un par de claves: una p{\'u}blica para \gls{cifrado} y otra privada para descifrado. Se usa ampliamente en comunicaciones seguras, firma digital y protocolos de intercambio de claves en \glspl{red}.},
		type=\glsdefaulttype,
		symbol={🔐},
		user1={Cifrado asim{\'e}trico},
		user2={Criptograf{\'i}a de clave p{\'u}blica},
		user3={Public key cryptography},
		see={criptografiasimetrica}
	}
	
	\newglossaryentry{criptografiasimetrica}
	{
		name={Criptograf{\'i}a sim{\'e}trica},
		text={criptograf{\'i}a sim{\'e}trica},
		sort={criptografia simetrica},
		plural={criptograf{\'i}as sim{\'e}tricas},
		first={criptograf{\'i}a sim{\'e}trica},
		firstplural={criptograf{\'i}as sim{\'e}tricas},
		description={La \emph{criptograf{\'i}a sim{\'e}trica} es una t{\'e}cnica criptogr{\'a}fica que utiliza la misma clave para \gls{cifrado} y descifrado de \glspl{dato}. Se caracteriza por ser r{\'a}pida y eficiente, aunque plantea retos en la distribuci{\'o}n segura de la clave compartida.},
		type=\glsdefaulttype,
		symbol={🔑},
		user1={Cifrado sim{\'e}trico},
		user2={Clave compartida},
		user3={Secret key cryptography},
		see={criptografiaasimetrica}
	}
	
	\newglossaryentry{crm}{
		name={CRM},
		text={CRM},
		plural={CRMs},
		first={Gesti{'o}n de Relaciones con Clientes (\textit{Customer Relationship Management}, CRM)},
		firstplural={CRMs},
		see={erp,gestion},
		sort={crm},
		description={Sistema de gesti{\'o}n de relaciones con clientes que integra procesos, datos y tecnolog{\'i}as para administrar interacciones comerciales, ventas, marketing y servicio al cliente. En un SIbIoT, el CRM consume resultados anal{\'i}ticos y alertas consolidadas para apoyar decisiones comerciales, personalizaci{\'o}n de ofertas y seguimiento del desempe{\~n}o del cliente}
	}
	
	\newglossaryentry{csmaca}{
		name={CSMA/CA},
		text={CSMA/CA},
		first={acceso m{\'u}ltiple con detecci{\'o}n de portadora y prevenci{\'o}n de colisiones (\textit{Carrier Sense Multiple Access with Collision Avoidance}, CSMA/CA)},
		firstplural={CSMAs/CA},
		sort={csma ca},
		description={\emph{CSMA/CA} es un protocolo de acceso al medio utilizado en \glspl{red} inal{\'a}mbricas, donde los \glspl{dispositivo} escuchan el canal antes de transmitir para evitar colisiones. A diferencia de CSMA/CD, intenta prevenir colisiones mediante espera aleatoria y confirmaciones, siendo fundamental en tecnolog{\'i}as como Wi-Fi (IEEE 802.11).}
	}
	
	\newglossaryentry{csmacd}{
		name={CSMA/CD},
		text={CSMA/CD},
		plural={CSMAs/CD},
		first={acceso m{\'u}ltiple con detecci{\'o}n de portadora y detecci{\'o}n de colisiones (\textit{Carrier Sense Multiple Access with Collision Detection}, CSMA/CD)},
		firstplural={CSMAs/CD},
		sort={csma cd},
		description={\emph{CSMA/CD} es un protocolo de acceso al medio utilizado en \glspl{red} Ethernet cl{\'a}sicas, donde los \glspl{dispositivo} escuchan el canal antes de transmitir y, en caso de detectar una colisi{\'o}n, interrumpen el env{\'i}o y reintentan tras un tiempo aleatorio. Forma parte del est{\'a}ndar IEEE 802.3.}
	}
	
	\newglossaryentry{csa}{
		name={Connectivity Standards Alliance},
		text={CSA},
		sort={connectivity standards alliance},
		plural={CSAs},
		first={alianza de est{\'a}ndares de conectividad (\textit{Connectivity Standards Alliance}, CSA)},
		firstplural={organizaciones CSA},
		description={Alianza que desarrolla est{\'a}ndares abiertos para interoperabilidad de \glspl{dispositivo} inteligentes; promueve \gls{matter}.},
		type=main,
		symbol={🌐},
		see={matter,interoperabilidad,estandar},
		user1={Ecosistema abierto},
		user2={Compatibilidad},
		user3={Gobernanza de est{\'a}ndares}
	}
	
	\newglossaryentry{css}
	{
		name={CSS},
		text={CSS},
		sort={css},
		plural={CSSs},
		first={hojas de estilo en cascada (\textit{Cascade Style Sheets}, CSS)},
		firstplural={CSSs},
		description={\emph{CSS} es un lenguaje utilizado para definir la apariencia y el dise{\~n}o de p{\'a}ginas web e \glspl{interfaz} de usuario escritas en \gls{html} o \gls{xhtml}. Permite controlar aspectos como colores, tipograf{\'i}a, distribuci{\'o}n y adaptabilidad del contenido. En conjunto con \gls{javascript}, se emplea para crear aplicaciones web y m{\'o}viles interactivas y responsivas.},
		type=\glsdefaulttype,
		symbol={🎨},
		user1={Hojas de estilo},
		user2={Cascade Style Sheets},
		user3={Lenguaje de dise{\~n}o web},
		see={html}
	}
	
	\newglossaryentry{dashboard}
	{
		name={Dashboard},
		text={dashboard},
		sort={dashboard},
		plural={dashboards},
		first={dashboard},
		firstplural={dashboards},
		description={Un \textit{\emph{dashboard}} es una interfaz visual que integra y presenta informaci{\'o}n relevante de un \gls{sistema}, permitiendo el \gls{monitorizacion}o y an{\'a}lisis en tiempo real. En los \glspl{sibiot}, los dashboards permiten visualizar datos recolectados por \glspl{sensor}, gestionar \glspl{dispositivo}, supervisar indicadores de desempe{\~n}o y facilitar la \gls{tomadecisiones} mediante gr{\'a}ficos, paneles e indicadores interactivos.},
		type=\glsdefaulttype,
		symbol={📊},
		user1={Tablero de control},
		user2={Panel de monitorizaci{\'o}n},
		user3={Cuadro de mando},
		see={servidorweb}
	}
	
	\newglossaryentry{datadriven}{
		name={Data-driven},
		text={\textit{data-driven}},
		sort={data driven},
		plural={\textit{data-driven}},
		first={\textit{Data-Driven}},
		firstplural={\textit{data-driven}},
		description={Enfoque de diseño, gestión o toma de decisiones basado en el análisis sistemático de datos objetivos en lugar de intuiciones o suposiciones. En sistemas \gls{iot} y \glspl{sibiot}, el enfoque data-driven implica recopilar, procesar y analizar datos generados por dispositivos y procesos organizacionales para optimizar operaciones, ajustar modelos y mejorar resultados mediante evidencia cuantificable. Este paradigma favorece decisiones trazables, medibles y adaptativas en entornos complejos.},
		type=main,
		symbol={📊},
		user1={Toma de decisiones basada en datos},
		user2={Optimización sustentada en evidencia},
		user3={Análisis cuantitativo sistemático},
		see={analitica,procesamientodatos,tomadecisiones,bi}
	}
	
	\newglossaryentry{dataflow}{
		name={Dataflow},
		text={\textit{Dataflow}},
		sort={dataflow},
		first={\textit{Dataflow}},
		description={Modelo y servicio de procesamiento distribuido de datos que permite la ejecuci{\'o}n de flujos de procesamiento tanto en modo continuo (\textit{stream processing}) como por lotes (\textit{batch processing}). En el contexto de Google Cloud, \emph{Dataflow} implementa el modelo de programaci{\'o}n Apache Beam y proporciona ejecuci{\'o}n escalable, tolerante a fallos y gestionada de \glspl{pipeline} de datos. En arquitecturas \glspl{sibiot}, Dataflow se emplea para el procesamiento avanzado de grandes vol{\'u}menes de eventos IoT, transformaci{\'o}n de flujos de datos, agregaci{\'o}n temporal, c{\'a}lculo de m{\'e}tricas y soporte a anal{\'i}tica operativa y predictiva},
		type=main,
		symbol={🔁},
		see={procesamientodistribuido,procesamientocontinuo,procesamientobatch,orientadoeventos,cloudcomputing,sibiot}
	}
	
	\newglossaryentry{datalake}{
		name={Data lake},
		text={data lake},
		sort={datalake},
		plural={data lakes},
		first={data lake},
		firstplural={data lakes},
		description={Repositorio centralizado dise{\~n}ado para almacenar grandes vol{\'u}menes de datos en su formato original, ya sean estructurados, semiestructurados o no estructurados. Un data lake permite la ingesta de datos a gran escala sin un esquema previo definido, posponiendo la transformaci{\'o}n hasta el momento del an{\'a}lisis. En arquitecturas SIbIoT, los data lakes se utilizan para conservar datos brutos provenientes de dispositivos IoT, eventos y registros hist{\'o}ricos, sirviendo como base para procesos batch, an{\'a}lisis avanzado y aprendizaje autom{\'a}tico},
		type=main,
		symbol={🏞️},
		see={almacenamiento,analiticadatos,procesamientobatch,procesamientodistribuido,sibiot}
	}
	
	\newglossaryentry{datamart}{
		name={Data Mart},
		text={data mart},
		sort={datamart},
		plural={data marts},
		first={almac{\'e}n de datos tem{\'a}tico, (\textit{Data Mart})},
		firstplural={almacenes de datos tem{\'a}ticos, (\textit{Data Marts})},
		description={Subsistema especializado de un \gls{basedatos} anal{\'i}tico o \gls{datawarehouse}, orientado a un {\'a}rea funcional espec{\'i}fica de la organizaci{\'o}n, como ventas, finanzas o log{\'i}stica. Un \emph{data mart} contiene un subconjunto de datos relevantes, estructurados para facilitar el an{\'a}lisis r{\'a}pido mediante \gls{olap}, consultas agregadas o cuadros de mando. En entornos de \gls{bi}, \gls{bigdata} y \glspl{sibiot}, los data marts permiten reducir la complejidad del procesamiento global, optimizar el acceso a informaci{\'o}n especializada y mejorar la eficiencia en la toma de decisiones sectoriales o departamentales.},
		type=main,
		user1={Subconjunto tem{\'a}tico de un almac{\'e}n de datos},
		user2={Soporte anal{\'i}tico para un {\'a}rea funcional espec{\'i}fica},
		user3={Facilita el uso de OLAP, consultas y visualizaci{\'o}n r{\'a}pida},
		user4={Componente esencial de arquitecturas BI y Big Data},
		see={datawarehouse,bi,olap,bigdata,etl,basedatos,sibiot},
		symbol={📦}
	}
	
	\newglossaryentry{dataset}{
		name={dataset},
		text={dataset},
		plural={datasets},
		first={dataset (conjunto de datos)},
		sort={dataset},
		see={dato,basedatos,analisisdatos,aprendizajeautomatico},
		description={Colecci{\'o}n estructurada de datos relacionados entre s{\'i}, recopilados y organizados con un prop{\'o}sito espec{\'i}fico para su an{\'a}lisis, procesamiento o entrenamiento de modelos computacionales. Un dataset puede estar compuesto por registros, atributos y etiquetas, y es un elemento fundamental en disciplinas como la ciencia de datos, el aprendizaje autom{\'a}tico, la miner{\'i}a de datos y los sistemas basados en IoT, donde sirve como base para la extracci{\'o}n de conocimiento y la toma de decisiones.}
	}
	
	\newglossaryentry{datawarehouse}{
		name={Data Warehouse},
		text={data warehouse},
		sort={datawarehouse},
		plural={data warehouses},
		first={almac{\'e}n de datos corporativo, (\textit{Data Warehouse})},
		firstplural={almacenes de datos corporativos, (\textit{Data Warehouses})},
		description={Sistema centralizado de almacenamiento de datos que integra informaci{\'o}n procedente de m{\'u}ltiples fuentes heterog{\'e}neas para su an{\'a}lisis y explotaci{\'o}n estrat{\'e}gica. Un \emph{data warehouse} consolida los procesos de extracci{\'o}n, transformaci{\'o}n y carga (\gls{etl}) en un modelo multidimensional que soporta consultas anal{\'i}ticas, \gls{olap}, generaci{\'o}n de informes y cuadros de mando. En el contexto de \gls{bi}, \gls{bigdata} y los \glspl{sibiot}, permite la integraci{\'o}n hist{\'o}rica de grandes vol{\'u}menes de datos provenientes de \glspl{sensor}, \glspl{dispositivo} y sistemas transaccionales, favoreciendo la toma de decisiones basada en evidencia y la optimizaci{\'o}n operativa.},
		type=main,
		user1={Almac{\'e}n centralizado de datos hist{\'o}ricos},
		user2={Integraci{\'o}n de datos mediante procesos ETL},
		user3={Soporte a an{\'a}lisis multidimensional y generaci{\'o}n de informes},
		user4={N{\'u}cleo de las arquitecturas BI y Big Data},
		see={bi,datamart,etl,olap,oltp,bigdata,sibiot},
		symbol={🗄️}
	}
	
	\newglossaryentry{dato}
	{
		name={Dato},
		text={dato},
		sort={dato},
		plural={datos},
		first={dato},
		firstplural={datos},
		description={Un \emph{dato} es una representaci{\'o}n simb{\'o}lica, num{\'e}rica, alfanum{\'e}rica o de otro tipo, que por s{\'i} mismo carece de significado completo. Los datos constituyen la unidad b{\'a}sica de la informaci{\'o}n y, al ser procesados, almacenados o transmitidos, adquieren valor para la \gls{tomadecisiones}. En el contexto de los \glspl{sibiot}, los datos son recolectados por \glspl{sensor}, transmitidos mediante \glspl{sistemacomunicacion}, procesados en \gls{edgecomputing} o \gls{nube}, y convertidos en informaci{\'o}n {\'u}til para optimizar procesos en diversos dominios.},
		type=\glsdefaulttype,
		symbol={💾},
		user1={Unidad de informaci{\'o}n},
		user2={Registro},
		user3={Elemento informativo},
		see={procesamientodatos}
	}
	
	\newglossaryentry{datobruto}{
		name={Dato bruto},
		text={dato bruto},
		sort={datobruto},
		plural={datos brutos},
		first={dato bruto},
		firstplural={datos brutos},
		description={Informaci{\'o}n primaria capturada en su forma original a trav{\'e}s de sensores, encuestas, registros o dispositivos, sin aplicar ning{\'u}n tipo de validaci{\'o}n, depuraci{\'o}n o procesamiento. Los datos brutos son el insumo inicial para el \gls{procesamientodatos}, donde posteriormente se convierten en informaci{\'o}n estructurada y significativa. En sistemas \glspl{sibiot}, los datos brutos recabados en tiempo real permiten la detecci{\'o}n de eventos, la generaci{\'o}n de patrones y el soporte para algoritmos de an{\'a}lisis, pero requieren \gls{preprocesamiento} para eliminar redundancias, errores o ruido},
		type=main,
		symbol={📊},
		user1={Dato sin procesar},
		user2={Registro primario},
		user3={Informaci{\'o}n en crudo},
		see={procesamientodatos,dato,informacion}
	}
	
	% ===================== Dato cr{\'i}tico =====================
	\newglossaryentry{datocritico}{
		name={Dato cr{\'i}tico},
		text={dato cr{\'i}tico},
		sort={dato critico},
		plural={datos cr{\'i}ticos},
		first={dato cr{\'i}tico},
		firstplural={datos cr{\'i}ticos},
		description={Un \emph{dato cr{\'i}tico} es informaci{\'o}n cuya disponibilidad, \gls{integridad} y \gls{confidencialidad} resultan esenciales para el correcto funcionamiento de un sistema o la seguridad de sus usuarios. La p{\'e}rdida, corrupci{\'o}n o acceso no autorizado a un dato cr{\'i}tico puede ocasionar fallos graves, interrupci{\'o}n de servicios, riesgos financieros o amenazas a la vida humana. En \glspl{sibiot}, ejemplos de datos cr{\'i}ticos incluyen se{\~n}ales biom{\'e}dicas en salud, datos de control industrial, informaci{\'o}n de ciberseguridad y registros de infraestructura de \gls{transporteinteligente}.},
		type=main,
		symbol={⚠️},
		user1={Alta dependencia de disponibilidad e integridad},
		user2={Riesgo elevado ante corrupci{\'o}n o p{\'e}rdida},
		user3={Clave en salud, industria, transporte e IoT},
		see={dato,integridad,seguridad,privacidad,sibiot}
	}
	
	\newglossaryentry{datocrudo}{
		name={Dato crudo},
		text={dato crudo},
		plural={datos crudos},
		first={dato crudo},
		firstplural={datos crudos},
		see={dato,datocurado},
		sort={datocrudo},
		description={Dato capturado directamente desde sensores o sistemas de origen sin haber sido sometido a procesos de limpieza, validaci{\'o}n o transformaci{\'o}n. En SIbIoT, los datos crudos pueden contener ruido, valores at{\'i}picos, duplicados o p{\'e}rdidas, y se conservan principalmente para fines de trazabilidad, auditor{\'i}a y reprocesamiento cuando cambian supuestos anal{\'i}ticos o modelos}
	}
	
	\newglossaryentry{datocurado}{
		name={Dato curado},
		text={dato curado},
		plural={datos curados},
		first={dato curado},
		firstplural={datos curados},
		see={dato,entrenamientomodelos},
		sort={datocurado},
		description={Dato que ha sido sometido a procesos sistem{\'a}ticos de limpieza, filtrado, validaci{\'o}n, normalizaci{\'o}n y enriquecimiento con metadatos. En SIbIoT, los datos curados se diferencian de los datos crudos por su calidad, consistencia y adecuaci{\'o}n para an{\'a}lisis y visualizaci{\'o}n, manteniendo trazabilidad hacia el origen}
	}
	
	% ===================== Dato heterog{\'e}neo =====================
	\newglossaryentry{datoheterogeneo}{
		name={Dato heterog{\'e}neo},
		text={dato heterog{\'e}neo},
		sort={dato heterogeneo},
		plural={datos heterog{\'e}neos},
		first={dato heterog{\'e}neo},
		firstplural={datos heterog{\'e}neos},
		description={Un \emph{dato heterog{\'e}neo} es aquel que proviene de distintas fuentes, formatos o estructuras, lo que dificulta su integraci{\'o}n y procesamiento directo. Puede abarcar datos num{\'e}ricos, textuales, multimedia, geoespaciales o de sensores, con diferentes niveles de calidad, frecuencia y sem{\'a}ntica. En el contexto de \glspl{sibiot}, los datos heterog{\'e}neos son comunes debido a la diversidad de \glspl{sensor}, \glspl{actuador}, protocolos y plataformas, lo que requiere mecanismos de \gls{interoperabilidad}, normalizaci{\'o}n y \gls{etl} para garantizar consistencia y valor anal{\'i}tico.},
		type=main,
		symbol={🔀},
		user1={Proveniencia de m{\'u}ltiples fuentes y formatos},
		user2={Diversidad en sem{\'a}ntica, calidad y frecuencia},
		user3={Requiere normalizaci{\'o}n e interoperabilidad en IoT},
		see={dato,heterogeneidad,etl,interoperabilidad,sibiot}
	}
	
	\newglossaryentry{decision}{
		name={Decisi{\'o}n},
		text={decisi{\'o}n},
		plural={decisiones},
		first={decisi{\'o}n},
		firstplural={decisiones},
		see={captura,ingesta,procesamiento,analisis},
		sort={decision},
		description={Acci{\'o}n seleccionada a partir del an{\'a}lisis de datos, indicadores o recomendaciones generadas por el sistema. En SIbIoT, las decisiones pueden ser humanas o automatizadas, y deben quedar registradas junto con su contexto y resultados para permitir evaluaci{\'o}n y mejora continua}
	}
	
	\newglossaryentry{decisionarquitectonica}{
		name={Decisi{\'o}n arquitect{\'o}nica},
		text={decisi{\'o}n arquitect{\'o}nica},
		sort={decision arquitectonica},
		plural={decisiones arquitect{\'o}nicas},
		first={decisi{\'o}n arquitect{\'o}nica},
		firstplural={decisiones arquitect{\'o}nicas},
		description={Elecci{\'o}n estrat{\'e}gica que define aspectos estructurales fundamentales del sistema, tales como patrones, estilos, tecnolog{\'i}as base y mecanismos de integraci{\'o}n, impactando directamente calidad, escalabilidad y mantenibilidad.},
		type=main,
		symbol={🧭},
		user1={Elecci{\'o}n estructural},
		user2={Impacto en calidad},
		user3={Base de arquitectura},
		see={estructura,restriccion,modelo}
	}
	
	\newglossaryentry{deformacion}
	{
		name={Deformaci{\'o}n},
		text={deformaci{\'o}n},
		sort={deformacion},
		plural={deformaciones},
		first={Deformaci{\'o}n},
		description={Cambio de forma o dimensi{\'o}n que experimenta un cuerpo cuando se somete a una fuerza o carga externa. 
			La deformaci{\'o}n puede ser el{\'a}stica o pl{\'a}stica, dependiendo de si el material recupera su forma original una vez retirada la fuerza. 
			En el contexto de los \glspl{sensor}, la deformaci{\'o}n es la base para medir magnitudes f{\'i}sicas como la presi{\'o}n o el peso mediante transductores.},
		see=[Relacionado con:]{fuerza,elasticidad,material}
	}
	
	\newglossaryentry{deformacionmecanica}
	{
		name={Deformaci{\'o}n mec{\'a}nica},
		text={deformaci{\'o}n mec{\'a}nica},
		sort={deformacion mecanica},
		plural={deformaciones mec{\'a}nicas},
		first={Deformaci{\'o}n mec{\'a}nica},
		description={Variaci{\'o}n estructural o dimensional que sufre un material como respuesta a una carga o esfuerzo aplicado. 
			Esta deformaci{\'o}n se cuantifica en funci{\'o}n del esfuerzo y la resistencia del material, siendo un par{\'a}metro fundamental en el dise{\~n}o de estructuras y dispositivos de medici{\'o}n. 
			En los sistemas de \gls{sibiot}, la deformaci{\'o}n mec{\'a}nica se traduce en una \gls{senal} el{\'e}ctrica a trav{\'e}s de celdas de carga o galgas extensiom{\'e}tricas para evaluar peso, tensi{\'o}n o presi{\'o}n en \gls{tiemporeal}.},
		see=[Relacionado con:]{esfuerzo,celdacarga,transductor}
	}
	
	\newglossaryentry{deduplicacion}{
		name={deduplicaci{\'o}n},
		text={deduplicaci{\'o}n},
		sort={deduplicacion},
		plural={deduplicaciones},
		first={deduplicaci{\'o}n},
		firstplural={deduplicaciones},
		description={La deduplicaci{\'o}n es una pr{\'a}ctica esencial en la \gls{gestioninformacion} y la \gls{calidaddatos}. Consiste en detectar entradas repetidas o equivalentes en bases de datos, asegurando la unicidad y coherencia de los registros. En los \glspl{sibiot}, evita la redundancia generada por m{\'u}ltiples \glspl{sensor} que reportan la misma informaci{\'o}n, optimizando el \gls{almacenamiento} y el rendimiento del sistema. Se implementa mediante algoritmos de coincidencia y comparaci{\'o}n de valores.},
		type=\glsdefaulttype,
		symbol={♻️},
		user1={Data deduplication},
		user2={Eliminaci{\'o}n de registros duplicados en un sistema de informaci{\'o}n},
		user3={M{\'e}todo para optimizar almacenamiento y consistencia de datos},
		see={depuracion,normalizacion,almacenamiento,calidaddatos}
	}
	
	\newglossaryentry{definitiondone}{
		name={Definition of Done},
		text={\textit{Definition of Done}},
		sort={\textit{definition of done}},
		plural={\textit{Definitions of Done}},
		first={\textit{Definition of Done}},
		firstplural={\textit{Definitions of Done}},
		description={Conjunto formal de criterios que determinan cu{\'a}ndo un incremento cumple los est{\'a}ndares de calidad requeridos y puede considerarse terminado. Garantiza transparencia, consistencia y calidad en cada entrega realizada por el equipo Scrum.},
		type=main,
		symbol={✅},
		user1={Criterios de finalizaci{\'o}n},
		user2={Est{\'a}ndar de calidad},
		user3={Base para liberar incrementos},
		see={incremento,sprint,sprintgoal}
	}
	
	\newglossaryentry{definitionofready}{
		name={Definition of Ready},
		text={\textit{Definition of Ready}},
		sort={definition of ready},
		plural={\textit{Definitions of Ready}},
		first={\textit{Definition of Ready}},
		firstplural={\textit{Definitions of Ready}},
		description={Conjunto de criterios que un elemento del Product Backlog debe cumplir antes de ser seleccionado para un Sprint. Busca asegurar claridad, estimabilidad y preparaci{\'o}n adecuada del trabajo antes de su desarrollo.},
		type=main,
		symbol={📦},
		user1={Criterios previos al Sprint},
		user2={Preparaci{\'o}n del trabajo},
		user3={Reducci{\'o}n de incertidumbre},
		see={productbacklog,sprint,sprintplanning}
	}
	
	\newglossaryentry{delete}{
		name={DELETE},
		text={DELETE},
		sort={delete},
		first={DELETE},
		description={M{\'e}todo del protocolo HTTP utilizado para eliminar un recurso identificado por una URL. En arquitecturas RESTful, \texttt{DELETE} es un m{\'e}todo idempotente y se emplea para remover entidades, registros o asociaciones existentes. En el contexto de sistemas de informaci{\'o}n basados en IoT, este m{\'e}todo se utiliza para la gesti{\'o}n del ciclo de vida de dispositivos, eliminaci{\'o}n de datos obsoletos o cancelaci{\'o}n de recursos l{\'o}gicos},
		type=main,
		see={http,rest,restful,api,gestion,seguridad}
	}
	
	\newglossaryentry{depuracion}{
		name={depuraci{\'o}n},
		text={depuraci{\'o}n},
		sort={depuracion},
		plural={depuraciones},
		first={depuraci{\'o}n},
		firstplural={depuraciones},
		description={La depuraci{\'o}n es una etapa fundamental en la \gls{gestioninformacion}, que tiene como objetivo eliminar datos redundantes, incompletos o err{\'o}neos provenientes de diversas fuentes. En los \glspl{sibiot}, la depuraci{\'o}n garantiza que los datos obtenidos de \glspl{sensor} y \glspl{dispositivo} reflejen el estado real del entorno, mejorando la \gls{calidaddatos} y la fiabilidad de los an{\'a}lisis posteriores. Este proceso suele combinar t{\'e}cnicas de validaci{\'o}n, \gls{deduplicacion}, \gls{normalizacion} y \gls{enriquecimientodatos}.},
		type=\glsdefaulttype,
		symbol={🧹},
		user1={Data cleansing},
		user2={Eliminaci{\'o}n de errores e inconsistencias en conjuntos de datos},
		user3={Procedimiento de limpieza y refinamiento de datos para an{\'a}lisis confiables},
		see={deduplicacion,normalizacion,enriquecimientodatos,calidaddatos,gestioninformacion}
	}
	
	\newglossaryentry{desacoplado}{
		name={Desacoplado},
		text={desacoplado},
		first={desacoplado},
		sort={desacoplado},
		plural={desacoplados},
		description={Que ha sido separado o independizado de otro elemento. En software, un sistema \emph{desacoplado} mantiene m{\'o}dulos o \glspl{servicio} con baja dependencia mutua. En \gls{iot}, se refiere a dispositivos o plataformas que interact{\'u}an mediante protocolos est{\'a}ndar sin requerir integraciones directas o personalizadas.},
		type=main,
		symbol={⚙️},
		user1={Sistema modular e independiente},
		user2={Servicios con baja dependencia},
		user3={Compatibilidad mediante est{\'a}ndares},
		see={desacoplar,interoperabilidad}
	}
	
	\newglossaryentry{desacoplamiento}{
		name={desacoplamiento},
		text={desacoplamiento},
		first={desacoplamiento},
		sort={desacoplamiento},
		description={Principio de dise{\~n}o arquitect{\'o}nico que consiste en reducir o eliminar las dependencias directas entre componentes de un sistema, de modo que cada uno pueda evolucionar, desplegarse y escalar de manera independiente. En sistemas distribuidos y SIbIoT, el desacoplamiento se logra mediante mecanismos como mensajer{\'i}a as{\'i}ncrona, arquitecturas orientadas a eventos, interfaces bien definidas y servicios independientes. Este enfoque mejora la escalabilidad, la tolerancia a fallos y la mantenibilidad del sistema},
		plural={desacoplamientos},
		firstplural={desacoplamientos},
		see={arquitecturaestratificada,enfoquereactivo,mensajeriatiemporeal,orientadoeventos}
	}
	
	\newglossaryentry{desacoplar}{
		name={Desacoplar},
		text={desacoplar},
		sort={desacoplar},
		description={Acci{\'o}n de separar o independizar elementos previamente acoplados. En ingenier{\'i}a de software, el \emph{desacoplamiento} implica reducir la dependencia entre \glspl{componente} o m{\'o}dulos, favoreciendo la \gls{modularidad}, la reutilizaci{\'o}n y el mantenimiento. En sistemas \gls{iot}, desacoplar significa que los dispositivos o servicios pueden operar de forma aut{\'o}noma y comunicarse mediante interfaces est{\'a}ndar.},
		plural={desacoplamientos},
		type=main,
		symbol={🔀},
		user1={Separaci{\'o}n de dependencias en software},
		user2={Independencia de servicios y dispositivos},
		user3={Mejora de modularidad y mantenibilidad},
		see={acoplar,acoplado,interoperabilidad}
	}
	
	\newglossaryentry{desafio}
	{
		name={Desaf{\'i}o},
		text={desaf{\'i}o},
		sort={desafio},
		plural={desaf{\'i}os},
		first={desaf{\'i}o},
		firstplural={desaf{\'i}os},
		description={Un \emph{desaf{\'i}o} es un conjunto de limitaciones, barreras t{\'e}cnicas, organizativas o contextuales que dificultan el dise{\~n}o, implementaci{\'o}n, operaci{\'o}n o escalamiento de sistemas \gls{iot}. Entre los principales desaf{\'i}os se encuentran la \gls{interoperabilidad}, la \gls{seguridad}, la \gls{privacidad} de los \glspl{dato}, el consumo energ{\'e}tico, la \gls{heterogeneidad} de los \glspl{dispositivoiot}, la \gls{escalabilidad} de la \gls{infraestructura} y la gesti{\'o}n eficiente de grandes vol{\'u}menes de \glspl{dato}.},
		type=\glsdefaulttype,
		symbol={⚠️},
		user1={Reto},
		user2={Limitaci{\'o}n},
		user3={Barrera},
		see={seguridad}
	}
	
	\newglossaryentry{desarrollo}{
		name={desarrollo},
		text={desarrollo},
		sort={desarrollo},
		plural={desarrollos},
		first={desarrollo},
		firstplural={desarrollos},
		description={El \emph{desarrollo} es el proceso sistem{\'a}tico mediante el cual se conciben, dise{\~n}an, implementan, prueban y perfeccionan sistemas, aplicaciones o componentes tecnol{\'o}gicos. En el contexto de los sistemas basados en \gls{iot}, el desarrollo implica la integraci{\'o}n de \glspl{hardware}, software, \glspl{protocolocomunicacion} y \glspl{servicio} en la \gls{nube} para crear soluciones funcionales, seguras y \glspl{escalable}. Este \gls{proceso} abarca desde la definici{\'o}n de \glspl{requisito} y la selecci{\'o}n de \glspl{tecnologia} hasta la \gls{validacion} y el \gls{mantenimiento} continuo, favoreciendo la \gls{interoperabilidad}, la \gls{eficiencia} y la \gls{automatizacion}.},
		type=main,
		symbol={🛠️},
		user1={Construcci{\'o}n y mejora de sistemas},
		user2={Integraci{\'o}n de hardware y software},
		user3={Proceso iterativo orientado a la calidad},
		see={impl,despliegue,metodologiaagil,integracion,solucioniot}
	}
	
	\newglossaryentry{desarrolloagilsoftware}
	{
		name={Desarrollo {\'a}gil de software},
		text={desarrollo {\'a}gil de software},
		sort={desarrollo agil software},
		plural={desarrollos {\'a}giles de software},
		first={desarrollo {\'a}gil de software},
		firstplural={desarrollos {\'a}giles de software},
		description={El \emph{desarrollo {\'a}gil de software} es un paradigma de ingenier{\'i}a del software que aplica los principios y valores del Manifiesto {\'A}gil para gestionar proyectos de forma flexible, colaborativa e incremental. Este enfoque se basa en ciclos cortos de desarrollo, entrega continua de versiones funcionales, retroalimentaci{\'o}n frecuente con los usuarios y adaptaci{\'o}n al cambio. En proyectos \gls{iot}, el desarrollo {\'a}gil permite integrar eficazmente los m{\'o}dulos l{\'o}gicos del \gls{sistema} con componentes f{\'i}sicos, facilitando la validaci{\'o}n temprana y la evoluci{\'o}n progresiva del producto.},
		type=\glsdefaulttype,
		symbol={⚡},
		user1={Agile software development},
		user2={Metodolog{\'i}a {\'a}gil},
		user3={Desarrollo incremental},
		see={cascada}
	}
	
	\newglossaryentry{desarrolloweb}{
		name={Desarrollo web},
		text={desarrollo web},
		sort={desarrolloweb},
		plural={desarrollos web},
		first={desarrollo web},
		firstplural={desarrollos web},
		description={Disciplina de la ingenier{\'\i}a de software orientada al dise{\~n}o, implementaci{\'o}n, despliegue y mantenimiento de aplicaciones y servicios accesibles a trav{\'e}s de navegadores web. En los \glspl{sibiot}, el desarrollo web resulta clave para la construcci{\'o}n de interfaces de usuario, paneles de monitoreo, servicios de integraci{\'o}n y APIs que permiten la visualizaci{\'o}n, gesti{\'o}n y explotaci{\'o}n de datos provenientes de dispositivos distribuidos en tiempo real o diferido.},
		type=main,
		symbol={🌐},
		user1={Aplicaciones web},
		user2={Interfaces de usuario},
		user3={Servicios y APIs},
		see={plataformaweb,servicioweb,arquitecturadistribuida,integracionsoftware}
	}
	
	\newglossaryentry{descubrimiento}{
		name={Descubrimiento},
		text={descubrimiento},
		sort={descubrimiento},
		plural={descubrimientos},
		first={descubrimiento},
		firstplural={descubrimientos},
		description={Proceso por el cual \glspl{dispositivo} y \glspl{servicio} se identifican y registran mutuamente para habilitar configuraci{\'o}n e interoperabilidad.},
		type=main,
		symbol={🧭},
		see={matter,provisionamiento},
		user1={Anuncio},
		user2={Registro},
		user3={Resoluci{\'o}n}
	}
	
	\newglossaryentry{desempeno}{
		name={Desempe{\~n}o},
		text={desempe{\~n}o},
		sort={desempeno},
		first={desempe{\~n}o},
		description={Medida del comportamiento de un sistema respecto a criterios como latencia, rendimiento, uso de recursos y capacidad de respuesta, bajo condiciones de carga definidas.},
		type=main,
		symbol={📈},
		user1={Latencia},
		user2={Uso de recursos},
		user3={Capacidad},
		see={rendimiento,anchobanda}
	}
	
	\newglossaryentry{despliegue}{
		name={despliegue},
		text={despliegue},
		sort={despliegue},
		plural={despliegues},
		first={despliegue},
		firstplural={despliegues},
		description={Proceso mediante el cual un sistema, aplicaci{\'o}n o componente de software es transferido desde su entorno de desarrollo hacia un entorno de prueba, preproducci{\'o}n o producci{\'o}n para su operaci{\'o}n real. El \emph{despliegue} implica automatizaci{\'o}n, configuraci{\'o}n de infraestructura, integraci{\'o}n con servicios existentes y verificaci{\'o}n de que el sistema funcione conforme a los requisitos operativos. En los \glspl{sibiot}, este proceso requiere considerar conectividad, seguridad, compatibilidad con dispositivos y garant{\'i}as de disponibilidad distribuida},
		type=main,
		symbol={🚀},
		user1={Entrega y activaci{\'o}n de software},
		user2={Automatizaci{\'o}n e infraestructura de ejecuci{\'o}n},
		user3={Operaci{\'o}n confiable en entornos distribuidos},
		see={ciclovida,sibiot,orquestacion}
	}
	
	\newglossaryentry{deteccion}{
		name={Detecci{\'o}n},
		text={detecci{\'o}n},
		sort={deteccion},
		plural={detecciones},
		first={detecci{\'o}n},
		firstplural={detecciones},
		description={Proceso mediante el cual un \gls{sistema} o dispositivo determina la presencia, ausencia o aparici{\'o}n de un objeto, se{\~n}al, evento o patr{\'o}n en el entorno. La \emph{detecci{\'o}n} no necesariamente implica reconocer la naturaleza o identidad del elemento observado, sino confirmar que existe una condici{\'o}n o fen{\'o}meno a partir de mediciones o evidencias. En \glspl{sibiot}, la detecci{\'o}n es realizada por \glspl{sensor} o algoritmos de an{\'a}lisis en \gls{tiemporeal}, y constituye la primera fase de procesos como el \gls{reconocimiento} o la \gls{identificacion}.},
		type=main,
		symbol={👁️},
		user1={Confirmaci{\'o}n de presencia o evento},
		user2={Etapa inicial de percepci{\'o}n},
		user3={Entrada a procesos de an{\'a}lisis avanzado},
		see={reconocimiento,identificacion,sensor,monitorizacion}
	}
	
	\newglossaryentry{deteccioncolisiones}
	{
		name={Detecci{\'o}n de colisiones},
		text={detecci{\'o}n de colisiones},
		sort={deteccion de colisiones},
		plural={detecciones de colisiones},
		first={detecci{\'o}n de colisiones},
		firstplural={detecciones de colisiones},
		description={La \emph{detecci{\'o}n de colisiones} es un mecanismo empleado en \glspl{red} de comunicaci{\'o}n compartidas para identificar la ocurrencia simult{\'a}nea de dos o m{\'a}s transmisiones en un mismo canal, lo que genera interferencia y p{\'e}rdida de \glspl{dato}. En sistemas \gls{iot}, es especialmente relevante en protocolos de acceso m{\'u}ltiple como \gls{csmacd} y \gls{csmaca}, cuyo objetivo es minimizar la congesti{\'o}n, optimizar el uso del canal y garantizar la eficiencia en entornos con numerosos \glspl{dispositivoconectado}.},
		type=\glsdefaulttype,
		symbol={📡},
		user1={Colision detection},
		user2={Detecci{\'o}n de interferencias},
		user3={Mecanismo de acceso al medio},
		see={csmacd}
	}
	
	
	\newglossaryentry{deteccioncolor}{
		name={Detecci{\'o}n de color},
		text={detecci{\'o}n de color},
		sort={deteccion de color},
		plural={detecciones de color},
		first={detecci{\'o}n de color},
		firstplural={detecciones de color},
		description={T{\'e}cnica de visi{\'o}n artificial que permite identificar y segmentar regiones de una imagen en funci{\'o}n de su informaci{\'o}n crom{\'a}tica. La \emph{detecci{\'o}n de color} utiliza modelos de representaci{\'o}n como RGB, HSV o LAB para filtrar y analizar pixeles seg{\'u}n rangos espec{\'i}ficos. En \glspl{sibiot}, se aplica en clasificaci{\'o}n de productos, se{\~n}alizaci{\'o}n industrial, monitorizaci{\'o}n ambiental y sistemas de guiado rob{\'o}tico, sirviendo como base para procesos de \gls{reconocimiento} y \gls{identificacion}.},
		type=main,
		symbol={🎨},
		user1={Procesamiento de im{\'a}genes},
		user2={Segmentaci{\'o}n basada en cromaticidad},
		user3={Soporte para reconocimiento e identificaci{\'o}n},
		see={deteccion,reconocimiento,visionartificial,sibiot}
	}
	
	\newglossaryentry{deteccionmovimiento}{
		name={Detecci{\'o}n de movimiento},
		text={detecci{\'o}n de movimiento},
		sort={deteccion de movimiento},
		plural={detecciones de movimiento},
		first={detecci{\'o}n de movimiento},
		firstplural={detecciones de movimiento},
		description={Proceso de identificar cambios en la posici{\'o}n de objetos o personas dentro de un espacio f{\'i}sico o digital, a partir de im{\'a}genes secuenciales, se{\~n}ales de radar, ultrasonido o sensores infrarrojos. La \emph{detecci{\'o}n de movimiento} puede basarse en diferencias de pixeles, an{\'a}lisis de flujo {\'o}ptico o algoritmos de aprendizaje. En los \glspl{sibiot}, se emplea en sistemas de seguridad, control de acceso, iluminaci{\'o}n autom{\'a}tica, monitorizaci{\'o}n de tr{\'a}fico y reconocimiento de patrones de actividad.},
		type=main,
		symbol={🏃},
		user1={An{\'a}lisis temporal de se{\~n}ales},
		user2={Visi{\'o}n artificial y sensores especializados},
		user3={Aplicaciones en seguridad y eficiencia energ{\'e}tica},
		see={deteccion,sensor,monitorizacion,visionartificial}
	}
	
	\newglossaryentry{deteccionobjetos}{
		name={Detecci{\'o}n de objetos},
		text={detecci{\'o}n de objetos},
		sort={deteccion de objetos},
		plural={detecciones de objetos},
		first={detecci{\'o}n de objetos},
		firstplural={detecciones de objetos},
		description={Proceso computacional orientado a identificar la presencia y localizaci{\'o}n de objetos dentro de una imagen o secuencia de video. La \emph{detecci{\'o}n de objetos} combina \glspl{tecnica} de \gls{visionartificial}, \gls{aprendizajeautomatico} y \glspl{redneuronal} convolucionales para ubicar regiones de inter{\'e}s y clasificarlas seg{\'u}n su tipo. En \glspl{sibiot}, esta t{\'e}cnica se emplea para aplicaciones como vigilancia inteligente, control de tr{\'a}fico, inventario automatizado, agricultura de precisi{\'o}n y asistencia en entornos dom{\'o}ticos.},
		type=main,
		symbol={📦},
		user1={Visi{\'o}n por computadora},
		user2={Localizaci{\'o}n y clasificaci{\'o}n de objetos},
		user3={Aplicaciones en seguridad y automatizaci{\'o}n},
		see={deteccion,reconocimiento,identificacion,aprendizajeautomatico}
	}
	
	\newglossaryentry{deteccionobstaculos}{
		name={Detecci{\'o}n de obst{\'a}culos},
		text={detecci{\'o}n de obst{\'a}culos},
		sort={deteccion obstaculos},
		plural={t{\'e}cnicas de detecci{\'o}n de obst{\'a}culos},
		first={detecci{\'o}n de obst{\'a}culos},
		firstplural={detecciones de obst{\'a}culos},
		description={Conjunto de t{\'e}cnicas y algoritmos utilizados para identificar la presencia de objetos que puedan interferir en el desplazamiento de un veh{\'i}culo, robot o sistema aut{\'o}nomo. La \emph{detecci{\'o}n de obst{\'a}culos} puede realizarse mediante sensores ultras{\'o}nicos, l{\'a}seres LiDAR, c{\'a}maras de visi{\'o}n artificial o combinaciones multimodales. En el contexto de los \glspl{sibiot}, es fundamental para aplicaciones de movilidad inteligente, veh{\'i}culos aut{\'o}nomos, drones y sistemas de seguridad industrial, asegurando navegaci{\'o}n eficiente y prevenci{\'o}n de colisiones.},
		type=main,
		symbol={🚧},
		user1={Prevenci{\'o}n de colisiones},
		user2={Sistemas de navegaci{\'o}n segura},
		user3={Tecnolog{\'i}a de visi{\'o}n artificial aplicada},
		see={visionartificial,reconocimientofacial,vehiculoautonomo}
	}
	
	\newglossaryentry{developer}{
		name={Developer},
		text={\textit{Developer}},
		sort={developer},
		plural={\textit{Developers}},
		first={\textit{Developer}},
		firstplural={\textit{Developers}},
		description={Miembro del Scrum Team responsable de crear el Incremento en cada Sprint. Los Developers poseen las habilidades necesarias para transformar los elementos del Product Backlog en funcionalidad terminada que cumple con la Definition of Done. Son responsables de la calidad t{\'e}cnica, la autoorganizaci{\'o}n, la planificaci{\'o}n del Sprint y la adaptaci{\'o}n diaria del trabajo para alcanzar el Sprint Goal.},
		type=main,
		symbol={👨‍💻},
		user1={Responsable de construir el Incremento},
		user2={Miembro autoorganizado del Scrum Team},
		user3={Cumplimiento de Definition of Done},
		see={sprint,incremento,productbacklog,sprintgoal,scrum}
	}
	
	\newglossaryentry{devops}{
		name={DevOps},
		text={DevOps},
		first={desarrollo y operaciones (\textit{Development and Operations}, DevOps)},
		firstplural={DevOps},
		sort={devops},
		type=main,
		description={\emph{DevOps} es una filosof{\'i}a y conjunto de pr{\'a}cticas que integran los equipos de desarrollo de software y operaciones de TI, con el objetivo de automatizar procesos, mejorar la colaboraci{\'o}n y acelerar la entrega continua de software confiable. Incluye metodolog{\'i}as como la integraci{\'o}n continua (CI), la entrega continua (CD), la monitorizaci{\'o}n automatizada y la infraestructura como c{\'o}digo (IaC). En el contexto de los \glspl{sibiot}, DevOps facilita la gesti{\'o}n de la \gls{infraestructura}, la \gls{escalabilidad} y el despliegue seguro de aplicaciones distribuidas.},
		see = {infraestructuracodigo,cicd}
	}
	
	\newglossaryentry{dht22}{
		name={DHT22},
		text={DHT22},
		short={DHT22},
		long={Sensor digital de temperatura y humedad (AM2302) de bus de un solo hilo},
		sort={dht22},
		plural={DHT22},
		first={DHT22},
		firstplural={DHT22},
		description={El DHT22 (AM2302) es un sensor digital que integra un elemento de medida de humedad y un termistor, entregando lecturas mediante un protocolo serial de un solo hilo y alimentaci{\'o}n t{\'i}pica entre 3.3V y 5.5V, con especificaciones de rango y precisi{\'o}n provistas por su hoja de datos.},
		see={sensor}
	}
	
	\newglossaryentry{diagrama}{
		name={Diagrama},
		text={diagrama},
		sort={diagrama},
		plural={diagramas},
		first={diagrama},
		firstplural={diagramas},
		description={Representaci{\'o}n gr{\'a}fica estructurada de elementos y relaciones de un sistema (p.\,ej., arquitectura, flujo de datos, casos de uso), utilizada para comunicar dise{\~n}os y facilitar el an{\'a}lisis.},
		type=main,
		symbol={🗺️},
		user1={Representaci{\'o}n gr{\'a}fica},
		user2={Comunicaci{\'o}n de dise{\~n}o},
		user3={Estructura},
		see={representacion,documentacion}
	}
	
	\newglossaryentry{diagramaflujo}{
		name={Diagrama de flujo},
		text={diagrama de flujo},
		sort={diagrama de flujo},
		plural={diagramas de flujo},
		first={diagrama de flujo},
		firstplural={diagramas de flujo},
		description={Representaci{\'o}n gr{\'a}fica de un \gls{algoritmo} o proceso, que utiliza s{\'i}mbolos normalizados (como rect{\'a}ngulos, rombos, flechas y paralelogramos) para mostrar las operaciones, decisiones, entradas y salidas, as{\'i} como la secuencia l{\'o}gica de ejecuci{\'o}n. Un \emph{diagrama de flujo} permite visualizar la l{\'o}gica de un sistema de manera clara e intuitiva, facilitando la comprensi{\'o}n, el an{\'a}lisis y la comunicaci{\'o}n entre usuarios, analistas y programadores. En el dise{\~n}o de \glspl{sibiot}, los diagramas de flujo se emplean para modelar procesos de monitorizaci{\'o}n, control y \gls{tomadecisiones}.},
		type=main,
		symbol={🔀},
		user1={Visualizaci{\'o}n de procesos},
		user2={Herramienta de dise{\~n}o y documentaci{\'o}n},
		user3={Representaci{\'o}n de algoritmos},
		see={algoritmo,pseudocodigo,modelado}
	}
	
	\newglossaryentry{diccionarioeventos}{
		name={Diccionario de eventos},
		text={diccionario de eventos},
		sort={diccionario eventos},
		plural={diccionarios de eventos},
		first={diccionario de eventos},
		firstplural={diccionarios de eventos},
		description={Documento o estructura formal que define y normaliza los eventos generados dentro de un sistema, especificando su nombre, tipo, significado, estructura de datos, origen, condiciones de activaci{\'o}n y sem{\'a}ntica asociada. En arquitecturas basadas en eventos y en \glspl{sibiot}, el diccionario de eventos garantiza consistencia sem{\'a}ntica, interoperabilidad y correcta interpretaci{\'o}n de mensajes entre componentes distribuidos, facilitando el \gls{procesamientoanalitico}, la auditor{\'i}a y la \gls{retroalimentacionoperacional}.},
		type=main,
		symbol={📖},
		user1={Normalizaci{\'o}n sem{\'a}ntica},
		user2={Base para interoperabilidad},
		user3={Soporte para arquitecturas orientadas a eventos},
		see={evento,interoperabilidad,procesamientoanalitico,retroalimentacionoperacional}
	}
	
	\newglossaryentry{digital}{
		name={Digital},
		text={digital},
		sort={digital},
		first={digital},
		description={Propiedad de una se{\~n}al, sistema o representaci{\'o}n basada en valores discretos, generalmente binarios. En sistemas \emph{digitales}, la informaci{\'o}n se representa mediante niveles bien definidos, como 0 y 1, lo que permite su procesamiento fiable por microcontroladores y computadoras. En sistemas embebidos y SIbIoT, las se{\~n}ales digitales son utilizadas para control, l{\'o}gica y comunicaci{\'o}n.},
		type=main,
		symbol={🔢},
		user1={Valores discretos},
		user2={Binario},
		user3={Procesamiento computacional},
		see={analogica,analogico,uart,spi}
	}
	
	\newglossaryentry{digitalizacion}{
		name={digitalizaci{\'o}n},
		text={digitalizaci{\'o}n},
		sort={digitalizacion},
		description={Proceso mediante el cual los objetos, se{\~n}ales, datos o procesos del mundo f{\'i}sico son transformados en representaciones digitales que pueden ser capturadas, almacenadas, procesadas y transmitidas por sistemas computacionales. En el contexto de los sistemas \gls{iot}, la digitalizaci{\'o}n se logra principalmente mediante el uso de \glspl{sensor} y \glspl{actuador}, permitiendo que los objetos cotidianos interact{\'u}en con infraestructuras de red, plataformas de procesamiento y servicios en la \gls{nube}.},
		type=main,
		see={iot,sensor,actuador,procesamientodatos,nube}
	}
	
	\newglossaryentry{disenosistemas}{
		name={Dise{\~n}o de sistemas},
		sort={diseno de sistemas},
		sort={diseno de sistemas},
		description={%
			Proceso sistem{\'a}tico mediante el cual se definen la estructura, los componentes, las interacciones y los principios de funcionamiento de un sistema con el fin de satisfacer un conjunto de requisitos funcionales y no funcionales.
			El dise{\~n}o de sistemas abarca decisiones relativas a la arquitectura, la distribuci{\'o}n de responsabilidades, los mecanismos de comunicaci{\'o}n, la gesti{\'o}n de datos y los controles de calidad, seguridad y rendimiento.
			En el contexto de los sistemas de informaci{\'o}n basados en IoT y SIbIoT, el dise{\~n}o de un sistema resulta cr{\'i}tico para integrar de manera coherente el entorno f{\'i}sico, el procesamiento distribuido, la mensajer{\'i}a en tiempo real y las capas de inteligencia y \gls{ux}.
		},
		first={dise{\~n}o de sistemas},
		plural={dise{\~n}os de sistemas},
		firstplural={dise{\~n}os de sistemas},
		see={arquitectura,arquitecturaestratificada,procesamientodistribuido,iot}
	}
	
	\newglossaryentry{displayport}{
		name={DisplayPort},
		text={DisplayPort},
		sort={displayport},
		plural={interfaces DisplayPort},
		first={\textit{DisplayPort}},
		firstplural={interfaces DisplayPort},
		description={Interfaz de comunicaci{\'o}n digital estandarizada para la transmisi{\'o}n de video, audio y datos auxiliares entre dispositivos de c{\'o}mputo y pantallas. DisplayPort fue desarrollada por la \emph{Video Electronics Standards Association} (VESA) con el objetivo de ofrecer mayores tasas de transferencia, soporte para resoluciones y frecuencias elevadas, y mecanismos avanzados de sincronizaci{\'o}n y gesti{\'o}n de m{\'u}ltiples pantallas. En contextos de sistemas de informaci{\'o}n y arquitecturas computacionales, DisplayPort permite la interconexi{\'o}n eficiente entre unidades de procesamiento gr{\'a}fico y dispositivos de visualizaci{\'o}n, siendo ampliamente utilizada en entornos profesionales y de alto rendimiento.},
		type=main,
		symbol={🖥️},
		user1={Interfaz de video digital},
		user2={Transmisi{\'o}n de audio y video},
		user3={Alta resoluci{\'o}n y frecuencia},
		see={hdmi,interfazvideo,visualizacion}
	}
	
	\newglossaryentry{disponibilidad}{
		name={Disponibilidad},
		text={disponibilidad},
		sort={disponibilidad},
		plural={disponibilidades},
		first={disponibilidad},
		firstplural={disponibilidades},
		description={Propiedad de un \gls{sistema}, servicio o componente que indica el grado en que est{\'a} operativo y accesible cuando se requiere su uso, expresada normalmente como un porcentaje del tiempo total en un per{\'i}odo determinado. La \emph{disponibilidad} combina factores como la \gls{fiabilidad} (ausencia de fallos) y la capacidad de recuperaci{\'o}n r{\'a}pida ante interrupciones, y est{\'a} influida por el \gls{rendimiento}, el mantenimiento preventivo y la \gls{escalabilidad}. En entornos \gls{iot} y \glspl{sibiot}, una alta disponibilidad es esencial para garantizar la continuidad de servicios cr{\'i}ticos, la integridad de los datos en \gls{tiemporeal} y la satisfacci{\'o}n del \gls{usuario} final.},
		type=main,
		symbol={},
		user1={M{\'e}trica de calidad de servicio},
		user2={Continuidad operativa},
		user3={Indicador de confiabilidad},
		see={fiabilidad,rendimiento,escalable}
	}
	
	\newglossaryentry{dispositivo}
	{
		name={Dispositivo},
		text={dispositivo},
		sort={dispositivo},
		plural={dispositivos},
		first={dispositivo},
		firstplural={dispositivos},
		description={Un \emph{dispositivo} es una unidad f{\'i}sica o electr{\'o}nica capaz de realizar funciones espec{\'i}ficas dentro de un \gls{sistema}. En el contexto de \gls{iot}, un dispositivo puede ser un \gls{sensor}, \gls{actuador}, \gls{microcontrolador}, m{\'o}dulo de \gls{conectividad} u otro componente que interact{\'u}a con el \gls{entorno}, recolecta \glspl{dato} o ejecuta acciones autom{\'a}ticas, formando parte integral de una \gls{arquitecturadistribuida}.},
		type=\glsdefaulttype,
		symbol={📱},
		user1={Equipo},
		user2={Terminal},
		user3={Componente electr{\'o}nico},
		see={dispositivoiot}
	}
	
	\newglossaryentry{dispositivobiometrico}{
		name={Dispositivo biom{\'e}trico},
		text={dispositivo biom{\'e}trico},
		sort={dispositivo biometrico},
		plural={dispositivos biom{\'e}tricos},
		first={dispositivo biom{\'e}trico},
		firstplural={dispositivos biom{\'e}tricos},
		description={Dispositivo electr{\'o}nico capaz de capturar, procesar y analizar caracter{\'i}sticas fisiol{\'o}gicas o conductuales {\'u}nicas de una persona con el fin de realizar procesos de identificaci{\'o}n o autenticaci{\'o}n. Entre las caracter{\'i}sticas biom{\'e}tricas m{\'a}s comunes se incluyen huellas dactilares, reconocimiento facial, iris, voz, geometr{\'i}a de la mano o patrones de comportamiento. En sistemas \gls{iot} y \glspl{sibiot}, los dispositivos biom{\'e}tricos se integran como mecanismos de \gls{control} de acceso, \gls{autenticacionmultifactor} y \gls{protecciondatospersonales}, garantizando mayores niveles de \gls{seguridad}, trazabilidad e integridad en entornos f{\'i}sicos y digitales interconectados.},
		type=main,
		symbol={🧬},
		user1={Captura de rasgos biom{\'e}tricos},
		user2={Autenticaci{\'o}n e identificaci{\'o}n segura},
		user3={Integraci{\'o}n en sistemas IoT y control de acceso},
		see={autenticacion,controlgranular,seguridad,protecciondatospersonales}
	}
	
	\newglossaryentry{dispositivoconectado}
	{
		name={Dispositivo conectado},
		text={dispositivo conectado},
		sort={dispositivo conectado},
		plural={dispositivos conectados},
		first={dispositivo conectado},
		firstplural={dispositivos conectados},
		description={Un \emph{dispositivo conectado} es un elemento f{\'i}sico que posee capacidad de comunicaci{\'o}n digital y est{\'a} integrado en una \gls{red}, permitiendo el intercambio de \glspl{dato} con otros sistemas o servicios. En entornos \gls{iot}, puede ser un \gls{sensor}, un \gls{actuador}, un \gls{microcontrolador} o un sistema embebido que env{\'i}a o recibe informaci{\'o}n mediante tecnolog{\'i}as como \gls{wifi}, \gls{bluetooth}, \gls{zigbee}, \gls{lora} o redes celulares, habilitando la \gls{automatizacion} y el \gls{controlinteligente} de procesos.},
		type=\glsdefaulttype,
		symbol={📡},
		user1={Nodo conectado},
		user2={Componente en red},
		user3={Connected device},
		see={dispositivoiot}
	}
	
	\newglossaryentry{dispositivoembebido}{
		name={Dispositivo embebido},
		text={dispositivo embebido},
		sort={dispositivo embebido},
		plural={dispositivos embebidos},
		first={dispositivo embebido},
		firstplural={dispositivos embebidos},
		description={Sistema computacional especializado integrado dentro de un producto o infraestructura f{\'i}sica, dise{\~n}ado para ejecutar funciones concretas con restricciones de recursos (energ{\'i}a, memoria, c{\'a}lculo y conectividad) y, con frecuencia, en condiciones de operaci{\'o}n en tiempo real. Un dispositivo embebido combina hardware y software (firmware) y suele incorporar perif{\'e}ricos de E/S, sensores y/o actuadores. En \gls{iot}/\glspl{sibiot}, los dispositivos embebidos operan como nodos de captura y acci{\'o}n en el \gls{mundofisico}, ejecutan preprocesamiento en el borde y participan en la comunicaci{\'o}n con \glspl{gateway} y servicios de \gls{nube}, por lo que su dise{\~n}o impacta directamente en latencia, consumo energ{\'e}tico, robustez y seguridad del sistema.},
		type=main,
		symbol={🧩},
		user1={Computaci{\'o}n integrada},
		user2={Restricciones de recursos},
		user3={Nodo de borde en IoT},
		see={microcontrolador,microprocesador,firmware,dispositivoiot,edgecomputing}
	}
	
	\newglossaryentry{dispositivointeligente}{
		name={Dispositivo inteligente},
		text={dispositivo inteligente},
		plural={dispositivos inteligentes},
		sort={dispositivo inteligente},
		first={dispositivo inteligente},
		firstplural={dispositivos inteligentes},
		description={Un \emph{dispositivo inteligente} es un elemento f{\'i}sico que combina capacidades de procesamiento, conectividad a redes digitales y, en muchos casos, algoritmos de inteligencia artificial o de automatizaci{\'o}n. Estos dispositivos pueden recopilar informaci{\'o}n del entorno a trav{\'e}s de \glspl{sensor}, procesarla localmente, tomar decisiones aut{\'o}nomas y ejecutar acciones mediante \glspl{actuador}. Ejemplos com{\'u}nes incluyen tel{\'e}fonos inteligentes, altavoces con asistentes virtuales, electrodom{\'e}sticos conectados y sistemas de dom{\'o}tica. En entornos \gls{iot}, los dispositivos inteligentes son nodos clave que habilitan la interacci{\'o}n continua entre el mundo f{\'i}sico y digital, aportando valor en sectores como la salud, la educaci{\'o}n, la agricultura, la industria y las \glspl{ciudadinteligente}.},
		type=main,
		symbol={🤖},
		user1={Capacidades de procesamiento y conectividad},
		user2={Uso de sensores, actuadores y algoritmos},
		user3={Base de aplicaciones en IoT y dom{\'o}tica},
		see={dispositivoconectado,sibiot,iot,actuador,sensor}
	}
	
	\newglossaryentry{dispositivoiot}
	{
		name={Dispositivo IoT},
		text={dispositivo IoT},
		sort={dispositivo iot},
		plural={dispositivos IoT},
		first={dispositivo IoT},
		firstplural={dispositivos IoT},
		description={Un \emph{dispositivo IoT} es un componente f{\'i}sico que forma parte de una \gls{red} \gls{iot}, dotado de capacidades para adquirir \glspl{dato} del \gls{entorno} (mediante \glspl{sensor}), procesarlos localmente (mediante \glspl{microcontrolador}) y/o ejecutar acciones (mediante \glspl{actuador}). Los dispositivos IoT pueden conectarse a trav{\'e}s de distintas tecnolog{\'i}as de \gls{conectividad} para enviar informaci{\'o}n a \glspl{servidor} remotos o recibir comandos, habilitando la \gls{automatizacion}, la monitorizaci{\'o}n y el \gls{controlinteligente} en \glspl{arquitecturadistribuida}.},
		type=\glsdefaulttype,
		symbol={📡},
		user1={Nodo IoT},
		user2={Componente IoT},
		user3={Device IoT},
		see={dispositivo}
	}
	
	\newglossaryentry{distribucion}{
		name={Distribuci{\'o}n},
		text={distribuci{\'o}n},
		sort={distribucion},
		first={distribuci{\'o}n},
		description={Propiedad por la cual componentes, datos o procesos se despliegan en m{\'u}ltiples nodos o ubicaciones, coordinados mediante mecanismos de comunicaci{\'o}n y control.},
		type=main,
		symbol={🌐},
		user1={Despliegue en nodos},
		user2={Coordinaci{\'o}n},
		user3={Escala},
		see={sistemadistribuido,basedatosdistribuida}
	}
	
	\newglossaryentry{divergencia}{
		name={Divergencia},
		text={divergencia},
		sort={divergencia},
		plural={divergencias},
		first={divergencia},
		firstplural={divergencias},
		description={Desacoplamiento o desalineaci{\'o}n entre artefactos que deber{\'\i}an corresponderse, por ejemplo entre modelos y c{\'o}digo, especificaci{\'o}n y comportamiento observado, o dise{\~n}o arquitect{\'o}nico y estructura real del sistema. La divergencia puede originarse por cambios no trazados, decisiones implementadas sin actualizar modelos, deuda t{\'e}cnica, presi{\'o}n de tiempo o falta de verificaci{\'o}n continua. En sistemas IoT y distribuidos, se agrava por la coexistencia de firmware, servicios en la nube, configuraciones de red y despliegues heterog{\'e}neos. Su control requiere trazabilidad, automatizaci{\'o}n de pruebas, revisiones y transformaciones controladas.},
		type=main,
		symbol={⚠️},
		user1={Desalineaci{\'o}n dise{\~n}o-implementaci{\'o}n},
		user2={Riesgo de inconsistencias},
		user3={Indicador de deuda t{\'e}cnica},
		see={consistencia,verificacion,transformacioncontrolada,evolucioncontrolada,trazabilidad}
	}
	
	\newglossaryentry{dns}
	{
		name={DNS},
		text={DNS},
		sort={dns},
		plural={DNSs},
		first={sistema de nombres de dominio (\textit{Domain Name System}, DNS)},
		firstplural={DNSs},
		description={\emph{DNS} es un sistema jer{\'a}rquico de nomenclatura que traduce nombres de dominio legibles por humanos (como \texttt{example.com}) en direcciones IP comprensibles por m{\'a}quinas. Permite a los \glspl{dispositivo} en una \gls{red} localizar y comunicarse con otros sistemas mediante nombres simb{\'o}licos. En entornos \gls{iot}, el uso de DNS facilita la integraci{\'o}n con servicios en la \gls{nube}, la identificaci{\'o}n remota de \glspl{dispositivoiot} y la gesti{\'o}n de recursos distribuidos.},
		type=\glsdefaulttype,
		symbol={🌐},
		user1={Sistema de nombres de dominio},
		user2={Domain Name System},
		user3={Servidor DNS},
		see={dominio}
	}
	
	\newglossaryentry{documentacion}{
		name={Documentaci{\'o}n},
		text={documentaci{\'o}n},
		sort={documentacion},
		first={documentaci{\'o}n},
		description={Conjunto de artefactos escritos, gr{\'a}ficos o digitales que describen, explican y respaldan el dise{\~n}o, funcionamiento, uso y mantenimiento de un sistema. La documentaci{\'o}n cumple un rol fundamental en la comunicaci{\'o}n entre los involucrados, la trazabilidad de decisiones, la transferencia de conocimiento y la garant{\'i}a de calidad a lo largo del ciclo de vida del software y de los sistemas de informaci{\'o}n. En \glspl{sibiot}, una documentaci{\'o}n adecuada facilita la integraci{\'o}n de componentes heterog{\'e}neos, la operaci{\'o}n segura y la evoluci{\'o}n controlada del sistema.},
		type=main,
		symbol={📚},
		user1={Artefactos descriptivos},
		user2={Trazabilidad},
		user3={Transferencia de conocimiento},
		see={analisisrequisitos,controlversiones,calidadsoftware}
	}
	
	\newglossaryentry{dominio}
	{
		name={Dominio},
		text={dominio},
		sort={dominio},
		plural={dominios},
		first={dominio},
		firstplural={dominios},
		description={Un \emph{dominio} es un campo o {\'a}rea de aplicaci{\'o}n definido en el que se desarrollan actividades, modelos o sistemas. En el contexto de \gls{iot}, un dominio puede abarcar sectores como agricultura, salud, ciudades inteligentes o industria, delimitando los objetivos, requisitos y caracter{\'i}sticas espec{\'i}ficas de cada aplicaci{\'o}n.},
		type=\glsdefaulttype,
		symbol={🏷️},
		user1={Campo de aplicaci{\'o}n},
		user2={Entorno tem{\'a}tico},
		user3={Area de especializaci{\'o}n},
		see={dns}
	}
	
	\newglossaryentry{domotica}{
		name={Dom{\'o}tica},
		text={dom{\'o}tica},
		sort={domotica},
		plural={dom{\'o}ticas},
		first={dom{\'o}tica},
		firstplural={dom{\'o}ticas},
		description={Conjunto de tecnolog{\'i}as aplicadas al control y la automatizaci{\'o}n inteligente de la vivienda, que permiten gestionar de forma eficiente la energ{\'i}a, la seguridad, el confort y la comunicaci{\'o}n entre dispositivos. La \emph{dom{\'o}tica} integra \glspl{sensor}, \glspl{actuador}, redes de comunicaci{\'o}n y plataformas de gesti{\'o}n para proporcionar un entorno habitable m{\'a}s seguro, accesible, sostenible y adaptado a las necesidades de sus ocupantes. En el marco de los \glspl{sibiot}, la dom{\'o}tica se concibe como un ecosistema de \glspl{objetointeligente} que cooperan en tiempo real mediante \gls{iot}, \gls{ia} y servicios en la nube.},
		type=main,
		symbol={🏠},
		user1={Automatizaci{\'o}n del hogar},
		user2={Integraci{\'o}n de tecnolog{\'i}as IoT},
		user3={Eficiencia, seguridad y confort residencial},
		see={iot,objetointeligente,automatizacion,sibiot}
	}
	
	\newglossaryentry{drawable}{
		name={Drawable},
		text={drawable},
		sort={drawable},
		plural={recursos drawable},
		first={drawable},
		firstplural={recursos drawable},
		description={Recurso gr{\'a}fico en Android utilizado para representar im{\'a}genes, formas o estados visuales (p.\,ej., \emph{bitmaps}, vectores o \emph{shape drawables}) consumidos por la interfaz de usuario.},
		type=main,
		symbol={🖼️},
		user1={Recurso gr{\'a}fico},
		user2={Android},
		user3={UI},
		see={assets,layout,values}
	}
	
	\newglossaryentry{dvi}{
		name={DVI},
		text={DVI},
		sort={dvi},
		plural={interfaces DVI},
		first={interfaz de video digital (\emph{Digital Visual Interface}, DVI)},
		firstplural={interfaces DVI},
		description={\emph{Digital Visual Interface} (DVI). Interfaz de video desarrollada para la transmisi{\'o}n de se{\~n}ales de imagen desde dispositivos de c{\'o}mputo hacia monitores digitales, con soporte para se{\~n}ales digitales (DVI-D), anal{\'o}gicas (DVI-A) o combinadas (DVI-I). DVI fue ampliamente utilizada durante la transici{\'o}n entre tecnolog{\'i}as de visualizaci{\'o}n anal{\'o}gica y digital, proporcionando mayor calidad de imagen y reducci{\'o}n de interferencias frente a interfaces anal{\'o}gicas tradicionales. En arquitecturas de sistemas de informaci{\'o}n y plataformas computacionales, DVI representa un est{\'a}ndar hist{\'o}rico relevante en la evoluci{\'o}n de las interfaces de video digitales.},
		type=main,
		symbol={🖥️},
		user1={Interfaz de video digital},
		user2={Transici{\'o}n anal{\'o}gica--digital},
		user3={Conexi{\'o}n gr{\'a}fica},
		see={interfazvideo,displayport,hdmi,vga}
	}
	
	\newglossaryentry{dynamodb}{
		name={Amazon DynamoDB},
		text={DynamoDB},
		sort={amazon dynamodb},
		first={DynamoDB},
		description={Servicio gestionado de base de datos NoSQL de tipo clave--valor y documento, ofrecido por Amazon Web Services (AWS), dise{\~n}ado para proporcionar latencia de un solo d{\'i}gito en milisegundos y escalabilidad horizontal autom{\'a}tica. Amazon DynamoDB abstrae completamente la administraci{\'o}n de infraestructura, incluyendo particionamiento, replicaci{\'o}n y tolerancia a fallos, permitiendo el manejo eficiente de grandes vol{\'u}menes de datos y altas tasas de operaciones concurrentes. En el contexto de los sistemas de informaci{\'o}n basados en IoT y SIbIoT, DynamoDB se utiliza para el almacenamiento de datos semiestructurados, eventos, estados de dispositivos y registros de alta frecuencia, complementando bases de datos relacionales en arquitecturas policl{\'i}nicas de persistencia},
		type=main,
		symbol={🗄️},
		see={basedatos,nosql,nube,escalabilidad,procesamientocontinuo,sibiot}
	}
	
	\newglossaryentry{ecosistema}{
		name={Ecosistema},
		text={ecosistema},
		sort={ecosistema},
		plural={ecosistemas},
		first={ecosistema},
		firstplural={ecosistemas},
		description={Conjunto de elementos interrelacionados que interact{\'u}an entre s{\'i} dentro de un entorno determinado, compartiendo recursos, informaci{\'o}n y objetivos comunes. En el {\'a}mbito tecnol{\'o}gico, un \emph{ecosistema} se refiere a la red de \glspl{dispositivointeligente}, \glspl{aplicacion}, \glspl{servicio} y plataformas que cooperan para generar valor de manera conjunta. En los \glspl{sibiot}, el ecosistema integra hardware, software, \glspl{sensor}, \glspl{actuador}, protocolos de comunicaci{\'o}n y mecanismos de \gls{interoperabilidad}, permitiendo la coordinaci{\'o}n de procesos en tiempo real, la escalabilidad y la sostenibilidad de las soluciones implementadas.},
		type=main,
		symbol={🌐},
		user1={Red de interacciones},
		user2={Entorno tecnol{\'o}gico colaborativo},
		user3={Integraci{\'o}n en SIbIoT},
		see={sibiot,iot,interoperabilidad,objetointeligente}
	}
	
	\newglossaryentry{ecosistemaiot}
	{
		name={Ecosistema IoT},
		text={ecosistema IoT},
		sort={ecosistema iot},
		plural={ecosistemas IoT},
		first={ecosistema IoT},
		firstplural={ecosistemas IoT},
		description={El \emph{ecosistema IoT} es el conjunto interrelacionado de componentes tecnol{\'o}gicos, actores, plataformas y procesos que permiten el funcionamiento integral de un sistema \gls{iot}. Incluye \glspl{dispositivoconectado} (como \glspl{sensor} y \glspl{actuador}), \glspl{redcomunicacion}, \glspl{infraestructura} de \gls{almacenamiento} y \gls{procesamientodatos} (edge, fog y cloud), servicios anal{\'i}ticos, \glspl{interfaz} de \gls{usuario} y normas de \gls{interoperabilidad}. Este ecosistema habilita soluciones distribuidas, escalables e inteligentes en sectores como salud, agricultura, industria y \glspl{ciudadinteligente}.},
		type=\glsdefaulttype,
		symbol={🌐},
		user1={Ecosistema de IoT},
		user2={IoT ecosystem},
		user3={Plataforma IoT integral},
		see={iot,sistemadistribuido,procesamientodistribuido,sibiot}
	}
	
	\newglossaryentry{edgecomputing}
	{
		name={Edge Computing},
		text={edge computing},
		sort={edge computing},
		plural={modelos de edge computing},
		first={computaci{\'o}n de borde o (\textit{Edge Computing})},
		firstplural={modelos de edge computing},
		description={El Computaci{\'o}n de borde o \emph{edge computing} es un modelo de computaci{\'o}n que procesa los \glspl{dato} cerca de su fuente de origen, en \glspl{dispositivo} perif{\'e}ricos o nodos intermedios, reduciendo la \gls{latencia} y la carga de la \gls{red} al evitar el env{\'i}o continuo de informaci{\'o}n a la \gls{nube}. Es clave en \glspl{sibiot} que requieren respuestas en \gls{tiemporeal}, como \gls{transporteinteligente}, \glspl{edificiointeligente} y aplicaciones industriales.},
		type=\glsdefaulttype,
		symbol={⚡},
		user1={Computaci{\'o}n en el borde},
		user2={Edge},
		user3={Procesamiento perimetral},
		see={fogcomputing}
	}
	
	\newglossaryentry{edhoc}{
		name={EDHOC},
		text={EDHOC},
		sort={edhoc},
		plural={protocolos EDHOC},
		first={Intercambio Diffie–Hellman ef{\'i}mero sobre COSE (\textit{Ephemeral Diffie--Hellman Over COSE}, EDHOC)},
		firstplural={protocolos EDHOC},
		description={\emph{EDHOC} es un protocolo de establecimiento de claves liviano dise{\~n}ado para dispositivos con recursos restringidos en entornos \gls{iot}. Permite realizar autenticaci{\'o}n mutua y generar claves criptogr{\'a}ficas seguras utilizando mensajes compactos basados en COSE (CBOR Object Signing and Encryption). Su dise{\~n}o minimalista reduce el consumo energ{\'e}tico, el ancho de banda y la latencia, facilitando su uso en redes \gls{lpwan}, sensores de bajo consumo, \glspl{dispositivointeligente} y sistemas embebidos. EDHOC tambi{\'e}n sirve como base para protocolos como OSCORE, permitiendo cifrado de extremo a extremo y protecci{\'o}n de datos en arquitecturas distribuidas y dispositivos heterog{\'e}neos},
		type=main,
		symbol={🔐},
		user1={Autenticaci{\'o}n mutua liviana para IoT},
		user2={Establecimiento eficiente de claves criptogr{\'a}ficas},
		user3={Base para OSCORE y protecci{\'o}n extremo a extremo},
		see={cose,oscore,protocoloseguridad,nbiot,lpwan,iot}
	}
	
	
	\newglossaryentry{edificiointeligente}
	{
		name={Edificio inteligente},
		text={edificio inteligente},
		sort={edificio inteligente},
		plural={edificios inteligentes},
		first={edificio inteligente},
		firstplural={edificios inteligentes},
		description={Un \emph{edificio inteligente} es una infraestructura que integra tecnolog{\'i}as de \gls{automatizacion}, \glspl{sensor}, sistemas de control y \gls{conectividad} para optimizar el uso de recursos como energ{\'i}a, agua, climatizaci{\'o}n, iluminaci{\'o}n y \gls{seguridad}. Utiliza \glspl{dato} en \gls{tiemporeal} para mejorar el confort de los ocupantes, la eficiencia operativa y la sostenibilidad ambiental. Constituye un componente fundamental de las \glspl{ciudadinteligente}.},
		type=\glsdefaulttype,
		symbol={🏢},
		user1={Smart building},
		user2={Infraestructura inteligente},
		user3={Edificio automatizado},
		see={ciudadinteligente}
	}
	
	\newglossaryentry{eeprom}{
		name={EEPROM},
		text={EEPROM},
		sort={eeprom},
		first={memoria de solo lectura programable y borrable el{\'e}ctricamente (\emph{Electrically Erasable Programmable Read-Only Memory}, EEPROM)},
		description={Memoria no vol{\'a}til dise{\~n}ada para almacenar peque{\~n}as cantidades de datos que deben conservarse entre reinicios del sistema, como configuraciones o par{\'a}metros de usuario. La \emph{EEPROM} permite escritura y lectura a nivel de byte, aunque con un n{\'u}mero limitado de ciclos de escritura.},
		type=main,
		symbol={🧾},
		user1={Persistencia de datos},
		user2={Configuraciones},
		user3={Almacenamiento por bytes},
		see={memoriaflash,sram}
	}
	
	\newglossaryentry{efectividad}
	{
		name={Efectividad},
		text={efectividad},
		sort={efectividad},
		plural={efectividades},
		first={efectividad},
		firstplural={efectividades},
		description={La efectividad es el grado en que un \gls{sistema} o \gls{proceso} logra los objetivos establecidos, cumpliendo con los \glspl{requisito} y las \glspl{condicion} definidas. En el contexto de los \glspl{sibiot}, la efectividad se mide por la capacidad del \gls{sistemainformatico} para ofrecer resultados confiables, precisos y oportunos, contribuyendo al desempe{\~n}o general y al logro de metas estrat{\'e}gicas.},
		type=\glsdefaulttype,
		symbol={},
		user1={Eficacia},
		user2={Rendimiento},
		user3={Cumplimiento de objetivos},
		see={requisito}
	}
	
	\newglossaryentry{eficacia}
	{
		name={Eficacia},
		text={eficacia},
		sort={eficacia},
		plural={eficacias},
		first={Eficacia},
		description={Capacidad de un \gls{sistema}, \gls{proceso} o \gls{dispositivo} para alcanzar los objetivos o resultados previstos en condiciones establecidas. 
			En el contexto de la ingenier{\'i}a y los \glspl{sibiot}, la eficacia se refiere al grado en que un sistema realiza correctamente las funciones para las cuales fue dise{\~n}ado, garantizando la detecci{\'o}n, medici{\'o}n o control esperados. 
			En gesti{\'o}n de proyectos y evaluaci{\'o}n de desempe{\~n}o, se asocia con la consecuci{\'o}n de metas y la capacidad de responder adecuadamente a los requerimientos del entorno. 
			A diferencia de la \gls{eficiencia}, que considera la relaci{\'o}n entre recursos empleados y resultados obtenidos, la eficacia se centra exclusivamente en el logro del resultado deseado, independientemente del costo o esfuerzo necesario.},
		see=[Relacionado con:]{eficiencia,desempeno,objetivo,resultado}
	}
	
	\newglossaryentry{eficiencia}{
		name={Eficiencia},
		text={eficiencia},
		sort={eficiencia},
		plural={eficiencias},
		first={eficiencia},
		firstplural={eficiencias},
		description={Relaci{\'o}n entre los resultados obtenidos y los recursos empleados para alcanzarlos, expresada como la capacidad de lograr un objetivo con el menor gasto posible de tiempo, energ{\'i}a o costes, sin sacrificar la calidad. La \emph{eficiencia} implica el uso {\'o}ptimo de medios disponibles y constituye un indicador clave en ingenier{\'i}a de sistemas y gesti{\'o}n organizacional. En el contexto de los \glspl{sibiot}, la eficiencia se refiere a la capacidad de los \glspl{sensor}, \glspl{actuador}, \glspl{aplicacion} y plataformas de procesar datos en \gls{tiemporeal}, optimizar recursos energ{\'e}ticos, reducir redundancias y mejorar la \gls{productividad}, contribuyendo a la \gls{sostenibilidad} y a la eficiencia operativa.},
		type=main,
		symbol={⚡},
		user1={Uso {\'o}ptimo de recursos},
		user2={Indicador de desempe{\~n}o y calidad},
		user3={Objetivo clave en SIbIoT},
		see={rendimiento,productividad,sostenibilidad,eficienciaoperativa}
	}
	
	\newglossaryentry{eficienciaanalitica}{
		name={Eficiencia anal{\'i}tica},
		text={eficiencia anal{\'i}tica},
		sort={eficiencia analitica},
		plural={eficiencias anal{\'i}ticas},
		first={eficiencia anal{\'i}tica},
		firstplural={eficiencias anal{\'i}ticas},
		description={Grado en que los m{\'e}todos, modelos o procesos de an{\'a}lisis de datos logran obtener resultados v{\'a}lidos, relevantes y oportunos empleando la menor cantidad posible de recursos computacionales, tiempo o coste. La \emph{eficiencia anal{\'i}tica} combina la precisi{\'o}n y validez de los resultados con la optimizaci{\'o}n de la capacidad de procesamiento, siendo fundamental en entornos con grandes vol{\'u}menes de informaci{\'o}n. En el contexto de los \glspl{sibiot}, la eficiencia anal{\'i}tica permite aprovechar de manera {\'o}ptima los datos generados por \glspl{sensor}, aplicando algoritmos de \gls{aprendizajeautomatico}, \glspl{modelopredictivo} o \glspl{redneuronal} en \gls{tiemporeal}, lo que facilita la \gls{eficienciaoperativa}, la \gls{tomadecisiones} inteligente y la \gls{sostenibilidad}.},
		type=main,
		symbol={📊},
		user1={Optimaci{\'o}n de recursos en an{\'a}lisis de datos},
		user2={Equilibrio entre precisi{\'o}n y coste computacional},
		user3={Soporte a la toma de decisiones en SIbIoT},
		see={eficiencia,rendimiento,eficienciaoperativa,modelopredictivo,redneuronal}
	}
	
	\newglossaryentry{eficienciacomputacional}{
		name={Eficiencia computacional},
		text={eficiencia computacional},
		sort={eficiencia computacional},
		plural={eficiencias computacionales},
		first={eficiencia computacional},
		firstplural={eficiencias computacionales},
		description={Capacidad de un sistema, algoritmo o arquitectura para realizar una tarea utilizando de manera {\'o}ptima los recursos computacionales disponibles, tales como tiempo de procesamiento, memoria, almacenamiento, energ{\'i}a y ancho de banda. La eficiencia computacional se eval{\'u}a habitualmente mediante m{\'e}tricas de complejidad temporal y espacial, consumo energ{\'e}tico y rendimiento bajo carga. En el contexto de los sistemas de informaci{\'o}n basados en IoT y SIbIoT, la eficiencia computacional es un factor cr{\'i}tico debido a la heterogeneidad y a las restricciones de recursos de los dispositivos, y se logra mediante t{\'e}cnicas como procesamiento en el borde, paralelizaci{\'o}n, desacoplamiento de componentes y selecci{\'o}n adecuada de modelos de comunicaci{\'o}n y procesamiento},
		type=main,
		symbol={⚡},
		see={procesamientodistribuido,procesamientoparalelo,procesamientocontinuo,eficienciaoperativa,escalabilidad,sibiot}
	}
	
	\newglossaryentry{eficienciaenergetica}{
		name={Eficiencia energ{\'e}tica},
		text={eficiencia energ{\'e}tica},
		sort={eficiencia energetica},
		plural={eficiencias energ{\'e}ticas},
		first={eficiencia energ{\'e}tica},
		firstplural={eficiencias energ{\'e}ticas},
		description={Relaci{\'o}n entre la energ{\'i}a consumida y el rendimiento obtenido por un sistema o proceso. Denota la capacidad de alcanzar un objetivo funcional con el menor consumo posible de energ{\'i}a, manteniendo la calidad del resultado. En tecnolog{\'i}a e ingenier{\'i}a, es criterio clave para dise{\~n}ar sistemas sostenibles y optimizados, reducir huella ambiental y costos operativos, y cumplir regulaciones de eficiencia. En \glspl{sibiot} y sistemas distribuidos, implica selecci{\'o}n de hardware eficiente, algoritmos y protocolos de bajo consumo, gesti{\'o}n inteligente de energ{\'i}a y monitorizaci{\'o}n continua de indicadores.},
		type=main,
		user1={Minimizaci{\'o}n del consumo manteniendo el rendimiento},
		user2={Dise{\~n}o, medici{\'o}n y optimizaci{\'o}n energ{\'e}tica de sistemas},
		user3={M{\'e}tricas: rendimiento por vatio, EER, PUE, kWh/operaci{\'o}n},
		user4={Estrategias: escalado din{\'a}mico, apagado selectivo, offloading a la nube/borde},
		see={computacionverde,optimizacion,rendimiento,sostenibilidad,consumoenergetico},
		symbol={⚡}
	}	
	
	\newglossaryentry{eficienciaoperativa}{
		name={Eficiencia operativa},
		text={eficiencia operativa},
		sort={eficiencia operativa},
		plural={eficiencias operativas},
		first={eficiencia operativa},
		firstplural={eficiencias operativas},
		description={Capacidad de una organizaci{\'o}n, proceso o \gls{sistema} para utilizar de manera {\'o}ptima sus recursos humanos, financieros, tecnol{\'o}gicos y materiales con el fin de maximizar el valor generado y minimizar desperdicios, costes o tiempos improductivos. La \emph{eficiencia operativa} implica un equilibrio entre \gls{rendimiento}, calidad, \gls{sostenibilidad} y \gls{productividad}, favoreciendo la competitividad y la resiliencia organizacional. En el contexto de los \glspl{sibiot}, la eficiencia operativa se alcanza mediante la integraci{\'o}n de \glspl{sensor}, \glspl{actuador}, algoritmos de an{\'a}lisis en \gls{tiemporeal} y plataformas inteligentes que permiten monitorizar procesos, anticipar fallos, optimizar recursos y automatizar decisiones.},
		type=main,
		symbol={⚙️},
		user1={Optimaci{\'o}n de recursos},
		user2={M{\'e}trica de desempe{\~n}o empresarial},
		user3={Objetivo de SIbIoT en organizaciones},
		see={rendimiento,productividad,sostenibilidad,sibiot}
	}
	
	\newglossaryentry{eficienciatermica}{
		name={Eficiencia t{\'e}rmica},
		text={eficiencia t{\'e}rmica},
		sort={eficiencia termica},
		plural={eficiencias t{\'e}rmicas},
		first={eficiencia t{\'e}rmica},
		firstplural={eficiencias t{\'e}rmicas},
		description={Relaci{\'o}n entre el trabajo o energ{\'i}a {\'u}til obtenida de un sistema t{\'e}rmico y la energ{\'i}a total suministrada en forma de calor. La \emph{eficiencia t{\'e}rmica} indica el grado en que una m{\'a}quina, motor o proceso convierte la energ{\'i}a t{\'e}rmica en energ{\'i}a mec{\'a}nica o el{\'e}ctrica aprovechable, siendo siempre menor al 100\% debido a las limitaciones de la segunda ley de la termodin{\'a}mica y a las p{\'e}rdidas inevitables. En el contexto de los \glspl{sibiot}, la eficiencia t{\'e}rmica puede monitorizarse en tiempo real mediante \glspl{sensor} de temperatura, presi{\'o}n y flujo, permitiendo optimizar procesos industriales, reducir el consumo energ{\'e}tico, mejorar la \gls{sostenibilidad} y anticipar fallos en sistemas de generaci{\'o}n o climatizaci{\'o}n.},
		type=main,
		symbol={🔥},
		user1={Conversi{\'o}n de calor en energ{\'i}a {\'u}til},
		user2={Indicador de eficiencia energ{\'e}tica},
		user3={Par{\'a}metro clave en IoT industrial},
		see={eficiencia,sostenibilidad,variablefisica,sensor}
	}
	
	\newglossaryentry{eficiente}{
		name={Eficiente},
		text={eficiente},
		sort={eficiente},
		plural={eficientes},
		first={eficiente},
		firstplural={eficientes},
		description={El t{\'e}rmino \emph{eficiente} se refiere a un \gls{sistema}, proceso, dispositivo o componente que logra sus objetivos utilizando la menor cantidad posible de recursos (tiempo, energ{\'i}a, memoria o costos), manteniendo un \gls{rendimiento} adecuado. A diferencia de la \gls{eficiencia}, que describe la propiedad abstracta, lo eficiente califica a un elemento espec{\'i}fico en su contexto de operaci{\'o}n. En \glspl{sibiot}, un dispositivo eficiente puede ser aquel que procesa datos en \gls{tiemporeal} con bajo consumo energ{\'e}tico, o una red que optimiza la \gls{latencia} y el \gls{anchobanda} para garantizar continuidad del servicio.},
		type=main,
		symbol={⚡},
		user1={Calificativo de sistemas u operaciones optimizadas},
		user2={Minimiza recursos como energ{\'i}a, tiempo o costos},
		user3={Clave en IoT para rendimiento sostenible},
		see={eficiencia,efectividad,rendimiento,latencia,tiemporeal}
	}
	
	\newglossaryentry{ejecutable}{
		name={Ejecutable},
		text={ejecutable},
		sort={ejecutable},
		plural={ejecutables},
		first={ejecutable},
		firstplural={ejecutables},
		description={Artefacto software que puede ser interpretado o ejecutado directamente por una m{\'a}quina o entorno de ejecuci{\'o}n. Representa la materializaci{\'o}n operativa de una especificaci{\'o}n o modelo previamente definido.},
		type=main,
		symbol={⚙️},
		user1={C{\'o}digo operativo},
		user2={Implementaci{\'o}n concreta},
		user3={Producto final},
		see={artefacto}
	}
	
	\newglossaryentry{elasticidad}{
		name={Elasticidad},
		text={elasticidad},
		sort={elasticidad},
		first={elasticidad},
		description={Capacidad de un sistema de ajustar din{\'a}micamente recursos o capacidad ante cambios de demanda. En sistemas distribuidos y nube, se materializa como escalado hacia arriba/abajo y aprovisionamiento automatizado.},
		type=main,
		symbol={🪢},
		user1={Ajuste din{\'a}mico},
		user2={Escalado},
		user3={Demanda variable},
		see={orquestacion,rendimiento}
	}
	\newglossaryentry{enclavamiento}{
		name={Enclavamiento},
		text={enclavamiento},
		sort={enclavamiento},
		plural={enclavamientos},
		first={enclavamiento},
		firstplural={enclavamientos},
		description={Un \emph{enclavamiento} es un mecanismo l{\'o}gico y/o f{\'\i}sico que impide ejecutar una acci{\'o}n si no se satisfacen condiciones previas definidas (p.\,ej., estado seguro, secuencia correcta, permisos, valores dentro de rango). En \glspl{cps}, los enclavamientos se usan para reducir riesgos en la actuaci{\'o}n y para asegurar que el lazo de control opere dentro de l{\'\i}mites aceptables, especialmente cuando hay conectividad, operaci{\'o}n remota y amenazas de seguridad.},
		type=main,
		symbol={🔒},
		user1={Prevenci{\'o}n de acciones inseguras},
		user2={Condiciones previas verificables},
		user3={Control con l{\'\i}mites de seguridad},
		see={seguridad,ciberseguridad,actuador,control}
	}
	
	\newglossaryentry{encriptacion}
	{
		name={Encriptaci{\'o}n},
		text={encriptaci{\'o}n},
		sort={encriptacion},
		plural={encriptaciones},
		first={encriptaci{\'o}n},
		firstplural={encriptaciones},
		description={La \emph{encriptaci{\'o}n} es el proceso de codificaci{\'o}n de informaci{\'o}n para proteger su confidencialidad, de modo que solo usuarios autorizados puedan acceder a ella mediante el uso de claves o algoritmos de \gls{cifrado}. Es un componente esencial en la \gls{ciberseguridad} y en la protecci{\'o}n de \glspl{dato} en \glspl{sistemainformatico} y \glspl{sibiot}.},
		type=\glsdefaulttype,
		symbol={🔒},
		user1={Codificaci{\'o}n},
		user2={Protecci{\'o}n de datos},
		user3={Encryption},
		see={cifrado}
	}
	
	\newglossaryentry{encriptaciondatos}{
		name={encriptaci{\'o}n de datos},
		text={encriptaci{\'o}n de datos},
		sort={encriptaciondatos},
		plural={encriptaciones de datos},
		first={encriptaci{\'o}n de datos},
		firstplural={encriptaciones de datos},
		description={La encriptaci{\'o}n de datos es una t{\'e}cnica de seguridad empleada para proteger la informaci{\'o}n frente a accesos no autorizados o interceptaciones. En los \glspl{sibiot}, donde los datos se transmiten constantemente entre \glspl{dispositivoconectado} y \glspl{servidor}, la encriptaci{\'o}n garantiza la privacidad del usuario y la integridad del flujo de informaci{\'o}n. Se implementa mediante algoritmos criptogr{\'a}ficos como AES, RSA o ECC, aplicados tanto a los datos en tr{\'a}nsito como a los almacenados.},
		type=\glsdefaulttype,
		symbol={🔒},
		user1={Data encryption},
		user2={Conversi{\'o}n de informaci{\'o}n en un formato cifrado para su protecci{\'o}n},
		user3={Mecanismo de seguridad para preservar la confidencialidad y autenticidad de los datos},
		see={seguridad,privacidad,integridad,calidaddatos,criptografia}
	}
	
	\newglossaryentry{energia}{
		name={Energ{\'i}a},
		text={energ{\'i}a},
		sort={energia},
		plural={energ{\'i}as},
		first={energ{\'i}a},
		firstplural={energ{\'i}as},
		description={Capacidad de un sistema f{\'i}sico para realizar trabajo o producir cambios, indispensable para el funcionamiento de dispositivos electr{\'o}nicos y sistemas computacionales. En sistemas de informaci{\'o}n y \glspl{sibiot}, la energ{\'i}a constituye un recurso cr{\'i}tico cuya gesti{\'o}n eficiente condiciona la autonom{\'i}a, disponibilidad y sostenibilidad operativa de los \glspl{dispositivo}. La optimizaci{\'o}n del uso de energ{\'i}as comprende estrategias de reducci{\'o}n de consumo, selecci{\'o}n de componentes de bajo consumo, gesti{\'o}n inteligente de potencia y aprovechamiento de fuentes renovables.},
		type=main,
		symbol={⚡},
		user1={Recurso energ{\'e}tico},
		user2={Consumo y eficiencia},
		user3={Sostenibilidad},
		see={consumoenergetico,bateria,computacionverde}
	}
	
	% Energy harvesting
	\newglossaryentry{energyharvesting}{
		name={Energy harvesting},
		text={energy harvesting},
		sort={energy harvesting},
		symbol={♻️},
		plural={t{\'e}cnicas de energy harvesting},
		first={energy harvesting},
		firstplural={t{\'e}cnicas de energy harvesting},
		description={T{\'e}cnicas que permiten recolectar energ{\'i}a ambiental, como vibraciones, temperatura, radiaci{\'o}n electromagn{\'e}tica o luz solar, para alimentar dispositivos IoT con baja demanda energ{\'e}tica}
	}
	
	\newglossaryentry{EnergiaSolar}{
		type=main,
		name={Energ{\'i}a solar},
		text={energ{\'i}a solar},
		sort={energia solar},
		plural={energ{\'i}as solares},
		first={energ{\'i}a solar},
		firstplural={energ{\'i}as solares},
		symbol={☀️},
		description={Forma de energ{\'i}a renovable obtenida de la radiaci{\'o}n solar. En sistemas IoT se utiliza mediante paneles fotovoltaicos para alimentar sensores y nodos inal{\'a}mbricos, prolongando la autonom{\'i}a y reduciendo la necesidad de mantenimiento},
		descriptionplural={Formas de energ{\'i}a renovable provenientes de la radiaci{\'o}n solar utilizadas para alimentar sistemas IoT}
	}
	
	\newglossaryentry{EnergiaCinetica}{
		type=main,
		name={Energ{\'i}a Cin{\'e}tica},
		text={energ{\'i}a cin{\'e}tica},
		sort={energia cinetica},
		plural={energ{\'i}as cin{\'e}ticas},
		first={energ{\'i}a cin{\'e}tica},
		firstplural={energ{\'i}as cin{\'e}ticas},
		symbol={⚙️},
		description={Energ{\'i}a asociada al movimiento de un cuerpo. En IoT puede ser recolectada mediante t{\'e}cnicas de \textit{\gls{energyharvesting}}, como piezoelectricidad, para producir electricidad y alimentar sensores de bajo consumo}
	}
	
	\newglossaryentry{enfoque}{
		name={Enfoque},
		text={enfoque},
		sort={enfoque},
		description={%
			Perspectiva conceptual o principio rector que gu{\'i}a el dise{\~n}o, la organizaci{\'o}n y la implementaci{\'o}n de un sistema, metodolog{\'i}a o soluci{\'o}n tecnol{\'o}gica.
			Un enfoque define c{\'o}mo se abordan los problemas, qu{\'e} criterios se priorizan y qu{\'e} supuestos subyacen a las decisiones de dise{\~n}o, influyendo directamente en la arquitectura, el comportamiento y la evoluci{\'o}n del sistema.
			En el contexto de los sistemas de informaci{\'o}n basados en IoT, el enfoque adoptado puede orientarse a principios como la reactividad, el desacoplamiento, la orientaci{\'o}n a eventos o el procesamiento distribuido.
		},
		first={enfoque},
		plural={enfoques},
		firstplural={enfoques},
		see={arquitectura,disenosistemas,sistemadistribuido}
	}
	
	\newglossaryentry{enfoquedesacoplado}{
		name={Enfoque desacoplado},
		text={enfoque desacoplado},
		sort={enfoque desacoplado},
		description={%
			Enfoque arquitect{\'o}nico que promueve la separaci{\'o}n l{\'o}gica y funcional de los componentes de un sistema, de modo que las interacciones entre ellos se realicen mediante interfaces bien definidas y contratos expl{\'i}citos.
			Un enfoque desacoplado reduce las dependencias directas entre componentes, facilitando la escalabilidad, la evoluci{\'o}n independiente y la integraci{\'o}n de nuevos servicios o tecnolog{\'i}as.
			En sistemas IoT y arquitecturas orientadas a eventos, el enfoque desacoplado suele materializarse mediante el uso de mensajer{\'i}a as{\'i}ncrona, buses de eventos y mecanismos de publicaci{\'o}n/suscripci{\'o}n.
		},
		first={enfoque desacoplado},
		plural={enfoques desacoplados},
		firstplural={enfoques desacoplados},
		see={desacoplado,mensajeriatiemporeal,orientadoeventos}
	}
	
	\newglossaryentry{enfoquereactivo}{
		name={Enfoque reactivo},
		text={enfoque reactivo},
		sort={enfoque reactivo},
		description={%
			Enfoque de dise{\~n}o de sistemas en el que el comportamiento global emerge como respuesta a la ocurrencia de eventos, cambios de estado o flujos de informaci{\'o}n generados en tiempo de ejecuci{\'o}n.
			En un enfoque reactivo, los componentes no siguen un control centralizado ni una secuencia fija de ejecuci{\'o}n, sino que reaccionan de forma as{\'i}ncrona y concurrente a los eventos del entorno.
			Este enfoque es ampliamente utilizado en sistemas IoT, arquitecturas orientadas a eventos y sistemas ciberf{\'i}sicos, donde la capacidad de respuesta en tiempo casi real resulta cr{\'i}tica.
		},
		first={enfoque reactivo},
		plural={enfoques reactivos},
		firstplural={enfoques reactivos},
		see={reactivo,orientadoeventos,procesamientodistribuido}
	}
	
	\newglossaryentry{enriquecimiento}{
		name={Enriquecimiento},
		text={enriquecimiento},
		sort={enriquecimiento},
		plural={enriquecimientos},
		first={enriquecimiento},
		firstplural={enriquecimientos},
		description={El enriquecimiento de datos busca aumentar el valor informativo mediante la incorporaci{\'o}n de fuentes externas o derivadas, como bases de datos complementarias, informaci{\'o}n geoespacial o variables calculadas. En el contexto de la gesti{\'o}n de calidad de datos, este proceso mejora la capacidad anal{\'i}tica y la \gls{tomadecisionesinformadas}, al proporcionar una visi{\'o}n m{\'a}s completa y contextual de los fen{\'o}menos observados. En los \glspl{sibiot}, el enriquecimiento puede incluir la combinaci{\'o}n de datos de m{\'u}ltiples \glspl{sensor} o la integraci{\'o}n con registros hist{\'o}ricos para generar conocimiento m{\'a}s profundo.},
		type=\glsdefaulttype,
		symbol={➕},
		user1={Data enrichment},
		user2={Ampliaci{\'o}n del valor informativo de los datos},
		user3={Integraci{\'o}n de fuentes complementarias para mejorar el an{\'a}lisis},
		see={validacion,limpieza,calidaddatos,tomadecisionesinformadas}
	}
	
	\newglossaryentry{enriquecimientodatos}{
		name={enriquecimiento de datos},
		text={enriquecimiento de datos},
		sort={enriquecimientodatos},
		plural={enriquecimientos de datos},
		first={enriquecimiento de datos},
		firstplural={enriquecimientos de datos},
		description={El enriquecimiento de datos permite aumentar el valor informativo de los registros al a{\~n}adir atributos complementarios, contextuales o derivados. En los \glspl{sibiot}, este proceso mejora la precisi{\'o}n de los an{\'a}lisis al combinar datos en \gls{tiemporeal} con hist{\'o}ricos o externos, como condiciones ambientales, ubicaci{\'o}n geogr{\'a}fica o informaci{\'o}n demogr{\'a}fica. Constituye un paso clave dentro de la \gls{calidaddatos} y la \gls{gestioninformacion}.},
		type=\glsdefaulttype,
		symbol={⚙️},
		user1={Data enrichment},
		user2={Ampliaci{\'o}n del contenido informativo de los datos mediante fuentes adicionales},
		user3={Proceso de mejora y contextualizaci{\'o}n de los datos recolectados},
		see={depuracion,normalizacion,calidaddatos,gestioninformacion,analisisbigdata}
	}
	
	\newglossaryentry{enrutador}
	{
		name={Enrutador},
		text={enrutador},
		sort={enrutador},
		plural={enrutadores},
		first={enrutador},
		firstplural={enrutadores},
		description={Un \emph{enrutador} es un dispositivo de \gls{red} que dirige el tr{\'a}fico de \glspl{dato} entre diferentes redes mediante el an{\'a}lisis de direcciones IP y el uso de tablas de encaminamiento. Determina la ruta m{\'a}s eficiente para transmitir paquetes hacia su destino. En sistemas \gls{iot}, los enrutadores permiten la \gls{conectividad} entre \glspl{dispositivo} locales y redes externas, e incluyen funciones de \gls{seguridad}, gesti{\'o}n de \gls{anchobanda} y soporte para protocolos inal{\'a}mbricos como \gls{wifi} o \gls{zigbee}.},
		type=\glsdefaulttype,
		symbol={📡},
		user1={Router},
		user2={Encaminador},
		user3={Gateway de red},
		see={red}
	}
	
	% ===================== Enrutamiento =====================
	\newglossaryentry{enrutamiento}{
		name={Enrutamiento},
		text={enrutamiento},
		sort={enrutamiento},
		plural={procesos de enrutamiento},
		first={enrutamiento},
		firstplural={procesos de enrutamiento},
		description={El \emph{enrutamiento} es el proceso mediante el cual se determina el camino m{\'a}s adecuado que deben seguir los paquetes de datos dentro de una \gls{red} para alcanzar su destino. Este proceso lo realizan los \glspl{enrutador} bas{\'a}ndose en tablas de encaminamiento y en \glspl{protocolocomunicacion}. En el contexto de \gls{iot} y \glspl{sibiot}, el enrutamiento debe considerar factores como \gls{latencia}, consumo energ{\'e}tico, \gls{fiabilidad}, movilidad de los nodos y la escalabilidad del sistema, siendo esencial para la eficiencia en la \gls{transmision} de datos en entornos distribuidos.},
		type=main,
		symbol={🗺️},
		user1={Selecci{\'o}n de rutas en redes de datos},
		user2={Basado en tablas y algoritmos de encaminamiento},
		user3={Crucial para IoT con movilidad y restricciones energ{\'e}ticas},
		see={protocolocomunicacion,red,enrutador,transmision,sibiot}
	}
	
	\newglossaryentry{enrutamientodinamico}{
		name={Enrutamiento din{\'a}mico},
		text={enrutamiento din{\'a}mico},
		sort={enrutamiento dinamico},
		plural={enrutamientos din{\'a}micos},
		first={enrutamiento din{\'a}mico},
		firstplural={enrutamientos din{\'a}micos},
		description={T{\'e}cnica de gesti{\'o}n de rutas en redes de comunicaci{\'o}n en la que los \glspl{enrutador} ajustan autom{\'a}ticamente sus tablas de encaminamiento en funci{\'o}n de cambios en la topolog{\'i}a de la red, el estado de los enlaces o el tr{\'a}fico existente. A diferencia del enrutamiento est{\'a}tico, que requiere configuraci{\'o}n manual, el \emph{enrutamiento din{\'a}mico} utiliza protocolos especializados como RIP, OSPF, EIGRP o BGP para descubrir y mantener rutas de forma autom{\'a}tica. En entornos \gls{iot}, esta capacidad resulta esencial para redes altamente distribuidas, m{\'o}viles y heterog{\'e}neas, garantizando la eficiencia, la escalabilidad y la resiliencia del sistema.},
		type=main,
		symbol={🌐},
		user1={Encaminamiento autom{\'a}tico},
		user2={Protocolos adaptativos de red},
		user3={Ajuste din{\'a}mico de rutas},
		see={enrutador,red,consistencia,wan}
	}
	
	\newglossaryentry{enterpriseapp}{
		name={Enterprise app},
		text={\textit{enterprise app}},
		sort={enterprise app},
		plural={\textit{enterprise apps}},
		first={\textit{enterprise app}},
		firstplural={\textit{enterprise apps}},
		description={Aplicaci{\'o}n orientada a procesos organizacionales que integra datos, flujos de trabajo, seguridad y servicios corporativos, con requisitos de disponibilidad, escalabilidad y mantenibilidad.},
		type=main,
		symbol={🏢},
		user1={Procesos organizacionales},
		user2={Integraci{\'o}n},
		user3={Requisitos empresariales},
		see={sistemainformacion,sgbd,plataformaweb}
	}
	
	\newglossaryentry{entorno}
	{
		name={Entorno},
		text={entorno},
		sort={entorno},
		plural={entornos},
		first={entorno},
		firstplural={entornos},
		description={Un \emph{entorno} es el conjunto de condiciones, elementos y circunstancias que rodean a un \gls{sistema}, proceso o entidad, influyendo en su funcionamiento o comportamiento. En el contexto de \gls{iot}, los entornos pueden ser f{\'i}sicos, virtuales, de desarrollo o distribuidos, dependiendo de la aplicaci{\'o}n.},
		type=\glsdefaulttype,
		symbol={🌍},
		user1={Ambiente},
		user2={Contexto},
		user3={Circunstancia},
		see={entornofisico}
	}
	
	\newglossaryentry{entornodesarrollo}
	{
		name={Entorno de desarrollo},
		text={entorno de desarrollo},
		sort={entorno de desarrollo},
		plural={entornos de desarrollo},
		first={entorno de desarrollo},
		firstplural={entornos de desarrollo},
		description={Un \emph{entorno de desarrollo} es un conjunto de herramientas, aplicaciones, bibliotecas y configuraciones utilizadas por los desarrolladores para escribir, probar, depurar y mantener software. Incluye editores de c{\'o}digo, compiladores o int{\'e}rpretes, sistemas de control de versiones y entornos de ejecuci{\'o}n o simulaci{\'o}n.},
		type=\glsdefaulttype,
		symbol={💻},
		user1={IDE},
		user2={Ambiente de desarrollo},
		user3={Development environment},
		see={entornodistribuido}
	}
	
	\newglossaryentry{entornodistribuido}
	{
		name={Entorno distribuido},
		text={entorno distribuido},
		sort={entorno distribuido},
		plural={entornos distribuidos},
		first={entorno distribuido},
		firstplural={entornos distribuidos},
		description={Un \emph{entorno distribuido} es una infraestructura computacional conformada por un conjunto de recursos de \gls{hardware} y software ubicados en distintos nodos interconectados, que cooperan entre s{\'i} para ejecutar procesos, compartir \glspl{dato} y ofrecer servicios de manera coordinada y transparente al usuario.},
		type=\glsdefaulttype,
		symbol={🖧},
		user1={Distributed environment},
		user2={Sistema distribuido},
		user3={Infraestructura distribuida},
		see={arquitecturadistribuida}
	}
	
	\newglossaryentry{entornofisico}
	{
		name={Entorno f{\'i}sico},
		text={entorno f{\'i}sico},
		sort={entorno fisico},
		plural={entornos f{\'i}sicos},
		first={entorno f{\'i}sico},
		firstplural={entornos f{\'i}sicos},
		description={El \emph{entorno f{\'i}sico} es el espacio tangible conformado por objetos materiales, condiciones ambientales y elementos naturales o artificiales que pueden ser percibidos, medidos o alterados mediante \glspl{sensor}. En sistemas \gls{iot}, constituye la fuente primaria de \glspl{dato} para la \gls{monitorizacion} y el \gls{controlinteligente}.},
		type=\glsdefaulttype,
		symbol={🏞️},
		user1={Ambiente f{\'i}sico},
		user2={Espacio tangible},
		user3={Entorno natural o artificial},
		see={entorno}
	}
	
	\newglossaryentry{entornoindustrial}{
		name={Entorno industrial},
		text={entorno industrial},
		sort={entorno industrial},
		plural={entornos industriales},
		first={entorno industrial},
		firstplural={entornos industriales},
		description={Contexto de trabajo caracterizado por la presencia de maquinaria, procesos de producci{\'o}n, sistemas de control y recursos humanos dedicados a la manufactura o transformaci{\'o}n de bienes. Un \emph{entorno industrial} integra infraestructuras f{\'i}sicas, tecnol{\'o}gicas y organizacionales, donde la automatizaci{\'o}n y la eficiencia son factores clave. En el marco de los \glspl{sibiot}, estos entornos incorporan \glspl{sensor}, \glspl{actuador}, sistemas de \gls{monitorizacion} y \gls{control} en \gls{tiemporeal}, con el fin de optimizar la producci{\'o}n, garantizar la \gls{seguridad}, mejorar la \gls{eficienciaoperativa} y favorecer la \gls{sostenibilidad} de los procesos industriales.},
		type=main,
		symbol={🏭},
		user1={Contexto de producci{\'o}n y manufactura},
		user2={Integraci{\'o}n de tecnolog{\'i}as IoT},
		user3={Base para la Industria 4.0},
		see={entorno,sibiot,automatizacion,eficienciaoperativa,sostenibilidad}
	}
	
	\newglossaryentry{entornointeligente}{
		name={Entorno inteligente},
		text={entorno inteligente},
		sort={entornointeligente},
		plural={entornos inteligentes},
		first={entorno inteligente},
		firstplural={entornos inteligentes},
		description={Un entorno inteligente integra \glspl{dispositivo} interconectados, sensores y sistemas de control capaces de recopilar, procesar y analizar \gls{informacion} en \gls{tiemporeal}. Estos entornos pueden encontrarse en hogares, ciudades, industrias o instituciones educativas, donde contribuyen a mejorar la eficiencia energ{\'e}tica, la seguridad, la accesibilidad y la experiencia de los usuarios. En el contexto de los \glspl{sibiot}, los entornos inteligentes representan una aplicaci{\'o}n directa de la \gls{automatizacion} y la \gls{tomadecisionesinformadas}.},
		type=\glsdefaulttype,
		symbol={🏙️},
		user1={Smart environment},
		user2={Espacio optimizado mediante tecnolog{\'i}as inteligentes},
		user3={Integraci{\'o}n de IoT, AI y automatizaci{\'o}n en espacios f{\'i}sicos o virtuales},
		see={iot,automatizacion,tiemporeal,dispositivo,sibiot}
	}
	
	\newglossaryentry{entornoiot}
	{
		name={Entorno IoT},
		text={entorno IoT},
		sort={entorno iot},
		plural={entornos IoT},
		first={entorno IoT},
		firstplural={entornos IoT},
		description={Un \emph{entorno IoT} es un ambiente compuesto por \glspl{dispositivoiot} interconectados, \glspl{red} de comunicaci{\'o}n, servicios digitales y usuarios que interact{\'u}an en \gls{tiemporeal} para capturar, transmitir y procesar \glspl{dato} en un sistema \gls{iot}. Estos entornos integran \gls{conectividad}, \gls{seguridad}, plataformas de \gls{almacenamiento} y \gls{procesamientodatos}, habilitando aplicaciones en dominios como salud, industria, agricultura y \glspl{ciudadinteligente}.},
		type=\glsdefaulttype,
		symbol={🌐},
		user1={Ambiente IoT},
		user2={Contexto IoT},
		user3={IoT environment},
		see={ecosistemaiot}
	}
	
	\newglossaryentry{entornooperativo}{
		name={Entorno operativo},
		text={entorno operativo},
		sort={entorno operativo},
		plural={entornos operativos},
		first={entorno operativo},
		firstplural={entornos operativos},
		description={Conjunto de condiciones, recursos y configuraciones en las que un \gls{sistema}, aplicaci{\'o}n o infraestructura se ejecuta y cumple sus funciones en la pr{\'a}ctica. Un \emph{entorno operativo} incluye el hardware, el software, las redes de comunicaci{\'o}n, las pol{\'i}ticas de seguridad, los procesos de administraci{\'o}n y las interacciones con los usuarios. En el contexto de los \glspl{sibiot}, el entorno operativo abarca los \glspl{sensor}, \glspl{actuador}, plataformas de \gls{monitorizacion}, sistemas de \gls{control} y servicios en la nube que trabajan en conjunto para garantizar \gls{disponibilidad}, \gls{fiabilidad}, \gls{eficienciaoperativa} y \gls{sostenibilidad} en escenarios reales de despliegue.},
		type=main,
		symbol={⚙️},
		user1={Condiciones de ejecuci{\'o}n real},
		user2={Infraestructura y configuraci{\'o}n del sistema},
		user3={Base de despliegue en SIbIoT},
		see={entorno,sibiot,sistema,monitorizacion,eficienciaoperativa}
	}
	
	\newglossaryentry{entornovirtualizado}
	{
		name={Entorno virtualizado},
		text={entorno virtualizado},
		sort={entorno virtualizado},
		plural={entornos virtualizados},
		first={entorno virtualizado},
		firstplural={entornos virtualizados},
		description={Un \emph{entorno virtualizado} es una infraestructura computacional en la que los recursos f{\'i}sicos, como \glspl{servidor}, \gls{almacenamiento} y \glspl{red}, son abstra{\'i}dos mediante tecnolog{\'i}as de virtualizaci{\'o}n. Esto permite la creaci{\'o}n de m{\'u}ltiples entornos l{\'o}gicos independientes que comparten el mismo hardware f{\'i}sico subyacente, optimizando la utilizaci{\'o}n de recursos y mejorando la \gls{escalabilidad}.},
		type=\glsdefaulttype,
		symbol={🖥️},
		user1={Virtualizaci{\'o}n},
		user2={Infraestructura virtual},
		user3={Entorno l{\'o}gico},
		see={entornodistribuido}
	}
	
	\newglossaryentry{entregable}{
		name={Entregable},
		text={entregable},
		sort={entregable},
		plural={entregables},
		first={entregable},
		firstplural={entregables},
		description={Producto verificable generado al finalizar una fase o iteraci{\'o}n de desarrollo. Puede consistir en c{\'o}digo, documentaci{\'o}n, modelo validado o unidad funcional operativa.},
		type=main,
		symbol={📄},
		user1={Resultado verificable},
		user2={Producto incremental},
		user3={Hito metodol{\'o}gico},
		see={unidadfuncional,ciclodesarrollo}
	}
	
	\newglossaryentry{entrenamientomodelos}{
		name={Entrenamiento de modelos},
		text={entrenamiento de modelos},
		plural={entrenamientos de modelos},
		first={Entrenamiento de modelos},
		firstplural={Entrenamientos de modelos},
		see={modelo,ia,aprendizajeautomatico},
		sort={entrenamientomodelos},
		description={Proceso mediante el cual un modelo anal{\'i}tico o de aprendizaje autom{\'a}tico ajusta sus par{\'a}metros a partir de datos hist{\'o}ricos (t{\'i}picamente curados), con el objetivo de capturar patrones o relaciones. En SIbIoT, el entrenamiento suele ejecutarse en infraestructura central o en la nube y debe registrarse con datos, par{\'a}metros y versiones para asegurar reproducibilidad y auditor{\'i}a}
	}
	
	\newglossaryentry{equipo}{
		name={Equipo},
		text={equipo},
		sort={equipo},
		plural={equipos},
		first={equipo},
		firstplural={equipos},
		description={Grupo organizado de personas que colaboran en el an{\'a}lisis, dise{\~n}o, implementaci{\'o}n y validaci{\'o}n de un sistema. En metodolog{\'i}as {\'a}giles, el equipo es autoorganizado y multidisciplinario.},
		type=main,
		symbol={👥},
		user1={Colaboraci{\'o}n},
		user2={Trabajo interdisciplinario},
		user3={Autoorganizaci{\'o}n},
		see={construccioncolaborativa}
	}
	
	\newglossaryentry{equipoagil}
	{
		name={Equipo {\'a}gil},
		text={equipo {\'a}gil},
		sort={equipo agil},
		plural={equipos {\'a}giles},
		first={equipo {\'a}gil},
		firstplural={equipos {\'a}giles},
		description={Un \emph{equipo {\'a}gil} es un grupo autoorganizado y multifuncional de personas que colaboran estrechamente para desarrollar soluciones de forma iterativa e incremental, siguiendo los principios del Manifiesto {\'A}gil. Se caracteriza por su adaptaci{\'o}n al cambio, comunicaci{\'o}n constante, responsabilidad compartida y entrega continua de valor. En proyectos \gls{iot}, los equipos {\'a}giles suelen incluir desarrolladores de hardware y software, arquitectos, ingenieros de datos y especialistas en despliegue y pruebas.},
		type=\glsdefaulttype,
		symbol={🤝},
		user1={Agile team},
		user2={Grupo {\'a}gil},
		user3={Equipo de desarrollo iterativo},
		see={desarrolloagilsoftware}
	}
	
	\newglossaryentry{erp}
	{
		name={ERP},
		text={ERP},
		sort={erp},
		firstplural={sistemas ERP},
		plural={sistemas ERP},
		first={planificaci{\'o}n de recursos empresariales (\textit{Enterprise Resource Planning}, ERP)},
		firstplural={sistemas ERP},
		description={Un \emph{ERP} es un sistema de informaci{\'o}n integrado que permite gestionar y automatizar los procesos clave de una organizaci{\'o}n, incluyendo finanzas, recursos humanos, log{\'i}stica, producci{\'o}n y ventas. Los sistemas ERP favorecen la eficiencia operativa, la trazabilidad de \glspl{dato} y la \gls{tomadecisiones} estrat{\'e}gica mediante una base de datos unificada y m{\'o}dulos especializados.},
		type=\glsdefaulttype,
		symbol={📊},
		user1={Enterprise Resource Planning},
		user2={Sistema de gesti{\'o}n empresarial},
		user3={Planificador de recursos},
		see={enterpriseapp}
	}
	
	\newglossaryentry{error}{
		name={Error},
		text={error},
		sort={error},
		plural={errores},
		first={error},
		firstplural={errores},
		description={Condici{\'o}n o defecto que conduce a un resultado incorrecto o comportamiento no deseado en un sistema. En ingenier{\'i}a de software se distingue de \emph{falla} como manifestaci{\'o}n observable.},
		type=main,
		symbol={❌},
		user1={Defecto},
		user2={Resultado incorrecto},
		user3={Comportamiento no deseado},
		see={falla,testabilidad}
	}
	
	\newglossaryentry{esb}
	{
		name={{ESB}},
		text={{ESB}},
		sort={esb},
		plural={{ESBs}},
		first={bus de servicios empresariales (\textit{Enterprise Service Bus}, ESB)},
		firstplural={{ESBs}},
		description={Un \emph{Enterprise Service Bus} (\gls{esb}) es un patr{\'o}n de integraci{\'o}n que provee mediaci{\'o}n (ruteo, transformaci{\'o}n, aplicaci{\'o}n de pol{\'\i}ticas) entre \glspl{servicio} para reducir acoplamiento y centralizar preocupaciones de integraci{\'o}n en sistemas empresariales.},
		type=\glsdefaulttype,
		symbol={🚌},
		user1={Mediaci{\'o}n},
		user2={Integraci{\'o}n},
		user3={Ruteo y transformaci{\'o}n},
		see={soa,api,mqtt}
	}
	
	\newglossaryentry{escahorizontal}
	{
		name={Escalabilidad horizontal},
		text={escalabilidad horizontal},
		sort={escalabilidad horizontal},
		plural={estrategias de escalabilidad horizontal},
		first={escalabilidad horizontal},
		firstplural={estrategias de escalabilidad horizontal},
		description={La \emph{escalabilidad horizontal} consiste en potenciar el rendimiento general de un \gls{sistema} a{\~n}adiendo m{\'a}s \glspl{nodo}, instancias o \glspl{componente} (de \gls{hardware} y/o \gls{software}) que trabajan en paralelo. Este enfoque permite distribuir la carga, mejorar la \gls{toleranciafallos} y aumentar la capacidad de procesamiento en \glspl{arquitecturadistribuida}.},
		type=\glsdefaulttype,
		symbol={➕},
		user1={Scale out},
		user2={Ampliaci{\'o}n horizontal},
		user3={Crecimiento por nodos},
		see={escalabilidad}
	}
	
	\newglossaryentry{escalabilidad}
	{
		name={Escalabilidad},
		text={escalabilidad},
		sort={escalabilidad},
		plural={escalabilidades},
		first={escalabilidad},
		firstplural={escalabilidades},
		description={La \emph{escalabilidad} es la capacidad de un \gls{sistema} para adaptarse en rendimiento a medida que aumenta la carga de trabajo, el n{\'u}mero de usuarios o los \glspl{dato}. Puede lograrse mediante escalabilidad horizontal (m{\'a}s nodos) o vertical (m{\'a}s recursos en un mismo nodo). En \glspl{sibiot}, la escalabilidad es fundamental para manejar el crecimiento de \glspl{dispositivoiot} y garantizar \gls{interoperabilidad}, rendimiento y continuidad operativa.},
		type=\glsdefaulttype,
		symbol={📈},
		user1={Scalability},
		user2={Capacidad de crecimiento},
		user3={Adaptabilidad de carga},
		see={escahorizontal,escavertical}
	}
	
	\newglossaryentry{escalabilidadanalitica}
	{
		name={Escalabilidad anal{\'i}tica},
		text={escalabilidad anal{\'i}tica},
		sort={escalabilidad analitica},
		plural={escalabilidades anal{\'i}ticas},
		first={escalabilidad anal{\'i}tica},
		firstplural={escalabilidades anal{\'i}ticas},
		description={La \emph{escalabilidad anal{\'i}tica} es la propiedad que permite ampliar almacenamiento y procesamiento conforme aumenta el volumen, velocidad y diversidad de datos sin degradar el rendimiento. En sistemas \gls{sibiot}, esta capacidad se implementa mediante arquitecturas distribuidas y procesamiento paralelo, asegurando continuidad operativa y an{\'a}lisis oportuno ante crecimiento progresivo de dispositivos y fuentes.},
		type=\glsdefaulttype,
		symbol={📈},
		user1={Procesamiento distribuido},
		user2={Elasticidad},
		user3={Arquitectura escalable},
		see={bigdata,procesamientoanalitico,infraestructura}
	}
	
	\newglossaryentry{escalable}{
		name={Escalable},
		text={escalable},
		sort={escalable},
		plural={escalables},
		first={escalable},
		firstplural={escalables},
		description={El t{\'e}rmino \emph{escalable} describe la capacidad de un \gls{sistema}, \gls{aplicacion} o componente de aumentar o disminuir su desempe{\~n}o de manera proporcional a la adici{\'o}n o reducci{\'o}n de recursos. Un dise{\~n}o escalable mantiene niveles aceptables de \gls{latencia}, \gls{confiabilidad} y coste a medida que crece la demanda. Puede lograrse mediante \gls{escahorizontal} (a{\~n}adiendo m{\'a}s instancias o nodos) y/o \gls{escavertical} (incrementando recursos en un nodo existente). En \glspl{sibiot}, la escalabilidad es esencial para absorber el crecimiento de \glspl{sensor} y \glspl{dispositivointeligente} heterog{\'e}neos, garantizando \gls{interoperabilidad}, continuidad y eficiencia operativa.},
		type=main,
		symbol={📈},
		user1={Crecimiento horizontal y vertical},
		user2={Mantenimiento de rendimiento y coste},
		user3={Clave en dise{\~n}o de SIbIoT},
		see={escalabilidad,escahorizontal,escavertical,interoperabilidad,sibiot}
	}
	
	\newglossaryentry{escavertical}
	{
		name={Escalabilidad vertical},
		text={escalabilidad vertical},
		sort={escalabilidad vertical},
		plural={estrategias de escalabilidad vertical},
		first={escalabilidad vertical},
		firstplural={estrategias de escalabilidad vertical},
		description={La \emph{escalabilidad vertical} consiste en mejorar el rendimiento de un \gls{sistema} a{\~n}adiendo recursos m{\'a}s potentes a una misma instancia, como mayor capacidad de CPU, memoria o almacenamiento. Es com{\'u}n en entornos tradicionales, aunque puede generar l{\'i}mites f{\'i}sicos y altos costos frente a la escalabilidad horizontal.},
		type=\glsdefaulttype,
		symbol={⬆️},
		user1={Scale up},
		user2={Ampliaci{\'o}n vertical},
		user3={Crecimiento por recursos},
		see={escalabilidad,escahorizontal}
	}
	
	\newglossaryentry{esfuerzo}
	{
		name={Esfuerzo},
		text={esfuerzo},
		sort={esfuerzo},
		description={Magnitud f{\'i}sica que representa la fuerza interna distribuida por unidad de {\'a}rea en el interior de un material sometido a cargas externas. 
			Se expresa como la respuesta de un cuerpo s{\'o}lido ante la aplicaci{\'o}n de una fuerza, y es fundamental para el an{\'a}lisis estructural y mec{\'a}nico. 
			En los \glspl{sibiot}, el esfuerzo se monitoriza mediante \glspl{sensor} piezorresistivos o piezoel{\'e}ctricos, lo que permite detectar deformaciones, tensiones y condiciones de sobrecarga en componentes mec{\'a}nicos, asegurando la integridad de los sistemas aut{\'o}nomos e industriales.},
		plural={esfuerzos},
		symbol={\ensuremath{\sigma}},
		first={Esfuerzo (\ensuremath{\sigma})},
		see=[Relacionado con:]{fuerza,presion,material,deformacion}
	}
	
	\newglossaryentry{esp}
	{
		name={ESP},
		text={ESP},
		sort={esp},
		firstplural={ESPs},
		plural={ESPs},
		first={plataforma de Espressif (\textit{Espressif Systems Platform}, ESP)},
		description={\emph{ESP} es una plataforma de desarrollo basada en \glspl{microcontrolador} de la empresa Espressif Systems. Se caracteriza por su bajo consumo, conectividad Wi-Fi/Bluetooth y soporte para entornos embebidos como ESP8266 y ESP32, ampliamente utilizada en proyectos \gls{iot}.},
		type=\glsdefaulttype,
		symbol={📡},
		user1={Espressif Systems Platform},
		user2={Plataforma ESP},
		user3={Microcontroladores ESP},
		see={esp32}
	}
	
	\newglossaryentry{esp32}
	{
		name={ESP32},
		text={ESP32},
		sort={esp32},
		plural={ESP32},
		first={ESP32},
		firstplural={ESP32},
		description={El \emph{ESP32} es un \gls{microcontrolador} de bajo costo y alta integraci{\'o}n desarrollado por Espressif Systems, que incluye conectividad Wi-Fi y Bluetooth. Es ampliamente utilizado en soluciones \gls{iot} por su eficiencia energ{\'e}tica y capacidad de procesamiento.},
		type=\glsdefaulttype,
		symbol={📶},
		user1={Microcontrolador ESP32},
		user2={Chip ESP32},
		user3={Dispositivo ESP32},
		see={esp}
	}
	
	\newglossaryentry{esp32cam}
	{
		name={ESP32-CAM},
		text={ESP32-CAM},
		sort={esp32cam},
		plural={ESP32-CAM},
		first={ESP32-CAM},
		firstplural={ESP32-CAM},
		description={El \emph{ESP32-CAM} es un m{\'o}dulo de desarrollo basado en el \gls{esp32}, que incorpora una c{\'a}mara digital (OV2640), conectividad Wi-Fi/Bluetooth y ranura para tarjeta microSD. Se utiliza en aplicaciones \gls{iot} que requieren visi{\'o}n artificial, videovigilancia, \gls{reconocimientofacial} o captura de im{\'a}genes en \gls{tiemporeal}, destacando por su bajo costo y versatilidad.},
		type=\glsdefaulttype,
		symbol={📷},
		user1={M{\'o}dulo ESP32-CAM},
		user2={C{\'a}mara IoT ESP32},
		user3={ESP32 con c{\'a}mara},
		see={esp32}
	}
	
	\newglossaryentry{esp8266}
	{
		name={ESP8266},
		text={ESP8266},
		sort={esp8266},
		plural={ESP8266},
		first={ESP8266},
		firstplural={ESP8266},
		description={El \emph{ESP8266} es un \gls{microcontrolador} de bajo costo fabricado por Espressif Systems, ampliamente utilizado en \glspl{sibiot} por su conectividad Wi-Fi integrada. Dispone de un procesador de 32 bits, pines GPIO programables y soporte para TCP/IP, lo que lo hace adecuado para aplicaciones embebidas que requieren comunicaci{\'o}n inal{\'a}mbrica eficiente y compacta.},
		type=\glsdefaulttype,
		symbol={📡},
		user1={Microcontrolador ESP8266},
		user2={Chip ESP8266},
		user3={Dispositivo ESP8266},
		see={esp}
	}
	
	\newglossaryentry{especificaciontecnica}{
		name={Especificaci{\'o}n t{\'e}cnica},
		text={especificaci{\'o}n t{\'e}cnica},
		sort={especificacion tecnica},
		plural={especificaciones t{\'e}cnicas},
		first={especificaci{\'o}n t{\'e}cnica},
		firstplural={especificaciones t{\'e}cnicas},
		description={Definici{\'o}n formal y precisa de requisitos, restricciones, estructuras y propiedades que debe cumplir un sistema. Constituye la referencia normativa para su construcci{\'o}n, validaci{\'o}n y mantenimiento.},
		type=main,
		symbol={📑},
		user1={Documento formal},
		user2={Base contractual},
		user3={Referencia de verificaci{\'o}n},
		see={modelo,restriccion,verificacion}
	}
	
	\newglossaryentry{estabilidad}
	{
		name={Estabilidad},
		text={estabilidad},
		sort={estabilidad},
		plural={estabilidades},
		first={estabilidad},
		firstplural={estabilidades},
		description={La \emph{estabilidad} es la capacidad de un \gls{sistema} para mantener su comportamiento funcional esperado frente a perturbaciones internas o externas. En sistemas \gls{iot}, implica que los \glspl{dispositivo}, \glspl{red} y servicios contin{\'u}en operando correctamente a lo largo del tiempo, incluso ante fluctuaciones de red, fallos de energ{\'i}a o variaciones ambientales.},
		type=\glsdefaulttype,
		symbol={⚖️},
		user1={Fiabilidad din{\'a}mica},
		user2={Robustez},
		user3={System stability},
		see={confiabilidad}
	}
	
	\newglossaryentry{estable}{
		name={Estable},
		text={estable},
		sort={estable},
		plural={estables},
		first={estable},
		firstplural={estables},
		description={Adjetivo que describe a un \gls{sistema}, componente o proceso que mantiene su funcionamiento correcto y predecible en el tiempo, sin experimentar fallos inesperados, fluctuaciones excesivas ni comportamientos err{\'a}ticos. Un sistema \emph{estable} conserva su desempe{\~n}o bajo condiciones normales de uso e incluso frente a perturbaciones moderadas. En el contexto de los \glspl{sibiot}, la estabilidad est{\'a} vinculada a la \gls{fiabilidad}, la \gls{resiliencia} y la \gls{disponibilidad}, ya que asegura que los datos capturados por \glspl{sensor}, los procesos de \gls{monitorizacion} y las acciones de \glspl{actuador} se ejecuten de manera coherente y continua, permitiendo la confianza en la \gls{tomadecisiones} en \gls{tiemporeal}.},
		type=main,
		symbol={⚖️},
		user1={Caracter{\'i}stica de consistencia operativa},
		user2={Resistencia a fallos o variaciones},
		user3={Base de la confianza en sistemas IoT},
		see={fiabilidad,resiliencia,disponibilidad,sibiot}
	}
	
	\newglossaryentry{estandar}{
		name={Est{\'a}ndar},
		text={est{\'a}ndar},
		sort={estandar},
		plural={est{\'a}ndares},
		first={est{\'a}ndar},
		firstplural={est{\'a}ndares},
		description={Un \emph{est{\'a}ndar} es un conjunto de reglas, especificaciones o criterios acordados por organismos de normalizaci{\'o}n, comunidades t{\'e}cnicas o la industria, cuyo objetivo es garantizar la \gls{interoperabilidad}, la \gls{compatibilidad}, la calidad y la seguridad en el dise{\~n}o, desarrollo y uso de sistemas. En \glspl{sibiot}, los est{\'a}ndares definen protocolos de \gls{comunicacion}, arquitecturas, formatos de datos y \glspl{protocoloseguridad}, permitiendo la integraci{\'o}n de dispositivos y plataformas heterog{\'e}neas de manera fiable y escalable. Ejemplos incluyen IEEE 802.15.4, MQTT, CoAP, ISO/IEC 30141 y los lineamientos de \gls{buenaspracticas}.},
		type=main,
		symbol={📏},
		user1={Reglas y especificaciones consensuadas},
		user2={Garantizan interoperabilidad y compatibilidad},
		user3={Fundamentales en IoT y sistemas distribuidos},
		see={protocolocomunicacion,interoperabilidad,compatibilidad,buenaspracticas,sibiot}
	}
	
	\newglossaryentry{estandarizacion}{
		name={Estandarizaci{\'o}n},
		text={estandarizaci{\'o}n},
		sort={estandarizacion},
		plural={estandarizaciones},
		first={estandarizaci{\'o}n},
		firstplural={estandarizaciones},
		description={Proceso mediante el cual se establecen normas, especificaciones y pr{\'a}cticas comunes para asegurar interoperabilidad, compatibilidad y calidad entre sistemas, dispositivos y plataformas. En el ecosistema \gls{iot}, la estandarizaci{\'o}n facilita la integraci{\'o}n de tecnolog{\'i}as heterog{\'e}neas, promueve la portabilidad de soluciones y reduce la dependencia de proveedores espec{\'i}ficos.},
		type=main,
		symbol={📏},
		user1={Interoperabilidad},
		user2={Compatibilidad tecnol{\'o}gica},
		user3={Normativas t{\'e}cnicas},
		see={interoperabilidad,iso,arquitectura}
	}
	
	\newglossaryentry{estructura}{
		name={Estructura},
		text={estructura},
		sort={estructura},
		plural={estructuras},
		first={estructura},
		firstplural={estructuras},
		description={Organizaci{\'o}n interna de componentes y relaciones que conforman un sistema. Define la disposici{\'o}n arquitect{\'o}nica que soporta propiedades funcionales y no funcionales.},
		type=main,
		symbol={🏛️},
		user1={Organizaci{\'o}n interna},
		user2={Arquitectura l{\'o}gica},
		user3={Relaci{\'o}n entre componentes},
		see={arquitectura,modelo}
	}
	
	\newglossaryentry{estructurametodologica}{
		name={Estructura metodol{\'o}gica},
		text={estructura metodol{\'o}gica},
		sort={estructurametodologica},
		plural={estructuras metodol{\'o}gicas},
		first={estructura metodol{\'o}gica},
		firstplural={estructuras metodol{\'o}gicas},
		description={Organizaci{\'o}n formal de un enfoque de trabajo que especifica fases, actividades, roles, artefactos, reglas de transici{\'o}n y criterios de control para guiar el desarrollo y la gesti{\'o}n de una soluci{\'o}n. Define qu{\'e} se produce, cu{\'a}ndo se produce, c{\'o}mo se valida y qui{\'e}n participa, proporcionando gobernanza t{\'e}cnica y trazabilidad del proceso. En proyectos IoT/SIbIoT, la estructura metodol{\'o}gica incorpora coordinaci{\'o}n entre desarrollo de software, integraci{\'o}n de hardware, pruebas sobre campo, despliegue y operaci{\'o}n continua.},
		type=main,
		symbol={🧭},
		user1={Fases, roles y artefactos},
		user2={Reglas y criterios de control},
		user3={Gobernanza del desarrollo},
		see={metodologia,metodologiadesarrollo,iterativa,incremental,ciclodesarrollo}
	}
	
	\newglossaryentry{ethernet}
	{
		name={Ethernet},
		text={Ethernet},
		sort={ethernet},
		plural={Ethernet},
		first={Ethernet},
		firstplural={Ethernet},
		description={\emph{Ethernet} es una tecnolog{\'i}a de \gls{red} por cable utilizada en \glspl{lan} para la transmisi{\'o}n de \glspl{dato}. Opera normalmente sobre cables de par trenzado o fibra {\'o}ptica y define protocolos para garantizar velocidad, confiabilidad y compatibilidad entre \glspl{dispositivoconectado}.},
		type=\glsdefaulttype,
		symbol={🔌},
		user1={LAN cableada},
		user2={IEEE 802.3},
		user3={Ethernet cable},
		see={lan}
	}
	
	\newglossaryentry{etiquetaelectronica}{
		name={etiqueta electr{\'o}nica},
		text={etiqueta electr{\'o}nica},
		sort={etiqueta electronica},
		plural={etiquetas electr{\'o}nicas},
		first={etiqueta electr{\'o}nica},
		firstplural={etiquetas electr{\'o}nicas},
		description={Dispositivo o componente digital que almacena y presenta informaci{\'o}n de manera din{\'a}mica mediante tecnolog{\'i}as como RFID, NFC, c{\'o}digos bidimensionales o pantallas de tinta electr{\'o}nica. Una \emph{etiqueta electr{\'o}nica} permite identificar, rastrear o actualizar datos sin intervenci{\'o}n manual, habilitando sincronizaci{\'o}n remota y automatizaci{\'o}n en procesos log{\'i}sticos, inventarios y sistemas \glspl{sibiot}. Su uso facilita la trazabilidad, la gesti{\'o}n inteligente de activos y la integraci{\'o}n en arquitecturas de monitorizaci{\'o}n distribuida},
		type=main,
		symbol={🏷️},
		user1={Identificaci{\'o}n digital y trazabilidad},
		user2={Actualizaci{\'o}n remota mediante RFID, NFC o tinta electr{\'o}nica},
		user3={Integraci{\'o}n en inventarios y sistemas SIbIoT},
		see={rfid,nfc,trazabilidad,sibiot}
	}
	
	\newglossaryentry{etl}
	{
		name={ETL},
		text={ETL},
		first={Extraer, Transformar y Cargar (\textit{Extract, Transform, Load}, ETL)},
		sort={etl},
		plural={procesos ETL},
		firstplural={procesos ETL},
		description={\emph{ETL} es un proceso de integraci{\'o}n de \glspl{dato} que consiste en tres fases: extracci{\'o}n desde fuentes heterog{\'e}neas, transformaci{\'o}n para adaptarlos a un formato com{\'u}n o aplicar reglas de negocio, y carga en un sistema de destino, generalmente una \gls{basedatos} o almac{\'e}n de datos. En \glspl{sibiot}, el ETL es fundamental para consolidar informaci{\'o}n generada por \glspl{sensor}, normalizarla y almacenarla en la \gls{nube}, garantizando \gls{consistencia}, \gls{calidad} y disponibilidad para \glspl{modelopredictivo} y \gls{tomadecisiones}.},
		type=\glsdefaulttype,
		symbol={🔄},
		user1={Extracci{\'o}n de datos},
		user2={Transformaci{\'o}n y limpieza},
		user3={Carga en almacenes de datos},
		see={basedatos}
	}
	
	\newglossaryentry{evento}
	{
		name={Evento},
		text={evento},
		sort={evento},
		plural={eventos},
		first={evento},
		firstplural={eventos},
		description={Los \emph{eventos} son unidades de informaci{\'o}n que representan la ocurrencia de un hecho significativo en un \gls{sistema} o en su \gls{entorno}, como un cambio de estado, una acci{\'o}n del usuario o una lectura de \gls{sensor}. En \glspl{sibiot} y arquitecturas reactivas, los eventos son fundamentales para activar flujos de \gls{procesamientodatos}, desencadenar respuestas autom{\'a}ticas y mantener la comunicaci{\'o}n eficiente y desacoplada entre componentes distribuidos.},
		type=\glsdefaulttype,
		symbol={📢},
		user1={Se{\~n}al},
		user2={Notificaci{\'o}n},
		user3={Event},
		see={arquitecturaeventos}
	}
		
	\newglossaryentry{evolucion}{
		name={Evoluci{\'o}n},
		text={evoluci{\'o}n},
		sort={evolucion},
		plural={procesos de evoluci{\'o}n},
		first={evoluci{\'o}n},
		firstplural={procesos de evoluci{\'o}n},
		description={Proceso gradual de transformaci{\'o}n y mejora de sistemas, tecnolog{\'i}as o paradigmas a lo largo del tiempo. En el contexto de IoT y \glspl{sibiot}, la \emph{evoluci{\'o}n} describe la transici{\'o}n desde sistemas cerrados de control hacia arquitecturas abiertas, distribuidas y orientadas a datos.},
		type=main,
		symbol={🔄},
		user1={Cambio progresivo},
		user2={Desarrollo tecnol{\'o}gico},
		user3={Madurez del sistema},
		see={arquitecturaabierta,integraciondatos}
	}
	
	\newglossaryentry{evolucioncontrolada}{
		name={Evoluci{\'o}n controlada},
		text={evoluci{\'o}n controlada},
		sort={evolucioncontrolada},
		plural={evoluciones controladas},
		first={evoluci{\'o}n controlada},
		firstplural={evoluciones controladas},
		description={Proceso de cambio planificado y trazable de un sistema, mediante el cual se introducen mejoras, correcciones o adaptaciones preservando coherencia arquitect{\'o}nica, compatibilidad y calidad. Implica gesti{\'o}n de versiones, evaluaci{\'o}n de impacto, mantenimiento de contratos de interfaz, verificaci{\'o}n automatizada y criterios de aceptaci{\'o}n. En \glspl{sibiot}, la evoluci{\'o}n controlada abarca firmware, configuraciones de red, servicios en la nube, esquemas de datos y pipelines, procurando continuidad operativa y mitigando riesgos por cambios en dispositivos heterog{\'e}neos desplegados.},
		type=main,
		symbol={🧩},
		user1={Cambio planificado y trazable},
		user2={Preservaci{\'o}n de coherencia},
		user3={Compatibilidad y calidad},
		see={gestioncambios,versionado,consistencia,verificacion,despliegue,infraestructuradistribuida}
	}
	
	\newglossaryentry{exactitud}{
		name={Exactitud},
		text={exactitud},
		sort={exactitud},
		plural={exactitudes},
		first={exactitud},
		firstplural={exactitudes},
		description={La \emph{exactitud} es el grado en que el valor medido o calculado por un sistema coincide con el valor real o de referencia. A diferencia de la \gls{precision}, que se refiere a la repetibilidad de las mediciones, la exactitud refleja la cercan{\'i}a al valor verdadero. En sistemas \gls{iot} y \glspl{sibiot}, la exactitud de los \glspl{sensor} y procesos de \gls{procesamientodatos} es fundamental para garantizar la \gls{fiabilidad}, la \gls{integridad} de la informaci{\'o}n y la correcta \gls{tomadecisiones} en aplicaciones como salud, agricultura de \gls{precision}, log{\'i}stica o transporte inteligente.},
		type=main,
		symbol={🎯},
		user1={Cercan{\'i}a al valor real},
		user2={Diferencia respecto a la precisi{\'o}n},
		user3={Clave en IoT para decisiones fiables},
		see={precision,fiabilidad,integridad,tomadecisiones,sibiot}
	}
	
	\newglossaryentry{extensibilidad}{
		name={Extensibilidad},
		text={extensibilidad},
		sort={extensibilidad},
		plural={extensibilidades},
		first={extensibilidad},
		firstplural={extensibilidades},
		description={Propiedad arquitect{\'o}nica que permite ampliar o incorporar nuevas funcionalidades a un sistema sin modificar de manera significativa su estructura central ni comprometer su estabilidad. La extensibilidad se logra mediante principios como bajo acoplamiento, alta cohesi{\'o}n, dise{\~n}o basado en interfaces, uso de contratos estables y mecanismos de extensi{\'o}n controlada (plugins, APIs versionadas, microservicios). En ecosistemas digitales e infraestructuras \gls{iot}, la extensibilidad es fundamental para facilitar integraci{\'o}n de nuevos participantes, evoluci{\'o}n funcional progresiva y adaptaci{\'o}n a cambios tecnol{\'o}gicos sin degradar \gls{consistencia}, \gls{portabilidad} ni \gls{mantenibilidad}.},
		type=main,
		symbol={🧩},
		user1={Capacidad de ampliaci{\'o}n},
		user2={Dise{\~n}o desacoplado},
		user3={Arquitectura evolutiva},
		see={modularidad,interoperabilidad,evolucioncontrolada,arquitectura}
	}
	
	\newglossaryentry{falla}{
		name={Falla},
		text={falla},
		sort={falla},
		plural={fallas},
		first={falla},
		firstplural={fallas},
		description={Manifestaci{\'o}n observable de un comportamiento incorrecto del sistema durante su ejecuci{\'o}n, usualmente causada por uno o m{\'a}s \glspl{error} o condiciones del entorno.},
		type=main,
		symbol={🛑},
		user1={Manifestaci{\'o}n},
		user2={Ejecuci{\'o}n incorrecta},
		user3={Observabilidad},
		see={error,recuperacion}
	}
	
	\newglossaryentry{fallo}{
		name={Fallo},
		text={fallo},
		sort={fallo},
		plural={fallos},
		first={fallo},
		firstplural={fallos},
		description={Un fallo es una alteraci{\'o}n en el funcionamiento esperado de un sistema, componente o proceso dentro de un \gls{sibiot}. Puede originarse por defectos de hardware, errores de software, p{\'e}rdida de conectividad, fallas energ{\'e}ticas, vulnerabilidades de seguridad o condiciones ambientales adversas. Los fallos comprometen la disponibilidad, confiabilidad y desempe{\~n}o del sistema, afectando la recolecci{\'o}n de \glspl{dato}, la operaci{\'o}n de \glspl{sensor}, la transmisi{\'o}n en redes distribuidas y la ejecuci{\'o}n de algoritmos en tiempo real. La gesti{\'o}n eficiente de fallos implica detecci{\'o}n temprana, diagn{\'o}stico, tolerancia y recuperaci{\'o}n para garantizar continuidad operativa y resiliencia},
		type=\glsdefaulttype,
		symbol={⚠️},
		user1={Failure},
		user2={Malfuncionamiento o interrupci{\'o}n operativa},
		user3={Eventos que afectan la disponibilidad y confiabilidad del sistema},
		see={error,falla,dispositivo,sibiot,seguridad}
	}
	
	\newglossaryentry{feedbackcontrol}
	{
		name={Feedback Control},
		text={feedback control},
		sort={feedback control},
		plural={feedback control},
		first={feedback control},
		firstplural={feedback control},
		description={El \emph{feedback control} es un mecanismo de control en el cual las salidas de un sistema (resultados anal{\'i}ticos, eventos o indicadores) se retroalimentan como entradas para ajustar din{\'a}micamente su comportamiento. En un \gls{sibiot}, este principio permite que decisiones generadas por la \gls{analiticadatos} y el \gls{motorreglas} se traduzcan en comandos hacia \glspl{actuador}, cerrando el ciclo entre el entorno f{\'i}sico y los procesos organizacionales con \gls{trazabilidadoperacional}.},
		type=\glsdefaulttype,
		symbol={🔁},
		user1={Control retroalimentado},
		user2={Cierre de ciclo},
		user3={Ajuste din{\'a}mico},
		see={actuador,analiticadatos,motorreglas,sibiot}
	}
	
	\newglossaryentry{fenomenofisico}{
		name={Fen{\'o}meno f{\'i}sico},
		text={fen{\'o}meno f{\'i}sico},
		sort={fenomeno fisico},
		plural={fen{\'o}menos f{\'i}sicos},
		first={fen{\'o}meno f{\'i}sico},
		firstplural={fen{\'o}menos f{\'i}sicos},
		description={Evento, proceso o manifestaci{\'o}n observable en la naturaleza que puede describirse, medirse o cuantificarse mediante leyes f{\'i}sicas. Ejemplos de \emph{fen{\'o}menos f{\'i}sicos} son el movimiento, la temperatura, la presi{\'o}n, la luz, la humedad, el magnetismo o la corriente el{\'e}ctrica. En el contexto de los \glspl{sibiot}, los fen{\'o}menos f{\'i}sicos constituyen la fuente primaria de informaci{\'o}n, ya que son captados por \glspl{sensor} y transformados en \glspl{variablefisica} digitales que permiten su procesamiento, monitorizaci{\'o}n y control en \gls{tiemporeal}.},
		type=main,
		symbol={🌍},
		user1={Evento observable en la naturaleza},
		user2={Fuente de variables f{\'i}sicas},
		user3={Dato primario en SIbIoT},
		see={variablefisica,sensor,sibiot,control}
	}
	
	\newglossaryentry{fiabilidad}{
		name={Fiabilidad},
		text={fiabilidad},
		sort={fiabilidad},
		plural={fiabilidades},
		first={fiabilidad},
		firstplural={fiabilidades},
		description={Probabilidad de que un \gls{sistema}, componente o servicio desempe{\~n}e sus funciones requeridas bajo condiciones especificadas y durante un periodo determinado, sin fallos. La \emph{fiabilidad} se relaciona con tasas de fallo, tiempo medio entre fallos (MTBF) y tiempo medio de reparaci{\'o}n (MTTR); contribuye de forma directa a la \gls{disponibilidad} percibida y al nivel de servicio. En entornos \gls{iot} y \glspl{sibiot}, depende tanto de dise{\~n}o y pruebas como de \gls{redundancia}, \gls{toleranciafallos} y buenas \glspl{buenaspracticas} de operaci{\'o}n.},
		type=main,
		symbol={},
		user1={Calidad y confiabilidad operativa},
		user2={Tolerancia a fallos y redundancia},
		user3={M{\'e}tricas MTBF/MTTR},
		see={disponibilidad,buenaspracticas,escalable}
	}
	
	\newglossaryentry{filosofia}{
		name={Filosof{\'\i}a},
		text={filosof{\'\i}a},
		sort={filosofia},
		plural={filosof{\'\i}as},
		first={filosof{\'\i}a},
		firstplural={filosof{\'\i}as},
		description={Conjunto de principios rectores y supuestos conceptuales que fundamentan un enfoque de dise{\~n}o, desarrollo o gesti{\'o}n de sistemas. En ingenier{\'\i}a de software e IoT, la filosof{\'\i}a define prioridades (por ejemplo, simplicidad, verificabilidad, automatizaci{\'o}n, calidad, seguridad, trazabilidad), orientando decisiones metodol{\'o}gicas y arquitect{\'o}nicas. No prescribe necesariamente pasos detallados, pero determina criterios de valoraci{\'o}n y de elecci{\'o}n entre alternativas t{\'e}cnicas.},
		type=main,
		symbol={🧠},
		user1={Principios rectores},
		user2={Criterios para decidir},
		user3={Base conceptual del enfoque},
		see={principio,orientacion,estructurametodologica,decisionarquitectonica}
	}
	
	
	
	\newglossaryentry{filtradodatos}{
		name={Filtrado de datos},
		text={filtrado de datos},
		sort={filtradodatos},
		plural={filtrados de datos},
		first={filtrado de datos},
		firstplural={filtrados de datos},
		description={Proceso mediante el cual se seleccionan, depuran o transforman conjuntos de datos con el fin de eliminar ruido, valores at{\'\i}picos o informaci{\'o}n irrelevante, preservando aquellos datos que resultan significativos para el an{\'a}lisis. En los \glspl{sibiot}, el filtrado de datos es esencial para mejorar la calidad de la informaci{\'o}n proveniente de sensores, reducir errores de medici{\'o}n, optimizar el procesamiento posterior y garantizar decisiones m{\'a}s precisas en sistemas distribuidos.},
		type=main,
		symbol={🧹},
		user1={Depuraci{\'o}n de informaci{\'o}n},
		user2={Reducci{\'o}n de ruido},
		user3={Calidad de datos},
		see={normalizaciondatos,analiticadatos,analisisstreaming,gestiondatos}
	}
	
	\newglossaryentry{filtro}
	{
		name={Filtro},
		text={filtro},
		sort={filtro},
		plural={filtros},
		first={filtro},
		firstplural={filtros},
		description={Un \emph{filtro} es un mecanismo que permite seleccionar, transformar, atenuar o descartar determinados elementos de una se{\~n}al, conjunto de \glspl{dato} o informaci{\'o}n, en funci{\'o}n de criterios definidos. Puede implementarse de diversas formas: 
			\begin{itemize}
				\item En \gls{hardware}, como los filtros anal{\'o}gicos que eliminan ruido en se{\~n}ales el{\'e}ctricas.
				\item En software, como los filtros digitales (media m{\'o}vil, Kalman, Butterworth) aplicados a \glspl{dato} capturados por \glspl{sensor}.
				\item En bases de datos o sistemas de informaci{\'o}n, como condiciones o restricciones l{\'o}gicas que permiten recuperar solo registros espec{\'i}ficos que cumplen ciertos criterios.
				\item En interfaces y aplicaciones, como mecanismos de b{\'u}squeda, clasificaci{\'o}n o segmentaci{\'o}n de informaci{\'o}n para los usuarios.
			\end{itemize}
			En sistemas \gls{iot}, los filtros son fundamentales para asegurar la calidad, relevancia y eficiencia del \gls{procesamientodatos}, reduciendo redundancia, eliminando ruido y mejorando la \gls{tomadecisiones}.},
		type=\glsdefaulttype,
		symbol={⚙️},
		user1={Condici{\'o}n},
		user2={Restricci{\'o}n},
		user3={Filter},
		see={filtrodatos}
	}
	
	\newglossaryentry{filtrodatos}
	{
		name={Filtro de datos},
		text={filtro de datos},
		sort={filtro de datos},
		plural={filtros de datos},
		first={filtro de datos},
		firstplural={filtros de datos},
		description={Un \emph{filtro de datos} es un algoritmo o procedimiento dise{\~n}ado para refinar, depurar o transformar conjuntos de \glspl{dato} con el fin de mejorar su calidad, reducir redundancias o eliminar inconsistencias. En el contexto de \glspl{sibiot}, los filtros de datos permiten optimizar el \gls{procesamientodatos} proveniente de \glspl{sensor}, facilitando un an{\'a}lisis m{\'a}s preciso para la \gls{tomadecisiones}.},
		type=\glsdefaulttype,
		symbol={🔍},
		user1={Depuraci{\'o}n de datos},
		user2={Data filter},
		user3={Filtro inform{\'a}tico},
		see={dato}
	}
	
	\newglossaryentry{filtroruido}
	{
		name={Filtros de ruido},
		text={filtros de ruido},
		sort={filtros de ruido},
		plural={filtros de ruido},
		first={filtros de ruido},
		firstplural={filtros de ruido},
		description={Los \emph{filtros de ruido} son algoritmos o dispositivos que permiten atenuar o eliminar componentes no deseados de una se{\~n}al, conocidos como ruido, que pueden alterar la \gls{precision} de los \glspl{dato} capturados. En sistemas \gls{iot}, los filtros de ruido son fundamentales para mejorar la calidad de las mediciones provenientes de \glspl{sensor}, y se implementan tanto a nivel de \gls{hardware} (filtros anal{\'o}gicos) como de software (filtros digitales como media m{\'o}vil, Kalman o Butterworth), permitiendo un an{\'a}lisis m{\'a}s fiable y robusto de la informaci{\'o}n.},
		type=\glsdefaulttype,
		symbol={🔉},
		user1={Filtro digital},
		user2={Filtro anal{\'o}gico},
		user3={Noise filter},
		see={sensor}
	}
	
	\newglossaryentry{firebase}
	{
		name={Firebase},
		text={Firebase},
		sort={firebase},
		plural={Firebase},
		first={Firebase},
		firstplural={Firebase},
		description={\emph{Firebase} es una plataforma de desarrollo de aplicaciones m{\'o}viles y web de Google que proporciona servicios como base de datos en \gls{tiemporeal}, autenticaci{\'o}n, hosting, almacenamiento y an{\'a}lisis. En el contexto de proyectos \gls{iot}, se utiliza para centralizar \glspl{dato}, gestionar usuarios y desplegar interfaces de monitorizaci{\'o}n.},
		type=\glsdefaulttype,
		symbol={🔥},
		user1={Plataforma Firebase},
		user2={Google Firebase},
		user3={Servicios m{\'o}viles y web},
		see={nube}
	}
	
	% ===================== Firestore =====================
	\newglossaryentry{firestore}{
		name={Firestore},
		text={Firestore},
		sort={firestore},
		first={Firestore},
		description={Forma abreviada y de uso com{\'u}n para referirse a \gls{cloudfirestore}. En contextos t{\'e}cnicos y acad{\'e}micos, el t{\'e}rmino \emph{Firestore} se emplea como denominaci{\'o}n corta una vez introducido formalmente el servicio Cloud Firestore, manteniendo el mismo modelo de datos NoSQL orientado a documentos, capacidades de sincronizaci{\'o}n en tiempo real y escalabilidad autom{\'a}tica},
		type=main,
		symbol={🔥},
		see={cloudfirestore,nosql,cloudcomputing,sibiot}
	}
	
	
	\newglossaryentry{firmadigital}
	{
		name={Firma digital},
		text={firma digital},
		sort={firma digital},
		plural={firmas digitales},
		first={firma digital},
		firstplural={firmas digitales},
		description={La \emph{firma digital} es un mecanismo criptogr{\'a}fico utilizado para verificar la autenticidad y la \gls{integridad} de los \glspl{dato}. Se basa en la \gls{criptografiaasimetrica}, donde una clave privada genera la firma y la clave p{\'u}blica permite su verificaci{\'o}n. Es ampliamente utilizada en transacciones electr{\'o}nicas, comunicaciones seguras y en la protecci{\'o}n de la \gls{privacidad}.},
		type=\glsdefaulttype,
		symbol={✍️},
		user1={Digital signature},
		user2={Firma electr{\'o}nica avanzada},
		user3={Autenticaci{\'o}n criptogr{\'a}fica},
		see={criptografiaasimetrica}
	}
	
	\newglossaryentry{firmware}{
		name={Firmware},
		text={firmware},
		sort={firmware},
		plural={firmwares},
		first={firmware},
		firstplural={firmwares},
		description={El \emph{firmware} es un tipo de software especializado grabado en memorias no vol{\'a}tiles de \glspl{dispositivo}, como ROM, Flash o EEPROM. Se encarga de controlar funciones b{\'a}sicas del \gls{hardware}, proporcionar instrucciones de arranque y coordinar la comunicaci{\'o}n entre componentes internos. En \glspl{dispositivoiot}, el firmware gestiona la adquisici{\'o}n de \glspl{dato}, la interacci{\'o}n con \glspl{sensor} y \glspl{actuador}, y la conexi{\'o}n con \glspl{red}, siendo esencial para garantizar \gls{fiabilidad}, \gls{seguridad} y capacidad de actualizaci{\'o}n remota (\textit{OTA}, \textit{Over-The-Air}).},
		type=\glsdefaulttype,
		symbol={💾},
		user1={Software embebido},
		user2={Control de hardware},
		user3={Actualizaci{\'o}n OTA},
		see={hardware,software,flash,dispositivo,dispositivoiot}
	}
	
	\newglossaryentry{flash}
	{
		name={Flash},
		text={flash},
		sort={flash},
		plural={memorias flash},
		first={flash},
		firstplural={memorias flash},
		description={\emph{Flash} es un tipo de memoria no vol{\'a}til utilizada para almacenar \glspl{dato} de manera persistente. Es com{\'u}n en \glspl{microcontrolador} y \glspl{dispositivoiot} para guardar firmware, configuraciones y registros, con la ventaja de poder ser reescrita el{\'e}ctricamente. Representa un componente clave en la \gls{placadesarrollo} y en \glspl{sistemaembebido} modernos.},
		type=\glsdefaulttype,
		symbol={💾},
		user1={Memoria flash},
		user2={Almacenamiento flash},
		user3={Flash memory},
		see={almacenamiento}
	}
	
	\newglossaryentry{flexibilidad}
	{
		name={Flexibilidad},
		text={flexibilidad},
		sort={flexibilidad},
		plural={flexibilidades},
		first={flexibilidad},
		firstplural={flexibilidades},
		description={La flexibilidad es la capacidad de un \gls{sistema} o \gls{componente} para adaptarse a cambios en los \glspl{requisito}, en el \gls{entorno} o en las \glspl{condicion} de operaci{\'o}n sin que se vea comprometido su desempe{\~n}o esencial. En el contexto de los \glspl{sibiot}, la flexibilidad se asocia con la posibilidad de reconfigurar, ampliar o integrar nuevos \glspl{dispositivo} y \glspl{servicio}, garantizando la continuidad y evoluci{\'o}n del \gls{sistemainformatico}.},
		type=\glsdefaulttype,
		symbol={},
		user1={Adaptabilidad},
		user2={Resiliencia},
		user3={Capacidad de reconfiguraci{\'o}n},
		see={sibiot,escalabilidad,robustez,fiabilidad}
	}
	
	\newglossaryentry{fluido}{
		name={Fluido},
		text={fluido},
		sort={fluido},
		plural={fluidos},
		first={fluido},
		firstplural={fluidos},
		description={Sustancia que puede fluir y deformarse continuamente (l{\'i}quidos o gases). En instrumentaci{\'o}n, variables como presi{\'o}n, caudal y temperatura de un fluido son magnitudes com{\'u}nmente monitorizadas.},
		type=main,
		symbol={💧},
		user1={L{\'i}quidos y gases},
		user2={Caudal/presi{\'o}n},
		user3={Instrumentaci{\'o}n},
		see={monitorizacion,transductor}
	}
	
	\newglossaryentry{flujodatos}{
		name={Flujo de datos},
		text={flujo de datos},
		plural={flujos de datos},
		first={flujo de datos},
		firstplural={flujos de datos},
		see={dato,procesamiento},
		sort={flujodatos},
		description={Secuencia estructurada de etapas mediante las cuales los datos son capturados, ingeridos, procesados, analizados, almacenados y consumidos para la toma de decisiones. En SIbIoT, el flujo de datos conecta sensores, sistemas empresariales, anal{\'i}tica, visualizaci{\'o}n y ejecuci{\'o}n, asegurando trazabilidad y control a lo largo de todo el ciclo informacional}
	}
	
	\newglossaryentry{fog}{
		name={Computaci{\'o}n en la niebla (Fog Computing)},
		text={computaci{\'o}n en la niebla},
		sort={fog},
		plural={arquitecturas de computaci{\'o}n en la niebla},
		first={computaci{\'o}n en la niebla},
		firstplural={arquitecturas de computaci{\'o}n en la niebla},
		description={Paradigma de procesamiento distribuido que extiende las capacidades de la \gls{edgecomputing} y de la computaci{\'o}n en la nube hacia los dispositivos y nodos intermedios cercanos al origen de los datos. La \emph{computaci{\'o}n en la niebla} reduce la latencia, optimiza el tr{\'a}fico hacia la nube y permite respuestas en \gls{tiemporeal} al ejecutar funciones de an{\'a}lisis, almacenamiento temporal, filtrado y toma de decisiones en una capa intermedia. En los \glspl{sibiot}, la niebla act{\'u}a como infraestructura habilitadora para coordinar \glspl{sensor}, gestionar cargas din{\'a}micas, ejecutar \glspl{modelopredictivo} localmente y mejorar la \gls{eficienciaoperativa}, especialmente en entornos que requieren fiabilidad, baja latencia y continuidad operativa incluso con conectividad limitada hacia la nube.},
		type=main,
		symbol={🌫️},
		user1={Procesamiento distribuido cercano al origen de datos},
		user2={Reducci{\'o}n de latencia y optimizaci{\'o}n de tr{\'a}fico},
		user3={Soporte a sistemas IoT cr{\'i}ticos y de tiempo real},
		see={edgecomputing,nube,sibiot,modelopredictivo,tiemporeal}
	}
	
	% ===================== Fog Computing =====================
	\newglossaryentry{fogcomputing}
	{
		name={Fog Computing},
		text={fog computing},
		sort={fog computing},
		plural={modelos de fog computing},
		first={Computaci{\'o}n en la niebla (\textit{Fog Computing}},
		firstplural={modelos de fog computing},
		description={El \emph{computaci{\'o}n en la niebla o \textit{Fog Computing}} es un modelo de computaci{\'o}n distribuida que extiende los servicios de la \gls{nube} hacia el borde de la \gls{red}, acercando \gls{procesamientodatos}, \gls{almacenamiento} y an{\'a}lisis a los \glspl{dispositivo} que generan \glspl{dato}. Reduce \gls{latencia}, ahorra ancho de banda y mejora la resiliencia en \glspl{sibiot} sensibles al \gls{tiemporeal}.},
		type=\glsdefaulttype,
		symbol={🌫️},
		user1={Computaci{\'o}n en niebla},
		user2={Plataforma perimetral},
		user3={Capa intermedia nube\textendash borde},
		see={edgecomputing,nube}
	}
	
	% ===================== Font =====================
	\newglossaryentry{font}
	{
		name={Font},
		text={font},
		sort={font},
		plural={fonts},
		first={font},
		firstplural={fonts},
		description={Una \emph{font} es el estilo o tipo de letra utilizado en una \gls{interfaz} gr{\'a}fica o documento, que define atributos como familia, peso, tama{\~n}o y espaciado. En aplicaciones y \glspl{dashboard} de \gls{iot} influye en usabilidad, accesibilidad y consistencia visual.},
		type=\glsdefaulttype,
		symbol={🔤},
		user1={Tipo de letra},
		user2={Familia tipogr{\'a}fica},
		user3={Fuente},
		see={interfaz}
	}
	
	\newglossaryentry{formalizacion}{
		name={Formalizaci{\'o}n},
		text={formalizaci{\'o}n},
		sort={formalizacion},
		plural={formalizaciones},
		first={formalizaci{\'o}n},
		firstplural={formalizaciones},
		description={Proceso mediante el cual un concepto, procedimiento o arquitectura se expresa de manera precisa y rigurosa utilizando modelos, notaciones formales o marcos normativos claramente definidos. En el contexto de sistemas \gls{iot} y arquitecturas Edge--Cloud, la formalizaci{\'o}n del borde implica definirlo como un entorno programable con reglas, interfaces, pol{\'i}ticas y mecanismos estandarizados que permiten su gesti{\'o}n, despliegue y evoluci{\'o}n controlada.},
		type=main,
		symbol={📘},
		user1={Rigurosidad metodol{\'o}gica},
		user2={Definici{\'o}n expl{\'i}cita de reglas},
		user3={Base para automatizaci{\'o}n},
		see={modelo,especificaciontecnica,estandarizacion}
	}
	
	\newglossaryentry{fragmentacion}{
		name={Fragmentaci{\'o}n},
		text={fragmentaci{\'o}n},
		sort={fragmentacion},
		plural={fragmentaciones},
		first={fragmentaci{\'o}n},
		firstplural={fragmentaciones},
		description={La fragmentaci{\'o}n se refiere a la divisi{\'o}n l{\'o}gica o f{\'i}sica de un sistema, estructura o conjunto de datos para mejorar la eficiencia en su gesti{\'o}n. En los \glspl{sibiot}, este concepto se aplica al almacenamiento y procesamiento de la \gls{informacion}, permitiendo distribuir la carga entre nodos o subsistemas y optimizar el rendimiento general. Tambi{\'e}n puede emplearse para favorecer la resiliencia, reducir la latencia y aumentar la escalabilidad de los \glspl{entornoiot}.},
		type=\glsdefaulttype,
		symbol={🧩},
		user1={Fragmentation},
		user2={Divisi{\'o}n de un conjunto en partes m{\'a}s peque{\~n}as},
		user3={Optimizaci{\'o}n estructural o funcional mediante segmentaci{\'o}n},
		see={fragmentaciondatos,basedatosdistribuida,eficiencia,entornoiot}
	}
	
	\newglossaryentry{fragmentaciondatos}{
		name={Fragmentaci{\'o}n de datos},
		text={fragmentaci{\'o}n de datos},
		sort={fragmentaciondatos},
		plural={fragmentaciones de datos},
		first={fragmentaci{\'o}n de datos},
		firstplural={fragmentaciones de datos},
		description={La fragmentaci{\'o}n de datos es una t{\'e}cnica utilizada en \glspl{basedatosdistribuida} y \glspl{sibiot} para optimizar la gesti{\'o}n y el acceso a la \gls{informacion}. Este proceso puede realizarse de manera horizontal, dividiendo los registros de una tabla, o vertical, separando columnas espec{\'i}ficas. En los \glspl{sibiot}, la fragmentaci{\'o}n contribuye a equilibrar la carga de trabajo entre nodos, reducir la latencia y aumentar la disponibilidad, garantizando que los \glspl{dispositivo} y \glspl{servidor} gestionen eficientemente grandes vol{\'u}menes de datos distribuidos.},
		type=\glsdefaulttype,
		symbol={🧩},
		user1={Data fragmentation},
		user2={Divisi{\'o}n de datos para optimizar rendimiento y distribuci{\'o}n},
		user3={Estrategia para equilibrar la carga en sistemas distribuidos e IoT},
		see={basedatosdistribuida,distribucion,almacenamiento,sibiot,eficiencia}
	}
	
	\newglossaryentry{frontend}{
		name={Front-end},
		text={front-end},
		sort={frontend},
		plural={front-ends},
		first={front-end},
		firstplural={front-ends},
		description={Parte de un sistema de software o aplicaci{\'o}n encargada de la interacci{\'o}n directa con el usuario. El \emph{front-end} incluye la interfaz visual, los elementos gr{\'a}ficos y la l{\'o}gica de presentaci{\'o}n que permiten al usuario comunicarse con el sistema. Com{\'u}nmente se desarrolla con tecnolog{\'i}as como \gls{html}, \gls{css} y \gls{javascript}, y puede estar complementado por frameworks como Angular, React o Vue. En sistemas \gls{iot} y \glspl{sibiot}, el front-end permite la visualizaci{\'o}n de datos en tiempo real, la configuraci{\'o}n de dispositivos y la gesti{\'o}n de procesos distribuidos.},
		type=main,
		symbol={🖥️},
		user1={Interfaz de usuario},
		user2={Capa de presentaci{\'o}n},
		user3={Cliente de aplicaci{\'o}n},
		see={backend,aplicacion,aplicacionweb}
	}
	
	% ===================== FTP =====================
	\newglossaryentry{ftp}
	{
		name={FTP},
		text={FTP},
		sort={ftp},
		first={Protocolo de Transferencia de Archivos (\textit{File Transfer Protocol}, FTP)},
		firstplural={FTPs},
		plural={FTPs},
		description={\emph{FTP} es un protocolo de comunicaci{\'o}n en arquitectura \gls{cliente}\textendash\gls{servidor} para transferir \glspl{dato} entre sistemas conectados a una \gls{red}. Requiere autenticaci{\'o}n y opera t{\'i}picamente sobre los \glspl{puerto} 20 y 21. En \gls{iot} puede emplearse para despliegue de firmware o recolecci{\'o}n de registros cuando se demanda compatibilidad heredada.},
		type=\glsdefaulttype,
		symbol={📤},
		user1={Transferencia de archivos},
		user2={Protocolo legado},
		user3={Cliente\textendash servidor},
		see={http,https}
	}
	
	\newglossaryentry{fuentedatos}{
		type=main,
		name={Fuente de datos},
		text={fuente de datos},
		sort={fuentedatos},
		plural={fuentes de datos},
		first={fuente de datos},
		firstplural={fuentes de datos},
		symbol={🛰️},
		see={dataset,sensor,basedatos},
		user1={Origen de informaci{\'o}n empleado para procesos computacionales},
		user2={Clasificaci{\'o}n: primaria o secundaria},
		user3={Naturaleza: estructurada o no estructurada},
		user4={Usada para an{\'a}lisis, monitorizaci{\'o}n y toma de decisiones},
		description={Entidad, mecanismo u origen que proporciona informaci{\'o}n relevante para un sistema computacional, anal{\'i}tico o inteligente. Una fuente de datos puede corresponder a sensores f{\'i}sicos, bases de datos locales o distribuidas, repositorios documentales, servicios web, flujos de \glspl{dato} en \gls{tiemporeal}, plataformas \gls{iot}, sistemas transaccionales u otros medios capaces de generar, almacenar o transmitir informaci{\'o}n para tareas de an{\'a}lisis, modelado, control, automatizaci{\'o}n y \gls{tomadecisiones}}
	}
	
	\newglossaryentry{fuerza}
	{
		name={Fuerza},
		text={fuerza},
		sort={fuerza},
		description={Magnitud vectorial que mide la interacci{\'o}n capaz de alterar el estado de movimiento o deformaci{\'o}n de un cuerpo. 
			Se expresa como el producto de la masa por la aceleraci{\'o}n (\ensuremath{F = m \cdot a}) y se mide en newtons (N). 
			En el contexto de los \glspl{sibiot}, la fuerza se monitoriza mediante \glspl{sensor} resistivos, capacitivos o piezoel{\'e}ctricos para evaluar el impacto, la carga o la tensi{\'o}n en componentes mec{\'a}nicos, permitiendo la detecci{\'o}n temprana de fallos estructurales y el control adaptativo de maquinaria. 
			En rob{\'o}tica, la medici{\'o}n de fuerza es fundamental para el control de agarre y la seguridad de la interacci{\'o}n humano-m{\'a}quina.},
		plural={fuerzas},
		symbol={F},
		first={Fuerza (F)},
		see=[Relacionado con:]{peso,masa,esfuerzo,tension}
	}
	
	\newglossaryentry{ganancia}{
		name={Ganancia},
		text={ganancia},
		sort={ganancia antena},
		plural={ganancias},
		first={ganancia de antena},
		firstplural={ganancias de antena},
		description={Magnitud que expresa la capacidad de una antena para concentrar y dirigir la energ{\'i}a electromagn{\'e}tica radiada en una direcci{\'o}n determinada, en comparaci{\'o}n con una antena de referencia ideal, usualmente una antena isot{\'o}pica. La ganancia no implica un aumento real de potencia, sino una redistribuci{\'o}n m{\'a}s eficiente de la energ{\'i}a disponible, y se expresa habitualmente en decibelios referidos a una antena isot{\'o}pica (dBi). En sistemas basados en \gls{rfid}, la ganancia influye directamente en el alcance de lectura y en el {\'a}rea de cobertura de la comunicaci{\'o}n, constituyendo un par{\'a}metro clave en el dise{\~n}o e implementaci{\'o}n de antenas.}
	}
	
	
	% ===================== Gateway =====================
	\newglossaryentry{gateway}
	{
		name={Gateway},
		text={gateway},
		first={pasarela o gateway},
		firstplural={gateways},
		description={Un \emph{gateway} es un dispositivo o componente intermedio que act{\'u}a como puente entre \glspl{red} o sistemas con protocolos diferentes. En \gls{iot} interconecta \glspl{sensor}, \glspl{actuador} y \glspl{microcontrolador} con redes de mayor escala, gestionando agregaci{\'o}n de \glspl{dato}, conversi{\'o}n de protocolos, filtrado y funciones de \gls{seguridad}.},
		type=\glsdefaulttype,
		symbol={🔗},
		user1={Puerta de enlace},
		user2={Traductor de protocolos},
		user3={Concentrador IoT},
		see={conectividad,infraestructurared}
	}
	
	\newglossaryentry{gcp}{
		name={Google Cloud Platform},
		text={GCP},
		sort={google cloud platform},
		first={Plataforma de servicios en la nube de Google (\textit{Google Cloud Platform}, GCP)},
		firstplural={GCPs},
		plural={GCPs},
		description={\emph{Google Cloud Platform (GCP)} es una plataforma de \gls{servicionube} que proporciona servicios de infraestructura (IaaS), plataforma (PaaS) y software (SaaS). \emph{GCP} incluye soluciones de \gls{almacenamiento}, \gls{procesamientodatos}, \gls{aprendizajeautomatico}, \gls{analisisdatos} y despliegue en \gls{tiemporeal}, lo que lo convierte en un ecosistema fundamental para proyectos \gls{iot}, \glspl{modelopredictivo} y arquitecturas \gls{escalable}.},
		type=main,
		symbol={☁️},
		user1={Plataforma en la nube de Google},
		user2={Servicios IaaS, PaaS y SaaS},
		user3={Soporte para IoT, Big Data y aprendizaje autom{\'a}tico},
		see={aws,azure,servicionube,iot}
	}
	
	\newglossaryentry{gdpr}{
		name={GDPR},
		text={GDPR},
		sort={gdpr},
		first={reglamento general de protecci{\'o}n de datos (\textit{General Data Protection Regulation}, GDPR)},
		firstplural={GDPRs},
		plural={GDPRs},
		description={El \emph{GDPR} es la normativa de la Uni{\'o}n Europea, en vigor desde 2018, que regula la protecci{\'o}n de datos personales y la \gls{privacidad} de los individuos en el entorno digital. Establece principios como el consentimiento expl{\'i}cito, el derecho al acceso y rectificaci{\'o}n, la portabilidad de los datos, el derecho al olvido y la obligaci{\'o}n de notificaci{\'o}n en caso de violaciones de seguridad. En el contexto de los \glspl{sibiot}, el GDPR resulta fundamental para garantizar que los \glspl{dispositivointeligente} y plataformas en la \gls{nube} gestionen informaci{\'o}n personal de forma segura, transparente y responsable, asegurando \gls{integridad}, \gls{confidencialidad} y control por parte del usuario.},
		type=main,
		symbol={⚖️},
		user1={Reglamento europeo de protecci{\'o}n de datos},
		user2={Derecho a la privacidad y al olvido},
		user3={Aplicaci{\'o}n en IoT y servicios en la nube},
		see={privacidad,seguridad,integridad,nube,sibiot}
	}
	
	\newglossaryentry{gemelodigital}{
		name={Gemelo digital},
		text={gemelo digital},
		sort={gemelo digital},
		plural={gemelos digitales},
		first={gemelo digital},
		firstplural={gemelos digitales},
		description={Un \emph{gemelo digital} es una representaci{\'o}n virtual de un objeto, proceso o sistema del \gls{mundofisico}, actualizada continuamente con datos en tiempo real provenientes de \glspl{sensor}, \glspl{dispositivointeligente} y \glspl{sibiot}. Permite monitorear, simular y predecir el comportamiento del sistema f{\'i}sico, optimizando el mantenimiento, la eficiencia operativa y la toma de decisiones. En el {\'a}mbito de \gls{iot}, los gemelos digitales conectan el \gls{mundofisico}, el \gls{mundodigital} y el \gls{mundovirtual} para habilitar soluciones inteligentes en manufactura, salud, agricultura y \glspl{ciudadinteligente}.},
		type=main,
		symbol={🧩},
		user1={Representaci{\'o}n virtual del mundo f{\'i}sico},
		user2={Actualizaci{\'o}n en tiempo real con datos IoT},
		user3={Base para predicci{\'o}n y optimizaci{\'o}n de sistemas},
		see={mundofisico,mundodigital,mundovirtual,realidadaumentada,realidadvirtual,sibiot}
	}
	
	% ===================== Generaci{\'o}n de valor =====================
	\newglossaryentry{generacionvalor}
	{
		name={Generaci{\'o}n de valor},
		text={generaci{\'o}n de valor},
		sort={generacion de valor},
		plural={procesos de generaci{\'o}n de valor},
		first={generaci{\'o}n de valor},
		firstplural={procesos de generaci{\'o}n de valor},
		description={La \emph{generaci{\'o}n de valor} es el proceso por el cual una organizaci{\'o}n, \gls{sistema} o tecnolog{\'i}a produce beneficios tangibles o intangibles (eficiencia, innovaci{\'o}n, satisfacci{\'o}n del cliente, rentabilidad, logro estrat{\'e}gico). En \gls{iot} se asocia a soluciones que impactan positivamente procesos y resultados del negocio.},
		type=\glsdefaulttype,
		symbol={💡},
		user1={Creaci{\'o}n de valor},
		user2={Impacto de negocio},
		user3={Propuesta de valor},
		see={gestionempresarial}
	}
	
	\newglossaryentry{gestion}{
		name={Gesti{\'o}n},
		text={gesti{\'o}n},
		sort={gestion},
		plural={gestiones},
		first={gesti{\'o}n},
		firstplural={gestiones},
		description={La \emph{gesti{\'o}n} es el conjunto de acciones, t{\'e}cnicas y estrategias orientadas a planificar, organizar, coordinar y controlar los recursos y actividades de un sistema, organizaci{\'o}n o proyecto, con el fin de alcanzar objetivos definidos de manera eficiente y sostenible. En entornos \gls{iot}, la gesti{\'o}n abarca desde la administraci{\'o}n de redes, \glspl{dispositivoconectado} y \glspl{servicionube}, hasta la \gls{seguridad}, el \gls{rendimiento} y la \gls{interoperabilidad}, garantizando la continuidad y el valor estrat{\'e}gico de la soluci{\'o}n.},
		type=main,
		symbol={📊},
		user1={Planificaci{\'o}n y organizaci{\'o}n de recursos},
		user2={Control y coordinaci{\'o}n de actividades},
		user3={Base de la administraci{\'o}n en IoT y software},
		see={proceso,gestionempresarial,infraestructura,sibiot}
	}
	
	\newglossaryentry{gestionconfianza}{
		name={Gesti{\'o}n de confianza},
		text={gesti{\'o}n de confianza},
		sort={gestion confianza},
		plural={m{\'e}todos de gesti{\'o}n de confianza},
		first={gesti{\'o}n de confianza},
		firstplural={m{\'e}todos de gesti{\'o}n de confianza},
		description={T{\'e}cnicas y pol{\'i}ticas para establecer y mantener relaciones seguras entre \glspl{dispositivo}, servicios y usuarios (autenticaci{\'o}n, reputaci{\'o}n, control de acceso).},
		type=main,
		symbol={🛡️},
		see={privacidad,seguridad,provisionamiento},
		user1={Autenticaci{\'o}n},
		user2={Reputaci{\'o}n},
		user3={Autorizaci{\'o}n}
	}
	
	\newglossaryentry{gestiondatos}{
		name={Gesti{\'o}n de datos},
		text={gesti{\'o}n de datos},
		sort={gestion de datos},
		first={gesti{\'o}n de datos},
		description={Pr{\'a}cticas y procesos para recolectar, almacenar, asegurar, gobernar y explotar datos, incluyendo calidad, integridad, seguridad y disponibilidad a lo largo del ciclo de vida.},
		type=main,
		symbol={🗂️},
		user1={Gobernanza},
		user2={Calidad e integridad},
		user3={Ciclo de vida},
		see={basedatos,sgbd,recuperacion}
	}
	
	\newglossaryentry{gestioneficienterecursos}{
		name={Gesti{\'o}n eficiente de recursos},
		text={gesti{\'o}n eficiente de recursos},
		sort={gestion eficiente de recursos},
		type={main},
		plural={gestiones eficientes de recursos},
		description={Conjunto de estrategias, t{\'e}cnicas y procedimientos destinados a optimizar el uso de recursos limitados, tales como energ{\'i}a, memoria, ancho de banda, capacidad de procesamiento o almacenamiento. En los sistemas basados en IoT, la gesti{\'o}n eficiente de recursos permite maximizar el rendimiento, prolongar la vida {\'u}til de los dispositivos, reducir el consumo energ{\'e}tico y asegurar la continuidad operativa incluso en entornos con restricciones significativas},
		first={gesti{\'o}n eficiente de recursos},
		firstplural={gestiones eficientes de recursos},
		symbol={📊},
		user1={Optimiza el uso de energ{\'i}a, memoria y procesamiento},
		user2={Esencial para dispositivos IoT con recursos limitados},
		user3={Contribuye a la continuidad operativa y resiliencia del sistema},
		user4={Incluye t{\'e}cnicas de programaci{\'o}n, priorizaci{\'o}n y administraci{\'o}n din{\'a}mica}
	}
	
	% ===================== Gesti{\'o}n empresarial =====================
	\newglossaryentry{gestionempresarial}
	{
		name={Gesti{\'o}n empresarial},
		text={gesti{\'o}n empresarial},
		sort={gestion empresarial},
		plural={gestiones empresariales},
		first={gesti{\'o}n empresarial},
		firstplural={gestiones empresariales},
		description={La \emph{gesti{\'o}n empresarial} comprende procesos, estrategias, t{\'e}cnicas y recursos para planificar, organizar, dirigir y controlar actividades de una organizaci{\'o}n con eficiencia y sostenibilidad. Integra \glspl{dato} y \gls{tomadecisiones} apoyadas en \gls{bigdata}, \gls{analisisdatos} y \gls{nube}.},
		type=\glsdefaulttype,
		symbol={📈},
		user1={Management},
		user2={Administraci{\'o}n de empresas},
		user3={Direcci{\'o}n estrat{\'e}gica},
		see={erp}
	}
	
	\newglossaryentry{gestionenergetica}{
		name={Gesti{\'o}n energ{\'e}tica},
		text={gesti{\'o}n energ{\'e}tica},
		sort={gestion energetica},
		plural={estrategias de gesti{\'o}n energ{\'e}tica},
		first={gesti{\'o}n energ{\'e}tica},
		firstplural={estrategias de gesti{\'o}n energ{\'e}tica},
		description={Medici{\'o}n, an{\'a}lisis y optimizaci{\'o}n de consumos y generaci{\'o}n (p.~ej., \gls{pv}, HVAC, bater{\'i}as) manteniendo confort y coste objetivo.},
		type=main,
		symbol={📊},
		see={hems,almacenamientoenergetico,respuestademanda},
		user1={Monitorizaci{\'o}n},
		user2={Optimizaci{\'o}n},
		user3={Ahorro}
	}
	
	\newglossaryentry{gestioninformacion}{
		name={gesti{\'o}n de la informaci{\'o}n},
		text={gesti{\'o}n de la informaci{\'o}n},
		sort={gestioninformacion},
		plural={gestiones de la informaci{\'o}n},
		first={gesti{\'o}n de la informaci{\'o}n},
		firstplural={gestiones de la informaci{\'o}n},
		description={La gesti{\'o}n de la informaci{\'o}n constituye una funci{\'o}n esencial en los \glspl{sibiot}, ya que permite transformar los datos capturados por \glspl{sensor} y \glspl{dispositivo} en conocimiento {\'u}til. Involucra tareas de \gls{almacenamiento}, clasificaci{\'o}n, acceso, seguridad, y mantenimiento de la \gls{calidaddatos}. Su prop{\'o}sito es garantizar que la informaci{\'o}n est{\'e} disponible, actualizada y protegida, de modo que respalde la \gls{tomadecisionesinformadas} y la \gls{eficiencia} organizacional.},
		type=\glsdefaulttype,
		symbol={🗂️},
		user1={Information management},
		user2={Administraci{\'o}n y control integral de los datos en un sistema},
		user3={Proceso que garantiza el acceso, la calidad y la seguridad de la informaci{\'o}n},
		see={calidaddatos,almacenamiento,seguridad,procesamiento,tomadecisionesinformadas}
	}
	
	\newglossaryentry{gestionrecursos}{
		name={Gesti{\'o}n de recursos},
		text={gesti{\'o}n de recursos},
		sort={gestion de recursos},
		type={main},
		plural={gestiones de recursos},
		description={Proceso mediante el cual se planifican, administran y optimizan los recursos disponibles en un sistema, tales como energ{\'i}a, memoria, ancho de banda, almacenamiento y capacidad de procesamiento. En los sistemas basados en IoT, la gesti{\'o}n de recursos es fundamental para garantizar el funcionamiento eficiente de los dispositivos, asegurar la continuidad operativa y mantener niveles adecuados de rendimiento, especialmente en entornos con limitaciones significativas},
		first={gesti{\'o}n de recursos},
		firstplural={gestiones de recursos},
		symbol={🗂️},
		user1={Implica planificaci{\'o}n y administraci{\'o}n de recursos limitados},
		user2={Clave para mantener la eficiencia y continuidad en sistemas IoT},
		user3={Abarca energ{\'i}a, memoria, procesamiento y ancho de banda},
		user4={Puede incluir estrategias autom{\'a}ticas o basadas en algoritmos de optimizaci{\'o}n},
		see={gestion,recurso,optimizacion,algoritmo}
	}
	
	\newglossaryentry{gestionservicio}{
		name={Gesti{\'o}n de servicios},
		text={gesti{\'o}n de servicios},
		sort={gestionservicio},
		plural={gestiones de servicios},
		first={gesti{\'o}n de servicios},
		firstplural={gestiones de servicios},
		description={Conjunto de procesos y pr{\'a}cticas orientadas a planificar, operar, supervisar y mejorar servicios de tecnolog{\'\i}a de la informaci{\'o}n a lo largo de su ciclo de vida. En los \glspl{sibiot}, la gesti{\'o}n de servicios abarca la administraci{\'o}n de plataformas, aplicaciones, flujos de datos y recursos distribuidos, garantizando niveles adecuados de disponibilidad, rendimiento, seguridad y calidad del servicio ofrecido a los usuarios.},
		type=main,
		symbol={🛠️},
		user1={Administraci{\'o}n de servicios},
		user2={Operaci{\'o}n y soporte},
		user3={Ciclo de vida del servicio},
		see={calidadservicio,orquestacion,arquitecturadistribuida}
	}
	
	\newglossaryentry{get}{
		name={GET},
		text={GET},
		sort={get},
		first={GET},
		description={M{\'e}todo del protocolo HTTP utilizado para solicitar la representaci{\'o}n de un recurso identificado por una URL, sin producir efectos secundarios sobre el estado del sistema. En arquitecturas RESTful, \texttt{GET} se emplea para la consulta de datos, lectura de estados y recuperaci{\'o}n de informaci{\'o}n, siendo un m{\'e}todo seguro e idempotente. En sistemas de informaci{\'o}n basados en IoT, \texttt{GET} es com{\'u}nmente utilizado para acceder a mediciones, estados de dispositivos y resultados de procesamiento almacenados en servicios backend},
		type=main,
		see={http,rest,restful,clienteservidor,api}
	}
	
	% ===================== GNSS =====================
	\newglossaryentry{gnss}
	{
		name={GNSS},
		text={GNSS},
		sort={gnss},
		first={sistema global de navegaci{\'o}n por sat{\'e}lite (\textit{Global Navigation Satellite System}, GNSS)},
		firstplural={GNSSs},
		plural={GNSSs},
		description={\emph{GNSS} es el conjunto de sistemas satelitales de navegaci{\'o}n que proveen posici{\'o}n, velocidad y tiempo globales (p.\,ej., \gls{gps}, GLONASS, Galileo, BeiDou). En \glspl{sibiot} soporta \glspl{sistemanavegacion} para movilidad, log{\'i}stica y agricultura de \gls{precision}, con alta \gls{disponibilidad} y \gls{fiabilidad} en \gls{tiemporeal}.},
		type=\glsdefaulttype,
		symbol={🛰️},
		user1={Constelaciones de navegaci{\'o}n},
		user2={Posicionamiento global},
		user3={Tiempo de referencia},
		see={gps}
	}
	
	\newglossaryentry{gobernanza}{
		name={Gobernanza},
		text={gobernanza},
		sort={gobernanza},
		plural={gobernanzas},
		first={gobernanza},
		firstplural={gobernanzas},
		description={Conjunto de principios, roles, normas y procesos que regulan la toma de decisiones, la rendici{\'o}n de cuentas y el uso responsable de recursos, datos y tecnolog{\'i}as dentro de una organizaci{\'o}n o sistema socio-t{\'e}cnico, garantizando alineaci{\'o}n estrat{\'e}gica, cumplimiento normativo y mejora continua},
		type=main,
		symbol={⚖},
		see={gestion,politicaorganizacional,control}
	}
	
	\newglossaryentry{googlecloudspanner}{
		name={Google Cloud Spanner},
		text={Google Cloud Spanner},
		sort={google cloud spanner},
		first={Google Cloud Spanner},
		description={Sistema de gesti{\'o}n de bases de datos distribuido, relacional y globalmente consistente, desarrollado por Google y ofrecido como servicio en la nube. Google Cloud Spanner combina escalabilidad horizontal propia de sistemas NoSQL con sem{\'a}nticas relacionales y soporte completo de transacciones ACID, proporcionando consistencia fuerte a escala global mediante relojes distribuidos y el mecanismo \emph{TrueTime}. En arquitecturas \glspl{sibiot} y sistemas empresariales distribuidos, Cloud Spanner resulta adecuado para datos transaccionales cr{\'i}ticos que requieren alta disponibilidad, baja latencia y coherencia estricta entre regiones geogr{\'a}ficas},
		type=main,
		symbol={🌐},
		see={cloudspanner,spanner,oltp,consistenciafuerte,arquitecturadistribuida,cloudcomputing}
	}
	
	% ===================== GPIO =====================
	\newglossaryentry{gpio}
	{
		name={GPIO},
		text={GPIO},
		sort={gpio},
		first={entrada/salida de prop{\'o}sito general (\textit{General Purpose Input/Output}, GPIO)},
		firstplural={GPIOs},
		plural={GPIOs},
		description={\emph{GPIO} refiere a \glspl{puerto} f{\'i}sicos en \glspl{microcontrolador} o microcomputadoras configurables por software como entradas o salidas digitales para interactuar con \glspl{sensor}, \glspl{actuador} y otros componentes electr{\'o}nicos.},
		type=\glsdefaulttype,
		symbol={📍},
		user1={Pin digital},
		user2={I/O de prop{\'o}sito general},
		user3={L{\'i}nea de E/S},
		see={placadesarrollo}
	}
	
	% ===================== GPS =====================
	\newglossaryentry{gps}
	{
		name={GPS},
		text={GPS},
		first={sistema de posicionamiento global (\textit{Global Positioning System}, GPS)},
		firstplural={GPSs},
		sort={gps},
		plural={GPSs},
		description={\emph{GPS} es un sistema satelital que permite determinar ubicaci{\'o}n geogr{\'a}fica y tiempo con alta \gls{precision}. En \gls{iot} se usa en movilidad, log{\'i}stica, agricultura de precisi{\'o}n, veh{\'i}culos aut{\'o}nomos y rastreo, integr{\'a}ndose con otros \glspl{gnss} para mejorar \gls{fiabilidad} y \gls{disponibilidad}.},
		type=\glsdefaulttype,
		symbol={📡},
		user1={Posicionamiento satelital},
		user2={Trilateraci{\'o}n},
		user3={Navegaci{\'o}n GNSS},
		see={gnss}
	}
	
	% ===================== GPU =====================
	\newglossaryentry{gpu}
	{
		name={GPU},
		text={GPU},
		sort={gpu},
		first={unidad de procesamiento gr{\'a}fico (\textit{Graphics Processing Unit}, GPU)},
		firstplural={GPUs},
		plural={GPUs},
		description={La \emph{GPU} es un procesador especializado en c{\'a}lculo gr{\'a}fico y paralelo (renderizado, video, c{\'a}lculo matricial). Adem{\'a}s de gr{\'a}ficos, se emplea en \gls{aprendizajeautomatico}, simulaci{\'o}n y \gls{procesamientodatos} intensivo para \glspl{sibiot}.},
		type=\glsdefaulttype,
		symbol={🧮},
		user1={C{\'o}mputo paralelo},
		user2={Acelerador gr{\'a}fico},
		user3={Procesamiento matricial},
		see={cpu}
	}
	
	% ===================== GPRS =====================
	\newglossaryentry{gprs}
	{
		name={GPRS},
		text={GPRS},
		sort={gprs},
		first={servicio general de radio por paquetes (\textit{General Packet Radio Service}, GPRS)},
		firstplural={tecnolog{\'i}as GPRS},
		plural={tecnolog{\'i}as GPRS},
		description={\emph{GPRS} es una extensi{\'o}n de \gls{gsm} que introduce conmutaci{\'o}n por paquetes para transmisi{\'o}n de \glspl{dato} m{\'a}s eficiente. Considerada 2.5G, habilit{\'o} acceso a Internet y aplicaciones m{\'o}viles de baja velocidad en \gls{iot}.},
		type=\glsdefaulttype,
		symbol={📶},
		user1={Conmutaci{\'o}n por paquetes},
		user2={Datos m{\'o}viles 2.5G},
		user3={Red celular heredada},
		see={gsm,lte}
	}
	
	% ===================== GSM =====================
	\newglossaryentry{gsm}
	{
		name={GSM},
		text={GSM},
		sort={gsm},
		first={sistema global de comunicaci{\'o}n m{\'o}vil (\textit{Global System for Mobile Communications}, GSM)},
		firstplural={tecnolog{\'i}as GSM},
		plural={tecnolog{\'i}as GSM},
		description={\emph{GSM} es un est{\'a}ndar de segunda generaci{\'o}n (2G) para comunicaci{\'o}n m{\'o}vil digital. Provee voz, SMS y datos b{\'a}sicos por conmutaci{\'o}n de circuitos; se usa en \gls{iot} para conectividad de baja velocidad y amplia cobertura.},
		type=\glsdefaulttype,
		symbol={📡},
		user1={2G},
		user2={Telefon{\'i}a m{\'o}vil digital},
		user3={Red celular},
		see={gprs,lte}
	}
	
	\newglossaryentry{iso}{
		name={ISO},
		text={ISO},
		sort={iso},
		first={Organizaci{\'o}n Internacional de Normalizaci{\'o}n (\textit{International Organization for Standardization}, ISO)},
		firstplural={ISO},
		plural={ISO},
		description={\emph{ISO} es una organizaci{\'o}n internacional independiente y no gubernamental que desarrolla y publica normas t{\'e}cnicas, industriales y comerciales reconocidas globalmente. Fundada en 1947, est{\'a} compuesta por organismos nacionales de estandarizaci{\'o}n de m{\'a}s de 160 pa{\'i}ses. En el {\'a}mbito de la inform{\'a}tica y las telecomunicaciones, ISO ha definido est{\'a}ndares fundamentales como el \gls{osi}, as{\'i} como normas para \gls{seguridad}, \gls{calidad} de software, \gls{interoperabilidad} y gesti{\'o}n de \glspl{sistemainformatico}. En el contexto de \glspl{sibiot}, sus est{\'a}ndares garantizan la compatibilidad, la eficiencia y la confianza en soluciones distribuidas.},
		type=main,
		symbol={🌐},
		user1={Organizaci{\'o}n de estandarizaci{\'o}n internacional},
		user2={Normas en tecnolog{\'i}a, seguridad y calidad},
		user3={Relevancia en IoT y sistemas distribuidos},
		see={osi,seguridad,calidad,interoperabilidad,sibiot}
	}
	
	\newglossaryentry{isr}{
		name={Rutina de servicio de interrupci{\'o}n (ISR)},
		text={rutina de servicio de interrupci{\'o}n},
		sort={rutina servicio interrupcion},
		first={rutina de servicio de interrupci{\'o}n (\textit{Interrupt Service Routine}, ISR)},
		description={Fragmento de c{\'o}digo que se ejecuta autom{\'a}ticamente en respuesta a la ocurrencia de una interrupci{\'o}n generada por hardware o software en un microcontrolador. Una \emph{rutina de servicio de interrupci{\'o}n} detiene temporalmente la ejecuci{\'o}n normal del programa, atiende el evento que origin{\'o} la interrupci{\'o}n y, una vez finalizada, devuelve el control al flujo principal. En sistemas embebidos y SIbIoT, las ISR permiten responder de forma inmediata a eventos cr{\'i}ticos, como cambios de estado, se{\~n}ales externas o temporizadores, garantizando un comportamiento reactivo y determinista.},
		type=main,
		symbol={⚡},
		user1={Atenci{\'o}n a eventos},
		user2={Ejecuci{\'o}n inmediata},
		user3={Manejo de interrupciones},
		see={interrupcion,arduinouno,temporizador}
	}
	
	\newglossaryentry{iteracion}{
		name={Iteraci{\'o}n},
		text={iteraci{\'o}n},
		sort={iteracion},
		plural={iteraciones},
		first={iteraci{\'o}n},
		firstplural={iteraciones},
		description={Ciclo de trabajo delimitado en el que se planifica, construye, verifica e incrementa un producto o componente, incorporando retroalimentaci{\'o}n para mejorar continuamente.},
		type=main,
		symbol={🔁},
		user1={Ciclo},
		user2={Incremento},
		user3={Retroalimentaci{\'o}n},
		see={agile,refactorizacion}
	}
	
	\newglossaryentry{iterativa}{
		name={Iterativa},
		text={iterativa},
		sort={iterativa},
		plural={iterativas},
		first={iterativa},
		firstplural={iterativas},
		description={Caracter{\'\i}stica de un proceso que avanza mediante ciclos repetidos de planificaci{\'o}n, construcci{\'o}n, evaluaci{\'o}n y ajuste, con el objetivo de refinar progresivamente requisitos, dise{\~n}o e implementaci{\'o}n. La iteraci{\'o}n permite aprendizaje temprano, incorporaci{\'o}n de retroalimentaci{\'o}n y detecci{\'o}n de defectos antes de que se consoliden. En IoT/SIbIoT, el car{\'a}cter iterativo es cr{\'\i}tico por la incertidumbre de condiciones de campo, heterogeneidad de dispositivos y variabilidad de redes, lo que exige ajustes continuos en configuraciones, modelos y pruebas.},
		type=main,
		symbol={🔄},
		user1={Refinamiento por ciclos},
		user2={Retroalimentaci{\'o}n temprana},
		user3={Aprendizaje y ajuste continuo},
		see={incremental,ciclodesarrollo,scrum,kanban,pruebaautomatizada}
	}
	
	\newglossaryentry{hadoop}{
		name={Hadoop},
		text={Hadoop},
		sort={hadoop},
		plural={ecosistemas Hadoop},
		first={Hadoop},
		firstplural={ecosistemas Hadoop},
		description={Plataforma de c{\'o}digo abierto desarrollada por la Apache Software Foundation para el almacenamiento y procesamiento distribuido de grandes vol{\'u}menes de datos. \emph{Hadoop} se compone principalmente del sistema de archivos distribuido HDFS (\textit{Hadoop Distributed File System}) y del modelo de programaci{\'o}n MapReduce, que permite ejecutar tareas de forma paralela en cl{\'u}steres de servidores. Su ecosistema incluye herramientas como Hive, Pig, HBase y Spark, que ampl{\'i}an sus capacidades de an{\'a}lisis y gesti{\'o}n de datos. En el contexto de los \glspl{sibiot}, Hadoop es fundamental para procesar informaci{\'o}n masiva proveniente de sensores y dispositivos, facilitando el an{\'a}lisis avanzado, la toma de decisiones en tiempo real y la implementaci{\'o}n de modelos predictivos.},
		type=main,
		symbol={🐘},
		user1={Marco de Big Data},
		user2={Procesamiento distribuido},
		user3={Apache Hadoop},
		see={bigdata,procesamientodatos,etl}
	}
	
	\newglossaryentry{halow}{
		name={Wi-Fi HaLow},
		text={Wi-Fi HaLow},
		sort={halow},
		first={Wi-Fi HaLow},
		description={Denominaci{\'o}n comercial de IEEE 802.11ah, orientada a IoT, que opera en bandas sub-1\,GHz para ampliar alcance y mejorar penetraci{\'o}n, con consumo energ{\'e}tico reducido respecto a Wi-Fi tradicional.},
		type=main,
		symbol={📡},
		user1={Sub-1\,GHz},
		user2={Largo alcance},
		user3={IoT},
		see={redinalambrica,anchobanda}
	}
	
	% ===================== Hardware =====================
	\newglossaryentry{hardware}
	{
		name={Hardware},
		text={hardware},
		sort={hardware},
		plural={componentes de hardware},
		first={hardware},
		firstplural={componentes de hardware},
		description={El \emph{hardware} comprende componentes el{\'e}ctricos, electr{\'o}nicos, electromec{\'a}nicos y mec{\'a}nicos: desde computadoras y \glspl{placadesarrollo}, hasta cables, memorias, discos, pantallas y perif{\'e}ricos necesarios para operar un \gls{sistemainformatico}.},
		type=\glsdefaulttype,
		symbol={🧰},
		user1={Soporte f{\'i}sico},
		user2={Plataforma},
		user3={Infraestructura f{\'i}sica},
		see={software,infraestructura}
	}
	
	\newglossaryentry{hbase}{
		name={Apache HBase},
		text={Apache HBase},
		sort={hbase},
		plural={Apache HBase},
		first={Apache HBase},
		firstplural={Apache HBase},
		description={Sistema de gesti{\'o}n de \glspl{basedatosnosql} distribuido, orientado a columnas, desarrollado como parte del ecosistema Apache Hadoop. \emph{HBase} est{\'a} inspirado en el modelo de Bigtable de Google y est{\'a} dise{\~n}ado para manejar grandes vol{\'u}menes de datos dispersos, proporcionando escalabilidad horizontal y acceso en \gls{tiemporeal} a datos almacenados en \gls{hdfs}. A diferencia de las \glspl{basedatosrelacional}, no utiliza \gls{sql} tradicional, sino un \gls{api} propia y soporte para consultas mediante Apache Phoenix. En el contexto de los \glspl{sibiot}, HBase es adecuado para almacenar y analizar datos masivos de \glspl{sensor}, integrar flujos de \gls{bigdata}, habilitar \glspl{modelopredictivo} y soportar aplicaciones de \gls{mantenimientopredictivo} y \glspl{ciudadinteligente}.},
		type=main,
		symbol={📊},
		user1={Base de datos NoSQL orientada a columnas},
		user2={Escalabilidad masiva sobre Hadoop},
		user3={Soporte a Big Data e IoT en tiempo real},
		see={basedatosnosql,bigdata,hdfs,modelopredictivo,mantenimientopredictivo,sibiot}
	}
	
	\newglossaryentry{hdd}{
		name={HDD},
		text={HDD},
		sort={hdd},
		first={HDD},
		description={\emph{Hard Disk Drive}. Dispositivo de almacenamiento magn{\'e}tico no vol{\'a}til para persistencia de grandes vol{\'u}menes de datos en sistemas computacionales.},
		type=main,
		symbol={💽},
		user1={Almacenamiento magn{\'e}tico},
		user2={Persistencia},
		user3={Capacidad},
		see={basedatos,recuperacion}
	}
	
	\newglossaryentry{hdfs}{
		name={HDFS},
		text={HDFS},
		sort={hdfs},
		first={sistema de archivos distribuido de Hadoop (\textit{Hadoop Distributed File System}, HDFS)},
		firstplural={HDFSs},
		plural={HDFSs},
		description={\emph{HDFS} es un sistema de archivos distribuido dise{\~n}ado para almacenar grandes vol{\'u}menes de datos en cl{\'u}steres.},
		type=\glsdefaulttype,
		symbol={🗄️},
		user1={Almacenamiento distribuido},
		user2={Bloques replicados},
		user3={Soporte Big Data},
		see={bigdata,hbase,etl,sibiot}
	}
	
	% ===================== HDMI =====================
	\newglossaryentry{hdmi}
	{
		name={HDMI},
		text={HDMI},
		sort={hdmi},
		first={interfaz multimedia de alta definici{\'o}n (\textit{High\textendash Definition Multimedia Interface}, HDMI)},
		firstplural={HDMI},
		plural={HDMIs},
		description={\emph{HDMI} es una interfaz digital para transmisi{\'o}n de video y audio de alta definici{\'o}n por un solo cable. Com{\'u}n en monitores, televisores y \glspl{sistemaembebido} usados como \glspl{dashboard} o estaciones de \gls{monitorizacion} en \gls{iot}.},
		type=\glsdefaulttype,
		symbol={🔌},
		user1={Video digital},
		user2={Audio HD},
		user3={Interfaz A/V},
		see={browser,interfaz}
	}
	
	\newglossaryentry{hems}{
		name={Sistema de gesti{\'o}n energ{\'e}tica del hogar},
		text={HEMS},
		sort={hems},
		plural={HEMSs},
		first={sistema de gesti{\'o}n energ{\'e}tica del hogar (\textit{Home Energy Management System}, HEMS)},
		firstplural={HEMSs},
		description={Hardware/software que supervisa, analiza y controla el consumo energ{\'e}tico en \glspl{hogarinteligente}, optimizando eficiencia y coste.},
		type=main,
		symbol={🧠},
		see={pv,almacenamientoenergetico,gestionenergetica,respuestademanda},
		user1={Medici{\'o}n fina},
		user2={Control de cargas},
		user3={Optimizaci{\'o}n multiobjetivo}
	}
	
	% ===================== Heterogeneidad =====================
	\newglossaryentry{heterogeneidad}
	{
		name={Heterogeneidad},
		text={heterogeneidad},
		sort={heterogeneidad},
		plural={heterogeneidades},
		first={heterogeneidad},
		firstplural={heterogeneidades},
		description={La \emph{heterogeneidad} caracteriza sistemas con diversidad de \gls{hardware}, software, protocolos, plataformas y formatos de \gls{dato}. En \gls{iot} los \glspl{dispositivoconectado} difieren en capacidades, S.O. e interfaces, lo que exige \gls{interoperabilidad} y estandarizaci{\'o}n.},
		type=\glsdefaulttype,
		symbol={🧩},
		user1={Diversidad tecnol{\'o}gica},
		user2={Variedad de protocolos},
		user3={Entorno mixto},
		see={interoperabilidad}
	}
	
	\newglossaryentry{heterogeneo}{
		name={heterog{\'e}neo},
		text={heterog{\'e}neo},
		first={heterog{\'e}neo},
		sort={heterogeneo},
		description={Propiedad de un sistema, entorno o conjunto de componentes caracterizada por la coexistencia de elementos con diferentes naturalezas, tecnolog{\'i}as, capacidades, arquitecturas, protocolos o formatos de datos. En el contexto de los sistemas de informaci{\'o}n basados en IoT y los sistemas distribuidos, la heterogeneidad se manifiesta en la integraci{\'o}n de dispositivos con distintos recursos computacionales, plataformas hardware, sistemas operativos, lenguajes de programaci{\'o}n y mecanismos de comunicaci{\'o}n. La gesti{\'o}n de entornos heterog{\'e}neos constituye un desaf{\'i}o central en t{\'e}rminos de interoperabilidad, escalabilidad y mantenimiento, y motiva el uso de est{\'a}ndares, arquitecturas estratificadas y mecanismos de abstracci{\'o}n},
		plural={heterog{\'e}neos},
		firstplural={heterog{\'e}neos},
		see={interoperabilidad,procesamientodistribuido,arquitecturaestratificada,clienteservidor}
	}
	
	\newglossaryentry{hipermedia}{
		name={Hipermedia},
		text={hipermedia},
		sort={hipermedia},
		plural={hipermedias},
		first={hipermedia},
		firstplural={hipermedias},
		description={Modelo de representaci{\'o}n y navegaci{\'o}n de informaci{\'o}n que integra m{\'u}ltiples tipos de medios (texto, im{\'a}genes, audio, video y datos estructurados) interconectados mediante enlaces sem{\'a}nticos que permiten transiciones din{\'a}micas entre recursos. En el contexto de arquitecturas \gls{restful}, la hipermedia constituye un principio fundamental expresado como \emph{Hypermedia as the Engine of Application State (HATEOAS)}, donde las respuestas del servidor incluyen enlaces que describen las acciones posibles y gu{\'i}an la evoluci{\'o}n del estado de la aplicaci{\'o}n. En sistemas de informaci{\'o}n basados en IoT y SIbIoT, la hipermedia favorece el desacoplamiento entre clientes y servicios, mejora la interoperabilidad y permite la evoluci{\'o}n independiente de las interfaces sin afectar a los consumidores},
		type=main,
		symbol={🌐},
		see={rest,restful,api,clienteservidor,acoplamiento,desacoplamiento,sibiot}
	}
	
	\newglossaryentry{hogarinteligente}{
		name={Hogar inteligente},
		text={hogar inteligente},
		sort={hogar inteligente},
		plural={hogares inteligentes},
		first={hogar inteligente (\textit{Smart Home}},
		firstplural={hogares inteligentes},
		description={Entorno dom{\'o}stico automatizado que integra \gls{iot}, \glspl{sensor} y \glspl{actuador} para controlar \glspl{dispositivo}, optimizar recursos y mejorar confort, seguridad y eficiencia energ{\'e}tica.},
		type=main,
		symbol={🏠},
		see={cps,gestionenergetica,hems,matter,gateway},
		user1={Dom{\'o}tica residencial},
		user2={Confort y seguridad},
		user3={Eficiencia energ{\'e}tica}
	}
	
	\newglossaryentry{hojatecnica}{
		name={Hoja t{\'e}cnica},
		text={hoja t{\'e}cnica},
		sort={hoja tecnica},
		plural={hojas t{\'e}cnicas},
		first={hoja t{\'e}cnica},
		firstplural={hojas t{\'e}cnicas},
		description={Documento elaborado por el fabricante de un componente, dispositivo o sistema, que contiene informaci{\'o}n detallada sobre sus caracter{\'i}sticas, especificaciones, condiciones de funcionamiento, limitaciones y recomendaciones de uso. Una \emph{hoja t{\'e}cnica} incluye par{\'a}metros como voltaje, consumo energ{\'e}tico, rango de operaci{\'o}n, tolerancias, interfaces de comunicaci{\'o}n y diagramas de conexi{\'o}n. En el contexto de los \glspl{sibiot}, las hojas t{\'e}cnicas de \glspl{sensor} y \glspl{actuador} son fundamentales para garantizar una correcta integraci{\'o}n de hardware y software, facilitar la \gls{interoperabilidad}, y asegurar la \gls{fiabilidad} y \gls{seguridad} de los sistemas desarrollados.},
		type=main,
		symbol={📄},
		user1={Documento de especificaciones t{\'e}cnicas},
		user2={Gu{\'i}a para instalaci{\'o}n y uso seguro},
		user3={Base para integraci{\'o}n en SIbIoT},
		see={sensor,actuador,interoperabilidad,fiabilidad,seguridad,sibiot}
	}
	
	\newglossaryentry{horizontalidad}{
		name={Horizontalidad},
		text={horizontalidad},
		sort={horizontalidad},
		plural={horizontalidades},
		first={horizontalidad},
		firstplural={horizontalidades},
		description={Principio organizativo y de trabajo que privilegia la distribuci{\'o}n equilibrada de participaci{\'o}n, responsabilidad y comunicaci{\'o}n dentro de un equipo, reduciendo dependencias jer{\'a}rquicas r{\'\i}gidas. Favorece la toma de decisiones informada, la transparencia del progreso y la colaboraci{\'o}n interfuncional. En contextos {\'a}giles y en proyectos IoT/SIbIoT, la horizontalidad facilita la coordinaci{\'o}n entre perfiles de software, hardware, redes, datos y operaci{\'o}n, mejorando la detecci{\'o}n temprana de riesgos y el alineamiento de interfaces.},
		type=main,
		symbol={↔️},
		user1={Participaci{\'o}n distribuida},
		user2={Colaboraci{\'o}n interfuncional},
		user3={Transparencia y responsabilidad},
		see={equipoagil,construccioncolaborativa,gestionempresarial,scrum,kanban}
	}
	
	\newglossaryentry{htap}{
		name={HTAP},
		text={HTAP},
		plural={HTAPs},
		first={procesamiento h{\'i}brido de transacciones y an{\'a}lisis (\textit{Hybrid transaccional/Analytical Processing}, HTAP)},
		firstplural={HTAPs},
		description={\emph{HTAP} es un paradigma de \gls{basedatos} que integra en una misma plataforma las capacidades de procesamiento transaccional en l{\'i}nea (OLTP) y de procesamiento anal{\'i}tico en l{\'i}nea (OLAP). Esto permite que los sistemas ejecuten operaciones de lectura y escritura en tiempo real, mientras realizan simult{\'a}neamente an{\'a}lisis complejos sobre los mismos datos sin necesidad de duplicarlos. En el contexto de los \glspl{sibiot}, HTAP resulta fundamental para habilitar la \gls{tomadecisiones} en tiempo real, soportar \glspl{modelopredictivo} y mejorar la eficiencia en la gesti{\'o}n de flujos masivos de informaci{\'o}n heterog{\'e}nea.},
		type=main,
		symbol={⚡},
		user1={Integraci{\'o}n de OLTP y OLAP en una sola plataforma},
		user2={Soporte a an{\'a}lisis y transacciones en tiempo real},
		user3={Clave en ecosistemas Big Data e IoT},
		see={basedatos,bigdata,analisisdatos,spanner,sibiot}
	}
	
	% ===================== HTML =====================
	\newglossaryentry{html}
	{
		name={HTML},
		text={HTML},
		sort={html},
		plural={HTML},
		first={lenguaje de marcado de hipertexto (\textit{HyperText Markup Language}, HTML)},
		firstplural={HTML},
		description={\emph{HTML} es el lenguaje de marcado para estructurar contenidos en la web. Define encabezados, p{\'a}rrafos, enlaces, im{\'a}genes, tablas y formularios. Junto con \gls{css} y \gls{javascript} conforma el n{\'u}cleo de \glspl{aplicacionweb} y \glspl{dashboard} de \gls{iot}.},
		type=\glsdefaulttype,
		symbol={🌐},
		user1={Marcado web},
		user2={Documento hipertexto},
		user3={Estructura de interfaz},
		see={css,javascript,browser}
	}
	
	% ===================== HTML5 =====================
	\newglossaryentry{html5}
	{
		name={HTML5},
		text={HTML5},
		sort={html5},
		plural={HTML5},
		first={HTML5},
		firstplural={HTML5},
		description={\emph{HTML5} expande \gls{html} con nuevos elementos sem{\'a}nticos y APIs de \gls{javascript} para multimedia, almacenamiento y acceso a hardware, convirtiendo la web en una plataforma completa para \glspl{aplicacionweb} y paneles \gls{iot}.},
		type=\glsdefaulttype,
		symbol={🧩},
		user1={Web moderna},
		user2={APIs de navegador},
		user3={Sem{\'a}ntica enriquecida},
		see={html,css,javascript}
	}
	
	% ===================== HTTP =====================
	\newglossaryentry{http}
	{
		name={HTTP},
		text={HTTP},
		sort={http},
		plural={HTTP},
		first={protocolo de transferencia de hipertexto (\textit{HyperText Transfer Protocol}, HTTP)},
		firstplural={HTTP},
		description={\emph{HTTP} es un protocolo \gls{cliente}\textendash\gls{servidor} para transferencia de \glspl{dato} en la web. Define m{\'e}todos, c{\'o}digos de estado y cabeceras para intercambio de recursos; base de \glspl{api} REST usadas por plataformas \gls{iot}.},
		type=\glsdefaulttype,
		symbol={🔗},
		user1={Transporte web},
		user2={Protocolo de aplicaci{\'o}n},
		user3={REST/JSON},
		see={https,api}
	}
	
	% ===================== HTTPS =====================
	\newglossaryentry{https}
	{
		name={HTTPS},
		text={HTTPS},
		sort={https},
		plural={HTTPS},
		first={protocolo de transferencia de hipertexto seguro (\textit{HyperText Transfer Protocol Secure}, HTTPS)},
		firstplural={HTTPS},
		description={\emph{HTTPS} es \gls{http} con cifrado (\gls{tls}) para garantizar confidencialidad, \gls{integridad} y autenticidad en la transmisi{\'o}n. Es esencial en \glspl{aplicacionweb} y servicios \gls{iot} que requieren protecci{\'o}n frente a intercepci{\'o}n o manipulaci{\'o}n de \glspl{dato}.},
		type=\glsdefaulttype,
		symbol={🔒},
		user1={HTTP seguro},
		user2={TLS/SSL},
		user3={Cifrado en tr{\'a}nsito},
		see={http,tls,ciberseguridad}
	}
	
	% ===================== IA =====================
	\newglossaryentry{ia}
	{
		name={AI},
		text={AI},
		sort={ai},
		plural={AIs},
		first={inteligencia artificial (\textit{Artificial Intelligence}, AI)},
		firstplural={AIs},
		description={La \emph{IA} estudia y construye sistemas capaces de realizar tareas que requieren inteligencia humana (reconocimiento de \glspl{patron}, \gls{tomadecisiones}, \gls{procesamientodatos} avanzado y \gls{aprendizajeautomatico}). En \gls{iot} potencia \glspl{modelopredictivo} y \gls{mantenimientopredictivo}.},
		type=\glsdefaulttype,
		symbol={🧠},
		user1={Artificial Intelligence},
		user2={Sistemas inteligentes},
		user3={Algoritmos de aprendizaje},
		see={aprendizajeautomatico}
	}
	
	\newglossaryentry{idempotente}{
		name={Idempotente},
		text={idempotente},
		sort={idempotente},
		plural={idempotentes},
		first={idempotente},
		firstplural={idempotentes},
		description={Propiedad de una operaci{\'o}n mediante la cual m{\'u}ltiples ejecuciones con los mismos par{\'a}metros producen el mismo resultado que una {\'u}nica ejecuci{\'o}n, sin generar efectos secundarios adicionales. En sistemas distribuidos y servicios web, una operaci{\'o}n idempotente garantiza consistencia frente a reintentos, fallos de red o duplicaci{\'o}n de mensajes, siendo especialmente relevante en m{\'e}todos HTTP como GET, PUT y DELETE.},
		type=main,
		symbol={♻️},
		user1={Operaci{\'o}n sin efectos acumulativos},
		user2={Consistencia ante reintentos},
		user3={Propiedad clave en servicios distribuidos},
		see={consistencia,transaccion,servicioweb,rest}
	}
	
	% ===================== Identificaci{\'o}n =====================
	\newglossaryentry{identificacion}
	{
		name={Identificaci{\'o}n},
		text={identificaci{\'o}n},
		sort={identificacion},
		plural={identificaciones},
		first={identificaci{\'o}n},
		firstplural={identificaciones},
		description={La \emph{identificaci{\'o}n} asigna una identidad {\'u}nica y espec{\'i}fica a un objeto, persona o \gls{patron}. En \gls{iot} se aplica en biometr{\'i}a, RFID, visi{\'o}n artificial y trazabilidad de activos, asegurando unicidad y confiabilidad para la \gls{seguridad}.},
		type=\glsdefaulttype,
		symbol={🆔},
		user1={Asignaci{\'o}n de identidad},
		user2={Identidad digital},
		user3={Etiquetado {\'u}nico},
		see={reconocimiento,deteccion}
	}
	
	\newglossaryentry{incremental}{
		name={Incremental},
		text={incremental},
		sort={incremental},
		plural={incrementales},
		first={incremental},
		firstplural={incrementales},
		description={Propiedad de un enfoque de desarrollo en el que la soluci{\'o}n se construye mediante entregas sucesivas que a{\~n}aden capacidad funcional y refinan artefactos existentes, manteniendo el sistema en un estado potencialmente operable. Cada incremento produce valor verificable y reduce incertidumbre t{\'e}cnica. En sistemas IoT, lo incremental permite integrar progresivamente dispositivos, servicios y anal{\'\i}tica, validando compatibilidad, rendimiento y seguridad por etapas y disminuyendo el riesgo de integraci{\'o}n tard{\'\i}a.},
		type=main,
		symbol={➕},
		user1={Entregas sucesivas},
		user2={Valor verificable por incremento},
		user3={Reducci{\'o}n de riesgo},
		see={iterativa,ciclodesarrollo,entregable,unidadfuncional,desarrolloagilsoftware}
	}
	
	\newglossaryentry{incremento}{
		name={Incremento},
		text={incremento},
		sort={incremento},
		plural={incrementos},
		first={incremento},
		firstplural={incrementos},
		description={Resultado tangible y funcional producido al finalizar un Sprint, que integra todos los elementos completados del \Gls{sprintbacklog} y los \emph{incrementos} previos. El incremento debe cumplir con la Definici{\'o}n de Terminado y estar en condiciones potenciales de ser liberado o desplegado. Representa un avance medible hacia el objetivo del producto y constituye evidencia concreta del progreso del equipo en cada iteraci{\'o}n.},
		type=main,
		symbol={📦},
		user1={Entrega funcional acumulativa},
		user2={Cumple definici{\'o}n de terminado},
		user3={Avance medible del producto},
		see={sprint,sprintbacklog,productbacklog,agil}
	}
	
	\newglossaryentry{indicador}{
		name={Indicador},
		text={indicador},
		plural={indicadores},
		first={indicador},
		firstplural={indicadores},
		see={kpi},
		sort={indicador},
		description={Medida cuantitativa o cualitativa utilizada para describir, evaluar o monitorear el estado, comportamiento o desempe{\~n}o de un proceso, sistema o fen{\'o}meno. En SIbIoT, los indicadores se derivan del an{\'a}lisis de datos IoT y de sistemas empresariales, y pueden representar variables operativas, resultados anal{\'i}ticos o m{\'e}tricas de control; cuando est{\'a}n alineados con objetivos estrat{\'e}gicos, se formalizan como KPIs}
	}
	
	% ===================== Industry 4.0 =====================
	\newglossaryentry{industry40}
	{
		name={Industry~4.0},
		text={Industry 4.0},
		sort={industry 4.0},
		plural={iniciativas Industry 4.0},
		first={Cuarta revoluci{\'o}n industrial (\textit{Industry 4.0})},
		firstplural={iniciativas Industry~4.0},
		description={\emph{Industry~4.0} describe la integraci{\'o}n de \gls{iot}, \gls{ia}, rob{\'o}tica y an{\'a}lisis de \gls{bigdata} en procesos industriales para fabricar con flexibilidad, trazabilidad y eficiencia.},
		type=\glsdefaulttype,
		symbol={🏭},
		user1={F{\'a}brica inteligente},
		user2={Transformaci{\'o}n digital},
		user3={Ciberf{\'i}sico},
		see={sibiot}
	}
	
	% ===================== Implementaci{\'o}n =====================
	\newglossaryentry{impl}
	{
		name={Implementaci{\'o}n},
		text={implementaci{\'o}n},
		sort={implementacion},
		plural={implementaciones},
		first={implementaci{\'o}n},
		firstplural={implementaciones},
		description={La \emph{implementaci{\'o}n} es la fase del ciclo de vida en que se construye la soluci{\'o}n y tambi{\'e}n el acto de materializar una interfaz mediante una clase concreta. En \gls{iot} integra \gls{hardware}, firmware y software.},
		type=\glsdefaulttype,
		symbol={🛠️},
		user1={Construcci{\'o}n},
		user2={Codificaci{\'o}n},
		user3={Despliegue},
		see={arquitectura,desarrolloagilsoftware}
	}
	
	\newglossaryentry{inferencia}{
		name={Inferencia},
		text={inferencia},
		sort={inferencia},
		plural={inferencias},
		first={inferencia},
		firstplural={inferencias},
		description={Proceso mediante el cual un sistema, algoritmo o entidad anal{\'i}tica deduce, estima o deriva nueva informaci{\'o}n a partir de datos observados, reglas predefinidas o modelos matem{\'a}ticos. En el contexto de los \glspl{sibiot}, la inferencia permite transformar \glspl{dato} crudos en conocimiento accionable mediante t{\'e}cnicas como el razonamiento l{\'o}gico, los modelos probabil{\'i}sticos, el aprendizaje autom{\'a}tico y los mecanismos de toma de decisiones basadas en patrones},
		symbol={∴},
		see={razonamiento,modelo},
		type=\glsdefaulttype
	}
	
	% ===================== Infraestructura =====================
	\newglossaryentry{infraestructura}
	{
		name={Infraestructura},
		text={infraestructura},
		sort={infraestructura},
		plural={infraestructuras},
		first={infraestructura},
		firstplural={infraestructuras},
		description={La \emph{infraestructura} comprende \glspl{dispositivo} terminales (\gls{sensor}, \gls{actuador}), \glspl{gateway}, \glspl{red}, \glspl{servidor}, \gls{nube}, \glspl{basedatos} y servicios de an{\'a}lisis que habilitan soluciones \gls{iot} escalables, seguras y eficientes.},
		type=\glsdefaulttype,
		symbol={🏗️},
		user1={Plataforma t{\'e}cnica},
		user2={Capas de soporte},
		user3={Base operativa},
		see={infraestructurared}
	}
	
	\newglossaryentry{infraestructuracomunicacion}{
		name={Infraestructura de comunicaci{\'o}n},
		text={infraestructura de comunicaci{\'o}n},
		sort={infraestructura comunicacion},
		plural={infraestructuras de comunicaci{\'o}n},
		first={infraestructura de comunicaci{\'o}n},
		firstplural={infraestructuras de comunicaci{\'o}n},
		description={Conjunto de medios f{\'i}sicos y l{\'o}gicos que permiten la transmisi{\'o}n de datos entre \glspl{dispositivo}, sistemas y plataformas de procesamiento. En los \glspl{sibiot}, la \emph{infraestructura de comunicaci{\'o}n} es esencial para garantizar conectividad, fiabilidad y escalabilidad en entornos distribuidos.},
		type=main,
		symbol={📡},
		user1={Transmisi{\'o}n de datos},
		user2={Conectividad},
		user3={Redes de comunicaci{\'o}n},
		see={arquitecturaabierta,redsensor,ip}
	}
	
	\newglossaryentry{infografia}{
		name={infograf{\'i}a},
		text={infograf{\'i}a},
		sort={infografia},
		plural={infograf{\'i}as},
		first={infograf{\'i}a},
		firstplural={infograf{\'i}as},
		description={Recurso visual que combina elementos gr{\'a}ficos, diagramas, texto y simbolog{\'i}a para comunicar informaci{\'o}n t{\'e}cnica de manera clara y sint{\'e}tica. En el contexto de los \glspl{sibiot}, las infograf{\'i}as permiten representar de forma estructurada conceptos complejos, como arquitecturas, flujos de datos, desaf{\'i}os de \gls{conectividad} y estrategias de \gls{mitigacion}, facilitando la comprensi{\'o}n y el an{\'a}lisis de los sistemas},
		symbol={🖼️},
		see={diagrama,representacion,visualizacion},
		type=main
	}
	
	\newglossaryentry{informacion}{
		name={Informaci{\'o}n},
		text={informaci{\'o}n},
		sort={informacion},
		plural={informaciones},
		first={informaci{\'o}n},
		firstplural={informaciones},
		description={Resultado de la transformaci{\'o}n de \glspl{dato} en un conocimiento {\'u}til y significativo para el usuario o sistema. La \emph{informaci{\'o}n} surge al aplicar procesos de an{\'a}lisis, clasificaci{\'o}n, interpretaci{\'o}n y contextualizaci{\'o}n sobre los datos. En el contexto de los \glspl{sibiot}, la informaci{\'o}n generada a partir de datos capturados en tiempo real permite optimizar procesos, apoyar la \gls{tomadecisiones}, detectar patrones y anticipar \glspl{evento}. Constituye un recurso estrat{\'e}gico para organizaciones y \glspl{sistemainteligente}.},
		type=main,
		symbol={ℹ️},
		user1={Conocimiento estructurado},
		user2={Datos procesados},
		user3={Resultado del an{\'a}lisis de datos},
		see={dato,procesamientodatos,mineriadatos}
	}
	
	\newglossaryentry{infraestructuracodigo}{
		name={Infraestructura como C{\'o}digo},
		text={infraestructura como c{\'o}digo},
		sort={infraestructuracomocodigo},
		type={main},
		plural={infraestructuras como c{\'o}digo},
		first={infraestructura como c{\'o}digo},
		firstplural={infraestructuras como c{\'o}digo},
		description={Enfoque para describir, aprovisionar y configurar infraestructura (c{\'o}mputo, red, almacenamiento y servicios) mediante especificaciones versionadas y ejecutables, permitiendo reproducibilidad, consistencia entre entornos y auditor{\'\i}a de cambios.},
		see={[see also]{devops,cicd}}
	}
	
	\newglossaryentry{infraestructuradistribuida}{
		name={Infraestructura distribuida},
		text={infraestructura distribuida},
		sort={infraestructura distribuida},
		plural={infraestructuras distribuidas},
		first={infraestructura distribuida},
		firstplural={infraestructuras distribuidas},
		description={Conjunto de recursos computacionales, tanto f{\'i}sicos como l{\'o}gicos, que se encuentran distribuidos en m{\'u}ltiples nodos interconectados y que colaboran de manera coordinada para proveer servicios, procesamiento y almacenamiento de datos. La \emph{infraestructura distribuida} es la base de arquitecturas modernas en la nube, \gls{edgecomputing}, \gls{fogcomputing} y sistemas \gls{iot}, ya que permite escalar horizontalmente, mejorar la tolerancia a fallos, optimizar la latencia y aprovechar la ubicuidad de los dispositivos conectados.},
		type=main,
		symbol={🌐},
		user1={Conjunto de recursos distribuidos en nodos},
		user2={Base de arquitecturas en la nube, edge y fog},
		user3={Soporte a escalabilidad y resiliencia en IoT},
		see={arquitecturadistribuida,infraestructura,edgecomputing,fogcomputing,sibiot}
	}
	
	% ===================== Infraestructura de red =====================
	\newglossaryentry{infraestructurared}
	{
		name={Infraestructura de red},
		text={infraestructura de red},
		sort={infraestructura de red},
		plural={infraestructuras de red},
		first={infraestructura de red},
		firstplural={infraestructuras de red},
		description={La \emph{infraestructura de red} incluye enrutadores, switches, puntos de acceso, cables, antenas, \glspl{servidor}, protocolos y software de control. En \gls{iot} garantiza \gls{conectividad}, \gls{escalabilidad}, eficiencia y \gls{seguridad} de las comunicaciones.},
		type=\glsdefaulttype,
		symbol={🖧},
		user1={Backbone},
		user2={Red de transporte},
		user3={Capa de comunicaci{\'o}n},
		see={ethernet,lan}
	}
	
	\newglossaryentry{ingenieria}{
		name={Ingenier{\'i}a},
		text={ingenier{\'i}a},
		sort={ingenieria},
		plural={ingenier{\'i}as},
		first={ingenier{\'i}a},
		firstplural={ingenier{\'i}as},
		description={Disciplina profesional y acad{\'e}mica orientada a la aplicaci{\'o}n de principios cient{\'i}ficos, matem{\'a}ticos y t{\'e}cnicos para la concepci{\'o}n, dise{\~n}o, construcci{\'o}n, evaluaci{\'o}n y mejora de sistemas, procesos y productos. La \emph{ingenier{\'i}a} abarca diversas ramas como la ingenier{\'i}a civil, electr{\'o}nica, industrial, mec{\'a}nica, qu{\'i}mica e inform{\'a}tica, entre otras. En el contexto de los \glspl{sibiot}, la ingenier{\'i}a integra hardware, software, comunicaciones y gesti{\'o}n de datos para crear soluciones robustas, eficientes y sostenibles que aporten valor a la sociedad.},
		type=main,
		symbol={⚙️},
		user1={Aplicaci{\'o}n de la ciencia a la tecnolog{\'i}a},
		user2={Dise{\~n}o y optimizaci{\'o}n de sistemas},
		user3={Disciplina profesional},
		see={ingeniero,ingenieriasoftware,sistemainformatico}
	}
		
	\newglossaryentry{ingenieriasistemas}{
		name={Ingenier{\'\i}a de sistemas},
		text={ingenier{\'\i}a de sistemas},
		sort={ingenieria de sistemas},
		plural={ingenier{\'\i}as de sistemas},
		first={ingenier{\'\i}a de sistemas},
		firstplural={ingenier{\'\i}as de sistemas},
		description={La \emph{ingenier{\'\i}a de sistemas} es la disciplina que coordina la especificaci{\'o}n, el dise{\~n}o, la verificaci{\'o}n y la operaci{\'o}n de sistemas complejos considerando de forma integrada componentes t{\'e}cnicos y restricciones del entorno. En el desarrollo de \glspl{cps}, esta disciplina resulta necesaria porque el comportamiento final depende de la co-dependencia entre proceso f{\'\i}sico, software, comunicaciones y tiempo; por ello, la gesti{\'o}n de interfaces, supuestos y evidencias de verificaci{\'o}n se vuelve parte del trabajo t{\'e}cnico, no un aspecto externo.},
		type=main,
		symbol={🧩},
		user1={Integraci{\'o}n multi-dominio},
		user2={Gesti{\'o}n de interfaces y supuestos},
		user3={Verificaci{\'o}n y operaci{\'o}n del sistema},
		see={arquitectura,verificacion,integracion,cps}
	}
	
	\newglossaryentry{ingenieriasoftware}{
		name={Ingenier{\'i}a de software},
		text={ingenier{\'i}a de software},
		sort={ingenieria de software},
		plural={ingenier{\'i}as de software},
		first={ingenier{\'i}a de software},
		firstplural={ingenier{\'i}as de software},
		description={Disciplina de la ingenier{\'i}a orientada a la aplicaci{\'o}n sistem{\'a}tica, disciplinada y cuantificable de principios cient{\'i}ficos, t{\'e}cnicos y metodol{\'o}gicos para el dise{\~n}o, desarrollo, validaci{\'o}n, despliegue, mantenimiento y evoluci{\'o}n del software. La \emph{ingenier{\'i}a de software} busca garantizar la calidad, la \gls{fiabilidad}, la \gls{seguridad}, la \gls{interoperabilidad} y la \gls{sostenibilidad} de las aplicaciones y sistemas, atendiendo tanto a requisitos funcionales como no funcionales. En el contexto de los \glspl{sibiot}, proporciona m{\'e}todos y herramientas para gestionar la complejidad del ciclo de vida de sistemas inteligentes distribuidos, asegurando la \gls{eficienciaoperativa}, la \gls{productividad} de los equipos y la satisfacci{\'o}n de los usuarios finales.},
		type=main,
		symbol={💻},
		user1={Disciplina de ingenier{\'i}a aplicada al software},
		user2={Metodolog{\'i}as y procesos de desarrollo},
		user3={Calidad y sostenibilidad de sistemas},
		see={ingeniero,ingenieria,software,sibiot,buenaspracticas}
	}

	\newglossaryentry{ingeniero}{
		name={Ingeniero},
		text={ingeniero},
		sort={ingeniero},
		plural={ingenieros},
		first={ingeniero},
		firstplural={ingenieros},
		description={Profesional formado en el campo de la ingenier{\'i}a, especializado en aplicar conocimientos cient{\'i}ficos, matem{\'a}ticos, tecnol{\'o}gicos y metodol{\'o}gicos para dise{\~n}ar, desarrollar, implementar, mantener y optimizar sistemas, procesos, productos o servicios. Un \emph{ingeniero} combina teor{\'i}a y pr{\'a}ctica para resolver problemas reales de manera eficiente, sostenible y segura. En el contexto de la ingenier{\'i}a de software, de los \glspl{sibiot} y de la ingenier{\'i}a de sistemas inteligentes, los ingenieros participan en todas las fases del ciclo de vida, desde la conceptualizaci{\'o}n y la arquitectura hasta las pruebas, la seguridad, la \gls{sostenibilidad} y la gesti{\'o}n del impacto social y {\'e}tico.},
		type=main,
		symbol={👷},
		user1={Profesional de la ingenier{\'i}a},
		user2={Dise{\~n}o y optimizaci{\'o}n de sistemas},
		user3={Responsabilidad social y t{\'e}cnica},
		see={ingenieria,sibiot,software,sostenibilidad}
	}
	
	\newglossaryentry{ingesta}{
		name={Ingesta},
		text={ingesta},
		plural={ingestas},
		first={ingesta},
		firstplural={ingestas},
		see={captura,procesamiento,analisis,decision},
		sort={ingesta},
		description={Proceso mediante el cual los datos son incorporados desde sus fuentes de origen hacia las plataformas de procesamiento y almacenamiento. En SIbIoT, la ingesta puede realizarse en tiempo real (\textit{streaming}) o por lotes (\textit{batch}), e incluye validaci{\'o}n inicial, control de duplicados y etiquetado con metadatos}
	}
	
	\newglossaryentry{ingestadatos}{
		name={Ingesta de datos},
		text={ingesta de datos},
		sort={ingesta de datos},
		plural={ingestas de datos},
		first={ingesta de datos},
		firstplural={ingestas de datos},
		description={Proceso mediante el cual los datos son capturados, recibidos y transferidos desde m{\'u}ltiples fuentes hacia un sistema de almacenamiento o procesamiento para su posterior an{\'a}lisis. La \emph{ingesta de datos} puede realizarse en modo batch (por lotes) o en tiempo real (\textit{streaming}), y constituye la fase inicial de un \gls{pipeline} de datos. En entornos \gls{iot} y \glspl{sibiot}, la ingesta implica la recepci{\'o}n de flujos provenientes de \glspl{sensor}, \glspl{dispositivoconectado}, \glspl{gateway} y sistemas externos, asegurando validaci{\'o}n, normalizaci{\'o}n, control de calidad y trazabilidad antes de su almacenamiento o procesamiento anal{\'i}tico.},
		type=main,
		symbol={📥},
		user1={Captura y recepci{\'o}n de datos},
		user2={Entrada en pipeline de datos},
		user3={Flujos batch y streaming},
		see={pipeline,procesamientodatos,etl,telemetria}
	}
	
	% ===================== Iniciativa {\'a}gil =====================
	\newglossaryentry{iniciativaagil}
	{
		name={Iniciativa {\'a}gil},
		text={iniciativa {\'a}gil},
		sort={iniciativa agil},
		plural={iniciativas {\'a}giles},
		first={iniciativa {\'a}gil},
		firstplural={iniciativas {\'a}giles},
		description={Una \emph{iniciativa {\'a}gil} es una acci{\'o}n estrat{\'e}gica orientada a la adopci{\'o}n, extensi{\'o}n o consolidaci{\'o}n de principios y pr{\'a}cticas del Manifiesto {\'A}gil en una organizaci{\'o}n o equipo. Puede materializarse como la transici{\'o}n desde modelos tradicionales hacia enfoques iterativos e incrementales, la implantaci{\'o}n de marcos como Scrum, Kanban o SAFe, o el impulso de la cultura colaborativa, la entrega continua de valor y la mejora adaptativa. En proyectos \gls{iot}, estas iniciativas facilitan integrar \gls{hardware} y software tempranamente, gestionar la complejidad y adaptarse r{\'a}pido a cambios tecnol{\'o}gicos o de mercado.},
		type=\glsdefaulttype,
		symbol={🚀},
		user1={Transformaci{\'o}n {\'a}gil},
		user2={Entrega continua},
		user3={Cultura de mejora},
		see={equipoagil,desarrolloagilsoftware,kanban}
	}
	
	\newglossaryentry{innovacion}{
		name={Innovaci{\'o}n},
		text={innovaci{\'o}n},
		sort={innovacion},
		first={innovaci{\'o}n},
		description={Introducci{\'o}n de mejoras o enfoques nuevos que generan valor, ya sea por nuevas funcionalidades, procesos m{\'a}s eficientes o mayor calidad de servicio en un contexto determinado.},
		type=main,
		symbol={💡},
		user1={Novedad con valor},
		user2={Mejora},
		user3={Transformaci{\'o}n},
		see={resultado,objetivo}
	}
	
	\newglossaryentry{integraciondatos}{
		name={Integraci{\'o}n de datos},
		text={integraci{\'o}n de datos},
		sort={integracion datos},
		plural={procesos de integraci{\'o}n de datos},
		first={integraci{\'o}n de datos},
		firstplural={procesos de integraci{\'o}n de datos},
		description={Proceso mediante el cual datos provenientes de m{\'u}ltiples fuentes heterog{\'e}neas se combinan en una visi{\'o}n coherente y unificada. En los \glspl{sibiot}, la \emph{integraci{\'o}n de datos} permite transformar mediciones aisladas en informaci{\'o}n estructurada para an{\'a}lisis y decisi{\'o}n.},
		type=main,
		symbol={🔗},
		user1={Unificaci{\'o}n de datos},
		user2={Datos heterog{\'e}neos},
		user3={Interoperabilidad},
		see={integraciondatosoperativos,sistemainformacioncorporativo}
	}
	
	\newglossaryentry{integraciondatosheterogeneos}
	{
		name={Integraci{\'o}n de datos heterog{\'e}neos},
		text={integraci{\'o}n de datos heterog{\'e}neos},
		sort={integracion datos heterogeneos},
		plural={integraciones de datos heterog{\'e}neos},
		first={integraci{\'o}n de datos heterog{\'e}neos},
		firstplural={integraciones de datos heterog{\'e}neos},
		description={La \emph{integraci{\'o}n de datos heterog{\'e}neos} es la capacidad arquitect{\'o}nica de consolidar informaci{\'o}n proveniente de m{\'u}ltiples fuentes con estructuras, formatos y frecuencias distintas bajo criterios de coherencia sem{\'a}ntica y alineaci{\'o}n temporal. En entornos \gls{sibiot}, esta propiedad permite vincular datos f{\'i}sicos de sensores con registros empresariales y pol{\'i}ticas organizacionales, garantizando consistencia anal{\'i}tica y trazabilidad operativa.},
		type=\glsdefaulttype,
		symbol={🔗},
		user1={Normalizaci{\'o}n de datos},
		user2={Metadatos},
		user3={Interoperabilidad},
		see={ingestadatos,procesamientoanalitico,arquitectura}
	}
	
	\newglossaryentry{integraciondatosoperativos}{
		name={Integraci{\'o}n de datos operativos},
		text={integraci{\'o}n de datos operativos},
		sort={integracion datos operativos},
		plural={procesos de integraci{\'o}n de datos operativos},
		first={integraci{\'o}n de datos operativos},
		firstplural={procesos de integraci{\'o}n de datos operativos},
		description={Incorporaci{\'o}n de datos generados por procesos f{\'i}sicos y operativos dentro de sistemas de informaci{\'o}n para su almacenamiento, an{\'a}lisis y explotaci{\'o}n. Este proceso es central en los \glspl{sibiot}, al conectar el mundo f{\'i}sico con la gesti{\'o}n digital de la informaci{\'o}n.},
		type=main,
		symbol={📥},
		user1={Datos de operaci{\'o}n},
		user2={Procesos f{\'i}sicos},
		user3={Sistemas informacionales},
		see={integraciondatos,scada}
	}
	
	\newglossaryentry{integracionsoftware}{
		name={Integraci{\'o}n de software},
		text={integraci{\'o}n de software},
		sort={integracionsoftware},
		plural={integraciones de software},
		first={integraci{\'o}n de software},
		firstplural={integraciones de software},
		description={Proceso mediante el cual diferentes aplicaciones, componentes o sistemas de software se interconectan para operar de manera coordinada, compartiendo datos, funcionalidades y flujos de control. En los \glspl{sibiot}, la integraci{\'o}n de software es esencial para enlazar dispositivos IoT, servicios backend, plataformas en la nube y aplicaciones web o m{\'o}viles, garantizando interoperabilidad, consistencia de datos y continuidad operativa en arquitecturas distribuidas.},
		type=main,
		symbol={🔗},
		user1={Interoperabilidad},
		user2={Conexi{\'o}n de sistemas},
		user3={Coordinaci{\'o}n de aplicaciones},
		see={arquitecturadistribuida,servicioweb,plataformaweb,gestiondatos}
	}
	
	% ===================== Integridad =====================
	\newglossaryentry{integridad}
	{
		name={Integridad},
		text={integridad},
		sort={integridad},
		plural={integridades},
		first={integridad},
		firstplural={integridades},
		description={La \emph{integridad} es la propiedad por la cual los \glspl{dato} y los \glspl{sistema} se mantienen completos, exactos y no alterados de forma no autorizada, ya sea en reposo, en tr{\'a}nsito o durante el \gls{procesamientodatos}. Se salvaguarda con sumas de verificaci{\'o}n, \textit{hashes} criptogr{\'a}ficos, \gls{firmadigital}, control de versiones y \glspl{api} con validaciones robustas. En \gls{iot}/\glspl{sibiot}, es esencial para la \gls{tomadecisiones} fiable y la auditor{\'i}a trazable.},
		type=\glsdefaulttype,
		symbol={🧷},
		user1={Calidad de datos},
		user2={Controles criptogr{\'a}ficos},
		user3={Trazabilidad},
		see={seguridad,disponibilidad,firmadigital,api}
	}
	
	% ===================== Intel Edison =====================
	\newglossaryentry{intedison}
	{
		name={Intel Edison},
		text={Intel Edison},
		sort={intel edison},
		plural={Intel Edison},
		first={Intel Edison},
		firstplural={Intel Edison},
		description={\emph{Intel Edison} es una plataforma de desarrollo compacta creada por Intel para aplicaciones \gls{iot} y computaci{\'o}n embebida. Integra Wi\textendash Fi, Bluetooth, almacenamiento y un procesador Intel Atom de doble n{\'u}cleo, facilitando prototipos conectados y \glspl{placadesarrollo} modulares.},
		type=\glsdefaulttype,
		symbol={🧩},
		user1={M{\'o}dulo IoT},
		user2={Conectividad embebida},
		user3={Prototipado r{\'a}pido},
		see={placadesarrollo,esp32,esp8266}
	}
	
	% ===================== Integraci{\'o}n =====================
	\newglossaryentry{integracion}{
		name={Integraci{\'o}n},
		text={integraci{\'o}n},
		sort={integracion},
		plural={integraciones},
		first={integraci{\'o}n},
		firstplural={integraciones},
		description={La \emph{integraci{\'o}n} es el proceso mediante el cual se combinan componentes, m{\'o}dulos o \glspl{sistema} independientes para que funcionen como una unidad coherente. En el contexto de \gls{iot}, la integraci{\'o}n abarca la conexi{\'o}n de \glspl{sensor}, \glspl{actuador}, \glspl{gateway}, servicios en la \gls{nube} y plataformas de gesti{\'o}n y \gls{procesamientodatos}, garantizando \gls{interoperabilidad}, sincronizaci{\'o}n y flujos eficientes de \glspl{dato}. Este proceso es clave para habilitar ecosistemas inteligentes, escalables y resilientes que soporten \gls{tomadecisiones} en \gls{tiemporeal} dentro de \glspl{sibiot}.},
		type=main,
		symbol={🔗},
		user1={Orquestaci{\'o}n extremo a extremo},
		user2={Acoplamiento controlado entre m{\'o}dulos},
		user3={Flujos de datos y control consistentes},
		see={interoperabilidad,arquitecturadistribuida,infraestructuradistribuida,api}
	}
	
	% ===================== Interacci{\'o}n =====================
	\newglossaryentry{interaccion}
	{
		name={Interacci{\'o}n},
		text={interacci{\'o}n},
		sort={interaccion},
		plural={interacciones},
		first={interacci{\'o}n},
		firstplural={interacciones},
		description={La \emph{interacci{\'o}n} es el proceso por el cual dos o m{\'a}s entidades (personas, \glspl{sistema} o componentes) se comunican o influyen mediante intercambio de \glspl{dato}, se{\~n}ales o acciones. En TIC incluye la \gls{interfaz} y la \gls{ux}, y en \gls{iot} tambi{\'e}n la comunicaci{\'o}n entre dispositivos y servicios.},
		type=\glsdefaulttype,
		symbol={🔄},
		user1={HCI/UX},
		user2={Se{\~n}ales y eventos},
		user3={Feedback},
		see={interfaz,evento,usuario}
	}
	
	\newglossaryentry{interconectado}{
		name={interconectado},
		text={interconectado},
		sort={interconectado},
		plural={interconectados},
		first={interconectado},
		firstplural={interconectados},
		description={Estado en el que dos o m{\'a}s sistemas, dispositivos o componentes mantienen enlaces de comunicaci{\'o}n que permiten el intercambio de informaci{\'o}n, sincronizaci{\'o}n de operaciones y coordinaci{\'o}n funcional. Un sistema \emph{interconectado} se caracteriza por su capacidad de interoperabilidad, flujo de datos bidireccional y dependencia colaborativa entre nodos. En el contexto de los \glspl{sibiot}, la interconexi{\'o}n es esencial para habilitar procesos distribuidos, agregaci{\'o}n de datos, toma de decisiones cooperativa y la articulaci{\'o}n de arquitecturas inteligentes basadas en m{\'u}ltiples fuentes de informaci{\'o}n},
		type=main,
		symbol={🔗},
		user1={Conexi{\'o}n funcional entre nodos},
		user2={Intercambio y sincronizaci{\'o}n de datos},
		user3={Soporte de interoperabilidad en arquitecturas distribuidas},
		see={sistemaconectado,interoperabilidad,redcomunicacion}
	}
	
	\newglossaryentry{interconexion}{
		name={Interconexi{\'o}n},
		text={interconexi{\'o}n},
		sort={interconexion},
		plural={interconexiones},
		first={interconexi{\'o}n},
		firstplural={interconexiones},
		description={La \emph{interconexi{\'o}n} es el proceso mediante el cual se enlazan distintos \glspl{sistema}, \glspl{dispositivo} o \glspl{red} para permitir el flujo de informaci{\'o}n y el uso compartido de recursos. Implica el establecimiento de enlaces f{\'i}sicos, l{\'o}gicos o virtuales soportados por \glspl{protocolocomunicacion}, hardware especializado (como \glspl{gateway} o \glspl{enrutador}) y est{\'a}ndares de \gls{interoperabilidad}. En \glspl{sibiot}, la interconexi{\'o}n asegura que sensores, actuadores y plataformas heterog{\'e}neas trabajen de manera coordinada, habilitando \gls{escalabilidad}, \gls{continuidad} y confiabilidad en los servicios distribuidos.},
		type=main,
		symbol={🔗},
		user1={Vinculaci{\'o}n de redes y sistemas heterog{\'e}neos},
		user2={Soporte en protocolos y hardware especializado},
		user3={Base de la interoperabilidad en SIbIoT},
		see={comunicacion,conectividad,protocolocomunicacion,interoperabilidad,sibiot}
	}
	
	\newglossaryentry{interdependiente}{
		name={Interdependiente},
		text={interdependiente},
		sort={interdependiente},
		plural={interdependientes},
		first={interdependiente},
		firstplural={interdependientes},
		description={Condici{\'o}n en la cual dos o m{\'a}s componentes, subsistemas u organizaciones mantienen una relaci{\'o}n de dependencia mutua, de modo que el funcionamiento, desempe{\~n}o o evoluci{\'o}n de uno influye directa o indirectamente en los dem{\'a}s. En arquitecturas distribuidas y en \glspl{sibiot}, la interdependencia implica coordinaci{\'o}n estructural, coherencia informacional y mecanismos de sincronizaci{\'o}n que permitan preservar estabilidad y resiliencia del sistema global.},
		type=main,
		symbol={🔗},
		user1={Dependencia mutua estructural},
		user2={Relaci{\'o}n funcional rec{\'i}proca},
		user3={Coordinaci{\'o}n entre subsistemas},
		see={interoperabilidad,integracion,infraestructuradistribuida,sistema}
	}
	
	% ===================== Interfaz =====================
	\newglossaryentry{interfaz}
	{
		name={Interfaz},
		text={interfaz},
		sort={interfaz},
		plural={interfaces},
		first={interfaz},
		firstplural={interfaces},
		description={La \emph{interfaz} es la din{\'a}mica f{\'i}sica y l{\'o}gica de interconexi{\'o}n (interacci{\'o}n) entre dos \glspl{sistema} independientes o entre un \gls{sistemainformatico} y la persona (usuario). Define protocolos, formatos y controles de entrada/salida que habilitan usabilidad y \gls{interoperabilidad}.},
		type=\glsdefaulttype,
		symbol={🖱️},
		user1={Interfaz de usuario},
		user2={Interfaz de sistema},
		user3={API/GUI},
		see={interaccion,api,browser}
	}
	
	\newglossaryentry{interfazvideo}{
		name={Interfaz de video},
		text={interfaz de video},
		sort={interfazvideo},
		plural={interfaces de video},
		first={interfaz de video},
		firstplural={interfaces de video},
		description={Conjunto de especificaciones f{\'i}sicas, el{\'e}ctricas y l{\'o}gicas que permiten la transmisi{\'o}n de se{\~n}ales de video, y en muchos casos audio y datos auxiliares, entre un dispositivo generador de imagen y un dispositivo de visualizaci{\'o}n. Una interfaz de video define aspectos como el tipo de se{\~n}al (anal{\'o}gica o digital), el protocolo de comunicaci{\'o}n, el ancho de banda disponible y los mecanismos de sincronizaci{\'o}n. En sistemas de informaci{\'o}n y arquitecturas computacionales, las interfaces de video son esenciales para la representaci{\'o}n visual de datos, la interacci{\'o}n humano--computadora y la integraci{\'o}n de subsistemas gr{\'a}ficos en plataformas de c{\'o}mputo.},
		type=main,
		symbol={🎞️},
		user1={Conexi{\'o}n de video},
		user2={Transmisi{\'o}n gr{\'a}fica},
		user3={Visualizaci{\'o}n de datos},
		see={displayport,hdmi,vga,dvi,visualizacion}
	}
	
	% ===================== Interferencia electromagn{\'e}tica =====================
	\newglossaryentry{interferenciaelectromagnetica}
	{
		name={Interferencia electromagn{\'e}tica},
		text={interferencia electromagn{\'e}tica},
		sort={interferencia electromagnetica},
		plural={interferencias electromagn{\'e}ticas},
		first={interferencia electromagn{\'e}tica},
		firstplural={interferencias electromagn{\'e}ticas},
		description={La \emph{interferencia electromagn{\'e}tica} (EMI) ocurre cuando una se{\~n}al el{\'e}ctrica o magn{\'e}tica no deseada afecta el funcionamiento de dispositivos o sistemas de comunicaci{\'o}n. En \gls{iot} puede degradar calidad de enlace, provocar p{\'e}rdida de \glspl{dato} o aumentar \gls{latencia}. Sus fuentes incluyen motores, fuentes de energ{\'i}a y otros equipos electr{\'o}nicos.},
		type=\glsdefaulttype,
		symbol={📡🚫},
		user1={EMI},
		user2={Ruido electromagn{\'e}tico},
		user3={Compatibilidad electromagn{\'e}tica},
		see={filtroruido,ethernet,conectividad}
	}
	
	% ===================== Internet =====================
	\newglossaryentry{internet}
	{
		name={Internet},
		text={Internet},
		sort={internet},
		plural={Internet},
		first={Internet},
		firstplural={Internet},
		description={\emph{Internet} es la red global de \glspl{red} interconectadas basada en protocolos como \gls{ip}, \gls{http}/\gls{https}, que permite el intercambio de \glspl{dato} y la provisi{\'o}n de servicios digitales. Es la base de las plataformas y ecosistemas \gls{iot}.},
		type=\glsdefaulttype,
		symbol={🌐},
		user1={Red global},
		user2={Protocolos TCP/IP},
		user3={Servicios web},
		see={ip,http,https,iot}
	}
	
	% ===================== Interoperabilidad =====================
	\newglossaryentry{interoperabilidad}
	{
		name={Interoperabilidad},
		text={interoperabilidad},
		sort={interoperabilidad},
		plural={interoperabilidades},
		first={interoperabilidad},
		firstplural={interoperabilidades},
		description={La \emph{interoperabilidad} es la capacidad de \glspl{sistema}, dispositivos, aplicaciones o componentes para comunicarse, intercambiar \glspl{dato} y utilizar la informaci{\'o}n compartida eficazmente sin intervenci{\'o}n adicional. Abarca aspectos t{\'e}cnicos, sem{\'a}nticos y organizativos. En \glspl{sibiot} permite que dispositivos heterog{\'e}neos, \glspl{protocolocomunicacion} y plataformas colaboren, reduciendo complejidad y mejorando eficiencia.},
		type=\glsdefaulttype,
		symbol={🔗},
		user1={Compatibilidad t{\'e}cnica},
		user2={Alineaci{\'o}n sem{\'a}ntica},
		user3={Est{\'a}ndares y APIs},
		see={integracion,compatibilidad,api}
	}
	
	\newglossaryentry{interrupcion}{
		name={Interrupci{\'o}n},
		text={interrupci{\'o}n},
		sort={interrupcion},
		plural={interrupciones},
		first={interrupci{\'o}n},
		firstplural={interrupciones},
		description={Mecanismo por el cual un microcontrolador suspende temporalmente el flujo normal de ejecuci{\'o}n para atender un evento, ejecutando una \gls{isr} y luego reanudando el programa principal.},
		type=main,
		symbol={⚡},
		user1={Eventos as{\'i}ncronos},
		user2={Prioridad},
		user3={Tiempo real},
		see={isr,temporizador,evento}
	}
	
	% ===================== I/O =====================
	\newglossaryentry{io}
	{
		name={I/O},
		text={I/O},
		sort={io},
		plural={Is/Os},
		first={entrada/salida (\textit{Input/Output}, I/O)},
		firstplural={Is/Os},
		description={\emph{I/O} designa la comunicaci{\'o}n entre un \gls{sistemainformatico} y el entorno: lectura desde \glspl{sensor}, escritura hacia \glspl{actuador}, o intercambio con perif{\'e}ricos y redes. En \gls{iot} se implementa v{\'i}a \gls{gpio}, buses y \glspl{protocolocomunicacion}.},
		type=\glsdefaulttype,
		symbol={↔️},
		user1={Entrada/salida},
		user2={Perif{\'e}ricos},
		user3={Buses/puertos},
		see={gpio,protocolocomunicacion}
	}
	
	% ===================== IoT =====================
	\newglossaryentry{iot}
	{
		name={IoT},
		text={IoT},
		sort={iot},
		plural={IoT},
		first={Internet de las cosas (\textit{Internet of Things}, IoT)},
		firstplural={IoT},
		description={\emph{IoT} es el paradigma que interconecta objetos f{\'i}sicos mediante \glspl{red} digitales para recolectar, intercambiar y procesar \glspl{dato}. Los \glspl{dispositivoconectado} con \glspl{sensor}, \glspl{actuador} y m{\'o}dulos de comunicaci{\'o}n forman \glspl{sibiot} que habilitan \gls{monitorizacion}, automatizaci{\'o}n y \gls{tomadecisiones} en \gls{tiemporeal}.},
		type=\glsdefaulttype,
		symbol={📡},
		user1={Internet de las Cosas},
		user2={Sistemas ciberf{\'i}sicos},
		user3={Conectividad de objetos},
		see={ecosistemaiot,edgecomputing,lpwan}
	}
	
	\newglossaryentry{iotstack}{
		name={IoT Stack},
		text={IoT Stack},
		sort={iotstack},
		plural={pilas IoT},
		first={pila de arquitectura de \gls{iot} (\textit{IoT Stack})},
		firstplural={pilas de arquitectura del \gls{iot} o \textit{IoT Stack}},
		description={Modelo estructural que define los niveles funcionales del ecosistema del \gls{iot}, describiendo la \gls{interaccion} entre \glspl{dispositivo} f{\'i}sicos, redes de comunicaci{\'o}n, \gls{procesamientodatos} y \glspl{aplicacion}. La \emph{IoT Stack} organiza los \glspl{componente} en \glspl{capa}, generalmente: \gls{capapercepcion} (sensado y actuaci{\'o}n), \gls{capared} (transmisi{\'o}n de datos), procesamiento (almacenamiento y an{\'a}lisis), y aplicaci{\'o}n (\glspl{servicio} inteligentes). Este enfoque facilita la \gls{interoperabilidad}, la \gls{modularidad} y la \gls{escalabilidad} de los \glspl{sibiot}. Adem{\'a}s, integra protocolos, est{\'a}ndares y \glspl{servicio} en la \gls{nube} que soportan la comunicaci{\'o}n y la \gls{tomadecisiones} autom{\'a}tica en \glspl{sistemadistribuido}.},
		type=main,
		user1={Modelo estructural de capas para sistemas IoT},
		user2={Organizaci{\'o}n de niveles: percepci{\'o}n, red, procesamiento y aplicaci{\'o}n},
		user3={Facilita interoperabilidad, modularidad y escalabilidad},
		user4={Base conceptual para arquitecturas de SIbIoT},
		see={arquitectura,capa,osi,iot,sibiot,protocolo},
		symbol={🏗️}
	}
	
	% ===================== IP =====================
	\newglossaryentry{ip}
	{
		name={IP},
		text={IP},
		sort={ip},
		plural={IP},
		first={protocolo de Internet (\textit{Internet Protocol}, IP)},
		firstplural={IP},
		description={\emph{IP} define el direccionamiento l{\'o}gico y el encaminamiento de paquetes en redes por conmutaci{\'o}n de paquetes. Versiones comunes: IPv4 e IPv6. En \gls{iot} permite identificar y comunicar nodos distribuidos en \glspl{lan} y en \gls{internet}.},
		type=main,
		symbol={🧭},
		user1={Direccionamiento},
		user2={Ruteo},
		user3={Capa de red},
		see={ethernet,lan,http}
	}
	
	\newglossaryentry{it}{
		name={IT},
		text={IT},
		sort={it},
		plural={IT},
		first={tecnolog{\'i}a de la Informaci{\'o}n (\textit{Information Technology}, IT)},
		firstplural={IT},
		description={Sigla de \emph{Information Technology}. Conjunto de tecnolog{\'\i}as orientadas a gesti{\'o}n de informaci{\'o}n y servicios digitales (sistemas empresariales, bases de datos, aplicaciones y redes corporativas). En arquitecturas modernas, la integraci{\'o}n IT--\gls{ot} habilita que datos y decisiones del proceso f{\'\i}sico se articulen con indicadores, auditor{\'\i}a y gobierno del \gls{sibiot}.},
		type=main,
		symbol={🖥},
		user1={Sistemas empresariales},
		user2={Datos y servicios},
		user3={Gobierno y auditor{\'\i}a},
		see={ot,governanza,kpi}
	}
	
	% ===================== Java =====================
	\newglossaryentry{java}
	{
		name={Java},
		text={Java},
		sort={java},
		plural={Java},
		first={Java},
		firstplural={Java},
		description={\emph{Java} es un lenguaje orientado a objetos y una plataforma que corre sobre la JVM (\textit{write once, run anywhere}). Ofrece tipado est{\'a}tico, gesti{\'o}n de memoria y un ecosistema amplio para aplicaciones de escritorio, web, m{\'o}viles y empresariales. En \gls{iot} se usa para \glspl{aplicacionweb}, \glspl{serviciodistribuido} y pasarelas.},
		type=\glsdefaulttype,
		symbol={☕},
		user1={JVM},
		user2={Servicios web},
		user3={Backend escalable},
		see={api,backend,arquitecturadistribuida}
	}
	
	\newglossaryentry{javascript}{
		name={JavaScript},
		text={JavaScript},
		sort={javascript},
		plural={JavaScript},
		first={JavaScript},
		firstplural={JavaScript},
		description={\emph{JavaScript} es un \gls{lenguajeprogramacion} interpretado, orientado a objetos y ampliamente utilizado en el desarrollo de \glspl{aplicacionweb}. Permite dotar de interactividad, dinamismo y l{\'o}gica de control a documentos \gls{html} y hojas de estilo \gls{css}. Es soportado por todos los navegadores modernos mediante motores de ejecuci{\'o}n como V8 o SpiderMonkey. En el contexto de los \glspl{sibiot}, \emph{JavaScript} se aplica para construir interfaces interactivas, gestionar \glspl{api} \gls{rest}, procesar datos en \gls{tiemporeal} y habilitar la comunicaci{\'o}n con plataformas en la \gls{nube} o en el \gls{edgecomputing}.},
		type=main,
		symbol={📜},
		user1={Lenguaje de programaci{\'o}n interpretado en navegadores},
		user2={Soporte a interactividad en aplicaciones web},
		user3={Uso en APIs, IoT y entornos en tiempo real},
		see={java,html,css,aplicacionweb,lenguajeprogramacion,sibiot}
	}
	
	\newglossaryentry{jdbc}{
		name={Java Database Connectivity},
		text={Java Database Connectivity},
		plural={interfaces Java Database Connectivity},
		first={conectividad de bases de datos en Java (\textit{Java Database Connectivity}, JDBC)},
		firstplural={interfaces Java Database Connectivity utilizadas para el acceso estructurado a bases de datos relacionales},
		description={Interfaz de programaci{\'o}n que define un conjunto de clases y m{\'e}todos para ejecutar consultas SQL, gestionar conexiones y manipular resultados desde aplicaciones Java},
		sort={java database connectivity},
		type=main,
		symbol={☕},
		see={java,basedatosrelacional}
	}
	
	\newglossaryentry{jitter}{
		name={Jitter},
		text={jitter},
		sort={jitter},
		plural={jitters},
		first={jitter},
		firstplural={jitters},
		description={Variaci{\'o}n del retardo (\gls{latencia}) en la entrega de mensajes o se{\~n}ales en una comunicaci{\'o}n. En \glspl{cps}, el jitter puede degradar la \gls{temporizacion} y afectar control y seguridad, porque altera el orden o el instante efectivo con que llegan mediciones y comandos.},
		type=main,
		symbol={📶},
		user1={Variaci{\'o}n de retardo},
		user2={Impacto en control},
		user3={Calidad temporal},
		see={latencia,temporizacion,tiemporeal}
	}
	
	\newglossaryentry{json}{
		name={JSON},
		text={JSON},
		sort={json},
		plural={JSON},
		first={notaci{\'o}n de objetos de JavaScript (\textit{JavaScript Object Notation}, JSON)},
		firstplural={JSON},
		description={Formato ligero de intercambio de datos basado en texto, derivado de la sintaxis de objetos de JavaScript, pero independiente de este lenguaje. Es f{\'a}cil de leer y escribir por humanos y sencillo de interpretar y generar por m{\'a}quinas. JSON utiliza una estructura de pares clave--valor y listas ordenadas, lo que lo convierte en una alternativa eficiente a \gls{xml}. Es ampliamente usado en \glspl{api}, servicios \gls{restful}, configuraciones y sistemas \gls{iot}, donde la simplicidad y el bajo costo de procesamiento son esenciales.},
		type=main,
		symbol={🗄️},
		user1={Formato de datos ligero basado en texto},
		user2={Estructura de pares clave--valor y listas},
		user3={Amplio uso en APIs, REST y sistemas IoT},
		see={xml,api,restful,iot}
	}
	
	% ===================== Kanban =====================
	\newglossaryentry{kanban}
	{
		name={Kanban},
		text={Kanban},
		sort={kanban},
		plural={Kanban},
		first={Kanban},
		firstplural={Kanban},
		description={\emph{Kanban} es un m{\'e}todo de gesti{\'o}n visual que emplea tableros y columnas para representar el estado de las tareas, limitar trabajo en curso (WIP) e identificar cuellos de botella. En \gls{iot} ayuda a coordinar hardware, firmware y software en entregas continuas.},
		type=\glsdefaulttype,
		symbol={🗂️},
		user1={Flujo tirado},
		user2={WIP limitado},
		user3={Mejora continua},
		see={equipoagil,desarrolloagilsoftware}
	}
	
	\newglossaryentry{kitdesarrollo}
	{
		name={Kit de desarrollo},
		text={kit de desarrollo},
		sort={kit de desarrollo},
		plural={kits de desarrollo},
		first={kit de desarrollo},
		firstplural={kits de desarrollo},
		description={Un kit de desarrollo es un conjunto de herramientas, bibliotecas, documentaci{\'o}n y recursos que permiten a los programadores crear, probar e implementar aplicaciones o sistemas. En el contexto de los \glspl{sibiot}, los kits de desarrollo incluyen \glspl{microcontrolador}, entornos de programaci{\'o}n, \glspl{sensor}, \glspl{actuador} y ejemplos de c{\'o}digo que facilitan la construcci{\'o}n de prototipos y la integraci{\'o}n de soluciones IoT de manera eficiente.},
		type=\glsdefaulttype,
		symbol={🧰},
		user1={SDK},
		user2={Kit de herramientas},
		user3={Paquete de desarrollo},
		see={microcontrolador}
	}
	
	\newglossaryentry{kpi}{
		name={KPI},
		text={KPI},
		plural={KPIs},
		first={Indicador Clave de Desempe{\~n}o (\textit{Key Performance Indicator}, KPI)},
		firstplural={KPIs},
		see={rendimiento},
		sort={kpi},
		description={Indicador cuantitativo utilizado para medir el desempe{\~n}o de un \gls{proceso}, sistema u objetivo estrat{\'e}gico. En \gls{sibiot}, los \emph{KPIs} sintetizan resultados consolidados derivados del \gls{analisis} de \gls{dato} \gls{iot} y de \glspl{sistemaempresarial}, permitiendo evaluar \gls{eficienciaoperativa}, cumplimiento de umbrales, impacto de \glspl{decision} y evoluci{\'o}n temporal del sistema}
	}
	
	\newglossaryentry{kubernetes}
	{
		name={Kubernetes},
		text={Kubernetes},
		sort={kubernetes},
		plural={plataformas Kubernetes},
		first={Kubernetes},
		firstplural={plataformas Kubernetes},
		description={Sistema de orquestaci{\'o}n de \glspl{contenedor} para despliegue, escalado y recuperaci{\'o}n de servicios, usado para administrar aplicaciones distribuidas en entornos de nube y, cuando aplica, en borde.},
		type=\glsdefaulttype,
		symbol={⚙️},
		user1={Orquestaci{\'o}n},
		user2={Escalado},
		user3={Recuperaci{\'o}n autom{\'a}tica},
		see={contenedor,cloudcomputing}
	}
	
	\newglossaryentry{lan}
	{
		name={LAN},
		text={LAN},
		sort={lan},
		plural={LANs},
		first={red de {\'a}rea local (\textit{Local Area Network}, LAN)},
		firstplural={LANs},
		description={Una LAN es un tipo de red de comunicaci{\'o}n que interconecta dispositivos dentro de un {\'a}rea geogr{\'a}fica limitada, como una vivienda, edificio o campus. Las redes LAN permiten compartir recursos como archivos, impresoras y conexi{\'o}n a Internet de forma eficiente. En entornos \gls{iot}, las LAN se utilizan para conectar nodos y \textit{gateways} que operan en redes dom{\'e}sticas, industriales o institucionales, ofreciendo baja latencia y alta velocidad.},
		type=main,
		symbol={🖧},
		user1={Red de {\'a}rea local},
		user2={Local Area Network},
		user3={Red privada interna},
		see={sistemacomunicacion}
	}
	
	\newglossaryentry{latencia}
	{
		name={Latencia},
		text={latencia},
		sort={latencia},
		plural={latencias},
		first={latencia},
		firstplural={latencias},
		description={La latencia es el tiempo que transcurre desde que se inicia una solicitud de datos o mensaje hasta que comienza a recibirse la respuesta. En redes de computadoras se mide com{\'u}nmente en milisegundos (ms), y su valor influye directamente en el desempe{\~n}o percibido por los usuarios y sistemas. En los \glspl{sibiot}, la latencia baja es esencial para aplicaciones en tiempo real como automatizaci{\'o}n industrial, salud conectada o veh{\'i}culos aut{\'o}nomos.},
		type=\glsdefaulttype,
		symbol={⏱},
		user1={Tiempo de respuesta},
		user2={Retraso de propagaci{\'o}n},
		user3={Delay},
		see={efectividad}
	}
	
	\newglossaryentry{layout}{
		name={Layout},
		text={\textit{layout}},
		sort={layout},
		plural={\textit{layouts}},
		first={\textit{layout}},
		firstplural={\textit{layouts}},
		description={Estructura que define la disposici{\'o}n de componentes visuales en una interfaz de usuario. En Android, un \emph{layout} describe jerarqu{\'i}as de vistas y su posicionamiento.},
		type=main,
		symbol={📐},
		user1={Disposici{\'o}n UI},
		user2={Jerarqu{\'i}a de vistas},
		user3={Android},
		see={drawable,values,assets}
	}
	
	\newglossaryentry{lcd}
	{
		name={LCD},
		text={LCD},
		sort={lcd},
		plural={LCDs},
		first={pantalla de cristal l{\'i}quido (\textit{Liquid Crystal Display}, LCD)},
		firstplural={LCDs},
		description={Un LCD es un dispositivo de visualizaci{\'o}n que utiliza propiedades de modulaci{\'o}n de la luz de cristales l{\'i}quidos para mostrar informaci{\'o}n en formato alfanum{\'e}rico o gr{\'a}fico. Son ampliamente usados en sistemas embebidos, paneles de control, dispositivos port{\'a}tiles y aplicaciones \gls{iot}, gracias a su bajo consumo de energ{\'i}a y buena visibilidad.},
		type=main,
		symbol={📺},
		user1={Pantalla de cristal l{\'i}quido},
		user2={Display de cristal l{\'i}quido},
		user3={Visualizador LCD},
		see={interfaz}
	}
	
	\newglossaryentry{lengscript}
	{
		name={Lenguaje scripting},
		text={lenguaje scripting},
		sort={lenguaje scripting},
		plural={lenguajes scripting},
		first={lenguaje scripting},
		firstplural={lenguajes scripting},
		description={Un lenguaje \textit{scripting} es una familia de lenguajes de programaci{\'o}n que permiten desarrollar r{\'a}pidamente soluciones a necesidades comunes. Los m{\'a}s extendidos se utilizan en el desarrollo web, tanto del lado del cliente (navegador) como del lado del servidor. Tambi{\'e}n existen lenguajes scripting aplicados en automatizaci{\'o}n de tareas, administraci{\'o}n de sistemas y desarrollo de prototipos en entornos \gls{iot}.},
		type=\glsdefaulttype,
		symbol={💻},
		user1={Lenguaje interpretado},
		user2={Script},
		user3={Lenguaje din{\'a}mico},
		see={programacion,lenguajeprogramacion}
	}
	
	\newglossaryentry{lenguajenatural}{
		name={Lenguaje natural},
		text={lenguaje natural},
		sort={lenguaje natural},
		plural={lenguajes naturales},
		first={lenguaje natural},
		firstplural={lenguajes naturales},
		description={Lengua utilizada de forma cotidiana por los seres humanos para comunicarse, como el espa{\~n}ol, el ingl{\'e}s o el franc{\'e}s. A diferencia de los lenguajes formales o de programaci{\'o}n, el \emph{lenguaje natural} presenta ambig{\"u}edades, redundancias y flexibilidad sem{\'a}ntica. En el contexto de la computaci{\'o}n y los \glspl{sibiot}, los lenguajes naturales son analizados mediante \gls{ia} y t{\'e}cnicas de procesamiento de lenguaje natural (NLP) para interpretar instrucciones humanas, traducirlas a \glspl{algoritmo} o interactuar con usuarios de forma m{\'a}s intuitiva.},
		type=main,
		symbol={🗣️},
		user1={Lenguas humanas},
		user2={Procesamiento autom{\'a}tico del lenguaje},
		user3={Interfaz hombre--m{\'a}quina},
		see={algoritmo,ia,sistemainteligente}
	}
	
	% ===================== Lenguaje de programaci{\'o}n =====================
	\newglossaryentry{lenguajeprogramacion}{
		name={Lenguaje de programaci{\'o}n},
		text={lenguaje de programaci{\'o}n},
		sort={lenguaje de programacion},
		plural={lenguajes de programaci{\'o}n},
		first={lenguaje de programaci{\'o}n},
		firstplural={lenguajes de programaci{\'o}n},
		description={Un \emph{lenguaje de programaci{\'o}n} es un sistema formal compuesto por reglas sint{\'a}cticas y sem{\'a}nticas que permite a los desarrolladores escribir \gls{codigofuente}. Facilita la comunicaci{\'o}n entre humanos y m{\'a}quinas, traduciendo instrucciones a operaciones ejecutables en un \gls{microprocesador} o \gls{microcontrolador}. Los lenguajes de programaci{\'o}n pueden clasificarse en lenguajes de bajo nivel (como ensamblador) y de alto nivel (como C, Java, Python). En \glspl{sibiot}, la elecci{\'o}n del lenguaje determina la eficiencia, portabilidad y facilidad de integraci{\'o}n de los sistemas desarrollados.},
		type=\glsdefaulttype,
		symbol={💬},
		user1={Sistema formal con sintaxis y sem{\'a}ntica},
		user2={Traducci{\'o}n de instrucciones humanas a m{\'a}quina},
		user3={Base del desarrollo de software IoT},
		see={codigofuente,algoritmo,microprocesador,microcontrolador}
	}
	
	\newglossaryentry{limpieza}{
		name={Limpieza},
		text={limpieza},
		sort={limpieza},
		plural={limpiezas},
		first={limpieza},
		firstplural={limpiezas},
		description={La limpieza de datos es una fase esencial en la gesti{\'o}n de calidad que busca depurar los conjuntos de datos mediante la detecci{\'o}n y correcci{\'o}n de errores, valores at{\'i}picos o ausentes. Este proceso mejora la precisi{\'o}n y confiabilidad de los an{\'a}lisis, permitiendo que los modelos de \gls{analisisdescriptivo}, \gls{analisispredictivo} y \gls{analisisprescriptivo} produzcan resultados m{\'a}s certeros. En entornos \gls{iot}, la limpieza resulta cr{\'i}tica para filtrar lecturas defectuosas de \glspl{sensor} y mantener la integridad del flujo de informaci{\'o}n.},
		type=\glsdefaulttype,
		symbol={🧹},
		user1={Data cleaning},
		user2={Depuraci{\'o}n de registros err{\'o}neos o inconsistentes},
		user3={Proceso para mejorar la calidad y precisi{\'o}n de los datos},
		see={validacion,enriquecimiento,calidaddatos,iot}
	}
	
	\newglossaryentry{logica}{
		name={L{\'o}gica},
		text={l{\'o}gica},
		sort={logica},
		plural={l{\'o}gicas},
		first={l{\'o}gica},
		firstplural={l{\'o}gicas},
		description={Conjunto de reglas y principios formales que permiten establecer relaciones v{\'a}lidas entre proposiciones, estructuras o transformaciones, asegurando inferencias correctas y ausencia de contradicciones. En ingenier{\'\i}a de software, la l{\'o}gica sustenta la especificaci{\'o}n formal, el dise{\~n}o de reglas de negocio, la verificaci{\'o}n de propiedades y la construcci{\'o}n de transformaciones preservadoras de significado. En IoT, tambi{\'e}n aparece en l{\'o}gica de control (reglas, condiciones, aut{\'o}matas) y en l{\'o}gica de procesamiento (validaciones, filtros, correlaciones y detecci{\'o}n de anomal{\'\i}as).},
		type=main,
		symbol={∴},
		user1={Reglas formales de validez},
		user2={Soporte a especificaci{\'o}n y verificaci{\'o}n},
		user3={Base de control y procesamiento},
		see={consistencia,verificacion,controlinteligente,procesamiento,regla,restriccion}
	}
	
	\newglossaryentry{logicadifusa}
	{
		name={L{\'o}gica difusa},
		text={l{\'o}gica difusa},
		sort={logica difusa},
		plural={l{\'o}gicas difusas},
		first={l{\'o}gica difusa},
		firstplural={l{\'o}gicas difusas},
		description={La \emph{l{\'o}gica difusa} es un paradigma matem{\'a}tico que extiende la l{\'o}gica cl{\'a}sica binaria al permitir valores de verdad intermedios entre "verdadero" y "falso". Se aplica en sistemas de \gls{controlinteligente}, \gls{automatizacion}, \glspl{sibiot} y procesos de \gls{tomadecisiones} donde existe incertidumbre, imprecisi{\'o}n o variabilidad en los \glspl{dato}. Permite modelar razonamientos aproximados similares al pensamiento humano, siendo clave en aplicaciones de \gls{aprendizajeautomatico} e inteligencia artificial.},
		type=\glsdefaulttype,
		symbol={➰},
		user1={Fuzzy logic},
		user2={L{\'o}gica borrosa},
		user3={L{\'o}gica aproximada},
		see={controlinteligente}
	}
	
	\newglossaryentry{logicanegocios}
	{
		name={L{\'o}gica de negocios},
		text={l{\'o}gica de negocios},
		sort={logica de negocios},
		plural={l{\'o}gicas de negocios},
		first={l{\'o}gica de negocios},
		firstplural={l{\'o}gicas de negocios},
		description={La \emph{l{\'o}gica de negocios} corresponde a la \gls{logicaprogramacion} de un sistema que codifica las reglas del negocio en el mundo real, determinando c{\'o}mo la informaci{\'o}n puede ser creada, almacenada, actualizada, eliminada e incluso compartida.},
		type=\glsdefaulttype,
		symbol={🏢},
		user1={Business logic},
		user2={Reglas de negocio},
		user3={Procesos de negocio},
		see={logicaprogramacion}
	}
	
	\newglossaryentry{logicapresentacion}
	{
		name={L{\'o}gica de presentaci{\'o}n},
		text={l{\'o}gica de presentaci{\'o}n},
		sort={logica de presentacion},
		plural={l{\'o}gicas de presentaci{\'o}n},
		first={l{\'o}gica de presentaci{\'o}n},
		firstplural={l{\'o}gicas de presentaci{\'o}n},
		description={La \emph{l{\'o}gica de presentaci{\'o}n} es la parte de la programaci{\'o}n dise{\~n}ada espec{\'i}ficamente para las pantallas de las aplicaciones inform{\'a}ticas. Esta l{\'o}gica gestiona la experiencia de navegaci{\'o}n del usuario, comunicando la informaci{\'o}n y capturando datos. Tambi{\'e}n se conoce como Interfaz Gr{\'a}fica de Usuario o en ingl{\'e}s GUI (\textit{Graphic User Interface}).},
		type=\glsdefaulttype,
		symbol={🖥},
		user1={Presentaci{\'o}n de datos},
		user2={Capa de interfaz},
		user3={GUI},
		see={interfaz}
	}
	
	\newglossaryentry{logicaprogramacion}
	{
		name={L{\'o}gica de programaci{\'o}n},
		text={l{\'o}gica de programaci{\'o}n},
		sort={logica de programacion},
		plural={l{\'o}gicas de programaci{\'o}n},
		first={l{\'o}gica de programaci{\'o}n},
		firstplural={l{\'o}gicas de programaci{\'o}n},
		description={La \emph{l{\'o}gica de programaci{\'o}n} es la organizaci{\'o}n y planificaci{\'o}n de instrucciones en un \gls{algoritmo}, escritas generalmente en un lenguaje de programaci{\'o}n ejecutable por una computadora, con el objetivo de resolver un problema. Se conoce tambi{\'e}n como \gls{logicanegocios} cuando est{\'a} vinculada a reglas de negocio.},
		type=\glsdefaulttype,
		symbol={💻},
		user1={Program logic},
		user2={Estructura algor{\'i}tmica},
		user3={Dise{\~n}o l{\'o}gico},
		see={algoritmo}
	}
	
	\newglossaryentry{logistica}{
		name={Log{\'i}stica},
		text={log{\'i}stica},
		sort={logistica},
		plural={log{\'i}sticas},
		first={log{\'i}stica},
		firstplural={log{\'i}sticas},
		description={Disciplina y conjunto de procesos encargados de planificar, implementar y controlar el flujo eficiente de bienes, servicios e informaci{\'o}n desde el punto de origen hasta el de consumo, con el fin de satisfacer los requerimientos del cliente. La \emph{log{\'i}stica} abarca transporte, almacenamiento, distribuci{\'o}n, gesti{\'o}n de inventarios y trazabilidad. En el contexto de los \glspl{sibiot}, la log{\'i}stica se potencia mediante \glspl{sensor}, \glspl{gnss}, \glspl{serviciodistribuido} y plataformas en la \gls{nube}, lo que permite optimizar rutas, reducir costes, mejorar la \gls{eficienciaoperativa} y contribuir a la \gls{sostenibilidad} de la cadena de suministro.},
		type=main,
		symbol={🚚},
		user1={Planificaci{\'o}n y gesti{\'o}n de flujos},
		user2={Optimizar transporte, inventarios y trazabilidad},
		user3={Aplicaciones en SIbIoT e Industria 4.0},
		see={cadenalogistica,eficienciaoperativa,sostenibilidad,gnss,sibiot}
	}
	
	\newglossaryentry{lora}
	{
		name={LoRa},
		text={LoRa},
		sort={lora},
		plural={LoRa},
		first={largo alcance (\textit{Long Range}, LoRa)},
		firstplural={LoRa},
		description={\emph{\gls{lora}} es una tecnolog{\'i}a de modulaci{\'o}n inal{\'a}mbrica basada en espectro ensanchado por chirp (\textit{Chirp Spread Spectrum}) que permite la transmisi{\'o}n de datos a larga distancia con bajo consumo energ{\'e}tico. Se utiliza en redes \gls{lpwan} y resulta especialmente adecuada para aplicaciones \gls{iot} que requieren \gls{conectividad} entre \glspl{dispositivo} distribuidos geogr{\'a}ficamente, como agricultura de \gls{precision}, monitorizaci{\'o}n ambiental y \glspl{ciudadinteligente}.},
		type=\glsdefaulttype,
		symbol={📡},
		user1={Long Range},
		user2={Modulaci{\'o}n chirp},
		user3={Transmisi{\'o}n IoT},
		see={lpwan}
	}
	
	\newglossaryentry{lorawan}
	{
		name={LoRaWAN},
		text={LoRaWAN},
		sort={lorawan},
		plural={LoRaWANs},
		first={red de {\'a}rea amplia de largo alcance (\textit{Long Range Wide Area Network}, LoRaWAN)},
		firstplural={LoRaWANs},
		description={\emph{\gls{lorawan}} es un protocolo de comunicaci{\'o}n de largo alcance y bajo consumo energ{\'e}tico basado en la tecnolog{\'i}a \gls{lora}. Dise{\~n}ado para \glspl{rediot} de gran escala, permite la transmisi{\'o}n de peque{\~n}os paquetes de datos a largas distancias con topolog{\'i}as tipo estrella y soporte para \glspl{dispositivo} alimentados por bater{\'i}a durante largos per{\'i}odos. Es ideal para aplicaciones como monitorizaci{\'o}n ambiental, agricultura de \gls{precision}, \gls{gestion} de recursos y \glspl{ciudadinteligente}.},
		type=\glsdefaulttype,
		symbol={🌐},
		user1={Protocolo LoRa},
		user2={Wide Area Network},
		user3={Est{\'a}ndar IoT},
		see={lora}
	}
	
	\newglossaryentry{lpwan}
	{
		name={LPWAN},
		text={LPWAN},
		sort={lpwan},
		plural={LPWANs},
		first={red de {\'a}rea amplia y bajo consumo energ{\'e}tico (\textit{Low Power Wide Area Network}, LPWAN)},
		firstplural={LPWANs},
		description={\emph{\gls{lpwan}} es una categor{\'i}a de tecnolog{\'i}as de comunicaci{\'o}n inal{\'a}mbrica dise{\~n}adas para ofrecer \gls{conectividad} a larga distancia con bajo consumo energ{\'e}tico. Estas redes permiten la transmisi{\'o}n de peque{\~n}os vol{\'u}menes de datos a velocidades reducidas, siendo ideales para \glspl{dispositivoiot} alimentados por bater{\'i}a. Tecnolog{\'i}as representativas incluyen \gls{lora}, \gls{lorawan}, Sigfox y \gls{nbiot}.},
		type=\glsdefaulttype,
		symbol={🔋},
		user1={Low Power WAN},
		user2={Red IoT de bajo consumo},
		user3={Tecnolog{\'i}a LPWA},
		see={lora}
	}
	
	\newglossaryentry{lte}
	{
		name={LTE},
		text={LTE},
		sort={lte},
		plural={LTE},
		first={evoluci{\'o}n a largo plazo (\textit{Long Term Evolution}, LTE)},
		firstplural={LTE},
		description={\emph{\gls{lte}} es un est{\'a}ndar de comunicaci{\'o}n inal{\'a}mbrica de alta velocidad para tel{\'e}fonos m{\'o}viles y \glspl{dispositivo} de datos. Surge de la evoluci{\'o}n de la tecnolog{\'i}a 3G hacia la cuarta generaci{\'o}n (4G), ofreciendo mayores tasas de transferencia, menor \gls{latencia} y mayor eficiencia espectral. Actualmente se encuentra ampliamente desplegado en \glspl{sistemacomunicacion} m{\'o}viles a nivel mundial.},
		type=\glsdefaulttype,
		symbol={📶},
		user1={4G},
		user2={Long Term Evolution},
		user3={Red m{\'o}vil de alta velocidad},
		see={ltem}
	}
	
	\newglossaryentry{ltem}
	{
		name={LTE-M},
		text={LTE-M},
		sort={ltem},
		plural={LTE-M},
		first={evoluci{\'o}n a largo plazo para m{\'a}quinas (\textit{Long Term Evolution for Machines}, LTE-M)},
		firstplural={LTE-M},
		description={\emph{\gls{ltem}} es una tecnolog{\'i}a de comunicaci{\'o}n de tipo \gls{lpwan} que aprovecha la infraestructura \gls{lte} existente para conectar \glspl{dispositivoiot} con requisitos de baja \gls{latencia} y movilidad. Permite mayores tasas de transferencia y soporte para voz sobre IP (VoLTE), siendo adecuada para aplicaciones industriales, de transporte y entornos m{\'o}viles.},
		type=\glsdefaulttype,
		symbol={📲},
		user1={LTE Cat-M1},
		user2={LTE for Machines},
		user3={IoT m{\'o}vil},
		see={lte}
	}
	
	\newglossaryentry{luminosidad}{
		name={Luminosidad},
		text={luminosidad},
		sort={luminosidad},
		plural={luminosidades},
		first={luminosidad},
		firstplural={luminosidades},
		description={Cantidad de luz percibida o medida en un entorno determinado, normalmente expresada en lux (lx). La \emph{luminosidad} es una \gls{variablefisica} fundamental para evaluar condiciones de iluminaci{\'o}n, confort visual y eficiencia energ{\'e}tica. En el contexto de los \glspl{sibiot}, se mide mediante \glspl{sensor} de luz o fotodiodos, permitiendo regular autom{\'a}ticamente sistemas de iluminaci{\'o}n, mejorar la eficiencia de consumo, optimizar entornos laborales y garantizar la seguridad en espacios interiores y exteriores. Adem{\'a}s, la luminosidad se integra con otras variables como temperatura o movimiento para habilitar aplicaciones avanzadas en dom{\'o}tica, agricultura de precisi{\'o}n, veh{\'i}culos aut{\'o}nomos y ciudades inteligentes.},
		type=main,
		symbol={💡},
		user1={Medici{\'o}n de luz ambiental},
		user2={Par{\'a}metro de confort y eficiencia},
		user3={Dato clave en aplicaciones IoT},
		see={variablefisica,sensor,domotica,sibiot}
	}
	
	% ===================== M2M =====================
	\newglossaryentry{m2m}{
		name={M2M},
		text={M2M},
		sort={m2m},
		plural={comunicaciones M2M},
		first={m{\'a}quina a m{\'a}quina, (\textit{Machine to Machine}, M2M)},
		firstplural={comunicaciones M2M},
		description={\emph{M2M} es un modelo de comunicaci{\'o}n que permite el intercambio directo de \glspl{dato} entre \glspl{dispositivo} sin intervenci{\'o}n humana. Constituye la base de procesos autom{\'a}ticos como la monitorizaci{\'o}n remota, la telemetr{\'i}a, el control industrial y la gesti{\'o}n de activos en sistemas \gls{iot}. Se implementa mediante redes celulares, \gls{lpwan} u otras interfaces de \gls{protocolocomunicacion}, habilitando la interacci{\'o}n eficiente entre \glspl{sensor}, \glspl{actuador} y \glspl{plataforma} centrales.},
		type=main,
		symbol={🤖},
		user1={Machine to Machine},
		user2={Comunicaci{\'o}n directa entre dispositivos},
		user3={Base de IoT automatizado},
		user4={Soporte en redes celulares y LPWAN},
		see={iot,lpwan,sensor,actuador,protocolocomunicacion}
	}
	
	\newglossaryentry{macroeconomia}{
		name={Macroeconomia},
		text={macroeconom{\'i}a},
		plural={macroeconom{\'i}as},
		first={macroeconom{\'i}a},
		firstplural={macroeconom{\'i}as},
		description={Rama de la econom{\'i}a que estudia el comportamiento agregado de un sistema econ{\'o}mico a gran escala, considerando variables globales como el producto interno bruto, la inflaci{\'o}n, el desempleo, el consumo, la inversi{\'o}n y la estabilidad financiera. La macroeconom{\'i}a proporciona el marco te{\'o}rico y metodol{\'o}gico para analizar tendencias econ{\'o}micas, formular pol{\'i}ticas p{\'u}blicas y apoyar la toma de decisiones estrat{\'e}gicas en sistemas productivos y organizacionales},
		type=main,
		see={variablemacroeconomica}
	}
	
	% ===================== MAN =====================
	\newglossaryentry{man}{
		name={MAN},
		text={MAN},
		sort={man},
		plural={MANs},
		first={red de {\'a}rea metropolitan (\textit{Metropolitan Area Network}, MAN)},
		firstplural={MANs},
		description={Una \emph{MAN} es un tipo de \gls{red} que abarca un {\'a}rea geogr{\'a}fica mayor que una \gls{lan} pero menor que una \gls{wan}, interconectando redes locales dentro de una ciudad o campus. Se implementa con tecnolog{\'i}as como fibra {\'o}ptica, enlaces inal{\'a}mbricos o redes metropolitanas \gls{ethernet}. En proyectos urbanos de \gls{iot}, las MAN son esenciales para conectar edificios, instituciones y \glspl{infraestructura} municipal, habilitando \glspl{red} inteligentes y ciudades conectadas.},
		type=\glsdefaulttype,
		symbol={🏙️},
		user1={Metropolitan Area Network},
		user2={Red urbana},
		user3={Conectividad a escala de ciudad},
		see={lan,wan,ethernet,infraestructura}
	}
	
	% ===================== Manifiesto Ágil =====================
	\newglossaryentry{manifiestoagil}{
		name={Manifiesto para el desarrollo {\'a}gil de software},
		text={manifiesto {\'a}gil},
		sort={manifiesto agil},
		plural={manifiestos {\'a}giles},
		first={Manifiesto {\'a}gil},
		firstplural={manifiestos {\'a}giles},
		description={El \emph{Manifiesto para el desarrollo {\'a}gil de software}, publicado en 2001, establece cuatro valores fundamentales y doce principios para mejorar eficiencia, flexibilidad y colaboraci{\'o}n en proyectos. Promueve entrega continua de valor, adaptaci{\'o}n al cambio, comunicaci{\'o}n fluida con el cliente y trabajo en equipo. En sistemas \gls{iot}, sus principios inspiran enfoques iterativos e incrementales para enfrentar complejidad t{\'e}cnica y organizativa en entornos distribuidos.},
		type=\glsdefaulttype,
		symbol={📜},
		user1={Agile Manifesto},
		user2={Valores y principios {\'a}giles},
		user3={Base de metodolog{\'i}as Agile},
		see={iniciativaagil,desarrolloagilsoftware,equipoagil,kanban}
	}
	
	\newglossaryentry{mantenibilidad}{
		name={Mantenibilidad},
		text={mantenibilidad},
		sort={mantenibilidad},
		plural={mantenibilidades},
		first={mantenibilidad},
		firstplural={mantenibilidades},
		description={Grado en que un producto o sistema de software puede ser modificado de forma eficaz y eficiente por los responsables del mantenimiento, con un riesgo controlado y un impacto acotado en la funcionalidad existente. Incluye atributos como la modularidad, la analizabilidad, la modificabilidad, la reusabilidad y la testabilidad (conforme a familias de est{\'a}ndares ISO/IEC 25010/25012). La mantenibilidad se ve favorecida por una arquitectura limpia, convenciones de c{\'o}digo, dise{\~n}o orientado a pruebas, automatizaci{\'o}n de \textit{build}, integraci{\'o}n/entrega continua y una documentaci{\'o}n pertinente. En el contexto de los \glspl{sibiot} y sistemas distribuidos, tambi{\'e}n depende de la observabilidad, el versionado de interfaces, la gesti{\'o}n de dependencias y la trazabilidad extremo a extremo.},
		type=main,
		user1={Facilidad de realizar correcciones y mejoras},
		user2={Capacidad de entender, analizar y localizar cambios},
		user3={Estructura modular con bajo acoplamiento y alta cohesi{\'o}n},
		user4={Soporte a pruebas y despliegues automatizados},
		see={calidadsoftware,modularidad,testabilidad,refactorizacion,documentacion},
		symbol={\ensuremath{\mathcal{M}}}
	}
	
	\newglossaryentry{mantenimiento}{
		name={Mantenimiento},
		text={mantenimiento},
		sort={mantenimiento},
		plural={mantenimientos},
		first={mantenimiento},
		firstplural={mantenimientos},
		description={Conjunto de actividades, t{\'e}cnicas y procedimientos destinados a preservar, restaurar o mejorar la funcionalidad de un sistema, equipo o componente, garantizando su disponibilidad, confiabilidad y seguridad a lo largo de su ciclo de vida. En ingenier{\'i}a de software y sistemas \gls{iot}, el \emph{mantenimiento} puede clasificarse en correctivo, preventivo, predictivo o evolutivo, seg{\'u}n el objetivo que persiga, y constituye una fase fundamental para la sostenibilidad y eficiencia operativa.},
		type=main,
		symbol={🛠️},
		user1={Acciones para preservar el funcionamiento de un sistema},
		user2={Incluye mantenimiento preventivo, correctivo, predictivo y evolutivo},
		user3={Fase esencial del ciclo de vida de sistemas e infraestructura},
		see={mantenimientopreventivo,mantenimientopredictivo,aseguramientocalidad}
	}
	
	\newglossaryentry{mantenimientopredictivo}{
		name={Mantenimiento predictivo},
		text={mantenimiento predictivo},
		sort={mantenimiento predictivo},
		plural={mantenimientos predictivos},
		first={mantenimiento predictivo},
		firstplural={mantenimientos predictivos},
		description={Estrategia de mantenimiento basada en el an{\'a}lisis continuo de datos y el uso de \glspl{modelopredictivo} para anticipar fallos o deterioros en equipos, m{\'a}quinas o sistemas antes de que ocurran. El \emph{mantenimiento predictivo} se apoya en \glspl{sensor}, algoritmos de \gls{aprendizajeautomatico} y t{\'e}cnicas de \gls{monitorizacion} en \gls{tiemporeal} para detectar patrones an{\'o}malos y estimar la vida {\'u}til restante de los componentes. En el contexto de los \glspl{sibiot}, esta pr{\'a}ctica contribuye a mejorar la \gls{eficienciaoperativa}, reducir costes de reparaci{\'o}n, minimizar tiempos de inactividad no planificada y garantizar la \gls{fiabilidad} y \gls{sostenibilidad} de procesos industriales, energ{\'e}ticos, de transporte y log{\'i}sticos.},
		type=main,
		symbol={🔧📊},
		user1={Anticipaci{\'o}n de fallos mediante datos},
		user2={Aplicaci{\'o}n de modelos predictivos e IA},
		user3={Reducci{\'o}n de costes y mejora de eficiencia},
		see={modelopredictivo,aprendizajeautomatico,monitorizacion,eficienciaoperativa,fiabilidad,sostenibilidad,sibiot}
	}
	
	\newglossaryentry{mantenimientopreventivo}{
		name={Mantenimiento preventivo},
		text={mantenimiento preventivo},
		sort={mantenimiento preventivo},
		plural={mantenimientos preventivos},
		first={mantenimiento preventivo},
		firstplural={mantenimientos preventivos},
		description={Estrategia de mantenimiento que consiste en realizar intervenciones programadas y peri{\'o}dicas en equipos, dispositivos o \glspl{componente} con el objetivo de reducir la probabilidad de fallos y alargar su vida {\'u}til. El \emph{mantenimiento preventivo} se basa en planes de revisi{\'o}n, limpieza, calibraci{\'o}n y sustituci{\'o}n de piezas seg{\'u}n calendarios o ciclos de uso definidos, sin necesidad de que se haya producido un fallo. En el contexto de los \glspl{sibiot}, este mantenimiento garantiza la \gls{fiabilidad} de \glspl{sensor}, \glspl{actuador} y nodos, contribuyendo a la \gls{eficienciaoperativa} y reduciendo los riesgos de inactividad no planificada, aunque puede implicar un mayor coste preventivo frente a otras estrategias como el mantenimiento predictivo.},
		type=main,
		symbol={🛠️},
		user1={Acciones peri{\'o}dicas de revisi{\'o}n},
		user2={Prevenci{\'o}n de fallos y prolongaci{\'o}n de vida {\'u}til},
		user3={Aplicado a sensores, actuadores y nodos en SIbIoT},
		see={mantenimientopredictivo,fiabilidad,eficienciaoperativa,sibiot}
	}
	
	% ===================== Manufactura =====================
	\newglossaryentry{manufactura}{
		name={Manufactura},
		text={manufactura},
		sort={manufactura},
		plural={manufacturas},
		first={manufactura},
		firstplural={manufacturas},
		description={La \emph{manufactura} es el proceso de transformaci{\'o}n de materias primas en productos terminados mediante m{\'a}quinas, herramientas, trabajo humano y tecnolog{\'i}as digitales. En el contexto de \gls{iot}, se ve potenciada por \glspl{sensor}, \glspl{red} y plataformas inteligentes que habilitan \gls{monitorizacion} en \gls{tiemporeal}, \gls{mantenimientopredictivo}, control automatizado y trazabilidad en las l{\'i}neas de producci{\'o}n, contribuyendo a la \gls{industry40}.},
		type=\glsdefaulttype,
		symbol={🏭},
		user1={Producci{\'o}n industrial},
		user2={Transformaci{\'o}n digital en la industria},
		user3={Base de Industry 4.0},
		see={industry40,mantenimientopredictivo,monitorizacion}
	}
	
	\newglossaryentry{mapeounidades}{
		name={Mapeo de unidades},
		text={mapeo de unidades},
		sort={mapeo unidades},
		plural={mapeos de unidades},
		first={mapeo de unidades},
		firstplural={mapeos de unidades},
		description={Proceso de estandarizaci{\'o}n mediante el cual se identifican, transforman y alinean diferentes unidades de medida provenientes de fuentes heterog{\'e}neas hacia un formato com{\'u}n y coherente. En sistemas \gls{iot}, el mapeo de unidades es esencial para garantizar exactitud, comparabilidad y consistencia de datos provenientes de m{\'u}ltiples sensores, evitando errores en el \gls{procesamiento} y en la \gls{tomadecisiones}. Puede implicar conversiones autom{\'a}ticas, validaciones y control de calidad de datos.},
		type=main,
		symbol={📏},
		user1={Normalizaci{\'o}n de magnitudes},
		user2={Conversi{\'o}n autom{\'a}tica de unidades},
		user3={Consistencia en an{\'a}lisis de datos},
		see={datoheterogeneo,exactitud,procesamiento}
	}
	
	\newglossaryentry{masa}
	{
		name={Masa},
		text={masa},
		sort={masa},
		description={Magnitud f{\'i}sica fundamental que cuantifica la cantidad de materia contenida en un cuerpo. 
			Se mantiene invariable independientemente de la ubicaci{\'o}n del objeto y constituye una propiedad intr{\'i}nseca relacionada con su inercia y su resistencia a los cambios de movimiento. 
			En el contexto industrial y de los \glspl{sibiot}, la masa se utiliza como variable de referencia para el \gls{pesaje} autom{\'a}tico y el control de procesos de producci{\'o}n, siendo medida indirectamente mediante \glspl{sensor} como las \glspl{celda} de carga y los \glspl{sensor} de fuerza.},
		plural={masas},
		symbol={⚖},
		first={Masa (m)},
		see=[Relacionado con:]{peso,fuerza,celda}
	}
	
	\newglossaryentry{material}{
		name={Material},
		text={material},
		sort={material},
		plural={materiales},
		first={material},
		firstplural={materiales},
		description={Sustancia o conjunto de sustancias con propiedades f{\'i}sicas y mec{\'a}nicas definidas. En sistemas IoT, los materiales condicionan resistencia, conductividad, durabilidad y compatibilidad con el entorno.},
		type=main,
		symbol={🧱},
		user1={Propiedades f{\'i}sicas},
		user2={Durabilidad},
		user3={Compatibilidad},
		see={elasticidad}
	}
	
	\newglossaryentry{matter}{
		name={Matter},
		text={Matter},
		sort={matter},
		plural={Matter},
		first={Matter},
		firstplural={Matter},
		description={Est{\'a}ndar de comunicaci{\'o}n para \gls{interoperabilidad} y seguridad multi-fabricante en hogares inteligentes, promovido por la \gls{csa}.},
		type=main,
		symbol={🔗},
		see={csa,interoperabilidad,provisionamiento,descubrimiento},
		user1={Modelo de datos com{\'u}n},
		user2={Compatibilidad multi-vendedor},
		user3={Onboarding seguro}
	}
	
	% ===================== MDE =====================
	\newglossaryentry{mde}{
		name={MDE},
		text={MDE},
		sort={mde},
		plural={MDE},
		first={ingenier{\'i}a dirigida por modelos (\textit{Model-Driven Engineering}, MDE)},
		firstplural={MDE},
		description={\emph{MDE} es un enfoque de desarrollo de software basado en la creaci{\'o}n y transformaci{\'o}n de modelos abstractos como principales artefactos del proceso. Permite separar niveles de abstracci{\'o}n, automatizar tareas mediante transformaciones de modelos y facilitar generaci{\'o}n de c{\'o}digo. Su objetivo es mejorar productividad, reutilizaci{\'o}n y calidad mediante mayor formalizaci{\'o}n del dise{\~n}o. En \glspl{sibiot}, MDE habilita dise{\~n}o sistem{\'a}tico de arquitecturas y la integraci{\'o}n de \glspl{modelopredictivo}.},
		type=main
	}
	
	\newglossaryentry{mec}
	{
		name={Multi-access Edge Computing (MEC)},
		text={MEC},
		sort={multi access edge computing},
		plural={sistemas MEC},
		first={Computaci{\'o}n perimetral (de borde) de acceso m{\'u}ltiple (Multi-access Edge Computing, MEC)},
		firstplural={sistemas MEC},
		description={Marco estandarizado para ejecutar aplicaciones y exponer servicios cerca del acceso de red, con arquitectura de referencia y gesti{\'o}n del ciclo de vida de aplicaciones para escenarios de baja \gls{latencia}.},
		type=\glsdefaulttype,
		symbol={📡},
		user1={Borde de red},
		user2={APIs estandarizadas},
		user3={Gesti{\'o}n de aplicaciones},
		see={edgecomputing,continuonubeperiferia}
	}
	
	\newglossaryentry{medioambiente}{
		name={Medio ambiente},
		text={medio ambiente},
		sort={medio ambiente},
		first={medio ambiente},
		description={Conjunto de condiciones naturales y antr{\'o}picas que rodean a un sistema (aire, agua, suelo, clima y actividad humana). En SIbIoT, es objeto de monitorizaci{\'o}n para gesti{\'o}n y mitigaci{\'o}n de impactos.},
		type=main,
		symbol={🌿},
		user1={Entorno natural},
		user2={Impacto},
		user3={Monitoreo},
		see={monitorizacion}
	}
	
	\newglossaryentry{memoria}{
		name={Memoria},
		text={memoria},
		sort={memoria},
		plural={memorias},
		first={memoria},
		firstplural={memorias},
		description={La \emph{memoria} es un subsistema fundamental de los computadores y \glspl{dispositivo} electr{\'o}nicos encargado de almacenar datos e instrucciones, ya sea de forma temporal o permanente. Incluye tecnolog{\'i}as vol{\'a}tiles como la \gls{ram}, utilizadas para el almacenamiento temporal durante la ejecuci{\'o}n de procesos, y tecnolog{\'i}as no vol{\'a}tiles como la \gls{rom}, \gls{flash} o unidades de estado s{\'o}lido, destinadas a conservar informaci{\'o}n incluso sin energ{\'i}a. En el contexto de los \glspl{sibiot}, la memoria determina la capacidad de procesamiento local, el almacenamiento de \glspl{dato}, la persistencia del \gls{firmware}, y la eficiencia en la ejecuci{\'o}n de algoritmos en el borde (\gls{edgecomputing}), influyendo directamente en el rendimiento, consumo energ{\'e}tico y confiabilidad del sistema.},
		type=main,
		symbol={🧠},
		user1={Almacenamiento temporal o permanente},
		user2={Recurso esencial para procesamiento y ejecuci{\'o}n},
		user3={Base del rendimiento en sistemas embebidos e IoT},
		see={ram,rom,flash,hardware,procesamientodatos,sibiot}
	}
	
	\newglossaryentry{memoriaflash}{
		name={Memoria Flash},
		text={memoria Flash},
		sort={memoria flash},
		first={memoria Flash},
		description={Tipo de memoria no vol{\'a}til utilizada para almacenar el programa principal y los datos constantes en un microcontrolador. La \emph{memoria Flash} conserva su contenido incluso cuando el dispositivo se encuentra sin alimentaci{\'o}n el{\'e}ctrica y permite un n{\'u}mero limitado de ciclos de escritura. En microcontroladores como el ATmega328P, esta memoria aloja el firmware cargado desde el entorno de desarrollo.},
		type=main,
		symbol={💾},
		user1={Memoria no vol{\'a}til},
		user2={Almacenamiento de programa},
		user3={Persistencia},
		see={sram,eeprom}
	}
	
	\newglossaryentry{mensajeriatiemporeal}{
		name={Mensajer{\'i}a en tiempo real},
		text={mensajer{\'i}a en tiempo real},
		sort={mensajeria en tiempo real},
		description={%
			Mecanismo de comunicaci{\'o}n as{\'i}ncrona que permite el intercambio inmediato de mensajes entre componentes de un sistema distribuido, con baja latencia y sin requerir una conexi{\'o}n directa o sincronizaci{\'o}n estricta entre emisor y receptor.
			La mensajer{\'i}a en tiempo real facilita la propagaci{\'o}n de eventos, notificaciones y cambios de estado de forma eficiente y escalable, constituyendo un elemento clave en arquitecturas reactivas y orientadas a eventos.
			En sistemas de informaci{\'o}n basados en IoT y SIbIoT, este tipo de mensajer{\'i}a se utiliza para desacoplar productores y consumidores de eventos, soportar flujos de datos continuos y habilitar respuestas oportunas ante situaciones din{\'a}micas del entorno f{\'i}sico y digital.
		},
		first={mensajer{\'i}a en tiempo real},
		plural={mensajer{\'i}as en tiempo real},
		firstplural={mensajer{\'i}as en tiempo real},
		see={enfoquereactivo,orientadoeventos,procesamientodistribuido,sistemadistribuido}
	}
	
	\newglossaryentry{metadato}{
		name={Metadato},
		text={metadato},
		plural={metadatos},
		first={metadato},
		firstplural={metadatos},
		see={dato},
		sort={metadato},
		description={Informaci{\'o}n estructurada que describe, contextualiza y aporta significado a otros datos. En SIbIoT, los metadatos incluyen atributos como tiempo de captura, ubicaci{\'o}n, unidad de medida, frecuencia de muestreo, estado de calibraci{\'o}n, fuente del sensor y versi{\'o}n del esquema, permitiendo trazabilidad, auditor{\'i}a, interoperabilidad y reproducibilidad del an{\'a}lisis}
	}
	
	% ===================== M{\'e}todo =====================
	\newglossaryentry{metodo}{
		name={M{\'e}todo},
		text={m{\'e}todo},
		sort={metodo},
		plural={m{\'e}todos},
		first={m{\'e}todo},
		firstplural={m{\'e}todos},
		description={Un \emph{m{\'e}todo} es un procedimiento sistem{\'a}tico, ordenado y replicable que se utiliza para alcanzar un objetivo espec{\'i}fico, resolver un problema o generar conocimiento. En proyectos \gls{iot}, los m{\'e}todos gu{\'i}an actividades como an{\'a}lisis de \glspl{dato}, evaluaci{\'o}n de desempe{\~n}o, pruebas de \glspl{algoritmo} o validaci{\'o}n de prototipos mediante pasos definidos y controlados.},
		type=\glsdefaulttype,
		symbol={🧾},
		user1={Procedimiento sistem{\'a}tico},
		user2={Gu{\'i}a de investigaci{\'o}n},
		user3={T{\'e}cnica replicable},
		see={metodologia,proceso,validacion}
	}
	
	% ===================== M{\'e}todo {\'a}gil =====================
	\newglossaryentry{metodoagil}{
		name={M{\'e}todo {\'a}gil},
		text={m{\'e}todo {\'a}gil},
		sort={metodo agil},
		plural={m{\'e}todos {\'a}giles},
		first={m{\'e}todo {\'a}gil},
		firstplural={m{\'e}todos {\'a}giles},
		description={Un \emph{m{\'e}todo {\'a}gil} es la implementaci{\'o}n concreta de una metodolog{\'i}a {\'a}gil que establece procesos, roles, artefactos y eventos espec{\'i}ficos para gestionar proyectos. Basado en los valores y principios del \gls{manifiestoagil}, prioriza la entrega continua, la iteraci{\'o}n frecuente, la colaboraci{\'o}n y la adaptabilidad. En \glspl{sibiot}, los m{\'e}todos {\'a}giles permiten coordinar equipos multidisciplinarios de \gls{hardware} y software, ajustando el desarrollo a condiciones cambiantes. Ejemplos incluyen Scrum, Kanban, XP y Crystal.},
		type=\glsdefaulttype,
		symbol={⚡},
		user1={Scrum},
		user2={Kanban},
		user3={Extreme Programming (XP)},
		see={manifiestoagil,equipoagil,iniciativaagil,kanban}
	}
	
	\newglossaryentry{metodocifrado}{
		name={M{\'e}todo de cifrado},
		text={m{\'e}todo de cifrado},
		sort={metodo de cifrado},
		plural={m{\'e}todos de cifrado},
		first={m{\'e}todo de cifrado},
		firstplural={m{\'e}todos de cifrado},
		description={T{\'e}cnica o algoritmo criptogr{\'a}fico empleado para transformar datos legibles (\textit{texto plano}) en un formato ilegible (\textit{texto cifrado}), con el fin de proteger su \gls{confidencialidad} y evitar accesos no autorizados. Los \emph{m{\'e}todos de cifrado} pueden clasificarse en sim{\'e}tricos (misma clave para cifrado y descifrado, como AES o DES) y asim{\'e}tricos (par de claves p{\'u}blica y privada, como RSA o ECC). En los \glspl{sibiot}, se utilizan para asegurar la transmisi{\'o}n de datos entre \glspl{sensor}, \glspl{actuador} y plataformas en la nube, garantizando integridad, autenticidad y protecci{\'o}n frente a ataques.},
		type=main,
		symbol={🧩},
		user1={Algoritmo de protecci{\'o}n de datos},
		user2={Clasificaci{\'o}n sim{\'e}trica y asim{\'e}trica},
		user3={Uso en IoT y seguridad de redes},
		see={confidencialidad,seguridad,privacidad,sibiot}
	}
	
	% ===================== Metodolog{\'i}a =====================
	\newglossaryentry{metodologia}{
		name={Metodolog{\'i}a},
		text={metodolog{\'i}a},
		sort={metodologia},
		plural={metodolog{\'i}as},
		first={metodolog{\'i}a},
		firstplural={metodolog{\'i}as},
		description={La \emph{metodolog{\'i}a} es el conjunto organizado de m{\'e}todos, principios y enfoques que estructuran y fundamentan un proceso de investigaci{\'o}n o desarrollo. A diferencia del \gls{metodo}, la metodolog{\'i}a constituye un marco conceptual m{\'a}s amplio que orienta la selecci{\'o}n, combinaci{\'o}n y justificaci{\'o}n de m{\'e}todos y t{\'e}cnicas. En proyectos \gls{iot}, permite integrar enfoques de hardware y software, equilibrando \gls{interoperabilidad}, \gls{seguridad} y objetivos de negocio.},
		type=\glsdefaulttype,
		symbol={📚},
		user1={Conjunto de m{\'e}todos},
		user2={Marco conceptual},
		user3={Fundamento de investigaci{\'o}n y desarrollo},
		see={metodo,metodoagil,metodologiadesarrollo}
	}
	
	% ===================== Metodolog{\'i}a {\'a}gil =====================
	\newglossaryentry{metodologiaagil}{
		name={Metodolog{\'i}a {\'a}gil},
		text={metodolog{\'i}a {\'a}gil},
		sort={metodologia agil},
		plural={metodolog{\'i}as {\'a}giles},
		first={metodolog{\'i}a {\'a}gil},
		firstplural={metodolog{\'i}as {\'a}giles},
		description={Una \emph{metodolog{\'i}a {\'a}gil} es un enfoque iterativo e incremental para la gesti{\'o}n y el desarrollo de proyectos, basado en los principios del \gls{manifiestoagil}. Prioriza la colaboraci{\'o}n con el cliente, la entrega continua de valor y la adaptaci{\'o}n al cambio. En sistemas \gls{iot}, facilita la coordinaci{\'o}n de equipos multidisciplinarios, la validaci{\'o}n temprana de prototipos y la gesti{\'o}n de la complejidad en entornos distribuidos y heterog{\'e}neos. Ejemplos incluyen Scrum, Kanban, XP y SAFe.},
		type=\glsdefaulttype,
		symbol={⚡},
		user1={Scrum},
		user2={Kanban},
		user3={Extreme Programming},
		see={manifiestoagil,metodoagil,equipoagil}
	}
	
	% ===================== Metodolog{\'i}a de desarrollo =====================
	\newglossaryentry{metodologiadesarrollo}{
		name={Metodolog{\'i}a de desarrollo},
		text={metodolog{\'i}a de desarrollo},
		sort={metodologia desarrollo},
		plural={metodolog{\'i}as de desarrollo},
		first={metodolog{\'i}a de desarrollo},
		firstplural={metodolog{\'i}as de desarrollo},
		description={Una \emph{metodolog{\'i}a de desarrollo} es un conjunto estructurado de principios, t{\'e}cnicas y pr{\'a}cticas que gu{\'i}an la planificaci{\'o}n, an{\'a}lisis, dise{\~n}o, implementaci{\'o}n, prueba y mantenimiento de \glspl{sistemainformatico}. En \gls{iot}, debe contemplar componentes de hardware y software, abordando conectividad, \gls{interoperabilidad}, tratamiento de \glspl{dato} y despliegue en \gls{nube} o \gls{edgecomputing}.},
		type=\glsdefaulttype,
		symbol={🛠️},
		user1={Ciclo de vida de software},
		user2={Procesos de ingenier{\'i}a},
		user3={Integraci{\'o}n hardware/software},
		see={metodologia,metodo,arquitectura}
	}
	
	% ===================== Metodolog{\'i}a Kanban =====================
	\newglossaryentry{metodologiakanban}{
		name={Metodolog{\'i}a Kanban},
		text={metodolog{\'i}a Kanban},
		sort={metodologia kanban},
		plural={metodolog{\'i}as Kanban},
		first={metodolog{\'i}a Kanban},
		firstplural={metodolog{\'i}as Kanban},
		description={La \emph{metodolog{\'i}a Kanban} es un enfoque de gesti{\'o}n del trabajo basado en principios lean y mejora continua. Se apoya en el uso del \gls{kanban} visual, incorporando pr{\'a}cticas como la limitaci{\'o}n del trabajo en curso (WIP), la medici{\'o}n de rendimiento (\textit{lead time}, \textit{throughput}) y la adaptaci{\'o}n progresiva. En entornos \gls{iot}, es eficaz en escenarios de alta variabilidad y necesidad de entrega continua.},
		type=\glsdefaulttype,
		symbol={📊},
		user1={Lean Kanban},
		user2={WIP},
		user3={Lead time / throughput},
		see={kanban,metodologiaagil,iniciativaagil}
	}
	
	% ===================== Microcontrolador =====================
	\newglossaryentry{microcontrolador}{
		name={Microcontrolador},
		text={microcontrolador},
		sort={microcontrolador},
		plural={microcontroladores},
		first={microcontrolador},
		firstplural={microcontroladores},
		description={Un \emph{microcontrolador} es un circuito integrado programable que contiene los componentes esenciales de un sistema embebido: unidad central de procesamiento (CPU), memoria y perif{\'e}ricos de entrada/salida. Son ampliamente utilizados en \glspl{dispositivo} \gls{iot} debido a su bajo consumo energ{\'e}tico, tama{\~n}o reducido y capacidad de ejecutar tareas espec{\'i}ficas de forma aut{\'o}noma. Ejemplos incluyen \gls{arduino}, STM32, MSP430, \gls{esp8266} y \gls{esp32}.},
		type=\glsdefaulttype,
		symbol={🔌},
		user1={Circuito integrado programable},
		user2={Control de hardware embebido},
		user3={Base de dispositivos IoT},
		see={microprocesador,arduino,esp32,esp8266}
	}
	
	% ===================== Microprocesador =====================
	\newglossaryentry{microprocesador}{
		name={Microprocesador},
		text={microprocesador},
		sort={microprocesador},
		plural={microprocesadores},
		first={microprocesador},
		firstplural={microprocesadores},
		description={Un \emph{microprocesador} es un circuito integrado que contiene la unidad central de procesamiento (CPU) de un sistema de c{\'o}mputo. Ejecuta instrucciones almacenadas en memoria, realiza operaciones aritm{\'e}tico-l{\'o}gicas y coordina perif{\'e}ricos externos mediante buses de datos, direcciones y control. Est{\'a} dise{\~n}ado para sistemas con mayores requerimientos de procesamiento, como computadoras, servidores o dispositivos embebidos avanzados. A diferencia de los \glspl{microcontrolador}, no incluye memoria ni perif{\'e}ricos integrados, lo que implica una arquitectura modular. Ejemplos incluyen procesadores Intel Core, AMD Ryzen y ARM Cortex-A.},
		type=\glsdefaulttype,
		symbol={🖥️},
		user1={CPU en chip},
		user2={Procesamiento de alto rendimiento},
		user3={Aplicaciones en PCs y servidores},
		see={cpu,microcontrolador,arquitectura}
	}
	
	% ===================== Microservicio =====================
	\newglossaryentry{microservicio}{
		name={Microservicio},
		text={microservicio},
		sort={microservicio},
		plural={microservicios},
		first={microservicio},
		firstplural={microservicios},
		description={Un \emph{microservicio} es una unidad funcional y aut{\'o}noma dentro de una arquitectura de software distribuida. Cada microservicio realiza una tarea espec{\'i}fica dentro de un sistema m{\'a}s amplio y se comunica con otros mediante interfaces bien definidas, usualmente \glspl{api} \gls{rest}. Pueden desarrollarse, desplegarse y mantenerse de forma independiente, lo que mejora modularidad, resiliencia y \gls{escalabilidad} de aplicaciones distribuidas, incluyendo plataformas \gls{iot}.},
		type=\glsdefaulttype,
		symbol={🧩},
		user1={Unidad funcional independiente},
		user2={Arquitectura modular distribuida},
		user3={Servicios escalables en IoT},
		see={api,rest,arquitecturadistribuida,escalabilidad}
	}
	
	% ===================== Microsoft SQL Server =====================
	\newglossaryentry{microsoftsqlserver}{
		name={Microsoft SQL Server},
		text={Microsoft SQL Server},
		sort={microsoft sql server},
		plural={Microsoft SQL Server},
		first={Microsoft SQL Server},
		firstplural={Microsoft SQL Server},
		description={\emph{Microsoft SQL Server}, en un sistema de gesti{\'o}n de \glspl{basedatosrelacional} que utiliza el lenguaje \gls{sql}, desarrollado por Microsoft, orientado a entornos corporativos. Adem{\'a}s de soportar SQL est{\'a}ndar, integra servicios de Azure, an{\'a}lisis de datos e inteligencia empresarial. En \glspl{sibiot}, se emplea para gestionar grandes vol{\'u}menes de \glspl{dato}, habilitar an{\'a}lisis en \gls{tiemporeal} y garantizar \gls{fiabilidad}, \gls{seguridad} y compatibilidad con ecosistemas Microsoft.},
		type=main,
		symbol={🗄️},
		user1={Base de datos corporativa},
		user2={Integraci{\'o}n con Azure y BI},
		user3={Aplicaciones en IoT},
		see={basedatosrelacional,fiabilidad,seguridad,sibiot}
	}
	
	\newglossaryentry{middleware}
	{
		name={Middleware},
		text={middleware},
		sort={middleware},
		plural={middleware},
		first={middleware},
		firstplural={middleware},
		description={El \emph{middleware} es una capa de software intermedia que proporciona servicios comunes tales como mensajer{\'i}a, gesti{\'o}n de \glspl{dato}, autenticaci{\'o}n, autorizaci{\'o}n, orquestaci{\'o}n y exposici{\'o}n de \glspl{api}, m{\'a}s all{\'a} de las capacidades ofrecidas por el sistema operativo. Su funci{\'o}n principal es facilitar la interoperabilidad, integraci{\'o}n y coordinaci{\'o}n entre \glspl{aplicacion}, \glspl{serviciodistribuido} y componentes heterog{\'e}neos dentro de una \gls{arquitecturadistribuida}.},
		type=main,
		symbol={🧬},
		user1={Capa intermedia},
		user2={Servicios de integraci{\'o}n},
		user3={Mensajer{\'i}a y seguridad},
		see={api,arquitecturadistribuida,coap,amqp}
	}
	
	\newglossaryentry{mineriadatos}{
		name={Miner{\'i}a de datos},
		text={miner{\'i}a de datos},
		sort={mineria de datos},
		plural={procesos de miner{\'i}a de datos},
		first={miner{\'i}a de datos},
		firstplural={procesos de miner{\'i}a de datos},
		description={Disciplina que combina estad{\'i}stica, inteligencia artificial y bases de datos para descubrir patrones, correlaciones y tendencias significativas en conjuntos masivos de informaci{\'o}n. La \emph{miner{\'i}a de datos} se aplica en clasificaci{\'o}n, agrupamiento, detecci{\'o}n de anomal{\'i}as, reglas de asociaci{\'o}n y predicci{\'o}n de comportamientos. En el contexto de los \glspl{sibiot}, permite analizar datos generados por sensores y dispositivos interconectados en tiempo real, facilitando la toma de decisiones, el mantenimiento predictivo, la optimizaci{\'o}n de recursos y el dise{\~n}o de servicios inteligentes.},
		type=main,
		symbol={⛏️},
		user1={\textit{Data Mining}},
		user2={Descubrimiento de conocimiento en bases de datos (KDD)},
		user3={Anal{\'i}tica avanzada de datos},
		see={procesamientodatos,bigdata,aprendizajeautomatico}
	}
	
	% ===================== Mitigaci{\'o}n =====================
	\newglossaryentry{mitigacion}{
		name={Mitigaci{\'o}n},
		text={mitigaci{\'o}n},
		sort={mitigacion},
		plural={mitigaciones},
		first={mitigaci{\'o}n},
		firstplural={mitigaciones},
		description={La \emph{mitigaci{\'o}n} es el conjunto de acciones, t{\'e}cnicas o estrategias destinadas a reducir impacto, probabilidad o severidad de un riesgo, amenaza o problema en un sistema. En \gls{iot}, se aplica frente a ataques de \gls{seguridad}, fallos de comunicaci{\'o}n, interferencias, p{\'e}rdida de energ{\'i}a o condiciones ambientales adversas, mediante \gls{redundancia}, cifrado, control de acceso o adaptaci{\'o}n din{\'a}mica del sistema.},
		type=\glsdefaulttype,
		symbol={🛡️},
		user1={Reducci{\'o}n de riesgos},
		user2={Acciones preventivas},
		user3={Protecci{\'o}n de sistemas IoT},
		see={seguridad,continuidad,estabilidad,redundancia}
	}
	
	\newglossaryentry{modelado}{
		name={Modelado},
		text={modelado},
		sort={modelado},
		plural={procesos de modelado},
		first={modelado},
		firstplural={procesos de modelado},
		description={Proceso sistem{\'a}tico mediante el cual se construyen representaciones conceptuales, matem{\'a}ticas o computacionales de un sistema, con el fin de comprender su funcionamiento, analizar su comportamiento, simular escenarios, predecir resultados y apoyar la toma de decisiones. En el contexto de \glspl{sibiot}, el \emph{modelado} integra modelos f{\'i}sicos, estad{\'i}sticos y de \gls{aprendizajeautomatico}, as{\'i} como modelos basados en agentes y de eventos discretos, para capturar din{\'a}micas de dispositivos, redes y procesos de negocio. Incluye actividades de definici{\'o}n de requisitos, identificaci{\'o}n de entidades y relaciones, formalizaci{\'o}n de restricciones, calibraci{\'o}n con \glspl{datobruto} y validaci{\'o}n/ verificaci{\'o}n del modelo (V\&V), adem{\'a}s de an{\'a}lisis de sensibilidad e incertidumbre.},
		type=main,
		symbol={🧩},
		user1={Representaci{\'o}n formal de sistemas},
		user2={Abstracci{\'o}n y simulaci{\'o}n},
		user3={Base para an{\'a}lisis y predicci{\'o}n},
		see={modelo,modelopredictivo,simulacion,diagramaflujo,pseudocodigo,aprendizajeautomatico,mineriadatos}
	}
	
	\newglossaryentry{modelo}{
		name={Modelo},
		text={modelo},
		sort={modelo},
		plural={modelos},
		first={modelo},
		firstplural={modelos},
		description={Representaci{\'o}n abstracta, simplificada o formal de un objeto, sistema o fen{\'o}meno, utilizada para comprender, analizar, predecir o controlar su comportamiento. Un \emph{modelo} puede ser matem{\'a}tico, computacional, f{\'i}sico o conceptual, y se apoya en datos, reglas o principios te{\'o}ricos. En el contexto de los \glspl{sibiot}, los modelos son esenciales para describir procesos de \gls{monitorizacion}, entrenar \glspl{redneuronal}, aplicar \glspl{modelopredictivo} y optimizar la gesti{\'o}n de recursos, contribuyendo a la \gls{eficienciaoperativa}, la \gls{resiliencia} y la \gls{sostenibilidad} de los sistemas inteligentes.},
		type=main,
		symbol={📐},
		user1={Representaci{\'o}n abstracta de un sistema},
		user2={Base para an{\'a}lisis, predicci{\'o}n y control},
		user3={Uso en IoT, IA y sistemas distribuidos},
		see={algoritmo,sistema,modelopredictivo,redneuronal,sibiot}
	}
	
	\newglossaryentry{modelopredictivo}{
		name={Modelo predictivo},
		text={modelo predictivo},
		sort={modelo predictivo},
		plural={modelos predictivos},
		first={modelo predictivo},
		firstplural={modelos predictivos},
		description={Representaci{\'o}n matem{\'a}tica, estad{\'i}stica o computacional dise{\~n}ada para anticipar el comportamiento futuro de un sistema, proceso o variable a partir de datos hist{\'o}ricos y actuales. Un \emph{modelo predictivo} utiliza t{\'e}cnicas como regresi{\'o}n, series temporales, \gls{aprendizajeautomatico} o \gls{ia} para identificar patrones y generar pron{\'o}sticos con un nivel de confianza determinado. En el contexto de los \glspl{sibiot}, los modelos predictivos permiten anticipar fallos en equipos, optimizar el mantenimiento, gestionar la demanda energ{\'e}tica, personalizar servicios o detectar anomal{\'i}as en \glspl{variablefisica} medidas por \glspl{sensor}, favoreciendo la toma de decisiones proactiva y la eficiencia operativa.},
		type=main,
		symbol={📈},
		user1={Anal{\'i}tica avanzada},
		user2={Pron{\'o}stico basado en datos},
		user3={Soporte a la toma de decisiones},
		see={aprendizajeautomatico,ia,eficienciaoperativa,sibiot}
	}
	
	\newglossaryentry{modeloreferencia}{
		name={Modelo de referencia},
		text={modelo de referencia},
		sort={modelo de referencia},
		plural={modelos de referencia},
		first={modelo de referencia},
		firstplural={modelos de referencia},
		description={Esquema conceptual que define un conjunto de funciones, componentes, relaciones y principios organizativos dentro de un dominio espec{\'i}fico, con el objetivo de servir como gu{\'i}a para el dise{\~n}o, an{\'a}lisis y comparaci{\'o}n de \glspl{sistema}. Un \emph{modelo de referencia} no prescribe implementaciones concretas, sino que proporciona un marco com{\'u}n para asegurar la coherencia, la \gls{interoperabilidad} y la aplicaci{\'o}n de buenas pr{\'a}cticas. En el contexto de los \glspl{sibiot}, los modelos de referencia, como el modelo OSI o el IoT-A, ayudan a estructurar capas de comunicaci{\'o}n, definir roles de \glspl{sensor} y \glspl{actuador}, y facilitar la integraci{\'o}n de servicios distribuidos y plataformas en la nube.},
		type=main,
		symbol={📐},
		user1={Esquema conceptual de gu{\'i}a},
		user2={Base para dise{\~n}o y comparaci{\'o}n de sistemas},
		user3={Soporte a interoperabilidad y est{\'a}ndares},
		see={arquitecturareferencia,arquitecturacapas,interoperabilidad,sibiot}
	}
	
	\newglossaryentry{modularidad}{
		name={Modularidad},
		text={modularidad},
		sort={modularidad},
		plural={modularidades},
		first={modularidad},
		firstplural={modularidades},
		description={La \emph{\gls{modularidad}} es la propiedad de un \gls{sistema} que permite dividir su estructura en componentes o m{\'o}dulos independientes, cohesivos y reutilizables, que pueden desarrollarse, entenderse, probarse y mantenerse de forma aislada. Esta caracter{\'i}stica facilita la \gls{escalabilidad}, el mantenimiento y la evoluci{\'o}n, constituyendo un principio fundamental en dise{\~n}o orientado a objetos, arquitecturas multicapa, \glspl{microservicio} y sistemas embebidos, as{\'i} como en ecosistemas \gls{iot}.},
		type=\glsdefaulttype,
		symbol={🧱},
		user1={Divisi{\'o}n en m{\'o}dulos cohesivos},
		user2={Reutilizaci{\'o}n y mantenibilidad},
		user3={Base de dise{\~n}o en OOP y microservicios},
		see={microservicio,escalabilidad,arquitecturacapas,arquitecturadistribuida}
	}
	
	\newglossaryentry{modulo}{
		name={M{\'o}dulo},
		text={m{\'o}dulo},
		sort={modulo},
		plural={m{\'o}dulos},
		first={m{\'o}dulo},
		firstplural={m{\'o}dulos},
		description={Unidad autocontenida de hardware o software con responsabilidad espec{\'i}fica e interfaz definida, dise{\~n}ada para integrarse en un sistema mayor.},
		type=main,
		symbol={🧩},
		user1={Componente},
		user2={Interfaz definida},
		user3={Reutilizaci{\'o}n},
		see={impl,integracionsoftware}
	}
	
	\newglossaryentry{mongodb}{
		name={MongoDB},
		text={MongoDB},
		sort={mongodb},
		plural={MongoDB},
		first={MongoDB},
		firstplural={MongoDB},
		description={Sistema de gesti{\'o}n de \glspl{basedatosnosql} orientado a documentos, de c{\'o}digo abierto, que almacena la informaci{\'o}n en formato BSON (\gls{bson}). \emph{MongoDB} permite manejar datos semi-estructurados y no estructurados de manera flexible, soportando escalabilidad horizontal mediante particionado (sharding) y alta disponibilidad con replicaci{\'o}n. En el contexto de los \glspl{sibiot}, MongoDB resulta adecuado para gestionar flujos masivos de datos generados por \glspl{sensor} en \gls{tiemporeal}, permitiendo su integraci{\'o}n con an{\'a}lisis en la \gls{nube}, aplicaciones de \gls{bigdata} y servicios distribuidos.},
		type=main,
		symbol={🍃},
		user1={Base de datos NoSQL orientada a documentos},
		user2={Uso de BSON para datos semi-estructurados},
		user3={Soporte a escalabilidad y datos IoT},
		see={basedatosnosql,bson,bigdata,sibiot}
	}
	
	\newglossaryentry{monitorizacion}{
		name={Monitorizaci{\'o}n},
		text={monitorizaci{\'o}n},
		sort={monitorizacion},
		plural={monitorizaciones},
		first={monitorizaci{\'o}n},
		firstplural={monitorizaciones},
		description={Proceso sistem{\'a}tico de observaci{\'o}n y medici{\'o}n continua de uno o varios par{\'a}metros de un \gls{sistema} o de su \gls{entorno}, con el fin de detectar desviaciones, anomal{\'i}as o tendencias relevantes para la \gls{tomadecisiones}. En el contexto de \glspl{sibiot}, la \emph{monitorizaci{\'o}n} implica la captura de \glspl{variablefisica} mediante \glspl{sensor}, la transmisi{\'o}n y el procesamiento de datos en \gls{tiemporeal}, as{\'i} como su correlaci{\'o}n con m{\'e}tricas de \gls{observabilidad} (m{\'e}tricas, logs, trazas) para garantizar \gls{disponibilidad}, \gls{fiabilidad} y \gls{eficienciaoperativa}. Puede habilitar alertas tempranas, \gls{mantenimientopredictivo} y \gls{optimizacion} continua de procesos.},
		type=main,
		symbol={📈},
		user1={Medici{\'o}n continua de par{\'a}metros},
		user2={Detecci{\'o}n de anomal{\'i}as y tendencias},
		user3={Base para decisiones y automatizaci{\'o}n en SIbIoT},
		see={sibiot,sensor,variablefisica,tiemporeal,observabilidad,disponibilidad,fiabilidad,eficienciaoperativa,mantenimientopredictivo,optimizacion}
	}
	
	\newglossaryentry{motordc}{
		name={Motor DC},
		text={motor DC},
		sort={motor dc},
		plural={motores DC},
		first={motor DC},
		firstplural={motores DC},
		description={Dispositivo electromec{\'a}nico que convierte energ{\'i}a el{\'e}ctrica en energ{\'i}a mec{\'a}nica mediante corriente continua (DC). Los \emph{motores DC} se emplean ampliamente en sistemas \gls{iot} y aplicaciones embebidas para generar movimiento rotativo en mecanismos de \gls{control}, \glspl{actuador}, ventiladores y plataformas m{\'o}viles. Se caracterizan por su sencillez de control (por ejemplo, modulaci{\'o}n por ancho de pulso, PWM), respuesta r{\'a}pida y compatibilidad con \glspl{microcontrolador}; en \glspl{sibiot}, son clave para la automatizaci{\'o}n y la eficiencia operativa.},
		type=main,
		symbol={⚙️},
		user1={Conversi{\'o}n de energ{\'i}a el{\'e}ctrica a mec{\'a}nica},
		user2={Control sencillo (p.\,ej., PWM) y respuesta r{\'a}pida},
		user3={Integraci{\'o}n con microcontroladores en IoT},
		see={actuador,microcontrolador,control,eficienciaoperativa,sibiot}
	}
	
	\newglossaryentry{motorreglas}{
		name={Motor de Reglas},
		text={motor de reglas},
		description={Componente de software especializado que permite definir, gestionar y ejecutar reglas de negocio de manera declarativa, separando la lógica de decisión del código de aplicación y facilitando la automatización, trazabilidad y modificabilidad de políticas organizacionales},
		sort={motorreglas},
		plural={motores de reglas},
		first={Motor de Reglas},
		symbol={⚙},
		type=main,
		see={reglanegocio,sistemadecision,automatizacion},
		user1={arquitectura empresarial},
		user2={sistemas basados en reglas},
		user3={gestión de políticas}
	}
	
	\newglossaryentry{movilidad}
	{
		name={Movilidad},
		text={movilidad},
		sort={movilidad},
		plural={movilidades},
		first={movilidad},
		firstplural={movilidades},
		description={\emph{\gls{movilidad}} es la capacidad de un \gls{dispositivo} o \gls{sistema} para desplazarse geogr{\'a}ficamente manteniendo la \gls{conectividad} y la continuidad de los \glspl{sistemacomunicacion}. En el contexto de los \glspl{sibiot}, la movilidad es esencial en aplicaciones como veh{\'i}culos aut{\'o}nomos, transporte inteligente, \glspl{ciudadinteligente} y soluciones industriales que requieren interacci{\'o}n constante aun cuando los \glspl{dispositivoiot} est{\'a}n en movimiento.},
		type=\glsdefaulttype,
		symbol={🚗},
		user1={Capacidad de desplazamiento},
		user2={Conectividad en movimiento},
		user3={Portabilidad},
		see={lte}
	}
	
	\newglossaryentry{mpls}{
		name={MPLS},
		text={MPLS},
		sort={mpls},
		plural={tecnolog{\'i}as MPLS},
		first={conmutaci{\'o}n de etiquetas multiprotocolo (\textit{Multiprotocol Label Switching}, MPLS)},
		firstplural={tecnolog{\'i}as MPLS},
		description={T{\'e}cnica de encaminamiento utilizada en redes de alta velocidad que dirige los paquetes de datos en funci{\'o}n de etiquetas asignadas en lugar de direcciones IP completas. \emph{MPLS} mejora la eficiencia del enrutamiento, reduce la latencia y permite la ingenier{\'i}a de tr{\'a}fico en redes complejas. Es ampliamente empleada en redes \gls{wan}, redes privadas virtuales (VPN) y servicios empresariales. En entornos \gls{iot}, \emph{MPLS} facilita la transmisi{\'o}n segura y priorizada de datos cr{\'i}ticos en infraestructuras distribuidas.},
		type=main,
		symbol={🔖},
		user1={Encaminamiento por etiquetas en lugar de direcciones IP},
		user2={Soporte en redes WAN y VPN},
		user3={Priorizaci{\'o}n y seguridad de tr{\'a}fico IoT},
		see={wan,red,infraestructurared,iot}
	}
	
	\newglossaryentry{mqtt}{
		name={MQTT},
		text={MQTT},
		sort={mqtt},
		plural={protocolos MQTT},
		first={Transporte de Telemetr{\'i}a de Cola de Mensajes (\textit{Message Queue Telemetry Transport}, MQTT)},
		firstplural={protocolos MQTT},
		description={Protocolo de comunicaci{\'o}n ligero de tipo \gls{m2m} (m{\'a}quina a m{\'a}quina) basado en el modelo de publicaci{\'o}n/suscripci{\'o}n. \emph{MQTT} est{\'a} dise{\~n}ado para redes con baja latencia, bajo consumo de energ{\'i}a y conectividad intermitente, lo que lo hace ideal para sistemas \gls{iot}. Funciona sobre TCP/IP y emplea un \textit{broker} central que gestiona los mensajes entre clientes, facilitando la \gls{interoperabilidad}, la \gls{escalabilidad} y la integraci{\'o}n de dispositivos distribuidos.},
		type=main,
		symbol={📩},
		user1={Protocolo de mensajer{\'i}a ligera para IoT},
		user2={Modelo publicaci{\'o}n/suscripci{\'o}n con broker},
		user3={Alta eficiencia en redes de baja potencia},
		see={m2m,iot,rediot,escalabilidad,interoperabilidad}
	}
	
	\newglossaryentry{msoa}{
		name={MSOA},
		text={MSOA},
		sort={msoa},
		plural={MSOAs},
		first={arquitectura orientada a microservicios (\textit{Microservices Oriented Architecture}, MSOA)},
		firstplural={MSOAs},
		description={Estilo de dise{\~n}o de software que organiza una aplicaci{\'o}n como un conjunto de servicios peque{\~n}os, independientes y desplegables de forma aut{\'o}noma, que se comunican entre s{\'i} a trav{\'e}s de interfaces bien definidas (normalmente \gls{api} \gls{rest}). Cada \gls{microservicio} se centra en una funcionalidad espec{\'i}fica y puede desarrollarse, escalarse y actualizarse de forma aislada. En sistemas \gls{iot}, esta arquitectura permite gestionar la complejidad, mejorar la \gls{escalabilidad} y facilitar la \gls{integracion} con servicios distribuidos.},
		type=main,
		symbol={🧩},
		user1={Estilo de dise{\~n}o basado en microservicios},
		user2={Independencia y despliegue aut{\'o}nomo de componentes},
		user3={Alta compatibilidad con arquitecturas IoT distribuidas},
		see={api,rest,microservicio,iot,escalabilidad,integracion}
	}
	
	\newglossaryentry{mundofisico}{
		name={Mundo f{\'i}sico},
		text={mundo f{\'i}sico},
		sort={mundo fisico},
		plural={mundos f{\'i}sicos},
		first={mundo f{\'i}sico},
		firstplural={mundos f{\'i}sicos},
		description={El mundo f{\'i}sico hace referencia al entorno material y tangible conformado por objetos, espacios, personas, recursos naturales y elementos construidos. Es el espacio en el que ocurren procesos naturales, actividades humanas y operaciones industriales. En el contexto de los \glspl{sibiot}, el mundo f{\'i}sico es la fuente de datos capturados mediante \glspl{sensor}, as{\'i} como el espacio donde act{\'u}an los \glspl{actuador} para modificar condiciones reales, lo que permite la interacci{\'o}n entre la tecnolog{\'i}a digital y la realidad concreta.},
		type=main,
		symbol={🌍},
		user1={Entorno material y tangible},
		user2={Fuente y destino de datos IoT},
		user3={Interacci{\'o}n con sensores y actuadores},
		see={entornofisico,entornoiot,realidad}
	}
	
	\newglossaryentry{mundodigital}{
		name={Mundo digital},
		text={mundo digital},
		sort={mundo digital},
		plural={mundos digitales},
		first={mundo digital},
		firstplural={mundos digitales},
		description={El \emph{mundo digital} es el {\'a}mbito compuesto por datos, \glspl{proceso} y servicios inform{\'a}ticos donde se adquiere, transporta, almacena y transforma informaci{\'o}n proveniente del \gls{mundofisico}. En el contexto de \gls{iot}/\glspl{sibiot}, abarca el \gls{procesamientodatos} en el \gls{borde} y en la \gls{nube}, las \glspl{basedatos}, los flujos \gls{etl}, las \glspl{api} y la \gls{infraestructuradistribuida} que habilitan la \gls{tomadecisiones} en \gls{tiemporeal}. Sirve de base para visualizaciones, \glspl{dashboard} y la construcci{\'o}n de \glspl{gemelodigital}.},
		type=main,
		symbol={💾},
		user1={Espacio de datos y c{\'o}mputo},
		user2={Procesos y servicios que integran IoT},
		user3={Soporte para anal{\'i}tica y decisiones},
		see={mundofisico,gemelodigital,etl,procesamientodatos,nube,borde,sibiot}
	}
	
	\newglossaryentry{mundovirtual}{
		name={Mundo virtual},
		text={mundo virtual},
		sort={mundo virtual},
		plural={mundos virtuales},
		first={mundo virtual},
		firstplural={mundos virtuales},
		description={El \emph{mundo virtual} es un entorno digital simulado que representa objetos, procesos y estados mediante modelos computacionales, permitiendo la interacci{\'o}n visual o inmersiva (por ejemplo, \gls{realidadvirtual} y \gls{realidadaumentada}). En \gls{iot}/\glspl{sibiot}, se materializa en \glspl{gemelodigital}, simuladores y \glspl{dashboard} que integran datos del \gls{mundofisico} y del \gls{mundodigital} para experimentar, predecir y optimizar el comportamiento de sistemas complejos antes de ejecutar acciones reales.},
		type=main,
		symbol={💻},
		user1={Entorno simulado e interactivo},
		user2={Gemelos digitales y simulaci{\'o}n},
		user3={Visualizaci{\'o}n para decisiones en IoT},
		see={mundofisico,mundodigital,gemelodigital,realidadvirtual,realidadaumentada,sibiot}
	}
	
	\newglossaryentry{mvc}{
		name={MVC},
		text={MVC},
		sort={mvc},
		plural={arquitecturas MVC},
		first={Modelo–Vista–Controlador (MVC)},
		firstplural={arquitecturas MVC},
		description={\emph{MVC} es un patr{\'o}n de arquitectura de software que separa una aplicaci{\'o}n en tres componentes principales: el \textbf{Modelo}, que gestiona los datos, la l{\'o}gica de negocio y las reglas del sistema; la \textbf{Vista}, que presenta la informaci{\'o}n al usuario; y el \textbf{Controlador}, que act{\'u}a como intermediario entre la vista y el modelo, procesando las solicitudes y coordinando las respuestas. Esta separaci{\'o}n facilita la modularidad, el mantenimiento, la reutilizaci{\'o}n de componentes y el desarrollo paralelo de interfaces y l{\'o}gica interna. En sistemas \gls{iot} y \glspl{sibiot}, MVC se utiliza para estructurar aplicaciones web, paneles de visualizaci{\'o}n, herramientas de gesti{\'o}n y servicios de \gls{monitorizacion}, garantizando escalabilidad, claridad en el c{\'o}digo e integraci{\'o}n eficaz con dispositivos conectados},
		type=main,
		symbol={🎛️},
		user1={Separaci{\'o}n clara entre modelo, vista y controlador},
		user2={Facilita mantenimiento y escalabilidad},
		user3={Patr{\'o}n ampliamente usado en aplicaciones web e IoT},
		see={arquitecturacapas,controller,aplicacionweb,modularidad}
	}
	
	\newglossaryentry{mysql}{
		name={MySQL},
		text={MySQL},
		sort={mysql},
		plural={MySQL},
		first={MySQL},
		firstplural={MySQL},
		description={Sistema de gesti{\'o}n de \glspl{basedatosrelacional} de c{\'o}digo abierto, ampliamente utilizado para aplicaciones web y empresariales. \emph{MySQL} emplea \gls{sql} como lenguaje principal y destaca por su simplicidad, rapidez y portabilidad en diversas plataformas. Es com{\'u}n en arquitecturas LAMP (Linux, Apache, MySQL, PHP/Python/Perl) y en entornos en la nube. En el contexto de los \glspl{sibiot}, se utiliza para gestionar datos estructurados de \glspl{sensor} y \glspl{actuador}, soportando aplicaciones que requieren integridad, \gls{fiabilidad} y escalabilidad moderada.},
		type=main,
		symbol={🐬},
		user1={Base de datos relacional de c{\'o}digo abierto},
		user2={Uso extensivo en aplicaciones web},
		user3={Soporte en arquitecturas IoT},
		see={basedatosrelacional,basedatos,sibiot}
	}
	
	\newglossaryentry{navegador}{
		name={Navegador},
		text={navegador},
		sort={navegador},
		plural={navegadores},
		first={navegador},
		firstplural={navegadores},
		description={\glslink{browser}{V{\'e}ase \emph{browser}}. Aplicaci{\'o}n cliente utilizada para acceder e interactuar con contenidos y \glspl{servicio} en la web; en \glspl{aplicacionweb} para \gls{sibiot} act{\'u}a como interfaz de usuario para \gls{monitorizacion}, \gls{control} y visualizaci{\'o}n de datos.},
		type=main,
		see={browser,aplicacionweb,sibiot}
	}
	
	\newglossaryentry{nbiot}{
		name={NB-IoT},
		text={NB-IoT},
		sort={nb-iot},
		first={banda estrecha para Internet de las cosas (\textit{Narrowband Internet of Things}, NB-IoT)},
		plural={tecnolog{\'i}as NB-IoT},
		firstplural={tecnolog{\'i}as NB-IoT},
		description={Tecnolog{\'i}a de comunicaci{\'o}n celular de tipo \gls{lpwan}, dise{\~n}ada para \glspl{dispositivoiot} que requieren bajo consumo energ{\'e}tico y conectividad confiable en entornos de cobertura limitada. Aprovecha infraestructuras de redes m{\'o}viles existentes (LTE) para transmitir peque{\~n}as cantidades de datos, ofreciendo gran cobertura en interiores, soporte para gran densidad de nodos y costos reducidos de despliegue. Es ideal para aplicaciones como medici{\'o}n remota de servicios, \gls{logistica}, \glspl{ciudadinteligente} y monitorizaci{\'o}n ambiental.},
		type=main,
		see={lorawan,nfc,bluetooth,zigbee,comunicacioninalambrica,iot}
	}
	
	\newglossaryentry{nfc}{
		name={NFC},
		text={NFC},
		sort={nfc},
		plural={tecnolog{\'i}as NFC},
		first={comunicaci{\'o}n de campo cercano (\textit{Near Field Communication}, NFC)},
		firstplural={tecnolog{\'i}as NFC},
		description={Tecnolog{\'i}a de comunicaci{\'o}n inal{\'a}mbrica de muy corto alcance (cent{\'i}metros) que permite el intercambio r{\'a}pido de \glspl{dato} cuando dos \glspl{dispositivo} est{\'a}n pr{\'o}ximos. Se utiliza en pagos m{\'o}viles, \gls{controlacceso}, identificaci{\'o}n y emparejamiento. En \gls{iot}, facilita configuraci{\'o}n inicial, \gls{autenticacion} y transmisi{\'o}n segura de informaci{\'o}n entre \gls{nodo}.},
		type=main,
		symbol={📳},
		user1={Enlace inal{\'a}mbrico de corto alcance},
		user2={Pagos, acceso e identificaci{\'o}n},
		user3={Onboarding y seguridad en IoT},
		see={comunicacioninalambrica,controlacceso,autenticacion,iot,nodo}
	}
	
	\newglossaryentry{nodo}{
		name={Nodo},
		text={nodo},
		sort={nodo},
		plural={nodos},
		first={nodo},
		firstplural={nodos},
		description={Elemento o punto de interconexi{\'o}n dentro de una \gls{red} o \gls{arquitecturadistribuida}, capaz de enviar, recibir o procesar informaci{\'o}n. Un \emph{nodo} puede ser un \gls{sensor}, un \gls{actuador}, un \gls{dispositivo} inteligente, un servidor o cualquier componente que participe activamente en la comunicaci{\'o}n y el procesamiento de datos. En los \glspl{sibiot}, los nodos conforman la infraestructura distribuida que permite la captura de \glspl{variablefisica}, la ejecuci{\'o}n de \glspl{algoritmo} de \gls{control} y la integraci{\'o}n con \glspl{servicio} en la \gls{nube}. La cantidad, ubicaci{\'o}n y funciones de los \glspl{nodo} determinan la \gls{escalabilidad}, la \gls{fiabilidad} y la \gls{resiliencia} del \gls{sistema}.},
		type=main,
		symbol={🔗},
		user1={Elemento de interconexi{\'o}n en redes},
		user2={Unidad de procesamiento o comunicaci{\'o}n},
		user3={Componente b{\'a}sico en SIbIoT},
		see={red,arquitecturadistribuida,sensor,actuador,sibiot}
	}
	
	\newglossaryentry{normalizacion}{
		name={normalizaci{\'o}n},
		text={normalizaci{\'o}n},
		sort={normalizacion},
		plural={normalizaciones},
		first={normalizaci{\'o}n},
		firstplural={normalizaciones},
		description={La normalizaci{\'o}n transforma los datos heterog{\'e}neos provenientes de distintas fuentes en formatos compatibles y comparables. En los \glspl{sibiot}, este proceso permite que la informaci{\'o}n recolectada por diversos \glspl{sensor} y plataformas sea integrada sin p{\'e}rdida de significado. La normalizaci{\'o}n es esencial para la \gls{interoperabilidad}, el \gls{analisispredictivo} y la \gls{automatizacion} de \glspl{proceso}, garantizando una \gls{basedatos} uniforme y f{\'a}cilmente procesable.},
		type=\glsdefaulttype,
		symbol={🧩},
		user1={Data normalization},
		user2={Estandarizaci{\'o}n de formatos y valores para mantener coherencia de datos},
		user3={Transformaci{\'o}n de datos a estructuras homog{\'e}neas para an{\'a}lisis y procesamiento},
		see={depuracion,deduplicacion,enriquecimientodatos,interoperabilidad,automatizacion}
	}
	
	\newglossaryentry{normalizaciondatos}{
		name={Normalizaci{\'o}n de datos},
		text={normalizaci{\'o}n de datos},
		sort={normalizaciondatos},
		plural={normalizaciones de datos},
		first={normalizaci{\'o}n de datos},
		firstplural={normalizaciones de datos},
		description={Proceso sistem{\'a}tico de organizaci{\'o}n y estructuraci{\'o}n de los datos con el objetivo de reducir redundancias, evitar inconsistencias y mejorar la integridad de la informaci{\'o}n. En sistemas de informaci{\'o}n y \glspl{sibiot}, la normalizaci{\'o}n de datos resulta fundamental para garantizar la coherencia entre fuentes heterog{\'e}neas, facilitar el almacenamiento eficiente, mejorar la calidad de los an{\'a}lisis y soportar procesos de integraci{\'o}n y toma de decisiones basadas en datos confiables.},
		type=main,
		symbol={📐},
		user1={Estructuraci{\'o}n de datos},
		user2={Reducci{\'o}n de redundancia},
		user3={Integridad de la informaci{\'o}n},
		see={gestiondatos,analiticadatos,basedatos,calidadsoftware}
	}
	
	\newglossaryentry{nosql}{
		name={NoSQL},
		text={NoSQL},
		sort={nosql},
		first={NoSQL},
		description={Familia de sistemas de gesti{\'o}n de bases de datos no relacionales dise{\~n}ados para manejar grandes vol{\'u}menes de datos de forma distribuida, escalable y con esquemas flexibles. Los sistemas NoSQL incluyen modelos de datos clave--valor, orientados a documentos, columnas anchas y grafos, y priorizan la escalabilidad horizontal y la disponibilidad frente a la consistencia estricta. En el contexto de los sistemas de informaci{\'o}n basados en IoT y SIbIoT, las bases de datos NoSQL se utilizan para almacenar eventos, telemetr{\'i}a, estados de dispositivos y datos semiestructurados generados a alta velocidad},
		type=main,
		symbol={📊},
		see={basedatos,procesamientocontinuo,procesamientodistribuido,escalabilidad,sibiot}
	}
	
	\newglossaryentry{notificacion}
	{
		name={Notificaci{\'o}n},
		text={notificaci{\'o}n},
		sort={notificacion},
		plural={notificaciones},
		first={Notificaci{\'o}n},
		description={Mensaje informativo generado por un \gls{sistema} para comunicar un evento, resultado o cambio de estado a un usuario o a otro componente del sistema. 
			En los \glspl{sibiot}, las notificaciones permiten la interacci{\'o}n entre el entorno f{\'i}sico y digital, facilitando la supervisi{\'o}n remota y la toma de decisiones. 
			Pueden emitirse mediante interfaces gr{\'a}ficas, servicios de mensajer{\'i}a o integraciones en plataformas web y m{\'o}viles. 
			A diferencia de las \glspl{alarma} o \glspl{alerta}, las notificaciones no implican necesariamente una condici{\'o}n de riesgo, sino la comunicaci{\'o}n de informaci{\'o}n relevante o de seguimiento de procesos.},
		symbol={🔔},
		user1={Comunicaci{\'o}n de eventos o cambios de estado del sistema},
		user2={Interacci{\'o}n informativa entre entorno f{\'i}sico y digital},
		user3={No implica riesgo, sino informaci{\'o}n contextual en \glspl{sibiot}},
		see={[Relacionado con:]{alerta,alarma,evento,comunicacion}}
	}
	
	\newglossaryentry{ntp}{
		name={NTP},
		text={NTP},
		sort={ntp},
		plural={NTP},
		first={protocolo de tiempo en red (\textit{Network Time Protocol}, NTP)},
		firstplural={NTP},
		description={\emph{NTP} es un protocolo de sincronizaci{\'o}n de tiempo dise{\~n}ado para mantener relojes de sistemas distribuidos con una gran precisi{\'o}n a trav{\'e}s de redes IP. Opera mediante una jerarqu{\'i}a de servidores organizados en niveles denominados \emph{strata}, que permiten minimizar la deriva del reloj local y compensar retardos de red. En sistemas \gls{iot}, NTP es fundamental para garantizar coherencia temporal en eventos, registros, algoritmos de \gls{procesamientodatos}, seguridad, auditor{\'i}as y coordinaci{\'o}n entre \glspl{dispositivoconectado}, especialmente cuando se requiere ordenamiento preciso de mediciones o ejecuci{\'o}n secuencial de procesos distribuidos.},
		type=main,
		symbol={⏱️},
		user1={Sincronizaci{\'o}n de tiempo mediante jerarqu{\'i}as},
		user2={Compensaci{\'o}n de retardos en redes IP},
		user3={Base temporal para sistemas distribuidos e IoT},
		see={ptp,protocoloseguridad,infraestructuradistribuida,sibiot}
	}
	
	\newglossaryentry{nube}{
		name={Nube},
		text={nube},
		sort={nube},
		plural={nubes},
		first={nube},
		firstplural={nubes},
		description={Modelo de prestaci{\'o}n de servicios de computaci{\'o}n que permite el acceso remoto a recursos como almacenamiento, procesamiento, redes o aplicaciones, a trav{\'e}s de Internet. La \emph{nube} ofrece caracter{\'i}sticas como escalabilidad bajo demanda, pago por uso y disponibilidad global. En el contexto de los \glspl{sibiot}, la nube es fundamental para gestionar grandes vol{\'u}menes de datos generados por \glspl{sensor}, facilitar la \gls{monitorizacion} en \gls{tiemporeal}, aplicar algoritmos de \gls{aprendizajeautomatico} o \glspl{modelopredictivo}, y asegurar \gls{interoperabilidad} mediante integraci{\'o}n de servicios distribuidos.},
		type=main,
		symbol={☁️},
		user1={Computaci{\'o}n como servicio en Internet},
		user2={Escalabilidad y disponibilidad global},
		user3={Soporte para SIbIoT y datos masivos},
		see={sibiot,serviciodistribuido,eficienciaoperativa,interoperabilidad}
	}
	
	\newglossaryentry{nucleo}{
		name={N{\'u}cleo},
		text={n{\'u}cleo},
		sort={nucleo},
		plural={n{\'u}cleos},
		first={n{\'u}cleo o core},
		firstplural={n{\'u}cleos},
		description={El n{\'u}cleo representa la parte central o \gls{componente} fundamental de un \gls{sistema}, desde el punto de vista del \gls{procesamientodatos}, coordinaci{\'o}n de tareas o \gls{tomadecisiones}. En el contexto de los \glspl{sibiot}, puede referirse al elemento encargado de ejecutar los \glspl{algoritmo} de \gls{control}, \gls{comunicacion} y \gls{gestion} de los flujos de \gls{informacion}.},
		type=\glsdefaulttype,
		symbol={🧠},
		user1={Core},
		user2={Componente central},
		user3={Procesamiento principal},
		see={sibiot,procesamiento,arquitectura}
	}
	
	\newglossaryentry{nucleoperativo}{
		name={N{\'u}cleo operativo},
		text={n{\'u}cleo operativo},
		sort={nucleoperativo},
		plural={n{\'u}cleos operativos},
		first={n{\'u}cleo operativo},
		firstplural={n{\'u}cleos operativos},
		description={El n{\'u}cleo operativo agrupa los componentes que garantizan el funcionamiento estable y continuo del sistema. En un entorno \gls{iot} o \gls{sibiot}, incluye los mecanismos de procesamiento, control de tareas, manejo de eventos y comunicaci{\'o}n entre \glspl{dispositivo} y \glspl{servicio} centralizados o distribuidos},
		type=\glsdefaulttype,
		symbol={⚙️},
		user1={Operational core},
		user2={Unidad funcional principal},
		user3={Gesti{\'o}n de procesos},
		see={nucleo,iot,sibiot}
	}
	
	\newglossaryentry{nvidia}{
		name={NVIDIA},
		text={NVIDIA},
		sort={nvidia},
		plural={NVIDIA},
		first={NVIDIA},
		firstplural={NVIDIA},
		description={Empresa multinacional estadounidense especializada en el desarrollo de \glspl{gpu}, plataformas de computaci{\'o}n acelerada y soluciones para \gls{ia}, visualizaci{\'o}n y procesamiento embebido. Las tecnolog{\'i}as de \emph{NVIDIA} se emplean ampliamente en sistemas \gls{iot}, \gls{aprendizajeautomatico} (deep learning), \gls{visionartificial}, \gls{vehiculoautonomo} y centros de datos, habilitando c{\'a}lculo paralelo masivo y \glspl{modelopredictivo} en \gls{tiemporeal}.},
		type=main,
		symbol={🎮},
		user1={L{\'i}der en GPU y c{\'o}mputo acelerado},
		user2={Soluciones para IA, visi{\'o}n y edge},
		user3={Aplicaciones en IoT y centros de datos},
		see={gpu,ia,aprendizajeautomatico,visionartificial,iot,vehiculoautonomo}
	}
	
	\newglossaryentry{nvidiajetson}{
		name={NVIDIA Jetson},
		text={NVIDIA Jetson},
		sort={nvidia jetson},
		plural={NVIDIA Jetson},
		first={NVIDIA Jetson},
		firstplural={NVIDIA Jetson},
		description={Familia de m{\'o}dulos de computaci{\'o}n de alto rendimiento para \gls{ia} en el borde (\emph{edge}) desarrollados por \gls{nvidia}. Los m{\'o}dulos \emph{Jetson} integran \gls{gpu} y CPU para procesamiento paralelo eficiente, acelerando tareas de \gls{visionartificial}, rob{\'o}tica, an{\'a}lisis en \gls{tiemporeal} y control de \glspl{actuador}/\glspl{sensor} en aplicaciones embebidas. Son comunes en \glspl{sibiot} por su balance entre rendimiento, consumo y ecosistema de software.},
		type=main,
		symbol={🧠},
		user1={C{\'o}mputo acelerado para IA en el borde},
		user2={Integraci{\'o}n GPU+CPU para visi{\'o}n y rob{\'o}tica},
		user3={Uso en sistemas embebidos e IoT},
		see={nvidia,gpu,visionartificial,iot,fogcomputing}
	}
	
	\newglossaryentry{objetivo}{
		name={Objetivo},
		text={objetivo},
		sort={objetivo},
		plural={objetivos},
		first={objetivo},
		firstplural={objetivos},
		description={Resultado deseado que orienta acciones y decisiones. En proyectos y sistemas, los objetivos gu{\'i}an requisitos, dise{\~n}o, implementaci{\'o}n y evaluaci{\'o}n.},
		type=main,
		symbol={🎯},
		user1={Meta},
		user2={Direcci{\'o}n},
		user3={Evaluaci{\'o}n},
		see={resultado,proceso}
	}
	
	\newglossaryentry{objetocotidiano}{
		name={Objeto cotidiano},
		text={objeto cotidiano},
		sort={objeto cotidiano},
		plural={objetos cotidianos},
		first={objeto cotidiano},
		firstplural={objetos cotidianos},
		description={Elemento del \gls{entornofisico} utilizado en la vida diaria (por ejemplo, electrodom{\'e}sticos, veh{\'i}culos, prendas o utensilios) que puede ser dotado de \glspl{sensor}, conectividad y capacidad de procesamiento en el contexto de \gls{iot}. Al incorporar \glspl{dispositivointeligente} o convertirlos en \glspl{dispositivoconectado}, los \emph{objetos cotidianos} posibilitan \gls{monitorizacion}, automatizaci{\'o}n y \gls{conectividad} con plataformas y servicios distribuidos.},
		type=main,
		symbol={🏠},
		user1={Objeto del entorno f{\'i}sico de uso diario},
		user2={Susceptible de sensorizar y conectar (IoT)},
		user3={Base de aplicaciones dom{\'o}ticas e industriales},
		see={entornofisico,dispositivointeligente,dispositivoconectado,iot,monitorizacion,serviciodistribuido}
	}
	
	\newglossaryentry{objetointeligente}{
		name={Objeto inteligente},
		text={objeto inteligente},
		sort={objeto inteligente},
		plural={objetos inteligentes},
		first={objeto inteligente},
		firstplural={objetos inteligentes},
		description={Dispositivo o artefacto cotidiano que integra sensores, actuadores, capacidad de procesamiento y conectividad, permitiendo interactuar de forma aut{\'o}noma o semiaut{\'o}noma con su entorno y con otros sistemas. Los \emph{objetos inteligentes} incluyen desde electrodom{\'e}sticos, veh{\'i}culos y dispositivos port{\'a}tiles, hasta sensores industriales y sistemas de control ambiental. En el contexto de los \glspl{sibiot}, los objetos inteligentes constituyen la base de la \gls{interaccion} entre el mundo f{\'i}sico y digital, facilitando la automatizaci{\'o}n, el an{\'a}lisis de datos en \gls{tiemporeal} y la creaci{\'o}n de servicios avanzados en sectores como la salud, la agricultura, la industria y las ciudades inteligentes.},
		type=main,
		symbol={📡},
		user1={Dispositivo inteligente},
		user2={Artefacto conectado},
		user3={Entidad IoT},
		see={dispositivo,dispositivointeligente,sibiot,visionartificial}
	}
	
	\newglossaryentry{observabilidad}{
		name={Observabilidad},
		text={observabilidad},
		sort={observabilidad},
		plural={observabilidades},
		first={observabilidad},
		firstplural={observabilidades},
		description={Propiedad de un \gls{sistema} que mide la capacidad de inferir su estado interno a partir de las salidas o datos observables que genera, como m{\'e}tricas, trazas y registros de eventos. La \emph{observabilidad} se basa en principios de ingenier{\'i}a de control y ha evolucionado como disciplina clave en el {\'a}mbito de las arquitecturas distribuidas, el \gls{devops} y los \glspl{sibiot}. A diferencia de la simple \gls{monitorizacion}, la observabilidad permite no s{\'o}lo detectar fallos, sino tambi{\'e}n diagnosticar sus causas y anticipar problemas en sistemas complejos y din{\'a}micos. En entornos \gls{iot}, una alta observabilidad es esencial para garantizar \gls{fiabilidad}, \gls{disponibilidad} y \gls{resiliencia}, al proporcionar visibilidad en tiempo real sobre el comportamiento de \glspl{sensor}, \glspl{actuador} y flujos de datos.},
		type=main,
		symbol={🔍},
		user1={Inferencia del estado interno de sistemas},
		user2={Diagn{\'o}stico en arquitecturas distribuidas},
		user3={Visibilidad en tiempo real en SIbIoT},
		see={monitorizacion,fiabilidad,disponibilidad,resiliencia,sibiot}
	}
	
	\newglossaryentry{olap}{
		name={OLAP},
		text={OLAP},
		sort={olap},
		plural={OLAPs},
		first={procesamiento anal{\'i}tico en l{\'i}nea (\textit{On-Line Analytical Processing}, OLAP)},
		firstplural={OLAPs},
		description={\emph{OLAP} es un enfoque de gesti{\'o}n de datos orientado al an{\'a}lisis multidimensional de grandes vol{\'u}menes de informaci{\'o}n. A diferencia de \gls{oltp}, que se centra en transacciones r{\'a}pidas y frecuentes, OLAP permite ejecutar consultas complejas, agregaciones y visualizaciones sobre datos hist{\'o}ricos para apoyar la \gls{tomadecisiones} estrat{\'e}gicas. En entornos de \gls{bigdata} y \glspl{sibiot}, OLAP se aplica al procesamiento de flujos masivos provenientes de \glspl{sensor} y \glspl{dispositivointeligente}, facilitando la detecci{\'o}n de patrones, anomal{\'i}as y correlaciones en tiempo cuasi real. Es la base de sistemas de \gls{bi} y cuadros de mando anal{\'i}ticos.},
		type=main,
		symbol={📊},
		user1={An{\'a}lisis multidimensional de datos},
		user2={Soporte a decisiones estrat{\'e}gicas},
		user3={Procesamiento anal{\'i}tico en SIbIoT y Big Data},
		see={oltp,htap,bigdata,sibiot,bi}
	}
	
	% ===================== OLED =====================
	\newglossaryentry{oled}{
		name={OLED},
		text={OLED},
		sort={oled},
		plural={OLEDs},
		first={diodo org{\'a}nico emisor de luz (\textit{Organic Light-Emitting Diode}, OLED)},
		firstplural={OLEDs},
		description={Un \textit{OLED} es una tecnolog{\'i}a de visualizaci{\'o}n basada en diodos org{\'a}nicos que emiten luz al ser excitados el{\'e}ctricamente. 
			A diferencia de las pantallas \gls{lcd}, los OLED no requieren retroiluminaci{\'o}n, lo que permite dise{\~n}os m{\'a}s delgados, mayor contraste, negros puros y mejor eficiencia energ{\'e}tica. 
			En \glspl{dispositivo} y sistemas \gls{iot}, las pantallas OLED son ampliamente utilizadas en \glspl{placadesarrollo}, wearables, sensores visuales y paneles de informaci{\'o}n debido a su bajo consumo y excelente visibilidad.},
		type=main,
		symbol={💡},
		user1={Pantallas de diodos org{\'a}nicos emisores de luz},
		user2={Mayor contraste y eficiencia que LCD},
		user3={Uso en IoT, wearables y microdispositivos},
		see={[Relacionado con:]{lcd,tft,dispositivo,placadesarrollo,iot}}
	}
	
	
	\newglossaryentry{oltp}{
		name={OLTP},
		text={OLTP},
		sort={oltp},
		plural={sistemas OLTP},
		first={procesamiento de transacciones en l{\'i}nea (\textit{Online Transaccion Processing}, OLTP)},
		firstplural={sistemas OLTP},
		description={Tipo de sistema de \gls{basedatos} y \glspl{aplicacion} orientadas a la gesti{\'o}n de grandes vol{\'u}menes de transacciones cortas en tiempo real. Los sistemas \emph{OLTP} se caracterizan por su alta concurrencia, baja latencia y capacidad de garantizar propiedades \gls{acid}. Son fundamentales en aplicaciones como banca, comercio electr{\'o}nico, telecomunicaciones y sistemas de reservas, donde se requiere exactitud y rapidez en la captura y procesamiento de datos. En el contexto de \glspl{sibiot}, OLTP puede emplearse para registrar interacciones en tiempo real entre \glspl{dispositivo}, \glspl{usuario} y \glspl{serviciodistribuido}.},
		type=main,
		symbol={💳},
		user1={Procesamiento transaccional},
		user2={Sistemas de transacciones en tiempo real},
		user3={Bases de datos operacionales},
		see={basedatos,consistencia,procesamientodatos}
	}
	
	\newglossaryentry{ondasonora}{
		name={Onda sonora},
		text={onda sonora},
		sort={onda sonora},
		plural={ondas sonoras},
		first={onda sonora},
		firstplural={ondas sonoras},
		description={Tipo de \gls{fenomenofisico} consistente en una perturbaci{\'o}n mec{\'a}nica que se propaga a trav{\'e}s de un medio material (como aire, agua o s{\'o}lidos) en forma de variaciones de presi{\'o}n y densidad, perceptibles por el o{\'i}do humano o registrables mediante \glspl{sensor} especializados. Una \emph{onda sonora} se caracteriza por propiedades como frecuencia, amplitud, longitud de onda y velocidad de propagaci{\'o}n. En el contexto de los \glspl{sibiot}, las ondas sonoras son \glspl{variablefisica} utilizadas en sistemas de reconocimiento de voz, detecci{\'o}n ac{\'u}stica de eventos, control por comandos verbales, monitorizaci{\'o}n ambiental y aplicaciones de seguridad o confort en entornos inteligentes.},
		type=main,
		symbol={🔊},
		user1={Perturbaci{\'o}n mec{\'a}nica en un medio material},
		user2={Se{\~n}al ac{\'u}stica perceptible o medible},
		user3={Dato de entrada en sistemas IoT},
		see={fenomenofisico,variablefisica,sensor,sibiot}
	}
	
	\newglossaryentry{optimizacion}{
		name={Optimizaci{\'o}n},
		text={optimizaci{\'o}n},
		sort={optimizacion},
		plural={optimizaciones},
		first={optimizaci{\'o}n},
		firstplural={optimizaciones},
		description={Proceso de mejorar un \gls{sistema}, algoritmo o recurso para alcanzar el mejor desempe{\~n}o posible seg{\'u}n criterios establecidos, como minimizar costes, reducir tiempos, maximizar la \gls{eficiencia} o incrementar el \gls{rendimiento}. La \emph{optimizaci{\'o}n} puede involucrar t{\'e}cnicas matem{\'a}ticas, heur{\'i}sticas o de \gls{ia}, y se aplica tanto en el dise{\~n}o de \glspl{algoritmo} como en la configuraci{\'o}n de procesos y recursos. En el contexto de los \glspl{sibiot}, la optimizaci{\'o}n se refleja en el uso eficiente de \glspl{sensor}, la distribuci{\'o}n de cargas en \glspl{nodo}, el ahorro energ{\'e}tico y la mejora de la \gls{eficienciaoperativa}, garantizando \gls{sostenibilidad} y \gls{escalabilidad}.},
		type=main,
		symbol={📈},
		user1={Mejora de recursos y procesos},
		user2={T{\'e}cnica para maximizar eficiencia y rendimiento},
		user3={Aplicaci{\'o}n en SIbIoT e ingenier{\'i}a de software},
		see={eficiencia,rendimiento,eficienciaoperativa,sostenibilidad,algoritmo,sibiot}
	}
	
		\newglossaryentry{optimizacionprocesos}{
		name={Optimizaci{\'o}n de procesos},
		text={optimizaci{\'o}n de procesos},
		sort={optimizacion procesos},
		plural={estrategias de optimizaci{\'o}n de procesos},
		first={optimizaci{\'o}n de procesos},
		firstplural={estrategias de optimizaci{\'o}n de procesos},
		description={Conjunto de acciones orientadas a mejorar la eficiencia, efectividad y rendimiento de procesos organizacionales. En los \glspl{sibiot}, la \emph{optimizaci{\'o}n de procesos} se sustenta en el an{\'a}lisis de datos provenientes de sensores y sistemas operativos.},
		type=main,
		symbol={⚙️},
		user1={Mejora continua},
		user2={Eficiencia operativa},
		user3={An{\'a}lisis de procesos},
		see={tomadecisionesestrategicas,integraciondatosoperativos}
	}
	
	\newglossaryentry{oracledatabase}{
		name={Oracle Database},
		text={Oracle Database},
		sort={oracle database},
		plural={Oracle Databases},
		first={Oracle Database},
		firstplural={Oracle Databases},
		description={Sistema de gesti{\'o}n de \glspl{basedatosrelacional} desarrollado por Oracle Corporation, ampliamente utilizado en entornos empresariales cr{\'i}ticos. \emph{Oracle Database} destaca por sus caracter{\'i}sticas avanzadas de \gls{seguridad}, replicaci{\'o}n, \gls{toleranciafallos} y optimizaci{\'o}n de rendimiento. Permite trabajar en arquitecturas distribuidas, con soporte para Big Data y an{\'a}lisis avanzado. En el contexto de los \glspl{sibiot}, se emplea en soluciones que requieren alta \gls{disponibilidad}, escalabilidad masiva y gesti{\'o}n de datos sensibles en sectores como salud, finanzas, industria o \glspl{ciudadinteligente}.},
		type=main,
		symbol={🗃️},
		user1={Base de datos empresarial robusta},
		user2={Alta seguridad y tolerancia a fallos},
		user3={Uso en aplicaciones IoT cr{\'i}ticas},
		see={basedatosrelacional,seguridad,disponibilidad,sibiot}
	}
	
	\newglossaryentry{ordentrabajo}{
		name={Orden de trabajo},
		text={orden de trabajo},
		sort={orden de trabajo},
		plural={\'ordenes de trabajo},
		first={orden de trabajo},
		firstplural={\'ordenes de trabajo},
		description={Una \emph{orden de trabajo} es un registro formal que organiza y documenta la ejecuci{\'o}n de tareas sobre activos y procesos (p.\,ej., inspecci{\'o}n, ajuste, reparaci{\'o}n, calibraci{\'o}n), incluyendo responsables, evidencia t{\'e}cnica, tiempos y resultados. En entornos donde un \gls{cps} se integra con capas empresariales, la orden de trabajo enlaza eventos del proceso (alarmas, degradaci{\'o}n, fallas) con procedimientos organizacionales, permitiendo auditor{\'\i}a, seguimiento y aprendizaje operacional a partir del historial de acciones y mediciones.},
		type=main,
		symbol={🧾},
		user1={Gesti{\'o}n formal de mantenimiento},
		user2={Evidencia t{\'e}cnica y seguimiento},
		user3={Vinculaci{\'o}n evento--acci{\'o}n},
		see={mantenimiento,auditoria,activo,incidente}
	}
	
	\newglossaryentry{organizacion}{
		name={Organizaci{\'o}n},
		text={organizaci{\'o}n},
		sort={organizacion},
		plural={organizaciones},
		first={organizaci{\'o}n},
		firstplural={organizaciones},
		description={Una organizaci{\'o}n representa un conjunto estructurado de personas, procesos y recursos orientados a la consecuci{\'o}n de metas espec{\'i}ficas. En el contexto de los \glspl{sibiot}, la organizaci{\'o}n integra sistemas inteligentes para optimizar la toma de decisiones, automatizar procesos y gestionar datos en tiempo real. Adem{\'a}s, su estructura favorece la interoperabilidad entre \glspl{dispositivo}, \glspl{servicio} y niveles de \gls{gestionempresarial}, impulsando la eficiencia y la adaptabilidad tecnol{\'o}gica.},
		type=\glsdefaulttype,
		symbol={🏢},
		user1={Organization},
		user2={Estructura coordinada de recursos y procesos},
		user3={Entidad que integra SIbIoT para optimizar la gesti{\'o}n},
		see={gestionempresarial,sibiot,procesamiento,automatizacion}
	}
	
	\newglossaryentry{orientacion}{
		name={Orientaci{\'o}n},
		text={orientaci{\'o}n},
		sort={orientacion},
		plural={orientaciones},
		first={orientaci{\'o}n},
		firstplural={orientaciones},
		description={Direcci{\'o}n conceptual que determina c{\'o}mo se priorizan y organizan las decisiones y actividades durante el desarrollo o an{\'a}lisis de un sistema. Una orientaci{\'o}n puede centrarse, por ejemplo, en modelos, servicios, datos, eventos, seguridad o calidad, y condiciona qu{\'e} artefactos se consideran primarios, qu{\'e} criterios de verificaci{\'o}n se aplican y c{\'o}mo se gestiona la evoluci{\'o}n. En \glspl{sibiot}, la orientaci{\'o}n influye en la integraci{\'o}n de componentes heterog{\'e}neos, el dise{\~n}o de interfaces y la trazabilidad desde el mundo f{\'\i}sico hacia servicios digitales.},
		type=main,
		symbol={🎛️},
		user1={Enfoque rector del proceso},
		user2={Prioridades de dise{\~n}o y calidad},
		user3={Impacto en artefactos y trazabilidad},
		see={filosofia,estructurametodologica,decisionarquitectonica,integracion,modelo}
	}
	
	\newglossaryentry{orientadoeventos}{
		name={orientado a eventos},
		sort={orientado a eventos},
		description={%
			Paradigma de dise{\~n}o de sistemas en el que la generaci{\'o}n, detecci{\'o}n y gesti{\'o}n de eventos constituye el eje central del comportamiento del sistema.
			En un sistema orientado a eventos, los componentes reaccionan a la ocurrencia de eventos significativos —tales como cambios de estado, acciones de usuario o se{\~n}ales del entorno— sin depender de flujos de control secuenciales o invocaciones directas.
			Este paradigma favorece el desacoplamiento, la escalabilidad y la adaptabilidad, siendo especialmente adecuado para entornos distribuidos, sistemas IoT y arquitecturas reactivas donde los eventos se producen de manera concurrente y no predecible.
		},
		first={orientado a eventos},
		plural={orientados a eventos},
		firstplural={orientados a eventos},
		see={enfoquereactivo,mensajeriatiemporeal,arquitecturadistribuida,iot}
	}
	
	\newglossaryentry{orquestacion}{
		name={Orquestaci{\'o}n},
		text={orquestaci{\'o}n},
		sort={orquestacion},
		first={orquestaci{\'o}n},
		description={Coordinaci{\'o}n automatizada de componentes y servicios para ejecutar flujos de trabajo, aprovisionar recursos y gestionar despliegues en sistemas distribuidos o en la nube.},
		type=main,
		symbol={🎼},
		user1={Coordinaci{\'o}n},
		user2={Automatizaci{\'o}n},
		user3={Servicios},
		see={elasticidad,impl}
	}
	
	\newglossaryentry{oscore}{
		name={OSCORE},
		text={OSCORE},
		plural={OSCORE},
		firstplural={OSCORE},
		sort={oscore},
		first={Seguridad de objetos para entornos REST limitados (\textit{Object Security for Constrained RESTful Environments}, OSCORE)},
		description={Mecanismo de seguridad definido por el IETF que proporciona confidencialidad, integridad y autenticaci{\'o}n a nivel de objeto para mensajes del protocolo CoAP. OSCORE protege el contenido de los mensajes de forma independiente de la capa de transporte, lo que lo hace especialmente adecuado para sistemas \gls{iot} y entornos con recursos limitados. Este mecanismo se basa en el uso de \gls{cbor} y \gls{cose} para la protecci{\'o}n criptogr{\'a}fica de los datos, y suele emplearse conjuntamente con protocolos de establecimiento de claves como \gls{edhoc}.},
		type=main,
		see={cbor,cose,edhoc,coap,iot}
	}
	
	\newglossaryentry{osi}{
		name={OSI},
		text={OSI},
		sort={osi},
		plural={OSI},
		first={Interconexi{\'o}n de Sistemas Abiertos (\textit{Open Systems Interconnection}, OSI)},
		firstplural={OSI},
		description={El \emph{modelo OSI} es un marco de referencia est{\'a}ndar propuesto por la ISO para describir la arquitectura de las \glspl{red} de comunicaci{\'o}n en siete capas jer{\'a}rquicas: f{\'i}sica, enlace de datos, red, transporte, sesi{\'o}n, presentaci{\'o}n y aplicaci{\'o}n. Cada \gls{capa} cumple funciones espec{\'i}ficas y se comunica con las capas adyacentes, lo que facilita la \gls{interoperabilidad}, el dise{\~n}o modular y la estandarizaci{\'o}n de protocolos de \gls{sistemacomunicacion}. En el contexto de \gls{iot}, el modelo OSI es utilizado como gu{\'i}a conceptual para organizar protocolos, garantizar la compatibilidad entre dispositivos heterog{\'e}neos y asegurar la \gls{escalabilidad} de \glspl{arquitecturadistribuida}.},
		type=main,
		symbol={📡},
		user1={Modelo de referencia en siete capas},
		user2={Base para dise{\~n}o e interoperabilidad de redes},
		user3={Gu{\'i}a conceptual en IoT},
		see={sistemacomunicacion,red,interoperabilidad,escalabilidad}
	}
	
	\newglossaryentry{ot}{
		name={OT},
		text={OT},
		sort={ot},
		plural={OT},
		first={OT},
		firstplural={OT},
		description={Sigla de \emph{Operational Technology}. Conjunto de tecnolog{\'\i}as (equipos, software y redes) usadas para supervisar y controlar procesos f{\'\i}sicos en entornos industriales e infraestructura cr{\'\i}tica. En la evoluci{\'o}n hacia \glspl{cps} y \glspl{sibiot}, OT se integra con capas de datos y gobierno, lo que hace necesario tratar seguridad y operaci{\'o}n de forma coordinada.},
		type=main,
		symbol={🏭},
		user1={Control industrial},
		user2={Procesos f{\'\i}sicos},
		user3={Infraestructura cr{\'\i}tica},
		see={it,control,ciberseguridad}
	}
		% Panel solar
	\newglossaryentry{panelsolar}{
		name={Panel solar},
		type=main,
		text={panel solar},
		sort={panel solar},
		symbol={🔋},
		plural={paneles solares},
		first={panel solar},
		firstplural={paneles solares},
		description={Dispositivo fotovoltaico que convierte la radiaci{\'o}n solar en energ{\'i}a el{\'e}ctrica utilizable para alimentar sistemas electr{\'o}nicos, incluyendo dispositivos IoT}
	}
	
	\newglossaryentry{pan}{
		name={PAN},
		text={PAN},
		sort={pan},
		plural={PANs},
		first={red de {\'a}rea personal (\textit{Personal Area Network}, PAN)},
		firstplural={PANs},
		description={Tipo de red de corto alcance que permite la interconexi{\'o}n de dispositivos electr{\'o}nicos personales, como tel{\'e}fonos m{\'o}viles, tabletas, computadoras port{\'a}tiles, \glspl{sensor} y perif{\'e}ricos. Las \emph{PAN} operan en distancias menores a 10 metros y utilizan tecnolog{\'i}as como \gls{bluetooth}, \gls{zigbee} o infrarrojo. Son fundamentales en aplicaciones de \gls{iot} y en entornos ubicuos donde se requiere conectividad inmediata, segura y de bajo consumo energ{\'e}tico.},
		type=main,
		symbol={📡},
		user1={Red de corto alcance},
		user2={Conexi{\'o}n de dispositivos personales},
		user3={Soporte en IoT y entornos ubicuos},
		see={bluetooth,zigbee,iot,rediot,infraestructurared}
	}
	
	\newglossaryentry{paradigma}{
		name={Paradigma},
		text={paradigma},
		sort={paradigma},
		plural={paradigmas},
		first={paradigma},
		firstplural={paradigmas},
		description={Modelo, esquema o enfoque conceptual que define una manera particular de comprender, estructurar y resolver problemas en un determinado campo del conocimiento. En programaci{\'o}n y \gls{ingenieriasoftware}, un \emph{paradigma} establece principios, reglas y estilos de dise{\~n}o, como la programaci{\'o}n orientada a objetos, funcional o reactiva. En investigaci{\'o}n cient{\'i}fica, orienta las metodolog{\'i}as y criterios de validaci{\'o}n del conocimiento.},
		type=main,
		symbol={🔎},
		user1={Modelo o enfoque conceptual},
		user2={Marco de referencia en programaci{\'o}n o ingenier{\'i}a},
		user3={Gu{\'i}a metodol{\'o}gica en ciencia y tecnolog{\'i}a},
		see={modelo,ingenieriasoftware}
	}
	
	\newglossaryentry{patron}
	{
		name={Patr{\'o}n},
		text={patr{\'o}n},
		sort={patron},
		plural={patrones},
		first={patr{\'o}n},
		firstplural={patrones},
		description={Un \emph{patr{\'o}n} es un modelo o forma recurrente que sirve como gu{\'i}a, referencia o criterio de organizaci{\'o}n. Puede entenderse en distintos niveles:
			\begin{itemize}
				\item En un contexto general, como un esquema o regularidad observable en datos, procesos o fen{\'o}menos naturales.
				\item En ingenier{\'i}a de software, como una soluci{\'o}n probada y reutilizable a un problema recurrente de dise{\~n}o o arquitectura, por ejemplo, los patrones de dise{\~n}o o los patrones arquitect{\'o}nicos.
				\item En \glspl{sibiot}, como las tendencias o comportamientos detectados en los \glspl{dato} capturados por \glspl{sensor}, que permiten identificar anomal{\'i}as, realizar predicciones y apoyar la \gls{tomadecisiones}.
				\item En bases de datos o consultas, como una condici{\'o}n que permite identificar coincidencias espec{\'i}ficas (ej. patrones de texto o expresiones regulares).
			\end{itemize}
			Su utilidad radica en la capacidad de simplificar la complejidad, sistematizar soluciones y generar conocimiento a partir de la repetici{\'o}n de estructuras o comportamientos.},
		type=\glsdefaulttype,
		symbol={🔎},
		user1={Modelo},
		user2={Esquema recurrente},
		user3={Pattern},
		see={filtro}
	}
	
	\newglossaryentry{persistencia}{
		name={Persistencia},
		text={persistencia},
		short={persistencia},
		long={Persistencia de datos},
		sort={persistencia},
		plural={persistencias},
		first={persistencia},
		firstplural={persistencias},
		description={Propiedad de un sistema de informaci{\'o}n que garantiza que los datos generados, procesados o recibidos se almacenen de forma estable y duradera en un medio de almacenamiento, de manera que puedan ser recuperados de forma consistente y confiable a lo largo del tiempo, incluso ante fallos del sistema, reinicios o interrupciones de energ{\'i}a. En el contexto de sistemas distribuidos e \gls{iot}, la persistencia es fundamental para asegurar la trazabilidad hist{\'o}rica, el an{\'a}lisis posterior, la auditor{\'i}a de eventos y el soporte a la \gls{tomadecisiones}, siendo t{\'i}picamente implementada mediante sistemas gestores de bases de datos como \gls{postgresql} y mecanismos transaccionales que preservan la integridad y consistencia de la informaci{\'o}n.},
		see={basedatos,almacenamiento,consistencia}
	}
	
	\newglossaryentry{personalizacion}{
		name={Personalizaci{\'o}n},
		text={personalizaci{\'o}n},
		sort={personalizacion},
		first={personalizaci{\'o}n},
		description={Adaptaci{\'o}n de un sistema a necesidades, preferencias o contexto de un usuario u organizaci{\'o}n, mediante configuraciones, perfiles o reglas de negocio.},
		type=main,
		symbol={🧑‍🎨},
		user1={Adaptaci{\'o}n},
		user2={Preferencias},
		user3={Contexto},
		see={configuracion,usabilidad}
	}
	
	\newglossaryentry{pesaje}
	{
		name={Pesaje},
		text={pesaje},
		sort={pesaje},
		description={Proceso de determinaci{\'o}n cuantitativa del \gls{peso} de un objeto o sustancia mediante el uso de \glspl{sensor} o instrumentos de medici{\'o}n, como las \glspl{celda} de carga o los \glspl{sensor} de presi{\'o}n. 
			En el contexto de los \glspl{sibiot}, el pesaje automatizado constituye una funci{\'o}n esencial para sistemas de control de inventario, gesti{\'o}n log{\'i}stica y mantenimiento predictivo, donde los datos se procesan en tiempo real para la toma de decisiones aut{\'o}nomas. 
			El pesaje inteligente permite optimizar el consumo de recursos, garantizar la trazabilidad de productos y mejorar la eficiencia operativa en entornos industriales conectados.},
		plural={pesajes},
		symbol={},
		first={Pesaje},
		see=[Relacionado con:]{peso,celda,fuerza,sensor}
	}
	
	\newglossaryentry{peso}
	{
		name={Peso},
		text={peso},
		sort={peso},
		description={Magnitud f{\'i}sica derivada que representa la fuerza con la que la masa de un cuerpo es atra{\'i}da por el campo gravitatorio de la Tierra u otro cuerpo celeste. 
			Se calcula como el producto de la masa y la aceleraci{\'o}n de la gravedad (\ensuremath{P = m \cdot g}). 
			En el contexto de los \glspl{sibiot}, el peso se mide o estima mediante \glspl{sensor} como las \glspl{celda} de carga, los \glspl{sensor} de presi{\'o}n y los \glspl{sensor} de fuerza, permitiendo el control de procesos industriales, la monitorizaci{\'o}n de inventarios y la automatizaci{\'o}n de sistemas log{\'i}sticos. 
			El peso constituye una variable fundamental para el dise{\~n}o de sistemas de pesaje inteligente, control de transporte y optimizaci{\'o}n energ{\'e}tica en dispositivos inteligentes.},
		plural={pesos},
		symbol={P},
		first={Peso (P)},
		see=[Relacionado con:]{masa,fuerza,presion,celda}
	}
	
	\newglossaryentry{pgadmin}{
		name={pgAdmin},
		text={pgAdmin},
		plural={instancias pgAdmin},
		first={pgAdmin, una plataforma gr{\'a}fica utilizada para la administraci{\'o}n y supervisi{\'o}n de bases de datos PostgreSQL},
		firstplural={instancias pgAdmin empleadas para la gesti{\'o}n y monitorizaci{\'o}n de servidores PostgreSQL},
		description={Herramienta de administraci{\'o}n que proporciona una interfaz visual para gestionar usuarios, esquemas, consultas y rendimiento en PostgreSQL},
		sort={pgadmin},
		type=main,
		symbol={🐘},
		see={postgresql,sgbd}
	}
	
	\newglossaryentry{php}{
		name={PHP},
		text={PHP},
		plural={aplicaciones PHP},
		first={PHP, un lenguaje de programaci{\'o}n interpretado orientado al desarrollo de aplicaciones web del lado del servidor},
		firstplural={aplicaciones PHP utilizadas en sistemas web y plataformas empresariales},
		description={Lenguaje de programaci{\'o}n dise{\~n}ado para la generaci{\'o}n din{\'a}mica de contenido web e integraci{\'o}n con bases de datos en entornos cliente--servidor},
		sort={php},
		type=main,
		symbol={🐘},
		see={desarrolloweb,servidor}
	}
	
	\newglossaryentry{pipeline}{
		name={Pipeline},
		text={pipeline},
		sort={pipeline},
		plural={pipelines},
		first={pipeline},
		firstplural={pipelines},
		description={Estructura de procesamiento compuesta por una secuencia ordenada de etapas o fases interconectadas, en la que la salida de cada etapa constituye la entrada de la siguiente. Un \emph{pipeline} permite dividir un proceso complejo en pasos m{\'a}s simples, facilitando la paralelizaci{\'o}n, la reutilizaci{\'o}n de componentes y la mejora del rendimiento global. En el contexto de \gls{iot} y de los \glspl{sibiot}, los pipelines se emplean para el \gls{procesamientodatos}, incluyendo fases como ingesta, filtrado, transformaci{\'o}n, an{\'a}lisis, almacenamiento y visualizaci{\'o}n de datos provenientes de \glspl{sensor} y \glspl{dispositivoconectado}. Tambi{\'e}n son fundamentales en procesos de integraci{\'o}n continua, despliegue continuo y flujos ETL, donde garantizan trazabilidad, automatizaci{\'o}n y consistencia operativa},
		type=main,
		symbol={🔗},
		user1={Procesamiento secuencial por etapas},
		user2={Flujos automatizados de datos o tareas},
		user3={Base de ETL, CI/CD y anal{\'i}tica en IoT},
		see={etl,procesamientodatos,integracion,automatizacion,sibiot}
	}
	
	\newglossaryentry{pki}{
		name={PKI},
		text={PKI},
		sort={pki},
		first={infraestructura de clave p{\'u}blica (\emph{Public Key Infrastructure}, PKI)},
		description={\emph{Public Key Infrastructure}. Conjunto de pol{\'i}ticas, roles y servicios criptogr{\'a}ficos para gestionar certificados y claves p{\'u}blicas, habilitando autenticaci{\'o}n y comunicaciones seguras.},
		type=main,
		symbol={🔑},
		user1={Certificados},
		user2={Claves p{\'u}blicas},
		user3={Autenticaci{\'o}n},
		see={gestionconfianza,seguridad}
	}
	
	\newglossaryentry{placa}{
		name={Placa},
		text={placa},
		sort={placa},
		plural={placas},
		first={placa},
		firstplural={placas},
		type=main,
		description={Soporte f{\'i}sico de material diel{\'e}ctrico que integra pistas conductoras y componentes electr{\'o}nicos interconectados, dise{\~n}ado para ejecutar funciones espec{\'i}ficas dentro de un circuito o \gls{sistema}. Las \emph{placas} pueden alojar \glspl{microcontrolador}, \glspl{sensor}, \glspl{actuador} o m{\'o}dulos de comunicaci{\'o}n, siendo fundamentales en dispositivos electr{\'o}nicos y \gls{sibiot}.},
		symbol={🧩},
		user1={Soporte de componentes y pistas},
		user2={Base de hardware para sistemas embebidos},
		user3={Integraci{\'o}n de sensores, actuadores y comunicaciones},
		see={placadesarrollo,placacontroladora,placaprocesamiento,microcontrolador,sensor,actuador,iot}
	}
	
	\newglossaryentry{placacontroladora}{
		name={Placa controladora},
		text={placa controladora},
		sort={placa controladora},
		plural={placas controladoras},
		first={placa controladora},
		firstplural={placas controladoras},
		type=main,
		description={Dispositivo electr{\'o}nico que integra circuiter{\'i}a, \glspl{microcontrolador}, entradas y salidas para gestionar y coordinar el funcionamiento de componentes externos, tales como \glspl{sensor}, \glspl{actuador}, motores o interfaces de comunicaci{\'o}n. Com{\'u}n en automatizaci{\'o}n, rob{\'o}tica y \gls{iot}. Ejemplos representativos incluyen \textit{\Gls{arduino}}, \textit{\Gls{raspberrypi}} y \textit{\gls{esp32}}.},
		symbol={🛠️},
		user1={Gesti{\'o}n de E/S y perif{\'e}ricos},
		user2={Control de procesos en tiempo real},
		user3={Papel central en automatizaci{\'o}n e IoT},
		see={placa,microcontrolador,sensor,actuador,esp32,arduino}
	}
	
	\newglossaryentry{placadesarrollo}{
		name={Placa de desarrollo},
		text={placa de desarrollo},
		sort={placa de desarrollo},
		plural={placas de desarrollo},
		first={placa de desarrollo},
		firstplural={placas de desarrollo},
		type=main,
		description={Dispositivo electr{\'o}nico programable utilizado para el dise{\~n}o, prueba y desarrollo de sistemas embebidos, especialmente en proyectos de \gls{iot}. Integra \glspl{microcontrolador} o \glspl{microprocesador} y facilita la conexi{\'o}n con \glspl{sensor}, \glspl{actuador} y perif{\'e}ricos. Ejemplos: \textit{Arduino}, \textit{Raspberry Pi}, \textit{ESP32}. Destaca por bajo costo, \gls{flexibilidad} y amplia documentaci{\'o}n.},
		symbol={🔬},
		user1={Prototipado r{\'a}pido de sistemas embebidos},
		user2={Conectividad sencilla a sensores y actuadores},
		user3={Amplio ecosistema de bibliotecas},
		see={placa,placacontroladora,placaprocesamiento,microcontrolador,iot}
	}
	
	\newglossaryentry{placaprocesamiento}{
		name={Placa de procesamiento},
		text={placa de procesamiento},
		sort={placa de procesamiento},
		plural={placas de procesamiento},
		first={placa de procesamiento},
		firstplural={placas de procesamiento},
		type=main,
		description={Placa que integra uno o varios elementos de c{\'o}mputo (p.\,ej., \glspl{microcontrolador}, \glspl{microprocesador} o SoC), junto con memorias, buses e interfaces digitales. Permite ejecutar instrucciones, controlar perif{\'e}ricos y procesar datos en \gls{tiemporeal}; esencial en \gls{iot}, rob{\'o}tica, automatizaci{\'o}n y computaci{\'o}n embebida.},
		symbol={⚙️},
		user1={Ejecuci{\'o}n y control en tiempo real},
		user2={Integraci{\'o}n de CPU/MCU, memoria y buses},
		user3={Base para aplicaciones embebidas e IoT},
		see={placa,placadesarrollo,microcontrolador,iot}
	}
	
	\newglossaryentry{planificacion}{
		name={Planificaci{\'o}n},
		text={planificaci{\'o}n},
		sort={planificacion},
		plural={planificaciones},
		first={planificaci{\'o}n},
		firstplural={planificaciones},
		description={La \gls{planificacion} constituye una funci{\'o}n esencial en la \gls{gestion} de \glspl{sistema} y organizaciones. En los \glspl{sibiot}, la planificaci{\'o}n implica coordinar la adquisici{\'o}n, procesamiento y an{\'a}lisis de datos para anticipar necesidades, optimizar recursos y garantizar la continuidad operativa. Asimismo, permite alinear los procesos automatizados con metas estrat{\'e}gicas mediante el uso de modelos predictivos y anal{\'i}ticos que facilitan la \gls{tomadecisionesinformadas}.},
		type=\glsdefaulttype,
		symbol={🗓️},
		user1={Planning},
		user2={Definici{\'o}n estructurada de objetivos y acciones},
		user3={Gesti{\'o}n anticipada de recursos y procesos en SIbIoT},
		see={gestionempresarial,tomadecisionesinformadas,sibiot,analisispredictivo}
	}
	
	\newglossaryentry{planificacionagil}{
		name={Planificaci{\'o}n {\'a}gil},
		text={planificaci{\'o}n {\'a}gil},
		sort={planificacion agil},
		plural={planificaciones {\'a}giles},
		first={planificaci{\'o}n {\'a}gil},
		firstplural={planificaciones {\'a}giles},
		type=main,
		description={Proceso iterativo y flexible de organizaci{\'o}n del trabajo que adapta continuamente alcance, tiempos y recursos en funci{\'o}n de retroalimentaci{\'o}n, colaboraci{\'o}n con el cliente y evoluci{\'o}n de requisitos. En proyectos \gls{iot}, facilita la integraci{\'o}n progresiva de hardware y \gls{software} en ciclos cortos de validaci{\'o}n, favoreciendo \gls{eficienciaoperativa} y \gls{adaptabilidad}.},
		symbol={🗂️},
		user1={Iteraciones cortas y entrega incremental},
		user2={Priorizaci{\'o}n por valor y feedback continuo},
		user3={Integraci{\'o}n de hardware y software en IoT},
		see={equipoagil,ingenieriasoftware,gestionempresarial}
	}
	
	\newglossaryentry{plataforma}{
		name={Plataforma},
		text={plataforma},
		sort={plataforma},
		plural={plataformas},
		first={plataforma},
		firstplural={plataformas},
		type=main,
		description={Entorno t{\'e}cnico o conjunto de tecnolog{\'i}as que proporciona una base para desarrollar, ejecutar o distribuir \glspl{aplicacion}, \glspl{servicio} o \glspl{sistema}. Puede referirse a sistemas operativos, entornos de ejecuci{\'o}n, infraestructuras en la \gls{nube}, o plataformas de desarrollo e integraci{\'o}n.},
		symbol={🧩},
		user1={Base tecnol{\'o}gica de ejecuci{\'o}n y desarrollo},
		user2={Servicios compartidos y APIs},
		user3={Soporte a integraci{\'o}n y despliegue},
		see={nube,aplicacion,servicio,infraestructurared}
	}
	
	\newglossaryentry{plataformaweb}{
		name={Plataforma web},
		text={plataforma web},
		sort={plataforma web},
		plural={plataformas web},
		first={plataforma web},
		firstplural={plataformas web},
		description={Conjunto de tecnolog{\'i}as y servicios que permiten desarrollar, desplegar y operar aplicaciones accesibles mediante navegadores, integrando interfaz, l{\'o}gica y datos a trav{\'e}s de la red.},
		type=main,
		symbol={🌍},
		user1={Aplicaciones web},
		user2={Servicios},
		user3={Navegador},
		see={sistemanavegacion,controller,servicioweb}
	}
	
	\newglossaryentry{plugin}{
		name={Plugin},
		text={plugin},
		sort={plugin},
		plural={plugins},
		first={plugin},
		firstplural={plugins},
		type=main,
		description={M{\'o}dulo o complemento de \gls{software} que extiende o a{\~n}ade funcionalidades a una aplicaci{\'o}n o \gls{plataforma} sin modificar su n{\'u}cleo. Facilita personalizaci{\'o}n, mantenibilidad y \gls{interoperabilidad} mediante interfaces bien definidas.},
		symbol={🧱},
		user1={Extensi{\'o}n de funcionalidades},
		user2={Integraci{\'o}n por interfaces o APIs},
		user3={Desacoplamiento del n{\'u}cleo de la aplicaci{\'o}n},
		see={aplicacion,plataforma,api}
	}
	
	\newglossaryentry{pmbok}{
		name={PMBOK},
		text={PMBOK},
		sort={pmbok},
		first={gu{\'i}a de los fundamentos para la direcci{\'o}n de proyectos (\textit{Project Management Body of Knowledge}, PWBOK)},
		plural={PMBOK},
		type=main,
		description={Est{\'a}ndar del \textit{Project Management Institute} (PMI) que recopila buenas pr{\'a}cticas, procesos, herramientas y t{\'e}cnicas para la gesti{\'o}n efectiva de proyectos. Proporciona una gu{\'i}a para planificar, ejecutar, monitorear y cerrar proyectos en ingenier{\'i}a, tecnolog{\'i}a y sistemas complejos como los basados en \gls{iot}, estandarizando procesos, mitigando riesgos y mejorando la eficiencia.},
		symbol={📘},
		user1={Marco de buenas pr{\'a}cticas de gesti{\'o}n},
		user2={Procesos para planificar y controlar proyectos},
		user3={Aplicable a proyectos IoT y multidisciplinarios},
		see={gestionempresarial,planificacionagil}
	}
	
	\newglossaryentry{politica}{
		name={Pol{\'i}tica},
		text={pol{\'i}tica},
		sort={politica},
		plural={pol{\'i}ticas},
		first={pol{\'i}tica},
		firstplural={pol{\'i}ticas},
		description={Conjunto de reglas, directrices o principios establecidos para orientar el comportamiento, la configuraci{\'o}n o la toma de decisiones en un sistema o instituci{\'o}n. En \glspl{sibiot}, las pol{\'i}ticas definen criterios de \gls{seguridad}, privacidad, \gls{gestionenergetica} y \gls{provisionamiento}, garantizando coherencia y cumplimiento normativo.},
		type=main,
		symbol={📜},
		user1={Directriz o regla formal},
		user2={Orientaci{\'o}n de gesti{\'o}n y control},
		user3={Marco de referencia organizativo},
		see={seguridad,privacidad,provisionamiento,gestionenergetica}
	}
	
	\newglossaryentry{politicaorganizacional}{
		name={Pol{\'i}tica organizacional},
		text={pol{\'i}tica organizacional},
		sort={politica organizacional},
		first={pol{\'i}tica organizacional},
		description={Conjunto de normas, directrices y criterios establecidos por una organizaci{\'o}n para regular procesos, responsabilidades, seguridad y toma de decisiones en sus sistemas.},
		type=main,
		symbol={🏛️},
		user1={Directrices},
		user2={Normas internas},
		user3={Gobernanza},
		see={gestion,seguridad}
	}
	
	\newglossaryentry{pop3}{
		name={POP3},
		text={POP3},
		sort={pop3},
		firstplural={protocolos POP3},
		plural={protocolos POP3},
		first={protocolo de oficina de correo (\textit{Post Office Protocol} versi{\'o}n 3, POP3)},
		type=main,
		description={Protocolo de aplicaci{\'o}n para recuperar mensajes de correo electr{\'o}nico desde un servidor remoto. \textit{POP3} permite descargar mensajes al cliente y, por defecto, eliminarlos del servidor. Aunque menos vers{\'a}til que IMAP, su simplicidad y bajo uso de recursos puede resultar adecuada en escenarios embebidos o de \gls{iot}.},
		symbol={✉️},
		user1={Recuperaci{\'o}n de correo desde servidor},
		user2={Modelo de descarga local de mensajes},
		user3={Bajo consumo de recursos},
		see={http,ftp,seguridad}
	}
	
	\newglossaryentry{portabilidad}{
		name={Portabilidad},
		text={portabilidad},
		sort={portabilidad},
		plural={portabilidades},
		first={portabilidad},
		firstplural={portabilidades},
		description={Propiedad de un software, modelo, componente o sistema que permite su traslado y ejecuci{\'o}n (o reutilizaci{\'o}n) en diferentes plataformas, entornos de ejecuci{\'o}n o configuraciones de infraestructura con cambios m{\'\i}nimos. La portabilidad puede referirse a niveles como: (i) portabilidad de c{\'o}digo (entre sistemas operativos, arquitecturas CPU o runtimes), (ii) portabilidad de modelos (entre herramientas MDE) y (iii) portabilidad de despliegue (entre nubes o entornos virtualizados/contenedorizados). En \gls{iot}/\glspl{sibiot}, la portabilidad se ve afectada por la heterogeneidad de \glspl{dispositivoembebido}, bibliotecas, drivers y protocolos, y se favorece mediante est{\'a}ndares, capas de abstracci{\'o}n, interfaces bien definidas y automatizaci{\'o}n del pipeline de construcci{\'o}n y despliegue.},
		type=main,
		symbol={🧳},
		user1={Independencia de plataforma},
		user2={Reutilizaci{\'o}n multi-entorno},
		user3={Mitigaci{\'o}n de dependencia tecnol{\'o}gica},
		see={interoperabilidad,compatibilidad,entornovirtualizado,devops,pipeline,escalabilidad}
	}
	
	\newglossaryentry{portalweb}{
		name={Portal web},
		text={portal web},
		sort={portal web},
		plural={portales web},
		first={portal web},
		firstplural={portales web},
		type=main,
		description={Punto de acceso unificado donde una \gls{red} inform{\'a}tica ofrece \glspl{servicio} y recursos de manera integrada a trav{\'e}s de un \gls{navegador}. Un \emph{portal web} centraliza contenidos, autenticaci{\'o}n y personalizaci{\'o}n para usuarios, y en \gls{sibiot} puede actuar como interfaz de \gls{monitorizacion}, \gls{control} y \gls{visualizacion} de datos.},
		symbol={🌐},
		user1={Acceso integrado a servicios y contenidos},
		user2={Interfaz web mediante navegador},
		user3={Paneles e integraciones para SIbIoT},
		see={aplicacionweb,navegador,sibiot}
	}
	
	\newglossaryentry{post}{
		name={POST},
		text={POST},
		sort={post},
		first={POST},
		description={M{\'e}todo del protocolo HTTP empleado para enviar datos al servidor con el fin de crear nuevos recursos o ejecutar operaciones que modifican el estado del sistema. A diferencia de \texttt{GET}, el m{\'e}todo \texttt{POST} no es idempotente y puede producir efectos secundarios. En el contexto de APIs RESTful y SIbIoT, se utiliza para registrar eventos, almacenar datos de sensores, iniciar procesos y enviar comandos que desencadenan acciones en sistemas distribuidos},
		type=main,
		see={http,rest,restful,api,evento,procesamientodistribuido}
	}
	
	\newglossaryentry{postgresql}{
		name={PostgreSQL},
		text={PostgreSQL},
		sort={postgresql},
		plural={PostgreSQL},
		first={PostgreSQL},
		firstplural={PostgreSQL},
		description={Sistema avanzado de gesti{\'o}n de \glspl{basedatosrelacional} y orientado a objetos, de c{\'o}digo abierto, reconocido por su robustez, extensibilidad y cumplimiento estricto de est{\'a}ndares. \emph{PostgreSQL} soporta tipos de datos personalizados, funciones avanzadas, transacciones ACID y procesamiento de consultas complejas. En el contexto de los \glspl{sibiot}, resulta adecuado para manejar grandes vol{\'u}menes de datos estructurados y semi-estructurados, integrarse con servicios distribuidos y garantizar \gls{fiabilidad}, \gls{consistencia} y \gls{escalabilidad}.},
		type=main,
		symbol={🐘},
		user1={Base de datos avanzada y extensible},
		user2={Soporte para datos estructurados y semi-estructurados},
		user3={Aplicaciones en sistemas IoT distribuidos},
		see={basedatosrelacional,basedatosdistribuida,consistencia,sibiot}
	}
	
	\newglossaryentry{practicaagil}{
		name={Pr{\'a}ctica {\'a}gil},
		text={pr{\'a}ctica {\'a}gil},
		sort={practica agil},
		plural={pr{\'a}cticas {\'a}giles},
		first={pr{\'a}ctica {\'a}gil},
		firstplural={pr{\'a}cticas {\'a}giles},
		description={Acci{\'o}n concreta, t{\'e}cnica o din{\'a}mica recurrente que implementa uno o varios \glspl{principioagil}. Las \emph{pr{\'a}cticas {\'a}giles} incluyen planificaci{\'o}n de iteraciones, reuniones diarias (\textit{daily stand-ups}), retrospectivas, integraci{\'o}n continua, pruebas automatizadas y entrega incremental. En proyectos \gls{iot}, permiten coordinar equipos de hardware y software, promoviendo la entrega frecuente de valor funcional y la mejora continua del producto.},
		type=main,
		symbol={⚡},
		user1={Acciones recurrentes de metodolog{\'i}as {\'a}giles},
		user2={Integraci{\'o}n continua y entrega incremental},
		user3={Mejora continua en proyectos IoT},
		see={principioagil,planificacionagil,equipoagil}
	}
	
	\newglossaryentry{precision}{
		name={Precisi{\'o}n},
		text={precisi{\'o}n},
		sort={precision},
		plural={precisiones},
		first={precisi{\'o}n},
		firstplural={precisiones},
		description={Grado de consistencia o repetibilidad de un conjunto de mediciones realizadas sobre una misma \gls{variablefisica} bajo condiciones constantes. La \emph{precisi{\'o}n} indica cu{\'a}nto se aproximan entre s{\'i} los valores obtenidos, aunque no necesariamente coincidan con el valor real, lo cual corresponde a la \gls{exactitud}. En el contexto de \glspl{sensor} y \glspl{sibiot}, la precisi{\'o}n es un par{\'a}metro fundamental para garantizar la fiabilidad de los datos, la validez de la \gls{monitorizacion} y la calidad de la \gls{tomadecisiones} en tiempo real, especialmente en aplicaciones cr{\'i}ticas como salud, industria o \gls{transporte}.},
		type=main,
		symbol={🎯},
		user1={Consistencia de mediciones},
		user2={Repetibilidad de resultados},
		user3={Par{\'a}metro clave en calidad de datos IoT},
		see={exactitud,variablefisica,sensor,sibiot}
	}
	
	\newglossaryentry{preciso}
	{
		name={Preciso},
		text={preciso},
		sort={preciso},
		plural={precisos},
		first={preciso},
		firstplural={precisos},
		description={T{\'e}rmino adjetivo que se refiere a aquello que posee o mantiene \gls{precision}. En el contexto de los \glspl{sibiot}, se aplica a un \gls{sensor} o \gls{sistema} cuya informaci{\'o}n refleja de forma fidedigna los valores reales.},
		type=\glsdefaulttype,
		symbol={},
		user1={Exacto},
		user2={Correcto},
		user3={Riguroso},
		see={precision}
	}
	
	\newglossaryentry{prediccion}{
		name={Predicci{\'o}n},
		text={predicci{\'o}n},
		sort={prediccion},
		plural={predicciones},
		first={predicci{\'o}n},
		firstplural={predicciones},
		description={Proceso de estimar o anticipar valores futuros de una variable con base en datos hist{\'o}ricos, patrones observados o modelos matem{\'a}ticos. En el contexto de la \gls{iot} y los \glspl{sibiot}, la \emph{predicci{\'o}n} se utiliza para prever comportamientos, demandas energ{\'e}ticas, condiciones ambientales o estados de \glspl{dispositivo}, permitiendo activar mecanismos de \gls{control} preventivo y optimizar la \gls{tomadecisiones}.},
		type=main,
		symbol={🔮},
		user1={Estimaci{\'o}n basada en datos},
		user2={Anticipaci{\'o}n de eventos futuros},
		user3={Soporte para la toma de decisiones},
		see={analitica,modelopredictivo,control,tomadecisiones}
	}
	
	\newglossaryentry{preprocesamiento}{
		name={Preprocesamiento},
		sort={preprocesamiento},
		text={preprocesamiento},
		description={%
			Conjunto de operaciones iniciales aplicadas a los datos brutos con el objetivo de mejorar su calidad, consistencia y utilidad antes de ser sometidos a etapas posteriores de procesamiento, an{\'a}lisis o toma de decisiones.
			El \emph{preprocesamiento} puede incluir tareas como validaci{\'o}n de rangos, eliminaci{\'o}n de valores at{\'i}picos, filtrado de ruido, normalizaci{\'o}n, agregaci{\'o}n temporal y detecci{\'o}n de inconsistencias.
			En sistemas de informaci{\'o}n basados en IoT y SIbIoT, el preprocesamiento se realiza frecuentemente en nodos de borde o pasarelas, permitiendo reducir la latencia, optimizar el uso del ancho de banda y garantizar que solo informaci{\'o}n relevante y confiable sea transmitida a las capas superiores del sistema.
		},
		first={preprocesamiento},
		plural={preprocesamientos},
		firstplural={preprocesamientos},
		see={procesamientodistribuido,edgecomputing,filtradodatos,normalizaciondatos,iot}
	}
	
	\newglossaryentry{presion}
	{
		name={Presi{\'o}n},
		text={presi{\'o}n},
		sort={presion},
		description={Magnitud f{\'i}sica que describe la fuerza ejercida por unidad de {\'a}rea sobre una superficie. 
			En contextos f{\'i}sicos, representa la acci{\'o}n de un fluido (l{\'i}quido o gas) sobre las paredes del recipiente que lo contiene. 
			En el {\'a}mbito industrial, se refiere a la variable cr{\'i}tica que permite controlar procesos hidr{\'a}ulicos, neum{\'a}ticos o t{\'e}rmicos. 
			En el contexto de la ingenier{\'i}a electr{\'o}nica y los \glspl{sensor}, la presi{\'o}n se convierte en una se{\~n}al el{\'e}ctrica mediante transductores, lo que posibilita la detecci{\'o}n de objetos, el control de fuerza o la medici{\'o}n de condiciones ambientales en sistemas inteligentes y en aplicaciones de \gls{sibiot}.},
		plural={presiones},
		symbol={Pₚ}
		first={Presi{\'o}n (P)},
		see=[Relacionado con:]{fuerza,fluido,sensor}
	}
	
	\newglossaryentry{principio}{
		name={Principio},
		text={principio},
		sort={principio},
		plural={principios},
		first={principio},
		firstplural={principios},
		description={Enunciado fundamental que act{\'u}a como regla de orientaci{\'o}n para dise{\~n}ar, construir, evaluar o gestionar sistemas. Los principios establecen criterios de decisi{\'o}n (por ejemplo, separaci{\'o}n de preocupaciones, modularidad, bajo acoplamiento, alta cohesi{\'o}n, seguridad por dise{\~n}o, automatizaci{\'o}n de pruebas) y permiten justificar elecciones t{\'e}cnicas de manera sistem{\'a}tica. En entornos IoT/SIbIoT, los principios ayudan a controlar complejidad, asegurar interoperabilidad y sostener la evoluci{\'o}n controlada en infraestructuras distribuidas.},
		type=main,
		symbol={📌},
		user1={Regla rectora},
		user2={Criterio para decisiones},
		user3={Justificaci{\'o}n t{\'e}cnica},
		see={filosofia,decisionarquitectonica,modularidad,interoperabilidad,seguridad,evolucioncontrolada}
	}
	
	\newglossaryentry{principioagil}{
		name={Principio {\'a}gil},
		text={principio {\'a}gil},
		sort={principio agil},
		plural={principios {\'a}giles},
		first={principio {\'a}gil},
		firstplural={principios {\'a}giles},
		description={Directriz fundamental que orienta la organizaci{\'o}n del trabajo y las relaciones dentro de una metodolog{\'i}a {\'a}gil. Los \emph{principios {\'a}giles}, derivados del Manifiesto {\'A}gil, promueven entrega temprana y continua de valor, colaboraci{\'o}n con el cliente, adaptabilidad al cambio, desarrollo sostenible, excelencia t{\'e}cnica y comunicaci{\'o}n efectiva. En entornos \gls{iot}, gu{\'i}an la construcci{\'o}n incremental de soluciones distribuidas, flexibles y centradas en el \gls{usuario}.},
		type=main,
		symbol={📜},
		user1={Directrices del Manifiesto {\'A}gil},
		user2={Orientaci{\'o}n para trabajo colaborativo},
		user3={Construcci{\'o}n incremental en IoT},
		see={practicaagil,planificacionagil,equipoagil}
	}
	
	\newglossaryentry{privacidad}{
		name={Privacidad},
		text={privacidad},
		sort={privacidad},
		plural={privacidades},
		first={privacidad},
		firstplural={privacidades},
		description={Derecho y capacidad de las personas y organizaciones para controlar la recolecci{\'o}n, uso, divulgaci{\'o}n y conservaci{\'o}n de sus datos personales o sensibles. En \gls{iot}/\glspl{sibiot}, la \emph{privacidad} exige minimizaci{\'o}n de datos, prop{\'o}sito espec{\'i}fico, transparencia, control de consentimiento y medidas t{\'e}cnicas como seudonimizaci{\'o}n, cifrado y control de acceso. Se relaciona con el cumplimiento normativo y con mecanismos de \gls{seguridad} para prevenir accesos o tratamientos no autorizados.},
		type=main,
		symbol={🕶️},
		user1={Protecci{\'o}n de datos},
		user2={Gobernanza y cumplimiento},
		user3={Controles t{\'e}cnicos y organizativos},
		see={seguridad,interoperabilidad,buenaspracticas}
	}
	
	\newglossaryentry{procesamiento}{
		name={Procesamiento},
		text={procesamiento},
		sort={procesamiento},
		plural={procesamientos},
		first={procesamiento},
		firstplural={procesamientos},
		description={Conjunto de operaciones l{\'o}gicas, aritm{\'e}ticas o de transformaci{\'o}n aplicadas a datos para generar informaci{\'o}n significativa o resultados {\'u}tiles. El procesamiento puede ser manual o autom{\'a}tico, centralizado o distribuido, en tiempo real o por lotes. En el {\'a}mbito de la ingenier{\'i}a de software y los \glspl{sibiot}, el procesamiento abarca la captura, filtrado, almacenamiento, an{\'a}lisis y visualizaci{\'o}n de datos provenientes de sensores, dispositivos o sistemas conectados. Incluye tambi{\'e}n la aplicaci{\'o}n de algoritmos de \gls{aprendizajeautomatico}, t{\'e}cnicas de \gls{cloudcomputing} o \gls{edgecomputing}, y modelos predictivos para optimizar la toma de decisiones autom{\'a}tica.},
		type=main,
		user1={Transformaci{\'o}n de datos en informaci{\'o}n significativa},
		user2={Operaciones l{\'o}gicas, aritm{\'e}ticas y de an{\'a}lisis sobre datos},
		user3={Ejecuci{\'o}n en tiempo real, por lotes o distribuida},
		user4={Base de los sistemas inteligentes y del an{\'a}lisis automatizado},
		see={dato,analisis,cloudcomputing,edgecomputing,aprendizajeautomatico,modelopredictivo},
		symbol={\ensuremath{\Pi}}
	}
	
	\newglossaryentry{procesamientoanalitico}{
		name={Procesamiento anal{\'i}tico},
		text={procesamiento anal{\'i}tico},
		sort={procesamiento analitico},
		plural={procesamientos anal{\'i}ticos},
		first={procesamiento anal{\'i}tico},
		firstplural={procesamientos anal{\'i}ticos},
		description={Conjunto de t{\'e}cnicas, m{\'e}todos y operaciones orientadas al an{\'a}lisis profundo de datos con el fin de descubrir patrones, tendencias, correlaciones y conocimiento relevante para la toma de decisiones. El procesamiento anal{\'i}tico puede involucrar consultas complejas, agregaciones, miner{\'i}a de datos, aprendizaje autom{\'a}tico y modelado estad{\'i}stico. En el contexto de sistemas \gls{iot} y \glspl{sibiot}, permite transformar grandes vol{\'u}menes de datos generados por \glspl{sensor} y \glspl{dispositivo} en informaci{\'o}n estrat{\'e}gica, soportando an{\'a}lisis descriptivo, predictivo y prescriptivo dentro de arquitecturas \gls{olap}, \gls{htap} o plataformas de \gls{bigdata}.},
		type=main,
		symbol={📊},
		user1={An{\'a}lisis de datos avanzado},
		user2={Descubrimiento de patrones y tendencias},
		user3={Soporte a decisiones estrat{\'e}gicas},
		see={olap,bigdata,htap,modelopredictivo,tomadecisiones}
	}
	
	\newglossaryentry{procesamientobatch}{
		name={Procesamiento batch},
		text={procesamiento batch},
		sort={procesamiento batch},
		plural={procesamientos batch},
		first={procesamiento batch},
		firstplural={procesamientos batch},
		description={Modelo de procesamiento de datos en el cual grandes vol{\'u}menes de informaci{\'o}n se recopilan durante un periodo de tiempo y se procesan de forma conjunta en ejecuciones programadas. El procesamiento batch no requiere respuesta inmediata y est{\'a} orientado al an{\'a}lisis hist{\'o}rico, la generaci{\'o}n de informes y el c{\'a}lculo de agregados. En sistemas de informaci{\'o}n basados en IoT y SIbIoT, este enfoque se utiliza principalmente en capas anal{\'i}ticas para evaluar patrones de comportamiento, realizar predicciones de demanda y soportar la toma de decisiones estrat{\'e}gicas},
		type=main,
		symbol={📦},
		see={procesamientocontinuo,procesamientoparalelo,analiticadatos,procesamientodistribuido,sibiot}
	}
	
	\newglossaryentry{procesamientocontinuo}{
		name={Procesamiento continuo},
		text={procesamiento continuo},
		sort={procesamientocontinuo},
		plural={procesamientos continuos},
		first={procesamiento continuo},
		firstplural={procesamientos continuos},
		description={Modelo de procesamiento de datos en el que la informaci{\'o}n es tratada de manera incremental y permanente a medida que se genera, permitiendo reacciones casi en tiempo real ante eventos o cambios del sistema. El procesamiento continuo es caracter{\'i}stico de arquitecturas orientadas a eventos y flujos de datos, y resulta esencial en sistemas IoT y SIbIoT para la monitorizaci{\'o}n en tiempo real, la detecci{\'o}n de anomal{\'i}as y la ejecuci{\'o}n de acciones autom{\'a}ticas con baja latencia},
		type=main,
		symbol={🔄},
		see={procesamientobatch,orientadoeventos,tiemporeal,procesamientodistribuido,sibiot}
	}
	
	\newglossaryentry{procesamientodatos}{
		name={Procesamiento de datos},
		text={procesamiento de datos},
		sort={procesamiento de datos},
		plural={procesamientos de datos},
		first={procesamiento de datos},
		firstplural={procesamientos de datos},
		description={El procesamiento de datos es el conjunto de operaciones que permiten recopilar, organizar, transformar y analizar datos con el fin de convertirlos en informaci{\'o}n {\'u}til y significativa. En el contexto de los \glspl{sibiot}, este procesamiento puede realizarse tanto en el \gls{borde} como en la \gls{nube}, permitiendo optimizar el desempe{\~n}o del \gls{sistemainformatico}, facilitar la \gls{tomadecisiones} en tiempo real y mejorar la eficiencia de los recursos.},
		type=\glsdefaulttype,
		symbol={🖥️},
		user1={Tratamiento de datos},
		user2={Gest{\'i}on de datos},
		user3={An{\'a}lisis de datos},
		see={dato}
	}
	
	\newglossaryentry{procesamientodistribuido}{
		name={Procesamiento distribuido},
		text={Procesamiento distribuido},
		sort={procesamiento distribuido},
		description={%
			Paradigma de computaci{\'o}n en el que las tareas de procesamiento de datos y ejecuci{\'o}n de l{\'o}gica se reparten entre m{\'u}ltiples nodos interconectados, los cuales cooperan para resolver un problema com{\'u}n sin depender de un {\'u}nico punto central de control.
			Este enfoque permite mejorar la escalabilidad, la tolerancia a fallos y el rendimiento del sistema, al distribuir la carga de trabajo y reducir los cuellos de botella asociados a arquitecturas centralizadas.
			En sistemas de informaci{\'o}n basados en IoT, el procesamiento distribuido se materializa mediante la ejecuci{\'o}n coordinada de funciones en dispositivos f{\'i}sicos, nodos de borde, pasarelas y servicios en la nube, posibilitando respuestas en tiempo casi real, reducci{\'o}n de latencia y uso eficiente de los recursos de red y c{\'o}mputo.},
		first={procesamiento distribuido},
		plural={procesamientos distribuidos},
		firstplural={procesamientos distribuidos},
		see={sistemadistribuido,edgecomputing,fogcomputing,cloudcomputing,iot}
	}
	
	\newglossaryentry{procesamientografico}{
		name={procesamiento gr{\'a}fico},
		text={procesamiento gr{\'a}fico},
		sort={procesamiento grafico},
		plural={procesamientos gr{\'a}ficos},
		first={procesamiento gr{\'a}fico},
		firstplural={procesamientos gr{\'a}ficos},
		description={Conjunto de t{\'e}cnicas, algoritmos y operaciones orientadas a generar, transformar, analizar o representar informaci{\'o}n visual mediante recursos computacionales especializados. El \emph{procesamiento gr{\'a}fico} emplea unidades de procesamiento gr{\'a}fico (GPU) y bibliotecas optimizadas para realizar c{\'a}lculos paralelos de alta velocidad, esenciales en renderizado 2D y 3D, visualizaci{\'o}n cient{\'i}fica, aprendizaje profundo, simulaciones y an{\'a}lisis de im{\'a}genes. En el contexto de los \glspl{sibiot}, permite interpretar datos visuales provenientes de c{\'a}maras, sensores de visi{\'o}n y dispositivos de borde, habilitando tareas como reconocimiento de objetos, monitorizaci{\'o}n inteligente y toma de decisiones basada en informaci{\'o}n visual},
		type=main,
		symbol={🎨},
		user1={C{\'o}mputo visual acelerado},
		user2={Uso de GPU, paralelismo y bibliotecas gr{\'a}ficas},
		user3={Aplicaciones en visi{\'o}n computacional y an{\'a}lisis visual en \glspl{sibiot}},
		see={gpu,visioncomputacional,procesamientodatos}
	}
	
	\newglossaryentry{procesamientoparalelo}{
		name={Procesamiento en paralelo},
		text={procesamiento en paralelo},
		sort={procesamiento parallo},
		plural={procesamientos en paralelo},
		first={procesamiento en paralelo},
		firstplural={procesamientos en paralelo},
		description={T{\'e}cnica de procesamiento en la que una tarea se divide en subprocesos que se ejecutan simult{\'a}neamente sobre m{\'u}ltiples unidades de c{\'o}mputo, con el objetivo de reducir el tiempo total de ejecuci{\'o}n y mejorar el rendimiento. En sistemas de informaci{\'o}n basados en IoT y SIbIoT, el procesamiento en paralelo es fundamental para manejar grandes vol{\'u}menes de datos, soportar cargas elevadas y escalar servicios anal{\'i}ticos y operacionales en infraestructuras distribuidas},
		type=main,
		symbol={⚙️},
		see={procesamientobatch,procesamientocontinuo,procesamientodistribuido,escalabilidad,sibiot}
	}
	
	\newglossaryentry{proceso}{
		name={Proceso},
		text={proceso},
		sort={proceso},
		plural={procesos},
		first={proceso},
		firstplural={procesos},
		description={Un \emph{proceso} es un conjunto organizado de actividades o pasos interrelacionados que se ejecutan de manera secuencial o concurrente para lograr un objetivo espec{\'i}fico. En ingenier{\'i}a de software, un proceso tambi{\'e}n se refiere a la instancia en ejecuci{\'o}n de un programa dentro de un sistema operativo. En el contexto de los \glspl{sibiot}, los procesos abarcan desde el \gls{procesamientodatos} en \glspl{dispositivoiot} hasta los flujos de \gls{gestion} que coordinan sensores, plataformas en la \gls{nube} y algoritmos de \gls{aprendizajeautomatico}, garantizando la trazabilidad, eficiencia y \gls{escalabilidad} del sistema.},
		type=main,
		symbol={🔄},
		user1={Conjunto de actividades organizadas},
		user2={Instancia en ejecuci{\'o}n de un programa},
		user3={Flujos en sistemas IoT y software},
		see={gestion,procesamientodatos,metodologia,sibiot}
	}
		
	\newglossaryentry{productbacklog}{
		name={Product Backlog},
		text={\textit{Product Backlog}},
		sort={product backlog},
		plural={\textit{Product Backlogs}},
		first={\textit{Product Backlog}},
		firstplural={\textit{Product Backlogs}},
		description={Lista priorizada y evolutiva de todos los requisitos, mejoras, correcciones y funcionalidades necesarias para el desarrollo de un producto. Es gestionado por el Product Owner y constituye la {\'u}nica fuente oficial de trabajo para el equipo de desarrollo. El Product Backlog se refina continuamente para asegurar claridad, orden y alineaci{\'o}n con los objetivos estrat{\'e}gicos del producto.},
		type=main,
		symbol={📋},
		user1={Lista priorizada de trabajo},
		user2={Fuente oficial de requisitos},
		user3={Gesti{\'o}n evolutiva del producto},
		see={sprintbacklog,sprint,agil,productowner}
	}
	
	\newglossaryentry{productgoal}{
		name={Product Goal},
		text={\textit{Product Goal}},
		sort={product goal},
		plural={\textit{Product Goals}},
		first={\textit{Product Goal}},
		firstplural={Product Goals},
		description={Objetivo estrat{\'e}gico de largo plazo del producto que define el estado futuro deseado y orienta la evoluci{\'o}n del Product Backlog. Proporciona direcci{\'o}n y coherencia al trabajo del equipo Scrum, sirviendo como referencia para priorizaci{\'o}n y toma de decisiones.},
		type=main,
		symbol={🎯},
		user1={Objetivo estrat{\'e}gico del producto},
		user2={Gu{\'i}a del Product Backlog},
		user3={Direcci{\'o}n a largo plazo},
		see={productbacklog,sprintgoal,incremento}
	}
	
	\newglossaryentry{productividad}{
		name={Productividad},
		text={productividad},
		sort={productividad},
		plural={productividades},
		first={productividad},
		firstplural={productividades},
		description={Medida de la eficiencia con la que se emplean los recursos (tiempo, energ{\'i}a, materiales, capital humano) para alcanzar resultados. La \emph{productividad} se expresa generalmente como una raz{\'o}n entre la producci{\'o}n obtenida y los insumos utilizados. En el contexto de los \glspl{sibiot}, la productividad se mejora mediante la automatizaci{\'o}n de procesos, la \gls{monitorizacion} en \gls{tiemporeal}, el an{\'a}lisis de datos y la optimizaci{\'o}n de recursos, lo que permite reducir costos, aumentar la \gls{eficienciaoperativa} y mejorar la \gls{sostenibilidad}.},
		type=main,
		symbol={📈},
		user1={Eficiencia en el uso de recursos},
		user2={Rendimiento de procesos},
		user3={Medida de resultados alcanzados},
		see={eficiencia,eficienciaoperativa,sostenibilidad}
	}
	
	\newglossaryentry{productowner}{
		name={Product Owner},
		text={Product Owner},
		sort={product owner},
		plural={Product Owners},
		first={Product Owner},
		firstplural={Product Owners},
		description={Rol dentro del marco de trabajo Scrum responsable de maximizar el valor del producto resultante del trabajo del equipo de desarrollo. El Product Owner gestiona y prioriza el Product Backlog, define criterios de aceptaci{\'o}n, clarifica requisitos y asegura que el trabajo desarrollado est{\'e} alineado con los objetivos estrat{\'e}gicos del producto y las necesidades de los interesados. Act{\'u}a como puente entre los stakeholders y el equipo, tomando decisiones sobre alcance, prioridades y evoluci{\'o}n del producto.},
		type=main,
		symbol={👤},
		user1={Responsable de maximizar valor},
		user2={Gesti{\'o}n y priorizaci{\'o}n del Product Backlog},
		user3={Representante de stakeholders},
		see={productbacklog,sprint,sprintbacklog,scrum,agil}
	}
	
	\newglossaryentry{programa}{
		name={Programa},
		text={programa},
		sort={programa},
		plural={programas},
		first={programa},
		firstplural={programas},
		description={Conjunto organizado de instrucciones, actividades o pol{\'i}ticas dise{\~n}adas para cumplir un objetivo espec{\'i}fico. En el contexto de la ingenier{\'i}a de software y los \glspl{sibiot}, un \emph{programa} puede referirse tanto a un componente de software ejecutable como a un plan estructurado de gesti{\'o}n o intervenci{\'o}n, como en los \glspl{programa} de respuesta a la demanda o gesti{\'o}n energ{\'e}tica.},
		type=main,
		symbol={💻},
		user1={Conjunto de instrucciones o acciones planificadas},
		user2={Implementaci{\'o}n software o plan de gesti{\'o}n},
		user3={Elemento estructural de un sistema o proceso},
		see={algoritmo,aplicacion,sibiot}
	}
	
	\newglossaryentry{programacion}{
		name={programaci{\'o}n},
		text={programaci{\'o}n},
		sort={programacion},
		plural={programaciones},
		first={programaci{\'o}n},
		firstplural={programaciones},
		description={Actividad fundamental de la ingenier{\'i}a de software que consiste en dise{\~n}ar, escribir, verificar y mantener c{\'o}digo fuente utilizando lenguajes formales para resolver problemas, ejecutar algoritmos y controlar sistemas computacionales. La \emph{programaci{\'o}n} constituye la base de la implementaci{\'o}n de los \glspl{sibiot}, integrando l{\'o}gica de negocio, comunicaci{\'o}n entre dispositivos, procesamiento de datos y operaciones distribuidas. Requiere rigor t{\'e}cnico, abstracci{\'o}n, dominio de estructuras algor{\'i}tmicas y comprensi{\'o}n de arquitecturas de software},
		type=main,
		symbol={💻},
		user1={Construcci{\'o}n y mantenimiento de c{\'o}digo},
		user2={Dise{\~n}o algor{\'i}tmico y resoluci{\'o}n de problemas},
		user3={Implementaci{\'o}n de funcionalidades en sistemas conectados},
		see={algoritmo,desarrollo,sibiot}
	}
	
	\newglossaryentry{progreso}{
		name={Progreso},
		text={progreso},
		sort={progreso},
		plural={progresos},
		first={progreso},
		firstplural={progresos},
		description={Evidencia objetiva y trazable del avance de un proyecto o proceso de desarrollo respecto de un plan, una l{\'\i}nea base o un conjunto de objetivos. El progreso se expresa mediante entregables verificables, incremento de funcionalidad, cierre de tareas, resultados de pruebas, m{\'e}tricas de calidad y cumplimiento de criterios de aceptaci{\'o}n. En contextos iterativos e incrementales, el progreso se eval{\'u}a por la disponibilidad de \glspl{unidadfuncional} operativas y por la reducci{\'o}n de riesgo t{\'e}cnico (por ejemplo, mediante \glspl{pruebaautomatizada} y verificaci{\'o}n temprana). En \gls{iot}/\glspl{sibiot}, adem{\'a}s, incluye evidencia de integraci{\'o}n hardware--software, conectividad estable, fiabilidad de datos y operaci{\'o}n en condiciones reales.},
		type=main,
		symbol={📈},
		user1={Avance medible},
		user2={Entregables verificables},
		user3={Reducci{\'o}n de riesgo},
		see={entregable,iterativa,incremental,ciclodesarrollo,verificacion,pruebaautomatizada}
	}
	
	\newglossaryentry{protecciondatospersonales}{
		name={Protecci{\'o}n de datos personales},
		text={protecci{\'o}n de datos personales},
		sort={proteccion datos personales},
		plural={protecciones de datos personales},
		first={protecci{\'o}n de datos personales},
		firstplural={protecciones de datos personales},
		description={Conjunto de principios, normas, mecanismos t{\'e}cnicos y medidas organizativas orientadas a garantizar el tratamiento l{\'i}cito, seguro y transparente de los datos personales de las personas f{\'i}sicas. La protecci{\'o}n de datos personales abarca aspectos como la \gls{confidencialidad}, \gls{integridad}, \gls{disponibilidad}, minimizaci{\'o}n, limitaci{\'o}n de finalidad, exactitud y responsabilidad proactiva en el tratamiento de la informaci{\'o}n. En sistemas \gls{iot} y \glspl{sibiot}, implica la implementaci{\'o}n de controles de acceso, cifrado, autenticaci{\'o}n robusta, anonimizaci{\'o}n, gesti{\'o}n de consentimientos y cumplimiento de marcos regulatorios como \gls{gdpr}, con el fin de preservar los derechos fundamentales de los titulares de los datos.},
		type=main,
		symbol={🔐},
		user1={Derecho fundamental},
		user2={Cumplimiento normativo},
		user3={Seguridad y privacidad en IoT},
		see={seguridad,privacidad,gdpr,integridad,confidencialidad}
	}
	
	\newglossaryentry{protocolo}{
		name={protocolo},
		text={protocolo},
		sort={protocolo},
		plural={protocolos},
		first={protocolo},
		firstplural={protocolos},
		description={Conjunto estructurado de reglas, normas y convenciones que regulan el intercambio de informaci{\'o}n entre entidades dentro de un sistema de \gls{comunicacion}. Un protocolo define el formato de los mensajes, los procedimientos de transmisi{\'o}n y recepci{\'o}n, los mecanismos de sincronizaci{\'o}n, el control de errores y las pol{\'i}ticas de gesti{\'o}n del flujo para garantizar una comunicaci{\'o}n segura, ordenada y fiable en entornos \gls{iot}.},
		symbol={PRT},
		see={estandar,comunicacion},
		type=main
	}
	
	\newglossaryentry{protocolocomunicacion}{
		name={Protocolo de comunicaci{\'o}n},
		text={protocolo de comunicaci{\'o}n},
		sort={protocolo comunicacion},
		plural={protocolos de comunicaci{\'o}n},
		first={protocolo de comunicaci{\'o}n},
		firstplural={protocolos de comunicaci{\'o}n},
		description={Un \emph{protocolo de comunicaci{\'o}n} es un conjunto de reglas, est{\'a}ndares y procedimientos que regulan c{\'o}mo se transmiten, reciben y procesan los datos entre dispositivos o sistemas dentro de una \gls{red}. Define aspectos como la estructura de los mensajes, la sincronizaci{\'o}n, la detecci{\'o}n de errores y los mecanismos de control de flujo. En sistemas \gls{iot}, los protocolos de comunicaci{\'o}n permiten la interoperabilidad entre \glspl{sensor}, \glspl{actuador}, \glspl{dispositivoiot} y plataformas en la \gls{nube}, asegurando \gls{fiabilidad}, baja \gls{latencia}, \gls{seguridad} y eficiencia en la \gls{transmision}. Ejemplos incluyen MQTT, CoAP, HTTP/HTTPS, AMQP, Zigbee y LoRaWAN.},
		type=main,
		symbol={📡},
		user1={Define reglas de intercambio de informaci{\'o}n},
		user2={Garantiza interoperabilidad y fiabilidad},
		user3={Fundamental en IoT y sistemas distribuidos},
		see={protocolo,transmision,comunicacion,iot,seguridad,latencia}
	}
	
	% ===================== Protocolo de enrutamiento =====================
	\newglossaryentry{protocoloenrutamiento}{
		name={Protocolo de enrutamiento},
		text={protocolo de enrutamiento},
		sort={protocolo enrutamiento},
		plural={protocolos de enrutamiento},
		first={protocolo de enrutamiento},
		firstplural={protocolos de enrutamiento},
		description={Un \emph{protocolo de enrutamiento} es el conjunto de reglas, m{\'e}todos y algoritmos que permiten a los \glspl{enrutador} y nodos de una \gls{red} intercambiar informaci{\'o}n sobre topolog{\'i}a, seleccionar las mejores rutas y adaptarse a cambios din{\'a}micos en la conectividad. Existen protocolos de enrutamiento est{\'a}ticos y din{\'a}micos, estos {\'u}ltimos pueden clasificarse en \textbf{interior} (RIP, OSPF, EIGRP) y \textbf{exterior} (BGP). En redes \gls{iot}, se emplean protocolos especializados como RPL o AODV, dise{\~n}ados para operar en entornos con recursos limitados, alta movilidad y gran cantidad de \glspl{dispositivoiot}.},
		type=main,
		symbol={🛰️},
		user1={Definen c{\'o}mo los nodos comparten rutas},
		user2={Incluyen est{\'a}ticos, din{\'a}micos, interiores y exteriores},
		user3={Soporte a IoT con protocolos como RPL y AODV},
		see={enrutamiento,protocolocomunicacion,red,sibiot}
	}
	
	% ===================== Protocolo de gesti{\'o}n de energ{\'i}a =====================
	\newglossaryentry{protocologestionenergia}{
		name={Protocolo de gesti{\'o}n de energ{\'i}a},
		text={protocolo de gesti{\'o}n de energ{\'i}a},
		sort={protocolo gestion energia},
		plural={protocolos de gesti{\'o}n de energ{\'i}a},
		first={protocolo de gesti{\'o}n de energ{\'i}a},
		firstplural={protocolos de gesti{\'o}n de energ{\'i}a},
		description={Un \emph{protocolo de gesti{\'o}n de energ{\'i}a} es el conjunto de mecanismos, algoritmos y reglas de comunicaci{\'o}n que permiten optimizar el consumo energ{\'e}tico de \glspl{dispositivo} dentro de una \gls{red}, sin comprometer la funcionalidad del sistema. En \glspl{sibiot}, estos protocolos son fundamentales para prolongar la vida {\'u}til de los nodos alimentados por bater{\'i}a, reducir el tr{\'a}fico innecesario, activar modos de bajo consumo y coordinar el uso eficiente del medio. Ejemplos incluyen IEEE 802.15.4e (TSCH), BLE Low Energy y mecanismos de duty cycling.},
		type=\glsdefaulttype,
		symbol={🔋},
		user1={Optimizaci{\'o}n del consumo energ{\'e}tico},
		user2={Extensi{\'o}n de vida {\'u}til de nodos IoT},
		user3={Ejemplos: IEEE 802.15.4e, BLE, duty cycling},
		see={sibiot,red,dispositivo,energia}
	}
	
	% ===================== Protocolo de seguridad =====================
	\newglossaryentry{protocoloseguridad}{
		name={Protocolo de seguridad},
		text={protocolo de seguridad},
		sort={protocolo seguridad},
		plural={protocolos de seguridad},
		first={protocolo de seguridad},
		firstplural={protocolos de seguridad},
		description={Un \emph{protocolo de seguridad} es el conjunto de reglas, algoritmos y procedimientos dise{\~n}ados para proteger la \gls{integridad}, la confidencialidad, la autenticidad y la disponibilidad de la informaci{\'o}n durante su transmisi{\'o}n a trav{\'e}s de \glspl{red}. En sistemas \gls{iot}, estos protocolos son esenciales debido a la naturaleza distribuida y expuesta de los \glspl{dispositivo}, e incluyen mecanismos como \gls{cifrado}, autenticaci{\'o}n mutua, gesti{\'o}n de llaves y \gls{control} de acceso. Ejemplos comunes son TLS, DTLS, IPsec y protocolos ligeros como OSCORE o EDHOC.},
		type=\glsdefaulttype,
		symbol={🔐},
		user1={Protecci{\'o}n de datos en tr{\'a}nsito},
		user2={Cifrado y autenticaci{\'o}n},
		user3={Protocolos: TLS, DTLS, IPsec, OSCORE},
		see={seguridad,integridad,cifrado,control,disponibilidad}
	}
	
	\newglossaryentry{prototipado}{
		name={Prototipado},
		text={prototipado},
		sort={prototipado},
		first={prototipado},
		description={Proceso iterativo de dise{\~n}o y construcci{\'o}n de versiones preliminares de un sistema, dispositivo o aplicaci{\'o}n con el objetivo de validar conceptos, funcionalidades y decisiones t{\'e}cnicas antes de su implementaci{\'o}n definitiva. En el contexto de sistemas embebidos y \glspl{sibiot}, el \emph{prototipado} permite evaluar arquitecturas, integrar sensores y microcontroladores, y detectar errores tempranos de dise{\~n}o con bajo costo y riesgo.},
		type=main,
		symbol={🛠},
		user1={Dise{\~n}o iterativo},
		user2={Validaci{\'o}n temprana},
		user3={Desarrollo experimental},
		see={prototipo,arduinonano,arduinopromini}
	}
	
	\newglossaryentry{prototipo}{
		name={Prototipo},
		text={prototipo},
		sort={prototipo},
		plural={prototipos},
		first={prototipo},
		firstplural={prototipos},
		description={Versi{\'o}n preliminar de un dispositivo, sistema o producto creada para evaluar su dise{\~n}o, funcionalidad y viabilidad antes de su producci{\'o}n definitiva. Un \emph{prototipo} permite identificar errores, validar requisitos y realizar pruebas experimentales con menor costo y riesgo. En sistemas \gls{iot} y \glspl{sibiot}, los prototipos se construyen frecuentemente con \glspl{placadesarrollo} como Arduino o Raspberry Pi, y sirven para demostrar conceptos, validar arquitecturas y acelerar el ciclo de innovaci{\'o}n.},
		type=main,
		symbol={🧪},
		user1={Modelo experimental},
		user2={Versi{\'o}n preliminar},
		user3={Dise{\~n}o de prueba},
		see={placadesarrollo,planificacionagil,iteracion}
	}
	
	\newglossaryentry{provisionamiento}{
		name={Provisionamiento},
		text={provisionamiento},
		sort={provisionamiento},
		plural={procesos de provisionamiento},
		first={provisionamiento},
		firstplural={procesos de provisionamiento},
		description={Alta y configuraci{\'o}n inicial de \glspl{dispositivo}/servicios, incluyendo credenciales, pol{\'i}ticas y asociaciones a redes y aplicaciones.},
		type=main,
		symbol={🧾},
		see={matter,gestionconfianza,gateway},
		user1={Alta de dispositivos},
		user2={Credenciales},
		user3={Pol{\'i}ticas}
	}
	
	\newglossaryentry{proyecto}{
		name={Proyecto},
		text={proyecto},
		sort={proyecto},
		plural={proyectos},
		first={proyecto},
		firstplural={proyectos},
		description={Iniciativa planificada y temporal orientada a desarrollar un sistema, producto o soluci{\'o}n espec{\'i}fica bajo restricciones de tiempo, costo y alcance definidos.},
		type=main,
		symbol={📊},
		user1={Iniciativa temporal},
		user2={Objetivo definido},
		user3={Restricciones controladas},
		see={ciclodesarrollo,entregable}
	}
	
	\newglossaryentry{prueba}{
		name={Prueba},
		text={prueba},
		sort={prueba},
		plural={pruebas},
		first={prueba},
		firstplural={pruebas},
		description={Actividad sistem{\'a}tica orientada a evaluar el comportamiento de un sistema, componente o artefacto con el fin de detectar defectos, validar requisitos y verificar conformidad con la especificaci{\'o}n t{\'e}cnica. Las pruebas pueden ser unitarias, de integraci{\'o}n, de sistema o de aceptaci{\'o}n, y constituyen un mecanismo fundamental para asegurar calidad, confiabilidad y trazabilidad en el desarrollo de sistemas.},
		type=main,
		symbol={🧪},
		user1={Validaci{\'o}n funcional},
		user2={Control de calidad},
		user3={Detecci{\'o}n de defectos},
		see={verificacion,consistencia,ciclodesarrollo}
	}
	
	\newglossaryentry{pruebaautomatizada}{
		name={Prueba automatizada},
		text={prueba automatizada},
		sort={prueba automatizada},
		plural={pruebas automatizadas},
		first={prueba automatizada},
		firstplural={pruebas automatizadas},
		description={Procedimiento de validaci{\'o}n ejecutado autom{\'a}ticamente mediante herramientas software para verificar que un sistema cumple su especificaci{\'o}n. Forma parte esencial de enfoques de integraci{\'o}n continua y desarrollo dirigido por pruebas.},
		type=main,
		symbol={🧪},
		user1={Validaci{\'o}n autom{\'a}tica},
		user2={Integraci{\'o}n continua},
		user3={Control de calidad},
		see={verificacion}
	}
	
	\newglossaryentry{pseudocodigo}{
		name={Pseudoc{\'o}digo},
		text={pseudoc{\'o}digo},
		sort={pseudocodigo},
		plural={pseudoc{\'o}digos},
		first={pseudoc{\'o}digo},
		firstplural={pseudoc{\'o}digos},
		description={Representaci{\'o}n informal y simplificada de un \gls{algoritmo}, escrita con una mezcla de lenguaje natural estructurado y convenciones propias de la programaci{\'o}n, pero sin apegarse a la sintaxis estricta de un lenguaje formal. El \emph{pseudoc{\'o}digo} permite describir la l{\'o}gica de un programa de manera clara y comprensible para humanos, sirviendo como puente entre la especificaci{\'o}n de requisitos y la implementaci{\'o}n en un lenguaje de programaci{\'o}n. En el contexto educativo y de ingenier{\'i}a de software, facilita la comunicaci{\'o}n entre equipos de desarrollo y el dise{\~n}o previo de \glspl{sistemainteligente} o \gls{sibiot}.},
		type=main,
		symbol={✍️},
		user1={Dise{\~n}o de algoritmos},
		user2={Lenguaje intermedio entre natural y formal},
		user3={Herramienta pedag{\'o}gica},
		see={algoritmo,diagramaflujo,lenguajeprogramacion}
	}
	
	\newglossaryentry{ptp}{
		name={PTP},
		text={PTP},
		sort={ptp},
		plural={PTP},
		first={protocolo de precisi{\'o}n de tiempo (\textit{Precision Time Protocol}, PTP)},
		firstplural={PTP},
		description={\emph{PTP} es un protocolo de sincronizaci{\'o}n temporal definido en el est{\'a}ndar IEEE~1588, dise{\~n}ado para lograr una precisi{\'o}n del orden de microsegundos o incluso nanosegundos en sistemas distribuidos. Utiliza la detecci{\'o}n y correcci{\'o}n activa de retardos dentro de la red, y se apoya en la figura del \emph{Grandmaster Clock}, que entrega la referencia temporal a los dem{\'a}s nodos. En entornos \gls{iot} industriales, vehiculares, rob{\'o}ticos o de \gls{automatizacion}, PTP resulta crucial para coordinar actuadores, \glspl{sensor} de alta frecuencia, procesos sincronizados, control en tiempo real y sistemas de borde donde la sincronizaci{\'o}n estricta es indispensable.},
		type=main,
		symbol={⏲️},
		user1={Sincronizaci{\'o}n con precisiones de micro/nanosegundos},
		user2={Adecuado para automatizaci{\'o}n y control en tiempo real},
		user3={Basado en IEEE~1588 con reloj maestro},
		see={ntp,infraestructuradistribuida,automatizacion,sibiot}
	}
	
	\newglossaryentry{pubsub}{
		name={Pub/Sub},
		text={Pub/Sub},
		sort={pubsub},
		first={Pub/Sub},
		description={Modelo y servicio de mensajer{\'i}a distribuida basado en el patr{\'o}n publicador--suscriptor, que permite la comunicaci{\'o}n as{\'i}ncrona y desacoplada entre productores y consumidores de eventos. En Google Cloud, \emph{Pub/Sub} proporciona entrega escalable, tolerante a fallos y de baja latencia de mensajes, siendo ampliamente utilizado como infraestructura de integraci{\'o}n orientada a eventos. En sistemas \glspl{sibiot}, Pub/Sub act{\'u}a como columna vertebral de la mensajer{\'i}a en tiempo real, facilitando la propagaci{\'o}n de eventos provenientes de dispositivos IoT, gateways y nodos de borde hacia capas de procesamiento, almacenamiento y aplicaciones},
		type=main,
		symbol={📨},
		see={mensajeriatiemporeal,orientadoeventos,desacoplamiento,escalabilidad,procesamientodistribuido,sibiot}
	}
	
	% ===================== Puerto (l{\'o}gico) =====================
	\newglossaryentry{puerto}{
		name={Puerto},
		text={puerto},
		sort={puerto},
		plural={puertos},
		first={puerto},
		firstplural={puertos},
		description={Un \emph{puerto} es un identificador num{\'e}rico l{\'o}gico utilizado en redes de computadoras para distinguir diferentes servicios o aplicaciones que se ejecutan en un mismo \gls{dispositivo}. En el modelo TCP/IP, los puertos permiten direccionar correctamente el tr{\'a}fico entre \glspl{cliente} y \glspl{servidor} mediante protocolos como TCP o UDP.},
		type=\glsdefaulttype,
		symbol={🔢},
		user1={Identificador num{\'e}rico l{\'o}gico},
		user2={Distinci{\'o}n de servicios en un dispositivo},
		user3={Uso en protocolos TCP/UDP},
		see={cliente,servidor,protocolocomunicacion}
	}
	
	% ===================== Puerto f{\'i}sico =====================
	\newglossaryentry{puertofisico}{
		name={Puerto (f{\'i}sico)},
		text={puerto (f{\'i}sico)},
		sort={puerto fisico},
		plural={puertos (f{\'i}sicos)},
		first={puerto f{\'i}sico},
		firstplural={puertos f{\'i}sicos},
		description={Un \emph{puerto f{\'i}sico} es un conector o interfaz tangible en un \gls{dispositivo} electr{\'o}nico que permite el intercambio de datos, energ{\'i}a o se{\~n}ales con otros dispositivos. Ejemplos comunes incluyen puertos HDMI, USB, VGA y Ethernet, que facilitan la comunicaci{\'o}n directa entre componentes de \gls{hardware}.},
		type=\glsdefaulttype,
		symbol={🧲},
		user1={Interfaz de hardware},
		user2={Conector para datos, energ{\'i}a o se{\~n}ales},
		user3={Ejemplos: HDMI, USB, Ethernet},
		see={hardware,ethernet,hdmi}
	}
	
	\newglossaryentry{put}{
		name={PUT},
		text={PUT},
		sort={put},
		first={PUT},
		description={M{\'e}todo del protocolo HTTP utilizado para crear o actualizar completamente un recurso identificado por una URL espec{\'i}fica. El m{\'e}todo \texttt{PUT} es idempotente, lo que significa que m{\'u}ltiples solicitudes id{\'e}nticas producen el mismo resultado. En sistemas RESTful e IoT, \texttt{PUT} se emplea para actualizar configuraciones, estados de dispositivos o representaciones completas de recursos gestionados por el sistema},
		type=main,
		see={http,rest,restful,api,configuracion,control}
	}
	
	\newglossaryentry{pv}{
		name={Energ{\'i}a fotovoltaica},
		text={energ{\'i}a fotovoltaica},
		sort={energia fotovoltaica},
		plural={tecnolog{\'i}as fotovoltaicas},
		first={energ{\'i}a fotovoltaica (\textit{PhotoVoltaic Energy},, PV)}
		firstplural={tecnolog{\'i}as fotovoltaicas},
		description={Conversi{\'o}n directa de radiaci{\'o}n solar en electricidad; base del autoconsumo residencial y de la integraci{\'o}n con \gls{hems}.},
		type=main,
		symbol={🔆},
		see={hems,almacenamientoenergetico,gestionenergetica},
		user1={Generaci{\'o}n distribuida},
		user2={Autoconsumo},
		user3={Sostenibilidad}
	}
	
	\newglossaryentry{python}{
		name={Python},
		text={Python},
		plural={lenguajes Python},
		first={Python, un lenguaje de programaci{\'o}n de alto nivel caracterizado por su sintaxis clara y su enfoque en la productividad},
		firstplural={lenguajes Python ampliamente utilizados en ciencia de datos, automatizaci{\'o}n y sistemas IoT},
		description={Lenguaje de programaci{\'o}n interpretado y multiparadigma empleado en anal{\'i}tica de datos, aprendizaje autom{\'a}tico, automatizaci{\'o}n y desarrollo de sistemas basados en IoT},
		sort={python},
		type=main,
		symbol={🐍},
		see={lenguajeprogramacion}
	}
	
	% ===================== RAM =====================
	\newglossaryentry{ram}{
		name={RAM},
		text={RAM},
		sort={ram},
		plural={RAMs},
		first={memoria de acceso aleatorio (\textit{Random Access Memory}, RAM)},
		firstplural={RAMs},
		description={\emph{RAM} (memoria de acceso aleatorio) es una memoria vol{\'a}til utilizada para almacenar temporalmente \glspl{dato} e instrucciones mientras un \gls{dispositivo} est{\'a} en funcionamiento. Permite acceso r{\'a}pido y aleatorio a cualquier direcci{\'o}n de memoria, siendo esencial para ejecutar programas y procesos en tiempo real. A diferencia de la memoria \gls{flash} o el \gls{almacenamiento}, la RAM se borra al apagar el sistema.},
		type=\glsdefaulttype,
		symbol={📖},
		user1={Memoria vol{\'a}til},
		user2={Acceso r{\'a}pido y aleatorio},
		user3={Ejecuci{\'o}n de programas y procesos},
		see={flash,almacenamiento,firmware}
	}
	
	\newglossaryentry{raspberrypi}{
		name={Raspberry~Pi},
		text={Raspberry~Pi},
		sort={raspberrypi},
		plural={Raspberry~Pi},
		description={Ordenador de placa reducida desarrollado por la \textit{Raspberry Pi Foundation}, ampliamente utilizado para ense{\~n}ar programaci{\'o}n, ciencias de la computaci{\'o}n y electr{\'o}nica. Se ha convertido en una plataforma de referencia en proyectos de \gls{iot}, rob{\'o}tica, automatizaci{\'o}n y sistemas embebidos, gracias a su bajo costo, tama{\~n}o compacto, bajo consumo energ{\'e}tico y gran comunidad de soporte. Incorpora un \gls{microprocesador} ARM, puertos GPIO, conectividad de red y soporte para diversos sistemas operativos como Linux, permitiendo el desarrollo de aplicaciones tanto educativas como profesionales.},
		type=\glsdefaulttype,
		symbol={🍓},
		user1={Computador de placa reducida},
		user2={Plataforma educativa y de desarrollo},
		user3={Hardware para proyectos IoT y embebidos},
		see={placadesarrollo,placacontroladora,placaprocesamiento,arduino,esp32}
	}
	
	\newglossaryentry{razonamiento}{
		name={Razonamiento},
		text={razonamiento},
		sort={razonamiento},
		plural={razonamientos},
		first={razonamiento},
		firstplural={razonamientos},
		description={Proceso cognitivo, l{\'o}gico o computacional mediante el cual se infieren conclusiones, se resuelven problemas o se toman decisiones a partir de informaci{\'o}n disponible, reglas o modelos. El \emph{razonamiento} puede ser deductivo (de lo general a lo particular), inductivo (de lo particular a lo general) o abductivo (generaci{\'o}n de la mejor explicaci{\'o}n posible). En el contexto de los \glspl{sistemainteligente} y los \glspl{sibiot}, el razonamiento se implementa mediante algoritmos de \gls{ia}, motores de reglas, l{\'o}gica difusa o modelos probabil{\'i}sticos, permitiendo interpretar datos de \glspl{sensor}, detectar patrones, anticipar situaciones y respaldar la toma de decisiones autom{\'a}tica o asistida.},
		type=main,
		symbol={🧩},
		user1={Inferencia y toma de decisiones},
		user2={Procesamiento l{\'o}gico o computacional},
		user3={Base de la inteligencia artificial},
		see={ia,sistemainteligente,aprendizajeautomatico,tomadecisiones}
	}
	
	\newglossaryentry{rds}{
		name={RDS},
		text={RDS},
		sort={rds},
		plural={RDS},
		first={servicio de base de datos relacional (\textit{Relational Database Service}, RDS)},
		firstplural={servicios de base de datos relacionales o \textit{Relational Database Services} (RDSs)},
		description={Servicio gestionado de bases de datos relacionales que abstrae la administraci{\'o}n de infraestructura, permitiendo el despliegue, escalado y mantenimiento de motores relacionales est{\'a}ndar como \gls{mysql}, \gls{postgresql}, MariaDB y otros. En arquitecturas \glspl{sibiot}, \emph{RDS} se emplea como componente de persistencia operacional (\gls{oltp}), soportando transacciones consistentes, integridad referencial y consultas estructuradas sobre datos operativos en tiempo casi real. Su uso resulta apropiado para el almacenamiento de eventos transaccionales, configuraciones del sistema, estados operativos y registros de control, integr{\'a}ndose habitualmente con capas anal{\'i}ticas y repositorios especializados dentro de una \gls{arquitecturapoliclinica}.},
		see={arquitecturapoliclinica,basedatosrelacional,cloudcomputing,consistencia,oltp,sibiot},
		type=main,
		symbol={{🗄️}}
	}
	
	
	\newglossaryentry{reactivo}{
		name={Reactivo},
		text={reactivo},
		sort={reactivo},
		description={%
			Caracter{\'i}stica de un sistema o componente cuya l{\'o}gica de comportamiento se activa como respuesta directa a eventos, cambios de estado o est{\'i}mulos provenientes de su entorno, en lugar de operar mediante ejecuciones peri{\'o}dicas o secuenciales predefinidas.
			Los sistemas reactivos priorizan la capacidad de respuesta oportuna, la adaptabilidad y la gesti{\'o}n eficiente de eventos concurrentes, siendo especialmente adecuados para entornos din{\'a}micos y distribuidos.
		},
		first={reactivo},
		plural={reactivos},
		firstplural={reactivos},
		see={enfoquereactivo,orientadoeventos,sistemadistribuido}
	}
	
	\newglossaryentry{realidad}{
		name={Realidad},
		text={realidad},
		sort={realidad},
		first={realidad},
		description={Condiciones y fen{\'o}menos del entorno f{\'i}sico y social observables, que constituyen el contexto de operaci{\'o}n de un sistema. En SIbIoT, la realidad es capturada mediante sensado y modelada como datos.},
		type=main,
		symbol={🌎},
		user1={Entorno},
		user2={Fen{\'o}menos},
		user3={Contexto},
		see={dato,monitorizacion}
	}
	
	\newglossaryentry{realidadaumentada}{
		name={Realidad aumentada},
		text={realidad aumentada},
		sort={realidad aumentada},
		firstplural={realidades aumentadas},
		plural={ARs},
		first={realidad aumentada (\textit{Augmented Reality}, AR)},
		description={La \emph{realidad aumentada} es una tecnolog{\'i}a que superpone informaci{\'o}n digital (im{\'a}genes, datos, modelos 3D) sobre el entorno f{\'i}sico percibido por el usuario, enriqueciendo su experiencia sensorial y cognitiva. En el contexto de \gls{iot}, la AR permite visualizar datos en tiempo real de \glspl{sensor}, controlar \glspl{dispositivointeligente} y guiar procesos industriales o educativos integrando el \gls{mundofisico} con el \gls{mundodigital}.},
		type=main,
		symbol={🕶️},
		user1={Superposici{\'o}n de datos digitales en el mundo real},
		user2={Visualizaci{\'o}n contextualizada de IoT},
		user3={Aplicaciones en industria, educaci{\'o}n y salud},
		see={realidadvirtual,gemelodigital,mundodigital,mundofisico}
	}
	
	\newglossaryentry{realidadvirtual}{
		name={Realidad virtual},
		text={realidad virtual},
		sort={realidad virtual},
		firstplural={realidades virtuales},
		plural={VRs},
		first={realidad virtual (\textit{Virtual Reality}, VR)},
		firstplural={realidades virtuales},
		description={La \emph{realidad virtual} es una tecnolog{\'i}a inmersiva que crea un entorno digital simulado, completamente independiente del mundo f{\'i}sico, al cual los usuarios acceden mediante dispositivos especializados (gafas, guantes h{\'a}pticos, interfaces de movimiento). En sistemas \gls{iot}, la VR se utiliza para simulaci{\'o}n de escenarios complejos, capacitaci{\'o}n en entornos seguros, experimentaci{\'o}n con \glspl{gemelodigital} y optimizaci{\'o}n de procesos industriales o urbanos antes de su implementaci{\'o}n en el mundo real.},
		type=main,
		symbol={🎮},
		user1={Entornos inmersivos simulados},
		user2={Capacitaci{\'o}n y simulaci{\'o}n seguras},
		user3={Soporte a gemelos digitales y an{\'a}lisis IoT},
		see={realidadaumentada,gemelodigital,mundovirtual,mundodigital}
	}
	
	\newglossaryentry{reconocimiento}{
		name={Reconocimiento},
		text={reconocimiento},
		sort={reconocimiento},
		plural={reconocimientos},
		first={reconocimiento},
		firstplural={reconocimientos},
		description={Proceso de clasificar o asociar un objeto, patr{\'o}n o se{\~n}al detectada con una categor{\'i}a conocida. A diferencia de la \gls{deteccion}, el \emph{reconocimiento} implica identificar caracter{\'i}sticas relevantes y compararlas con modelos, bases de datos o descripciones previas, sin necesariamente individualizar al elemento observado. En \glspl{sibiot}, se aplica en tareas como reconocimiento de voz, reconocimiento de im{\'a}genes o patrones de actividad, utilizando \glspl{redneuronal} y algoritmos de \gls{aprendizajeautomatico}.},
		type=main,
		symbol={🔎},
		user1={Clasificaci{\'o}n de patrones},
		user2={Vinculaci{\'o}n a una categor{\'i}a conocida},
		user3={Segunda etapa tras la detecci{\'o}n},
		see={deteccion,identificacion,aprendizajeautomatico,redneuronal}
	}
	
	\newglossaryentry{reconocimientofacial}{
		name={Reconocimiento facial},
		text={reconocimiento facial},
		sort={reconocimiento facial},
		plural={reconocimientos faciales},
		first={reconocimiento facial},
		firstplural={reconocimientos faciales},
		description={Tecnolog{\'i}a de \gls{visionartificial} y biometr{\'i}a que permite identificar o verificar la identidad de una persona a partir de las caracter{\'i}sticas y patrones de su rostro. El \emph{reconocimiento facial} combina procesos de \gls{deteccion} de rostros, \gls{reconocimiento} de rasgos faciales y \gls{identificacion} en bases de datos para establecer coincidencias. En el contexto de los \glspl{sibiot}, se utiliza en sistemas de seguridad, control de accesos, monitorizaci{\'o}n de asistencia, personalizaci{\'o}n de servicios y an{\'a}lisis de comportamiento, apoy{\'a}ndose en \glspl{redneuronal} convolucionales y t{\'e}cnicas de \gls{aprendizajeautomatico} para operar en \gls{tiemporeal} con alta \gls{precision}.},
		type=main,
		symbol={👤},
		user1={Biometr{\'i}a basada en rostros},
		user2={Identificaci{\'o}n y verificaci{\'o}n de identidad},
		user3={Aplicaciones en seguridad y personalizaci{\'o}n},
		see={visionartificial,deteccion,reconocimiento,identificacion,aprendizajeautomatico}
	}
	
	\newglossaryentry{recuperacion}{
		name={Recuperaci{\'o}n},
		text={recuperaci{\'o}n},
		sort={recuperacion},
		first={recuperaci{\'o}n},
		description={Conjunto de t{\'e}cnicas para restaurar o reestablecer el funcionamiento y los datos tras fallas, errores o incidentes, incluyendo copias de seguridad, restauraci{\'o}n y mecanismos de tolerancia a fallos.},
		type=main,
		symbol={🛟},
		user1={Restauraci{\'o}n},
		user2={Continuidad},
		user3={Resiliencia},
		see={falla,error,basedatos}
	}
	
	\newglossaryentry{recurso}{
		name={Recurso},
		text={recurso},
		sort={recurso},
		plural={recursos},
		first={recurso},
		firstplural={recursos},
		description={Elemento f{\'i}sico, l{\'o}gico, humano o informacional que puede ser utilizado, gestionado o consumido por un sistema para cumplir un objetivo determinado. Un \emph{recurso} puede incluir hardware (procesamiento, memoria, almacenamiento, energ{\'i}a), software (servicios, bibliotecas, licencias), datos, tiempo de ejecuci{\'o}n o capacidad de red. En el contexto de \gls{iot} y de los \glspl{sibiot}, los recursos abarcan tanto los componentes de los \glspl{dispositivoconectado} como los servicios de \gls{nube}, plataformas de \gls{procesamientodatos} y capacidades anal{\'i}ticas, cuya gesti{\'o}n eficiente resulta esencial para garantizar \gls{escalabilidad}, \gls{rendimiento}, \gls{disponibilidad} y \gls{continuidad} operativa del sistema},
		type=main,
		symbol={📦},
		user1={Elemento utilizable por un sistema},
		user2={Capacidad f{\'i}sica o l{\'o}gica},
		user3={Objeto de gesti{\'o}n y asignaci{\'o}n},
		see={infraestructura,servicionube,procesamientodatos,escalabilidad,sibiot}
	}
	
	\newglossaryentry{red}{
		name={Red},
		text={red},
		sort={red},
		plural={redes},
		first={red},
		firstplural={redes},
		description={Conjunto de equipos denominados \glspl{nodo} y software conectados entre s{\'i} por medio de dispositivos f{\'i}sicos que env{\'i}an y reciben impulsos el{\'e}ctricos, ondas electromagn{\'e}ticas u otros medios para el transporte de datos, con el fin de compartir informaci{\'o}n, recursos y ofrecer servicios. Las \emph{redes} pueden clasificarse seg{\'u}n su alcance (LAN, WAN, MAN), su topolog{\'i}a (estrella, malla, {\'a}rbol) o su tecnolog{\'i}a (al{\'a}mbrica o inal{\'a}mbrica).},
		type=main,
		symbol={🌐},
		user1={Conjunto de nodos y enlaces},
		user2={Soporte para compartir recursos e informaci{\'o}n},
		user3={Infraestructura de comunicaci{\'o}n de datos},
		see={rediot,redinalambrica,redcelular,redelectrica}
	}
	
	\newglossaryentry{redcelular}{
		name={Red celular},
		text={red celular},
		sort={redcelular},
		plural={redes celulares},
		first={red celular},
		firstplural={redes celulares},
		description={Tipo de infraestructura de telecomunicaciones que divide un territorio en celdas geogr{\'a}ficas, cada una con su propia estaci{\'o}n base, permitiendo la \gls{comunicacioninalambrica} mediante tecnolog{\'i}as m{\'o}viles como \gls{gsm}, 3G, 4G y 5G. En sistemas \gls{iot}, las \emph{redes celulares} son fundamentales para garantizar la conectividad ubicua de dispositivos remotos, especialmente en aplicaciones como agricultura de precisi{\'o}n, log{\'i}stica, salud m{\'o}vil y veh{\'i}culos conectados.},
		type=main,
		symbol={📱},
		user1={Red divida en celdas geogr{\'a}ficas},
		user2={Soporte para tecnolog{\'i}as m{\'o}viles},
		user3={Conectividad IoT ubicua y remota},
		see={gsm,gprs,iot,rediot}
	}
	
	\newglossaryentry{redcomunicacion}{
		name={Red de comunicaci{\'o}n},
		text={red de comunicaci{\'o}n},
		sort={red de comunicacion},
		first={red de comunicaci{\'o}n},
		first={redes de comunicaci{\'o}n},
		description={Infraestructura y protocolos que permiten el intercambio de datos entre dispositivos y sistemas, incluyendo medios cableados e inal{\'a}mbricos, direccionamiento y servicios de transporte.},
		type=main,
		symbol={🕸️},
		user1={Interconexi{\'o}n},
		user2={Protocolos},
		user3={Transporte de datos},
		see={anchobanda,redinalambrica,websocket}
	}
	
	\newglossaryentry{redelectrica}{
		name={Red el{\'e}ctrica},
		text={red el{\'e}ctrica},
		sort={redelectrica},
		plural={redes el{\'e}ctricas},
		first={red el{\'e}ctrica},
		firstplural={redes el{\'e}ctricas},
		description={Infraestructura interconectada que tiene como prop{\'o}sito suministrar electricidad desde los generadores y proveedores hasta los consumidores finales. Incluye centrales de producci{\'o}n, subestaciones, sistemas de transmisi{\'o}n y distribuci{\'o}n, as{\'i} como dispositivos de control y protecci{\'o}n. En aplicaciones modernas, las \emph{redes el{\'e}ctricas inteligentes} integran sensores, sistemas de monitorizaci{\'o}n y \glspl{tecnologia} \gls{iot} para mejorar la \gls{eficiencia}, \gls{confiabilidad} y \gls{sostenibilidad} del suministro energ{\'e}tico.},
		type=main,
		symbol={⚡},
		user1={Infraestructura de generaci{\'o}n y distribuci{\'o}n de energ{\'i}a},
		user2={Sistema de transporte de electricidad},
		user3={Base de redes el{\'e}ctricas inteligentes (Smart Grids)},
		see={energia,sostenibilidad,rediot}
	}
	
	\newglossaryentry{redinalambrica}{
		name={Red inal{\'a}mbrica},
		text={red inal{\'a}mbrica},
		sort={redinalambrica},
		plural={redes inal{\'a}mbricas},
		first={red inal{\'a}mbrica},
		firstplural={redes inal{\'a}mbricas},
		description={Tipo de \gls{infraestructurared} que permite la transmisi{\'o}n de datos sin necesidad de cables f{\'i}sicos, mediante ondas electromagn{\'e}ticas. Las \emph{redes inal{\'a}mbricas} abarcan tecnolog{\'i}as como \gls{wifi}, \gls{zigbee}, \gls{lorawan}, \gls{nbiot} y \gls{bluetooth}. En el contexto de \gls{iot}, estas redes facilitan la \gls{conectividad} de \glspl{dispositivo} distribuidos, reducen costos de instalaci{\'o}n, permiten movilidad y son clave en escenarios que requieren bajo consumo energ{\'e}tico o cobertura en entornos remotos o urbanos.},
		type=main,
		symbol={📶},
		user1={Infraestructura de comunicaci{\'o}n sin cables},
		user2={Soporte de protocolos inal{\'a}mbricos},
		user3={Movilidad y flexibilidad en IoT},
		see={wifi,zigbee,lorawan,nbiot,bluetooth,rediot}
	}
	
	\newglossaryentry{redinteligente}{
		name={Red el{\'e}ctrica inteligente},
		text={\textit{Smart Grid}},
		sort={red electrica inteligente},
		plural={redes el{\'e}ctricas inteligentes},
		first={red el{\'e}ctrica inteligente (\textit{Smart Grid}},
		firstplural={\textit{Smart Grids}},
		description={Infraestructura digitalizada con medici{\'o}n avanzada, sensores y comunicaci{\'o}n bidireccional que optimiza eficiencia y confiabilidad del suministro.},
		type=main,
		symbol={⚡},
		see={hems,respuestademanda,gestionenergetica},
		user1={Bidireccionalidad},
		user2={Flexibilidad de demanda},
		user3={Integraci{\'o}n renovable}
	}
	
	\newglossaryentry{rediot}{
		name={Red \glsentrytext{iot}},
		text={red \glsentrytext{iot}},
		sort={red iot},
		plural={redes \glsentrytext{iot}},
		first={red \glsentrytext{iot}},
		firstplural={redes \glsentrytext{iot}},
		description={Infraestructura de comunicaci{\'o}n compuesta por \glspl{dispositivoconectado} que intercambian datos entre s{\'i} y con sistemas externos mediante protocolos heterog{\'e}neos. Una \emph{red \gls{iot}} integra \glspl{sensor}, \glspl{actuador}, \glspl{microcontrolador}, \glspl{gateway} y servicios en la \gls{nube} para habilitar la \gls{conectividad}, la \gls{interoperabilidad}, la \gls{monitorizacion} y el control remoto. Estas redes pueden emplear tecnolog{\'i}as como \gls{wifi}, \gls{zigbee}, \gls{lorawan}, \gls{nbiot}, \gls{bluetooth} o 5G, y resultan esenciales en aplicaciones distribuidas de salud, manufactura, \gls{agriculturainteligente}, \gls{transporte} y \glspl{ciudadinteligente}.},
		type=main,
		symbol={🌐},
		user1={Conjunto de dispositivos y enlaces IoT},
		user2={Soporte de protocolos y pasarelas heterog{\'e}neas},
		user3={Base para monitorizaci{\'o}n y control distribuido},
		see={iot,conectividad,interoperabilidad,gateway,lorawan,zigbee,bluetooth,nube,sibiot}
	}
	
	\newglossaryentry{redneuronal}{
		name={Red neuronal},
		text={red neuronal},
		sort={red neuronal},
		plural={redes neuronales},
		first={red neuronal (\textit{Neuronal Netwaork}, NN)},
		firstplural={redes neuronales},
		description={Modelo computacional inspirado en la estructura y funcionamiento del cerebro humano, compuesto por nodos o \emph{neuronas artificiales} interconectadas organizadas en capas (entrada, ocultas y salida). Una \emph{red neuronal} procesa informaci{\'o}n ajustando los pesos de sus conexiones durante el entrenamiento, lo que le permite aprender patrones y realizar tareas como clasificaci{\'o}n, regresi{\'o}n, detecci{\'o}n de anomal{\'i}as o generaci{\'o}n de predicciones. En el contexto de la \gls{ia} y el \gls{aprendizajeautomatico}, las redes neuronales constituyen la base del \textit{deep learning}. Aplicadas en los \glspl{sibiot}, permiten analizar grandes vol{\'u}menes de datos de \glspl{sensor} en \gls{tiemporeal}, optimizar procesos, detectar eventos cr{\'i}ticos y habilitar capacidades de \gls{modelopredictivo} y \gls{tomadecisiones} inteligente.},
		type=main,
		symbol={🧠},
		user1={Modelo inspirado en el cerebro humano},
		user2={Aprendizaje de patrones y representaciones},
		user3={Base del aprendizaje profundo (deep learning)},
		see={ia,aprendizajeautomatico,modelopredictivo,sibiot}
	}
	
		\newglossaryentry{redsensor}{
		name={Red de sensores},
		text={red de sensores},
		sort={red sensores},
		plural={redes de sensores},
		first={red de sensores},
		firstplural={redes de sensores},
		description={Conjunto de \glspl{sensor} interconectados que cooperan para recolectar y transmitir datos del entorno f{\'i}sico. Las \emph{redes de sensores} constituyen la capa de percepci{\'o}n de los \glspl{sibiot}, proporcionando datos primarios para su procesamiento posterior.},
		type=main,
		symbol={🌐},
		user1={Sensores distribuidos},
		user2={Captura de datos},
		user3={Percepci{\'o}n del entorno},
		see={infraestructuracomunicacion,integraciondatos}
	}
	
	\newglossaryentry{redsocial}{
		name={Red social},
		text={red social},
		sort={red social},
		plural={redes sociales},
		first={red social},
		firstplural={redes sociales},
		description={Plataforma digital que permite a los usuarios crear perfiles, compartir contenidos e interactuar con otros, facilitando la comunicaci{\'o}n, la colaboraci{\'o}n y la difusi{\'o}n de informaci{\'o}n en l{\'i}nea. Las \emph{redes sociales} pueden adoptar diferentes enfoques, como la conexi{\'o}n personal, profesional, acad{\'e}mica o tem{\'a}tica, y desempe{\~n}an un papel clave en la generaci{\'o}n de comunidades virtuales. En el contexto de los \glspl{sibiot}, las redes sociales tambi{\'e}n se utilizan para integrar flujos de datos, fomentar la interacci{\'o}n con usuarios finales y apoyar la difusi{\'o}n de informaci{\'o}n en \gls{tiemporeal}.},
		type=main,
		symbol={🌐},
		user1={Plataforma digital de interacci{\'o}n},
		user2={Difusi{\'o}n y compartici{\'o}n de contenidos},
		user3={Soporte a comunidades virtuales e IoT},
		see={servicio,entorno,sibiot}
	}
	
	\newglossaryentry{redundancia}{
		name={Redundancia},
		text={redundancia},
		sort={redundancia},
		plural={redundancias},
		first={redundancia},
		firstplural={redundancias},
		description={Dise{\~n}o e inclusi{\'o}n deliberada de \glspl{componente}, recursos o caminos alternativos que duplican funciones cr{\'i}ticas de un \gls{sistema} para mantener el servicio ante fallos parciales. La \emph{redundancia} puede aplicarse a nivel de hardware (m{\'u}ltiples \glspl{nodo}, fuentes de energ{\'i}a, \glspl{sistemarespaldoenergia}), software (instancias redundantes, \glspl{serviciodistribuido}), datos (\gls{replicacion} y copias de seguridad) y comunicaciones (enlaces y rutas alternas). En el contexto de \glspl{sibiot}, la redundancia mejora la \gls{disponibilidad}, la \gls{fiabilidad} y la \gls{toleranciafallos}, y suele coordinarse con \gls{sincronizacion} y pol{\'i}ticas de \gls{consistencia} para evitar p{\'e}rdidas o divergencias de informaci{\'o}n, manteniendo a la vez una \gls{eficienciaoperativa} adecuada frente al coste adicional.},
		type=main,
		symbol={♻️},
		user1={Duplicaci{\'o}n planificada de recursos cr{\'i}ticos},
		user2={Continuidad operativa frente a fallos},
		user3={Aplicable a hardware, software, datos y redes},
		see={toleranciafallos,disponibilidad,fiabilidad,replicacion,sincronizacion,consistencia,eficienciaoperativa,sistemarespaldoenergia,sibiot}
	}
	
	\newglossaryentry{redundante}{
		name={Redundante},
		text={redundante},
		sort={redundante},
		description={%
			Caracter{\'i}stica de un dato, evento, mensaje o componente que se repite innecesariamente sin aportar informaci{\'o}n adicional relevante para el funcionamiento o la toma de decisiones de un sistema.
			La \gls{redundancia} puede manifestarse como la transmisi{\'o}n reiterada de valores invariables, la generaci{\'o}n de eventos duplicados o la replicaci{\'o}n no controlada de informaci{\'o}n, lo que incrementa el consumo de recursos de red, almacenamiento y procesamiento.
			En sistemas de informaci{\'o}n basados en IoT y SIbIoT, la gesti{\'o}n de informaci{\'o}n redundante es un aspecto cr{\'i}tico, abordado mediante t{\'e}cnicas de \gls{preprocesamiento}, filtrado y agregaci{\'o}n de eventos, con el fin de optimizar la eficiencia operativa y preservar la escalabilidad del sistema.
		},
		first={redundante},
		plural={redundantes},
		firstplural={redundantes},
		see={preprocesamiento,filtradodatos,mensajeriatiemporeal,procesamientodistribuido}
	}
	
	\newglossaryentry{refactorizacion}{
		name={Refactorizaci{\'o}n},
		text={refactorizaci{\'o}n},
		sort={refactorizacion},
		first={refactorizaci{\'o}n},
		description={Transformaci{\'o}n interna del c{\'o}digo para mejorar su estructura, legibilidad y mantenibilidad sin alterar el comportamiento externo. Es una pr{\'a}ctica clave para reducir deuda t{\'e}cnica.},
		type=main,
		symbol={🧹},
		user1={Mejora interna},
		user2={Mantenibilidad},
		user3={Deuda t{\'e}cnica},
		see={controlversiones,calidadsoftware,testabilidad}
	}
	
	\newglossaryentry{regla}{
		name={Regla},
		text={regla},
		plural={reglas},
		first={regla},
		firstplural={reglas},
		see={entrenamientomodelos,modelo},
		sort={regla},
		description={Condici{\'o}n o conjunto de condiciones l{\'o}gicas que determinan una acci{\'o}n cuando se cumplen ciertos criterios. En SIbIoT, las reglas se emplean para generar alertas, activar respuestas autom{\'a}ticas o guiar decisiones, y pueden ser ajustadas en funci{\'o}n del desempe{\~n}o observado y del contexto empresarial}
	}
	
	\newglossaryentry{reglavalidacion}{
		name={regla de validaci{\'o}n},
		text={regla de validaci{\'o}n},
		sort={reglavalidacion},
		plural={reglas de validaci{\'o}n},
		first={regla de validaci{\'o}n},
		firstplural={reglas de validaci{\'o}n},
		description={Las reglas de validaci{\'o}n son directrices utilizadas para garantizar que los datos sean precisos, coherentes y completos. Pueden incluir l{\'i}mites de valores, formatos esperados, tipos de datos v{\'a}lidos o relaciones l{\'o}gicas entre campos. En los \glspl{sibiot}, estas reglas se aplican autom{\'a}ticamente durante la \gls{capturadatos} o la \gls{limpieza} para prevenir errores y asegurar que la informaci{\'o}n recolectada sea confiable. La correcta definici{\'o}n de las reglas de validaci{\'o}n es fundamental para mantener la \gls{integridad} y la \gls{consistencia} de los conjuntos de datos.},
		type=\glsdefaulttype,
		symbol={🧩},
		user1={Validation rule},
		user2={Condici{\'o}n o criterio de control de calidad de los datos},
		user3={Conjunto de par{\'a}metros que definen la validez de un registro},
		see={validacion,limpieza,integridad,consistencia,calidaddatos}
	}
	
	\newglossaryentry{referencia}
	{
		name={Referencia},
		text={referencia},
		sort={referencia},
		plural={referencias},
		first={referencia},
		firstplural={referencias},
		description={Una \emph{referencia} es un elemento que remite o enlaza a otro contenido, recurso o entidad. Puede entenderse en distintos contextos:
			\begin{itemize}
				\item En el {\'a}mbito acad{\'e}mico, como la cita o menci{\'o}n de una fuente bibliogr{\'a}fica que respalda un argumento o informaci{\'o}n.
				\item En ingenier{\'i}a de software, como el puntero o enlace a un objeto, variable o recurso dentro de un \gls{sistemainformatico}.
				\item En \glspl{sibiot}, como la relaci{\'o}n entre \glspl{dato}, dispositivos y procesos que permiten trazar la procedencia, la integridad o el contexto de la informaci{\'o}n.
				\item En documentaci{\'o}n t{\'e}cnica, como el conjunto de directrices o especificaciones que sirven de gu{\'i}a para dise{\~n}o, implementaci{\'o}n o verificaci{\'o}n de sistemas.
			\end{itemize}
			Su funci{\'o}n principal es garantizar la trazabilidad, la claridad y la confiabilidad en la comunicaci{\'o}n de la informaci{\'o}n.},
		type=\glsdefaulttype,
		symbol={📑},
		user1={Cita},
		user2={Referencia bibliogr{\'a}fica},
		user3={Enlace},
		see={dato}
	}
	
	\newglossaryentry{releaseplan}{
		name={Release Plan},
		text={\textit{Release Plan}},
		sort={release plan},
		plural={\textit{Release Plans}},
		first={\textit{Release Plan}},
		firstplural={\textit{Release Plans}},
		description={Plan que define c{\'o}mo y cu{\'a}ndo se liberar{\'a}n incrementos del producto hacia usuarios finales. Integra previsiones de alcance, dependencias y cronograma de entregas.},
		type=main,
		symbol={🚀},
		user1={Planificaci{\'o}n de liberaciones},
		user2={Coordinaci{\'o}n de versiones},
		user3={Gesti{\'o}n de despliegues},
		see={incremento,roadmap,sprint}
	}
	
	\newglossaryentry{rendimiento}{
		name={Rendimiento},
		text={rendimiento},
		sort={rendimiento},
		plural={rendimientos},
		first={rendimiento (\textit{throughput})},
		firstplural={rendimientos (\textit{throughput})},
		description={Medida de la cantidad de datos o trabajo que un \gls{sistema} es capaz de procesar, transmitir o completar en un intervalo de tiempo determinado, expresada habitualmente en unidades como bits por segundo (bps), transacciones por segundo (TPS), mensajes por segundo o unidades por hora, seg{\'u}n el contexto. En telecomunicaciones y \gls{iot}, el \emph{rendimiento} refleja la \emph{tasa de transferencia efectiva}, es decir, la velocidad real alcanzada, que puede ser inferior al \gls{anchobanda} te{\'o}rico debido a factores como latencia, congesti{\'o}n, p{\'e}rdida de paquetes o sobrecarga de protocolo. Es un indicador clave para evaluar la \gls{eficiencia} y capacidad operativa de \glspl{red}, \glspl{sistema} de \glspl{procesamientodatos} y \glspl{plataforma} distribuidas.},
		type=main,
		symbol={},
		user1={M{\'e}trica de desempe{\~n}o},
		user2={Tasa de transferencia efectiva},
		user3={Indicador clave de eficiencia},
		see={anchobanda,latencia,escalable}
	}
	
	\newglossaryentry{rendimientodelared}{
		name={Rendimiento de la red},
		text={rendimiento de la red},
		sort={rendimiento de la red},
		plural={rendimientos de la red},
		first={rendimiento de la red},
		firstplural={rendimientos de la red},
		description={Medida de la eficiencia con que una \gls{red} transmite datos entre \glspl{nodo} o sistemas. Se eval{\'u}a mediante indicadores como el \gls{anchobanda}, la latencia, la p{\'e}rdida de paquetes, el \textit{jitter} y la \gls{disponibilidad}. En entornos \gls{iot}, un buen \emph{rendimiento de la red} es crucial para garantizar la entrega oportuna y fiable de informaci{\'o}n, especialmente en aplicaciones sensibles al tiempo como \gls{monitorizacion}, \gls{control} remoto o automatizaci{\'o}n distribuida.},
		type=main,
		symbol={📶},
		user1={Indicador de eficiencia de transmisi{\'o}n},
		user2={Eval{\'u}a latencia, jitter y p{\'e}rdida},
		user3={Fundamental para aplicaciones IoT cr{\'i}ticas},
		see={red,anchobanda,disponibilidad,monitorizacion,control}
	}
	
	\newglossaryentry{replica}{
		name={R{\'e}plica},
		text={r{\'e}plica},
		sort={replica},
		plural={r{\'e}plicas},
		first={r{\'e}plica},
		firstplural={r{\'e}plicas},
		description={Copia exacta de datos, procesos o servicios almacenada en un \gls{sistema} para garantizar \gls{disponibilidad}, \gls{redundancia} y \gls{consistencia}. En bases de datos distribuidas o sistemas en la \gls{nube}, una \emph{r{\'e}plica} permite tolerancia a fallos, escalabilidad de lectura y protecci{\'o}n frente a p{\'e}rdidas.},
		type=main,
		symbol={📂},
		user1={Copia exacta en sistemas distribuidos},
		user2={Soporte para redundancia y disponibilidad},
		user3={Escalabilidad de lectura y tolerancia a fallos},
		see={replicacion,redundancia,consistencia,nube}
	}
	
	\newglossaryentry{replicacion}{
		name={Replicaci{\'o}n},
		text={replicaci{\'o}n},
		sort={replicacion},
		plural={replicaciones},
		first={replicaci{\'o}n},
		firstplural={replicaciones},
		description={Proceso mediante el cual los datos, procesos o servicios de un \gls{sistema} se copian y mantienen en m{\'u}ltiples ubicaciones o nodos para garantizar \gls{disponibilidad}, \gls{fiabilidad} y continuidad operativa. La \emph{replicaci{\'o}n} puede ser sincr{\'o}nica (los cambios se reflejan inmediatamente en todas las copias) o asincr{\'o}nica (los cambios se propagan con cierto retardo). En el contexto de las \glspl{basedatosdistribuida} y de los \glspl{sibiot}, la replicaci{\'o}n asegura que los datos generados por \glspl{sensor} y procesados en tiempo real est{\'e}n disponibles incluso ante fallos de un \gls{nodo} o de la red, contribuyendo a la \gls{redundancia}, la tolerancia a fallos y la \gls{resiliencia} del sistema.},
		type=main,
		symbol={🔁},
		user1={Copia y distribuci{\'o}n de datos o servicios},
		user2={Soporte a disponibilidad y continuidad},
		user3={Mecanismo de redundancia en SIbIoT},
		see={redundancia,fiabilidad,disponibilidad,nodo,basedatosdistribuida,sibiot}
	}
	
	\newglossaryentry{repositorio}
	{
		name={Repositorio},
		text={repositorio},
		sort={repositorio},
		plural={repositorios},
		first={Repositorio},
		description={Entorno o sistema de almacenamiento estructurado destinado a conservar, organizar y gestionar informaci{\'o}n digital, c{\'o}digo fuente, documentos o datos de manera centralizada y segura. 
			En el contexto de la ingenier{\'i}a del software, un repositorio almacena versiones de proyectos y facilita la colaboraci{\'o}n mediante sistemas de control de versiones como \textit{Git}, garantizando la trazabilidad, la integridad y la recuperaci{\'o}n de los cambios. 
			En los \glspl{sibiot}, los repositorios se utilizan para almacenar datos generados por \glspl{sensor}, modelos de \gls{aprendizajeautomatico}, configuraciones de \glspl{dispositivo} y registros hist{\'o}ricos, permitiendo el an{\'a}lisis y la reutilizaci{\'o}n de la informaci{\'o}n. 
			Tambi{\'e}n pueden funcionar como plataformas de datos abiertos o federados que facilitan la interoperabilidad y el intercambio de conocimiento entre sistemas inteligentes.},
		symbol={🗄️},
		user1={Sistema estructurado de almacenamiento de informaci{\'o}n o c{\'o}digo},
		user2={Permite trazabilidad, control de versiones y colaboraci{\'o}n},
		user3={Soporte para datos, modelos y configuraciones en SIbIoT},
		see={[Relacionado con:]{basedatos,dato,controlversiones,sibiot}}
	}
	
	\newglossaryentry{representacion}{
		name={Representaci{\'o}n},
		text={representaci{\'o}n},
		sort={representacion},
		first={representaci{\'o}n},
		description={Forma en que informaci{\'o}n, conceptos o estructuras se expresan mediante modelos, diagramas, c{\'o}digo o notaciones, para su comprensi{\'o}n y comunicaci{\'o}n.},
		type=main,
		symbol={🧠},
		user1={Modelos},
		user2={Notaciones},
		user3={Comunicaci{\'o}n},
		see={diagrama,documentacion}
	}
	
	\newglossaryentry{reprocesamiento}{
		name={Reprocesamiento},
		text={reprocesamiento},
		sort={reprocesamiento},
		description={Proceso mediante el cual un sistema vuelve a ejecutar operaciones de an{\'a}lisis, transformaci{\'o}n o validaci{\'o}n sobre datos previamente capturados y almacenados. En sistemas de informaci{\'o}n basados en IoT y arquitecturas orientadas a eventos, el reprocesamiento se utiliza para corregir errores, aplicar nuevos modelos anal{\'i}ticos, recalcular m{\'e}tricas hist{\'o}ricas o responder a cambios en reglas de negocio sin necesidad de reacquirir los datos originales. Este mecanismo resulta fundamental para la auditor{\'i}a, la trazabilidad y la evoluci{\'o}n de los modelos de decisi{\'o}n en sistemas distribuidos},
		plural={reprocesamientos},
		first={reprocesamiento},
		firstplural={reprocesamientos},
		see={procesamientodistribuido,orientadoeventos,gestiondatos,trazabilidad}
	}
	
	\newglossaryentry{requisito}
	{
		name={Requisito},
		text={requisito},
		sort={requisito},
		plural={requisitos},
		first={requisito},
		firstplural={requisitos},
		description={Un requisito es una necesidad, condici{\'o}n o funci{\'o}n que un \gls{sistema} o \gls{componente} debe satisfacer para garantizar su aceptaci{\'o}n y correcto funcionamiento. En la ingenier{\'i}a de software y en los \glspl{sibiot}, los requisitos se clasifican en funcionales y no funcionales, siendo fundamentales para el dise{\~n}o, desarrollo y evaluaci{\'o}n del \gls{sistemainformatico}.},
		type=\glsdefaulttype,
		symbol={},
		user1={Necesidad},
		user2={Exigencia},
		user3={Especificaci{\'o}n},
		see={condicion}
	}
	
	\newglossaryentry{res}{
		name={Res (res)},
		text={res},
		sort={res},
		plural={carpetas res},
		description={Carpeta de recursos en proyectos de software, especialmente en entornos de desarrollo de aplicaciones Android. Dentro de \texttt{res} se organizan elementos como im{\'a}genes (\texttt{drawable}), dise{\~n}os de interfaz de usuario (\texttt{layout}), definiciones de cadenas de texto (\texttt{values}), estilos, temas, archivos XML y otros recursos est{\'a}ticos necesarios para la ejecuci{\'o}n de la aplicaci{\'o}n. Esta carpeta permite separar la l{\'o}gica del c{\'o}digo de los activos visuales y de configuraci{\'o}n, facilitando la mantenibilidad, la internacionalizaci{\'o}n y la modularidad del sistema.},
		type=\glsdefaulttype,
		symbol={📁},
		user1={Carpeta de recursos en Android},
		user2={Directorio de im{\'a}genes, XML y dise{\~n}os},
		user3={Soporte para internacionalizaci{\'o}n y configuraci{\'o}n},
		see={xml,layout,drawable,values,assets}
	}
	
	\newglossaryentry{resiliencia}{
		name={Resiliencia},
		text={resiliencia},
		sort={resiliencia},
		plural={resiliencias},
		first={resiliencia},
		firstplural={resiliencias},
		description={Capacidad de un \gls{sistema}, organizaci{\'o}n o tecnolog{\'i}a para anticipar, resistir, adaptarse y recuperarse r{\'a}pidamente frente a perturbaciones, incidentes o fallos, manteniendo un nivel aceptable de servicio. La \emph{resiliencia} integra conceptos de \gls{fiabilidad}, \gls{disponibilidad}, \gls{seguridad} y \gls{adaptabilidad}, y se alcanza mediante estrategias como redundancia, tolerancia a fallos, elasticidad, monitorizaci{\'o}n proactiva y planes de recuperaci{\'o}n ante desastres. En los \glspl{sibiot}, la resiliencia resulta esencial para garantizar la continuidad de servicios cr{\'i}ticos, la protecci{\'o}n de datos en \gls{tiemporeal} y la confianza de los usuarios en entornos distribuidos y cambiantes.},
		type=main,
		symbol={🛡️},
		user1={Propiedad de calidad de sistemas},
		user2={Continuidad y recuperaci{\'o}n},
		user3={Tolerancia a perturbaciones},
		see={fiabilidad,disponibilidad,seguridad,adaptabilidad,resiliente}
	}
	
	\newglossaryentry{resiliente}{
		name={Resiliente},
		text={resiliente},
		sort={resiliente},
		plural={resilientes},
		first={resiliente},
		firstplural={resilientes},
		description={Adjetivo que describe a un \gls{sistema}, componente u organizaci{\'o}n que posee \gls{resiliencia}, es decir, que puede adaptarse, recuperarse o continuar funcionando ante perturbaciones o fallos.},
		type=main,
		symbol={🛡️},
		user1={Caracter{\'i}stica derivada de resiliencia},
		user2={Atributo de sistemas tolerantes a fallos},
		user3={Adjetivo t{\'e}cnico de calidad},
		see={resiliencia}
	}
	
	\newglossaryentry{respaldo}{
		name={respaldo},
		text={respaldo},
		sort={respaldo},
		plural={respaldos},
		first={respaldo},
		firstplural={respaldos},
		description={El respaldo es un procedimiento esencial dentro de la gesti{\'o}n de la informaci{\'o}n que busca asegurar la disponibilidad y recuperaci{\'o}n de los datos en caso de incidentes. En los \glspl{sibiot}, los respaldos peri{\'o}dicos garantizan la continuidad operativa frente a interrupciones, errores humanos o ciberataques. Pueden realizarse de forma local, remota o en la nube, y se acompa{\~n}an de pol{\'i}ticas de restauraci{\'o}n y verificaci{\'o}n para garantizar su eficacia. La combinaci{\'o}n de respaldo y \gls{encriptaciondatos} conforma una estrategia integral de protecci{\'o}n de la informaci{\'o}n.},
		type=\glsdefaulttype,
		symbol={💾},
		user1={Backup},
		user2={Copia de seguridad de la informaci{\'o}n},
		user3={Medio de protecci{\'o}n y recuperaci{\'o}n de datos ante incidentes o fallos},
		see={seguridad,encriptaciondatos,calidaddatos,recuperacion}
	}
	
	\newglossaryentry{responsabilidadorganizacional}{
		name={Responsabilidad organizacional},
		text={responsabilidad organizacional},
		sort={responsabilidad organizacional},
		plural={responsabilidades organizacionales},
		first={responsabilidad organizacional},
		firstplural={responsabilidades organizacionales},
		description={Compromiso asumido por una organizaci{\'o}n para gestionar sus actividades de forma {\'e}tica, transparente y sostenible, considerando el impacto econ{\'o}mico, social y ambiental de sus decisiones y operaciones. La \emph{responsabilidad organizacional} implica garantizar la protecci{\'o}n de los derechos de empleados, clientes y comunidades, promover la innovaci{\'o}n responsable y cumplir con normativas y est{\'a}ndares internacionales. En el contexto de los \glspl{sibiot}, se traduce en el dise{\~n}o e implementaci{\'o}n de sistemas que respeten la \gls{privacidad}, aseguren la \gls{seguridad}, fomenten la \gls{sostenibilidad} y generen confianza en los usuarios y la sociedad.},
		type=main,
		symbol={🏛️},
		user1={Compromiso {\'e}tico y sostenible},
		user2={Gest{\'i}on del impacto social y ambiental},
		user3={Marco de actuaci{\'o}n para SIbIoT},
		see={sostenibilidad,privacidad,seguridad,ingeniero,sibiot}
	}
	
	\newglossaryentry{respuestademanda}{
		name={Respuesta a la demanda},
		text={respuesta a la demanda},
		sort={respuesta demanda},
		plural={programas de respuesta a la demanda},
		first={respuesta a la demanda},
		firstplural={programas de respuesta a la demanda},
		description={Ajuste de consumos el{\'e}ctricos seg{\'u}n precios o se{\~n}ales de red, desplazando cargas para reducir picos y costes.},
		type=main,
		symbol={🔄},
		see={redinteligente,hems},
		user1={Tarifas din{\'a}micas},
		user2={Desplazamiento de carga},
		user3={Reducci{\'o}n de picos}
	}
	
	\newglossaryentry{rest}{
		name={REST},
		text={REST},
		sort={rest},
		plural={estilos REST},
		first={transferencia de estado representacional (\textit{Representational State Transfer}, REST)},
		firstplural={estilos REST},
		description={Estilo arquitect{\'o}nico para sistemas distribuidos que define un conjunto de restricciones (cliente--servidor, sin estado, capacidad de \textit{cache}, interfaz uniforme, sistema por capas y, opcionalmente, c{\'o}digo bajo demanda) para construir servicios interoperables en la web. En la pr{\'a}ctica, los servicios \emph{REST} exponen recursos identificados por \texttt{\gls{uri}} y manipulados con m{\'e}todos de \gls{http} (p.\,ej., \texttt{GET}, \texttt{POST}, \texttt{PUT}, \texttt{DELETE}), devolviendo representaciones en formatos como \gls{json} o \gls{xml}. Su adopci{\'o}n favorece la \gls{interoperabilidad}, la \gls{escalabilidad} y el bajo acoplamiento.},
		type=main,
		symbol={🔗},
		user1={Estilo arquitect{\'o}nico para servicios web},
		user2={Recursos via URI y m{\'e}todos HTTP},
		user3={Representaciones en JSON o XML},
		see={http,api,interoperabilidad,escalabilidad}
	}
	
	\newglossaryentry{restful}{
		name={RESTful},
		text={RESTful},
		sort={restful},
		plural={RESTful},
		first={RESTful},
		firstplural={RESTful},
		description={Calificativo de una \gls{api} o servicio web que implementa de forma coherente los principios del estilo \gls{rest}. Una \emph{API RESTful} organiza sus recursos en \texttt{URI} bien dise{\~n}ados, usa m{\'e}todos de \gls{http} sem{\'a}nticos, c{\'o}digos de estado est{\'a}ndar, cabeceras para \textit{cache}/negociaci{\'o}n de contenido y representaciones comunes como \gls{json} o \gls{xml}. En entornos empresariales e \gls{iot}, facilita la integraci{\'o}n segura entre aplicaciones heterog{\'e}neas y la evoluci{\'o}n independiente de clientes y servidores.},
		type=main,
		symbol={🧭},
		user1={Implementaci{\'o}n pr{\'a}ctica de REST},
		user2={Recursos, m{\'e}todos HTTP y c{\'o}digos de estado},
		user3={Integraci{\'o}n ligera entre sistemas heterog{\'e}neos},
		see={rest,http,api,interoperabilidad}
	}
	
	\newglossaryentry{restriccion}{
		name={Restricci{\'o}n},
		text={restricci{\'o}n},
		sort={restriccion},
		plural={restricciones},
		first={restricci{\'o}n},
		firstplural={restricciones},
		description={Condici{\'o}n o limitaci{\'o}n que delimita comportamiento, estructura, recursos o desempe{\~n}o de un sistema. Puede ser t{\'e}cnica, normativa, arquitect{\'o}nica o de dominio, y debe cumplirse durante todo el ciclo de vida del sistema.},
		type=main,
		symbol={⛓️},
		user1={Limitaci{\'o}n formal},
		user2={Condici{\'o}n obligatoria},
		user3={Control de cumplimiento},
		see={especificaciontecnica,decisionarquitectonica}
	}
	
	\newglossaryentry{resultado}{
		name={Resultado},
		text={resultado},
		sort={resultado},
		plural={resultados},
		first={resultado},
		firstplural={resultados},
		description={Producto, efecto o medici{\'o}n obtenida tras ejecutar un proceso o actividad. Los resultados se contrastan con \glspl{objetivo} y criterios de evaluaci{\'o}n.},
		type=main,
		symbol={🏁},
		user1={Evidencia},
		user2={Medici{\'o}n},
		user3={Evaluaci{\'o}n},
		see={objetivo,desempeno}
	}
	
	\newglossaryentry{resultadoconsolidado}{
		name={Resultado consolidado},
		text={resultado consolidado},
		plural={resultados consolidados},
		first={resultado consolidado},
		firstplural={resultados consolidados},
		see={procesamiento,dato,informacion,metodo},
		sort={resultadoconsolidado},
		description={Salida anal{\'i}tica derivada de la agregaci{\'o}n, integraci{\'o}n y s{\'i}ntesis de m{\'u}ltiples fuentes de datos y/o modelos. En SIbIoT, los resultados consolidados suelen materializarse como indicadores, reportes o tablas de servicio, y constituyen la base para la visualizaci{\'o}n y la toma de decisiones organizacionales}
	}
	
	\newglossaryentry{retroalimentacion}
	{
		name={Retroalimentaci{\'o}n},
		text={retroalimentaci{\'o}n},
		sort={retroalimentacion},
		plural={retroalimentaciones},
		first={retroalimentaci{\'o}n},
		firstplural={retroalimentaciones},
		description={La retroalimentaci{\'o}n es el proceso mediante el cual un \gls{sistema} recibe informaci{\'o}n sobre su propio desempe{\~n}o o resultados, con el fin de ajustar, corregir o mejorar su funcionamiento. En los \glspl{sibiot}, la retroalimentaci{\'o}n permite que \glspl{actuador} y \glspl{dispositivo} modifiquen su comportamiento en funci{\'o}n de los datos obtenidos por los \glspl{sensor}, garantizando mayor precisi{\'o}n, estabilidad y control en aplicaciones como la automatizaci{\'o}n o la rob{\'o}tica.},
		type=\glsdefaulttype,
		symbol={🔄},
		user1={Feedback},
		user2={Se{\~n}al de control},
		user3={Correcci{\'o}n de sistema},
		see={servomotor}
	}
	
	\newglossaryentry{retroalimentacionoperacional}{
		name={Retroalimentaci{\'o}n operacional},
		text={retroalimentaci{\'o}n operacional},
		sort={retroalimentacion operacional},
		plural={retroalimentaciones operacionales},
		first={retroalimentaci{\'o}n operacional},
		firstplural={retroalimentaciones operacionales},
		description={Tipo de retroalimentaci{\'o}n orientada espec{\'i}ficamente al ajuste de procesos operativos en tiempo real o cuasi real, basada en datos obtenidos durante la ejecuci{\'o}n de actividades productivas o funcionales. En sistemas basados en IoT y SIbIoT, la retroalimentaci{\'o}n operacional permite modificar par{\'a}metros de operaci{\'o}n, redistribuir recursos, activar alertas o ajustar configuraciones autom{\'a}ticamente, contribuyendo a la eficiencia, resiliencia y estabilidad del sistema.},
		type=main,
		symbol={⚙️},
		user1={Ajuste en tiempo real},
		user2={Optimizaci{\'o}n operativa},
		user3={Control autom{\'a}tico},
		see={retroalimentacion,eficiencia,infraestructuradistribuida}
	}
	
	\newglossaryentry{reuse}{
		name={Reutilizaci{\'o}n},
		text={reutilizaci{\'o}n},
		sort={reutilizacion},
		plural={reutilizaciones},
		first={reutilizaci{\'o}n},
		firstplural={reutilizaciones},
		description={Uso sistem{\'a}tico de \glspl{componente}, artefactos o conocimientos existentes (c{\'o}digo fuente, bibliotecas, \glspl{servicio}, \glspl{modelo}, \glspl{patron}, casos de prueba, documentaci{\'o}n y plantillas) para construir nuevos \glspl{sistema} o ampliar los ya existentes. La \emph{reutilizaci{\'o}n} reduce costos y tiempo de desarrollo, mejora la \gls{calidad} y favorece la \gls{mantenibilidad}, al promover menor duplicaci{\'o}n y mayor estandarizaci{\'o}n. En \glspl{sibiot}, impulsa la integraci{\'o}n de m{\'o}dulos de \gls{hardware}/\gls{software} y \glspl{api} comunes, as{\'i} como la adopci{\'o}n de buenas pr{\'a}cticas.},
		type=main,
		symbol={♻️},
		user1={Aprovechamiento de artefactos existentes},
		user2={Reducci{\'o}n de costo y tiempo de desarrollo},
		user3={Estandarizaci{\'o}n y calidad del producto},
		see={buenaspracticas,ingenieriasoftware,componente,api,servicio,sibiot}
	}
	
	\newglossaryentry{rfid}{
		name={RFID},
		text={RFID},
		sort={rfid},
		plural={tecnolog{\'i}as RFID},
		first={identificaci{\'o}n por radiofrecuencia (\textit{Radio Frequency Identification}, RFID)},
		firstplural={tecnolog{\'i}as RFID},
		description={\Gls{tecnologia} de \gls{comunicacioninalambrica} que permite la identificaci{\'o}n autom{\'a}tica de objetos, animales o personas mediante etiquetas (tags) que responden a un lector a trav{\'e}s de se{\~n}ales de radio. Un sistema \emph{RFID} se compone de etiqueta/transpondedor, lector e antena; puede operar en modos pasivo, activo o semipasivo y en diferentes bandas (LF, HF, UHF). Se utiliza ampliamente en \gls{logistica}, inventarios, \gls{trazabilidad}, \gls{controlacceso} y \gls{iot}, al posibilitar captura de datos en \gls{tiemporeal} sin contacto f{\'i}sico ni l{\'i}nea de visi{\'o}n directa.},
		type=main,
		symbol={📡},
		user1={Identificaci{\'o}n sin contacto por radiofrecuencia},
		user2={Etiquetas pasivas/activas y lectores dedicados},
		user3={Aplicaciones en trazabilidad, acceso e IoT},
		see={nfc,bluetooth,comunicacioninalambrica,logistica,controlacceso,iot}
	}
	
	\newglossaryentry{risc}{
		name={RISC},
		text={RISC},
		plural={RISCs},
		firstplural={RISCs},
		first={computador con conjunto de instrucciones reducido (\textit{Reduced Instruction Set Computer}, RISC)},
		sort={risc},
		see={arm,cpu,arquitecturaprocesador},
		description={Paradigma de arquitectura de procesadores basado en un conjunto reducido y altamente optimizado de instrucciones simples, dise{\~n}adas para ejecutarse r{\'a}pidamente y de manera eficiente. La arquitectura RISC prioriza la ejecuci{\'o}n en uno o pocos ciclos de reloj, el uso intensivo de registros y la segmentaci{\'o}n de instrucciones (pipeline), lo que permite un alto rendimiento y eficiencia energ{\'e}tica. Este enfoque es la base de arquitecturas ampliamente utilizadas como ARM y es com{\'u}n en microcontroladores, sistemas embebidos, dispositivos IoT y sistemas ciberf{\'i}sicos.}
	}
	
	\newglossaryentry{riesgo}{
		name={Riesgo},
		text={riesgo},
		sort={riesgo},
		plural={riesgos},
		first={riesgo},
		firstplural={riesgos},
		description={Probabilidad de que una amenaza explote una vulnerabilidad causando un impacto negativo en activos, procesos o personas. En ingenier{\'i}a de software y sistemas \gls{iot}, el \emph{riesgo} se relaciona con la seguridad, la fiabilidad y la disponibilidad de los sistemas, y requiere identificaci{\'o}n, an{\'a}lisis, evaluaci{\'o}n y tratamiento. La gesti{\'o}n de riesgos incluye medidas preventivas, de mitigaci{\'o}n y de respuesta para reducir la probabilidad o el impacto de incidentes.},
		type=main,
		symbol={⚠️},
		user1={Probabilidad de p{\'e}rdida o da{\~n}o},
		user2={Gest{\'i}on de riesgos},
		user3={An{\'a}lisis de riesgos},
		see={amenaza,vulnerabilidad,seguridad,fiabilidad}
	}
	
	\newglossaryentry{roadmap}{
		name={Roadmap},
		text={\textit{Roadmap}},
		sort={roadmap},
		plural={\textit{Roadmaps}},
		first={\textit{Roadmap}},
		firstplural={\textit{Roadmaps}},
		description={Plan de alto nivel que describe la evoluci{\'o}n prevista del producto en el tiempo, alineando objetivos estrat{\'e}gicos, versiones y prioridades. Sirve como instrumento de comunicaci{\'o}n entre stakeholders y equipo de desarrollo.},
		type=main,
		symbol={🗺️},
		user1={Plan estrat{\'e}gico del producto},
		user2={Visi{\'o}n de evoluci{\'o}n},
		user3={Comunicaci{\'o}n con stakeholders},
		see={productgoal,productbacklog}
	}
	
	% ===================== Robustez =====================
	\newglossaryentry{robustez}{
		name={Robustez},
		text={robustez},
		sort={robustez},
		plural={robusteces},
		first={robustez},
		firstplural={robusteces},
		description={La \emph{robustez} es la propiedad de un \gls{sistema} o componente que le permite seguir funcionando de manera correcta aun cuando se enfrenta a perturbaciones, fallos parciales, condiciones extremas o entradas inesperadas. En sistemas \gls{iot}, la robustez implica tolerancia a fallos de \glspl{dispositivoiot}, variaciones en la \gls{latencia}, interferencias en la \gls{comunicacion} o cambios en el \gls{entorno}. Una alta robustez asegura continuidad operativa, fiabilidad y confianza en la \gls{tomadecisiones} basada en datos.},
		type=main,
		symbol={🛡️},
		user1={Tolerancia a fallos y perturbaciones},
		user2={Mantenimiento del funcionamiento esperado},
		user3={Clave en IoT frente a entornos variables},
		see={robusto,fiabilidad,continuidad,estabilidad,sibiot}
	}
	
	% ===================== Robusto =====================
	\newglossaryentry{robusto}{
		name={Robusto},
		text={robusto},
		sort={robusto},
		plural={robustos},
		first={robusto},
		firstplural={robustos},
		description={El adjetivo \emph{robusto} califica a un \gls{sistema}, componente o proceso que se caracteriza por su alta \gls{robustez}, es decir, por mantener un funcionamiento confiable y estable ante fallos, sobrecargas o condiciones inesperadas. En \glspl{sibiot}, un nodo robusto puede continuar operando incluso con p{\'e}rdida de conectividad intermitente, interferencias de se{\~n}al o fallos parciales de \glspl{sensor}, garantizando \gls{fiabilidad} y \gls{continuidad} del servicio.},
		type=main,
		symbol={💪},
		user1={Calificativo de sistemas o componentes resistentes},
		user2={Mantiene operaci{\'o}n confiable en IoT},
		user3={Vinculado directamente con la robustez},
		see={robustez,fiabilidad,continuidad,estabilidad,sibiot}
	}
	
	\newglossaryentry{router}{
		name={Router},
		text={router},
		sort={router},
		plural={routers},
		first={router},
		firstplural={routers},
		description={Dispositivo de red encargado de dirigir el tr{\'a}fico de datos entre diferentes \glspl{red}, seleccionando las rutas m{\'a}s adecuadas para el env{\'i}o de paquetes. Un \emph{router} opera en la capa de red del modelo OSI y puede incluir funciones de encaminamiento din{\'a}mico, traducci{\'o}n de direcciones (NAT), asignaci{\'o}n de direcciones IP (DHCP) y medidas de \gls{seguridad}. En el contexto de los \glspl{sibiot}, los routers facilitan la conectividad de \glspl{sensor}, \glspl{actuador} y \glspl{nodo} a infraestructuras de comunicaci{\'o}n inal{\'a}mbricas y cableadas, integr{\'a}ndose con tecnolog{\'i}as como Wi-Fi, Zigbee, LoRaWAN o redes celulares, garantizando \gls{disponibilidad}, \gls{interoperabilidad} y \gls{eficienciaoperativa}.},
		type=main,
		symbol={📡},
		user1={Dispositivo de interconexi{\'o}n de redes},
		user2={Encaminamiento de paquetes y gesti{\'o}n de tr{\'a}fico},
		user3={Soporte a conectividad IoT y sistemas distribuidos},
		see={red,nodo,comunicacioninalambrica,seguridad,sibiot}
	}
	
	\newglossaryentry{rom}{
		name={ROM},
		text={ROM},
		sort={rom},
		firstplural={ROMs},
		plural={ROMs},
		first={memoria de s{\'o}lo lectura (\textit{Read-Only Memory}, ROM)},
		description={\emph{ROM} es un tipo de memoria no vol{\'a}til que contiene datos permanentes o semipermanentes, utilizada com{\'u}nmente para almacenar el \gls{firmware} o el sistema b{\'a}sico de arranque de un \gls{dispositivo}. Su contenido no se pierde al apagar el equipo y, en muchos casos, no puede ser modificado por el \gls{usuario} final.},
		type=main,
		symbol={💾},
		user1={Memoria no vol{\'a}til},
		user2={Almacenamiento de firmware y arranque},
		user3={Base para sistemas embebidos e IoT},
		see={firmware,cpu,memoria}
	}
		
	\newglossaryentry{rpl}{
		name={RPL},
		text={RPL},
		first={protocolo de enrutamiento para redes de baja potencia y con p{\'e}rdidas (\textit{Routing Protocol for Low-Power and Lossy Networks}, RPL)},
		firstplural={RPLs},
		plural={protocolos RPL},
		description={\emph{RPL} es un protocolo de enrutamiento dise{\~n}ado espec{\'i}ficamente para redes de sensores y \glspl{dispositivoiot} con recursos limitados y enlaces poco fiables. Organiza la red en una estructura jer{\'a}rquica denominada DODAG (\textit{Destination Oriented Directed Acyclic Graph}), optimizada para eficiencia energ{\'e}tica y confiabilidad en la \gls{transmision}. Se utiliza ampliamente en aplicaciones de monitorizaci{\'o}n ambiental, dom{\'o}tica y \glspl{ciudadinteligente}.},
		see={enrutamiento,protocoloenrutamiento,iot}
	}
	
	\newglossaryentry{ruido}{
		name={Ruido},
		text={ruido},
		sort={ruido},
		plural={ruidos},
		first={ruido},
		firstplural={ruidos},
		description={El \emph{ruido} es una perturbaci{\'o}n aleatoria o conjunto de se{\~n}ales no deseadas que interfieren con la \gls{transmision} o recepci{\'o}n de informaci{\'o}n en un \gls{sistemacomunicacion}. En el contexto de \gls{iot}, puede provenir de interferencias electromagn{\'e}ticas, fluctuaciones t{\'e}rmicas o equipos electr{\'o}nicos cercanos, afectando la \gls{precision} de las mediciones y la \gls{fiabilidad} de los datos. Para su tratamiento se emplean \glspl{filtro}, t{\'e}cnicas de modulaci{\'o}n robustas y algoritmos de correcci{\'o}n de errores.},
		type=main,
		symbol={📡},
		user1={Perturbaciones aleatorias en se{\~n}ales},
		user2={Causa de p{\'e}rdida de calidad en datos IoT},
		user3={Mitigable con filtros y modulaci{\'o}n robusta},
		see={filtro,precision,fiabilidad,transmision}
	}
	
	\newglossaryentry{ruta}{
		name={Ruta},
		text={ruta},
		sort={ruta},
		first={ruta},
		plural={rutas},
		firstplural={rutas},
		description={Trayectoria o camino definido para desplazamiento o transporte. En sistemas de log{\'i}stica conectada, la gesti{\'o}n de rutas integra datos de ubicaci{\'o}n, tiempos y restricciones operativas.},
		type=main,
		symbol={🛣️},
		user1={Trayectorias},
		user2={Log{\'i}stica},
		user3={Planificaci{\'o}n},
		see={vehiculo,sistemanavegacion}
	}
	
	\newglossaryentry{salud}{
		name={Salud},
		text={salud},
		sort={salud},
		description={Estado de completo bienestar f{\'i}sico, mental y social, y no solamente la ausencia de enfermedades. En el contexto de \gls{iot}, la \emph{salud} se transforma mediante dispositivos biom{\'e}dicos conectados, telemedicina, monitorizaci{\'o}n remota de pacientes, detecci{\'o}n temprana de \glspl{anomalia} y gesti{\'o}n de datos cl{\'i}nicos en \gls{tiemporeal}. Estas aplicaciones mejoran la atenci{\'o}n m{\'e}dica, la prevenci{\'o}n de riesgos y la eficiencia de los sistemas sanitarios.},
		plural={aplicaciones de salud},
		type=main,
		symbol={🏥},
		user1={Monitoreo biom{\'e}dico remoto},
		user2={Telemedicina y gesti{\'o}n cl{\'i}nica},
		user3={Prevenci{\'o}n y diagn{\'o}stico en IoT},
		see={iot,anomalia,tiemporeal}
	}
	
	\newglossaryentry{scada}{
		name={SCADA},
		text={SCADA},
		sort={scada},
		plural={sistemas SCADA},
		first={Control Supervisor y Adquisici{\'o}n de Datos (\textit{Supervisory Control and Data Acquisition}, SCADA)},
		firstplural={sistemas SCADA},
		description={Acr{\'o}nimo de \emph{Supervisory Control and Data Acquisition}, refiere a sistemas utilizados para supervisar y controlar procesos industriales a distancia. Los sistemas \emph{SCADA} representan un antecedente directo de los \glspl{sibiot}, al integrar adquisici{\'o}n de datos, control y visualizaci{\'o}n en tiempo real.},
		type=main,
		symbol={🖥️},
		user1={Supervisi{\'o}n remota},
		user2={Control industrial},
		user3={Adquisici{\'o}n de datos},
		see={automatizacionindustrial,integraciondatosoperativos}
	}
	
	\newglossaryentry{scm}{
		name={SCM},
		text={SCM},
		sort={scm},
		plural={SCM},
		first={gesti{\'o}n de la cadena de subministros o Supply Chain Management, SCM)},
		firstplural={Supply Chain Management (SCM)},
		description={Conjunto de procesos, metodolog{\'i}as y sistemas de informaci{\'o}n orientados a la planificaci{\'o}n, coordinaci{\'o}n y control integral de la cadena de suministro, desde la adquisici{\'o}n de materias primas hasta la entrega final de productos al cliente. En sistemas de informaci{\'o}n basados en IoT y \glspl{sibiot}, los sistemas \emph{SCM} se integran mediante \glspl{api} \gls{restful}, mensajer{\'i}a orientada a eventos y flujos de datos en tiempo real, permitiendo la sincronizaci{\'o}n de inventarios, la trazabilidad de productos, la predicci{\'o}n de la demanda y la coordinaci{\'o}n autom{\'a}tica con proveedores. La incorporaci{\'o}n de datos IoT en SCM mejora la visibilidad de la cadena, reduce incertidumbre operativa y habilita decisiones proactivas basadas en datos},
		type=main,
		symbol={🔗},
		see={logistica,integracion,api,restful,trazabilidad,orientadoeventos,sibiot}
	}
	
	\newglossaryentry{scrum}{
		name={Scrum},
		text={Scrum},
		sort={scrum},
		description={\emph{Scrum} es un marco de trabajo \gls{metodoagil} utilizado para la gesti{\'o}n y desarrollo de productos complejos, especialmente software. Se basa en iteraciones llamadas \textit{sprints}, en roles definidos (\textit{Product Owner}, \textit{Scrum Master} y equipo de desarrollo) y en eventos estructurados como reuniones diarias, revisiones y retrospectivas. En proyectos \gls{iot}, Scrum facilita la entrega incremental de soluciones funcionales, promoviendo la adaptabilidad, la colaboraci{\'o}n continua y la integraci{\'o}n temprana de componentes de hardware y software.},
		plural={implementaciones de Scrum},
		type=main,
		symbol={📋},
		user1={Marco de trabajo {\'a}gil basado en sprints},
		user2={Promueve colaboraci{\'o}n e iteraci{\'o}n},
		user3={Soporte a IoT en desarrollo incremental},
		see={metodoagil,kanban,metodologiaagil}
	}
	
	\newglossaryentry{scrummaster}{
		name={Scrum Master},
		text={Scrum Master},
		sort={scrum master},
		plural={Scrum Masters},
		first={Scrum Master},
		firstplural={Scrum Masters},
		description={Rol del Scrum Team responsable de establecer y promover la correcta aplicaci{\'o}n del marco de trabajo Scrum. El Scrum Master act{\'u}a como l{\'i}der servicial, facilitando eventos, eliminando impedimentos y ayudando al equipo y a la organizaci{\'o}n a comprender y adoptar pr{\'a}cticas {\'a}giles. Asimismo, fomenta la mejora continua, la autoorganizaci{\'o}n de los Developers y la transparencia en los artefactos y compromisos del equipo.},
		type=main,
		symbol={🧭},
		user1={Facilitador de Scrum},
		user2={L{\'i}der servicial},
		user3={Promotor de mejora continua},
		see={developer,productowner,sprint,agil,scrum}
	}
	
	\newglossaryentry{seguridad}{
		name={Seguridad},
		text={seguridad},
		sort={seguridad},
		plural={seguridades},
		first={seguridad},
		firstplural={seguridades},
		description={Conjunto de principios, pr{\'a}cticas y controles destinados a proteger sistemas, redes y datos frente a amenazas y riesgos, garantizando propiedades como \emph{confidencialidad} (acceso s{\'o}lo para quienes est{\'a}n autorizados), \gls{integridad} (no alteraci{\'o}n no autorizada) y \gls{disponibilidad} (servicios operativos cuando se requieren). En \gls{iot} y \glspl{sibiot} incluye gesti{\'o}n de identidades y accesos, segmentaci{\'o}n de redes, cifrado extremo a extremo, actualizaciones seguras, monitorizaci{\'o}n continua y respuesta a incidentes.},
		type=main,
		symbol={},
		user1={Ciberseguridad y gesti{\'o}n de riesgos},
		user2={Protecci{\'o}n de datos y servicios},
		user3={Controles preventivos y reactivos},
		see={privacidad,integridad,disponibilidad,buenaspracticas}
	}
	
	\newglossaryentry{seguridadvial}{
		name={Seguridad vial},
		text={seguridad vial},
		sort={seguridad vial},
		first={seguridad vial},
		description={Conjunto de medidas t{\'e}cnicas, normativas y conductuales orientadas a prevenir siniestros de tr{\'a}nsito y mitigar sus consecuencias, mediante control, se{\~n}alizaci{\'o}n y asistencia al conductor.},
		type=main,
		symbol={🚦},
		user1={Prevenci{\'o}n},
		user2={Normativa},
		user3={Asistencia},
		see={vehiculo,monitorizacion}
	}
	
	\newglossaryentry{senal}
	{
		name={Se{\~n}al},
		text={se{\~n}al},
		sort={senal},
		plural={se{\~n}ales},
		first={se{\~n}al},
		firstplural={se{\~n}ales},
		symbol={📶},
		see={[v{\'e}ase tambi{\'e}n]{senalelectrica}},
		description={Magnitud f{\'i}sica o simb{\'o}lica que porta informaci{\'o}n sobre un fen{\'o}meno o variable, susceptible de medici{\'o}n, procesamiento y transmisi{\'o}n. En \glspl{sibiot}, una se{\~n}al puede ser continua o discreta, y codifica estados o cambios detectados por \glspl{sensor} para su an{\'a}lisis por \glspl{sistemaelectronico} y \glspl{algoritmo}.}
	}
	
	\newglossaryentry{senalelectrica}
	{
		name={Se{\~n}al el{\'e}ctrica},
		text={se{\~n}al el{\'e}ctrica},
		sort={senalelectrica},
		plural={se{\~n}ales el{\'e}ctricas},
		first={se{\~n}al el{\'e}ctrica},
		firstplural={se{\~n}ales el{\'e}ctricas},
		symbol={∿},
		see={[relacionada]{sensor}},
		description={Variaci{\'o}n temporal o espacial de una magnitud el{\'e}ctrica (tensi{\'o}n, corriente, potencia o carga) que representa \gls{informacion} proveniente de un \gls{sensor} o generada por un \gls{sistemaelectronico}. En \glspl{sibiot}, las se{\~n}ales el{\'e}ctricas posibilitan la adquisici{\'o}n, acondicionamiento y \gls{conectividad} de \glspl{dato} hacia m{\'o}dulos de \gls{procesamiento} y toma de decisiones.}
	}
	
	% ===================== Sensibilidad =====================
	\newglossaryentry{sensibilidad}{
		name={Sensibilidad},
		text={sensibilidad},
		sort={sensibilidad},
		plural={sensibilidades},
		first={sensibilidad},
		firstplural={sensibilidades},
		description={La \emph{sensibilidad} es la magnitud m{\'i}nima en la se{\~n}al de entrada requerida para producir una determinada salida, bajo un criterio especificado como relaci{\'o}n se{\~n}al/ruido. En \glspl{sensor}, la sensibilidad determina su capacidad para detectar cambios peque{\~n}os en el entorno, siendo esencial para aplicaciones de \gls{iot} que requieren \gls{precision} en la medici{\'o}n.},
		type=main,
		symbol={🎚️},
		user1={Magnitud m{\'i}nima detectable},
		user2={Factor clave en sensores IoT},
		user3={Relaci{\'o}n con precisi{\'o}n y ruido},
		see={sensor,precision,ruido}
	}
	
	% ===================== Sensor =====================
	\newglossaryentry{sensor}{
		name={Sensor},
		text={sensor},
		sort={sensor},
		plural={sensores},
		first={sensor},
		firstplural={sensores},
		description={Un \emph{sensor} es un dispositivo electr{\'o}nico capaz de detectar acciones o est{\'i}mulos externos, transformando magnitudes f{\'i}sicas o qu{\'i}micas en magnitudes el{\'e}ctricas. En \gls{iot}, los sensores son la base para la \gls{deteccion}, la \gls{monitorizacion} y la recolecci{\'o}n de datos en \glspl{ecosistemaiot}.},
		type=main,
		symbol={📟},
		user1={Dispositivo de detecci{\'o}n de magnitudes},
		user2={Conversi{\'o}n f{\'i}sica/qu{\'i}mica a el{\'e}ctrica},
		user3={Base de la adquisici{\'o}n de datos IoT},
		see={sensibilidad,actuador,ecosistemaiot}
	}
	
	% ===================== Sensor de fuerza =====================
	\newglossaryentry{sensordefuerza}{
		name={Sensor de fuerza},
		text={sensor de fuerza},
		sort={sensor de fuerza},
		plural={sensores de fuerza},
		first={sensor de fuerza},
		firstplural={sensores de fuerza},
		description={Un \emph{sensor de fuerza} detecta y mide la magnitud de una fuerza mec{\'a}nica aplicada, convirti{\'e}ndola en una se{\~n}al el{\'e}ctrica proporcional. Puede basarse en galgas extensiom{\'e}tricas, piezoresistencias, efectos piezoel{\'e}ctricos o capacitivos. En \glspl{sibiot}, se emplea para monitorizaci{\'o}n de presi{\'o}n, \gls{control} de actuadores, rob{\'o}tica, salud y manufactura inteligente.},
		type=main,
		symbol={⚖️},
		user1={Medici{\'o}n de fuerzas mec{\'a}nicas},
		user2={Conversi{\'o}n fuerza–se{\~n}al el{\'e}ctrica},
		user3={Uso en IoT: rob{\'o}tica y control industrial},
		see={sensor,control,manufactura,sibiot}
	}
	
	% ===================== Sensor de peso =====================
	\newglossaryentry{sensordepeso}{
		name={Sensor de peso},
		text={sensor de peso},
		sort={sensor de peso},
		plural={sensores de peso},
		first={sensor de peso},
		firstplural={sensores de peso},
		description={Un \emph{sensor de peso}, tambi{\'e}n conocido como celda de carga, mide la fuerza ejercida por un objeto debido a su masa y la convierte en una se{\~n}al el{\'e}ctrica. Se aplica en balanzas digitales, \gls{control} de inventarios y sistemas industriales. En \gls{iot}, permite monitorizaci{\'o}n remota, an{\'a}lisis de patrones de carga y detecci{\'o}n de sobrepeso en \gls{tiemporeal}.},
		type=main,
		symbol={🏋️},
		user1={Medici{\'o}n de peso mediante celdas de carga},
		user2={Aplicaci{\'o}n en balanzas e inventarios},
		user3={Soporte IoT: monitorizaci{\'o}n remota en tiempo real},
		see={celdacarga,sensor,control}
	}
	
	% ===================== Sensor de presi{\'o}n =====================
	\newglossaryentry{sensordepresion}{
		name={Sensor de presi{\'o}n},
		text={sensor de presi{\'o}n},
		sort={sensor de presion},
		plural={sensores de presi{\'o}n},
		first={sensor de presi{\'o}n},
		firstplural={sensores de presi{\'o}n},
		description={Un \emph{sensor de presi{\'o}n} mide la presi{\'o}n de gases o l{\'i}quidos y la convierte en una se{\~n}al el{\'e}ctrica proporcional. Puede usar principios piezorresistivos, capacitivos, inductivos o piezoel{\'e}ctricos. En \gls{iot}, se emplea en monitorizaci{\'o}n ambiental, \gls{control} de procesos industriales, sistemas hidr{\'a}ulicos y salud (p.\,ej., medici{\'o}n de presi{\'o}n arterial).},
		type=main,
		symbol={💨},
		user1={Medici{\'o}n de presi{\'o}n de gases o l{\'i}quidos},
		user2={Principios: piezorresistivos, capacitivos, inductivos},
		user3={Aplicaciones IoT: industria, salud, ambiente},
		see={sensor,control,precision}
	}
	
	% ===================== Sensor de visi{\'o}n =====================
	\newglossaryentry{sensorvision}{
		name={Sensor de visi{\'o}n},
		text={sensor de visi{\'o}n},
		sort={sensor de vision},
		plural={sensores de visi{\'o}n},
		first={sensor de visi{\'o}n},
		firstplural={sensores de visi{\'o}n},
		description={Un \emph{sensor de visi{\'o}n} captura informaci{\'o}n visual mediante im{\'a}genes o v{\'i}deo, que son procesados con algoritmos de visi{\'o}n artificial para reconocer patrones o tomar decisiones. En sistemas \gls{iot}, permiten aplicaciones como detecci{\'o}n de objetos, \gls{seguridad}, \gls{reconocimientofacial}, inspecci{\'o}n industrial y control de calidad.},
		type=main,
		symbol={📷},
		user1={Captura visual del entorno},
		user2={Soporte a visi{\'o}n artificial},
		user3={Aplicaciones IoT: seguridad, inspecci{\'o}n, calidad},
		see={sensor,seguridad,reconocimientofacial}
	}
	
	% ===================== Sensor infrarrojo =====================
	\newglossaryentry{sensorinfrarrojo}{
		name={Sensor infrarrojo},
		text={sensor infrarrojo},
		sort={sensor infrarrojo},
		plural={sensores infrarrojos},
		first={sensor infrarrojo},
		firstplural={sensores infrarrojos},
		description={Un \emph{sensor infrarrojo} detecta radiaci{\'o}n en la regi{\'o}n infrarroja del espectro, normalmente asociada a la temperatura de los objetos. Se emplea para detecci{\'o}n de movimiento, medici{\'o}n de temperatura sin contacto y visi{\'o}n t{\'e}rmica. En \gls{iot}, se usa en \gls{seguridad}, automatizaci{\'o}n dom{\'o}tica, agricultura de \gls{precision} y control industrial.},
		type=main,
		symbol={🌡️},
		user1={Detecci{\'o}n de radiaci{\'o}n infrarroja},
		user2={Aplicaci{\'o}n en movimiento y temperatura},
		user3={Uso IoT: dom{\'o}tica, agricultura y seguridad},
		see={sensor,precision,seguridad}
	}
	
	% ===================== GlassFish =====================
	\newglossaryentry{serGlaFi}{
		name={GlassFish},
		text={GlassFish},
		sort={glassfish},
		plural={GlassFish},
		first={GlassFish},
		firstplural={GlassFish},
		description={\emph{GlassFish} es una implementaci{\'o}n de referencia de la plataforma Java EE, de c{\'o}digo abierto. Funciona como servidor de aplicaciones para ejecutar aplicaciones empresariales Java, brindando soporte a servicios web, APIs, transacciones y escalabilidad. En proyectos \gls{iot}, se utiliza para gestionar servicios distribuidos, APIs REST y la integraci{\'o}n con bases de datos y sistemas externos.},
		type=main,
		symbol={☕},
		user1={Servidor de aplicaciones Java EE},
		user2={Plataforma de c{\'o}digo abierto},
		user3={Uso en servicios distribuidos e IoT},
		see={java,servicioweb,aplicacionempresarial}
	}
	
	\newglossaryentry{servicio}{
		name={Servicio},
		text={servicio},
		sort={servicio},
		plural={servicios},
		first={servicio},
		firstplural={servicios},
		description={Funcionalidad, recurso o conjunto de operaciones ofrecidas por un \gls{sistema}, componente o aplicaci{\'o}n, que puede ser consumida por usuarios, otros sistemas o \glspl{aplicacion} a trav{\'e}s de interfaces definidas. En la ingenier{\'i}a de software y los entornos distribuidos, un \emph{servicio} se caracteriza por su independencia, su reutilizaci{\'o}n y la posibilidad de integrarse con otros servicios para componer soluciones m{\'a}s complejas. En el contexto de los \glspl{sibiot}, los servicios abarcan desde el procesamiento de datos de \glspl{sensor}, la monitorizaci{\'o}n en \gls{tiemporeal} y el \gls{control} de \glspl{actuador}, hasta servicios en la nube como anal{\'i}tica avanzada, \glspl{modelopredictivo} o integraci{\'o}n con ecosistemas externos, asegurando \gls{interoperabilidad}, \gls{seguridad} y \gls{escalabilidad}.},
		type=main,
		symbol={🔧},
		user1={Unidad funcional ofrecida a usuarios o sistemas},
		user2={Componente reutilizable e integrable},
		user3={Base de arquitecturas distribuidas y SIbIoT},
		see={aplicacion,sibiot,interoperabilidad,msoa,api}
	}
	
	\newglossaryentry{serviciodistribuido}{
		name={Servicio distribuido},
		text={servicio distribuido},
		sort={servicio distribuido},
		plural={servicios distribuidos},
		first={servicio distribuido},
		firstplural={servicios distribuidos},
		description={Funcionalidad o conjunto de operaciones implementadas en una \gls{arquitecturadistribuida}, de manera que su ejecuci{\'o}n est{\'a} repartida en varios nodos interconectados por una red. Un \emph{servicio distribuido} se caracteriza por su capacidad de escalado, tolerancia a fallos y disponibilidad, al no depender de un {\'u}nico punto de falla. En el contexto de los \glspl{sibiot}, los servicios distribuidos permiten procesar datos de \glspl{sensor} en tiempo real, coordinar \glspl{actuador}, integrar plataformas en la nube y soportar aplicaciones cr{\'i}ticas en dominios como la salud, la industria, las ciudades inteligentes o la log{\'i}stica. Su dise{\~n}o se apoya en principios de \gls{msoa}, \gls{api} y est{\'a}ndares de \gls{interoperabilidad}.},
		type=main,
		symbol={🧩},
		user1={Operaciones ejecutadas en m{\'u}ltiples nodos},
		user2={Escalabilidad y tolerancia a fallos},
		user3={Soporte a aplicaciones distribuidas en SIbIoT},
		see={arquitecturadistribuida,msoa,api,sibiot,microservicio}
	}
	
	\newglossaryentry{servicionube}{
		name={Servicio en la nube},
		text={servicio en la nube},
		sort={servicionube},
		first={servicio en la nube},
		plural={servicios en la nube},
		firstplural={servicios en la nube},
		description={Un \emph{servicio en la nube} es una soluci{\'o}n inform{\'a}tica ofrecida a trav{\'e}s de \gls{cloudcomputing}, que permite a los usuarios acceder bajo demanda a recursos como \gls{almacenamiento}, \gls{procesamiento} de datos, bases de datos, \glspl{servicioweb}, aplicaciones o infraestructura completa, sin necesidad de gestionar directamente hardware o software local. Estos servicios se clasifican com{\'u}nmente en \textbf{SaaS} (Software como Servicio), \textbf{PaaS} (Plataforma como Servicio) e \textbf{IaaS} (Infraestructura como Servicio). En el contexto de \gls{iot}, los servicios en la nube son esenciales para consolidar datos de \glspl{dispositivoiot}, habilitar \glspl{dashboard} de monitorizaci{\'o}n, ejecutar \glspl{modelopredictivo} y facilitar la \gls{escalabilidad} y \gls{interoperabilidad} de soluciones distribuidas.},
		type=main,
		symbol={☁️},
		user1={Acceso bajo demanda a recursos inform{\'a}ticos},
		user2={Modelos SaaS, PaaS, IaaS},
		user3={Soporte a IoT y big data},
		see={nube,cloudcomputing,servicioweb,sibiot,escalabilidad}
	}
	
	\newglossaryentry{servicioweb}{
		name={Servicio web},
		text={servicio web},
		sort={servicioweb},
		plural={servicios web},
		first={servicio web},
		firstplural={servicios web},
		description={Un \emph{servicio web} es una aplicaci{\'o}n o componente de software que utiliza protocolos estandarizados como \gls{http}, \gls{https} o SOAP para intercambiar datos entre sistemas heterog{\'e}neos a trav{\'e}s de Internet o de una \gls{red}. Los servicios web suelen exponer interfaces en forma de \glspl{api}, permitiendo la interoperabilidad entre aplicaciones desarrolladas en diferentes lenguajes o plataformas. En el contexto de \gls{iot}, los servicios web son fundamentales para integrar \glspl{dispositivoiot} con servidores, plataformas en la \gls{nube} y aplicaciones m{\'o}viles, habilitando funcionalidades como monitorizaci{\'o}n remota, control inteligente y an{\'a}lisis en \gls{tiemporeal}.},
		type=main,
		symbol={🌐},
		user1={Interoperabilidad mediante APIs},
		user2={Uso de HTTP/HTTPS y SOAP},
		user3={Integraci{\'o}n con IoT y servicios en la nube},
		see={api,aplicacionweb,nube,iot,tiemporeal}
	}
	
	\newglossaryentry{servidor}
	{
		name={Servidor},
		text={servidor},
		sort={servidor},
		plural={servidores},
		first={servidor},
		firstplural={servidores},
		description={Un servidor es un programa o sistema que se ejecuta generalmente en una m{\'a}quina con mayores prestaciones que las dem{\'a}s de la red (clientes). Su funci{\'o}n principal es satisfacer las demandas de los sistemas cliente, suministrando datos o informaci{\'o}n solicitada. En el contexto de los \glspl{sibiot}, los servidores permiten gestionar la interacci{\'o}n entre \glspl{dispositivo}, \glspl{usuario} y \glspl{serviciodistribuido}.},
		type=\glsdefaulttype,
		symbol={🖥},
		user1={Host},
		user2={Sistema servidor},
		user3={Proveedor de datos},
		see={servidorweb}
	}
	
	\newglossaryentry{servidorweb}
	{
		name={Servidor web},
		text={servidor web},
		sort={servidor web},
		plural={servidores web},
		first={servidor web},
		firstplural={servidores web},
		description={Un servidor web es un sistema inform{\'a}tico que almacena, procesa y entrega p{\'a}ginas o servicios web mediante los protocolos HTTP o HTTPS. Puede alojar sitios, aplicaciones, APIs o interfaces de monitorizaci{\'o}n. En el contexto de los \glspl{sibiot}, los servidores web permiten visualizar datos recolectados por \glspl{dispositivo}, gestionar configuraciones remotas, habilitar \textit{dashboards} interactivos y facilitar la comunicaci{\'o}n entre usuarios y \glspl{dispositivoconectado}.},
		type=\glsdefaulttype,
		symbol={🌐},
		user1={Host web},
		user2={Servidor HTTP},
		user3={Servidor HTTPS},
		see={servidor}
	}
	
	\newglossaryentry{serverless}
	{
		name={Computaci{\'o}n sin servidor},
		text={computaci{\'o}n sin servidor},
		sort={computacion sin servidor},
		plural={plataformas sin servidor},
		first={computaci{\'o}n sin servidor},
		firstplural={plataformas sin servidor},
		description={Modelo de ejecuci{\'o}n en el que funciones o servicios se ejecutan bajo demanda y por eventos, con gesti{\'o}n de infraestructura asumida por la plataforma, com{\'u}n en tareas de ingesta, validaci{\'o}n y automatizaci{\'o}n en nube.},
		type=\glsdefaulttype,
		symbol={⚡},
		user1={Ejecuci{\'o}n por eventos},
		user2={Menos operaci{\'o}n},
		user3={Servicios gestionados},
		see={cloudcomputing}
	}
	
	\newglossaryentry{servomotor}
	{
		name={Servomotor},
		text={servomotor},
		sort={servomotor},
		plural={servomotores},
		first={servomotor},
		firstplural={servomotores},
		description={Un servomotor es un dispositivo electromec{\'a}nico que permite controlar de forma precisa la posici{\'o}n angular, velocidad y aceleraci{\'o}n de un eje o brazo mec{\'a}nico. Funciona mediante un sistema de retroalimentaci{\'o}n (\textit{feedback}) que ajusta el movimiento seg{\'u}n la diferencia entre la posici{\'o}n deseada y la real. Es ampliamente utilizado en rob{\'o}tica, automatizaci{\'o}n y sistemas de control embebido.},
		type=\glsdefaulttype,
		symbol={⚙},
		user1={Motor de control},
		user2={Motor con retroalimentaci{\'o}n},
		user3={Actuador preciso},
		see={actuador}
	}
	
	\newglossaryentry{setpoint}{
		name={Setpoint},
		text={setpoint},
		short={setpoint},
		long={Valor objetivo de una variable de control que el sistema intenta mantener mediante retroalimentaci{\'o}n},
		sort={setpoint},
		plural={setpoints},
		first={setpoint},
		firstplural={setpoints},
		description={Un \emph{setpoint} es el valor objetivo asignado a una variable controlada (p.\,ej., temperatura, presi{\'o}n, velocidad) que un sistema intenta mantener o alcanzar mediante un lazo de retroalimentaci{\'o}n. En \glspl{cps}, el setpoint se define dentro de restricciones de \gls{temporizacion}, seguridad y estabilidad del proceso; adem{\'a}s, puede derivarse de pol{\'\i}ticas operativas (horarios, prioridades, umbrales) y debe quedar registrado cuando se ajusta, para soportar trazabilidad y auditor{\'\i}a en contextos organizacionales.},
		type=main,
		symbol={🎯},
		user1={Valor objetivo de control},
		user2={Pol{\'\i}ticas y restricciones},
		user3={Registro y trazabilidad del ajuste},
		see={control,temporizacion,trazabilidadoperacional}
	}
	
	\newglossaryentry{sgbd}{
		name={DBMS},
		text={DBMS},
		sort={dbms},
		first={Sistema Gestor de Base de Datos (\textit{Database Management System}, DBMS)},
		firstplural={DBMSs},
		plural={DBMSs},
		description={Sistema Gestor de Base de Datos. Software que define, crea, consulta, actualiza y administra una \gls{basedatos}, proporcionando lenguajes de consulta, control de concurrencia, seguridad y mecanismos de \gls{recuperacion}.},
		type=main,
		symbol={🗄️},
		user1={Administraci{\'o}n de datos},
		user2={Consultas},
		user3={Seguridad},
		see={basedatos,tsql,recuperacion}
	}
	
	\newglossaryentry{sibiot}{
		name={SIbIoT},
		text={SIbIoT},
		sort={¡sibiot},
		first={sistema de informaci{\'o}n basado en IoT (SIbIoT)},
		firstplural={sistemas de informaci{\'o}n basados en IoT},
		plural={SIbIoT},
		description={Un Sistema de informaci{\'o}n basado en IoT (SIbIoT) es un sistema de informaci{\'o}n empresarial que, adem{\'a}s de gestionar datos transaccionales tradicionales, integra datos provenientes de sensores f{\'i}sicos y mecanismos de actuaci{\'o}n, permitiendo la monitorizaci{\'o}n, el control y la toma de decisiones sobre procesos organizacionales en tiempo real o casi real.},
		type=main,
		symbol={🌐},
		see={sistema,sistemainformacion,sistemaconectado,sistemadistribuido}
	}
	\newglossaryentry{simulacion}{
		name={Simulaci{\'o}n},
		text={simulaci{\'o}n},
		sort={simulacion},
		first={simulaci{\'o}n},
		description={Representaci{\'o}n y ejecuci{\'o}n de un modelo de sistema para estudiar su comportamiento bajo escenarios hipot{\'e}ticos, sin intervenir directamente el entorno real.},
		type=main,
		symbol={🧪},
		user1={Modelado},
		user2={Escenarios},
		user3={An{\'a}lisis},
		see={realidad,representacion}
	}
	
	\newglossaryentry{sincronizacion}{
		name={Sincronizaci{\'o}n},
		text={sincronizaci{\'o}n},
		sort={sincronizacion},
		plural={sincronizaciones},
		first={sincronizaci{\'o}n},
		firstplural={sincronizaciones},
		description={Proceso mediante el cual dos o m{\'a}s componentes, procesos o \glspl{nodo} de un \gls{sistema} coordinan su estado, informaci{\'o}n o ejecuci{\'o}n para mantener coherencia y funcionamiento conjunto. La \emph{sincronizaci{\'o}n} puede referirse a relojes (tiempo), datos (consistencia en \glspl{basedatosdistribuida}), se{\~n}ales (comunicaci{\'o}n) o procesos paralelos. En el contexto de los \glspl{sibiot}, la sincronizaci{\'o}n es esencial para asegurar que las mediciones de \glspl{sensor}, las acciones de \glspl{actuador} y los servicios distribuidos operen de forma coordinada, garantizando \gls{fiabilidad}, \gls{precision} temporal y \gls{resiliencia} frente a fallos o retardos. Su implementaci{\'o}n puede lograrse mediante protocolos de tiempo (NTP, PTP), \glspl{algoritmo} de consenso o mecanismos de \gls{replicacion}.},
		type=main,
		symbol={⏲️},
		user1={Coordinaci{\'o}n de estados y procesos},
		user2={Garant{\'i}a de consistencia y coherencia},
		user3={Base de operaci{\'o}n en sistemas distribuidos e IoT},
		see={replicacion,nodo,basedatosdistribuida,fiabilidad,resiliencia,sibiot}
	}
	
	\newglossaryentry{sincronizada}{
		name={Sincronizada},
		text={sincronizada},
		sort={sincronizada},
		plural={sincronizadas},
		first={sincronizada},
		firstplural={sincronizadas},
		description={Condici{\'o}n en la que dos o m{\'a}s procesos coordinan su ejecuci{\'o}n en el tiempo, respetando reglas de orden, dependencia o eventos comunes. En sistemas embebidos, la sincronizaci{\'o}n es esencial para evitar inconsistencias y garantizar un comportamiento determinista.},
		type=main,
		symbol={🔗},
		user1={Coordinaci{\'o}n temporal},
		user2={Orden de ejecuci{\'o}n},
		user3={Consistencia},
		see={tareasincronizada,concurrente}
	}
	
		\newglossaryentry{sincrona}{
		name={S{\'i}ncrona},
		text={s{\'i}ncrona},
		sort={sincrona},
		plural={s{\'i}ncronas},
		first={s{\'i}ncrona},
		firstplural={s{\'i}ncronas},
		description={Femenino de \gls{sincrono}. Se aplica a se{\~n}ales, procesos o comunicaciones que est{\'a}n coordinados en el tiempo, de modo que sus eventos significativos ocurren de manera simult{\'a}nea o en correspondencia temporal estricta. En sistemas \gls{iot}, la comunicaci{\'o}n s{\'i}ncrona asegura que la transmisi{\'o}n y recepci{\'o}n de datos est{\'e}n alineadas en intervalos definidos, garantizando predictibilidad y control en la interacci{\'o}n entre dispositivos.},
		type=main,
		symbol={⏳},
		user1={Comunicaci{\'o}n o proceso alineado en el tiempo},
		user2={Propiedad de sincronizaci{\'o}n estricta},
		user3={Aplicable a se{\~n}ales, procesos y comunicaciones},
		see={sincrono,asincrono}
	}
	
	\newglossaryentry{sincrono}{
		name={S{\'i}ncrono},
		text={s{\'i}ncrono},
		sort={sincrono},
		plural={s{\'i}ncronos},
		description={Adjetivo que describe procesos, se{\~n}ales o sistemas coordinados en el tiempo. En computaci{\'o}n, un sistema \emph{s{\'i}ncrono} ejecuta operaciones de forma secuencial y bloqueante. En comunicaciones, implica que emisor y receptor comparten un mismo reloj de referencia.},
		type=main,
		symbol={⏱️},
		user1={Procesos coordinados en el tiempo},
		user2={Operaciones secuenciales y bloqueantes},
		user3={Comunicaci{\'o}n con reloj compartido},
		see={asincrono}
	}
	
	\newglossaryentry{sistema}{
		name={Sistema},
		text={sistema},
		plural={sistemas},
		sort={sistema},
		first={sistema},
		firstplural={sistemas},
		description={Conjunto organizado de elementos interrelacionados (personas, \glspl{dispositivo}, \glspl{proceso}, datos o componentes de software y hardware) que interact{\'u}an para alcanzar un objetivo com{\'u}n. Un \emph{sistema} puede ser f{\'i}sico, l{\'o}gico o socio-t{\'e}cnico, y se caracteriza por poseer entradas, procesos, salidas y retroalimentaci{\'o}n. En el contexto de \gls{iot} y \glspl{sibiot}, los sistemas integran \glspl{dispositivointeligente}, plataformas de comunicaci{\'o}n, \gls{almacenamiento} en la \gls{nube}, servicios anal{\'i}ticos y aplicaciones de usuario para lograr automatizaci{\'o}n, \gls{monitorizacion} y \gls{tomadecisiones} inteligentes en \gls{tiemporeal}.},
		type=main,
		symbol={⚙️},
		user1={Conjunto organizado de elementos},
		user2={Interacci{\'o}n con entradas, procesos y salidas},
		user3={Base de automatizaci{\'o}n y control en IoT},
		see={sibiot,dispositivointeligente,arquitecturadistribuida}
	}
		
	\newglossaryentry{sistemacifrado}{
		name={Sistema de cifrado},
		text={sistema de cifrado},
		sort={sistemacifrado},
		plural={sistemas de cifrado},
		first={sistema de cifrado},
		firstplural={sistemas de cifrado},
		description={Conjunto de m{\'e}todos criptogr{\'a}ficos que permiten codificar datos para garantizar su confidencialidad, integridad y autenticidad. Los sistemas de cifrado pueden ser de clave sim{\'e}trica (como AES o DES), donde la misma clave se usa para cifrar y descifrar, o de clave asim{\'e}trica (como RSA o ECC), que emplean un par de claves p{\'u}blica y privada. En \glspl{sibiot}, los sistemas de cifrado son esenciales para proteger la comunicaci{\'o}n entre sensores, actuadores y servicios en la nube, evitando accesos no autorizados o manipulaci{\'o}n de datos.},
		type=main,
		symbol={🔐},
		user1={Criptograf{\'i}a},
		user2={Protecci{\'o}n de datos},
		user3={Seguridad de la informaci{\'o}n},
		see={encriptacion,seguridad,privacidad}
	}
	
	\newglossaryentry{sistemacomunicacion}
	{
		name={Sistema de comunicaci{\'o}n},
		text={sistema de comunicaci{\'o}n},
		sort={sistema de comunicacion},
		plural={sistemas de comunicaci{\'o}n},
		first={sistema de comunicaci{\'o}n},
		firstplural={sistemas de comunicaci{\'o}n},
		description={Un sistema de comunicaci{\'o}n es un conjunto organizado de medios, dispositivos y protocolos que permiten la transmisi{\'o}n de informaci{\'o}n entre un emisor y uno o varios receptores. En el contexto de los \glspl{sibiot}, los sistemas de comunicaci{\'o}n integran \glspl{sensor}, \glspl{actuador}, redes inal{\'a}mbricas y plataformas digitales, garantizando la conectividad, interoperabilidad y confiabilidad de los datos intercambiados.},
		type=\glsdefaulttype,
		symbol={🌐},
		user1={Red de comunicaci{\'o}n},
		user2={Infraestructura de transmisi{\'o}n},
		user3={Canal de comunicaci{\'o}n},
		see={conectividad}
	}
	
	\newglossaryentry{sistemaconectado}{
		name={sistema conectado},
		text={sistema conectado},
		sort={sistema conectado},
		plural={sistemas conectados},
		first={sistema conectado},
		firstplural={sistemas conectados},
		description={Conjunto de componentes f{\'i}sicos o l{\'o}gicos que mantienen comunicaci{\'o}n continua o peri{\'o}dica mediante redes cableadas o inal{\'a}mbricas, permitiendo el intercambio de datos en tiempo real y la coordinaci{\'o}n de procesos. Un \emph{sistema conectado} se caracteriza por su capacidad de interoperar con otros dispositivos, acceder a servicios remotos, enviar telemetr{\'i}a y recibir instrucciones centralizadas. En el contexto de los \glspl{sibiot}, constituye la base operativa para habilitar monitorizaci{\'o}n distribuida, automatizaci{\'o}n, sincronizaci{\'o}n de estados y toma de decisiones basada en datos provenientes de m{\'u}ltiples nodos heterog{\'e}neos},
		type=main,
		symbol={🔗},
		user1={Intercambio continuo de informaci{\'o}n},
		user2={Conectividad mediante redes cableadas o inal{\'a}mbricas},
		user3={Integraci{\'o}n funcional dentro de arquitecturas distribuidas},
		see={dispositivoconectado,redcomunicacion,sibiot}
	}
	
	\newglossaryentry{sistemadistribuido}{
		name={Sistema distribuido},
		text={sistema distribuido},
		sort={sistema distribuido},
		plural={sistemas distribuidos},
		first={sistema distribuido},
		firstplural={sistemas distribuidos},
		description={Conjunto de \glspl{computador} o \glspl{nodo} interconectados mediante una \gls{red}, que cooperan para ofrecer un servicio o ejecutar tareas como si se tratara de un sistema {\'u}nico y coherente. En un \emph{sistema distribuido}, los recursos de hardware y \gls{software} se coordinan de manera descentralizada, compartiendo datos y procesos para alcanzar objetivos comunes. Sus ventajas incluyen \gls{escalabilidad}, \gls{toleranciafallos}, disponibilidad y rendimiento mejorado, aunque tambi{\'e}n plantean desaf{\'i}os de \gls{sincronizacion}, \gls{consistencia} y \gls{seguridad}.},
		type=main,
		symbol={🌐},
		user1={Colaboraci{\'o}n de m{\'u}ltiples nodos},
		user2={Recursos distribuidos y coordinados},
		user3={Escalabilidad y tolerancia a fallos},
		see={arquitecturadistribuida,red,nodo,consistencia,toleranciafallos}
	}
	
	\newglossaryentry{sistemaelectronico}{
		name={Sistema electr{\'o}nico},
		text={sistema electr{\'o}nico},
		sort={sistema electronico},
		plural={sistemas electr{\'o}nicos},
		first={sistema electr{\'o}nico},
		firstplural={sistemas electr{\'o}nicos},
		description={Conjunto de dispositivos y circuitos electr{\'o}nicos interconectados que procesan, transmiten o transforman se{\~n}ales el{\'e}ctricas con el fin de realizar una funci{\'o}n determinada. Un \emph{sistema electr{\'o}nico} puede incluir \glspl{sensor}, \glspl{actuador}, \glspl{microcontrolador}, memorias y componentes de comunicaci{\'o}n. En el contexto de los \glspl{sibiot}, los sistemas electr{\'o}nicos son esenciales para la captura de datos, el control de dispositivos y la ejecuci{\'o}n de procesos autom{\'a}ticos en \gls{tiemporeal}.},
		type=main,
		symbol={⚙️},
		user1={Circuito integrado de funciones},
		user2={Arquitectura electr{\'o}nica},
		user3={Sistema de hardware electr{\'o}nico},
		see={sistemainformatico,placa,placaprocesamiento}
	}
	
	% Entrada principal
	\newglossaryentry{sistemaembebido}{
		name={Sistema embebido},
		text={sistema embebido},
		sort={sistema embebido},
		plural={sistemas embebidos},
		first={sistema embebido},
		firstplural={sistemas embebidos},
		description={\gls{sistema} computacional especializado que forma parte de un \gls{dispositivo} mayor y est{\'a} dise{\~n}ado para realizar funciones espec{\'i}ficas con recursos limitados. Integra \gls{hardware} y \gls{software} embebidos, opera frecuentemente en \gls{tiemporeal} y se encuentra com{\'u}nmente en electrodom{\'e}sticos, autom{\'o}viles, dispositivos m{\'e}dicos, \glspl{sensor} \gls{iot} y sistemas industriales. Tambi{\'e}n denominado \emph{sistema incrustado}.},
		type=main,
		symbol={🔧},
		user1={Computaci{\'o}n especializada integrada en un dispositivo},
		user2={Operaci{\'o}n en tiempo real y recursos limitados},
		user3={Aplicaciones en IoT, automoci{\'o}n y m{\'e}dica},
		see={iot,hardware,software,tiemporeal,microcontrolador,placacontroladora,placadesarrollo}
	}
	
	\newglossaryentry{sistemaempresarial}{
		name={Sistema empresarial},
		text={sistema empresarial},
		plural={sistemas empresariales},
		first={sistema empresarial},
		firstplural={sistemas empresariales},
		see={erp,crm,sistemainformacion,sibiot},
		sort={sistemaempresarial},
		description={Plataforma de informaci{\'o}n utilizada para soportar procesos organizacionales como gesti{\'o}n comercial, log{\'i}stica, financiera o de producci{\'o}n. En SIbIoT, los sistemas empresariales consumen resultados anal{\'i}ticos y alertas para coordinar acciones operativas y estrat{\'e}gicas}
	}
	
	\newglossaryentry{sistemaescalable}{
		name={Sistema escalable},
		text={sistema escalable},
		sort={sistema escalable},
		plural={sistemas escalables},
		first={sistema escalable},
		firstplural={sistemas escalables},
		description={Un \emph{sistema escalable} es aquel que mantiene su funcionamiento correcto y eficiente incluso cuando aumenta el n{\'u}mero de usuarios, el volumen de datos o las funcionalidades. En \glspl{sibiot}, la escalabilidad implica que el sistema pueda crecer horizontalmente (m{\'a}s nodos o instancias) o verticalmente (mayores recursos por nodo), garantizando \gls{rendimiento}, \gls{fiabilidad} y \gls{interoperabilidad} a medida que se expande.},
		type=main,
		symbol={📈},
		user1={Capacidad de crecimiento horizontal y vertical},
		user2={Adaptaci{\'o}n a m{\'a}s usuarios o funciones},
		user3={Factor esencial en arquitecturas IoT},
		see={escalable,escalabilidad,adaptabilidad,sibiot}
	}
	
	\newglossaryentry{sistemaidentificacion}{
		name={Sistema de identificaci{\'o}n},
		text={sistema de identificaci{\'o}n},
		sort={sistema de identificacion},
		plural={sistemas de identificaci{\'o}n},
		first={sistema de identificaci{\'o}n},
		firstplural={sistemas de identificaci{\'o}n},
		description={Conjunto de tecnolog{\'i}as, \glspl{dispositivo} y procedimientos destinados a reconocer y verificar la identidad de una persona, objeto o \gls{dispositivo}. Los \emph{sistemas de identificaci{\'o}n} pueden basarse en credenciales f{\'i}sicas (tarjetas, c{\'o}digos), biometr{\'i}a (huellas digitales, \gls{reconocimientofacial}, iris, voz) o identificadores electr{\'o}nicos (\gls{rfid}, \gls{nfc}). En el contexto de los \glspl{sibiot}, estos sistemas permiten gestionar el acceso a \glspl{sensor}, \glspl{actuador} y \glspl{plataforma} de \glspl{dato}, garantizando la \gls{seguridad}, la \gls{privacidad} y la \gls{trazabilidad} de las interacciones. Adem{\'a}s, se integran con mecanismos de \gls{autenticacion} y \gls{autorizacion} dentro de los esquemas de \gls{controlacceso}, favoreciendo la confianza y la \gls{fiabilidad} en aplicaciones industriales, urbanas y \gls{domotica}.},
		type=main,
		symbol={🆔},
		user1={Reconocimiento de identidades humanas o digitales},
		user2={Uso de biometr{\'i}a y credenciales},
		user3={Aplicaciones en seguridad y SIbIoT},
		see={controlacceso,autenticacion,autorizacion,seguridad,privacidad,sibiot}
	}
	
	% Alias/sin{\'o}nimo que remite a la entrada principal
	\newglossaryentry{sistemaincrustado}{
		name={Sistema incrustado},
		text={sistema incrustado},
		sort={sistema incrustado},
		plural={sistemas incrustados},
		description={V{\'e}ase sistemaembebido.},
		type=main,
		symbol={🔧},
		see={sistemaembebido}
	}
	
	\newglossaryentry{sistemainformacion}{
		name={Sistema de informaci{\'o}n},
		text={IS},
		sort={si},
		plural={IS},
		first={sistema de informaci{\'o}n (\textit{Information System}, IS)},
		firstplural={IS},
		description={Conjunto organizado de recursos humanos, tecnol{\'o}gicos, materiales y procedimientos que permiten recopilar, procesar, almacenar y distribuir informaci{\'o}n para apoyar la toma de decisiones, la coordinaci{\'o}n y el control en una organizaci{\'o}n. Un \emph{sistema de informaci{\'o}n} integra \gls{hardware}, \gls{software}, \gls{dato} y procesos de negocio para transformar datos en conocimiento {\'u}til. En el contexto de los \glspl{sibiot}, estos sistemas se potencian mediante la incorporaci{\'o}n de \glspl{sensor}, \glspl{actuador} y servicios en la nube, facilitando la monitorizaci{\'o}n en \gls{tiemporeal}, la \gls{automatizacion} y la generaci{\'o}n de valor estrat{\'e}gico para distintos sectores como salud, industria, \gls{agriculturainteligente} o \glspl{ciudadinteligente}.},
		type=main,
		symbol={💻📊},
		user1={Conjunto de recursos para gestionar informaci{\'o}n},
		user2={Soporte a la toma de decisiones},
		user3={Base de los SIbIoT y sistemas inteligentes},
		see={sibiot,software,dato,automatizacion}
	}
	
	\newglossaryentry{sistemainformacioncorporativo}{
		name={Sistema de informaci{\'o}n corporativo},
		text={sistema de informaci{\'o}n corporativo},
		sort={sistema informacion corporativo},
		plural={sistemas de informaci{\'o}n corporativos},
		first={sistema de informaci{\'o}n corporativo},
		firstplural={sistemas de informaci{\'o}n corporativos},
		description={Sistema destinado a apoyar la gesti{\'o}n y toma de decisiones dentro de una organizaci{\'o}n mediante el procesamiento de informaci{\'o}n estructurada. En los \glspl{sibiot}, estos sistemas integran datos operativos provenientes del entorno f{\'i}sico.},
		type=main,
		symbol={🏢},
		user1={Gesti{\'o}n organizacional},
		user2={Soporte a decisiones},
		user3={Informaci{\'o}n empresarial},
		see={tomadecisionesestrategicas,integraciondatos}
	}
	
	\newglossaryentry{sistemainformatico}{
		name={Sistema inform{\'a}tico},
		text={sistema inform{\'a}tico},
		sort={sistema informatico},
		plural={sistemas inform{\'a}ticos},
		first={sistema inform{\'a}tico},
		firstplural={sistemas inform{\'a}ticos},
		description={Conjunto integrado de \gls{hardware}, \gls{software}, datos y recursos humanos que interact{\'u}an para procesar, almacenar y transmitir informaci{\'o}n de forma autom{\'a}tica. Un \emph{sistema inform{\'a}tico} se orienta principalmente al tratamiento electr{\'o}nico de datos, sirviendo como base para aplicaciones de negocio, servicios en la \gls{nube}, plataformas web y sistemas de \gls{control}. En el contexto de los \glspl{sibiot}, los sistemas inform{\'a}ticos act{\'u}an como la infraestructura de soporte que conecta \glspl{sensor}, \glspl{actuador} y dispositivos inteligentes con algoritmos de \gls{monitorizacion}, \gls{optimizacion} y \gls{tomadecisiones}, garantizando \gls{eficienciaoperativa}, \gls{seguridad} y \gls{fiabilidad}.},
		type=main,
		symbol={💻},
		user1={Infraestructura de hardware y software},
		user2={Procesamiento y gesti{\'o}n de datos autom{\'a}ticos},
		user3={Soporte a SIbIoT y sistemas de informaci{\'o}n},
		see={sistemainformacion,hardware,software,sibiot}
	}
	
	\newglossaryentry{sistemaintegrado}{
		name={Sistema integrado},
		text={sistema integrado},
		sort={sistema integrado},
		plural={sistemas integrados},
		first={sistema integrado},
		firstplural={sistemas integrados},
		description={Conjunto de componentes f{\'i}sicos, l{\'o}gicos y organizacionales que operan de manera coordinada y cohesionada para cumplir una finalidad com{\'u}n, compartiendo informaci{\'o}n, recursos y mecanismos de control. Un sistema integrado elimina \glspl{redundancia}, reduce inconsistencias y asegura interoperabilidad entre sus partes. En el contexto de sistemas \gls{iot} y \glspl{sibiot}, un sistema integrado articula \glspl{sensor}, \glspl{actuador}, \glspl{dispositivoembebido}, plataformas en la \gls{nube}, \glspl{serviciodistribuido} y mecanismos de \gls{procesamiento} y \gls{tomadecisiones}, garantizando coherencia arquitect{\'o}nica, \gls{trazabilidad} y \gls{evolucioncontrolada}.},
		type=main,
		symbol={🔗},
		user1={Integraci{\'o}n funcional y estructural},
		user2={Coordinaci{\'o}n extremo a extremo},
		user3={Interoperabilidad y cohesi{\'o}n sist{\'e}mica},
		see={integracion,interoperabilidad,infraestructuradistribuida,sibiot}
	}
	
	\newglossaryentry{sistemaintegral}{
		name={Sistema integral},
		text={sistema integral},
		sort={sistema integral},
		plural={sistemas integrales},
		first={sistema integral},
		firstplural={sistemas integrales},
		description={Conjunto organizado de \glspl{componente}, \glspl{proceso} y recursos que funcionan de manera coordinada y unificada para alcanzar objetivos comunes, garantizando coherencia y \gls{eficiencia} en su operaci{\'o}n. Un \emph{sistema integral} se caracteriza por la interconexi{\'o}n de subsistemas heterog{\'e}neos (hardware, software, datos, personas y procedimientos) bajo un marco de gesti{\'o}n y control compartido. En el contexto de los \glspl{sibiot}, un sistema integral combina \glspl{sensor}, \glspl{actuador}, \glspl{red} de comunicaci{\'o}n, \glspl{plataforma} de an{\'a}lisis y \glspl{aplicacion} para ofrecer soluciones completas de \gls{monitorizacion}, \gls{control} y \gls{optimizacion} en \gls{tiemporeal}, asegurando \gls{interoperabilidad}, \gls{seguridad} y \gls{sostenibilidad}.},
		type=main,
		symbol={🔗},
		user1={Conjunto unificado de subsistemas},
		user2={Enfoque hol{\'i}stico de gesti{\'o}n},
		user3={Soluci{\'o}n completa en SIbIoT},
		see={sistema,sibiot,interoperabilidad,automatizacion}
	}
	
	\newglossaryentry{sistemainteligente}{
		name={Sistema inteligente},
		text={sistema inteligente},
		sort={sistema inteligente},
		plural={sistemas inteligentes},
		first={sistema inteligente},
		firstplural={sistemas inteligentes},
		description={Conjunto de componentes de hardware y software capaz de percibir su entorno mediante \glspl{sensor}, procesar informaci{\'o}n con algoritmos avanzados (incluyendo \gls{ia} y \gls{aprendizajeautomatico}), y actuar sobre dicho entorno a trav{\'e}s de \glspl{actuador} de manera aut{\'o}noma o semi-aut{\'o}noma para cumplir objetivos espec{\'i}ficos. Un \emph{sistema inteligente} integra capacidades de \gls{razonamiento}, adaptaci{\'o}n y aprendizaje, lo que le permite mejorar su desempe{\~n}o con la experiencia y reaccionar a condiciones cambiantes. En el contexto de los \glspl{sibiot}, los sistemas inteligentes facilitan la toma de decisiones en tiempo real, optimizan recursos, detectan anomal{\'i}as y ofrecen servicios personalizados en dominios como salud, industria, agricultura, ciudades inteligentes y hogares dom{\'o}ticos.},
		type=main,
		symbol={🤖},
		user1={Sistemas aut{\'o}nomos y adaptativos},
		user2={Integraci{\'o}n de IA y aprendizaje autom{\'a}tico},
		user3={Soporte a la toma de decisiones en tiempo real},
		see={sibiot,iot,ia,aprendizajeautomatico,automatizacion}
	}
	
	\newglossaryentry{sistemaoperativo}{
		name={Sistema Operativo},
		text={sistema operativo},
		sort={sistemaoperativo},
		plural={sistemas operativos},
		first={sistema operativo},
		firstplural={sistemas operativos},
		description={Software fundamental que act{\'u}a como intermediario entre el hardware de un computador o dispositivo electr{\'o}nico y los programas de usuario. Se encarga de la gesti{\'o}n de recursos como la memoria, el procesador, los dispositivos de entrada/salida y los sistemas de archivos. Adem{\'a}s, define una interfaz para la ejecuci{\'o}n de aplicaciones, controla la seguridad y garantiza la estabilidad del sistema. Ejemplos com{\'u}nes incluyen Windows, Linux, macOS y Android. En el contexto de los \glspl{sibiot}, los sistemas operativos especializados (como sistemas en tiempo real o versiones embebidas) permiten la coordinaci{\'o}n eficiente de sensores, actuadores y comunicaciones distribuidas.},
		type=main,
		symbol={💻},
		user1={Software base},
		user2={Gestor de recursos},
		user3={Interfaz entre hardware y aplicaciones},
		see={sistemainformatico,software}
	}
	
	\newglossaryentry{sistemanavegacion}{
		name={Sistema de navegaci{\'o}n},
		text={sistema de navegaci{\'o}n},
		sort={sistema de navegacion},
		plural={sistemas de navegaci{\'o}n},
		first={sistema de navegaci{\'o}n},
		firstplural={sistemas de navegaci{\'o}n},
		description={Conjunto de tecnolog{\'i}as, dispositivos y algoritmos dise{\~n}ados para determinar la posici{\'o}n, orientaci{\'o}n y trayectoria de un objeto o usuario, y guiarlo de manera precisa hacia un destino. Un \emph{sistema de navegaci{\'o}n} combina \glspl{sensor} como aceler{\'o}metros, giroscopios y receptores GNSS (GPS, GLONASS, Galileo, BeiDou), con \glspl{algoritmo} de fusi{\'o}n de datos y visualizaci{\'o}n en \glspl{aplicacion}. En el contexto de los \glspl{sibiot}, estos sistemas se emplean en veh{\'i}culos aut{\'o}nomos, drones, robots m{\'o}viles, log{\'i}stica inteligente y aplicaciones de movilidad urbana, garantizando \gls{precision}, eficiencia y seguridad en la toma de decisiones en \gls{tiemporeal}.},
		type=main,
		symbol={🧭},
		user1={Determinaci{\'o}n de posici{\'o}n y trayectoria},
		user2={Integraci{\'o}n de sensores y GNSS},
		user3={Aplicaciones en movilidad e IoT},
		see={sensor,algoritmo,aplicacion,precision,sibiot}
	}
	
	\newglossaryentry{sistemapeltier}{
		name={Sistema Peltier},
		text={sistema Peltier},
		sort={sistema peltier},
		plural={sistemas Peltier},
		first={sistema Peltier},
		firstplural={sistemas Peltier},
		description={Dispositivos termoel{\'e}ctricos que funcionan en base al efecto Peltier, por el cual el paso de corriente el{\'e}ctrica a trav{\'e}s de la uni{\'o}n de dos materiales semiconductores genera un flujo de calor, provocando enfriamiento en un lado y calentamiento en el otro. Los \emph{sistemas Peltier} se utilizan en aplicaciones de refrigeraci{\'o}n localizada, control de temperatura y generaci{\'o}n de energ{\'i}a. En el contexto de los \glspl{sibiot}, son empleados para mantener condiciones t{\'e}rmicas estables en equipos sensibles, mejorar la \gls{eficienciatermica}, controlar \glspl{variablefisica} como temperatura y humedad, e integrarse con \glspl{sensor} y algoritmos de \gls{control} para aplicaciones en salud, agricultura de precisi{\'o}n, electr{\'o}nica y entornos industriales.},
		type=main,
		symbol={❄️🔥},
		user1={Dispositivos de efecto termoel{\'e}ctrico},
		user2={Refrigeraci{\'o}n y control de temperatura},
		user3={Aplicaciones en IoT y sistemas inteligentes},
		see={eficienciatermica,variablefisica,sensor,control,sibiot}
	}
	
	\newglossaryentry{sistemapersonalizado}{
		name={sistema personalizado},
		text={sistema personalizado},
		sort={sistema personalizado},
		plural={sistemas personalizados},
		first={sistema personalizado},
		firstplural={sistemas personalizados},
		description={Soluci{\'o}n de software o \gls{sibiot} dise{\~n}ada y configurada espec{\'i}ficamente para satisfacer requisitos particulares de un usuario, organizaci{\'o}n o dominio funcional. Un \emph{sistema personalizado} adapta sus componentes, flujos de datos, interfaces y reglas de negocio a necesidades concretas, ofreciendo flexibilidad y mayor adecuaci{\'o}n al contexto de uso. Este enfoque es fundamental cuando se requiere alineaci{\'o}n precisa con procesos internos, restricciones de operaci{\'o}n o integraci{\'o}n con infraestructura existente},
		type=main,
		symbol={⚙️},
		user1={Adaptaci{\'o}n a requisitos espec{\'i}ficos},
		user2={Configuraci{\'o}n ajustada a procesos organizacionales},
		user3={Integraci{\'o}n con ecosistemas tecnol{\'o}gicos existentes},
		see={personalizacion,disenosistemas,sibiot}
	}
	
	\newglossaryentry{sistemarespaldoenergia}{
		name={Sistema de respaldo de energ{\'i}a},
		text={sistema de respaldo de energ{\'i}a},
		sort={sistema de respaldo de energia},
		plural={sistemas de respaldo de energ{\'i}a},
		first={sistema de respaldo de energ{\'i}a},
		firstplural={sistemas de respaldo de energ{\'i}a},
		description={Conjunto de dispositivos y tecnolog{\'i}as dise{\~n}ados para garantizar el suministro el{\'e}ctrico continuo en caso de fallos, cortes o inestabilidad de la red principal. Los \emph{sistemas de respaldo de energ{\'i}a} incluyen bater{\'i}as, generadores, sistemas de alimentaci{\'o}n ininterrumpida (UPS) y soluciones basadas en energ{\'i}as renovables. En el contexto de los \glspl{sibiot}, son fundamentales para asegurar la \gls{disponibilidad} de \glspl{sensor}, \glspl{actuador}, \glspl{nodo} y plataformas en la \gls{nube}, garantizando \gls{fiabilidad} y continuidad operativa en aplicaciones cr{\'i}ticas como salud, transporte, industria y centros de datos. Adem{\'a}s, contribuyen a la \gls{resiliencia} y a la \gls{sostenibilidad} al integrarse con tecnolog{\'i}as inteligentes de gesti{\'o}n energ{\'e}tica.},
		type=main,
		symbol={🔋⚡},
		user1={Garant{\'i}a de continuidad el{\'e}ctrica},
		user2={Prevenci{\'o}n de fallos en sistemas cr{\'i}ticos},
		user3={Soporte a la resiliencia en IoT},
		see={fiabilidad,disponibilidad,resiliencia,sostenibilidad,sibiot}
	}
	
	\newglossaryentry{sistematransaccional}{
		name={Sistema Transaccional},
		text={sistema transaccional},
		sort={sistema transaccional},
		plural={sistemas transaccionales},
		first={sistema transaccional},
		firstplural={sistemas transaccionales},
		description={Sistema de informaci{\'o}n dise{\~n}ado para registrar, procesar y controlar transacciones operativas de manera confiable y consistente, garantizando integridad, atomicidad y trazabilidad de los datos en actividades rutinarias como ventas, pagos, registros o actualizaciones},
		type=main,
		symbol={🔄},
		see={sistemainformacion,transaccion,basedatos}
	}
	
	% ===================== Sitio web =====================
	\newglossaryentry{sitioweb}{
		name={Sitio web},
		text={sitio web},
		sort={sitio web},
		plural={sitios web},
		first={sitio web},
		firstplural={sitios web},
		description={Un \emph{sitio web} es un conjunto de p{\'a}ginas web relacionadas, agrupadas bajo un mismo dominio, que ofrecen informaci{\'o}n sobre una marca, organizaci{\'o}n, empresa o producto. Cada sitio web est{\'a} estructurado mediante \gls{html}, \gls{css} y \gls{javascript}, y es accesible mediante un \gls{browser}. En el contexto de \glspl{sibiot}, los sitios web sirven para presentar paneles de \gls{dashboard}, \gls{servicioweb} y aplicaciones de monitorizaci{\'o}n y \gls{control}.},
		type=main,
		symbol={🌐},
		user1={Conjunto de p{\'a}ginas bajo un dominio},
		user2={Soporte mediante HTML, CSS y JavaScript},
		user3={Interfaz de acceso a servicios IoT},
		see={browser,servicioweb,aplicacionweb,sibiot}
	}
	
	% ===================== Six Sigma =====================
	\newglossaryentry{sixsigma}{
		name={Six Sigma},
		text={Six Sigma},
		sort={six sigma},
		plural={metodolog{\'i}as Six Sigma},
		first={Six Sigma},
		firstplural={metodolog{\'i}as Six Sigma},
		description={\emph{Six Sigma} es una metodolog{\'i}a de gesti{\'o}n de calidad orientada a la mejora continua de procesos mediante la identificaci{\'o}n, an{\'a}lisis y reducci{\'o}n de la variabilidad. Se fundamenta en herramientas estad{\'i}sticas y en las fases DMAIC (Definir, Medir, Analizar, Mejorar y Controlar). En sistemas \gls{iot}, Six Sigma permite optimizar la fabricaci{\'o}n inteligente, el \gls{mantenimientopredictivo} y el control de calidad automatizado en l{\'i}neas de producci{\'o}n conectadas.},
		type=main,
		symbol={➿},
		user1={Metodolog{\'i}a de calidad DMAIC},
		user2={Reducci{\'o}n de variabilidad en procesos},
		user3={Aplicaci{\'o}n en manufactura e IoT},
		see={manufactura,mantenimientopredictivo,calidad}
	}
	
	\newglossaryentry{sla}{
		name={Acuerdos de nivel de servicio},
		text={SLA},
		plural={SLA},
		first={acuerdos de nivel de servicio o Service Level Agreements, SLA)},
		firstplural={acuerdos de nivel de servicio que definen responsabilidades, m{\'e}tricas y penalizaciones asociadas al cumplimiento del servicio},
		description={Instrumentos contractuales que especifican indicadores cuantificables como disponibilidad, tiempo de respuesta, latencia y soporte, utilizados para garantizar y evaluar la calidad de servicios tecnol{\'o}gicos y empresariales},
		sort={acuerdos de nivel de servicio},
		type=main,
		symbol={📄},
		see={calidadservicio,gestionservicio}
	}
	
	\newglossaryentry{smtp}{
		name={SMTP},
		text={SMTP},
		sort={smtp},
		plural={protocolos SMTP},
		first={protocolo simple de transferencia de correo (\textit{Simple Mail Transfer Protocol}, SMTP)},
		firstplural={protocolos SMTP},
		description={\emph{SMTP} es un protocolo est{\'a}ndar de la capa de aplicaci{\'o}n usado para el env{\'i}o de mensajes de correo electr{\'o}nico entre servidores. Opera com{\'u}nmente junto con POP o IMAP para la recepci{\'o}n de mensajes. En \glspl{sibiot}, se utiliza para notificaciones autom{\'a}ticas, alertas de \gls{seguridad} o comunicaci{\'o}n entre plataformas de monitorizaci{\'o}n y usuarios finales.},
		type=main,
		symbol={✉️},
		user1={Protocolo de env{\'i}o de correos electr{\'o}nicos},
		user2={Operaci{\'o}n junto con POP/IMAP},
		user3={Uso en notificaciones IoT},
		see={protocolocomunicacion,seguridad}
	}
	
	\newglossaryentry{soa}{
		name={SOA},
		text={SOA},
		sort={soa},
		plural={arquitecturas orientadas a servicios},
		first={Arquitectura orientada a servicios (\textit{Service-Oriented Architecture}, SOA)},
		firstplural={arquitecturas SOA},
		description={\emph{SOA} es un modelo de dise{\~n}o de software en el que los componentes ofrecen servicios a trav{\'e}s de una red mediante protocolos estandarizados, como \gls{http}, \gls{https} o \gls{soap}. Cada servicio es modular, reutilizable y accesible mediante interfaces definidas, favoreciendo la \gls{interoperabilidad} y la \gls{escalabilidad}. En \glspl{sibiot}, SOA habilita la integraci{\'o}n de \glspl{servicioweb}, \glspl{aplicacionweb} y \gls{sistemadistribuido} heterog{\'e}neos.},
		type=main,
		symbol={🧩},
		user1={Arquitectura basada en servicios independientes},
		user2={Uso de protocolos estandarizados (HTTP, SOAP)},
		user3={Soporte para escalabilidad e interoperabilidad en IoT},
		see={servicioweb,aplicacionweb,arquitecturadistribuida,interoperabilidad}
	}
	
	\newglossaryentry{soap}{
		name={SOAP},
		text={SOAP},
		sort={soap},
		plural={protocolos SOAP},
		first={protocolo simple de acceso a objetos (\textit{Simple Object Access Protocol}, SOAP)},
		firstplural={protocolos SOAP},
		description={\emph{SOAP} es un protocolo estandarizado basado en XML para el intercambio estructurado de mensajes entre aplicaciones a trav{\'e}s de redes inform{\'a}ticas. Se utiliza principalmente en la implementaci{\'o}n de \glspl{servicioweb}, ofreciendo extensibilidad, neutralidad e independencia de la plataforma y el lenguaje de programaci{\'o}n. Aunque m{\'a}s pesado que alternativas como \gls{rest}, \gls{soap} proporciona mecanismos robustos de seguridad, transacciones y confiabilidad en la comunicaci{\'o}n. En sistemas \gls{iot}, puede emplearse en entornos corporativos y cr{\'i}ticos que requieren \gls{integridad}, \gls{interoperabilidad} con sistemas legados y cumplimiento de est{\'a}ndares.},
		type=main,
		symbol={🧼},
		user1={Protocolo basado en XML},
		user2={Servicios web interoperables},
		user3={Soporte a seguridad y transacciones},
		see={servicioweb,soa,api,iot}
	}
	
	\newglossaryentry{sobrecarga}{
		name={Sobrecarga},
		text={sobrecarga},
		sort={sobrecarga},
		plural={sobrecargas},
		first={sobrecarga},
		firstplural={sobrecargas},
		description={Condici{\'o}n en la que un \gls{sistema}, componente o \gls{red} recibe una demanda de recursos superior a su capacidad de procesamiento, almacenamiento o comunicaci{\'o}n, lo que provoca degradaci{\'o}n del \gls{rendimiento}, aumento de la \gls{latencia} o incluso fallos. La \emph{sobrecarga} puede ocurrir por exceso de peticiones, saturaci{\'o}n de \glspl{nodo}, complejidad en el procesamiento de datos o ineficiencia en los algoritmos. En el contexto de los \glspl{sibiot}, la sobrecarga afecta a la \gls{monitorizacion} en \gls{tiemporeal}, la coordinaci{\'o}n de \glspl{actuador} y la integraci{\'o}n de servicios distribuidos, siendo necesario aplicar \gls{optimizacion}, balanceo de cargas, \gls{toleranciafallos} y estrategias de \gls{escalabilidad} para mitigar su impacto.},
		type=main,
		symbol={⚠️},
		user1={Exceso de demanda sobre recursos},
		user2={Degradaci{\'o}n del rendimiento o fallos},
		user3={Riesgo en sistemas distribuidos e IoT},
		see={rendimiento,latencia,toleranciafallos,optimizacion,escalabilidad,sibiot}
	}
	
	% ===================== Sockets =====================
	\newglossaryentry{sockets}{
		name={Sockets},
		text={sockets},
		sort={sockets},
		plural={sockets},
		first={sockets},
		firstplural={sockets},
		description={\emph{Sockets} son puntos finales (\textit{endpoints}) de enlaces de \gls{comunicacion} entre procesos. Se tratan como descriptores de ficheros, lo que permite a los procesos intercambiar datos a trav{\'e}s de conexiones basadas en sockets. En sistemas distribuidos e \glspl{iot}, los sockets habilitan la comunicaci{\'o}n entre \glspl{dispositivo}, servidores y aplicaciones, soportando protocolos como \gls{tcp} o \gls{udp}.},
		type=main,
		symbol={🔌},
		see={protocolocomunicacion,cliente,servidor}
	}
	
	% ===================== Software =====================
	\newglossaryentry{software}{
		name={Software},
		text={software},
		sort={software},
		plural={softwares},
		first={software},
		firstplural={softwares},
		description={\emph{Software} es un concepto inform{\'a}tico que hace referencia a un programa o conjunto de programas de c{\'o}mputo, as{\'i} como a los datos asociados que permiten realizar distintas tareas en un \gls{sistemainformatico}. El software se encuentra no solo en computadoras, tablets y \textit{smartphones}, sino tambi{\'e}n en dispositivos electr{\'o}nicos como televisores, microondas o refrigeradores. En el contexto de \glspl{sibiot}, el software integra la \gls{logicaprogramacion}, el \gls{firmware} y las aplicaciones que habilitan la \gls{interoperabilidad}, el \gls{control} y la \gls{automatizacion}.},
		type=main,
		symbol={💾},
		user1={Conjunto de programas y datos},
		user2={Soporte a sistemas inform{\'a}ticos e IoT},
		user3={Complemento esencial del hardware},
		see={hardware,firmware,aplicacion,sibiot}
	}
	
	% ===================== Solicitud HTTP =====================
	\newglossaryentry{solicitudhttp}{
		name={Solicitud HTTP},
		text={solicitud HTTP},
		sort={solicitud http},
		plural={solicitudes HTTP},
		first={solicitud HTTP},
		firstplural={solicitudes HTTP},
		description={Una \emph{solicitud HTTP} es el mensaje que un \gls{cliente} env{\'i}a a un \gls{servidorweb} para obtener un recurso o ejecutar una acci{\'o}n. Incluye un m{\'e}todo (GET, POST, PUT, DELETE, etc.), la URI/URL del recurso, encabezados con metadatos y, opcionalmente, un cuerpo con datos adicionales. En aplicaciones \gls{iot}, las solicitudes HTTP permiten el \gls{control} remoto, la transferencia de datos desde \glspl{dispositivoiot} y la integraci{\'o}n con \glspl{servicioweb}.},
		type=main,
		symbol={📡},
		user1={Mensaje de petici{\'o}n de cliente a servidor},
		user2={Incluye m{\'e}todo, URI, encabezados y cuerpo},
		user3={Usado en aplicaciones web e IoT},
		see={http,https,solicitudpost,servicioweb}
	}
	
	% ===================== Solicitud POST =====================
	\newglossaryentry{solicitudpost}{
		name={Solicitud POST},
		text={solicitud POST},
		sort={solicitud post},
		plural={solicitudes POST},
		first={solicitud POST},
		firstplural={solicitudes POST},
		description={Una \emph{solicitud POST} es un tipo de \gls{solicitudhttp} que env{\'i}a datos en el cuerpo del mensaje al servidor para crear o procesar recursos. A diferencia de GET, no incluye los datos en la URL y no es idempotente, lo que implica que puede modificar el estado del sistema en cada ejecuci{\'o}n. En entornos \gls{iot}, se utiliza para registrar datos de sensores, enviar comandos a actuadores o almacenar informaci{\'o}n en la \gls{nube}.},
		type=main,
		symbol={📤},
		user1={Tipo de solicitud HTTP no idempotente},
		user2={Env{\'i}a datos en el cuerpo del mensaje},
		user3={Aplicaci{\'o}n en IoT para comandos y registros},
		see={solicitudhttp,http,https}
	}
	
	\newglossaryentry{solucion}{
		name={Soluci{\'o}n},
		text={soluci{\'o}n},
		sort={solucion},
		plural={soluciones},
		first={soluci{\'o}n},
		firstplural={soluciones},
		description={Respuesta estructurada a un problema espec{\'i}fico, materializada mediante una combinaci{\'o}n coherente de modelos, componentes, tecnolog{\'i}as y decisiones arquitect{\'o}nicas. En sistemas IoT, una soluci{\'o}n integra elementos f{\'i}sicos y digitales para generar valor funcional.},
		type=main,
		symbol={🧩},
		user1={Respuesta estructurada},
		user2={Integraci{\'o}n tecnol{\'o}gica},
		user3={Valor funcional},
		see={modelo,construccion}
	}
	
	% ===================== Soluci{\'o}n IoT =====================
	\newglossaryentry{solucioniot}{
		name={Soluci{\'o}n IoT},
		text={soluci{\'o}n IoT},
		sort={solucion iot},
		plural={soluciones IoT},
		first={soluci{\'o}n IoT},
		firstplural={soluciones IoT},
		description={Una \emph{soluci{\'o}n IoT} es un conjunto integrado de componentes tecnol{\'o}gicos dise{\~n}ados para resolver un problema espec{\'i}fico o automatizar un proceso mediante el uso de tecnolog{\'i}as \gls{iot}. Una soluci{\'o}n IoT puede incluir \glspl{sensor}, \glspl{actuador}, \glspl{microcontrolador}, redes de \gls{comunicacion}, \glspl{gateway}, plataformas de \gls{procesamientodatos}, algoritmos de \glspl{modelopredictivo} e interfaces de usuario. Su objetivo es habilitar la \gls{monitorizacion}, el \gls{control} y la \gls{tomadecisiones} inteligente en sectores como salud, manufactura, transporte o agricultura.},
		type=main,
		symbol={🛠️},
		user1={Conjunto integrado de hardware, software y redes},
		user2={Soporte a monitorizaci{\'o}n, control y automatizaci{\'o}n},
		user3={Aplicaciones en salud, manufactura y transporte},
		see={iot,sensor,actuador,gateway,procesamientodatos,sibiot}
	}
	
	\newglossaryentry{sostenibilidad}{
		name={Sostenibilidad},
		text={sostenibilidad},
		sort={sostenibilidad},
		plural={sostenibilidades},
		first={sostenibilidad},
		firstplural={sostenibilidades},
		description={Capacidad de un sistema, organizaci{\'o}n o actividad para mantenerse en el tiempo garantizando el equilibrio entre el desarrollo econ{\'o}mico, la protecci{\'o}n ambiental y el bienestar social. La \emph{sostenibilidad} implica la gesti{\'o}n responsable de los recursos naturales y humanos, evitando su agotamiento o degradaci{\'o}n, y promoviendo pr{\'a}cticas que aseguren la satisfacci{\'o}n de las necesidades presentes sin comprometer las de las generaciones futuras. En el contexto de la ingenier{\'i}a, las \gls{tic} y los \glspl{sibiot}, la sostenibilidad incluye el dise{\~n}o y uso de tecnolog{\'i}as eficientes energ{\'e}ticamente, el fomento de la econom{\'i}a circular, la reducci{\'o}n de residuos y emisiones, y la implementaci{\'o}n de soluciones que equilibren la rentabilidad con el impacto social y ambiental positivo.},
		type=main,
		symbol={},
		user1={Responsabilidad ambiental y social},
		user2={Principio rector de proyectos tecnol{\'o}gicos},
		user3={Enfoque integral econ{\'o}mico--social--ambiental},
		see={medioambiente,eficienciaenergetica,sibiot}
	}
	
	\newglossaryentry{sostenible}{
		name={Sostenible},
		text={sostenible},
		sort={sostenible},
		plural={sostenibles},
		first={sostenible},
		firstplural={sostenibles},
		description={Adjetivo que describe a un \gls{sistema}, tecnolog{\'i}a, proceso u organizaci{\'o}n capaz de mantener su funcionamiento en el tiempo sin agotar los recursos disponibles ni causar da{\~n}os irreversibles al medio ambiente o a la sociedad. Algo \emph{sostenible} equilibra factores econ{\'o}micos, ambientales y sociales, asegurando viabilidad a largo plazo. En el contexto de los \glspl{sibiot}, se refiere a soluciones que optimizan el consumo de energ{\'i}a, reducen emisiones, prolongan la vida {\'u}til de los \glspl{dispositivo}, fomentan la \gls{eficienciaoperativa} y garantizan beneficios para usuarios y comunidades sin comprometer el futuro.},
		type=main,
		symbol={♻️},
		user1={Capacidad de mantenerse en el tiempo},
		user2={Equilibrio econ{\'o}mico, ambiental y social},
		user3={Objetivo en el dise{\~n}o de SIbIoT},
		see={sostenibilidad,eficiencia,eficienciaoperativa,sibiot}
	}
	
	% ===================== Spanner =====================
	\newglossaryentry{spanner}{
		name={Spanner},
		text={Spanner},
		sort={spanner},
		first={Spanner},
		description={Nombre corto y conceptual utilizado para referirse a \gls{googlecloudspanner}, especialmente en literatura acad{\'e}mica sobre sistemas distribuidos. El t{\'e}rmino \emph{Spanner} se asocia al modelo de base de datos relacional distribuida con consistencia fuerte global, escalabilidad horizontal y sincronizaci{\'o}n temporal mediante relojes distribuidos},
		type=main,
		symbol={⏱️},
		see={googlecloudspanner,cloudspanner,sistemadistribuido}
	}
	
	\newglossaryentry{spark}{
		name={Apache Spark},
		text={Spark},
		sort={spark},
		plural={ecosistemas Spark},
		first={Apache Spark},
		firstplural={ecosistemas Spark},
		description={Plataforma de c{\'o}digo abierto dise{\~n}ada para el procesamiento distribuido de datos a gran escala, desarrollada por la Apache Software Foundation. \emph{Spark} permite realizar an{\'a}lisis de datos en memoria de manera r{\'a}pida y eficiente, soportando tareas de procesamiento por lotes, flujos de datos en \gls{tiemporeal}, consultas \gls{sql}, aprendizaje autom{\'a}tico y an{\'a}lisis de grafos. En el contexto de los \glspl{sibiot}, se utiliza para gestionar y analizar informaci{\'o}n masiva proveniente de \glspl{sensor} y \glspl{dispositivo}, habilitando modelos predictivos y decisiones basadas en datos.},
		type=main,
		symbol={⚡},
		user1={Procesamiento distribuido en memoria},
		user2={Plataforma de Big Data},
		user3={Apache Spark},
		see={hadoop,mineriadatos,bigdata}
	}
	
	\newglossaryentry{spi}{
		name={SPI},
		text={SPI},
		sort={spi},
		first={interfaz de comunicaci{\'o}n serie sincr{\'o}nica (\emph{Serial Peripheral Interface}, SPI)},
		description={Interfaz de comunicaci{\'o}n serie sincr{\'o}nica o \emph{Serial Peripheral Interface}, SPI) utilizada para la comunicaci{\'o}n r{\'a}pida entre un microcontrolador maestro y uno o m{\'a}s dispositivos esclavos. \emph{SPI} emplea l{\'i}neas dedicadas para reloj, datos y selecci{\'o}n de dispositivo, ofreciendo alta velocidad y baja latencia en la transferencia de datos.},
		type=main,
		symbol={🚀},
		user1={Alta velocidad},
		user2={Comunicaci{\'o}n sincr{\'o}nica},
		user3={Perif{\'e}ricos},
		see={uart,capared}
	}
	
	\newglossaryentry{sprint}{
		name={Sprint},
		text={\textit{Sprint}},
		sort={sprint},
		plural={\textit{Sprints}},
		first={\textit{Sprint}},
		firstplural={\textit{Sprints}},
		description={Iteraci{\'o}n de duraci{\'o}n fija (generalmente entre una y cuatro semanas) en la cual el equipo Scrum desarrolla un incremento funcional del producto. Cada Sprint incluye planificaci{\'o}n, desarrollo, revisi{\'o}n y retrospectiva, y finaliza con un incremento potencialmente desplegable.},
		type=main,
		symbol={🏃},
		user1={Iteraci{\'o}n temporal fija},
		user2={Incremento funcional},
		user3={Ciclo completo de desarrollo},
		see={productbacklog,sprintbacklog,agil,incremento}
	}
	
	\newglossaryentry{sprintbacklog}{
		name={Sprint Backlog},
		text={\textit{Sprint Backlog}},
		sort={sprint backlog},
		plural={\textit{Sprint Backlogs}},
		first={\textit{Sprint Backlog}},
		firstplural={\textit{Sprint Backlogs}},
		description={Conjunto de elementos seleccionados del \gls{productbacklog} para ser desarrollados durante un \gls{sprint}, junto con un plan detallado para entregar el \gls{incremento} comprometido. Es gestionado por el equipo de desarrollo y se actualiza diariamente conforme avanza el trabajo.},
		type=main,
		symbol={🗂️},
		user1={Trabajo comprometido en el Sprint},
		user2={Plan detallado de ejecuci{\'o}n},
		user3={Gesti{\'o}n diaria por el equipo},
		see={productbacklog,sprint,agil,incremento}
	}
	
	\newglossaryentry{sprintgoal}{
		name={Sprint Goal},
		text={Sprint Goal},
		sort={sprint goal},
		plural={Sprint Goals},
		first={Sprint Goal},
		firstplural={Sprint Goals},
		description={Meta concreta establecida para un Sprint que describe el prop{\'o}sito funcional del incremento a desarrollar. Proporciona foco y coherencia al equipo durante la iteraci{\'o}n, permitiendo cierta flexibilidad en la implementaci{\'o}n siempre que se cumpla el objetivo definido.},
		type=main,
		symbol={🏁},
		user1={Meta del Sprint},
		user2={Prop{\'o}sito del incremento},
		user3={Foco del equipo},
		see={sprint,sprintbacklog,incremento}
	}
	
	\newglossaryentry{sql}{
		name={SQL},
		text={SQL},
		sort={sql},
		plural={SQLs},
		first={lenguaje de consulta estructurado (\textit{Structured Query Language}, SQL)},
		firstplural={SQLs},
		description={Lenguaje est{\'a}ndar utilizado para gestionar y manipular datos en \glspl{basedatosrelacional}. \emph{SQL} permite definir estructuras de datos (DDL), manipular registros (DML), controlar el acceso (DCL) y realizar consultas (SELECT) de forma eficiente. Est{\'a} normalizado por ANSI/ISO y es compatible con la mayor{\'i}a de sistemas de gesti{\'o}n como \gls{mysql}, \gls{postgresql}, \gls{oracledatabase} o \gls{microsoftsqlserver}. En el contexto de los \glspl{sibiot}, SQL se utiliza para gestionar datos estructurados provenientes de \glspl{sensor}, soportar an{\'a}lisis en \gls{tiemporeal}, generar reportes, y garantizar \gls{consistencia}, \gls{fiabilidad} y \gls{trazabilidad} en \glspl{aplicacion} cr{\'i}ticas.},
		type=main,
		symbol={📑},
		user1={Lenguaje est{\'a}ndar de gesti{\'o}n de datos},
		user2={Soporte en bases de datos relacionales},
		user3={Aplicaciones en IoT y Big Data},
		see={basedatosrelacional,mysql,postgresql,oracledatabase,microsoftsqlserver,sibiot}
	}
	
	\newglossaryentry{sqlserver}{
		name={SQL Server},
		text={SQL Server},
		sort={sql server},
		plural={instancias de SQL Server},
		first={Sistema de gesti{\'o}n de bases de datos relacional de Microsoft, (\textit{Microsoft SQL Server}},
		firstplural={instancias de SQL Server},
		description={Plataforma de gesti{\'o}n de datos relacional (SGBDR) desarrollada por Microsoft para cargas de trabajo \textit{OLTP} y \textit{OLAP}. Proporciona motor transaccional conforme a propiedades \gls{acid}, lenguaje \textit{Transact{-}SQL} (T{-}SQL), procedimientos almacenados, disparadores, vistas e {\'i}ndices avanzados (columnstore, particionado). Ofrece mecanismos de seguridad (cifrado, control de acceso basado en roles, auditor{\'i}a), alta disponibilidad y recuperaci{\'o}n ante desastres (Always On, replicaci{\'o}n, \textit{log shipping}), as{\'i} como integraci{\'o}n con herramientas del ecosistema: SSMS, SSIS, SSRS y SSAS. Est{\'a} disponible en ediciones \textit{Express}, \textit{Developer}, \textit{Standard} y \textit{Enterprise}, con despliegue local (\textit{on{-}premises}) o administrado en la nube mediante \textit{Azure SQL}.},
		type=main,
		user1={SGBDR de Microsoft para OLTP/OLAP},
		user2={Lenguaje T{-}SQL y objetos de base de datos},
		user3={Alta disponibilidad: Always On, replicaci{\'o}n, particionado},
		user4={Ecosistema: SSMS, SSIS, SSRS, SSAS; opciones en Azure},
		see={basedatos,sgbd,acid,tsql,azure,oltp,olap,replicacion},
		symbol={\ensuremath{\mathsf{DB}}}
	}
	
	\newglossaryentry{sram}{
		name={SRAM},
		text={SRAM},
		sort={sram},
		first={memoria de acceso aleatorio est{\'a}tica (\emph{Static Random-Access Memory}, SRAM)},
		description={Memoria de acceso aleatorio vol{\'a}til utilizada para almacenar variables, pilas y datos temporales durante la ejecuci{\'o}n de un programa. La \emph{SRAM} ofrece acceso r{\'a}pido, pero pierde su contenido cuando se interrumpe la alimentaci{\'o}n. En sistemas embebidos, su gesti{\'o}n eficiente es cr{\'i}tica debido a su capacidad limitada.},
		type=main,
		symbol={⚡},
		user1={Memoria vol{\'a}til},
		user2={Datos temporales},
		user3={Alta velocidad},
		see={memoriaflash,eeprom}
	}
	
	\newglossaryentry{ssd}{
		name={SSD},
		text={SSD},
		sort={ssd},
		first={unidad de estado s{\'o}lido (\textit{Solid State Drive}, SSD)},
		plural={unidades SSD},
		description={Dispositivo de almacenamiento de datos que utiliza memoria flash en lugar de discos magn{\'e}ticos como los discos duros tradicionales (HDD). Los \glspl{ssd} ofrecen velocidades de lectura y escritura superiores, menor latencia, mayor resistencia a impactos y menor consumo energ{\'e}tico. En sistemas \gls{iot}, los SSD se utilizan en \glspl{gateway}, servidores de borde o unidades de procesamiento intensivo que requieren almacenamiento local de alto \gls{rendimiento}.},
		type=main,
		symbol={💾},
		user1={Almacenamiento en memoria flash},
		user2={Alta velocidad y baja latencia},
		user3={Uso en IoT y servidores de borde},
		see={hdd,almacenamiento,iot}
	}
	
	\newglossaryentry{ssl}{
		name={SSL},
		text={SSL},
		first={capa de conexi{\'o}n segura (\textit{Secure Sockets Layer}, SSL)},
		sort={ssl},
		plural={protocolos SSL},
		description={Protocolo criptogr{\'a}fico dise{\~n}ado para comunicaciones seguras mediante cifrado de datos transmitidos entre cliente y servidor. Aunque SSL ha sido reemplazado por \gls{tls} debido a vulnerabilidades, el t{\'e}rmino sigue utiliz{\'a}ndose de forma gen{\'e}rica. En sistemas \gls{iot}, SSL/\gls{tls} se aplica en servicios como \gls{https}, \gls{mqtt} o correo electr{\'o}nico seguro para garantizar \gls{confidencialidad} e \gls{integridad}.},
		type=main,
		symbol={🔐},
		user1={Protocolo criptogr{\'a}fico legado},
		user2={Base de HTTPS y servicios IoT},
		user3={Reemplazado por TLS},
		see={tls,https,seguridad,confidencialidad,integridad}
	}
	
	\newglossaryentry{streaming}{
		name={Streaming},
		text={streaming},
		sort={streaming},
		plural={servicios de streaming},
		description={T{\'e}cnica que permite la transmisi{\'o}n continua de datos en \gls{tiemporeal}, posibilitando su procesamiento o consumo inmediato sin necesidad de descarga completa previa. Se aplica en video, audio, telemetr{\'i}a y flujos de datos IoT, donde la inmediatez es esencial.},
		type=main,
		symbol={📡},
		user1={Transmisi{\'o}n continua en tiempo real},
		user2={Video, audio y datos IoT},
		user3={Procesamiento inmediato sin descarga completa},
		see={tiemporeal,dato,iot}
	}
	
	\newglossaryentry{stm32}{
		name={STM32},
		text={STM32},
		first={microcontroladores de STMicroelectronics con arquitectura de 32 bits (\textit{STMicroelectronics Microcontroller of 32 bits}, STM32)}
		sort={stm32},
		plural={microcontroladores STM32},
		description={Familia de \glspl{microcontrolador} desarrollada por STMicroelectronics, basada en la \gls{armcortexm}. Los STM32 son reconocidos por su alto rendimiento, eficiencia energ{\'e}tica y amplia variedad de perif{\'e}ricos. Se utilizan en sistemas embebidos e \gls{iot} para tareas de procesamiento en \gls{tiemporeal}, gesti{\'o}n de sensores, comunicaciones industriales y aplicaciones cr{\'i}ticas en salud, manufactura, transporte y dom{\'o}tica.},
		type=main,
		symbol={🔌},
		user1={Microcontroladores de alto rendimiento},
		user2={Basados en ARM Cortex-M},
		user3={Uso en IoT y sistemas embebidos},
		see={armcortexm,microcontrolador,iot}
	}
	
	\newglossaryentry{subdominio}{
		name={Subdominio},
		text={subdominio},
		sort={subdominio},
		plural={subdominios},
		description={Subclasificaci{\'o}n de un nombre de dominio, definida con fines de gesti{\'o}n, organizaci{\'o}n o seguridad. Un subdominio aparece como una palabra previa al dominio principal, separada por un punto (por ejemplo: \texttt{api.ejemplo.com}). Facilita la segmentaci{\'o}n de servicios, control de accesos y escalabilidad en sistemas web e \gls{iot}.},
		type=main,
		symbol={🌐},
		user1={Clasificaci{\'o}n de nombres de dominio},
		user2={Organizaci{\'o}n y seguridad en servicios web},
		user3={Uso en APIs y despliegues IoT},
		see={dominio,plataformaweb}
	}
	
	\newglossaryentry{tabla}{
		name={Tabla},
		text={tabla},
		plural={tablas},
		first={tabla},
		firstplural={tablas},
		see={basedatos,dato,informacion},
		sort={tabla},
		description={Estructura de \glspl{dato} organizada en filas y columnas utilizada para almacenar, consultar y compartir informaci{\'o}n de manera sistem{\'a}tica. En \gls{sibiot}, las \emph{tablas} suelen materializar resultados consolidados, KPIs o salidas anal{\'i}ticas que son consumidas por aplicaciones, visualizaciones o sistemas empresariales}
	}
	
	\newglossaryentry{tareaconcurrente}{
		name={Tarea concurrente},
		text={tarea concurrente},
		sort={tarea concurrente},
		plural={tareas concurrentes},
		first={tarea concurrente},
		firstplural={tareas concurrentes},
		description={Unidad de ejecuci{\'o}n l{\'o}gica que se desarrolla de manera simult{\'a}nea o intercalada con otras tareas dentro de un sistema computacional. En sistemas embebidos y \glspl{sibiot}, una \emph{tarea concurrente} permite gestionar m{\'u}ltiples procesos, como temporizaci{\'o}n, adquisici{\'o}n de datos y comunicaci{\'o}n, sin bloquear la ejecuci{\'o}n global del sistema.},
		type=main,
		symbol={⏱},
		user1={Ejecuci{\'o}n simult{\'a}nea},
		user2={Multitarea},
		user3={Procesos independientes},
		see={concurrente,tareasincronizada}
	}
	
	\newglossaryentry{tareasincronizada}{
		name={Tarea sincronizada},
		text={tarea sincronizada},
		sort={tarea sincronizada},
		plural={tareas sincronizadas},
		first={tarea sincronizada},
		firstplural={tareas sincronizadas},
		description={Tarea cuya ejecuci{\'o}n est{\'a} coordinada con otros procesos mediante eventos, temporizadores o condiciones espec{\'i}ficas. En sistemas SIbIoT, una \emph{tarea sincronizada} permite asegurar que las acciones se ejecuten en momentos precisos o tras la ocurrencia de eventos relevantes.},
		type=main,
		symbol={⏲},
		user1={Ejecuci{\'o}n coordinada},
		user2={Control por eventos},
		user3={Determinismo},
		see={sincronizada,tareaconcurrente}
	}
	
	\newglossaryentry{taskboard}{
		name={Task Board},
		text={\textit{Task Board}},
		sort={task board},
		plural={\textit{Task Boards}},
		first={\textit{Task Board}},
		firstplural={\textit{Task Boards}},
		description={Tablero visual utilizado para representar el estado de las tareas del Sprint, generalmente organizado en columnas como Pendiente, En progreso y Terminado. Facilita la transparencia y la autoorganizaci{\'o}n del equipo.},
		type=main,
		symbol={📌},
		user1={Tablero visual de tareas},
		user2={Gesti{\'o}n del flujo de trabajo},
		user3={Soporte para autoorganizaci{\'o}n},
		see={sprintbacklog,sprint,agil}
	}
	
	\newglossaryentry{tcp}{
		name={TCP},
		text={TCP},
		sort={tcp},
		plural={TCP},
		first={protocolo de control de transmisi{\'o}n (\textit{Transmission Control Protocol}, TCP)},
		firstplural={TCP},
		description={\emph{TCP} es un \gls{protocolocomunicacion} de transporte orientado a conexi{\'o}n que garantiza la entrega fiable, ordenada y sin duplicados de los datos entre aplicaciones que se comunican sobre redes \gls{ip}. Implementa mecanismos como el establecimiento de conexi{\'o}n (handshake de tres pasos), control de flujo, control de congesti{\'o}n y retransmisi{\'o}n autom{\'a}tica. En sistemas \gls{iot}, TCP se utiliza cuando se requiere confiabilidad estricta, integridad de datos y consistencia, como en \glspl{servicioweb}, \glspl{comunicacion} con \glspl{servidor} en la \gls{nube}, configuraci{\'o}n remota de \glspl{dispositivo} o transferencia segura de registros cr{\'i}ticos, aunque su \gls{sobrecarga} lo hace menos adecuado para \glspl{dispositivo} de muy baja potencia o \glspl{red} inestables.},
		type=main,
		symbol={📦},
		user1={Transporte orientado a conexi{\'o}n y fiable},
		user2={Mecanismos de control de flujo y congesti{\'o}n},
		user3={Uso en aplicaciones IoT que requieren integridad},
		see={udp,protocolocomunicacion,ip,protocoloseguridad}
	}
	
	\newglossaryentry{tdd}{
		name={TDD},
		text={TDD},
		first={desarrollo guiado por pruebas (\textit{Test-Driven Development}, TDD)},
		sort={tdd},
		description={\Gls{metodologiadesarrollo} de software en la que las pruebas se escriben antes del c{\'o}digo funcional. El ciclo de TDD implica crear una prueba que falle, implementar el c{\'o}digo necesario para superarla y luego refactorizar el c{\'o}digo manteniendo la funcionalidad. Esta pr{\'a}ctica fomenta c{\'o}digo modular, verificable y de alta \gls{calidad}. Es especialmente efectiva en sistemas cr{\'i}ticos, incluyendo aplicaciones \gls{iot}, donde la fiabilidad y validaci{\'o}n temprana son esenciales.},
		type=main,
		symbol={🧪},
		user1={Pruebas antes del c{\'o}digo},
		user2={C{\'i}clo escribir--probar--refactorizar},
		user3={Alta calidad y fiabilidad en IoT},
		see={basedatos,calidad,iot}
	}
	
	\newglossaryentry{tddm4iots}{
		name={TDDM4IoTS},
		text={TDDM4IoTS},
		sort={tddm4iots},
		plural={TDDM4IoTS},
		first={Metodolog{\'i}a de Desarrollo Guiado por Pruebas para Sistemas Basados en Internet de las Cosas (\textit{Test-Driven Development Methodology for IoT-Based Systems}, TDDM4IoTS)},
		firstplural={TDDM4IoTS},
		description={Metodolog{\'\i}a de desarrollo para sistemas \gls{iot} basada en un enfoque dirigido por pruebas y apoyado en modelado, orientada a reducir la divergencia entre lo dise{\~n}ado y lo implementado. TDDM4IoTS estructura el trabajo mediante artefactos de modelado y su refinamiento hacia implementaciones, integrando \glspl{pruebaautomatizada} desde etapas tempranas para validar comportamiento, restricciones y calidad de manera continua. En \glspl{sibiot}, su aplicaci{\'o}n resulta pertinente por la presencia de \glspl{dispositivoembebido}, comunicaciones heterog{\'e}neas y despliegues distribuidos, donde la verificaci{\'o}n temprana y la evoluci{\'o}n controlada del sistema son determinantes para mantener consistencia y trazabilidad.},
		type=main,
		symbol={🧪},
		user1={Desarrollo guiado por pruebas},
		user2={Modelado y refinamiento},
		user3={Trazabilidad dise{\~n}o--implementaci{\'o}n},
		see={metodologia,pruebaautomatizada,modelo,transformacioncontrolada,verificacion,evolucioncontrolada}
	}
	
	\newglossaryentry{tecnica}{
		name={T{\'e}cnica},
		text={t{\'e}cnica},
		sort={tecnica},
		plural={t{\'e}cnicas},
		description={Procedimiento espec{\'i}fico y pr{\'a}ctico que implementa un paso dentro de un m{\'e}todo. Las t{\'e}cnicas son herramientas operativas que pueden incluir algoritmos, diagramas, procesos de medici{\'o}n o mecanismos de control. En \gls{iot}, se aplican t{\'e}cnicas como filtrado de datos, compresi{\'o}n, autenticaci{\'o}n o direccionamiento de paquetes.},
		type=main,
		symbol={🛠️},
		user1={Procedimiento espec{\'i}fico de un m{\'e}todo},
		user2={Incluye algoritmos y procesos},
		user3={Uso en control y optimizaci{\'o}n IoT},
		see={metodo,algoritmo,iot}
	}
	
	\newglossaryentry{tecnicacorreccionerrores}{
		name={T{\'e}cnica de correcci{\'o}n de errores},
		text={t{\'e}cnica de correcci{\'o}n de errores},
		sort={tecnicacorreccionerrores},
		plural={t{\'e}cnicas de correcci{\'o}n de errores},
		description={Conjunto de m{\'e}todos de comunicaci{\'o}n digital que detectan y corrigen errores introducidos durante la transmisi{\'o}n de datos. En \gls{iot}, se utilizan c{\'o}digos de detecci{\'o}n (CRC) y de correcci{\'o}n (Hamming, Reed-Solomon, LDPC) para mejorar la \gls{fiabilidad} de datos en entornos con ruido, interferencias o limitaciones energ{\'e}ticas.},
		type=main,
		symbol={✔️},
		user1={Detecci{\'o}n y correcci{\'o}n en transmisi{\'o}n},
		user2={C{\'o}digos CRC, Hamming, Reed-Solomon, LDPC},
		user3={Garantiza fiabilidad en IoT},
		see={fiabilidad,seguridad,comunicacioninalambrica}
	}
	
	\newglossaryentry{tecnicamitigacion}{
		name={T{\'e}cnica de mitigaci{\'o}n},
		text={t{\'e}cnica de mitigaci{\'o}n},
		sort={tecnicamitigacion},
		plural={t{\'e}cnicas de mitigaci{\'o}n},
		description={Conjunto de procedimientos espec{\'i}ficos orientados a reducir la exposici{\'o}n a riesgos, vulnerabilidades o fallos en sistemas tecnol{\'o}gicos. En \gls{iot}, las t{\'e}cnicas de mitigaci{\'o}n incluyen \gls{cifrado} de \glspl{dato}, autenticaci{\'o}n robusta, aislamiento de redes, control de acceso, redundancia de hardware, balanceo de carga y \gls{gestionenergetica} \gls{eficiente}, con el fin de preservar \gls{disponibilidad}, \gls{integridad} y \gls{seguridad}.},
		type=main,
		symbol={🛡️},
		user1={Reducci{\'o}n de riesgos y vulnerabilidades},
		user2={Incluye cifrado, redundancia y control de acceso},
		user3={Preserva seguridad y resiliencia IoT},
		see={seguridad,fiabilidad,disponibilidad}
	}
	
	\newglossaryentry{tecnicasdemodulacion}{
		name={T{\'e}cnicas de modulaci{\'o}n},
		text={t{\'e}cnicas de modulaci{\'o}n},
		sort={tecnicasdemodulacion},
		plural={t{\'e}cnicas de modulaci{\'o}n},
		description={Procesos que modifican una se{\~n}al portadora para transmitir informaci{\'o}n. En \gls{iot}, se emplean modulaciones como FSK, PSK, QAM y LoRa (modulaci{\'o}n por chirp de espectro ensanchado), seg{\'u}n el alcance, el consumo energ{\'e}tico y los requisitos del canal.},
		type=main,
		symbol={📶},
		user1={Modificaci{\'o}n de se{\~n}al portadora},
		user2={FSK, PSK, QAM, LoRa},
		user3={Adaptaci{\'o}n a entorno IoT},
		see={comunicacioninalambrica,iot}
	}
	
	\newglossaryentry{tecnologia}{
		name={Tecnolog{\'i}a},
		text={tecnolog{\'i}a},
		sort={tecnologia},
		plural={tecnolog{\'i}as},
		description={Conjunto de conocimientos, habilidades, \glspl{tecnica} y procesos para dise{\~n}ar, crear y aplicar herramientas, sistemas o dispositivos que resuelvan problemas o mejoren condiciones. En \gls{iot}, la tecnolog{\'i}a comprende tanto el \gls{hardware} (sensores, microcontroladores) como el software (plataformas, algoritmos, redes de comunicaci{\'o}n) que habilitan la interconexi{\'o}n y automatizaci{\'o}n de sistema.},
		type=main,
		symbol={⚙️},
		user1={Conocimientos y habilidades aplicadas},
		user2={Incluye hardware y software},
		user3={Base de sistemas IoT},
		see={tecnica,hardware,software,iot}
	}
	
	\newglossaryentry{tecnologiainalambrica}{
		name={Tecnolog{\'i}a inal{\'a}mbrica},
		text={tecnolog{\'i}a inal{\'a}mbrica},
		plural={tecnolog{\'i}as inal{\'a}mbricas},
		sort={tecnologiainalambrica},
		first={tecnolog{\'i}a inal{\'a}mbrica},
		firstplural={tecnolog{\'i}as inal{\'a}mbricas},
		description={Conjunto de medios de comunicaci{\'o}n que permiten el intercambio de informaci{\'o}n entre dispositivos a distancia sin cables f{\'i}sicos, utilizando ondas electromagn{\'e}ticas. Incluye Wi-Fi, Bluetooth, Zigbee, LoRa, NB-IoT, entre otras.},
		type=main,
		symbol={📡},
		see={wifi,zigbee,bluetooth,lora,nbiot}
	}
	
	\newglossaryentry{tecnologiaiot}{
		name={Tecnolog{\'i}a \glsentrytext{iot}},
		text={tecnolog{\'i}a \glsentrytext{iot}},
		sort={tecnologia iot},
		plural={tecnolog{\'i}as \glsentrytext{iot}},
		first={tecnolog{\'i}a \glsentrytext{iot}},
		firstplural={tecnolog{\'i}as \glsentrytext{iot}},
		description={Conjunto de dispositivos, arquitecturas, protocolos, plataformas y \glspl{servicio} que habilitan la interconexi{\'o}n de objetos f{\'i}sicos mediante redes digitales para \gls{captura}, procesamiento, transmisi{\'o}n y an{\'a}lisis de datos del \gls{entorno} en \gls{tiemporeal}. La tecnolog{\'i}a \gls{iot} integra \glspl{sensor}, \glspl{actuador}, \glspl{microcontrolador}, \glspl{gateway}, \gls{nube}, \gls{procesamientodatos} y \glspl{modelopredictivo}, permitiendo automatizaci{\'o}n inteligente y \gls{tomadecisiones} en sectores como salud, agricultura, manufactura y transporte, con criterios de \gls{seguridad}, \gls{interoperabilidad}, \gls{escalabilidad} y \gls{sostenibilidad}.},
		type=main,
		symbol={🌐},
		user1={Conjunto de componentes y protocolos para IoT},
		user2={Integraci{\'o}n de captura, an{\'a}lisis y acci{\'o}n},
		user3={Automatizaci{\'o}n inteligente y soporte a decisiones},
		see={iot,sensor,actuador,microcontrolador,gateway,nube,procesamientodatos,interoperabilidad,escalabilidad,seguridad,sostenibilidad,sibiot}
	}
	
	\newglossaryentry{tecnologiasatelital}{
		name={Tecnolog{\'i}a satelital},
		text={tecnolog{\'i}a satelital},
		sort={tecnologia satelital},
		plural={tecnolog{\'i}as satelitales},
		first={tecnolog{\'i}a satelital},
		firstplural={tecnolog{\'i}as satelitales},
		description={La \emph{tecnolog{\'i}a satelital} abarca el conjunto de infraestructuras, dispositivos y sistemas que utilizan sat{\'e}lites artificiales para la \gls{comunicacion}, la observaci{\'o}n de la Tierra, la navegaci{\'o}n global y la monitorizaci{\'o}n ambiental. Incluye \glspl{gnss} como \gls{gps}, Galileo, GLONASS y BeiDou, as{\'i} como sat{\'e}lites de comunicaciones y teledetecci{\'o}n. En el contexto de \glspl{sibiot}, estas tecnolog{\'i}as permiten la \gls{conectividad} en zonas remotas, la integraci{\'o}n de datos geoespaciales, la gesti{\'o}n de recursos naturales y el soporte a \glspl{transporteinteligente} y \glspl{ciudadinteligente}.},
		type=main,
		symbol={🛰️},
		user1={Sat{\'e}lites de comunicaci{\'o}n y observaci{\'o}n},
		user2={Sistemas de navegaci{\'o}n global GNSS},
		user3={Aplicaciones IoT en movilidad y monitorizaci{\'o}n ambiental},
		see={gnss,gps,conectividad,transporteinteligente,sibiot}
	}
	
	\newglossaryentry{telemetria}{
		name={Telemetr{\'i}a},
		text={telemetr{\'i}a},
		sort={telemetria},
		plural={telemetr{\'i}as},
		first={telemetr{\'i}a},
		firstplural={telemetr{\'i}as},
		description={La \emph{telemetr{\'i}a} es el proceso mediante el cual se obtienen mediciones de un \gls{entorno} o \glspl{dispositivo} de forma remota y autom{\'a}tica, para luego transmitirlas a un sistema central de \gls{monitorizacion} o an{\'a}lisis. Implica la adquisici{\'o}n continua de \glspl{dato}, su codificaci{\'o}n y env{\'i}o a trav{\'e}s de \glspl{protocolocomunicacion} como MQTT, CoAP, UDP o HTTP. En sistemas \gls{iot}, la telemetr{\'i}a es fundamental para habilitar la monitorizaci{\'o} en \gls{tiemporeal}, la detecci{\'o}n de \glspl{anomalia}, el mantenimiento predictivo, la toma automatizada de decisiones y la gesti{\'o}n remota de \glspl{dispositivointeligente}.},
		type=main,
		symbol={📡},
		user1={Medici{\'o}n remota y autom{\'a}tica de datos},
		user2={Base de la monitorizaci{\'o}n en tiempo real},
		user3={Esencial para mantenimiento predictivo y control IoT},
		see={sensor,m2m,protocolocomunicacion,monitorizacion,sibiot}
	}
	
	\newglossaryentry{temporizacion}{
		name={Temporizaci{\'o}n},
		text={temporizaci{\'o}n},
		sort={temporizacion},
		plural={temporizaciones},
		first={temporizaci{\'o}n},
		firstplural={temporizaciones},
		description={Tratamiento expl{\'\i}cito de requisitos de tiempo en el dise{\~n}o y operaci{\'o}n de un sistema (plazos, orden de eventos, sincronizaci{\'o}n, latencia y variaci{\'o}n temporal). En \glspl{cps}, la temporizaci{\'o}n forma parte de la correcci{\'o}n del sistema, porque una decisi{\'o}n puede ser correcta en l{\'o}gica pero inaceptable si llega fuera de la ventana temporal requerida por el proceso f{\'\i}sico.},
		type=main,
		symbol={⏱},
		user1={Plazos y ventanas},
		user2={Sincronizaci{\'o}n},
		user3={Orden de eventos},
		see={tiemporeal,latencia,control}
	}
	
	\newglossaryentry{temporizador}{
		name={Temporizador},
		text={temporizador},
		sort={temporizador},
		plural={temporizadores},
		first={temporizador},
		firstplural={temporizadores},
		description={M{\'o}dulo de hardware interno del microcontrolador encargado de medir intervalos de tiempo, generar retardos precisos y sincronizar eventos. Un \emph{temporizador} es fundamental para la implementaci{\'o}n de tareas peri{\'o}dicas, control de tiempos y mecanismos de concurrencia en sistemas embebidos e IoT.},
		type=main,
		symbol={⏱},
		user1={Control temporal},
		user2={Eventos peri{\'o}dicos},
		user3={Sincronizaci{\'o}n},
		see={tareaconcurrente,tareasincronizada}
	}
	
	\newglossaryentry{tendencia}{
		name={Tendencia},
		text={tendencia},
		sort={tendencia},
		plural={tendencias},
		first={tendencia},
		firstplural={tendencias},
		description={Una tendencia refleja la evoluci{\'o}n o desplazamiento progresivo de un fen{\'o}meno medido o analizado. En los \glspl{sibiot}, las tendencias se identifican a partir de la recopilaci{\'o}n y el an{\'a}lisis de datos en \gls{tiemporeal}, permitiendo detectar patrones recurrentes, predecir eventos y anticipar decisiones estrat{\'e}gicas. El estudio de tendencias es esencial en procesos de \gls{analisisdescriptivo} y \gls{analisispredictivo}, apoyando la mejora continua de los sistemas inteligentes.},
		type=\glsdefaulttype,
		symbol={📈},
		user1={Trend},
		user2={Evoluci{\'o}n o patr{\'o}n de cambio en los datos},
		user3={Identificaci{\'o}n y predicci{\'o}n de comportamientos en SIbIoT},
		see={comportamiento,analisisdescriptivo,analisispredictivo,sibiot,tiemporeal}
	}
	
	\newglossaryentry{tension}
	{
		name={Tensi{\'o}n},
		text={tensi{\'o}n},
		sort={tension},
		description={Magnitud f{\'i}sica que representa la fuerza interna distribuida a lo largo de un cuerpo o elemento estructural cuando se somete a una carga que tiende a estirarlo. 
			En el campo de la ingenier{\'i}a y los \glspl{sibiot}, la tensi{\'o}n se mide mediante \glspl{sensor} de deformaci{\'o}n o galgas extensiom{\'e}tricas integradas en \glspl{celda} de carga, permitiendo la supervisi{\'o}n de esfuerzos en estructuras y maquinaria. 
			En sistemas el{\'e}ctricos, el t{\'e}rmino tensi{\'o}n tambi{\'e}n se utiliza para describir la diferencia de potencial entre dos puntos, expresada en voltios (V), siendo una variable clave en la alimentaci{\'o}n y control de dispositivos IoT.
			Su monitorizaci{\'o}n continua garantiza la estabilidad operativa, la seguridad del sistema y la prevenci{\'o}n de fallos cr{\'i}ticos.},
		plural={tensiones},
		symbol={\ensuremath{\sigma}},
		first={Tensi{\'o}n (\ensuremath{\sigma})},
		see=[Relacionado con:]{fuerza,esfuerzo,deformacion,presion}
	}
	
	\newglossaryentry{testabilidad}{
		name={Testabilidad},
		text={testabilidad},
		sort={testabilidad},
		plural={testabilidades},
		first={testabilidad},
		firstplural={testabilidades},
		description={Grado en que un sistema o componente de software permite ser probado de manera efectiva para verificar su correcto funcionamiento, detectar defectos y validar el cumplimiento de requisitos. En los \glspl{sibiot}, la testabilidad es un atributo clave de calidad, ya que facilita la validaci{\'o}n de sistemas distribuidos, la verificaci{\'o}n de interacciones entre dispositivos y servicios, y la detecci{\'o}n temprana de fallos en entornos heterog{\'e}neos y din{\'a}micos.},
		type=main,
		symbol={🧪},
		user1={Facilidad de prueba},
		user2={Calidad del software},
		user3={Verificaci{\'o}n del sistema},
		see={calidadsoftware,refactorizacion,disenosistemas}
	}
	
	\newglossaryentry{tft}{
		name={TFT},
		text={TFT},
		sort={tft},
		plural={TFT},
		first={transistor de pel{\'i}cula fina (\textit{Thin Film Transistor}, TFT)},
		firstplural={TFT},
		description={Un \emph{TFT} es una tecnolog{\'i}a de visualizaci{\'o}n que forma parte de las pantallas \gls{lcd}, en la cual cada p{\'i}xel es controlado por un transistor de pel{\'i}cula fina. Esto permite mejorar la nitidez, el tiempo de respuesta y la reproducci{\'o}n de color en comparaci{\'o}n con LCD tradicionales. En sistemas \gls{iot}, las pantallas TFT son empleadas en interfaces gr{\'a}ficas interactivas, paneles de control industrial y dispositivos port{\'a}tiles que requieren resoluciones medias o altas.},
		type=main,
		symbol={🖥️},
		user1={Pantallas LCD con transistores de pel{\'i}cula fina},
		user2={Mayor nitidez y tiempo de respuesta},
		user3={Uso en paneles de control e interfaces IoT},
		see={lcd,oled,interfaz,dispositivo}
	}
	
	\newglossaryentry{throughput}{
		name={Throughput},
		text={throughput},
		sort={throughput},
		plural={throughputs},
		first={throughput},
		firstplural={throughputs},
		description={V{\'e}ase \gls{rendimiento} por su sin{\'o}nimo.},
		type=main,
		symbol={🚀},
		user1={Tasa de transferencia},
		user2={Alias en ingl{\'e}s de rendimiento},
		user3={Terminolog{\'i}a de desempe{\~n}o},
		see={rendimiento}
	}
	
	\newglossaryentry{tic}{
		name={Tecnolog{\'i}as de la Informaci{\'o}n y la Comunicaci{\'o}n},
		text={TIC},
		sort={tic},
		plural={TIC},
		first={Tecnolog{\'i}as de la Informaci{\'o}n y la Comunicaci{\'o}n (\textit{Information and Comunication Technologies}, TIC)},
		firstplural={TIC},
		description={Conjunto de recursos, infraestructuras, dispositivos, servicios y metodolog{\'i}as orientados a la adquisici{\'o}n, almacenamiento, procesamiento, transmisi{\'o}n y visualizaci{\'o}n de informaci{\'o}n en formatos digitales. Las TIC abarcan redes de datos, c{\'o}mputo, software, plataformas en la nube, servicios web y medios de comunicaci{\'o}n sincr{\'o}nica y asincr{\'o}nica. En educaci{\'o}n y gesti{\'o}n organizacional, habilitan procesos de automatizaci{\'o}n, colaboraci{\'o}n, anal{\'i}tica y toma de decisiones. En el contexto de los \glspl{sibiot} y de los \glspl{sistemainteligente}, las TIC proporcionan la columna vertebral para la conectividad, la interoperabilidad y la seguridad extremo a extremo, facilitando integraci{\'o}n de datos, trazabilidad y servicios digitales escalables.},
		type=main,
		user1={Infraestructura digital y servicios de red},
		user2={Procesos de gesti{\'o}n de informaci{\'o}n y comunicaci{\'o}n},
		user3={Habilitadores de automatizaci{\'o}n, anal{\'i}tica y colaboraci{\'o}n},
		user4={Soporte a \textit{e{-}learning}, gobierno digital y transformaci{\'o}n digital},
		see={ia,aprendizajeautomatico,sistemainteligente,sibiot},
		symbol={\textsc{tic}}
	}
	
	\newglossaryentry{tidb}{
		name={TiDB},
		text={TiDB},
		sort={tidb},
		plural={TiDB},
		first={TiDB},
		firstplural={TiDB},
		description={\emph{TiDB} es un sistema de gesti{\'o}n de \glspl{basedatosrelacional} distribuido, de c{\'o}digo abierto, compatible con MySQL y orientado a soportar cargas h{\'i}bridas de procesamiento transaccional (OLTP) y anal{\'i}tico (\gls{olap}), siguiendo el paradigma \gls{htap}. Dise{\~n}ado por PingCAP, TiDB ofrece \gls{escalabilidad} horizontal autom{\'a}tica, replicaci{\'o}n distribuida, tolerancia a fallos y consistencia fuerte. En entornos \gls{sibiot}, \emph{TiDB} permite manejar flujos masivos de \glspl{dato} heterog{\'e}neos en tiempo real, integrando transacciones cr{\'i}ticas con an{\'a}lisis complejos sobre los mismos datos.},
		type=main,
		symbol={🐯},
		user1={Base de datos distribuida de c{\'o}digo abierto},
		user2={Soporte a cargas HTAP},
		user3={Compatibilidad con MySQL},
		see={htap,basedatosrelacional,bigdata,sibiot}
	}
	
	\newglossaryentry{tiemporeal}{
		name={Tiempo real},
		text={tiempo real},
		sort={tiempo real},
		plural={sistemas de tiempo real},
		first={tiempo real},
		firstplural={sistemas de tiempo real},
		description={Propiedad de un \gls{sistema} de c{\'o}mputo o control que exige que sus operaciones/respuestas ocurran dentro de l{\'i}mites temporales estrictos y predecibles. Un \emph{sistema de tiempo real} debe garantizar no solo la correcci{\'o}n l{\'o}gica del resultado, sino tambi{\'e}n su entrega dentro de un plazo espec{\'i}fico (deadlines). Se aplica en control industrial, sistemas embebidos, \glspl{vehiculoautonomo}, monitorizaci{\'o}n m{\'e}dica y \gls{iot} cr{\'i}tico.},
		type=main,
		symbol={⏱️},
		user1={Cumplimiento de plazos deterministas},
		user2={Correcci{\'o}n funcional y temporal},
		user3={Aplicaciones cr{\'i}ticas de control y IoT},
		see={latencia,monitorizacion,observabilidad,eficienciaoperativa,iot}
	}
	
	\newglossaryentry{tiemporetardo}{
		name={Tiempo de retardo},
		text={tiempo de retardo},
		sort={tiempo de retardo},
		plural={tiempos de retardo},
		first={tiempo de retardo},
		firstplural={tiempos de retardo},
		description={Intervalo de tiempo transcurrido entre la emisi{\'o}n de una se{\~n}al, instrucci{\'o}n o evento, y la respuesta efectiva del \gls{sistema} o del dispositivo receptor. El \emph{tiempo de retardo} puede deberse a factores de procesamiento, comunicaci{\'o}n, propagaci{\'o}n de se{\~n}ales o sincronizaci{\'o}n entre componentes. En el contexto de los \glspl{sibiot}, un \emph{tiempo de retardo configurable} permite ajustar dicho intervalo de forma deliberada para sincronizar acciones de \glspl{actuador}, establecer tolerancias en protocolos de comunicaci{\'o}n o adaptar el comportamiento del \gls{sistema} a distintos \glspl{entorno}, y necesidades de \gls{monitorizacion} y \gls{control}. Su gesti{\'o}n es fundamental para garantizar \gls{fiabilidad}, \gls{rendimiento} y \gls{disponibilidad} en aplicaciones de \gls{tiemporeal}.},
		type=main,
		symbol={⏱️},
		user1={Intervalo entre emisi{\'o}n y respuesta},
		user2={Ajuste configurable en sistemas IoT},
		user3={Par{\'a}metro de rendimiento en SIbIoT},
		see={latencia,rendimiento,fiabilidad,disponibilidad,sibiot}
	}
	
	\newglossaryentry{tls}{
		name={TLS},
		text={TLS},
		sort={tls},
		plural={protocolos TLS},
		first={seguridad de la capa de transporte (\textit{Transport Layer Security}, TLS)},
		firstplural={protocolos TLS},
		description={Protocolo criptogr{\'a}fico que proporciona comunicaciones seguras sobre una \gls{red}. \emph{TLS} garantiza \gls{confidencialidad}, \gls{integridad} y \gls{autenticacion} de los datos intercambiados entre \glspl{nodo} o \glspl{servicio} mediante cifrado y certificados digitales. Es ampliamente utilizado en \glspl{aplicacionweb} (p.\,ej., \gls{https}) y en sistemas \gls{iot} que requieren mecanismos robustos de \gls{seguridad} en la \gls{transmision} de \gls{informacion}.},
		type=main,
		symbol={🔒},
		user1={Cifrado y autenticaci{\'o}n de extremo a extremo},
		user2={Protege integridad y confidencialidad de datos},
		user3={Base de \gls{https} y seguridad en IoT},
		see={https,seguridad,confidencialidad,integridad,autenticacion,iot}
	}
	
	\newglossaryentry{token}{
		name={Token},
		text={token},
		sort={token},
		plural={tokens},
		first={token},
		firstplural={tokens},
		description={Credencial digital utilizada para autenticar, autorizar y validar identidades o permisos dentro de un sistema de informaci{\'o}n distribuido. Los \emph{tokens} encapsulan informaci{\'o}n de identidad, roles, permisos y tiempo de validez, y suelen transmitirse de forma segura en encabezados HTTP o mensajes cifrados. En sistemas IoT y \glspl{sibiot}, los tokens permiten la autenticaci{\'o}n ligera y escalable de dispositivos, usuarios y servicios, facilitando el control de acceso a \glspl{api}, servicios WebSocket y recursos en la nube sin necesidad de compartir credenciales sensibles. Su uso es fundamental en arquitecturas desacopladas y orientadas a servicios},
		type=main,
		symbol={🎫},
		see={autenticacion,autorizacion,seguridad,api,restful,desacoplamiento,sibiot}
	}
	
	\newglossaryentry{toleranciafallos}{
		name={Tolerancia a fallos},
		text={tolerancia a fallos},
		sort={tolerancia a fallos},
		plural={tolerancias a fallos},
		first={tolerancia a fallos},
		firstplural={tolerancias a fallos},
		description={Capacidad de un \gls{sistema}, componente o infraestructura para continuar operando correctamente, aunque uno o varios de sus elementos fallen o se comporten de manera inesperada. La \emph{tolerancia a fallos} se logra mediante t{\'e}cnicas como redundancia, replicaci{\'o}n, detecci{\'o}n y correcci{\'o}n de errores, y mecanismos de recuperaci{\'o}n autom{\'a}tica. En el contexto de los \glspl{sibiot}, esta propiedad garantiza la \gls{disponibilidad} y la \gls{fiabilidad} en escenarios cr{\'i}ticos, permitiendo que los servicios contin{\'u}en activos incluso ante fallos en \glspl{sensor}, \glspl{nodo} o canales de comunicaci{\'o}n, favoreciendo adem{\'a}s la \gls{resiliencia} y la confianza del usuario final.},
		type=main,
		symbol={🛠️},
		user1={Continuidad operativa ante fallos},
		user2={Redundancia y recuperaci{\'o}n autom{\'a}tica},
		user3={Mecanismo de confiabilidad en SIbIoT},
		see={fiabilidad,resiliencia,disponibilidad,nodo,sibiot}
	}
	
	% ===================== Toma de decisiones =====================
	\newglossaryentry{tomadecisiones}{
		name={Toma de decisiones},
		text={toma de decisiones},
		sort={toma de decisiones},
		plural={procesos de toma de decisiones},
		first={toma de decisiones},
		firstplural={procesos de toma de decisiones},
		description={La \emph{toma de decisiones} es el proceso mediante el cual un agente, humano o artificial, selecciona una acci{\'o}n o alternativa entre varias posibles, con base en objetivos, restricciones y criterios definidos. En sistemas computacionales, puede estar respaldada por algoritmos, \glspl{modelopredictivo}, reglas heur{\'i}sticas o \gls{ia}, con aplicaciones en automatizaci{\'o}n, diagn{\'o}stico, planificaci{\'o}n, rob{\'o}tica y gesti{\'o}n estrat{\'e}gica. En el contexto de \glspl{sibiot}, la toma de decisiones es fundamental para optimizar recursos, responder en tiempo real a datos de \glspl{sensor} y garantizar \gls{efectividad}.},
		type=main,
		symbol={🎯},
		user1={Selecci{\'o}n de acciones entre alternativas},
		user2={Soporte por algoritmos e IA},
		user3={Aplicaciones en IoT y sistemas inteligentes},
		see={ia,modelopredictivo,efectividad,sibiot}
	}
	
	\newglossaryentry{tomadecisionesestrategicas}{
		name={Toma de decisiones estrat{\'e}gicas},
		text={toma de decisiones estrat{\'e}gicas},
		sort={toma decisiones estrategicas},
		plural={procesos de toma de decisiones estrat{\'e}gicas},
		first={toma de decisiones estrat{\'e}gicas},
		firstplural={procesos de toma de decisiones estrat{\'e}gicas},
		description={Proceso mediante el cual se seleccionan cursos de acci{\'o}n que afectan a largo plazo el desempe{\~n}o de una organizaci{\'o}n. En los \glspl{sibiot}, este proceso se apoya en informaci{\'o}n derivada del an{\'a}lisis de datos operativos y contextuales.},
		type=main,
		symbol={📊},
		user1={Estrategia organizacional},
		user2={An{\'a}lisis de informaci{\'o}n},
		user3={Planeaci{\'o}n},
		see={sistemainformacioncorporativo,optimizacionprocesos}
	}
	
	\newglossaryentry{tomadecisionesinformadas}{
		name={toma de decisiones informadas},
		text={toma de decisiones informadas},
		sort={tomadecisionesinformadas},
		plural={tomas de decisiones informadas},
		first={toma de decisiones informadas},
		firstplural={tomas de decisiones informadas},
		description={La toma de decisiones informadas implica seleccionar la mejor alternativa posible a partir de evidencias concretas y resultados anal{\'i}ticos. En los \glspl{sibiot}, este proceso se apoya en la recopilaci{\'o}n y el an{\'a}lisis de datos en \gls{tiemporeal}, lo que permite responder con precisi{\'o}n ante condiciones cambiantes. Su aplicaci{\'o}n abarca {\'a}mbitos como la \gls{gestionempresarial}, la \gls{planificacion}, el \gls{mantenimiento} predictivo y la \gls{optimizacion} de recursos en \glspl{sistemainteligente}.},
		type=\glsdefaulttype,
		symbol={🧭},
		user1={Informed decision-making},
		user2={Selecci{\'o}n basada en evidencia y an{\'a}lisis de datos},
		user3={Proceso de elecci{\'o}n inteligente en SIbIoT},
		see={tomadecisiones,analisisdescriptivo,analisispredictivo,analisisprescriptivo,sibiot}
	}
	
	\newglossaryentry{topologia}{
		name={topolog{\'i}a},
		text={topolog{\'i}a},
		sort={topologia},
		plural={topolog{\'i}as},
		first={topolog{\'i}a},
		firstplural={topolog{\'i}as},
		type=main,
		description={Estructura l{\'o}gica y organizativa que define la forma en que los nodos, enlaces y \glspl{dispositivo} se interconectan dentro de una \gls{red}. La topolog{\'i}a influye en el flujo de datos, la eficiencia del \gls{enrutamiento}, la tolerancia a fallos y el desempe{\~n}o general del sistema. En el contexto de los \glspl{sibiot}, las topolog{\'i}as determinan c{\'o}mo se distribuyen los \glspl{dispositivoiot}, c{\'o}mo intercambian informaci{\'o}n y c{\'o}mo se gestiona la escalabilidad y la resiliencia operativa},
		symbol={🕸️},
		see={red,protocolo,protocolocomunicacion,enrutamiento}
	}
	
	% ===================== Transacci{\'o}n =====================
	\newglossaryentry{transaccion}{
		name={Transacci{\'o}n},
		text={transacci{\'o}n},
		sort={transaccion},
		plural={transacciones},
		first={transacci{\'o}n},
		firstplural={transacciones},
		description={Una \emph{transacci{\'o}n} es una unidad de interacci{\'o}n con una estructura de datos compleja, compuesta por varias operaciones que deben ejecutarse en su totalidad o no ejecutarse en absoluto, siguiendo el principio de atomicidad. Las transacciones garantizan que el sistema permanezca en un estado consistente, aun ante fallos, y son fundamentales en \glspl{basedatos}, sistemas distribuidos y aplicaciones financieras. En \glspl{sibiot}, las transacciones aseguran la coherencia en el procesamiento de datos de \glspl{sensor} y en la integraci{\'o}n con \glspl{servicionube}.},
		type=main,
		symbol={🔄},
		user1={Ejecuci{\'o}n at{\'o}mica de operaciones},
		user2={Garant{\'i}a de consistencia y aislamiento},
		user3={Base en bases de datos y sistemas distribuidos},
		see={transaccional,basedatos,sibiot}
	}
	
	% ===================== Transaccional =====================
	\newglossaryentry{transaccional}{
		name={Transaccional},
		text={transaccional},
		sort={transaccional},
		plural={transaccionales},
		first={transaccional},
		firstplural={transaccionales},
		description={El t{\'e}rmino \emph{transaccional} se refiere a todo aquello que pertenece o est{\'a} relacionado con una \gls{transaccion}. Se dice que un sistema o proceso es transaccional cuando soporta transacciones, garantizando las propiedades ACID (Atomicidad, Consistencia, Aislamiento y Durabilidad). En \glspl{sibiot}, un enfoque transaccional asegura que las operaciones de \gls{almacenamiento}, \gls{actualizacion} y \gls{comunicacion} de \glspl{dato} se realicen de manera fiable y segura.},
		type=\glsdefaulttype,
		symbol={✅},
		user1={Relaci{\'o}n directa con las transacciones},
		user2={Garant{\'i}a de propiedades ACID},
		user3={Aplicaci{\'o}n en IoT y sistemas distribuidos},
		see={transaccion,basedatos,sibiot}
	}
	
	\newglossaryentry{transductor}{
		name={Transductor},
		text={transductor},
		sort={transductor},
		plural={transductores},
		first={transductor},
		firstplural={transductores},
		description={Dispositivo que convierte una magnitud f{\'i}sica (p.\,ej., fuerza, temperatura, presi{\'o}n) en una se{\~n}al utilizable (normalmente el{\'e}ctrica) para medici{\'o}n, control o adquisici{\'o}n de datos.},
		type=main,
		symbol={🎛️},
		user1={Conversi{\'o}n f{\'i}sica--el{\'e}ctrica},
		user2={Instrumentaci{\'o}n},
		user3={Medici{\'o}n},
		see={celdacarga,monitorizacion}
	}
	
	\newglossaryentry{transformacioncontrolada}{
		name={Transformaci{\'o}n controlada},
		text={transformaci{\'o}n controlada},
		sort={transformacioncontrolada},
		plural={transformaciones controladas},
		first={transformaci{\'o}n controlada},
		firstplural={transformaciones controladas},
		description={Proceso sistem{\'a}tico de conversi{\'o}n de un artefacto a otro (por ejemplo, de modelos a c{\'o}digo, de especificaciones a configuraciones, o de modelos abstractos a modelos m{\'a}s concretos) bajo reglas definidas, verificables y trazables. Su prop{\'o}sito es preservar el significado (sem{\'a}ntica) y mantener correspondencia entre niveles de abstracci{\'o}n, reduciendo errores humanos y evitando divergencia. En enfoques dirigidos por modelos y en \glspl{sibiot}, habilita automatizaci{\'o}n, consistencia arquitect{\'o}nica y verificaci{\'o}n temprana, integr{\'a}ndose con pipelines de construcci{\'o}n, prueba y despliegue.},
		type=main,
		symbol={🧬},
		user1={Conversi{\'o}n trazable entre artefactos},
		user2={Preservaci{\'o}n sem{\'a}ntica},
		user3={Automatizaci{\'o}n con control},
		see={modelo,especificaciontecnica,consistencia,verificacion,divergencia,pipeline}
	}
	
	\newglossaryentry{transformaciondigital}{
		name={Transformaci{\'o}n digital},
		text={transformaci{\'o}n digital},
		sort={transformacion digital},
		plural={transformaciones digitales},
		first={transformaci{\'o}n digital},
		firstplural={transformaciones digitales},
		description={La \emph{transformaci{\'o}n digital} es el proceso mediante el cual las organizaciones incorporan tecnologías digitales para redefinir sus productos, servicios, modelos operativos y formas de generar valor. Este proceso trasciende la simple automatizaci{\'o}n, abarcando cambios culturales, estructurales y estrat{\'e}gicos que impulsan la innovaci{\'o}n continua. En contextos de \gls{iot} y \glspl{sibiot}, la transformaci{\'o}n digital implica integrar \glspl{dispositivointeligente}, anal{\'i}tica avanzada, \gls{bigdata}, \gls{cloudcomputing} y \gls{automatizacion} para optimizar procesos, habilitar toma de decisiones en \gls{tiemporeal} y mejorar la \gls{ux} y organizaciones. Constituye un pilar fundamental para la competitividad y la evoluci{\'o}n tecnol{\'o}gica \gls{sostenible}.},
		type=main,
		symbol={🚀},
		user1={Adopci{\'o}n de tecnolog{\'i}as digitales},
		user2={Cambio organizacional y cultural},
		user3={Integraci{\'o}n IoT, datos y automatizaci{\'o}n},
		see={innovacion,automatizacion,gestionempresarial,sibiot,bigdata}
	}
	
	\newglossaryentry{transicion}{
		name={Transici{\'o}n},
		text={transici{\'o}n},
		sort={transicion},
		plural={transiciones},
		first={transici{\'o}n},
		firstplural={transiciones},
		description={La \emph{transici{\'o}n} es el proceso de cambio o paso gradual de un estado, fase, tecnolog{\'i}a o metodolog{\'i}a a otra, manteniendo la \gls{continuidad} operativa y reduciendo riesgos. En el contexto de \glspl{sibiot}, puede referirse a la evoluci{\'o}n de infraestructuras tradicionales hacia arquitecturas inteligentes, a la adopci{\'o}n de nuevas \glspl{metodologia} de desarrollo como las {\'a}giles, o al reemplazo progresivo de dispositivos y protocolos por alternativas m{\'a}s eficientes. La transici{\'o}n requiere planificaci{\'o}n, \gls{adaptabilidad} y mecanismos de mitigaci{\'o}n para garantizar \gls{estabilidad} y \gls{confiabilidad} durante el proceso.},
		type=main,
		symbol={🔄},
		user1={Cambio de estado o fase en sistemas y procesos},
		user2={Evoluci{\'o}n tecnol{\'o}gica y metodol{\'o}gica},
		user3={Gestión planificada para asegurar continuidad},
		see={continuidad,adaptabilidad,metodologia,sibiot}
	}
	
	\newglossaryentry{transmision}{
		name={Transmisi{\'o}n},
		text={transmisi{\'o}n},
		sort={transmision},
		plural={transmisiones},
		first={transmisi{\'o}n},
		firstplural={transmisiones},
		description={La \emph{transmisi{\'o}n} es el proceso mediante el cual los datos o se{\~n}ales se env{\'i}an desde un origen hasta un destino a trav{\'e}s de un medio f{\'i}sico o inal{\'a}mbrico. Puede ser anal{\'o}gica o digital, unidireccional (simplex), bidireccional alterna (half-duplex) o bidireccional simult{\'a}nea (full-duplex). En el contexto de \gls{iot}, la transmisi{\'o}n asegura que los \glspl{sensor}, \glspl{actuador} y \glspl{dispositivoconectado} comuniquen informaci{\'o}n en tiempo real a trav{\'e}s de \glspl{protocolocomunicacion}, preservando \gls{integridad}, \gls{latencia} m{\'i}nima y \gls{fiabilidad}.},
		type=main,
		symbol={📡},
		user1={Env{\'i}o de datos entre sistemas},
		user2={Puede ser anal{\'o}gica o digital},
		user3={Modalidades simplex, half-duplex y full-duplex},
		see={protocolocomunicacion,integridad,latencia,fiabilidad,dispositivoconectado}
	}
	
	\newglossaryentry{transparencia}{
		name={Transparencia},
		text={transparencia},
		sort={transparencia},
		plural={transparencias},
		first={transparencia},
		firstplural={transparencias},
		description={Propiedad de un \gls{sistema}, organizaci{\'o}n o proceso que permite que sus operaciones, decisiones y resultados sean comprensibles, visibles y accesibles para los interesados, reduciendo la opacidad y facilitando la confianza. En ingenier{\'i}a de sistemas y \glspl{sibiot}, la \emph{transparencia} tambi{\'e}n hace referencia a la capacidad de ocultar la complejidad interna de un sistema distribuido, de modo que el usuario o las aplicaciones accedan a los recursos como si fueran {\'u}nicos y centralizados, garantizando consistencia, \gls{fiabilidad} y usabilidad. Asimismo, en el plano organizacional, implica responsabilidad {\'e}tica y cumplimiento normativo respecto al tratamiento de datos, la \gls{privacidad} y la \gls{seguridad}.},
		type=main,
		symbol={🔍},
		user1={Claridad y accesibilidad de procesos},
		user2={Ocultamiento de complejidad en sistemas distribuidos},
		user3={Responsabilidad {\'e}tica y confianza},
		see={fiabilidad,privacidad,seguridad,sibiot}
	}
	
	\newglossaryentry{transporte}{
		name={Transporte},
		text={transporte},
		sort={transporte},
		plural={transportes},
		first={transporte},
		firstplural={transportes},
		description={Conjunto de medios, infraestructuras y operaciones destinadas al desplazamiento de personas o mercanc{\'i}as. En entornos basados en \gls{iot}, el transporte se optimiza mediante veh{\'i}culos inteligentes, \glspl{sensor} distribuidos, sistemas de posicionamiento (p.\,ej., \gls{gps}), an{\'a}lisis de tr{\'a}fico en \gls{tiemporeal} y gesti{\'o}n predictiva de rutas, mejorando la \gls{eficiencia}, la \gls{seguridad} y la \gls{sostenibilidad} de la \gls{movilidad}.},
		type=main,
		user1={Movilidad de personas y carga},
		user2={Infraestructura, medios y operaciones log{\'i}sticas},
		user3={Optimizable con IoT, an{\'a}lisis en tiempo real y rutas predictivas},
		user4={Objetivos: eficiencia, seguridad y sostenibilidad},
		see={iot,vehiculo,sensor,gps,logistica,ruta,seguridadvial,tiemporeal},
		symbol={\ensuremath{\mathsf{T}}}
	}
	
	\newglossaryentry{transporteinteligente}
	{
		name={Transporte inteligente},
		text={transporte inteligente},
		sort={transporte inteligente},
		plural={transportes inteligentes},
		first={transporte inteligente},
		firstplural={transportes inteligentes},
		description={\emph{\gls{transporteinteligente}} es el uso de \glspl{dispositivoiot}, \glspl{sensor}, \glspl{actuador} y \glspl{sistemacomunicacion} avanzados para mejorar la eficiencia, seguridad y sostenibilidad de los sistemas de transporte. Se integra con la \gls{movilidad} urbana y las \glspl{ciudadinteligente}, permitiendo la gesti{\'o}n de tr{\'a}fico en tiempo real, la monitorizaci{\'o}n de veh{\'i}culos aut{\'o}nomos, la optimizaci{\'o}n de rutas y la reducci{\'o}n de emisiones contaminantes.},
		type=\glsdefaulttype,
		symbol={🚌},
		user1={Smart transport},
		user2={Sistema de transporte inteligente},
		user3={ITS (Intelligent Transport Systems)},
		see={movilidad}
	}
	
	\newglossaryentry{trazabilidad}{
		name={Trazabilidad},
		text={trazabilidad},
		sort={trazabilidad},
		plural={mecanismos de trazabilidad},
		first={trazabilidad},
		firstplural={mecanismos de trazabilidad},
		description={Capacidad de identificar, documentar y seguir el historial, la ubicaci{\'o}n y la trayectoria de un elemento, proceso o dato a lo largo de su ciclo de vida. En ingenier{\'i}a de software, la \emph{trazabilidad} permite vincular \glspl{requisito} con sus implementaciones, pruebas y validaciones. En sistemas f{\'i}sicos o \gls{iot}, garantiza el seguimiento de productos, eventos o activos mediante tecnolog{\'i}as como \gls{rfid}, \glspl{sensor} o registros digitales. Es esencial para la transparencia, la auditor{\'i}a, la \gls{calidad} y la \gls{seguridad}.},
		type=main,
		symbol={🧭},
		user1={Seguimiento de requisitos en software},
		user2={Monitoreo de productos y activos f{\'i}sicos},
		user3={Soporte para auditor{\'i}a y transparencia},
		see={calidad,seguridad,rfid,sibiot}
	}
	
	\newglossaryentry{trazabilidaddecisional}
	{
		name={Trazabilidad decisional},
		text={trazabilidad decisional},
		sort={trazabilidad decisional},
		plural={trazabilidades decisionales},
		first={trazabilidad decisional},
		firstplural={trazabilidades decisionales},
		description={La \emph{trazabilidad decisional} es la capacidad de reconstruir el origen de una decisi{\'o}n automatizada mediante el registro estructurado de datos de entrada, modelo aplicado y acci{\'o}n ejecutada. Esta propiedad es necesaria en organizaciones donde las decisiones afectan activos cr{\'i}ticos o cumplimiento normativo, y constituye un mecanismo de auditor{\'i}a y control dentro del \gls{sibiot}.},
		type=\glsdefaulttype,
		symbol={🧾},
		user1={Auditor{\'i}a},
		user2={Gobernanza de datos},
		user3={Registro hist{\'o}rico},
		see={retroalimentacionoperacional,modelopredictivo,control}
	}
	
	\newglossaryentry{trazabilidadoperacional}{
		name={Trazabilidad operacional},
		text={trazabilidad operacional},
		sort={trazabilidad operacional},
		plural={trazabilidades operacionales},
		first={trazabilidad operacional},
		firstplural={trazabilidades operacionales},
		description={Capacidad de registrar y reconstruir, con contexto verificable, la secuencia de mediciones, decisiones y acciones ejecutadas por un sistema a lo largo del tiempo. En \glspl{cps} y en sistemas industriales, la trazabilidad operacional respalda auditor{\'\i}a, an{\'a}lisis de incidentes, mantenimiento y cumplimiento, al permitir asociar cambios y actuaci{\'o}n con reglas, configuraciones, identidad de actores y evidencia t{\'e}cnica.},
		type=main,
		symbol={🧾},
		user1={Registro verificable},
		user2={Auditor{\'\i}a e incidentes},
		user3={Evidencia de decisiones},
		see={auditoria,integridad,ciberseguridad}
	}
	
	\newglossaryentry{tsql}{
		name={T-SQL},
		text={T-SQL},
		sort={tsql},
		first={Transact-SQL (T-SQL)},
		description={Extensi{\'o}n de SQL utilizada en Microsoft SQL Server para consultas, definici{\'o}n de datos y programaci{\'o}n de procedimientos, funciones y disparadores.},
		type=main,
		symbol={🧮},
		user1={SQL extendido},
		user2={Procedimientos},
		user3={SQL Server},
		see={sgbd,basedatos}
	}
	
	\newglossaryentry{uart}{
		name={UART},
		text={UART},
		sort={uart},
		first={Interfaz de comunicaci{\'o}n serie as{\'i}ncrona (\emph{Universal Asynchronous Receiver/Transmitter}, UART)},
		description={Interfaz de comunicaci{\'o}n serie as{\'i}ncrona o \emph{Universal Asynchronous Receiver/Transmitter}, UART) que permite el intercambio de datos entre dispositivos mediante l{\'i}neas de transmisi{\'o}n (TX) y recepci{\'o}n (RX). La \emph{UART} es ampliamente utilizada para depuraci{\'o}n, configuraci{\'o}n y comunicaci{\'o}n con computadoras o m{\'o}dulos externos en sistemas embebidos y IoT.},
		type=main,
		symbol={🔌},
		user1={Comunicaci{\'o}n serie},
		user2={TX/RX},
		user3={Depuraci{\'o}n},
		see={spi,capared}
	}
	
	\newglossaryentry{uddi}
	{
		name={{UDDI}},
		text={{UDDI}},
		sort={uddi},
		plural={{UDDI}},
		first={descripci{\'o}n universal, descubrimiento e integraci{\'o}n (\textit{Universal Description, Discovery and Integration}, UDDI)},
		firstplural={{UDDI}},
		description={\emph{Universal Description, Discovery and Integration} (\gls{uddi}) es una especificaci{\'o}n para publicar y localizar descripciones de servicios mediante un registro con metadatos, taxonom{\'\i}as y mecanismos de consulta, usada para soportar descubrimiento y administraci{\'o}n de servicios en entornos orientados a servicios.},
		type=\glsdefaulttype,
		symbol={📚},
		user1={Registro de servicios},
		user2={Descubrimiento},
		user3={Metadatos},
		see={soa,wsdl,api}
	}
	
	\newglossaryentry{udp}{
		name={UDP},
		text={UDP},
		sort={udp},
		plural={UDP},
		first={protocolo de datagrama de usuario (\textit{User Datagram Protocol}, UDP)},
		firstplural={UDP},
		description={\emph{UDP} es un protocolo de transporte no orientado a conexi{\'o}n que permite el env{\'i}o de datagramas sin garantizar su entrega, orden o integridad. Su dise{\~n}o ligero y la ausencia de mecanismos complejos de control lo hacen especialmente eficiente para comunicaciones r{\'a}pidas y en tiempo real. En sistemas \gls{iot}, UDP es ampliamente utilizado en aplicaciones que toleran p{\'e}rdidas parciales de datos o requieren baja \gls{latencia}, como \gls{transmision} de \gls{telemetria}, sensores de alta frecuencia, descubrimiento de servicios, protocolos como \gls{coap} y comunicaciones en redes de baja potencia.},
		type=main,
		symbol={📨},
		user1={Transporte ligero sin conexi{\'o}n},
		user2={Id{\'o}neo para baja latencia y alto rendimiento},
		user3={Base de protocolos IoT como CoAP},
		see={tcp,coap,protocolocomunicacion,ip}
	}
	
	\newglossaryentry{uefi}{
		name={UEFI},
		text={UEFI},
		sort={uefi},
		plural={interfaces UEFI},
		first={interfaz de firmware extensible unificada (\textit{Unified Extensible Firmware Interface}, UEFI)},
		firstplural={interfaces UEFI},
		description={Especificaci{\'o}n que define la interfaz entre el sistema operativo y el firmware de un dispositivo de c{\'o}mputo. La \emph{UEFI} reemplaza al tradicional \gls{bios}, ofreciendo un entorno m{\'a}s flexible, seguro y con soporte para discos duros de gran capacidad, arranque seguro (\textit{secure boot}), configuraciones avanzadas y una mayor velocidad de inicializaci{\'o}n. En sistemas embebidos e IoT, UEFI proporciona mecanismos modernos para la gesti{\'o}n del arranque, la \gls{seguridad} y la \gls{interoperabilidad} de \gls{hardware}.},
		type=main,
		symbol={🖥️},
		user1={Firmware moderno},
		user2={Reemplazo del BIOS},
		user3={Interfaz de arranque avanzado},
		see={bios,hardware,seguridad}
	}
	
	\newglossaryentry{uhf}{
		name={UHF},
		text={UHF},
		sort={uhf},
		plural={bandas UHF},
		first={frecuencia ultraalta (\textit{Ultra High Frequency}, UHF)},
		firstplural={UHFs},
		description={Rango del espectro electromagn{\'e}tico comprendido entre 300~MHz y 3~GHz. Las se{\~n}ales en \emph{UHF} se utilizan ampliamente en comunicaciones \gls{comunicacioninalambrica}, incluyendo televisi{\'o}n digital, telefon{\'i}a m{\'o}vil, \gls{wifi}, \gls{bluetooth} y tecnolog{\'i}as \gls{rfid}. En aplicaciones \gls{iot}, la banda UHF es valorada por su capacidad para transmitir datos a distancias medias, con buena penetraci{\'o}n en entornos urbanos y consumo energ{\'e}tico moderado, lo que la hace ideal para redes de \glspl{sensor}, identificaci{\'o}n remota y gesti{\'o}n de activos.},
		type=main,
		symbol={📡},
		user1={Rango 300~MHz -- 3~GHz},
		user2={Aplicaciones en TV digital, m{\'o}vil y RFID},
		user3={Transmisi{\'o}n media distancia en IoT},
		see={rfid,wifi,bluetooth,comunicacioninalambrica}
	}
	
	\newglossaryentry{unidadfuncional}{
		name={Unidad funcional},
		text={unidad funcional},
		sort={unidadfuncional},
		plural={unidades funcionales},
		first={unidad funcional},
		firstplural={unidades funcionales},
		description={Conjunto cohesivo de funcionalidades y artefactos asociados que, de manera integrada, proporcionan una capacidad completa y verificable dentro de un sistema. Una unidad funcional incluye comportamiento observable, interfaces, datos requeridos y criterios de aceptaci{\'o}n, y puede evolucionar mediante incrementos sin perder operatividad. En \glspl{sibiot}, una unidad funcional puede abarcar desde la adquisici{\'o}n de datos de sensores y su procesamiento, hasta la exposici{\'o}n de servicios, almacenamiento, anal{\'\i}tica y mecanismos de control o actuaci{\'o}n, con pruebas que validen el extremo a extremo.},
		type=main,
		symbol={🧩},
		user1={Funcionalidad completa y verificable},
		user2={Cohesi{\'o}n y criterio de aceptaci{\'o}n},
		user3={Bloque para desarrollo incremental},
		see={incremental,iterativa,ciclodesarrollo,entregable,pruebaautomatizada,integracion}
	}
	
	\newglossaryentry{uri}{
		name={Identificador Uniforme de Recurso},
		text={URI},
		sort={uri},
		plural={URIs},
		first={Identificador Uniforme de Recurso (\textit{Uniform Resource Identifier}, URI)},
		firstplural={URIs},
		description={Un identificador uniforme de recurso o \textit{Uniform Resource Identifier} (URI) es un esquema estandarizado que permite localizar, nombrar o acceder a recursos en Internet o en redes privadas. En los \glspl{sibiot}, los URI se emplean para identificar dispositivos, servicios, sensores o flujos de datos, facilitando la interoperabilidad entre componentes distribuidos. Su estructura general puede incluir un esquema (como \texttt{http}, \texttt{mqtt}, \texttt{coap}), autoridad, ruta y par{\'a}metros, permitiendo una gesti{\'o}n precisa y universal de los recursos digitales.},
		type=\glsdefaulttype,
		symbol={🌐},
		user1={Uniform Resource Identifier},
		user2={Cadena que identifica de forma {\'u}nica un recurso en una red},
		user3={Estructura que facilita la interoperabilidad entre dispositivos y servicios IoT},
		see={url,api,protocolo,interoperabilidad,sibiot}
	}
	
	\newglossaryentry{url}{
		name={URL},
		text={URL},
		sort={url},
		first={Localizador Universal de Recursos (\emph{Uniform Resource Locator}, URL)},
		description={Localizador Universal de Recursos o \emph{Uniform Resource Locator}, URL). Identificador que especifica la ubicaci{\'o}n y el mecanismo de acceso a un recurso en una red, t{\'i}picamente en la Web.}, 
		type=main,
		symbol={🔗},
		user1={Direcci{\'o}n de recurso},
		user2={Web},
		user3={Acceso},
		see={plataformaweb,websocket}
	}
	
	\newglossaryentry{usabilidad}{
		name={Usabilidad},
		text={usabilidad},
		sort={usabilidad},
		plural={m{\'e}tricas de usabilidad},
		first={usabilidad},
		firstplural={m{\'e}tricas de usabilidad},
		description={Grado en que un sistema, producto o servicio puede ser utilizado por usuarios determinados para alcanzar objetivos con efectividad, eficiencia y satisfacci{\'o}n en un contexto de uso definido. En \glspl{sibiot}, la \emph{usabilidad} se relaciona con dise{\~n}os accesibles, interfaces intuitivas y experiencias centradas en el \gls{usuario}.},
		type=main,
		symbol={🧩},
		user1={Facilidad de uso},
		user2={Eficiencia e intuitividad},
		user3={Experiencia del usuario},
		see={interfaz,accesibilidad,calidadexperiencia}
	}
	
	\newglossaryentry{usb}{
		name={USB},
		text={USB},
		sort={usb},
		plural={interfaces USB},
		first={Bus Serial Universal (\textit{Universal Serial Bus}, USB)},
		firstplural={interfaces USB},
		description={Est{\'a}ndar industrial para la conexi{\'o}n, comunicaci{\'o}n y suministro de energ{\'i}a entre computadoras y dispositivos perif{\'e}ricos. En sistemas \gls{iot}, el \emph{USB} se emplea para la programaci{\'o}n de \glspl{microcontrolador}, transferencia de datos, depuraci{\'o}n de sistemas embebidos y alimentaci{\'o}n el{\'e}ctrica de nodos de baja potencia. Sus variantes incluyen USB 2.0, 3.0 y USB-C, que difieren en velocidad y capacidad de energ{\'i}a.},
		type=main,
		symbol={🔌},
		user1={Conexi{\'o}n universal de perif{\'e}ricos},
		user2={Programaci{\'o}n y depuraci{\'o}n en IoT},
		user3={Alimentaci{\'o}n de nodos de baja potencia},
		see={hardware,comunicacioninalambrica,microcontrolador}
	}
	
	\newglossaryentry{usuario}{
		name={usuario},
		text={usuario},
		sort={usuario},
		plural={usuarios},
		first={usuario},
		firstplural={usuarios},
		description={En el contexto de los \glspl{sibiot} y de la \gls{ingenieriasoftware}, el t{\'e}rmino usuario se refiere a la persona, grupo u organizaci{\'o}n que utiliza o se beneficia de las funcionalidades proporcionadas por un sistema tecnol{\'o}gico. El usuario puede ser final, cuando emplea directamente la interfaz del sistema; administrador, cuando gestiona la infraestructura y el acceso; o desarrollador, cuando interviene en la creaci{\'o}n y mantenimiento del sistema. La identificaci{\'o}n de los tipos de usuarios y sus necesidades constituye una fase esencial del \gls{analisisrequisitos} y del dise{\~n}o centrado en el usuario, favoreciendo la \gls{usabilidad}, \gls{accesibilidad} y satisfacci{\'o}n en el uso del \gls{sistema}.},
		type=\glsdefaulttype,
		symbol={👤},
		user1={User},
		user2={Persona que interact{\'u}a con un sistema o dispositivo},
		user3={Agente humano o institucional que utiliza un sistema tecnol{\'o}gico},
		see={interfaz,accesibilidad,usabilidad,ingenieriasoftware,analisisrequisitos}
	}
	
	\newglossaryentry{ux}{
		name={UX},
		short={UX},
		long={Experiencia de Usuario (\textit{User eXperience}, UX)},
		text={UX},
		sort={ux},
		first={Experiencia de Usuario (UX)},
		firstplural={Experiencias de Usuario (UX)},
		plural={Experiencias de Usuario},
		description={Disciplina que se enfoca en el dise{\~n}o, evaluaci{\'o}n y mejora de la experiencia que una persona tiene al interactuar con un sistema, producto o servicio digital. La UX abarca aspectos como usabilidad, accesibilidad, percepci{\'o}n, satisfacci{\'o}n, eficiencia y respuesta emocional del usuario. En sistemas \gls{iot} y plataformas digitales, una buena UX garantiza que la interacci{\'o}n con dispositivos, aplicaciones y servicios sea intuitiva, coherente y alineada con las necesidades reales de los usuarios, facilitando la adopci{\'o}n del sistema y la toma de decisiones informadas},
		type=\glsdefaulttype,
		symbol={🧠},
		user1={Experiencia centrada en el usuario},
		user2={Interacci{\'o}n humano--computador},
		user3={Dise{\~n}o orientado a la usabilidad},
		see={usabilidad,accesibilidad}
	}
	
	\newglossaryentry{validacion}{
		name={Validaci{\'o}n},
		text={validaci{\'o}n},
		sort={validacion},
		plural={validaciones},
		first={validaci{\'o}n},
		firstplural={validaciones},
		description={En la gesti{\'o}n de calidad de datos, la validaci{\'o}n garantiza que la informaci{\'o}n almacenada o procesada cumpla con criterios de integridad, formato y consistencia. Este proceso incluye la verificaci{\'o}n de tipos de datos, rangos permitidos, unicidad, completitud y relaciones l{\'o}gicas entre registros. Una validaci{\'o}n adecuada previene errores anal{\'i}ticos, evita decisiones basadas en informaci{\'o}n incorrecta y mejora la \gls{fiabilidad} de los sistemas \gls{sibiot} y de inteligencia organizacional.},
		type=\glsdefaulttype,
		symbol={✔️},
		user1={Data validation},
		user2={Verificaci{\'o}n de la exactitud y coherencia de los datos},
		user3={Control de calidad de los datos antes de su almacenamiento o an{\'a}lisis},
		see={calidaddatos,limpieza,enriquecimiento,integridad}
	}
	
	\newglossaryentry{valor}{
		name={Valor},
		text={valor},
		sort={valor},
		first={valor},
		description={Beneficio o utilidad que un sistema, producto o servicio aporta a uno o varios interesados, en relaci{\'o}n con sus objetivos, necesidades y contexto de uso. En sistemas de informaci{\'o}n y \glspl{sibiot}, el \emph{valor} se manifiesta a trav{\'e}s de la mejora de procesos, la optimizaci{\'o}n de recursos, el apoyo a la toma de decisiones, la reducci{\'o}n de riesgos o la generaci{\'o}n de ventajas competitivas, y no se limita exclusivamente a aspectos econ{\'o}micos.},
		type=main,
		symbol={💎},
		user1={Beneficio},
		user2={Utilidad},
		user3={Impacto},
		see={resultado,objetivo,innovacion}
	}
	
	\newglossaryentry{valoragregado}{
		name={Valor agregado},
		text={valor agregado},
		sort={valoragregado},
		plural={valores agregados},
		first={valor agregado},
		firstplural={valores agregados},
		description={El valor agregado se refiere al aporte diferencial que resulta de la incorporaci{\'o}n de innovaci{\'o}n, inteligencia o eficiencia en los procesos. En los \glspl{sibiot}, este valor se genera mediante el uso de datos procesados, an{\'a}lisis predictivos y automatizaci{\'o}n, lo que permite optimizar decisiones, reducir costos y aumentar la satisfacci{\'o}n del usuario. Representa el impacto tangible e intangible que las tecnolog{\'i}as inteligentes aportan a la \gls{gestion} y \gls{transformaciondigital} de las organizaciones.},
		type=\glsdefaulttype,
		symbol={💡},
		user1={Added value},
		user2={Beneficio adicional que mejora utilidad o competitividad},
		user3={Resultado del uso inteligente de datos en SIbIoT},
		see={sibiot,analisis,innovacion,gestionempresarial,eficiencia}
	}
	
	\newglossaryentry{values}{
		name={Values},
		text={values},
		sort={values},
		first={values},
		description={Conjunto de recursos declarativos en Android (cadenas, colores, dimensiones, estilos) que soportan internacionalizaci{\'o}n y consistencia visual mediante definiciones centralizadas.},
		type=main,
		symbol={🏷️},
		user1={Recursos declarativos},
		user2={Android},
		user3={Consistencia},
		see={assets,layout,drawable}
	}
	
	\newglossaryentry{variabilidad}{
		name={Variabilidad},
		text={variabilidad},
		sort={variabilidad},
		first={variabilidad},
		description={Grado de cambio u oscilaci{\'o}n de una variable o comportamiento. En sistemas, puede referirse a variaciones de carga, se{\~n}ales, demanda o condiciones operativas que afectan dise{\~n}o y control.},
		type=main,
		symbol={📉},
		user1={Oscilaci{\'o}n},
		user2={Incertidumbre},
		user3={Condiciones cambiantes},
		see={analisisstreaming,desempeno}
	}
	
	\newglossaryentry{variablemacroeconomica}{
		name={variable macroecon{\'o}mica},
		plural={variables macroecon{\'o}micas},
		description={Indicador cuantitativo o cualitativo que describe el estado, la evoluci{\'o}n o el comportamiento agregado de una econom{\'i}a a nivel nacional, regional o sectorial. Las variables macroecon{\'o}micas reflejan fen{\'o}menos de gran escala, tales como crecimiento econ{\'o}mico, inflaci{\'o}n, empleo, consumo, inversi{\'o}n y estabilidad financiera, y se utilizan para el an{\'a}lisis, la planificaci{\'o}n y la toma de decisiones estrat{\'e}gicas},
		long={variable macroecon{\'o}mica},
		short={VM},
		see={macroeconomia}
	}
	
	\newglossaryentry{variablefisica}{
		name={Variable f{\'i}sica},
		text={variable f{\'i}sica},
		sort={variable fisica},
		plural={variables f{\'i}sicas},
		first={variable f{\'i}sica},
		firstplural={variables f{\'i}sicas},
		description={Magnitud medible del mundo real que describe una propiedad o condici{\'o}n del entorno natural o artificial, susceptible de ser registrada mediante un \gls{sensor}. Ejemplos de \emph{variables f{\'i}sicas} incluyen temperatura, presi{\'o}n, humedad, velocidad, aceleraci{\'o}n, nivel de luz, caudal o tensi{\'o}n el{\'e}ctrica. En el contexto de los \glspl{sibiot}, la medici{\'o}n de variables f{\'i}sicas es la base para la monitorizaci{\'o}n, el control y la \gls{tomadecisiones} en tiempo real, ya que permiten transformar fen{\'o}menos f{\'i}sicos en datos digitales procesables por sistemas inteligentes.},
		type=main,
		symbol={📏},
		user1={Magnitud medible},
		user2={Dato de entrada en SIbIoT},
		user3={Base para la monitorizaci{\'o}n y el control},
		see={sensor,monitorizacion,sibiot,control}
	}
	
	\newglossaryentry{vehiculo}{
		name={Veh{\'i}culo},
		text={veh{\'i}culo},
		sort={vehiculo},
		plural={veh{\'i}culos},
		first={veh{\'i}culo},
		firstplural={veh{\'i}culos},
		description={Medio de transporte motorizado o no motorizado destinado al desplazamiento de personas o mercanc{\'i}as. En IoT automotriz, incorpora sensores y comunicaciones para seguridad, mantenimiento y servicios conectados.},
		type=main,
		symbol={🚗},
		user1={Transporte},
		user2={Sensores},
		user3={Conectividad},
		see={seguridadvial,sistemanavegacion}
	}
	
	\newglossaryentry{vehiculoautonomo}{
		name={Veh{\'i}culo aut{\'o}nomo},
		text={veh{\'i}culo aut{\'o}nomo},
		sort={vehiculo autonomo},
		plural={veh{\'i}culos aut{\'o}nomos},
		first={veh{\'i}culo aut{\'o}nomo},
		firstplural={veh{\'i}culos aut{\'o}nomos},
		description={Medio de transporte equipado con \glspl{sensor}, \glspl{actuador}, \glspl{sistemainteligente} y algoritmos de \gls{ia} que le permiten desplazarse y tomar decisiones de manera aut{\'o}noma, sin intervenci{\'o}n humana directa. Un \emph{veh{\'i}culo aut{\'o}nomo} integra tecnolog{\'i}as como \gls{visionartificial}, \glspl{redneuronal}, \glspl{sistemanavegacion} GNSS y sistemas de comunicaci{\'o}n V2X (veh{\'i}culo a todo) para percibir su entorno, planificar rutas y ejecutar maniobras seguras. En el contexto de los \glspl{sibiot}, los veh{\'i}culos aut{\'o}nomos se aplican en movilidad urbana inteligente, log{\'i}stica, agricultura de precisi{\'o}n y transporte sostenible, optimizando recursos y reduciendo riesgos de accidente.},
		type=main,
		symbol={🚗},
		user1={Transporte aut{\'o}nomo e inteligente},
		user2={Integraci{\'o}n de sensores y sistemas de navegaci{\'o}n},
		user3={Aplicaciones en movilidad y log{\'i}stica IoT},
		see={sistemanavegacion,visionartificial,sensor,sibiot,ia}
	}
	
	\newglossaryentry{velocidad}{
		name={velocidad},
		text={velocidad},
		sort={velocidad},
		plural={velocidades},
		first={velocidad},
		firstplural={velocidades},
		description={En la gesti{\'o}n de datos, la velocidad hace referencia a la rapidez con la que la informaci{\'o}n fluye desde su origen hasta su an{\'a}lisis o almacenamiento. En los \glspl{sibiot}, este factor resulta esencial para la \gls{capturadatos} en \gls{tiemporeal} y el soporte de decisiones r{\'a}pidas. La velocidad influye en la arquitectura del sistema, el rendimiento de las bases de datos y la eficiencia de los procesos de \gls{analisistreaming}. Una alta velocidad exige mecanismos de procesamiento inmediato, compresi{\'o}n de datos y priorizaci{\'o}n de flujos cr{\'i}ticos.},
		type=\glsdefaulttype,
		symbol={⚡},
		user1={Velocity},
		user2={Rapidez en la generaci{\'o}n o procesamiento de datos},
		user3={Frecuencia con que los sistemas IoT actualizan o transmiten informaci{\'o}n},
		see={volumen,analisisstreaming,tiemporeal,almacenamiento}
	}
	
	\newglossaryentry{velocitychart}{
		name={Velocity Chart},
		text={\textit{Velocity Chart}},
		sort={velocity chart},
		plural={\textit{Velocity Charts}},
		first={\textit{Velocity Chart}},
		firstplural={\textit{Velocity Charts}},
		description={Representaci{\'o}n gr{\'a}fica de la cantidad de trabajo completado por el equipo en cada Sprint. Se utiliza para estimar capacidad futura y mejorar la planificaci{\'o}n basada en datos hist{\'o}ricos.},
		type=main,
		symbol={⚙️},
		user1={Medici{\'o}n de rendimiento},
		user2={Capacidad hist{\'o}rica del equipo},
		user3={Soporte para planificaci{\'o}n},
		see={sprint,burndownchart,productbacklog}
	}
	
	\newglossaryentry{verificacion}{
		name={Verificaci{\'o}n},
		text={verificaci{\'o}n},
		sort={verificacion},
		plural={verificaciones},
		first={verificaci{\'o}n},
		firstplural={verificaciones},
		description={Actividad mediante la cual se comprueba, con evidencia objetiva, que un artefacto o sistema cumple lo especificado (requisitos, restricciones y propiedades) sin evaluar necesariamente su adecuaci{\'o}n al uso final. Se orienta a responder si ``se construy{\'o} correctamente'' respecto de la especificaci{\'o}n. Incluye revisiones, an{\'a}lisis est{\'a}ticos, pruebas unitarias e integraci{\'o}n, verificaci{\'o}n de modelos y comprobaciones autom{\'a}ticas en pipelines. En IoT/SIbIoT, la verificaci{\'o}n incorpora adem{\'a}s aspectos de comunicaciones, interoperabilidad, latencia, consumo energ{\'e}tico y seguridad, debido a la naturaleza distribuida del sistema.},
		type=main,
		symbol={🔍},
		user1={Conformidad con especificaci{\'o}n},
		user2={Evidencia objetiva},
		user3={Automatizable en pipeline},
		see={validacion,consistencia,pruebaautomatizada,especificaciontecnica,pipeline}
	}
	
	\newglossaryentry{versatilidad}{
		name={Versatilidad},
		text={versatilidad},
		sort={versatilidad},
		plural={versatilidades},
		first={versatilidad},
		firstplural={versatilidades},
		description={Capacidad de un \gls{sistema}, componente o tecnolog{\'i}a para desempe{\~n}ar diversas funciones, adaptarse a diferentes entornos o cumplir m{\'u}ltiples prop{\'o}sitos sin necesidad de modificaciones profundas. La \emph{versatilidad} implica flexibilidad en el dise{\~n}o, interoperabilidad y facilidad de integraci{\'o}n con otros sistemas. En el contexto de los \glspl{sibiot}, la versatilidad se refleja en plataformas capaces de gestionar distintos tipos de \glspl{sensor}, soportar variadas \glspl{aplicacion} y adaptarse a necesidades cambiantes en dominios como salud, agricultura, industria o ciudades inteligentes, favoreciendo la \gls{adaptabilidad} y la \gls{sostenibilidad}.},
		type=main,
		symbol={🔀},
		user1={Capacidad multifuncional},
		user2={Adaptaci{\'o}n a diversos entornos},
		user3={Facilidad de integraci{\'o}n en SIbIoT},
		see={adaptabilidad,sostenibilidad,sibiot,sistemainteligente}
	}
	
	\newglossaryentry{vga}{
		name={VGA},
		text={VGA},
		plural={VGA},
		sort={vga},
		first={matriz de gr{\'a}ficos de video (\textit{Video Graphics Array}, VGA)},
		description={Est{\'a}ndar de visualizaci{\'o}n anal{\'o}gico desarrollado por IBM para la transmisi{\'o}n de se{\~n}ales de video desde un dispositivo a una pantalla. Aunque ha sido sustituido por tecnolog{\'i}as digitales modernas como HDMI o DisplayPort, el conector VGA todav{\'i}a se emplea en algunos sistemas embebidos y entornos \gls{iot} para fines de depuraci{\'o}n visual o interfaces gr{\'a}ficas b{\'a}sicas.},
		type=main,
		symbol={🖥️},
		see={hdmi,displayport}
	}
	
	\newglossaryentry{viabilidad}{
		name={Viabilidad},
		text={viabilidad},
		sort={viabilidad},
		plural={viabilidades},
		first={viabilidad},
		firstplural={viabilidades},
		description={Grado en que un proyecto, sistema o soluci{\'o}n puede desarrollarse y mantenerse de manera exitosa dentro de unas condiciones determinadas de recursos, tiempo, tecnolog{\'i}a y contexto organizacional. La \emph{viabilidad} se eval{\'u}a desde diferentes perspectivas: t{\'e}cnica (capacidad de implementaci{\'o}n), econ{\'o}mica (costes y beneficios), operativa (factibilidad de uso y gesti{\'o}n) y ambiental/social (impacto en el entorno y sostenibilidad). En el contexto de los \glspl{sibiot}, la viabilidad implica determinar si una infraestructura con \glspl{sensor}, \glspl{actuador} y servicios en la \gls{nube} puede desplegarse eficientemente, integrarse con sistemas existentes y garantizar \gls{eficienciaoperativa}, seguridad y \gls{sostenibilidad} a largo plazo.},
		type=main,
		symbol={✔️},
		user1={Factibilidad de implementaci{\'o}n y mantenimiento},
		user2={Evaluaci{\'o}n t{\'e}cnica, econ{\'o}mica y operativa},
		user3={Determinaci{\'o}n de sostenibilidad en SIbIoT},
		see={sostenibilidad,eficienciaoperativa,sibiot}
	}
	
	\newglossaryentry{viabilidadtecnica}{
		name={Viabilidad t{\'e}cnica},
		text={viabilidad t{\'e}cnica},
		sort={viabilidad tecnica},
		plural={viabilidades t{\'e}cnicas},
		first={viabilidad t{\'e}cnica},
		firstplural={viabilidades t{\'e}cnicas},
		description={Evaluaci{\'o}n espec{\'i}fica de la \gls{viabilidad} de un proyecto o \gls{sistema} en funci{\'o}n de las capacidades tecnol{\'o}gicas disponibles, los recursos humanos especializados y la infraestructura requerida. La \emph{viabilidad t{\'e}cnica} determina si es posible dise{\~n}ar, implementar e integrar los \glspl{componente} necesarios (hardware, software, redes y plataformas) para alcanzar los objetivos planteados. En el contexto de los \glspl{sibiot}, implica analizar si la combinaci{\'o}n de \glspl{sensor}, \glspl{actuador}, protocolos de \gls{comunicacioninalambrica}, servicios en la \gls{nube} y \glspl{modelopredictivo} puede operar de forma segura, \gls{eficiente} y \gls{escalable}, garantizando la \gls{eficienciaoperativa} y la \gls{sostenibilidad}.},
		type=main,
		symbol={🛠️},
		user1={Factibilidad tecnol{\'o}gica de un sistema},
		user2={Disponibilidad de recursos y capacidades t{\'e}cnicas},
		user3={Aplicaci{\'o}n en proyectos IoT y SIbIoT},
		see={viabilidad,componente,eficienciaoperativa,sostenibilidad,sibiot}
	}
	
	\newglossaryentry{vidaindependiente}{
		name={Vida independiente},
		text={vida independiente},
		sort={vida independiente},
		plural={estrategias de vida independiente},
		first={vida independiente},
		firstplural={estrategias de vida independiente},
		description={Modelo que promueve la autonom{\'i}a y calidad de vida de las personas, especialmente mayores o con discapacidad, mediante el uso de \gls{iot}, \glspl{sensor} y \glspl{wearable} para la \gls{monitorizacion} remota, la \gls{seguridad} y el bienestar en el hogar.},
		type=main,
		symbol={🏡},
		user1={Autonom{\'i}a asistida tecnol{\'o}gicamente},
		user2={Apoyo a personas mayores o con discapacidad},
		user3={Aplicaciones IoT orientadas al bienestar},
		see={hogarinteligente,wearable,seguridad,monitorizacion}
	}
	
	\newglossaryentry{virtualizacion}{
		name={Virtualizaci{\'o}n},
		text={virtualizaci{\'o}n},
		sort={virtualizacion},
		plural={virtualizaciones},
		first={virtualizaci{\'o}n},
		firstplural={virtualizaciones},
		description={T{\'e}cnica que permite abstraer recursos f{\'i}sicos de computaci{\'o}n, almacenamiento o red para presentarlos como recursos l{\'o}gicos independientes. La virtualizaci{\'o}n posibilita ejecutar m{\'u}ltiples entornos aislados sobre una misma infraestructura f{\'i}sica. En arquitecturas modernas de \gls{iot} y Edge Computing, la virtualizaci{\'o}n ligera mediante contenedores permite desplegar servicios en el borde con eficiencia, aislamiento y portabilidad.},
		type=main,
		symbol={🖥️},
		user1={Abstracci{\'o}n de recursos},
		user2={Aislamiento l{\'o}gico},
		user3={Contenedores y microservicios},
		see={infraestructuradistribuida,contenedor,orquestacion}
	}
	
	\newglossaryentry{visionartificial}{
		name={Visi{\'o}n artificial},
		text={visi{\'o}n artificial},
		sort={vision artificial},
		plural={visiones artificiales},
		first={visi{\'o}n artificial},
		firstplural={visiones artificiales},
		description={Campo de la \gls{ia} y de la ingenier{\'i}a computacional que desarrolla t{\'e}cnicas y sistemas para que las m{\'a}quinas puedan adquirir, procesar y comprender informaci{\'o}n visual a partir de im{\'a}genes o secuencias de video. La \emph{visi{\'o}n artificial} incluye procesos como \gls{deteccion} de objetos, \gls{deteccioncolor}, \gls{deteccionmovimiento}, \gls{reconocimiento} de patrones y \gls{identificacion} biom{\'e}trica. En el contexto de los \glspl{sibiot}, la visi{\'o}n artificial se aplica a vigilancia inteligente, control de calidad industrial, agricultura de precisi{\'o}n, diagn{\'o}stico m{\'e}dico asistido y veh{\'i}culos aut{\'o}nomos, aprovechando \glspl{redneuronal} y algoritmos de \gls{aprendizajeautomatico} para operar en \gls{tiemporeal}.},
		type=main,
		symbol={👁️‍🗨️},
		user1={Percepci{\'o}n visual computacional},
		user2={Procesamiento de im{\'a}genes y video},
		user3={Aplicaciones inteligentes en SIbIoT},
		see={deteccion,reconocimiento,identificacion,aprendizajeautomatico,redneuronal,sibiot}
	}
	
	\newglossaryentry{visioncomputacional}{
		name={Visi{\'o}n computacional},
		text={visi{\'o}n computacional},
		sort={vision computacional},
		first={visi{\'o}n computacional},
		description={Campo que desarrolla m{\'e}todos para extraer informaci{\'o}n de im{\'a}genes y video, permitiendo detecci{\'o}n, clasificaci{\'o}n, seguimiento y estimaci{\'o}n de atributos del entorno.},
		type=main,
		symbol={👁️},
		user1={Im{\'a}genes y video},
		user2={Detecci{\'o}n},
		user3={Percepci{\'o}n},
		see={analisisstreaming,monitorizacion}
	}
	
	\newglossaryentry{visioncomputador}{
		name={Visi{\'o}n por computador},
		text={visi{\'o}n por computador},
		sort={vision computador},
		plural={t{\'e}cnicas de visi{\'o}n por computador},
		first={visi{\'o}n por computador},
		firstplural={t{\'e}cnicas de visi{\'o}n por computador},
		description={Campo interdisciplinario que desarrolla algoritmos y modelos para que los sistemas inform{\'a}ticos puedan adquirir, procesar y comprender informaci{\'o}n visual proveniente del entorno. La \emph{visi{\'o}n por computador} integra matem{\'a}ticas, estad{\'i}stica, aprendizaje autom{\'a}tico y procesamiento de im{\'a}genes para tareas como reconocimiento de objetos, detecci{\'o}n de rostros, clasificaci{\'o}n de escenas, seguimiento de movimiento y reconstrucci{\'o}n 3D. En el contexto de los \glspl{sibiot}, es esencial para veh{\'i}culos aut{\'o}nomos, vigilancia inteligente, sistemas biom{\'e}tricos y rob{\'o}tica avanzada.},
		type=main,
		symbol={👁️‍🗨️},
		user1={Computer Vision},
		user2={An{\'a}lisis autom{\'a}tico de im{\'a}genes},
		user3={Percepci{\'o}n visual computacional},
		see={visionartificial,reconocimientofacial,deteccionobstaculos}
	}
	
	\newglossaryentry{visualizacion}{
		name={Visualizaci{\'o}n},
		text={visualizaci{\'o}n},
		sort={visualizacion},
		plural={visualizaciones},
		first={visualizaci{\'o}n},
		firstplural={visualizaciones},
		description={Proceso de representaci{\'o}n gr{\'a}fica o simb{\'o}lica de datos e informaci{\'o}n con el objetivo de facilitar su comprensi{\'o}n, an{\'a}lisis y toma de decisiones. En el contexto de \glspl{sibiot}, la \emph{visualizaci{\'o}n} se realiza mediante paneles interactivos, dashboards o gr{\'a}ficas que muestran datos recolectados por \glspl{sensor} en \gls{tiemporeal}, permitiendo a \glspl{usuario} y sistemas identificar patrones, anomal{\'i}as y tendencias.},
		type=main,
		symbol={📊},
		user1={Representaci{\'o}n gr{\'a}fica de datos},
		user2={Soporte a toma de decisiones},
		user3={Dashboards y paneles en sistemas IoT},
		see={monitorizacion,tiemporeal,sibiot}
	}
	
	\newglossaryentry{vlc}{
		name={VLC},
		text={VLC},
		plural={VLC},
		sort={vlc},
		first={comunicaci{\'o}n por luz visible (\textit{Visible Light Communication}, VLC)},
		description={Tecnolog{\'i}a de comunicaci{\'o}n inal{\'a}mbrica que utiliza el espectro de luz visible para transmitir datos. Aprovecha fuentes LED para enviar informaci{\'o}n mediante parpadeos imperceptibles, ofreciendo una alternativa a la radiofrecuencia en entornos donde las interferencias electromagn{\'e}ticas deben minimizarse. Se emplea en aplicaciones de IoT, comunicaciones interiores, veh{\'i}culos aut{\'o}nomos y sistemas de posicionamiento en interiores.},
		type=main,
		symbol={💡},
		see={wifi,bluetooth}
	}
	
	\newglossaryentry{volatilidad}{
		name={Volatilidad},
		text={volatilidad},
		sort={volatilidad},
		plural={volatilidades},
		first={volatilidad},
		firstplural={volatilidades},
		description={Grado de variabilidad, inestabilidad o fluctuaci{\'o}n en los datos, procesos, condiciones operativas o cargas de trabajo que afectan el comportamiento de un sistema. En los \glspl{sibiot}, la \emph{volatilidad} se manifiesta en cambios bruscos en las se{\~n}ales capturadas por \glspl{sensor}, en oscilaciones del tr{\'a}fico de red, en variaciones ambientales no previstas o en patrones din{\'a}micos del entorno f{\'i}sico. Comprender y gestionar la volatilidad es esencial para garantizar la robustez de los \glspl{modelopredictivo}, mejorar la \gls{tomadecisiones} en \gls{tiemporeal} y asegurar que la infraestructura distribuida---incluida la \gls{edgecomputing} y la \gls{fog}---mantenga niveles adecuados de \gls{eficienciaoperativa} y resiliencia.},
		type=main,
		symbol={🌪️},
		user1={Fluctuaciones en datos, carga o condiciones del entorno},
		user2={Factor cr{\'i}tico en modelos predictivos y procesamientos IoT},
		user3={Indicador de estabilidad y resiliencia en sistemas distribuidos},
		see={variabilidad,tiemporeal,modelopredictivo,eficienciaoperativa,fog}
	}
	
	\newglossaryentry{volumen}{
		name={volumen},
		text={volumen},
		sort={volumen},
		plural={vol{\'u}menes},
		first={volumen},
		firstplural={vol{\'u}menes},
		description={En el contexto de los \glspl{sibiot} y la \gls{calidaddatos}, el volumen se refiere al tama{\~n}o o magnitud de los datos producidos en un entorno determinado. Este factor representa uno de los principales retos en la administraci{\'o}n de informaci{\'o}n masiva, especialmente cuando se trata de flujos continuos provenientes de \glspl{sensor} y \glspl{dispositivoconectado}. Un alto volumen de datos requiere soluciones de almacenamiento escalables, procesamiento distribuido y t{\'e}cnicas de \gls{analisisbigdata} para garantizar la eficiencia en la \gls{tomadecisionesinformadas}.},
		type=\glsdefaulttype,
		symbol={📦},
		user1={Volume},
		user2={Cantidad total de informaci{\'o}n generada o almacenada},
		user3={Dimensi{\'o}n cuantitativa de los datos en sistemas IoT o Big Data},
		see={velocidad,calidaddatos,almacenamiento,analisisbigdata}
	}
	
	\newglossaryentry{vulnerabilidad}{
		name={Vulnerabilidad},
		text={vulnerabilidad},
		sort={vulnerabilidad},
		plural={vulnerabilidades},
		first={vulnerabilidad},
		firstplural={vulnerabilidades},
		description={Condici{\'o}n inherente a un activo tecnol{\'o}gico que lo expone a posibles ataques o fallos. Puede originarse por errores de dise{\~n}o, implementaci{\'o}n, configuraci{\'o}n deficiente, falta de actualizaciones o debilidades en los controles de seguridad. Una \gls{vulnerabilidad} no siempre implica un da{\~n}o inmediato, pero constituye un riesgo potencial de \gls{vulneracion}. En sistemas \glspl{sibiot}, las vulnerabilidades pueden comprometer sensores, actuadores, protocolos de comunicaci{\'o}n o servicios en la nube, generando riesgos para la \gls{disponibilidad}, \gls{integridad}, \gls{confidencialidad} y \gls{privacidad} de los \glspl{dato}.},
		type=main,
		symbol={🛡️},
		user1={Debilidad de seguridad},
		user2={Riesgo potencial},
		user3={Explotaci{\'o}n posible},
		see={vulneracion,seguridad,riesgo,privacidad}
	}
	
	\newglossaryentry{vulneracion}{
		name={Vulneraci{\'o}n},
		text={vulneraci{\'o}n},
		sort={vulneracion},
		plural={vulneraciones},
		first={vulneraci{\'o}n},
		firstplural={vulneraciones},
		description={Incidente o acci{\'o}n que compromete la protecci{\'o}n de un recurso tecnol{\'o}gico o informativo, derivado de la explotaci{\'o}n de una \gls{vulnerabilidad}, la ausencia de controles adecuados o la acci{\'o}n deliberada de un atacante. Una vulneraci{\'o}n puede implicar acceso no autorizado, modificaci{\'o}n, divulgaci{\'o}n, destrucci{\'o}n o interrupci{\'o}n de la informaci{\'o}n o de los servicios. En el contexto de los \glspl{sibiot}, la vulneraci{\'o}n de la seguridad puede tener consecuencias cr{\'i}ticas al afectar infraestructuras, datos sensibles o procesos de control en \gls{tiemporeal}.},
		type=main,
		symbol={⚠️},
		user1={Incidente de seguridad},
		user2={Quebrantamiento de protecciones},
		user3={Compromiso de sistemas},
		see={vulnerabilidad,seguridad,privacidad,integridad}
	}
	
	\newglossaryentry{wan}{
		name={WAN},
		text={WAN},
		plural={WAN},
		sort={wan},
		first={red de {\'a}rea amplia (\textit{Wide Area Network} , WAN)},
		description={Tipo de red de comunicaci{\'o}n que abarca grandes extensiones geogr{\'a}ficas, como regiones, pa{\'i}ses o continentes. Interconecta m{\'u}ltiples \gls{lan} y \gls{man}, utilizando enlaces de larga distancia provistos por operadores de telecomunicaciones. Se apoya en tecnolog{\'i}as como \gls{mpls}, fibra {\'o}ptica, enlaces satelitales y redes celulares. En sistemas \gls{iot}, permite la comunicaci{\'o}n remota entre nodos distribuidos, servidores en la nube y plataformas de gesti{\'o}n centralizadas.},
		type=main,
		symbol={🌐},
		see={lan,man,red}
	}
	
	\newglossaryentry{wearable}{
		name={Tecnolog{\'i}a vestible},
		text={tecnolog{\'i}a vestible},
		sort={tecnologia vestible},
		plural={dispositivos \textit{wearables}},
		first={tecnolog{\'i}a vestible (\textit{wearable})},
		firstplural={dispositivos \textit{wearables}},
		description={Dispositivo electr{\'o}nico que se lleva en el cuerpo y que registra, analiza y transmite datos biom{\'e}tricos, ambientales o de actividad. En los \glspl{sibiot}, los \textit{wearables} permiten monitorizaci{\'o}n de salud, seguimiento de actividad y asistencia personalizada.},
		type=main,
		symbol={⌚},
		user1={Dispositivo corporal inteligente},
		user2={Captura de datos fisiol{\'o}gicos y ambientales},
		user3={Monitoreo de salud y actividad},
		see={sensor,vidaindependiente,monitorizacion,seguridad}
	}
	
	\newglossaryentry{web}{
		name={Web},
		text={web},
		plural={web},
		sort={web},
		first={web},
		description={Sistema de distribuci{\'o}n de documentos hipermedia interconectados y accesibles v{\'i}a Internet mediante el protocolo \gls{http}. Permite a los usuarios acceder a informaci{\'o}n y servicios a trav{\'e}s de navegadores, integrando texto, im{\'a}genes, v{\'i}deo y aplicaciones interactivas. Es fundamental en arquitecturas cliente-servidor y ampliamente usada como interfaz en plataformas \gls{iot}, sistemas distribuidos y aplicaciones en la nube.},
		type=main,
		symbol={🌍},
		see={http,https,browser}
	}
	
	\newglossaryentry{webhook}{
		name={Webhook},
		text={webhook},
		sort={webhook},
		plural={webhooks},
		first={webhook},
		firstplural={webhooks},
		description={Mecanismo de comunicaci{\'o}n basado en eventos mediante el cual un sistema notifica autom{\'a}ticamente a otro sistema a trav{\'e}s de una llamada HTTP cuando ocurre un evento espec{\'i}fico. A diferencia de los modelos de consulta peri{\'o}dica (\emph{polling}), los \emph{webhooks} permiten una integraci{\'o}n reactiva y eficiente, reduciendo latencia y consumo de recursos. En arquitecturas IoT y \glspl{sibiot}, los webhooks se emplean para notificar eventos relevantes como cambios de inventario, alertas operativas, confirmaciones de transacciones o anomal{\'i}as detectadas, facilitando el desacoplamiento entre sistemas, la integraci{\'o}n con \glspl{erp}, \glspl{scm} y servicios externos, y la construcci{\'o}n de flujos de automatizaci{\'o}n orientados a eventos},
		type=main,
		symbol={🔔},
		see={mensajeriatiemporeal,orientadoeventos,desacoplamiento,integracion,api,restful,sibiot}
	}
	
	\newglossaryentry{webservices}{
		name={Servicio web},
		text={servicio web},
		plural={servicios web},
		sort={webservices},
		first={servicio web},
		firstplural={servicios web},
		description={Tecnolog{\'i}a que utiliza un conjunto de protocolos y est{\'a}ndares para intercambiar datos entre aplicaciones heterog{\'e}neas. Facilita la integraci{\'o}n de sistemas distribuidos y soporta \gls{interoperabilidad} en entornos \gls{iot}.},
		type=main,
		symbol={🔗},
		see={api,rest,restful}
	}
	
	\newglossaryentry{websocket}{
		name={WebSocket},
		text={WebSocket},
		plural={WebSockets},
		sort={websocket},
		first={WebSocket},
		firstplural={WebSockets},
		description={Tecnolog{\'i}a que permite abrir una sesi{\'o}n de comunicaci{\'o}n interactiva entre el navegador y un servidor, enviando y recibiendo datos en ambos sentidos en tiempo real. Gracias a la API WebSockets, la comunicaci{\'o}n es persistente y orientada a eventos, lo que mejora la eficiencia frente a peticiones HTTP repetitivas.},
		type=main,
		symbol={🔌},
		see={http,api,rest}
	}
	
	\newglossaryentry{wss}{
		name={WebSocket Secure},
		text={WSS},
		sort={wss},
		plural={WSS},
		first={WebSocket seguro o WebSocket Secure, WSS)},
		firstplural={WSSs},
		description={Extensi{\'o}n segura del protocolo \gls{websocket} que utiliza cifrado \gls{tls} para proporcionar comunicaci{\'o}n bidireccional persistente con confidencialidad, integridad y autenticaci{\'o}n entre cliente y servidor. \emph{WSS} opera sobre \gls{https} y permite mantener canales de comunicaci{\'o}n en \gls{tiemporeal} protegidos frente a interceptaci{\'o}n, manipulaci{\'o}n de datos y ataques de intermediario. En arquitecturas IoT y \glspl{sibiot}, WSS es ampliamente utilizado para la transmisi{\'o}n segura de telemetr{\'i}a, eventos y comandos de control, especialmente en escenarios que requieren baja latencia, alta frecuencia de mensajes y supervisi{\'o}n continua de dispositivos distribuidos},
		type=main,
		symbol={🔐},
		see={websocket,https,tls,seguridad,integridad,autenticacion,mensajeriatiemporeal,sibiot}
	}
	
	
	\newglossaryentry{wifi}{
		name={Wi-Fi},
		text={Wi-Fi},
		plural={Wi-Fi},
		sort={wifi},
		first={fidelidad inal{\'a}mbrica (\textit{Wireless Fidelity}, Wi-Fi)},
		firstplural={Wi-Fi},
		description={Tecnolog{\'i}a de comunicaci{\'o}n inal{\'a}mbrica que conecta dispositivos a redes de datos mediante ondas de radio. Basada en IEEE 802.11, es esencial para la conectividad local en hogares, empresas e industrias. En \gls{iot}, facilita la comunicaci{\'o}n entre nodos, gateways y la nube, ofreciendo alta velocidad aunque con mayor consumo energ{\'e}tico que otras tecnolog{\'i}as como \gls{ble} o \gls{lora}.},
		type=main,
		symbol={📶},
		see={ble,lora,zigbee}
	}
	
	\newglossaryentry{wifihalow}{
		name={Wi-Fi HaLow},
		text={Wi-Fi HaLow},
		sort={wifi halow},
		plural={versiones Wi-Fi HaLow},
		first={Wi-Fi HaLow (IEEE~802.11ah)},
		firstplural={versiones Wi-Fi HaLow},
		description={Variante 802.11ah para \gls{iot}, con bajo consumo, gran alcance y alta densidad de conexiones, adecuada para sensores interiores.},
		type=main,
		symbol={🪪},
		see={wifi,redinalambrica,capaenlace}
	}
	
	\newglossaryentry{workbench}{
		name={MySQL Workbench},
		text={MySQL Workbench},
		plural={MySQL Workbench},
		first={MySQL Workbench},
		firstplural={MySQL Workbench},
		description={Entorno visual que permite modelar esquemas, ejecutar consultas \gls{sql} y administrar servidores \gls{mysql} desde una interfaz integrada.},
		sort={mysql workbench},
		type=main,
		symbol={🛠},
		see={mysql,basedatos}
	}
	
	\newglossaryentry{workflow}{
		name={Workflow},
		text={workflow},
		sort={workflow},
		plural={workflows},
		first={workflow o flujo de trabajo},
		firstplural={workflows o flujos de trabajo},
		description={Secuencia definida de tareas, actividades o procesos que se ejecutan de forma coordinada para alcanzar un objetivo espec{\'i}fico. Un \emph{workflow} puede ser manual, autom{\'a}tico o mixto, e involucra la asignaci{\'o}n de roles, recursos, reglas y condiciones de ejecuci{\'o}n. En ingenier{\'i}a de software y en los \glspl{sibiot}, los workflows se utilizan para orquestar la interacci{\'o}n entre componentes de hardware, software y servicios distribuidos, optimizando la eficiencia operativa, la trazabilidad de las acciones y la automatizaci{\'o}n de procesos.},
		type=main,
		symbol={🔄},
		user1={Flujo de trabajo},
		user2={Proceso coordinado},
		user3={Orquestaci{\'o}n de tareas},
		see={proceso,automatizacion,gestionempresarial}
	}
	
	\newglossaryentry{wsdl}
	{
		name={{WSDL}},
		text={{WSDL}},
		sort={wsdl},
		plural={{WSDL}},
		first={lenguaje de descripci{\'o}n de servicios web (\textit{Web Services Description Language}, WSDL)},
		firstplural={{WSDL}},
		description={El \emph{Web Services Description Language} (\gls{wsdl}) es un lenguaje basado en \gls{xml} para describir servicios de \glspl{red} como colecciones de \emph{endpoints}, operaciones y mensajes, con el fin de publicar un contrato t{\'e}cnico consumible por herramientas y clientes.},
		type=\glsdefaulttype,
		symbol={🧾},
		user1={Contrato de servicio},
		user2={Descripci{\'o}n de interfaz},
		user3={Servicios web},
		see={soa,soap,api}
	}
	
	\newglossaryentry{wsn}{
		name={WSN},
		text={WSN},
		plural={WSNs},
		sort={wsn},
		first={red de sensores inal{\'a}mbricos (\textit{Wireless Sensor Network}, WSN)},
		firstplural={WSNs},
		description={Conjunto distribuido de nodos sensores conectados de manera inal{\'a}mbrica para monitorear condiciones ambientales como temperatura, humedad, presi{\'o}n o movimiento. Cada nodo combina sensores, procesamiento, comunicaci{\'o}n y energ{\'i}a. Usadas en agricultura inteligente, salud, ciudades inteligentes y sistemas \gls{iot}, permiten recolectar y transmitir datos en \gls{tiemporeal}.},
		type=main,
		symbol={🌱},
		see={sensor,iot,redinalambrica}
	}
	
	\newglossaryentry{www}{
		name={WWW},
		text={WWW},
		plural={WWW},
		sort={www},
		first={red inform{\'a}tica mundial (\textit{World Wide Web}, WWW)},
		firstplural={WWW},
		description={Acr{\'o}nimo com{\'u}n para \gls{web}. Ver \gls{web} para la descripci{\'o}n detallada.},
		type=main,
		symbol={🌎},
		see={web,http}
	}
	
	\newglossaryentry{xhtml}{
		name={XHTML},
		text={XHTML},
		plural={XHTML},
		sort={xhtml},
		first={lenguaje de marcado de hipertexto extensible (\textit{Extensible HyperText Markup Language}, XHTML)},
		firstplural={XHTML},
		description={Lenguaje de marcado que combina la estructura de \gls{html} con la rigurosidad sint{\'a}ctica de \gls{xml}. Desarrollado como reformulaci{\'o}n de HTML 4.01 con reglas XML 1.0, permite validaci{\'o}n estricta e interoperabilidad. Aunque hoy ha sido reemplazado en gran parte por \gls{html5}, sigue us{\'a}ndose en entornos que requieren alta compatibilidad y consistencia.},
		type=main,
		symbol={📄},
		see={html,html5,xml}
	}
	
	\newglossaryentry{xml}{
		name={XML},
		text={XML},
		plural={XML},
		sort={xml},
		first={lenguaje de marcado extensible (\textit{Extensible Markup Language}, XML)},
		firstplural={XML},
		description={Metalenguaje desarrollado por el W3C que permite definir lenguajes de marcas. Se usa para almacenar e intercambiar datos en forma legible y estructurada, independiente del lenguaje de programaci{\'o}n empleado. XML es ampliamente usado en integraci{\'o}n de sistemas, configuraciones y servicios web.},
		type=main,
		symbol={🗂️},
		see={json,xhtml}
	}
	
	\newglossaryentry{xp}{
		name={XP},
		text={XP},
		sort={xp},
		first={Programaci{\'o}n Extrema (\textit{eXtrem Programming}, XP)},
		firstplural={Programaciones Extremas (XP)},
		description={Metodolog{\'i}a {\'a}gil de desarrollo de software que enfatiza la mejora continua, la retroalimentaci{\'o}n r{\'a}pida, el desarrollo incremental, las pruebas constantes y la colaboraci{\'o}n estrecha con el cliente. XP se centra en valores como simplicidad, comunicaci{\'o}n, retroalimentaci{\'o}n y valent{\'i}a, promoviendo pr{\'a}cticas como programaci{\'o}n en pareja, integraci{\'o}n continua, desarrollo guiado por pruebas y entregas frecuentes. Es especialmente relevante en proyectos que requieren alta adaptabilidad y calidad de software.},
		type=main,
		symbol={⚙️},
		see={scrum,agile,tdd},
		user1={Metodolog{\'i}a {\'a}gil},
		user2={Desarrollo de software},
		user3={Buenas pr{\'a}cticas de ingenier{\'i}a}
	}
	
	\newglossaryentry{zigbee}{
		name={Zigbee},
		text={Zigbee},
		sort={zigbee},
		plural={Zigbee},
		first={Zigbee},
		firstplural={Zigbee},
		description={Est{\'a}ndar de comunicaci{\'o}n inal{\'a}mbrica basado en el protocolo IEEE 802.15.4, dise{\~n}ado para redes de {\'a}rea personal inal{\'a}mbrica (WPAN) de baja potencia y bajo consumo de energ{\'i}a. \emph{Zigbee} se utiliza en aplicaciones de automatizaci{\'o}n del hogar, entornos industriales y sistemas de control remoto. Ofrece una comunicaci{\'o}n segura, soporta redes de tipo malla y opera en distancias cortas (hasta unos 100 metros). Se destaca por su eficiencia energ{\'e}tica, capacidad para formar redes grandes, baja latencia y simplicidad en su implementaci{\'o}n. En el contexto de los \glspl{sibiot}, Zigbee es ampliamente empleado en \glspl{sensor} y \glspl{actuador} para tareas de \gls{domotica}, monitorizaci{\'o}n y control de procesos, integr{\'a}ndose con otros protocolos de \gls{comunicacioninalambrica} y plataformas en la nube.},
		type=main,
		symbol={📡},
		user1={Protocolo de comunicaci{\'o}n inal{\'a}mbrica de baja potencia},
		user2={Soporte para redes malladas},
		user3={Uso en dom{\'o}tica, industria y SIbIoT},
		see={comunicacioninalambrica,domotica,sensor,actuador,sibiot}
	}
	
	\newglossaryentry{zwave}{
		name={Z-Wave},
		text={Z-Wave},
		plural={Z-Wave},
		sort={zwave},
		first={Z-Wave},
		firstplural={Z-Wave},
		description={\Gls{protocolocomunicacion} de comunicaci{\'o}n inal{\'a}mbrica de baja potencia y corto alcance dise{\~n}ado para la \gls{automatizacion} del hogar y \glspl{aplicacion} \glspl{domotica}. Opera en bandas sub-GHz y utiliza topolog{\'i}a en malla para asegurar fiabilidad de transmisi{\'o}n. En \gls{iot}, se usa para controlar luces, cerraduras, sensores, termostatos y \glspl{dispositivointeligente}, destacando por su \gls{interoperabilidad}, \gls{seguridad} y bajo consumo energ{\'e}tico.},
		type=main,
		symbol={🏠},
		see={zigbee,bluetooth,wifi}
	}